\documentclass[a4paper,UTF8,fontset=windows,AutoFakeBold]{ctexart}
\pagestyle{headings}
\title{\textbf{高等数学A\ 习题课讲义}}
\author{原生生物}
\date{}
\setcounter{tocdepth}{3}
\usepackage{amsmath,amssymb,amsthm,enumerate,framed,geometry,hyperref,paralist,pgfplots,ulem,unicode-math}

\geometry{left = 2.0cm, right = 2.0cm, top = 2.0cm, bottom = 2.0cm}
\ctexset{section={number=\zhnum{section}}}
\ctexset{subsection={name={\S},number=\arabic{section}.\arabic{subsection}}}

\newcommand*{\dr}{\hspace{0.07em}\mathrm{d}}
\newcommand*{\cm}{\mathcal{C}}
\newcommand*{\er}{\mathrm{e}}
\newcommand*{\Fv}{\vec{F}}
\newcommand*{\ir}{\mathrm{i}}
\newcommand*{\nm}{\mathcal{N}}
\newcommand*{\nv}{\vec{n}}
\newcommand*{\rb}{\mathbb{R}}
\newcommand*{\rv}{\vec{r}}
\newcommand*{\rot}{\mathrm{rot}}
\newcommand*{\suminf}[1][n=1]{\sum_{#1}^\infty}

\newcommand*{\bm}{\boldsymbol}

\newcommand*{\extra}{\textbf{附加}}

\newcommand*{\note}{\noindent *}

\newtheoremstyle{ex}% name
{3pt}% Space above
{3pt}% Space below
{\songti}% Body font
{}% Indent amount
{\bfseries}% Theorem head font
{.}% Punctuation after theorem head
{.5em}% Space after theorem head
{}% Theorem head spec (can be left empty, meaning ‘normal’)

\theoremstyle{ex}
\newtheorem{exercise}{题}
\newcommand{\exref}[1]{\textbf{题\;\ref{#1}\;}}

\newcommand{\sol}[1]{{\vspace{5pt}\kaishu\noindent\textbf{解答}：\vspace{-3pt}
\begin{compactitem}
    \item[] #1
\end{compactitem}
}}

\pgfplotsset{compat=1.18}

\setcounter{section}{17}

\begin{document}

\maketitle

\note 高等数学A [周蜀林老师班]习题课讲义。

\note 由于高等数学是一个需要大量练习的学科，因此习题课的主要组织结构将会是通过\textbf{题目}进行，根据每道题的编号，大家可以看到一学期的知识究竟需要多少个练习来掌握。题目的解答将以{\kaishu\textbf{楷书}}展示，对一些相对复杂的定理，如果时间不够可仅了解结论，跳过证明细节。

\note 讲义中提到的$n$默认为正整数，$n$趋于无穷时的极限均指\textbf{数列极限}。$C_m^n$代表组合数$\frac{m!}{n!(m-n)!}$，且当$m<n$时定义为0。

\note 讲义中的$\partial_1f$完全等同于上册教材中的$f'_1$，代表函数$f$对第一个分量的偏导数，$\partial_2f$、$\partial_3f$等同理；$\partial_1^2f$代表$f''_{11}$，也即对第一个分量作两次偏导数，其他同理。$\nabla u$表示多元函数$u$的梯度向量，$\nabla\cdot\Fv$表示多元映射$\Fv$的散度，$\nabla\times\Fv$表示多元映射$\Fv$的旋度，$\triangle u=\nabla\cdot\nabla u$为$u$的各二阶偏导求和。

\note 讲义中$L^+$代表平面简单闭曲线$L$沿\textbf{逆时针方向}，$S^+$代表封闭曲面$S$的\textbf{外侧}。

\tableofcontents

\newpage
\section{得失}
\subsection{期末考试}
本次习题课主要介绍了关于高等数学与学习状态相关的建议，不存在需要掌握的知识性内容，但仍然\textbf{非常推荐大家阅读}。

\subsubsection{试题}
\begin{enumerate}
    \item 设三维空间中四点$ABCD$满足
        $$\overrightarrow{AB}=(2,a,0),\quad\overrightarrow{AC}=(0,-1,2),\quad\overrightarrow{AD}=(2,b,1)$$
        其中$a>0$、$b<\frac{1}{2}$。已知$\overrightarrow{AB}$、$\overrightarrow{AC}$作为相邻两边形成的平行四边形面积为$2\sqrt6$，$\overrightarrow{AB}$、$\overrightarrow{AC}$、$\overrightarrow{AD}$作为共顶点三棱形成的平行六面体体积为2。
    \begin{enumerate}
        \item 求$a$、$b$；
        \item 求平行于向量$\overrightarrow{AB}$、$\overrightarrow{AD}$且过点$(2,3,4)$的平面方程；
        \item 求函数$f(x,y,z)=x^2+y^2+z^2+xyz$在点$(1,2,1)$的梯度以及在该点沿$\overrightarrow{AB}$方向的方向导数。
    \end{enumerate}
    
    \item 设函数$f(x)$在闭区间$[2,4]$上连续，在开区间$(2,4)$可导且导数大于0，极限$\lim_{x\to2}\frac{f(2x-2)}{x-2}$存在，证明：
    \begin{enumerate}
        \item $f(x)$在$(2,4)$上取值恒大于0；
        \item 存在$\xi\in(2,4)$使得
        $$\frac{6f(\xi)}{\xi}=\int_2^4f(x)\dr x$$
        \item 对上一问中的$\xi$，存在$\eta\in(2,4)$、$\eta\ne\xi$使得
        $$6f'(\eta)=\frac{\xi}{\xi-2}\int_2^4f(x)\dr x$$
    \end{enumerate}

    \item 解答下列各题：
    \begin{enumerate}
        \item 写出$f(x)=\ln(1+x)$在$x=0$处的$n$阶带Peano余项的泰勒公式；
        \item 计算极限
        $$I=\lim_{x\to0}\frac{1}{x^2}\big(2(1+x)^{\frac{1}{x}}+\er(x-2)\big)$$
    \end{enumerate}

    \item 设二元函数$f(x,y)$有连续偏导数，且
    $$f(1,2)=0,\quad\frac{\partial f(x,y)}{\partial x}\bigg|_{(1,2)}=5$$
    已知对任意实数$t$、$x$、$y$有$f(tx,ty)=t^2f(x,y)$：
    \begin{enumerate}
        \item 计算$\frac{\partial f(x,y)}{\partial y}\big|_{(1,2)}$；
        \item 计算
        $$\lim_{x\to0}\frac{1}{x}\int_0^x\big(1+f(2t-2\sin t+1,\sqrt[3]{1+t^3}+1)\big)^{\frac{1}{\ln(1+t^3)}}\dr t$$
    \end{enumerate}

    \item 设区域$D=(0,+\infty)\times(-\infty,+\infty)$：
    \begin{enumerate}
        \item 若定义在$D$上的二元函数$u=h(x,y)$具有连续的二阶偏导数且满足
        $$\forall(x,y)\in D,\quad\frac{\partial^2u}{\partial x\partial y}=0$$
        求证存在$(0,+\infty)$上二阶连续可导的函数$f$与$\rb$上二阶连续可导的函数$g$使得$u(x,y)=f(x)+g(y)$对任何$(x,y)\in D$成立；
        \item 若定义在$D$上的二元正值函数$u=h(x,y)$具有连续的二阶偏导数且满足
        $$\forall(x,y)\in D,\quad u\frac{\partial^2u}{\partial x\partial y}=\frac{\partial u}{\partial x}\frac{\partial u}{\partial y}$$
        验证$z=\ln h(x,y)$满足
        $$\forall(x,y)\in D,\quad\frac{\partial^2z}{\partial x\partial y}=0$$
        \item 若定义在$D$上的二元正值函数$u=h(x,y)$具有连续的二阶偏导数且满足
        $$\begin{cases}\forall(x,y)\in D,\quad u\dfrac{\partial^2u}{\partial x\partial y}=\dfrac{\partial u}{\partial x}\dfrac{\partial u}{\partial y}\\\forall x>0,\quad h(x,0)=x\er^{-\frac{x^2}{2}}\\\forall y\in\rb,\quad h(1,y)=\er^{-\frac{1+y^2}{2}}\end{cases}$$
        求函数$h(x,y)$的表达式。
    \end{enumerate}

    \item 若函数$f(x,y)$在$\rb^2$上有二阶连续偏导数，且
    $$f(x,y)=1-x-y+o(\sqrt{(x-1)^2+y^2})\quad (x,y)\to(1,0)$$
    设
    $$g(x,y)=f(\er^{xy},x^2+y^2)$$
    证明$g$在$(0,0)$处取极值，计算极值并判断其为极大值或极小值。

    \item 设定义域为$\rb^2$，值域为区间$I$的二元函数$u=f(x,y)$有二阶连续偏导数且对$y$偏导处处非零，证明以下两个命题等价：
    \begin{itemize}
        \item 对任意实数$C\in I$，曲线$f(x,y)=C$是$\rb^2$上的一条直线；
        \item 对任何$(x,y)\in\rb^2$有
        $$\bigg(\frac{\partial u}{\partial y}\bigg)^2\frac{\partial^2u}{\partial x^2}-2\frac{\partial u}{\partial x}\frac{\partial u}{\partial y}\frac{\partial^2u}{\partial x\partial y}+\bigg(\frac{\partial u}{\partial x}\bigg)^2\frac{\partial^2u}{\partial y^2}=0$$
    \end{itemize}
\end{enumerate}

\subsubsection{解答}
\begin{enumerate}
    \item
    \begin{enumerate}
        \item 
        利用向量\textbf{外积的几何意义}可知平行四边形面积为$2\sqrt6$等价于
        $$|(2,a,0)\times(0,-1,2)|=2\sqrt6$$
        直接计算可知外积结果为$(2a,-4,-2)$，从而得到$a=\pm1$，由范围可知$a=1$。

        利用\textbf{混合积几何意义}(或直接通过行列式)可知平行六面体体积为2等价于
        $$|(2,-4,-2)\cdot(2,b,1)|=2$$
        从而得到$b=0$或1，由范围可知$b=0$。

        \item 由于平面平行于$(2,1,0)$与$(2,0,1)$，计算外积可得到平面的一个法向量$(1,-2,-2)$，由此得到点法式平面方程
        $$x-2-2(y-3)-2(z-4)=0$$
        化简得
        $$x-2y-2z+12=0$$

        \item
        直接计算可知
        $$\frac{\partial f(x,y,z)}{\partial x}=2x+yz,\quad\frac{\partial f(x,y,z)}{\partial y}=2y+xz,\quad\frac{\partial f(x,y,z)}{\partial z}=2z+xy$$
        从而由定义可知梯度为$(4,5,4)$。
        
        将方向向量单位化得到$\frac{1}{\sqrt5}(2,1,0)$，从而直接利用方向导数公式可知方向导数结果为
        $$\frac{1}{\sqrt5}(2,1,0)\cdot(4,5,4)=\frac{13\sqrt5}{5}$$
    \end{enumerate}

    \item
    \begin{enumerate}
        \item 
        由开区间导数大于0可知$f(x)$在$[2,4]$严格单调增。另一方面，若$f(2)\ne0$，换元$t=2x-2$有
        $$\lim_{x\to2}\frac{f(2x-2)}{x-2}=\lim_{t\to2}\frac{f(t)}{t-2}$$
        由于$t\to2$时分母极限为0，分子极限非零，极限应不存在，与条件矛盾，从而必然$f(2)=0$。

        结合$f(2)=0$与$f$在$[2,4]$严格单调增即得到$x\in(2,4)$时$f(x)>f(2)=0$，得证。

        \item 由于出现了商中的$\xi$，想到利用\textbf{柯西中值定理}进行处理。为使取导函数中值后出现$\frac{f(\xi)}{\xi}$的形式，构造
        $$F(x)=\int_2^xf(t)\dr t,\quad g(x)=\frac{1}{2}x^2$$
        则$g'(x)=x$在$[2,4]$恒非零，且由变上限积分的导数结论，$F'(x)=f(x)$在$(2,4)$成立。

        对$F(x)$、$g(x)$应用柯西中值定理可得存在$\xi\in(2,4)$使得
        $$\frac{F'(\xi)}{g'(\xi)}=\frac{F(4)-F(2)}{g(4)-g(2)}$$
        代入计算可知
        $$\frac{f(\xi)}{\xi}=\frac{\int_2^4f(x)\dr x}{6}$$
        变形即得要证的结果。

        \item
        将(b)中$\xi$满足的条件改写为
        $$\xi\int_2^4f(x)\dr x=6f(\xi)$$
        代入即可将要证的式子化为存在$\eta\in(2,4)$、$\eta\ne\xi$使得
        $$f'(\eta)=\frac{f(\xi)}{\xi-2}$$
        由(a)中已经证明的$f(2)=0$，右侧即为$\frac{f(\xi)-f(2)}{\xi-2}$，直接通过\textbf{拉格朗日中值定理}可知存在$\eta\in(2,\xi)$符合要求。

        \note 本题的(c)并不依赖(b)的过程，且实际上比(b)更加简单。因此，考试时前面的小问没做出来仍然可以尝试后面的小问。
    \end{enumerate}

    \item
    \begin{enumerate}
        \item 
        根据教材可直接写出
        $$\ln(1+x)=\sum_{k=1}^n(-1)^{k-1}\frac{x^k}{k}+o(x^n)$$

        \item
        将$x$的指数写成对数可知原极限为
        $$\lim_{x\to0}\frac{2\er^{\frac{\ln(1+x)}{x}}+\er(x-2)}{x^2}$$
        由于分母为二阶，我们想要将分子\textbf{也展开到二阶}。由于$\frac{\ln(1+x)}{x}$在0处极限为1，我们记$t=\frac{\ln(1+x)}{x}$，先写出$\er^t$在1处的二阶泰勒展开
        $$\er^t=\er+\er(t-1)+\frac{\er}{2}(t-1)^2+o((t-1)^2)\quad(t\to1)$$
        利用(a)中的展开，我们又可以写出$\frac{\ln(1+x)}{x}$在0处的二阶泰勒展开
        $$\frac{\ln(1+x)}{x}=1-\frac{x}{2}+\frac{x^2}{3}+o(x^2)\quad(x\to0)$$
        将后一个泰勒展开式\textbf{代入}前一个泰勒展开式，并将超过二次的项合并到$o(x^2)$中，可得到(详细过程可参考上学期讲义11.3.2)，最终得到展开结果
        $$\er^{\frac{\ln(1+x)}{x}}=\er+\er\bigg(-\frac{x}{2}+\frac{x^2}{3}\bigg)+\frac{\er}{2}\bigg(-\frac{x}{2}\bigg)^2+o(x^2)=\er-\frac{\er}{2}x+\frac{11\er}{24}x^2+o(x^2)$$
        代入到分子即得原极限为(由$\frac{o(x^2)}{x^2}$在0处极限为0)
        $$\lim_{x\to0}\frac{1}{x^2}\bigg(2\er-\er x+\frac{11\er}{12}x^2+o(x^2)+\er(x-2)\bigg)=\frac{11}{12}\er$$
    \end{enumerate}

    \item
    \begin{enumerate}
        \item 
        将条件两侧的$x$、$y$、$t$\textbf{匹配}(这是为了保证表达式偏导的严谨性，详见上学期讲义14.4.3)，对$t$求偏导，由链式法则可知
        $$\frac{\partial(tx)}{\partial t}f_1'(tx,ty)+\frac{\partial(ty)}{\partial t}f_2'(tx,ty)=2tf(x,y)$$
        也即
        $$xf_1'(tx,ty)+yf_2'(tx,ty)=2tf(x,y)$$
        由于我们已知$(1,2)$处的函数值与对$x$偏导值，需要知道$(1,2)$处的对$y$偏导值，代入$t=1$、$x=1$、$y=2$可知
        $$f_1'(1,2)+2f_2'(1,2)=2f(1,2)$$
        利用表达式偏导的定义
        $$f_1'(1,2)=\frac{\partial f(x,y)}{\partial x}\bigg|_{(1,2)},\quad f_2'(1,2)=\frac{\partial f(x,y)}{\partial y}\bigg|_{(1,2)}$$
        代入即解得
        $$\frac{\partial f(x,y)}{\partial y}\bigg|_{(1,2)}=-\frac{5}{2}$$

        \item
        利用初等函数与复合函数的连续性，验证条件并使用\textbf{洛必达法则}(利用变上限积分的导数结论)可得只需计算
        $$\lim_{x\to0}\big(1+f(2x-2\sin x+1,\sqrt[3]{1+x^3}+1)\big)^{\frac{1}{\ln(1+x^3)}}$$
        由指数函数连续性取$\ln$得只需计算
        $$\lim_{x\to0}\frac{\ln\big(1+f(2x-2\sin x+1,\sqrt[3]{1+x^3}+1)\big)}{\ln(1+x^3)}$$
        由于$t\to0$时$\ln(1+t)$与$t$等价，而$x\to0$时$x^3$与$f(2x-2\sin x+1,\sqrt[3]{1+x^3}+1)$极限均为0，从而可以进行推广的等价无穷小替换得到原极限为
        $$\lim_{x\to0}\frac{f(2x-2\sin x+1,\sqrt[3]{1+x^3}+1)}{x^3}$$
        \note 严格意义上这步过程需要保证$f(2x-2\sin x+1,\sqrt[3]{1+x^3}+1)$在0的某个邻域恒非零，这可以由两偏导均非零证明，或根据上学期讲义15.2.2的讨论跳过这部分证明。

        接下来我们分为四部分说明：
        \begin{itemize}
            \item 先分析$2x-2\sin x+1$，直接利用$\sin x$的泰勒展开可知
            $$2x-2\sin x+1=1+\frac{1}{3}x^3+o(x^3)$$
            \item 再分析$\sqrt[3]{1+x^3}+1$，利用推广的等价无穷小替换可知$\sqrt[3]{1+x^3}-1$与$\frac{1}{3}x^3$等价，从而进一步有(也可直接泰勒展开得到)
            $$\sqrt[3]{1+x^3}+1=2+\frac{1}{3}x^3+o(x^3)$$
            \item 利用$f$有连续偏导数，其可微，从而在$(1,2)$处有
            $$f(1+t,2+s)=f(1,2)+f_1'(1,2)t+f_2'(1,2)s+o(\sqrt{s^2+t^2})\quad(t,s)\to(0,0)$$
            利用(a)中的结论可得
            $$f(1+t,2+s)=5t-\frac{5}{2}s+o(\sqrt{s^2+t^2})\quad(t,s)\to(0,0)$$
            代入前两部分即得到$f(2x-2\sin x+1,\sqrt[3]{1+x^3}+1)$为
            $$5\cdot\frac{1}{3}x^3-\frac{5}{2}\cdot\frac{1}{3}x^3+\sqrt{o(x^3)^2+o(x^3)^2}=\frac{5}{6}x^3+\sqrt{o(x^3)^2+o(x^3)^2}$$
            \item 由定义与连续性可验证
            $$\lim_{x\to0}\frac{\sqrt{o(x^3)^2+o(x^3)^2}}{|x^3|}=\lim_{x\to0}\sqrt{\bigg(\frac{o(x^3)}{x^3}\bigg)^2+\bigg(\frac{o(x^3)}{x^3}\bigg)^2}=0$$
            从而可得$\sqrt{o(x^3)^2+o(x^3)^2}=o(x^3)$，进一步得到
            $$f(2x-2\sin x+1,\sqrt[3]{1+x^3}+1)=\frac{5}{6}x^3+o(x^3)$$
        \end{itemize}
        综合以上，最终得到
        $$\lim_{x\to0}\frac{f(2x-2\sin x+1,\sqrt[3]{1+x^3}+1)}{x^3}=\frac{5}{6}$$
        于是原极限为
        $$\er^{\frac{5}{6}}$$
    \end{enumerate}

    \item
    \begin{enumerate}
        \item 
        由于$\frac{\partial u}{\partial y}=h_2'(x,y)$对$x$的偏导为0，固定任何$y_0$，必然存在唯一的$C$使得
        $$h_2'(x,y_0)=C,\quad\forall x\in(0,+\infty)$$
        由于$C$与$y_0$有关，存在函数$C(y_0)$使得
        $$\forall y_0\in\rb,\quad h_2'(x,y_0)=C(y_0),\quad\forall x\in(0,+\infty)$$
        由条件可知$h_2'(x,y_0)$对$y_0$连续，从而$C(y_0)$也对$y_0$连续，两侧对$y_0$从0到$y$积分，利用微积分基本定理可知
        $$h(x,y)-h(x,0)=\int_0^yC(y_0)\dr y_0$$
        从而记
        $$f(x)=h(x,0),\quad g(y)=\int_0^yC(y_0)\dr y_0$$
        即得到$D$中恒有$h(x,y)=f(x)+g(y)$。

        由于(注意之前推导可发现$u$对$y$的一阶偏导数与$x$无关，因此二阶偏导数也与$x$无关，在任何$x$处相同)
        $$f''(x)=\frac{\partial^2u}{\partial x^2}\bigg|_{(x,0)},\quad g''(y)=C'(y)=\frac{\partial h_2'(x,y)}{\partial y}=\frac{\partial^2u}{\partial y^2}\bigg|_{(x,y)}$$
        利用$h$具有连续二阶偏导数可知$f$、$g$均二阶连续可导。从而得证。。

        \item
        由于$u>0$对任何$x$、$y$成立，直接计算可知
        $$\frac{\partial \ln u}{\partial x}=\frac{1}{u}\frac{\partial u}{\partial x}$$
        进一步得到
        $$\frac{\partial^2z}{\partial x\partial y}=\frac{\partial}{\partial y}\bigg(\frac{1}{u}\frac{\partial u}{\partial x}\bigg)=-\frac{1}{u^2}\frac{\partial u}{\partial y}\frac{\partial u}{\partial x}+\frac{1}{u}\frac{\partial^2u}{\partial x\partial y}$$
        右侧化简即为
        $$\frac{1}{u^2}\bigg(u\frac{\partial^2u}{\partial x\partial y}-\frac{\partial u}{\partial y}\frac{\partial u}{\partial x}\bigg)=0$$
        从而得证。

        \item
        记$z=\ln h(x,y)$，由(b)与(a)可知存在$(0,+\infty)$上二阶连续可导的函数$f$与$\rb$上二阶连续可导的函数$g$使得
        $$\ln h(x,y)=f(x)+g(y)$$
        代入两个条件取$\ln$可发现
        $$\forall x>0,\quad\ln h(x,0)=\ln x-\frac{x^2}{2}$$
        $$\forall y\in\rb,\quad\ln h(1,y)=-\frac{1+y^2}{2}$$
        从而在$D$中有
        $$f(x)+g(0)=\ln x-\frac{x^2}{2},\quad f(1)+g(y)=-\frac{1+y^2}{2}$$
        在第一式中代入$x=1$可发现$f(1)+g(0)=-\frac{1}{2}$，求和可得
        $$f(x)+g(y)=\ln x-\frac{x^2}{2}-\frac{1+y^2}{2}-f(1)-g(0)=\ln x-\frac{x^2}{2}-\frac{1+y^2}{2}+\frac{1}{2}$$
        由于这就是$\ln h(x,y)$，化简即知
        $$h(x,y)=x\er^{-\frac{x^2+y^2}{2}}$$
        可验证其符合条件，从而这就是唯一满足要求的$h(x,y)$。
    \end{enumerate}

    \item
    分为以下三步：
    \begin{itemize}
        \item $f$\textbf{相关条件化简}
        
        由$f$有二阶连续偏导数，其在$(1,0)$可微，将条件化为
        $$f(x,y)=-(x-1)-y+o(\sqrt{(x-1)^2+y^2})\quad(x,y)\to(1,0)$$
        根据可微函数性质可知
        $$f(1,0)=0,\quad f_1'(1,0)=-1,\quad f_2'(1,0)=-1$$

        \item \textbf{链式法则计算}
        
        代入可得
        $$g(0,0)=f(1,0)=0$$
        由于$f$有二阶连续偏导数，直接通过链式法则代入计算可知
        $$g_1'(x,y)=y\er^{xy}f_1'(\er^{xy},x^2+y^2)+2xf_2'(\er^{xy},x^2+y^2)$$
        $$g_2'(x,y)=x\er^{xy}f_1'(\er^{xy},x^2+y^2)+2yf_2'(\er^{xy},x^2+y^2)$$
        代入可得
        $$g_1'(0,0)=g_2'(0,0)=0$$
        进一步得到
        $$\begin{aligned}g_{11}'(x,y)&=y\er^{xy}\frac{\partial f_1'(\er^{xy},x^2+y^2)}{\partial x}+\frac{\partial(y\er^{xy})}{\partial x}f_1'(\er^{xy},x^2+y^2)\\ &+2x\frac{\partial f_2'(\er^{xy},x^2+y^2)}{\partial x}+\frac{\partial(2x)}{\partial x}f_2'(\er^{xy},x^2+y^2)\end{aligned}$$
        由于$(0,0)$处$y\er^{xy}$与$2x$为0，第一、三两项必然为0，从而只剩二、四两项，可算出
        $$g_{11}'(0,0)=\frac{\partial(y\er^{xy})}{\partial x}\bigg|_{(0,0)}f_1'(1,0)+\frac{\partial(2x)}{\partial x}\bigg|_{(0,0)}f_2'(1,0)=2f_2'(1,0)=-2$$
        同理可知
        $$g_{12}'(0,0)=\frac{\partial(y\er^{xy})}{\partial y}\bigg|_{(0,0)}f_1'(1,0)+\frac{\partial(2x)}{\partial y}\bigg|_{(0,0)}f_2'(1,0)=f_1'(1,0)=-1$$
        $$g_{21}'(0,0)=\frac{\partial(x\er^{xy})}{\partial x}\bigg|_{(0,0)}f_1'(1,0)+\frac{\partial(2y)}{\partial x}\bigg|_{(0,0)}f_2'(1,0)=f_1'(1,0)=-1$$
        $$g_{22}'(0,0)=\frac{\partial(x\er^{xy})}{\partial y}\bigg|_{(0,0)}f_1'(1,0)+\frac{\partial(2y)}{\partial y}\bigg|_{(0,0)}f_2'(1,0)=2f_2'(1,0)=-2$$

        \item \textbf{结论整理}
        
        根据以上计算，可得$g$在$(0,0)$处各一阶偏导为0，且在$(0,0)$处$A=-2<0$、$B^2-AC=-3<0$，从而根据教材结论可知$(0,0)$为极大值点，极大值为$g(0,0)=0$。
    \end{itemize}

    \item 
    从本题条件中可以想到与隐函数定理有关。由于$\frac{\partial f(x,y)}{\partial y}\ne0$，任何$f(x,y)=C$都确定了$y$对$x$的隐函数$y=\varphi_C(x)$\ (这里下标$C$是由于函数$\varphi_C$与$C$相关)。

    由此，对任何$C$，$f(x,y)=C$都是直线当且仅当$\varphi_C''(x)=0$对任何$C$成立。直接利用隐函数定理计算可知(由于$f(x,y)-C$的导数与$f(x,y)$相同)
    $$\varphi_C'(x)=-\frac{\frac{\partial f(x,y)}{\partial x}}{\frac{\partial f(x,y)}{\partial y}}=-\frac{\frac{\partial u}{\partial x}}{\frac{\partial u}{\partial y}}$$
    进一步计算可知
    $$\varphi_C''(x)=-\frac{\dr}{\dr x}\frac{\frac{\partial u}{\partial x}}{\frac{\partial u}{\partial y}}=-\frac{1}{\big(\frac{\partial u}{\partial y}\big)^2}\bigg(\frac{\dr}{\dr x}\frac{\partial u}{\partial x}\cdot\frac{\partial u}{\partial y}-\frac{\dr}{\dr x}\frac{\partial u}{\partial y}\cdot\frac{\partial u}{\partial x}\bigg)$$
    利用链式法则可得(由$f$有二阶连续偏导数，其混合偏导数可交换)
    $$\frac{\dr}{\dr x}\frac{\partial u}{\partial x}=\frac{\partial^2u}{\partial x^2}+\frac{\partial^2u}{\partial x\partial y}\frac{\dr y}{\dr x}=\frac{\partial^2u}{\partial x^2}+\frac{\partial^2u}{\partial x\partial y}\varphi_C'(x)=\frac{1}{\frac{\partial u}{\partial y}}\bigg(\frac{\partial u}{\partial y}\frac{\partial^2u}{\partial x^2}-\frac{\partial^2u}{\partial x\partial y}\frac{\partial u}{\partial x}\bigg)$$
    同理
    $$\frac{\dr}{\dr x}\frac{\partial u}{\partial y}=\frac{\partial^2u}{\partial x\partial y}+\frac{\partial^2u}{\partial y^2}\frac{\dr y}{\dr x}=\frac{\partial^2u}{\partial x\partial y}+\frac{\partial^2u}{\partial y^2}\varphi_C'(x)=\frac{1}{\frac{\partial u}{\partial y}}\bigg(\frac{\partial u}{\partial y}\frac{\partial^2u}{\partial x\partial y}-\frac{\partial^2u}{\partial y^2}\frac{\partial u}{\partial x}\bigg)$$
    进一步化简得到
    $$\varphi_C''(x)=-\frac{1}{\big(\frac{\partial u}{\partial y}\big)^3}\bigg(\bigg(\frac{\partial u}{\partial y}\bigg)^2\frac{\partial^2u}{\partial x^2}-2\frac{\partial u}{\partial x}\frac{\partial u}{\partial y}\frac{\partial^2u}{\partial x\partial y}+\bigg(\frac{\partial u}{\partial x}\bigg)^2\frac{\partial^2u}{\partial y^2}\bigg)$$
    由此即证明了$\varphi_C''(x)$对任何$C$为0当且仅当
    $$\bigg(\frac{\partial u}{\partial y}\bigg)^2\frac{\partial^2u}{\partial x^2}-2\frac{\partial u}{\partial x}\frac{\partial u}{\partial y}\frac{\partial^2u}{\partial x\partial y}+\bigg(\frac{\partial u}{\partial x}\bigg)^2\frac{\partial^2u}{\partial y^2}=0$$

    \note 注意$f(x,y)$对每个$C$均满足$f(x,y)=C$未必有$f(x,y)=\alpha x+\beta y+\gamma$，例如$f(x,y)=\sqrt[3]{x}$也满足要求。此外，证明中必须考虑到对每个$C$均满足，若设$f(x,y)=C$的斜率为$k$，需要先证明对每个$C$有共同的$k$\ (否则$f(x,y)=C_1$与$f(x,y)=C_2$将相交，只能推出$C_1=C_2$)。
\end{enumerate}

\subsection{下册建议}
在经历了一学期的高等数学学习后，相信大家都总结了一些自己的经验与学习方法。因此，本节主要是给大家介绍一些常见的误区，分为三个方面：做题角度、理解角度与学习状态角度。

\subsubsection{做题与背题}
首先是做题角度，一个常见的问题是认为自己来自``题做少了''或``做题花的时间少了''。以\textbf{做题为主}(而非听课或看书)的学习模式确实是值得推荐的模式，不过，学习阶段将教材上的题目都做一遍是\textbf{完全够用}的，至于复习阶段，只要有上学期与本学期习题课讲义中的复习题也足够了。

事实上，之所以会觉得问题来自做题的量不够，往往是由于做题中倾向于了一种类似\textbf{背题}的模式。所谓的背题并不是指纯粹机械记忆过程，而是指将题中涉及的当前还不会的部分都当作\textbf{技巧的积累}，且不去在乎\textbf{想到技巧的方式}。

做题与背题的本质区别

往年题

高低观点

\note 值得一提的是，如果需要通过很短的时间得到不太差的成绩，背题是一种\textbf{行之有效}的方式，只是它存在很大的\textbf{风险}：背题能取得不错的结果非常依赖出题模式在背过的范围内，只要出题相对灵活，背题的收益就会降低了。就高等数学A来说，``不让背题取得好成绩''的确是出题的趋势。

\subsubsection{直观与严谨}
\subsubsection{状态}

\newpage
\section{重积分}
本章介绍了从矩形区域出发定义重积分的思路，与此思路下重积分化为累次积分的方法(尤其是如何从代数而非几何出发研究问题)，并进一步介绍了换元法的使用，知识基础为定积分的定义与计算。

\subsection{习题解答}
\subsubsection{作业解答}
\begin{enumerate}
    \item (习题7.1.2)设平面区域
    $$D=\{(x,y)\mid-1\le x\le1,0\le y\le1\}$$
    由定义证明
    $$\iint_Dx\dr\sigma=0$$

    \sol{
        \note 本题已默认连续函数可积，证明以\textbf{可积}为基础。

        对正整数$n$，考虑划分：以$\frac{1}{n}$为单位划分$x$轴为$2n$个区间、$y$轴为$n$个区间，将$x$取$\big[-1+\frac{j}{n},-1+\frac{j+1}{n}\big]$、$y$取$\big[\frac{k}{n},\frac{k+1}{n}\big]$构成的矩形记为
        $$D_{jk},\quad j=0,\dots,2n-1,\quad k=0,\dots,n-1$$
        每个矩形的面积均为$\frac{1}{n^2}$。在$D_{jk}$中取点
        $$(x_{jk},y_{jk})=\bigg(-1+\frac{j}{n}+\frac{1}{2n},\frac{k}{n}\bigg)$$
        以此进行估算可得到黎曼和(对$j$直接使用等差数列求和公式可发现为0)
        $$\sum_{k=0}^{n-1}\sum_{j=0}^{2n-1}x_{jk}\cdot\frac{1}{n^2}=\frac{1}{n^2}\sum_{k=0}^{n-1}\sum_{j=0}^{2n-1}\bigg(-1+\frac{j}{n}+\frac{1}{2n}\bigg)=0$$
        由于此划分的直径即为正方形对角线长度$\frac{\sqrt2}{n}$，当$n\to\infty$时极限为0。
        
        由于已知可积，若积分值为$c\ne0$，存在$\delta>0$使得只要分割半径小于$\delta$，任意取点，得到的黎曼和应与$c$差距在$\frac{|c|}{2}$以内，不可能为0，矛盾，因此积分值只能为0。
    }

    \item (习题7.1.3)设函数$f(x,y)$在有界闭区域$D$上连续，$g(x,y)$在$D$上非负且可积，且$f(x,y)g(x,y)$在$D$上也可积。证明存在$(x_0,y_0)\in D$使得
    $$f(x_0,y_0)\iint_Dg(x,y)\dr\sigma=\iint_Df(x,y)g(x,y)\dr\sigma$$

    \sol{
        利用有界闭区域的\textbf{最值定理}(证明可见上学期讲义14.3.2)，$f$在$D$上存在最大值与最小值，设为$M$与$m$，则由$g$非负性
        $$\forall(x,y)\in D,\quad mg(x,y)\le f(x,y)g(x,y)\le Mg(x,y)$$
        利用重积分性质有
        $$m\iint_Dg(x,y)\dr\sigma\le\iint_Df(x,y)g(x,y)\dr\sigma\le M\iint_Dg(x,y)\dr\sigma$$

        由$g$非负可知
        $$\iint_Dg(x,y)\dr\sigma\ge0$$

        分类讨论：
        \begin{itemize}
            \item 若$\iint_Dg(x,y)\dr\sigma=0$，则上方不等式左侧右侧均为0，因此中间也必然为0，任取$(x_0,y_0)\in D$，要证明的等式左右均为0，即得相等。 
            
            \note 注意条件并未给出$g$连续，因此$g$的积分为0无法推出$g$恒为0。

            \item 若$\iint_Dg(x,y)\dr\sigma>0$，可将不等式改写为
            $$m\le\frac{\iint_Df(x,y)g(x,y)\dr\sigma}{\iint_Dg(x,y)\dr\sigma}\le M$$
            利用有界闭区域的\textbf{介值定理}(证明可见上学期讲义14.3.2)，存在$(x_0,y_0)\in D$使得
            $$f(x_0,y_0)=\frac{\iint_Df(x,y)g(x,y)\dr\sigma}{\iint_Dg(x,y)\dr\sigma}$$
            变形即得结论。
        \end{itemize}
    }

    \item (习题7.2.2)改变下列累次积分的积分次序：
    \begin{enumerate}
        \item[(3)] $$\int_0^2\dr x\int_{2x}^{6-x}f(x,y)\dr y$$
        
        \sol{
            分为如下步骤：
            \begin{itemize}
                \item \textbf{区域描述}
                
                根据条件，区域可描述为满足
                $$0\le x\le 2,\quad 2x\le y\le 6-x$$
                的$(x,y)$集合。

                \item \textbf{累次积分化}
                
                为了化为先$x$后$y$的累次积分，我们需要将集合描述成
                $$\alpha\le y\le\beta,\quad \phi_1(y)\le x\le\phi_2(y)$$
                当$x$可任取时，$y$最小对应$x$为0时，为0；$y$最大对应$x$为0时，为6。综合得第一个不等式为
                $$0\le y\le 6$$
                当$y$给定时，关于$x$的不等式可改写为
                $$0\le x\le 2,\quad x\le 6-y,\quad x\le\frac{y}{2}$$
                $x$必须不超过三个上限中最小的，进一步估算可发现$6-y$与$\frac{y}{2}$中较小的一定不超过2，因此范围即
                $$0\le x\le\min\bigg\{6-y,\frac{y}{2}\bigg\}$$
                最终写出积分
                $$\int_0^6\dr y\int_0^{\min\{6-y,y/2\}}f(x,y)\dr x$$
                或讨论$6-y$与$\frac{y}{2}$何者较小而拆分为
                $$\int_0^4\dr y\int_0^{y/2}f(x,y)\dr x+\int_4^6\dr y\int_0^{6-y}f(x,y)\dr x$$
            \end{itemize}
        }

        \item[(5)] $$\int_1^\er\dr x\int_0^{\ln x}f(x,y)\dr y$$
        
        \sol{
            分为如下步骤：
            \begin{itemize}
                \item \textbf{区域描述}
                
                根据条件，区域可描述为满足
                $$1\le x\le\er,\quad 0\le y\le \ln x$$
                的$(x,y)$集合。

                \item \textbf{累次积分化}
                
                为了化为先$x$后$y$的累次积分，我们需要将集合描述成
                $$\alpha\le y\le\beta,\quad \phi_1(y)\le x\le\phi_2(y)$$
                当$x$可任取时，$y$最小一定能取到0；$y$最大对应$x$为$\er$时，为1。综合得第一个不等式为
                $$0\le y\le1$$
                当$y$给定时，关于$x$的不等式可改写为
                $$0\le x\le\er,\quad x\ge\er^y$$
                $x$必须大于等于两个下限中最大的，也即$\er^y$，因此范围即
                $$\er^y\le x\le1$$
                最终写出积分
                $$\int_0^1\dr y\int_{\er^y}^\er f(x,y)\dr x$$
            \end{itemize}
        }
    \end{enumerate}

    \item (习题7.2.3)计算
    $$\iint_Dy\dr x\dr y$$
    其中$D$是由$y=0$与$y=\sin x$在$x\in[0,\pi]$中围成的区域。

    \sol{
        分为如下步骤：
        \begin{itemize}
            \item \textbf{区域描述}
                
            根据条件，区域可描述为满足
            $$0\le x\le\pi,\quad 0\le y\le\sin x$$
            的$(x,y)$集合。

            \item \textbf{累次积分化}
                
            根据此形式可直接将重积分化为累次积分
            $$\int_0^\pi\dr x\int_0^{\sin x}y\dr y$$

            \item \textbf{积分计算}
                
            对$y$积分可将累次积分化为
            $$\frac{1}{2}\int_0^\pi\sin^2x\dr x$$
            拆分为0到$\frac{\pi}{2}$与$\frac{\pi}{2}$到$\pi$两部分积分，第二部分利用对称性换元$t=\pi-x$可发现值与第一部分相同，从而上述积分为
            $$\int_0^{\pi/2}\sin^2x\dr x$$
            此时可直接得出结果为$\frac{\pi}{4}$\ (可见上册教材3.4节例2)。
        \end{itemize}        
    }

    \item (习题7.2.8)计算
    $$\int_0^\pi\dr x\int_x^\pi\frac{\sin y}{y}\dr y$$

    \sol{
        由于$\frac{\sin y}{y}$不存在初等原函数，我们需要考虑更换累次积分顺序。分为如下步骤：
        \begin{itemize}
            \item \textbf{区域描述}
                
            根据条件，区域可描述为满足
            $$0\le x\le\pi,\quad x\le y\le\pi$$
            的$(x,y)$集合。

            \item \textbf{累次积分化}
                
            为了化为先$x$后$y$的累次积分，我们需要将集合描述成
            $$\alpha\le y\le\beta,\quad \phi_1(y)\le x\le\phi_2(y)$$
            当$x$可任取时，$y$最小对应$x$为0时，为0；$y$最大一定可取到$\pi$。综合得第一个不等式为
            $$0\le y\le\pi$$
            此时再考虑对$x$的约束可写出第二个不等式
            $$0\le x\le y$$
            最终写出积分
            $$\int_0^\pi\dr y\int_0^y\frac{\sin y}{y}\dr x$$

            \item \textbf{积分计算}
                
            对$x$积分可将累次积分化为
            $$\int_0^\pi\sin y\dr y=2$$
        \end{itemize}
    }

    \item (习题7.2.14)利用极坐标计算
    $$\int_{-1}^0\dr x\int_{-\sqrt{1-x^2}}^0\frac{2}{1+\sqrt{x^2+y^2}}\dr y$$

    \sol{
        分为如下步骤：
        \begin{itemize}
            \item \textbf{区域描述}
                
            根据条件，区域可描述为满足
            $$-1\le x\le0,\quad -\sqrt{1-x^2}\le y\le0$$
            的$(x,y)$集合。

            \item \textbf{旋转对称性换元}
            
            利用旋转对称性，进行极坐标换元$x=r\cos\theta$、$y=r\sin\theta$。

            此时，原本的被积函数可写为$\frac{2}{1+r}$，且计算Jacobi行列式得$\dr x\dr y=r\dr r\dr\theta$。区域的两个条件可写为
            $$-1\le r\cos\theta\le0,\quad-\sqrt{1-r^2\cos^2\theta}\le r\sin\theta\le0$$
            将上述条件与极坐标自然范围确定的$(r,\theta)$区域记为$D_0$，原积分化为
            $$\iint_{D_0}\frac{2}{1+r}r\dr r\dr\theta$$

            \item \textbf{累次积分化}
                
            由于被积函数对于$\theta$为常数，我们希望化为先$\theta$后$r$的累次积分，也即需要将集合描述成
            $$\alpha\le r\le\beta,\quad\phi_1(r)\le\theta\le\phi_2(r)$$

            \note 注意极坐标\textbf{自然带有范围}$r\ge0$、$\theta\in[0,2\pi]$。

            将第二个约束两边同平方，由于已知$r\sin\theta\le0$，同平方后不等号反号，也即得到
            $$1-r^2\cos\theta\ge r^2\sin^2\theta$$
            移项后结合极坐标自然范围有$0\le r\le1$。此时，进一步化简约束条件可发现只剩$\cos\theta\le0$、$\sin\theta\le0$，从而结合极坐标自然范围有
            $$\pi\le\theta\le\frac{3\pi}{2}$$
            这已经是符合要求的范围，从而可最终写出积分
            $$\int_0^1\dr r\int_\pi^{3\pi/2}\frac{2r}{1+r}\dr\theta$$

            \item \textbf{积分计算}
                
            对$\theta$积分可将累次积分化为
            $$\frac{\pi}{2}\int_0^1\frac{2r}{1+r}\dr r=\pi\int_0^1\bigg(1-\frac{1}{1+r}\bigg)\dr r=\pi\cdot(1-(\ln 2-\ln1))=\pi(1-\ln2)$$
        \end{itemize}
    }

    \item (习题7.2.16)利用极坐标计算
    $$\int_0^R\dr x\int_0^{\sqrt{R^2-x^2}}\ln(1+x^2+y^2)\dr y$$
    这里参数$R>0$。

    \sol{
        分为如下步骤：
        \begin{itemize}
            \item \textbf{区域描述}
                
            根据条件，区域可描述为满足
            $$0\le x\le R,\quad 0\le y\le\sqrt{R^2-x^2}$$
            的$(x,y)$集合。

            \item \textbf{旋转对称性换元}
            
            利用旋转对称性，进行极坐标换元$x=r\cos\theta$、$y=r\sin\theta$。

            此时，原本的被积函数可写为$\ln(1+r^2)$，且计算Jacobi行列式得$\dr x\dr y=r\dr r\dr\theta$。区域的两个条件可写为
            $$0\le r\cos\theta\le R,\quad 0\le r\sin\theta\le\sqrt{R^2-r^2\cos^2\theta}$$
            将上述条件与极坐标自然范围确定的$(r,\theta)$区域记为$D_0$，原积分化为
            $$\iint_{D_0}\ln(1+r^2)r\dr r\dr\theta$$

            \item \textbf{累次积分化}
                
            由于被积函数对于$\theta$为常数，我们希望化为先$\theta$后$r$的累次积分，也即需要将集合描述成
            $$\alpha\le r\le\beta,\quad\phi_1(r)\le\theta\le\phi_2(r)$$

            将第二个约束两边同平方，由于已知$r\sin\theta\ge0$，同平方后不等号不变，也即得到
            $$r^2\sin^2\theta\le R^2-r^2\cos^2\theta$$
            移项后结合极坐标自然范围有$0\le r\le R$。此时，进一步化简约束条件可发现只剩$\cos\theta\ge0$、$\sin\theta\ge0$，从而结合极坐标自然范围有
            $$0\le\theta\le\frac{\pi}{2}$$
            这已经是符合要求的范围，从而可最终写出积分
            $$\int_0^R\dr r\int_0^{\pi/2}\ln(1+r^2)r\dr\theta$$

            \item \textbf{积分计算}
                
            对$\theta$积分可将累次积分化为(换元$s=r^2$)
            $$\frac{\pi}{2}\int_0^Rr\ln(1+r^2)\dr r=\frac{\pi}{4}\int_0^{R^2}\ln(1+s)\dr s$$
            分部积分得到(最后一部分计算类似前一题)
            $$\int_0^{R^2}\ln(1+s)\dr s=R^2\ln(1+R^2)-0\ln(1+0)-\int_0^R\frac{s}{s+1}\dr s=R^2\ln(1+R^2)-(R^2-\ln(1+R^2))$$
            从而化简得原题最终结果为
            $$\frac{\pi}{4}\big((R^2+1)\ln(R^2+1)-R^2\big)$$
        \end{itemize}
    }

    \item (习题7.2.18)计算
    $$\iint_Dr\dr\sigma$$
    其中$D$是在极坐标下的心脏线$r=a(1+\cos\theta)$内侧、圆周$r=a$外侧的部分，这里参数$a>0$。

    \sol{
        分为如下步骤：
        \begin{itemize}
            \item \textbf{区域描述与换元}
                
            我们将条件确定的$(r,\theta)$区域记为$D_0$\ (注意它与积分区域$D$不同，因为$D$是针对$x$、$y$的)，其为满足
            $$r\ge0,\quad0\le\theta\le2\pi,\quad a\le r\le a(1+\cos\theta)$$
            的$(r,\theta)$集合。

            计算极坐标Jacobi行列式得$\dr\sigma=\dr x\dr y=r\dr r\dr\theta$，原积分化为
            $$\iint_{D_0}r^2\dr r\dr\theta$$

            \note 注意$\dr\sigma$就是$\dr x\dr y$。

            \item \textbf{累次积分化}
                
            由于极坐标下心脏线与圆$r$被$\theta$限制，我们希望化为先$r$后$\theta$的累次积分，也即需要将集合描述成
            $$\alpha\le\theta\le\beta,\quad \phi_1(\theta)\le r\le\phi_2(\theta)$$

            由于第三个约束需要满足$a\le a(1+\cos\theta)$，也即$\cos\theta\ge0$，对应
            $$0\le\theta\le\frac{\pi}{2},\quad\frac{3\pi}{2}\le\theta\le2\pi$$
            两部分，对每部分$r$的约束均为
            $$a\le r\le a(1+\cos\theta)$$
            从而可最终写出积分
            $$\int_0^{\pi/2}\dr\theta\int_a^{a(1+\cos\theta)}r^2\dr r+\int_{3\pi/2}^{2\pi}\dr\theta\int_a^{a(1+\cos\theta)}r^2\dr r$$

            \item \textbf{积分计算}
                
            由于$\cos\theta=\cos(2\pi-\theta)$，在上述第二个积分中将$\theta$换元为$2\pi-\theta$可发现值与第一个积分相等，从而结果可化为
            $$2\int_0^{\pi/2}\dr\theta\int_a^{a(1+\cos\theta)}r^2\dr r$$

            \note 事实上两个积分相等也可以通过区域的\textbf{对称性}进行先行的化简。

            直接对$r$积分得其为
            $$\frac{2a^3}{3}\int_0^{\pi/2}(\cos^3\theta+3\cos^2\theta+3\cos\theta)\dr\theta$$
            由上册教材3.4节例2可计算出结果为
            $$\frac{2a^3}{3}\bigg(\frac{2}{3}+3\cdot\frac{\pi}{4}+3\cdot1\bigg)=\bigg(\frac{22}{9}+\frac{\pi}{2}\bigg)a^3$$
        \end{itemize}
    }

    \item (习题7.2.22)计算
    $$\iint_D\bigg(\sqrt{\frac{y}{x}}+\sqrt{xy}\bigg)\dr x\dr y$$
    其中$D$是第一象限内$xy=1$、$xy=9$、$y=x$、$y=4x$围成的区域。

    \sol{
        分为如下步骤：
        \begin{itemize}
            \item \textbf{区域描述}
            
            由于``围成''需要确定一个\textbf{有界}的区域，必然$xy$在1与9之间，$y$在$x$与$4x$之间，结合第一象限条件也即区域可描述为满足
            $$1\le xy\le 9,\quad 0\le x\le y\le 4x$$
            的$(x,y)$集合。

            \item \textbf{换元与累次积分化}
            
            由于直接以$x$、$y$描述区域相对复杂，我们希望区域能拥有更简单的形式，而最简单的形式即\textbf{矩形}，从而换元
            $$\xi=xy,\quad\eta=\frac{y}{x}$$
            这样区域即成为
            $$1\le\xi\le 9,\quad 1\le\eta\le 4$$
            将上述条件确定的$(\xi,\eta)$区域记为$D_0$，被积函数可写为$\sqrt{\xi}+\sqrt{\eta}$。由于在第一象限，可反解出
            $$x=\sqrt{\frac{\xi}{\eta}},\quad y=\sqrt{\xi\eta}$$
            直接计算Jacobi行列式可知
            $$\begin{vmatrix}\frac{1}{2\sqrt{\xi\eta}}&\frac{\sqrt{\eta}}{2\sqrt{\xi}}\\-\frac{\sqrt{\xi}}{2\sqrt{\eta^3}}&\frac{\sqrt{\xi}}{2\sqrt{\eta}}\end{vmatrix}=\frac{1}{2\eta}$$
            由于其在$D_0$中恒正，换元是合理的，且
            $$\dr x\dr y=\frac{1}{2\eta}\dr\xi\dr\eta$$
            于是原积分化为(由于被积函数中$\xi$的形式更简单，先对$\xi$积分)
            $$\int_1^4\dr\eta\int_1^9\frac{\sqrt{\xi}+\sqrt{\eta}}{2\eta}\dr\xi$$

            \item \textbf{积分计算}
                
            对$\xi$积分，直接利用多项式积分结论可将累次积分化为
            $$\int_1^4\frac{\frac{2}{3}(27-1)+8\sqrt{\eta}}{2\eta}\dr\eta=\int_1^4\bigg(\frac{26}{3\eta}+\frac{4}{\sqrt\eta}\bigg)\dr\eta$$
            进一步利用多项式积分结论得结果为
            $$\frac{26}{3}\ln 4+8=\frac{52}{3}\ln2+8$$
        \end{itemize}

        \note 原题并未说明第一象限，此时四条曲线无法确定唯一的\textbf{区域}。
    }

    \item (习题7.2.24)计算
    $$\iint_\Omega(x^2+y^2)\dr x\dr y$$
    其中$\Omega$为给定参数$a>0$、$b>0$的椭圆
    $$\bigg\{(x,y)\ \bigg|\ \frac{x^2}{a^2}+\frac{y^2}{b^2}\le1\bigg\}$$

    \sol{
        分为如下步骤：
        \begin{itemize}
            \item \textbf{换元与累次积分化}
            
            由于直接以$x$、$y$描述区域相对复杂，为了将区域换元为矩形，我们考虑\textbf{广义极坐标}换元
            $$x=ar\cos\theta,\quad y=br\sin\theta$$
            这样约束条件即化为了$r^2\le1$，结合极坐标自然范围有
            $$0\le r\le1,\quad0\le\theta\le2\pi$$
            将上述条件确定的$(r,\theta)$区域记为$D_0$，被积函数可写为$a^2r^2\cos^2\theta+b^2r^2\sin^2\theta$。 

            直接计算Jacobi行列式可知
            $$\begin{vmatrix}a\cos\theta&b\sin\theta\\-ar\cos\theta&br\sin\theta\end{vmatrix}=abr$$

            类似极坐标可说明换元合理(即使Jacobi行列式在区域中某点为0，详见本讲义19.3.2)，且
            $$\dr x\dr y=abr\dr r\dr\theta$$
            于是原积分化为(由于被积函数中$r$的形式更简单，先对$r$积分)
            $$\int_0^{2\pi}\dr\theta\int_0^1abr(a^2r^2\cos^2\theta+b^2r^2\sin^2\theta)\dr r$$

            \item \textbf{积分计算}
            
            对$r$积分，直接利用多项式积分结论可将累次积分化为
            $$\frac{ab}{4}\int_0^{2\pi}(a^2\cos^2\theta+b^2\sin^2\theta)\dr\theta$$
            类似本部分第4题的对称性处理可知
            $$\int_0^{2\pi}\cos^2\theta\dr\theta=\int_0^{2\pi}\sin^2\theta\dr\theta=4\int_0^{\pi/2}\sin^2\theta\dr\theta=\pi$$
            由此得到最终结果为
            $$\frac{\pi ab}{4}(a^2+b^2)$$        
        \end{itemize}
    }

    \item (总练习题7.1(4))改变如下累次积分的积分次序：
    $$\int_0^1\dr y\int_{-\sqrt{1-y^2}}^{\sqrt{1-y^2}}3y\dr x$$
    
    \sol{
        分为如下步骤：
        \begin{itemize}
            \item \textbf{区域描述}
                
            根据条件，区域可描述为满足
            $$0\le y\le1,\quad-\sqrt{1-y^2}\le x\le\sqrt{1-y^2}$$
            的$(x,y)$集合。

            \item \textbf{累次积分化}
                
            为了化为先$y$后$x$的累次积分，我们需要将集合描述成
            $$\alpha\le x\le\beta,\quad \phi_1(x)\le y\le\phi_2(x)$$
            当$y$可任取时，$x$最小对应$y$为0时，为$-1$；$x$最大对应$y$为0时，为1。综合得第一个不等式为
            $$-1\le x\le1$$
            当$x$给定时，关于$y$的不等式可改写为(考虑$x$正负的情况)
            $$0\le y\le1,\quad x^2+y^2\le1$$
            进一步估算可发现这确定了范围
            $$0\le y\le\sqrt{1-x^2}$$
            最终写出积分
            $$\int_{-1}^1\dr x\int_0^{\sqrt{1-x^2}}3y\dr y$$
        \end{itemize}
    }

    \item (总练习题7.2(2))求曲线$x=y^2$与$x=2y-y^2$围成区域的面积。
    
    \sol{
        分为如下步骤：
        \begin{itemize}
            \item \textbf{区域描述}
                
            由于``围成''需要确定一个有界的区域，从$x$的有界性可知必然为
            $$y^2\le x\le 2y-y^2$$
            或
            $$2y-y^2\le x\le y^2$$
            直接求解可发现$y^2\le 2y-y^2$等价于$y\in[0,1]$，有界，而另一边$y$无界，因此最终确定出区域是满足
            $$y^2\le x\le 2y-y^2$$
            的$(x,y)$集合。

            \item \textbf{累次积分化}
            
            我们已经求解出$0\le y\le1$，这结合$2y-y^2\le x\le y^2$已经是符合要求的范围，从而可最终写出积分(面积即为对1积分)
            $$\int_0^1\dr y\int_{y^2}^{2y-y^2}1\dr x$$
            
            \item \textbf{积分计算}
            
            直接利用多项式积分结论可将累次积分化为
            $$\int_0^1(2y-y^2-y^2)\dr y=\frac{1}{3}$$
        \end{itemize}
    }

    \item (总练习题7.4(1))计算
    $$\int_0^1\dr y\int_{2y}^24\cos(x^2)\dr x$$

    \sol{
        由于$\cos(x^2)$不存在初等原函数，我们需要考虑更换累次积分顺序。分为如下步骤：
        \begin{itemize}
            \item \textbf{区域描述}
                
            根据条件，区域可描述为满足
            $$0\le y\le1,\quad 2y\le x\le2$$
            的$(x,y)$集合。

            \item \textbf{累次积分化}
                
            为了化为先$y$后$x$的累次积分，我们需要将集合描述成
            $$\alpha\le x\le\beta,\quad \phi_1(x)\le y\le\phi_2(x)$$
            当$y$可任取时，$x$最小对应$y$为0时，为0；$x$最大一定可取到2。综合得第一个不等式为
            $$0\le x\le2$$
            此时再考虑对$y$的约束可写出第二个不等式
            $$0\le y\le\frac{x}{2}$$
            最终写出积分
            $$\int_0^2\dr x\int_0^{x/2}4\cos(x^2)\dr y$$

            \item \textbf{积分计算}
                
            对$y$积分可将累次积分化为(换元$t=x^2$)
            $$\int_0^22x\cos(x^2)\dr x=\int_0^4\cos t\dr t=\sin4$$
        \end{itemize}
    }

    \item (总练习题7.8)求极坐标表示的双扭线$r^2=4\cos(2\theta)$围成部分的面积。
    
    \sol{
        将双扭线围成部分记为$D$，所求即 
        $$\iint_D1\dr x\dr y$$
        分为如下步骤：
        \begin{itemize}
            \item \textbf{区域描述与换元}
                
            我们将条件确定的$(r,\theta)$区域记为$D_0$\ (注意它与积分区域$D$不同，因为$D$是针对$x$、$y$的)，其为满足
            $$r\ge0,\quad0\le\theta\le2\pi,\quad r^2\le4\cos(2\theta)$$
            的$(r,\theta)$集合。

            计算极坐标Jacobi行列式得$\dr\sigma=\dr x\dr y=r\dr r\dr\theta$，原积分化为
            $$\iint_{D_0}r\dr r\dr\theta$$

            \note 注意$\dr\sigma$就是$\dr x\dr y$。

            \item \textbf{累次积分化}
                
            由于极坐标下双扭线$r$被$\theta$限制，我们希望化为先$r$后$\theta$的累次积分，也即需要将集合描述成
            $$\alpha\le\theta\le\beta,\quad\phi_1(\theta)\le r\le\phi_2(\theta)$$

            由于$r^2\le4\cos(2\theta)$需要能确定合理的$r$，必然有
            $$\cos(2\theta)\ge0$$
            因此结合极坐标自然范围知$\theta$的范围对应
            $$0\le\theta\le\frac{\pi}{4},\quad\frac{3\pi}{4}\le\theta\le\frac{5\pi}{4},\quad\frac{7\pi}{4}\le\theta\le2\pi$$
            三部分，对每部分$r$的约束均为
            $$0\le r\le2\sqrt{\cos(2\theta)}$$
            从而可最终写出积分
            $$\int_0^{\pi/4}\dr\theta\int_0^{2\sqrt{\cos(2\theta)}}r\dr r+\int_{3\pi/4}^{5\pi/4}\dr\theta\int_0^{2\sqrt{\cos(2\theta)}}r\dr r+\int_{7\pi/4}^{2\pi}\dr\theta\int_0^{2\sqrt{\cos(2\theta)}}r\dr r$$

            \item \textbf{积分计算}
            
            由于$\cos(2\theta)=\cos(2(2\pi-\theta))$，在上述第三个积分中将$\theta$换元为$2\pi-\theta$可发现值与第一个积分相等。将第二个积分拆分为$3\pi/4$到$\pi$与$\pi$到$5\pi/4$两段，分别将$\theta$换元为$\pi-\theta$与$\theta-\pi$可发现值与第一个积分相等。综合以上，结果可化为
            $$4\int_0^{\pi/4}\dr\theta\int_0^{2\sqrt{\cos(2\theta)}}r\dr r$$

            \note 这部分讨论从\textbf{几何直观}来看非常自然，但仍然需要知道代数上处理的方法。

            进一步计算即得原积分为
            $$2\int_0^{\pi/4}(2\sqrt{\cos(2\theta)})^2\dr\theta=8\int_0^{\pi/4}\cos(2\theta)\dr\theta=4$$
        \end{itemize}
    }
\end{enumerate}
\subsubsection{补充题解答}
\begin{enumerate}
    \item (习题7.1.4)设函数$f(x,y)$在有界闭区域$D$上非负连续，且
    $$\iint_Df(x,y)\dr\sigma=0$$
    证明$f(x,y)$在$D$上恒为0。

    \note 本题的实际严谨证明难度远高于教材放这题时对难度的估算，可当作\extra。
    
    \sol{
        为了保证严谨性，我们将过程分为两次尝试：
        \begin{itemize}
            \item \textbf{直接尝试}
            
            若否，由非负性可知存在$(x_0,y_0)\in D$使得设$f(x_0,y_0)=c>0$。利用连续性(注意一般集合的连续性定义，可见上学期讲义14.3.1)，存在$\delta>0$使得
            $$\forall (x,y)\in D,\quad\sqrt{(x-x_0)^2+(y-y_0)^2}<\delta\Longrightarrow|f(x,y)-f(x_0,y_0)|<\frac{c}{2}$$
            将满足$\sqrt{(x-x_0)^2+(y-y_0)^2}<\delta$的$(x,y)\in D$构成的集合记作$E$，我们构造函数$g(x,y)$满足
            $$g(x,y)=\begin{cases}\frac{c}{2}&(x,y)\in E\\0&(x,y)\notin E\end{cases}$$
            可以发现对任何$(x,y)\in D$均有$f(x,y)\ge g(x,y)$。

            但是，以此构造的$g$\textbf{无法保证可积}(因为$D$的边界可能性质很差，而$E$的形状可能包含$D$的边界)，即使可积，\textbf{也难以说明积分非零}(因为涉及$E$的面积，而这是难以确定的)。因此，这个方法无法进行进一步的估算，需要更好的方式。

            \item \textbf{引理证明}
            
            进一步思考可以发现，我们事实上希望构造过程不受$D$的\textbf{边界}的干扰。

            因此，我们先证明引理：若连续函数$f$在有界闭区域$D$中不恒为0，一定存在\textbf{内点}$(x_0,y_0)$使得
            $$f(x_0,y_0)>0$$
            根据有界闭区域的定义，它是某个开区域$D_0$的\textbf{闭包}，而$D_0$就是它的内点集合。
            
            若$f$在$D_0$中恒为0，根据闭包的\textbf{等价定义}(可见上学期讲义13.3.2)，$D$上的任何一点$(x,y)$都存在$D_0$中的点列$(x_n,y_n)$使得
            $$\lim_{n\to\infty}(x_n,y_n)=(x,y)$$
            再利用连续函数的\textbf{归结原理}(可见上学期讲义14.2.2)即得到
            $$f(x,y)=\lim_{n\to\infty}f(x_n,y_n)=0$$
            与$f$在$D$上不恒为0矛盾。

            \note 如果对这部分证明不熟悉，可以参考上学期讲义13.3进行复习，不过这并非本学期重点。

            \item \textbf{最终构造}
            
            若否，利用引理与非负性可知存在\textbf{内点}$(x_0,y_0)\in D$使得设$f(x_0,y_0)=c>0$。利用连续性，存在$\delta>0$使得
            $$\forall (x,y)\in D,\quad\sqrt{(x-x_0)^2+(y-y_0)^2}<\delta\Longrightarrow|f(x,y)-f(x_0,y_0)|<\frac{c}{2}$$
            再利用内点定义，存在$\delta_0>0$使得
            $$\sqrt{(x-x_0)^2+(y-y_0)^2}<\delta_0\Longrightarrow(x,y)\in D$$

            由此，记$\gamma=\min\{\delta,\delta_0\}$，将满足$\sqrt{(x-x_0)^2+(y-y_0)^2}<\gamma$的$(x,y)$构成的集合记作$E$，我们构造函数$g(x,y)$满足
            $$g(x,y)=\begin{cases}\frac{c}{2}&(x,y)\in E\\0&(x,y)\notin E\end{cases}$$
            可以发现对任何$(x,y)\in D$均有$f(x,y)\ge g(x,y)$。

            此时，$g$在$D$中的某半径为$\gamma$的小圆上为$\frac{c}{2}$，其他为0，根据圆的光滑性可知其可积(一些讨论可见本讲义19.2.3)，积分即为$\frac{c}{2}$乘圆面积，即
            $$\iint_Dg(x,y)\dr\sigma=\frac{c\pi\gamma^2}{2}>0$$
            再利用重积分性质可知
            $$\iint_Df(x,y)\dr\sigma\ge\iint_Dg(x,y)\dr\sigma>0$$
            与条件积分为0矛盾。
        \end{itemize}
        \note 本题的核心是$g$的构造思路，可与上学期讲义6.3.1一元函数时的情况类比。
    }

    \item (习题7.2.1(5))将二重积分
    $$\iint_Df(x,y)\dr\sigma$$
    化为累次积分，其中
    $$D=\{(x,y)\mid x^2+y^2\le y\}$$

    \sol{
        分为如下步骤：
        \begin{itemize}
            \item \textbf{区域描述}
                
            由于区域是一个圆，我们改写为更符合圆的形式，即满足
            $$x^2+\bigg(y-\frac{1}{2}\bigg)^2\le\frac{1}{2^2}$$
            的$(x,y)$集合。

            \item \textbf{累次积分化}
                
            由于$x$用$y$表示相对简单，我们试着将它化为先$x$后$y$的累次积分，即
            $$\alpha\le y\le\beta,\quad \phi_1(y)\le x\le\phi_2(y)$$

            \note 另一边的复杂度差别并不大，大家可以自行尝试。

            由积分区域为圆且半径$\frac{1}{2}$，圆心纵坐标$\frac{1}{2}$，可以给出$y$的范围限制
            $$0\le y\le1$$
            直接从区域原表达式化简可得到$x$的范围
            $$-\sqrt{y-y^2}\le x\le\sqrt{y-y^2}$$
            
            最终写出积分
            $$\int_0^1\dr y\int_{-\sqrt{y-y^2}}^{\sqrt{y-y^2}}f(x,y)\dr x$$
        \end{itemize}
    }

    \item (习题7.2.2(6))改变如下累次积分的积分次序：
    $$\int_0^5\dr y\int_y^{10-y}f(x,y)\dr x$$

    \sol{
        分为如下步骤：
        \begin{itemize}
            \item \textbf{区域描述}
                
            根据条件，区域可描述为满足
            $$0\le y\le 5,\quad y\le x\le10-y$$
            的$(x,y)$集合。

            \item \textbf{累次积分化}
                
            为了化为先$y$后$x$的累次积分，我们需要将集合描述成
            $$\alpha\le x\le\beta,\quad \phi_1(x)\le y\le\phi_2(x)$$
            当$y$可任取时，$x$最小对应$y$为0时，为0；$x$最大对应$y$为0时，为10。综合得第一个不等式为
            $$0\le x\le10$$
            当$x$给定时，关于$y$的不等式可改写为
            $$0\le y\le 5,\quad y\le x,\quad y\le10-x$$
            $x$必须不超过三个上限中最小的，进一步估算可发现$10-x$与$x$中较小的一定不超过5，因此范围即
            $$0\le y\le\min\bigg\{10-x,x\bigg\}$$
            最终写出积分
            $$\int_0^{10}\dr x\int_0^{\min\{10-x,x\}}f(x,y)\dr y$$
            或讨论$10-x$与$x$何者较小而拆分为
            $$\int_0^5\dr x\int_0^xf(x,y)\dr y+\int_5^{10}\dr x\int_0^{10-x}f(x,y)\dr y$$
        \end{itemize}
    }

    \item (习题7.2.9)计算
    $$\int_0^2\dr x\int_x^22y^2\sin(xy)\dr y$$

    \sol{
        由于$y^2\sin(xy)$对$y$的积分相对复杂，我们需要考虑更换累次积分顺序(事实上本题不交换次序也能分部积分得出结果)。分为如下步骤：
        \begin{itemize}
            \item \textbf{区域描述}
                
            根据条件，区域可描述为满足
            $$0\le x\le2,\quad x\le y\le2$$
            的$(x,y)$集合。

            \item \textbf{累次积分化}
                
            为了化为先$x$后$y$的累次积分，我们需要将集合描述成
            $$\alpha\le y\le\beta,\quad \phi_1(y)\le x\le\phi_2(y)$$
            当$x$可任取时，$y$最小对应$x$为0时，为0；$y$最大一定可取到2。综合得第一个不等式为
            $$0\le y\le2$$
            此时再考虑对$x$的约束可写出第二个不等式
            $$0\le x\le y$$
            最终写出积分
            $$\int_0^2\dr y\int_0^y2y^2\sin(xy)\dr x$$

            \item \textbf{积分计算}
                
            对$x$积分可将累次积分化为
            $$\int_0^22y(1-\cos(y^2))\dr y=4-2\int_0^2y\cos(y^2)\dr y$$
            换元$t=y^2$即得其为
            $$4-\int_0^4\cos t\dr t=4-\sin4$$
        \end{itemize}
    }

    \item (习题7.2.11)计算
    $$\iint_D(|x|+y)\dr\sigma$$
    其中
    $$D=\{(x,y)\mid |x|+|y|\le1\}$$

    \sol{
        分为如下步骤：
        \begin{itemize}
            \item \textbf{反射对称化简}
            
            由于区域$D$满足$(x,y)\in D$当且仅当$(-x,y)\in D$，我们可以按照对称轴$y$轴进行翻折，得到的新区域为
            $$D_0=\{(x,y)\mid|x|+|y|\le1,\ x\ge0\}$$
            新的函数为旧的函数与换元$x$为$-x$的累加，即将积分化为
            $$\iint_{D_0}(|x|+y+|-x|+y)\dr\sigma=2\iint_{D_0}(|x|+y)\dr\sigma$$
            由于$D_0$满足$(x,y)\in D_0$当且仅当$(x,-y)\in D_0$，我们可以按照对称轴$x$轴进行翻折，得到的新区域为
            $$D_1=\{(x,y)\mid|x|+|y|\le1,\ x\ge0,\ y\ge0\}$$
            新的函数为旧的函数与换元$x$为$-x$的累加，即将积分化为
            $$2\iint_{D_1}(|x|+y+|x|+(-y))\dr\sigma=4\iint_{D_1}|x|\dr\sigma$$
            由于$D_1$已保证$x$、$y$为正，可描述成
            $$D_1=\{(x,y)\mid x+y\le1,\ x\ge0,\ y\ge0\}$$
            积分也即
            $$4\iint_{D_1}x\dr\sigma$$

            \note 本段更详细的逻辑可参考本讲义20.3.2。
            
            \item \textbf{累次积分化}
            
            由于被积函数形式简单，我们任取顺序即可，如将它化为先$x$后$y$的累次积分，需要将集合描述成
            $$a\le y\le v,\quad\phi_1(y)\le x\le\phi_2(y)$$
            当$x$任取时$y$的最大范围是
            $$0\le y\le1$$
            而$y$给定时$x$的范围被限制在
            $$0\le x\le1-y$$
            最终写出积分
            $$4\int_0^1\dr y\int_0^{1-y}x\dr x$$

            \item \textbf{积分计算}
            
            对$x$积分可将累次积分化为
            $$4\int_0^1\frac{1}{2}(1-y)^2\dr y=\frac{2}{3}$$
        \end{itemize}
    }

    \item (总练习题7.1(3))改变如下累次积分的积分次序：
    $$\int_0^{3/2}\dr x\int_0^{9-4x^2}16x\dr y$$

    \sol{
        分为如下步骤：
        \begin{itemize}
            \item \textbf{区域描述}
                
            根据条件，区域可描述为满足
            $$0\le x\le\frac{3}{2},\quad 0\le y\le 9-4x^2$$
            的$(x,y)$集合。

            \item \textbf{累次积分化}
                
            为了化为先$x$后$y$的累次积分，我们需要将集合描述成
            $$\alpha\le y\le\beta,\quad \phi_1(y)\le x\le\phi_2(y)$$
            当$x$可任取时，$y$最小一定能取到0；$y$最大对应$x$为0时，为9。综合得第一个不等式为
            $$0\le y\le 9$$
            当$y$给定时，关于$x$的不等式可改写为(由$x$非负可同取平方根)
            $$0\le x\le\frac{3}{2},\quad x\le\frac{\sqrt{9-y}}{2}$$
            进一步估算可发现范围即
            $$0\le x\le\frac{\sqrt{9-y}}{2}$$
            最终写出积分
            $$\int_0^9\dr y\int_0^{\sqrt{9-y}/2}16x\dr x$$
        \end{itemize}
    }

    \item (总练习题7.3(2))描述如下累次积分对应的重积分区域，并计算面积：
    $$\int_{-1}^0\dr x\int_{-2x}^{1-x}\dr y+\int_0^2\dr x\int_{-x/2}^{1-x}\dr y$$

    \sol{
        可以验证$x\in[-1,0]$时$-2x\le1-x$、$x\in[0,2]$时$-\frac{x}{2}\le1-x$，因此两累次积分均可化为重积分，且由于$x$的范围不同，两个积分区域无重叠，可拼成一个完整的部分。这样，整个区域的面积即为对1的重积分，直接计算上方累次积分结果可知面积为$\frac{3}{2}$，下面进行区域的描述。

        由于两区域无重叠，看作重积分的区域为两个区域的并集
        $$\{(x,y)\mid-1\le x\le0,\ -2x\le y\le1-x\}\cup\bigg\{(x,y)\ \bigg|\ 0\le x\le2,\ -\frac{x}{2}\le y\le1-x\bigg\}$$
        由于约束均为线性函数，两区域均为多边形，直接考虑$x$与$y$取左端/右端的情况可以得到第一个区域的端点分别为
        $$x=-1,\quad y=-2x\Longrightarrow(-1,2)$$
        $$x=-1,\quad y=1-x\Longrightarrow(-1,2)$$
        $$x=0,\quad y=-2x\Longrightarrow(0,0)$$
        $$x=0,\quad y=1-x\Longrightarrow(0,1)$$
        从而它是$(0,0)$、$(0,1)$、$(-1,2)$围成的三角形。同理，第二个区域的端点分别为
        $$x=0,\quad y=-\frac{x}{2}\Longrightarrow(0,0)$$
        $$x=0,\quad y=1-x\Longrightarrow(0,1)$$
        $$x=2,\quad y=\frac{x}{2}\Longrightarrow(2,-1)$$
        $$x=2,\quad y=1-x\Longrightarrow(2,-1)$$
        从而它是$(0,0)$、$(0,1)$、$(2,-1)$围成的三角形。

        将两者拼接，即可知最终区域为$(0,0)$、$(-1,2)$、$(0,1)$、$(2,-1)$围成的四边形。进一步观察可发现$(-1,2)$、$(0,1)$、$(2,-1)$都是由$y=1-x$确定的，这条边界实际是平的，由此区域实际上是$(0,0)$、$(-1,2)$、$(2,-1)$围成的三角形。

        \note 本题原题为画出区域，这里提供一种\textbf{只从代数角度}分析多边形区域的方法。
    }
\end{enumerate}

\subsection{重积分的定义}
\subsubsection{矩形区域的重积分}
矩形区域类比定积分定义重积分

可积函数

重积分的几何意义

\subsubsection{累次积分与分离变量}
从定义证明连续函数的矩形区域重积分等于累次积分


\begin{framed}
\begin{exercise}\label{separation}
    考虑定义在$[a,b]$上的连续函数$g(x)$与定义在$[c,d]$上的连续函数$h(y)$，证明
    $$\iint_{[a,b]\times[c,d]}g(x)h(y)\dr x\dr y=\int_a^bg(x)\dr x\int_c^dh(y)\dr y$$
\end{exercise}
\sol{
    利用上学期知识可验证$g(x)h(y)$也为连续函数，从而可写出累次积分
    $$\iint_{[a,b]\times[c,d]}g(x)h(y)\dr x\dr y=\int_a^b\dr x\int_c^dg(x)h(y)\dr y$$
    由于对$y$积分时$g(x)$为常数，利用定积分性质可得上式化为(注意到这步并未做完，无法直接去除括号)
    $$\int_a^b\dr x\bigg(g(x)\int_c^dh(y)\dr y\bigg)$$
    最后，由于$\int_c^dh(y)\dr y$是一个常数，再次利用定积分性质可以将它从对$x$的积分中提取出，最终化成
    $$\int_a^bg(x)\dr x\int_c^dh(y)\dr y$$
    这就是等式的右端。
}
\end{framed}

$f(x,y)=g(x)h(y)$的分离变量情况

\subsubsection{一般区域的重积分}
用矩形包住定义重积分

与教材定义的对比

有直边界情况

\subsubsection{区域的描述}
分割区域进行计算

不同方向描述区域进行计算

脱离直观的代数化理解

\begin{framed}
\begin{exercise}
    改变累次积分
    $$\int_0^1\dr x\int_0^{1-x}\dr y\int_0^{x+y}f(x,y,z)\dr z$$
    的积分次序为：先对$y$、再对$x$、最后对$z$积分。
\end{exercise}
\end{framed}

\subsection{重积分换元法}
\subsubsection{再谈Jacobi行列式}
Jacobi行列式的几何意义

是否非零的要求

局部面积放大倍数

\subsubsection{从局部到整体}
一一对应换元的直观理解

Jacobi行列式与代数角度一一对应

排除部分非一一对应点的合理性

教材定义的意义

换元法成立性的几何理解

\subsubsection{换元的意义}
区域角度决定换元方式

函数角度决定换元方式

\newpage
\section{重积分进阶}
\subsection{习题解答}
\subsubsection{作业解答}
\begin{enumerate}
    \item (习题7.3.5)计算
    $$\iiint_\Omega(x^2-y^2-z^2)\dr V$$
    其中
    $$\Omega=\{(x,y,z)\mid x^2+y^2+z^2\le a^2\}$$
    这里参数$a>0$。

    \sol{
        分为如下步骤：
        \begin{itemize}
            \item \textbf{反射对称化简}
            
            由于区域$\Omega$满足$(x,y,z)\in\Omega$当且仅当$(y,x,z)\in\Omega$，将$x$与$y$交换换元不会改变区域，考虑被积函数的改变可以得到
            $$\iiint_\Omega x^2\dr V=\iiint_\Omega y^2\dr V$$

            \note 这也是\textbf{反射对称性}的一种使用，交换$x$、$y$本质上是关于平面$y=x$对称，详见本讲义20.3.2。

            同理可以得到
            $$\iiint_\Omega x^2\dr V=\iiint_\Omega z^2\dr V$$

            考虑到区域形式，接下来应采用\textbf{球坐标换元}，利用上述结论，我们可以将原积分化为
            $$\iiint_\Omega(x^2-y^2-z^2)\dr V=-\iiint_\Omega z^2\dr V=-\frac{1}{3}\iiint_\Omega(x^2+y^2+z^2)\dr V$$
            使球坐标换元后更加简单。

            \note 即使不进行这一步处理，直接换元计算也可以得到结果，只是计算过程相对更\textbf{复杂}。

            \item \textbf{换元与累次积分化}
            
            由于直接以$x$、$y$、$z$描述区域相对复杂，为了将区域换元为\textbf{长方体}，我们考虑\textbf{球坐标}换元
            $$x=\rho\sin\varphi\cos\theta,\quad y=\rho\sin\varphi\sin\theta,\quad z=\rho\cos\varphi$$
            这样约束条件即化为了$\rho\le a$，结合球坐标\textbf{自然范围}有
            $$0\le\rho\le a,\quad0\le\varphi\le\pi,\quad0\le\theta\le2\pi$$

            \note 注意球坐标的自然范围是$\rho\ge0$、$\theta\in[0,2\pi]$、$\varphi\in[0,\pi]$。

            将上述条件确定的$(\rho,\varphi,\theta)$区域记为$\Omega_0$，被积函数可写为$-\frac{1}{3}\rho^2$，直接计算Jacobi行列式可知
            $$\dr x\dr y\dr z=\rho^2\sin\varphi\dr\rho\dr\varphi\dr\theta$$
            于是原积分化为(按照被积函数中的复杂程度，先对$\theta$再对$\varphi$积分更加简单)
            $$-\frac{1}{3}\iiint_{\Omega_0}\rho^4\sin\varphi\dr\rho=-\frac{1}{3}\int_0^a\dr\rho\int_0^\pi\dr\varphi\int_0^{2\pi}\rho^4\sin\varphi\dr\theta$$

            \item \textbf{积分计算}

            直接按照顺序计算可知积分为
            $$-\frac{2\pi}{3}\int_0^a\dr\rho\int_0^\pi\rho^4\sin\varphi\dr\varphi=-\frac{4\pi}{3}\int_0^a\rho^4\dr\rho=-\frac{4\pi a^5}{15}$$
        \end{itemize}
        
        \note 注意教材原题并无$a>0$的条件，因此原题结果实际上需要加上绝对值。
    }

    \item (习题7.3.10)计算
    $$\iiint_\Omega(1+xy+yz+zx)\dr V$$
    其中$\Omega$是曲面$x^2+y^2=2z$、$x^2+y^2+z^2=8$在$z\ge0$部分围成的区域。

    \sol{
        分为如下步骤：
        \begin{itemize}
            \item \textbf{区域描述}
            
            根据``围成''的有界性要求，球面必然对应内部
            $$x^2+y^2+z^2\le 8$$
            结合条件在$z\ge0$部分，可知必然(``在$z\ge0$部分''应当理解为，$z=0$不为区域的边界，但区域恒满足$z\ge0$)
            $$2z\ge x^2+y^2$$
            由此$\Omega$的代数描述最终为满足如上两个不等式的$(x,y,z)$集合。

            \item \textbf{反射对称化简}
            
            由于区域$\Omega$满足$(x,y,z)\in\Omega$当且仅当$(-x,y,z)\in\Omega$，类似本讲义19.1.2第5题可知
            $$\iiint_\Omega(xy+zx)\dr V=\iiint_\Omega(-xy-zx)\dr V$$
            移项即得
            $$\iiint_\Omega(xy+zx)=0$$
            同理，从$(x,y,z)\in\Omega$当且仅当$(x,-y,z)\in\Omega$可得到
            $$\iiint_\Omega yz\dr V=0$$
            由此可将原积分化为
            $$\iiint_\Omega1\dr V$$
            
            \item \textbf{旋转对称性换元}
            
            利用区域的旋转对称性(详见本讲义20.3.3)，想到进行\textbf{柱坐标}换元
            $$x=r\cos\theta,\quad y=r\sin\theta,\quad z=z$$
            此时，被积函数仍为1，且计算Jacobi行列式得$\dr x\dr y\dr z=r\dr r\dr\theta\dr z$。

            区域的两个条件可写为
            $$r^2+z^2\le8,\quad2z\ge r^2$$
            将上述条件与柱坐标自然范围确定的$(r,\theta,z)$区域记为$\Omega_0$，原积分化为
            $$\iiint_{\Omega_0}r\dr r\dr\theta\dr z$$

            \item \textbf{累次积分化}
            
            由于被积函数对$\theta$与$z$为常数，而约束可以将$z$用$r$约束，因此选择先对$\theta$、再对$z$、最后对$r$积分。此时，我们需要将集合描述成
            $$\alpha\le r\le\beta,\quad \phi_1(r)\le z\le\phi_2(r),\quad\psi_1(r,z)\le\theta\le\psi_2(r,z)$$
            当$z$、$\theta$可任取时($\theta$实际无影响)，$r$最小一定能取到0；$r$最大同时受$\sqrt{8-z^2}$与$\sqrt{2z}$限制，可发现两者相等时能取到最大，也即$z=2$时取到，为2。综合以上讨论得第一个不等式为
            $$0\le r\le 2$$
            当$r$给定时，由$z$非负，关于$z$的不等式可改写为
            $$\frac{r^2}{2}\le z\le\sqrt{8-r^2}$$
            这就是符合要求的约束，实际上与$\theta$无关，而$\theta$的范围即为柱坐标自然范围$[0,2\pi]$。

            综合以上，可最终写出积分
            $$\int_0^2\dr r\int_{r^2/2}^{\sqrt{8-r^2}}\dr z\int_0^{2\pi}r\dr\theta$$

            \item \textbf{积分计算}
            
            对$\theta$积分可将累次积分化为
            $$2\pi\int_0^2\dr r\int_{r^2/2}^{\sqrt{8-r^2}}r\dr z$$
            再对$z$积分也即
            $$2\pi\int_0^2\bigg(\sqrt{8-r^2}-\frac{r^2}{2}\bigg)r\dr r$$
            换元$s=r^2$将积分化为
            $$\pi\int_0^4\bigg(\sqrt{8-s}-\frac{s}{2}\bigg)\dr s$$
            此时已可直接多项式积分结论得到结果为
            $$\pi\bigg(\frac{2}{3}((8-0)^{3/2}-(8-4)^{3/2})-\frac{4^2-0^2}{4}\bigg)=\frac{32\sqrt2-28}{3}\pi$$
        \end{itemize}

        \note 本题的五个步骤是重积分问题的一个\textbf{标准}的处理流程，本质上之后的曲线、曲面积分也是在这五个步骤的基础上进行的。
    }

    \item (习题7.3.16)计算
    $$\iiint_\Omega\frac{\dr V}{\sqrt{x^2+y^2+(z-2)^2}}$$
    其中
    $$\Omega=\{(x,y,z)\mid x^2+y^2+z^2\le1\}$$

    \sol{
        分为如下步骤：
        \begin{itemize}
            \item \textbf{换元与累次积分化}
            
            为了将区域换元为\textbf{长方体}，我们考虑\textbf{球坐标}换元
            $$x=\rho\sin\varphi\cos\theta,\quad y=\rho\sin\varphi\sin\theta,\quad z=\rho\cos\varphi$$
            这样约束条件即化为了$\rho^2\le1$，结合球坐标\textbf{自然范围}有
            $$0\le\rho\le1,\quad0\le\varphi\le\pi,\quad0\le\theta\le2\pi$$

            将上述条件确定的$(\rho,\varphi,\theta)$区域记为$\Omega_0$，被积函数可写为$\frac{1}{\sqrt{\rho^2+4-4\rho\cos\varphi}}$，直接计算Jacobi行列式可知
            $$\dr x\dr y\dr z=\rho^2\sin\varphi\dr\rho\dr\varphi\dr\theta$$
            于是原积分化为(按照被积函数中的复杂程度，先对$\theta$再对$\varphi$积分更加简单)
            $$\iiint_{\Omega_0}\frac{\rho^2\sin\varphi}{\sqrt{\rho^2+4-4\rho\cos\varphi}}\dr\rho=\int_0^1\dr\rho\int_0^\pi\dr\varphi\int_0^{2\pi}\frac{\rho^2\sin\varphi}{\sqrt{\rho^2+4-4\rho\cos\varphi}}\dr\theta$$

            \item \textbf{积分计算}
            
            对$\theta$积分可将累次积分化为
            $$2\pi\int_0^1\dr\rho\int_0^\pi\frac{\rho^2\sin\varphi}{\sqrt{\rho^2+4-4\rho\cos\varphi}}\dr\varphi$$
            换元$t=\cos\varphi$可将积分化为(注意直接换元后积分限为1到$-1$，利用$\dr\cos\varphi=-\sin\varphi\dr\varphi$交换上下限)
            $$2\pi\int_0^1\dr\rho\int_{-1}^1\frac{\rho^2}{\sqrt{\rho^2+4-4\rho t}}\dr t$$
            这对$t$是幂函数复合一次函数，可直接写出原函数
            $$\rho^2\frac{2\sqrt{\rho^2+4-4\rho t}}{-4\rho}+C=-\frac{\rho}{2}\sqrt{\rho^2+4-4\rho t}+C$$
            由此对$t$的积分结果为
            $$\pi\int_0^1\rho(\sqrt{\rho^2+4+4\rho}-\sqrt{\rho^2+4-4\rho})\dr\rho$$
            由$\rho$的范围可知$\rho+2>0$、$2-\rho>0$，从而最终化为
            $$\pi\int_0^1\rho((\rho+2)-(2-\rho))\dr\rho=\frac{2\pi}{3}$$
        \end{itemize}
    }

    \item (习题7.3.22)将累次积分
    $$\int_0^a\dr x\int_0^x\dr y\int_0^yf(z)\dr z$$
    化为定积分。

    \sol{
        \note 注意条件并未给出$a>0$，我们的对它进行重积分的转化是建立在$a>0$的，因此需要分别讨论两种情况。

        我们先假定$a>0$，分为如下步骤：
        \begin{itemize}
            \item \textbf{区域描述}
            
            根据条件，区域可描述为满足
            $$0\le x\le a,\quad0\le y\le x,\quad 0\le z\le y$$
            的$(x,y,z)$集合。

            \item \textbf{累次积分化}
            
            为了将原题化为定积分，$f(z)$必须放在最后进行积分，因此需要考虑$x$、$y$、$z$的积分顺序(由于原区域$x$约束$y$、$y$约束$z$，$x$、$y$、$z$的顺序比$y$、$x$、$z$更加自然)。

            此时，我们需要将集合描述成
            $$\alpha\le z\le\beta,\quad\phi_1(z)\le y\le\phi_2(z),\quad\psi_1(y,z)\le x\le\psi_2(y,z)$$
            当$x$、$y$可任取时，$z$最小一定能取到0；$z$最大对应$x=y=a$时，为$a$。综合得第一个不等式为
            $$0\le z\le a$$
            当$z$给定、$x$可任取时，$y$最小受$y\ge0$与$y\ge z$约束，由$z\ge0$为$z$；$y$最大对应$x=a$时，为$a$。综合得第二个不等式为
            $$z\le y\le a$$
            当$z$、$y$给定时，关于$x$的不等式可改写为
            $$0\le x\le a,\quad y\le x$$
            $x$必须大于等于两个下限中最大的，也即$y$，因此范围即
            $$y\le x\le a$$
            最终写出积分
            $$\int_0^a\dr z\int_z^a\dr y\int_y^af(z)\dr x$$

            \item \textbf{积分计算}
            
            由于$f(z)$对$x$、$y$均为常数，直接计算积分可得结果为
            $$\int_0^a\dr z\int_z^a(a-y)f(z)\dr y=\frac{1}{2}\int_0^a(a-z)^2f(z)\dr z$$
        \end{itemize}

        若$a<0$，利用定积分交换上下限积分反号，可以将原式重写为
        $$-\int_a^0\dr x\int_x^0\dr y\int_y^0f(z)\dr z$$
        由此与上方类似操作可得到相同形式的结论。
    }

    \item (习题7.4.3)计算三个圆柱面$x^2+y^2=R^2$、$y^2+z^2=R^2$、$x^2+z^2=R^2$围成的区域$\Omega$的表面积，这里参数$R>0$。
    
    \sol{
        见\exref{three column s}。
    }

    \item (习题7.4.4)计算三个圆柱面$x^2+y^2=R^2$、$y^2+z^2=R^2$、$x^2+z^2=R^2$围成的区域$\Omega$的体积，这里参数$R>0$。
    
    \sol{
        见\exref{three column v}。
    }

    \item (习题7.4.7)给定参数$a>0$、$b>0$、$c>0$，求如下区域
    $$\Omega=\bigg\{(x,y,z)\ \bigg|\ \frac{x^2}{a^2}+\frac{y^2}{b^2}+\frac{z^2}{c^2}\le1,\ x\ge0,\ y\ge0,\ z\ge0\bigg\}$$
    的质心坐标，假定体密度均匀。

    \sol{
        利用质心坐标公式，设质心坐标为$(x_0,y_0,z_0)$，应满足
        $$x_0=\frac{1}{m}\iiint_\Omega x\dr V,\quad y_0=\frac{1}{m}\iiint_\Omega y\dr V,\quad z_0=\frac{1}{m}\iiint_\Omega z\dr V,\quad m=\iiint_\Omega1\dr V$$

        先计算对$x$的三重积分：
        \begin{itemize}
            \item \textbf{换元与累次积分化}
            
            根据区域本身为球的部分进行伸缩，为了将区域换元为矩形，我们考虑\textbf{广义球坐标}换元
            $$x=a\rho\sin\varphi\cos\theta,\quad y=b\rho\sin\varphi\sin\theta,\quad z=c\rho\cos\varphi$$
            这样约束条件即化为了
            $$\rho^2\le1,\quad a\rho\sin\varphi\cos\theta\ge0,\quad b\rho\sin\varphi\sin\theta\ge0,\quad c\rho\cos\varphi\ge0$$

            结合$\rho$的自然范围$\rho\ge0$有
            $$0\le\rho\le1$$
            由此在$\rho>0$时第三个式子可化为$\cos\varphi\ge0$，结合$\varphi$的自然范围有
            $$\varphi\in\bigg[0,\frac{\pi}{2}\bigg]$$
            再结合$\theta$的自然范围得到
            $$\theta\in\bigg[0,\frac{\pi}{2}\bigg]$$

            将上述条件确定的$(\rho,\varphi,\theta)$区域记为$\Omega_0$，被积函数可写为$a\rho\sin\varphi\cos\theta$。 

            直接类似球坐标计算Jacobi行列式可知
            $$\dr x\dr y\dr z=abc\rho^2\sin\varphi\dr\rho\dr\varphi\dr\theta$$
            于是原积分化为(这里积分顺序任取，事实上无影响)
            $$\int_0^{\pi/2}\dr\varphi\int_0^{\pi/2}\dr\theta\int_0^1a^2bc\rho^3\sin^2\varphi\cos\theta\dr\rho$$

            \note 这里说明换元合理的理由类似球坐标(即使Jacobi行列式在$z$轴上为0)，进行范围确定时本质上也忽略了$z$轴上的情况，如确定$\theta$的范围时实际上未考虑$\rho=0$或$\varphi=0$时$\theta$可任取。此过程的合理性详见本讲义19.3.2。

            \item \textbf{积分计算}
            
            直接依次积分可得结果为(具体计算可使用上册教材3.4节例2)
            $$\frac{a^2bc}{4}\int_0^{\pi/2}\dr\varphi\int_0^{\pi/2}\sin^2\varphi\cos\theta\dr\theta=\frac{a^2bc}{4}\int_0^{\pi/2}\sin^2\varphi\dr\varphi=\frac{\pi a^2bc}{16}$$
        \end{itemize}
        同理，用相同的换元可算出
        $$\iiint_\Omega y\dr V=\frac{\pi ab^2c}{16},\quad\iiint_\Omega z\dr V=\frac{\pi abc^2}{16},\quad\iiint_\Omega1\dr V=\frac{\pi abc}{6}$$
        由此可最终得出质心坐标
        $$\bigg(\frac{3a}{8},\frac{3b}{8},\frac{3c}{8}\bigg)$$
    }

    \item (习题7.4.11)设引力常数为$k$，求密度为$\rho$的均匀圆柱体
    $$\Omega=\{(x,y,z)\mid x^2+y^2\le a^2,\ 0\le z\le b\}$$
    对位于$(0,0,h)$处的单位质点的引力，这里参数$h>b>0$、$a>0$、$\rho>0$。

    \sol{
        利用引力公式，设引力为$(F_x,F_y,F_z)$，应满足
        $$F_x=k\iiint_\Omega\frac{x}{(x^2+y^2+(z-h)^2)^{3/2}}\rho\dr V$$
        $$F_y=k\iiint_\Omega\frac{y}{(x^2+y^2+(z-h)^2)^{3/2}}\rho\dr V$$
        $$F_z=k\iiint_\Omega\frac{z-h}{(x^2+y^2+(z-h)^2)^{3/2}}\rho\dr V$$
        由于$(x,y,z)\in\Omega$当且仅当$(-x,y,z)\in\Omega$，类似本部分第2题的分析可得$F_x=0$，同理$F_y=0$，只需计算$F_z$。

        分为如下步骤：
        \begin{itemize}
            \item \textbf{换元与累次积分化}
            
            由于区域为圆柱，为了将区域换元为长方体，我们考虑柱坐标换元
            $$x=r\cos\theta,\quad y=r\sin\theta,\quad z=z$$
            这样约束条件即化为了
            $$r^2\le a^2,\quad 0\le z\le b$$
            结合柱坐标的自然范围可得
            $$0\le r\le a,\quad 0\le\theta\le2\pi,\quad 0\le z\le b$$
            将上述条件确定的$(r,\theta,z)$区域记为$\Omega_0$，被积函数可写为
            $$\frac{z-h}{(r^2+(z-h)^2)^{3/2}}$$
            直接计算Jacobi行列式可知
            $$\dr x\dr y\dr z=r\dr r\dr\theta\dr z$$
            于是(这里对$\theta$为常数因此先积分，对$z$、$r$积分复杂度类似，先后区别不大)
            $$F_z=k\rho\iiint_{\Omega_0}\frac{(z-h)r}{(r^2+(z-h)^2)^{3/2}}\dr r\dr\theta\dr z=k\rho\int_0^a\dr r\int_0^b\dr z\int_0^{2\pi}\frac{r(z-h)}{(r^2+(z-h)^2)^{3/2}}\dr\theta$$
            
            \item \textbf{积分计算}
            
            对$\theta$积分可知
            $$F_z=2\pi k\rho\int_0^a\dr r\int_0^b\frac{r(z-h)}{(r^2+(z-h)^2)^{3/2}}\dr z$$
            根据形式，换元$s=(z-h)^2$有(注意条件$h>b>0$)
            $$F_z=\pi k\rho\int_0^a\dr r\int_{h^2}^{(b-h)^2}\frac{r\dr s}{(r^2+s)^{3/2}}$$
            此为$s$的幂函数复合一次函数，可直接写出被积函数原函数
            $$\frac{-2r}{\sqrt{r^2+s}}+C$$
            由此代入算得对$s$积分为
            $$F_z=2\pi k\rho\int_0^a\bigg(\frac{r}{\sqrt{r^2+h^2}}-\frac{r}{\sqrt{r^2+(h-b)^2}}\bigg)\dr r$$
            根据形式，换元$t=r^2$有
            $$F_z=\pi k\rho\int_0^{a^2}\bigg(\frac{1}{\sqrt{t+h^2}}-\frac{1}{\sqrt{t+(h-b)^2}}\bigg)\dr s$$
            由于两部分都是$t$的幂函数复合一次函数，类似写出原函数并计算可知
            $$F_z=2\pi k\rho\big((\sqrt{a^2+h^2}-h)-(\sqrt{a^2+(h-b)^2}-(h-b))\big)$$
            化简可得
            $$F_z=2\pi k\rho(\sqrt{a^2+h^2}-\sqrt{a^2+(h-b)^2}-b)$$
        \end{itemize}
        综合$F_x$、$F_y$、$F_z$可得结果。
    }

    \item (总练习题7.12)改变如下累次积分的积分次序：
    $$\int_{-1}^1\dr x\int_{x^2}^1\dr y\int_0^{1-y}\dr z$$

    \begin{enumerate}[(1)]
        \item 改变为按$y$、$z$、$x$的顺序积分。
        
        \sol{
            分为如下步骤：
            \begin{itemize}
                \item \textbf{区域描述}
                
                根据条件，区域可描述为满足
                $$-1\le x\le1,\quad x^2\le y\le1,\quad 0\le z\le1-y$$
                的$(x,y,z)$集合。

                \item \textbf{累次积分化}
                
                为了化为$y$、$z$、$x$的累次积分，我们需要将集合描述成
                $$\alpha\le x\le\beta,\quad\phi_1(x)\le z\le\phi_2(x),\quad\psi_1(x,z)\le y\le\psi_2(x,z)$$
                原本的$x$的范围已经满足第一个不等式要求。当$x$给定、$y$可任取时，$z$最小一定可取到0；$z$最大要求$y$尽量小，对应$y=x^2$时，为$1-x^2$。综合得第二个不等式为
                $$0\le z\le1-x^2$$
                当$x$、$z$给定时，关于$y$的不等式可改写为
                $$x^2\le y\le1,\quad y\le1-z$$
                $y$必须小于等于两个上限中最小的，也即$1-z$，因此范围即
                $$x^2\le y\le1-z$$
                最终写出积分
                $$\int_{-1}^1\dr x\int_0^{1-x^2}\dr z\int_{x^2}^{1-z}\dr y$$
            \end{itemize}
        }

        \item 改变为按$x$、$z$、$y$的顺序积分。
        
        \sol{
            区域描述的步骤与前一问相同，我们只需进行累次积分化过程。
                
            为了化为$x$、$z$、$y$的累次积分，我们需要将集合描述成
            $$\alpha\le y\le\beta,\quad\phi_1(y)\le z\le\phi_2(y),\quad\psi_1(y,z)\le x\le\psi_2(y,z)$$
            当$x$、$z$可任取时，$y$最小对应$x=0$时，为0；$y$最大对应$z=0$时，为1。综合得第一个不等式为
            $$0\le y\le1$$

            当$y$给定、$x$可任取时，$z$最小一定可取到0；$z$最大即为$1-y$。综合得第二个不等式为
            $$0\le z\le1-y$$
            当$y$、$z$给定时，关于$x$的不等式可改写为
            $$-1\le x\le1,\quad x^2\le y$$
            由于$y\in[0,1]$，两者确定的实际$x$的范围是
            $$-\sqrt{y}\le x\le\sqrt{y}$$
            最终写出积分
            $$\int_0^1\dr y\int_0^{1-y}\dr z\int_{-\sqrt{y}}^{\sqrt{y}}\dr x$$
        }
    \end{enumerate}

    \item (总练习题7.15(3))将柱坐标表示下的三重积分
    $$\iiint_\Omega f(r,\theta,z)r\dr r\dr\theta\dr z$$
    化为柱坐标下的累次积分，其中区域$\Omega$对应的直角坐标区域由平面$y=0$、$y=x$、$x=1$、$z=0$、$z=2-y$围成。

    \sol{
        分为如下步骤：
        \begin{itemize}
            \item \textbf{区域描述}
            
            根据``围成''的有界性要求与围成的区域必然以每个平面为边界，我们逐个进行分析。
            
            若$y=0$对应的约束为$y\le0$，则必然有$y\ge x$\ (否则$y=x$不再为边界)，但这可以得到$x\le0$，与$x=1$为边界矛盾。由此，对应的约束必然为$y\ge0$，此时类似可得第二个约束必然为$y\le x$、第三个约束必然为$x\le1$。

            同理，后两个约束为$0\le z\le 2-y$，这就得到了完整的区域描述，即$\Omega$是满足
            $$0\le y\le x\le1,\quad0\le z\le 2-y$$
            的$(x,y,z)$构成的集合。

            代入柱坐标即得$\Omega$对应$(r,\theta,z)$的约束
            $$0\le r\sin\theta\le r\cos\theta\le1,\quad0\le z\le2-r\sin\theta$$

            \item \textbf{累次积分化}
            
            为了将三重积分化为累次积分，我们先尝试化简约束。

            第一个约束中可提取出$\theta$的要求(忽略$r=0$的特殊点)
            $$0\le\sin\theta\le\cos\theta$$
            结合柱坐标$\theta$的自然范围，可得必然
            $$\theta\in\bigg[0,\frac{\pi}{4}\bigg]$$
            这时结合$r$的自然范围即可得到第一个约束确定的$r$的范围
            $$0\le r\le\frac{1}{\cos\theta}$$
            由于$r\sin\theta\le1$时，$2-r\sin\theta\ge0$一定成立，因此对满足如上两条件的$\theta$、$r$，一定存在$z$满足第二个约束，无需进一步缩减范围。综合以上讨论，我们得到了区域的等价表述
            $$0\le\theta\le\frac{\pi}{4},\quad0\le r\le\frac{1}{\cos\theta},\quad 0\le z\le2-r\sin\theta$$
            这已经是符合累次积分要求的范围，从而可最终写出积分
            $$\int_0^{\pi/4}\dr\theta\int_0^{1/\cos\theta}\dr r\int_0^{2-r\sin\theta}f(r,\theta,z)r\dr z$$
        \end{itemize}
    }

    \item (总练习题7.18)考虑柱坐标下的累次积分
    $$\int_0^{2\pi}\dr\theta\int_0^{\sqrt2}\dr r\int_r^{\sqrt{4-r^2}}3r\dr z$$
    \begin{enumerate}[(1)]
        \item 将它化作直角坐标下的累次积分。
        
        \sol{
            分为如下步骤：
            \begin{itemize}
                \item \textbf{区域描述}
                
                根据条件，区域可描述为满足
                $$0\le\theta\le2\pi,\quad 0\le r\le\sqrt2,\quad r\le z\le\sqrt{4-r^2}$$
                的$(r,\theta,z)$集合。

                \item \textbf{换元回直角坐标}
                
                转换为直角坐标后，利用$r=\sqrt{x^2+y^2}$，$\theta\in[0,2\pi]$为柱坐标自然约束，约束条件即化为了
                $$0\le\sqrt{x^2+y^2}\le\sqrt2,\quad\sqrt{x^2+y^2}\le z\le\sqrt{4-x^2-y^2}$$
                将上述条件确定的$(x,y,z)$区域记为$\Omega_0$，
                直接计算Jacobi行列式可知
                $$\dr x\dr y\dr z=r\dr r\dr\theta\dr z$$
                由此直角坐标下的被积函数即为3，原积分化为
                $$\iiint_{\Omega_0}3\dr V$$

                \item \textbf{累次积分化}

                我们可以化简$\Omega_0$的约束条件为
                $$x^2+y^2\le2,\quad\sqrt{x^2+y^2}\le z\le\sqrt{4-x^2-y^2}$$
                由于第二个约束条件自然表示了$z$被$x$、$y$约束，考虑先对$z$、再对$y$、最后对$x$积分，前两个约束即需要写成
                $$\alpha\le x\le\beta,\quad\phi_1(x)\le y\le\phi_2(x)$$
                当$y$可任取时，$x$最小对应$y=0$时，为$-\sqrt2$；$x$最大对应$y=0$时，为$\sqrt2$。综合得第一个不等式为
                $$-\sqrt2\le x\le\sqrt2$$
                当$x$给定时，关于$y$的不等式即可改写为
                $$-\sqrt{2-x^2}\le y\le\sqrt{2-x^2}$$
                由于$x^2+y^2\le2$时，$\sqrt{x^2+y^2}\le\sqrt{4-x^2-y^2}$一定成立，因此对满足如上两不等式的$x$、$y$，一定存在$z$满足第二个约束，无需进一步缩减范围。综合以上讨论，我们可以写出累次积分
                $$\int_{-\sqrt2}^{\sqrt2}\dr x\int_{-\sqrt{2-x^2}}^{\sqrt{2-x^2}}\dr y\int_{\sqrt{x^2+y^2}}^{\sqrt{4-x^2-y^2}}3\dr z$$
                
            \end{itemize}
        }
        
        \item 将它化为球坐标下的累次积分。
        
        \sol{
            分为如下步骤：
            \begin{itemize}
                \item \textbf{区域描述}
                
                我们采用上一问的直角坐标描述区域为满足
                $$x^2+y^2\le2,\quad\sqrt{x^2+y^2}\le z\le\sqrt{4-x^2-y^2}$$
                的$(x,y,z)$集合。

                \item \textbf{旋转对称性换元}
                
                考虑球坐标换元
                $$x=\rho\sin\varphi\cos\theta,\quad y=\rho\sin\varphi\sin\theta,\quad z=\rho\cos\varphi$$
                此时，被积函数仍为3，且计算Jacobi行列式得$\dr x\dr y\dr z=r^2\sin\varphi\dr\rho\dr\varphi\dr\theta$。

                区域的两个条件可写为(注意自然范围保证了$\sin\varphi\ge0$)
                $$\rho^2\sin^2\varphi\le2,\quad\rho\sin\varphi\le\rho\cos\varphi\le\sqrt{4-\rho^2\sin^2\varphi}$$
                将上述条件与球坐标自然范围确定的$(\rho,\varphi,\theta)$区域记为$\Omega_1$，原积分化为
                $$\iiint_{\Omega_1}3\rho^2\sin\varphi\dr\rho\dr\varphi\dr\theta$$
                
                \item \textbf{累次积分化}
                
                为了将三重积分化为累次积分，我们先尝试化简约束。

                由$\varphi$的自然范围保证了$\sin\varphi\ge0$，约束$\rho\sin\varphi\le\rho\cos\varphi$可以化简为(忽略$z$轴上的特殊情况)
                $$0\le\sin\varphi\le\cos\varphi$$
                也即
                $$\varphi\in\bigg[0,\frac{\pi}{4}\bigg]$$
                此时，由于$\cos\varphi$非负，同平方即得
                $$\rho^2\cos^2\varphi\le 4-\rho^2\sin^2\varphi$$
                结合自然范围也即得到
                $$0\le\rho\le2$$
                此时，进一步计算可发现$\rho^2\sin^2\varphi\le2$必然满足，结合$\theta$的自然范围最终得到约束
                $$0\le\varphi\le\frac{\pi}{4},\quad0\le\rho\le2,\quad0\le\theta\le2\pi$$
                从而可以写出累次积分(这里积分顺序任取，事实上无影响)
                $$\int_0^{\pi/4}\dr\varphi\int_0^2\dr\rho\int_0^{2\pi}3\rho^2\sin\varphi\dr\theta$$
            \end{itemize}
        }

        \item 计算它的结果。
        
        \sol{
            利用上一问的累次积分表示，对$\theta$积分得到其为
            $$6\pi\int_0^{\pi/4}\dr\varphi\int_0^2\rho^2\sin\varphi\dr\rho$$
            再对$\rho$、$\varphi$积分即得到其为
            $$16\pi\int_0^{\pi/4}\sin\varphi\dr\varphi=(16-8\sqrt2)\pi$$
        }
    \end{enumerate}

    \item (总练习题7.20)考虑柱坐标下的累次积分
    $$\int_0^{\pi/2}\dr\theta\int_1^{\sqrt3}\dr r\int_1^{\sqrt{4-r^2}}r^3\sin\theta\cos\theta\cdot z^2\dr z$$
    将它化作直角坐标下的累次积分。

    \sol{
        分为如下步骤：
        \begin{itemize}
            \item \textbf{区域描述}
                
            根据条件，区域可描述为满足
            $$0\le\theta\le\frac{\pi}{2},\quad 1\le r\le\sqrt3,\quad 1\le z\le\sqrt{4-r^2}$$
            的$(r,\theta,z)$集合。

            \item \textbf{换元回直角坐标}
                
            转换为直角坐标后，利用$r=\sqrt{x^2+y^2}$，后两个约束条件即化为了
            $$1\le\sqrt{x^2+y^2}\le\sqrt3,\quad1\le z\le\sqrt{4-x^2-y^2}$$
            第一个约束条件对应第一象限，因此可转化为
            $$x\ge0,\quad y\ge0$$

            将上述条件确定的$(x,y,z)$区域记为$\Omega_0$，
            直接计算Jacobi行列式可知
            $$\dr x\dr y\dr z=r\dr r\dr\theta\dr z$$
            由此直角坐标下的被积函数即为$r^2\sin\theta\cos\theta z^2$，再利用$x=r\cos\theta$、$y=r\sin\theta$即能将原积分化为
            $$\iiint_{\Omega_0}xyz^2\dr V$$

            \item \textbf{累次积分化}

            我们可以化简$\Omega_0$的约束条件为
            $$1\le x^2+y^2\le3,\quad1\le z\le\sqrt{4-x^2-y^2},\quad x\ge0,\quad y\ge0$$
            由于第二个约束条件自然表示了$z$被$x$、$y$约束，考虑先对$z$、再对$y$、最后对$x$积分，前两个约束即需要写成
            $$\alpha\le x\le\beta,\quad\phi_1(x)\le y\le\phi_2(x)$$
            当$y$可任取时，$x$最小对应$y=1$时，为0；$x$最大对应$y=0$时，为$\sqrt3$。综合得第一个不等式为
            $$0\le x\le\sqrt3$$
            当$x$给定时，由于已知$x$、$y$为正，关于$y$的不等式即可改写为
            $$y^2\ge1-x^2,\quad y\le\sqrt{3-x^2},\quad y\ge0$$
            从中间的式子可看出$y$的上限一定可取到$\sqrt{3-x^2}$；当$x\in[1,\sqrt3]$时，$1-x^2\le0$，第一个条件自动满足，对应下限为0；当$x\in[0,1]$时，第一个条件对应的下限比0大，实际下限为$\sqrt{1-x^2}$。由此，令
            $$\phi(x)=\begin{cases}\sqrt{1-x^2}&x\in[0,1]\\0&x\in[1,3]\end{cases}$$
            则$x$、$y$的约束可最终写为
            $$0\le x\le\sqrt3,\quad\phi(x)\le y\le\sqrt{3-x^2}$$

            由于$1\le x^2+y^2\le3$时，$1\le\sqrt{4-x^2-y^2}$一定成立，因此对满足如上两不等式的$x$、$y$，一定存在$z$满足第二个约束，无需进一步缩减范围。综合以上讨论，我们可以写出累次积分
            $$\int_0^{\sqrt3}\dr x\int_{\phi(x)}^{\sqrt{3-x^2}}\dr y\int_1^{\sqrt{4-x^2-y^2}}xyz^2\dr z$$
            或讨论$\phi(x)$的两种情况而拆分为
            $$\int_0^1\dr x\int_{\sqrt{1-x^2}}^{\sqrt{3-x^2}}\dr y\int_1^{\sqrt{4-x^2-y^2}}xyz^2\dr z+\int_1^{\sqrt3}\dr x\int_0^{\sqrt{3-x^2}}\dr y\int_1^{\sqrt{4-x^2-y^2}}xyz^2\dr z$$
        \end{itemize}
    }
\end{enumerate}
\subsubsection{补充题解答}
\begin{enumerate}
    \item (习题7.2.13)利用极坐标计算
    $$\int_0^1\dr x\int_0^{\sqrt{1-x^2}}(x^2+y^2)\dr y$$

    \sol{
        分为如下步骤：
        \begin{itemize}
            \item \textbf{区域描述}
                
            根据条件，区域可描述为满足
            $$0\le x\le1,\quad 0\le y\le\sqrt{1-x^2}$$
            的$(x,y)$集合。

            \item \textbf{旋转对称性换元}
            
            利用旋转对称性，进行极坐标换元$x=r\cos\theta$、$y=r\sin\theta$。

            此时，原本的被积函数可写为$r^2$，且计算Jacobi行列式得$\dr x\dr y=r\dr r\dr\theta$。区域的两个条件可写为
            $$0\le r\cos\theta\le1,\quad0\le r\sin\theta\le\sqrt{1-r^2\cos\theta}$$
            将上述条件与极坐标自然范围确定的$(r,\theta)$区域记为$D_0$，原积分化为
            $$\iint_{D_0}r^3\dr r\dr\theta$$

            \item \textbf{累次积分化}
                
            由于被积函数对于$\theta$为常数，我们希望化为先$\theta$后$r$的累次积分，也即需要将集合描述成
            $$\alpha\le r\le\beta,\quad\phi_1(r)\le\theta\le\phi_2(r)$$

            将第二个约束两边同平方，由于已知$r\sin\theta\ge0$，同平方后不等号不变，也即得到
            $$r^2\sin^2\theta\le1-r^2\cos\theta$$
            移项后结合极坐标自然范围有$0\le r\le1$。此时，进一步化简约束条件可发现只剩$\cos\theta\ge0$、$\sin\theta\ge0$，从而结合极坐标自然范围有
            $$0\le\theta\le\frac{\pi}{2}$$
            这已经是符合要求的范围，从而可最终写出积分
            $$\int_0^1\dr r\int_0^{\pi/2}r^3\dr\theta$$

            \item \textbf{积分计算}
                
            对$\theta$积分可将累次积分化为
            $$\frac{\pi}{2}\int_0^1r^3\dr r=\frac{\pi}{8}$$
        \end{itemize}
    }

    \item (习题7.2.17)利用极坐标计算
    $$\iint_D\frac{1}{x^2}\dr x\dr y$$
    其中$D$是由曲线$y=\alpha x$、$y=\beta x$、$x^2+y^2=a^2$、$x^2+y^2=b^2$在第一象限围成的区域，这里参数$\beta>\alpha>0$、$b>a>0$。

    \sol{
        分为如下步骤：
        \begin{itemize}
            \item \textbf{区域描述}
                
            根据``围成''的有界性要求，类似本讲义20.1.1第10题可得到区域为满足
            $$\alpha x\le y\le\beta x,\quad a^2\le x^2+y^2\le b^2$$
            的$(x,y)$集合。

            \item \textbf{旋转对称性换元}
            
            利用旋转对称性，进行极坐标换元$x=r\cos\theta$、$y=r\sin\theta$。

            此时，原本的被积函数可写为$\frac{1}{r^2\cos^2\theta}$，且计算Jacobi行列式得$\dr x\dr y=r\dr r\dr\theta$。区域的两个条件可写为
            $$r\alpha\cos\theta\le r\sin\theta\le r\beta\cos\theta,\quad a^2\le r^2\le b^2$$
            忽略$r=0$的特殊点，利用$r$、$a$、$b$均为正可进一步化简为
            $$\alpha\cos\theta\le\sin\theta\le\beta\cos\theta,\quad a\le r\le b$$

            将上述条件与极坐标自然范围确定的$(r,\theta)$区域记为$D_0$，原积分化为
            $$\iint_{D_0}\frac{1}{r\cos^2\theta}\dr r\dr\theta$$

            \item \textbf{累次积分化}
                
            为了化为累次积分，我们需要进一步化简约束。若$\cos\theta=0$，由第一个约束可推出$\sin\theta=0$，矛盾，从而其必然非零，又由区域在第一象限中得其为正，第一个约束同除以$\cos\theta$得
            $$\alpha\le\tan\theta\le\beta$$
            根据第一象限条件可知$\theta\in[0,\frac{\pi}{2}]$，再结合$0<\alpha<\beta$与$\tan$的单调性即得
            $$\arctan\alpha\le\theta\le\arctan\beta$$

            这结合$a\le r\le b$已经是符合要求的范围，从而可最终写出积分
            $$\int_a^b\dr r\int_{\arctan\alpha}^{\arctan\beta}\frac{1}{r\cos^2\theta}\dr\theta$$

            \item \textbf{积分计算}
                
            由于$\frac{1}{\cos^2\theta}$的一个原函数为$\tan\theta$，对$\theta$积分可将累次积分化为
            $$(\beta-\alpha)\int_a^b\frac{\dr r}{r}=(\beta-\alpha)\ln\frac{b}{a}$$
        \end{itemize}
    }

    \item (习题7.2.19)利用重积分的几何意义证明：极坐标下的射线$\theta=\alpha$、$\theta=\beta$与曲线$r=f(\theta)$围成区域的面积可以表示为
    $$\frac{1}{2}\int_\alpha^\beta f^2(\theta)\dr\theta$$
    这里参数$\alpha\le\beta$，$f(\theta)$在$[\alpha,\beta]$有定义且非负连续。

    \sol{
        根据重积分几何意义，设区域为$D$，其面积即为
        $$\iint_D1\dr\sigma$$
        计算分为如下步骤：
        \begin{itemize}
            \item \textbf{换元与累次积分化}
            
            利用旋转对称性，进行极坐标换元$x=r\cos\theta$、$y=r\sin\theta$。

            此时，被积函数仍为1，计算Jacobi行列式得$\dr x\dr y=r\dr r\dr\theta$。根据``围成''的要求与极坐标自然条件，区域的两个条件可写为
            $$\alpha\le\theta\le\beta,\quad0\le r\le f(\theta)$$
            将上述条件确定的$(r,\theta)$区域记为$D_0$，此描述已经符合累次积分的要求，从而可将原积分写为
            $$\iint_{D_0}r\dr r\dr\theta=\int_\alpha^\beta\int_0^{f(\theta)}r\dr r$$

            \item \textbf{积分计算}
            
            对$r$积分可将累次积分化为
            $$\int_\alpha^\beta\frac{1}{2}f^2(\theta)\dr\theta$$
            这就是所求的结果。
        \end{itemize}
    }

    \item (习题7.2.21)计算
    $$\iint_D(2x^2-xy-y^2)\dr x\dr y$$
    其中区域$D$由直线$y=-2x+4$、$y=-2x+7$、$y=x-2$、$y=x+1$围成。

    \sol{
        分为如下步骤：
        \begin{itemize}
            \item \textbf{区域描述}
            
            根据``围成''的有界性要求与围成的区域必然以每个平面为边界，由于$-2x+4<-2x+7$、$x-2<x+1$，范围只能为
            $$-2x+4\le y\le-2x+7,\quad x-2\le y\le x+1$$
            的$(x,y)$集合。

            \item \textbf{换元与累次积分化}
            
            由于直接以$x$、$y$描述区域相对复杂，为了将区域换元为矩形，我们进行换元
            $$\xi=2x+y,\quad\eta=y-x$$
            这样区域即成为
            $$4\le\xi\le7,\quad-2\le\eta\le1$$
            将上述条件确定的$(\xi,\eta)$区域记为$D_0$，反解出
            $$x=\frac{1}{3}(\xi-\eta),\quad y=\frac{1}{3}(\xi+2\eta)$$
            从而被积函数可展开化简为$-\xi\eta$，直接计算Jacobi行列式可知
            $$\begin{vmatrix}\frac{1}{3}&\frac{1}{3}\\-\frac{1}{3}&\frac{2}{3}\end{vmatrix}=\frac{1}{3}$$
            由于其在$D_0$中恒正，换元是合理的，且
            $$\dr x\dr y=\frac{1}{3}\dr\xi\dr\eta$$
            于是原积分化为
            $$\int_{-2}^1\dr\eta\int_4^7-\frac{\xi\eta}{3}\dr\xi$$

            \item \textbf{积分计算}
                
            对$\xi$积分可将累次积分化为
            $$\int_{-2}^1\frac{11}{2}(-\eta)\dr\eta=\frac{33}{4}$$
        \end{itemize}
    }

    \item (习题7.2.26)设$a>0$，并考虑
    $$I(a)=\int_0^a\er^{-x^2}\dr x,\quad J(a)=\iint_{D_a}\er^{-x^2-y^2}\dr x\dr y$$
    其中
    $$D_a=\{(x,y)\mid x^2+y^2\le a^2,\ x\ge0,\ y\ge0\}$$
    \begin{enumerate}[(1)]
        \item 证明
        $$I^2(a)=\iint_{R_a}\er^{-x^2-y^2}\dr x\dr y$$
        其中
        $$R_a=\{(x,y)\mid 0\le x\le a,\ 0\le y\le a\}$$

        \sol{
            见\exref{gauss}第一部分证明。
        }

        \item 证明$J(a)\le I^2(a)\le J(\sqrt2 a)$。
        
        \sol{
            见\exref{gauss}第二部分证明。
        }

        \item 证明
        $$\lim_{a\to+\infty}\int_0^a\er^{-x^2}\dr x=\frac{\sqrt\pi}{2}$$
        此结果可以记作
        $$\int_0^{+\infty}\er^{-x^2}\dr x=\frac{\sqrt\pi}{2}$$

        \sol{
            见\exref{gauss}第三部分证明。
        }
    \end{enumerate}

    \item (总练习题7.9(1))将如下累次积分化为极坐标下的累次积分并计算：
    $$\int_0^1\dr y\int_0^{\sqrt{1-y^2}}(x^2+y^2)\dr x$$

    \sol{
        分为如下步骤：
        \begin{itemize}
            \item \textbf{区域描述}
                
            根据条件，区域可描述为满足
            $$0\le y\le1,\quad 0\le x\le\sqrt{1-y^2}$$
            的$(x,y)$集合。

            \item \textbf{换元与累次积分化}
            
            利用旋转对称性，进行极坐标换元$x=r\cos\theta$、$y=r\sin\theta$。

            此时，原本的被积函数可写为$r^2$，且计算Jacobi行列式得$\dr x\dr y=r\dr r\dr\theta$。区域的两个条件可写为
            $$0\le r\sin\theta\le1,\quad 0\le r\cos\theta\le\sqrt{1-r^2\sin^2\theta}$$
            忽略$r=0$的特殊点，由$r\sin\theta$与$r\cos\theta$均非负，结合极坐标自然范围得
            $$0\le\theta\le\frac{\pi}{2}$$
            将第二个约束两边同平方，由于已知$r\sin\theta\le0$，同平方后不等号反号，也即得到
            $$1-r^2\cos\theta\ge r^2\sin^2\theta$$
            移项后结合极坐标自然范围有$0\le r\le1$。进一步代入其他约束可以验证
            $$0\le\theta\le\frac{\pi}{2},\quad 0\le r\le1$$
            是原区域的等价描述。将上述条件确定的$(r,\theta)$区域记为$D_0$，即可以写出累次积分(这里积分顺序任取，事实上无影响)

            $$\iint_{D_0}r^3\dr r\dr\theta=\int_0^{\pi/2}\dr\theta\int_0^1r^3\dr r$$

            \item \textbf{积分计算}
                
            直接对$r$积分可将累次积分化为
            $$\int_0^{\pi/2}\frac{1}{4}\dr\theta=\frac{\pi}{8}$$
        \end{itemize}
    }

    \item (习题7.3.2)计算
    $$\iiint_\Omega x^2y^2z\dr V$$
    其中区域$\Omega$由曲面$2z=x^2+y^2$、$z=2$围成。

    \sol{
        分为如下步骤：
        \begin{itemize}
            \item \textbf{区域描述}
            
            根据``围成''的有界性要求，若$z=2$对应的约束为$z\ge2$，无论要求$2z\le x^2+y^2$还是$2z\ge x^2+y^2$，$x$、$y$均可趋于无穷，由此第二个约束必然为$z\le2$。进一步讨论第一个约束即可得到区域为满足
            $$\frac{1}{2}(x^2+y^2)\le z\le 2$$
            的$(x,y,z)$集合。

            \item \textbf{旋转对称性换元}
            
            利用区域关于$z$轴的旋转对称性(详见本讲义20.3.3)，我们考虑柱坐标换元
            $$x=r\cos\theta,\quad y=r\sin\theta,\quad z=z$$
            这样约束条件即化为了
            $$\frac{1}{2}r^2\le z\le 2$$
            此时被积函数可写为
            $$zr^4\cos^2\theta\sin^2\theta$$
            且直接计算Jacobi行列式可知
            $$\dr x\dr y\dr z=r\dr r\dr\theta\dr z$$
            我们将约束条件与柱坐标自然范围确定的$(r,\theta,z)$区域记为$\Omega_0$，原积分化为
            $$\iiint_{\Omega_0}zr^5\cos^2\theta\sin^2\theta\dr r\dr\theta\dr z$$

            \item \textbf{累次积分化}
            
            由于约束中$z$被$r$自然限制，且不涉及$\theta$，可以选择先对$\theta$、再对$z$、最后对$r$积分。此时，我们需要将集合描述成
            $$\alpha\le r\le\beta,\quad \phi_1(r)\le z\le\phi_2(r),\quad\psi_1(r,z)\le\theta\le\psi_2(r,z)$$
            当$z$、$\theta$可任取时($\theta$实际无影响)，$r$最小一定能取到0；$r$最大受$z$限制，可发现$z=2$时能取到最大，为2。综合以上讨论得第一个不等式为
            $$0\le r\le 2$$
            当$r$给定、$\theta$任取时，$z$的范围限制即为之前的
            $$\frac{1}{2}r^2\le z\le 2$$
            最后，由于$\theta$不受到约束，其范围即为自然范围$[0,2\pi]$。

            综合以上，可最终写出积分
            $$\int_0^2\dr r\int_{r^2/2}^2\dr z\int_0^{2\pi}zr^5\cos^2\theta\sin^2\theta\dr\theta$$

            \item \textbf{积分计算}
            
            对于$\theta$的积分实际上也即只需计算(换元$t=2\theta$)
            $$\int_0^{2\pi}\cos^2\theta\sin^2\theta\dr\theta=\frac{1}{4}\int_0^{2\pi}\sin^2(2\theta)\dr\theta=\frac{1}{8}\int_0^{4\pi}\sin^2t\dr t$$
            类似本讲义19.1.1第4题进行拆分可将上述积分化为$\sin^2t$在0到$\frac{\pi}{2}$的积分，再由上册教材3.4节例2可计算出结果为$\frac{\pi}{4}$，从而原累次积分化为
            $$\frac{\pi}{4}\int_0^2\dr r\int_{r^2/2}^2zr^5\dr z=\frac{\pi}{8}\int_0^2\bigg(4-\frac{r^4}{4}\bigg)r^5\dr r=\frac{\pi}{8}\bigg(\frac{4\cdot 2^6}{6}-\frac{2^{10}}{4\cdot10}\bigg)=\frac{32}{15}\pi$$
        \end{itemize}
    }

    \item (习题7.3.8)计算
    $$\iiint_\Omega(x^2+z^2)\dr V$$
    其中
    $$\Omega=\{(x,y,z)\mid x^2+y^2\le z\le1\}$$

    \sol{
        分为如下步骤：
        \begin{itemize}
            \item \textbf{旋转对称性换元}
            
            利用区域关于$z$轴的旋转对称性(详见本讲义20.3.3)，我们考虑柱坐标换元
            $$x=r\cos\theta,\quad y=r\sin\theta,\quad z=z$$
            这样约束条件即化为了
            $$r^2\le z\le1$$
            此时被积函数可写为
            $$r^2\cos^2\theta+z^2$$
            且直接计算Jacobi行列式可知
            $$\dr x\dr y\dr z=r\dr r\dr\theta\dr z$$
            我们将约束条件与柱坐标自然范围确定的$(r,\theta,z)$区域记为$\Omega_0$，原积分化为
            $$\iiint_{\Omega_0}(r^2\cos^2\theta+z^2)r\dr\theta\dr z$$

            \note 本题也可以先通过反射对称性将被积函数改写为$\frac{1}{2}(x^2+y^2)+z^2$再进行操作，后续积分将更加简单。

            \item \textbf{累次积分化}
            
            由于约束中$z$被$r$自然限制，且不涉及$\theta$，可以选择先对$\theta$、再对$z$、最后对$r$积分。此时，我们需要将集合描述成
            $$\alpha\le r\le\beta,\quad \phi_1(r)\le z\le\phi_2(r),\quad\psi_1(r,z)\le\theta\le\psi_2(r,z)$$
            当$z$、$\theta$可任取时($\theta$实际无影响)，$r$最小一定能取到0；$r$最大受$z$限制，可发现$z=1$时能取到最大，为1。综合以上讨论得第一个不等式为
            $$0\le r\le1$$
            当$r$给定、$\theta$任取时，$z$的范围限制即为之前的
            $$r^2\le z\le1$$
            最后，由于$\theta$不受到约束，其范围即为自然范围$[0,2\pi]$。

            综合以上，可最终写出积分
            $$\int_0^1\dr r\int_{r^2}^1\dr z\int_0^{2\pi}(r^3\cos^2\theta+rz^2)\dr\theta$$

            \item \textbf{积分计算}
            
            对于$\theta$积分可将累次积分化为(对$\cos^2\theta$部分的积分过程类似上题，下方计算虽然较为复杂但都是多项式积分，本质上是容易的)
            $$\int_0^1\dr r\int_{r^2}^1(\pi r^3+2\pi rz^2)\dr z=\int_0^1\bigg(\pi r^3(1-r^2)+2\pi r\cdot\frac{1-r^6}{3}\bigg)\dr r=\frac{\pi}{4}-\frac{\pi}{6}+\frac{\pi}{3}-\frac{\pi}{12}=\frac{\pi}{3}$$
        \end{itemize}
    }

    \item (习题7.3.19)计算
    $$\iiint_\Omega(x+1)(y+1)\dr V$$
    其中
    $$\Omega=\bigg\{(x,y,z)\ \bigg|\ \frac{x^2}{a^2}+\frac{y^2}{b^2}+\frac{z^2}{c^2}\le1\bigg\}$$
    这里参数$a>0$、$b>0$、$c>0$。

    \sol{
        分为如下步骤：
        \begin{itemize}
            \item \textbf{反射对称化简}
            
            由于区域$\Omega$满足$(x,y,z)\in\Omega$当且仅当$(-x,y,z)\in\Omega$，类似本讲义19.1.2第5题可知
            $$\iiint_\Omega(x+1)(y+1)\dr V=\iiint_\Omega(-x+1)(y+1)\dr V$$
            两者平均即得原积分化为
            $$\iiint_\Omega(y+1)\dr V$$
            同理，再利用$\Omega$满足$(x,y,z)\in\Omega$当且仅当$(x,-y,z)\in\Omega$，对称并平均即得原积分化为
            $$\iiint_\Omega1\dr V$$
            
            \item \textbf{换元与累次积分化}
            
            根据区域本身为球的部分进行伸缩，为了将区域换元为矩形，我们考虑\textbf{广义球坐标}换元
            $$x=a\rho\sin\varphi\cos\theta,\quad y=b\rho\sin\varphi\sin\theta,\quad z=c\rho\cos\varphi$$
            这样约束条件即化为了$\rho^2\le1$，结合球坐标自然范围有
            $$0\le\rho\le1,\quad0\le\varphi\le\pi,\quad0\le\theta\le2\pi$$
            将上述条件确定的$(\rho,\varphi,\theta)$区域记为$\Omega_0$，被积函数仍为1。 

            直接类似球坐标计算Jacobi行列式可知
            $$\dr x\dr y\dr z=abc\rho^2\sin\varphi\dr\rho\dr\varphi\dr\theta$$
            于是原积分化为(这里积分顺序任取，事实上无影响)
            $$\int_0^{\pi}\dr\varphi\int_0^{2\pi}\dr\theta\int_0^1abc\rho^2\sin\varphi\dr\rho$$

            \item \textbf{积分计算}
            
            对于$\rho$积分可将累次积分化为
            $$\int_0^{\pi}\dr\varphi\int_0^{2\pi}\frac{abc}{3}\sin\varphi\dr\theta=\int_0^\pi\frac{2\pi abc}{3}\sin\varphi\dr\varphi=\frac{4\pi}{3}abc$$
        \end{itemize}
    }

    \item (习题7.3.21)用柱坐标与球坐标分别将三重积分
    $$\iiint_\Omega f(\sqrt{x^2+y^2+z^2})\dr V$$
    化为累次积分，其中$\Omega$是球体$x^2+y^2+z^2\le z$在锥面$z=\sqrt{3x^2+3y^2}$上方的部分。

    \sol{
        我们在进行区域描述后分类处理：
        \begin{itemize}
            \item \textbf{区域描述}
            
            由于``上方''对应$z$更大，区域可描述为满足
            $$x^2+y^2+z^2\le z,\quad z\ge\sqrt{3x^2+3y^2}$$
            的$(x,y,z)$集合。

            \item \textbf{柱坐标换元}
            
            考虑柱坐标换元
            $$x=r\cos\theta,\quad y=r\sin\theta,\quad z=z$$
            这样约束条件即化为了
            $$r^2+z^2\le z,\quad z\ge\sqrt{3r^2}$$
            此时被积函数可写为
            $$f(\sqrt{r^2+z^2})$$
            且直接计算Jacobi行列式可知
            $$\dr x\dr y\dr z=r\dr r\dr\theta\dr z$$
            我们将约束条件与柱坐标自然范围确定的$(r,\theta,z)$区域记为$\Omega_0$，原积分化为
            $$\iiint_{\Omega_0}f(\sqrt{r^2+z^2})r\dr r\dr\theta\dr z$$
            
            \item \textbf{柱坐标累次积分化}
            
            利用柱坐标的自然范围可将第二个约束化简为$z\ge\sqrt{3}r$，约束中$r$被$z$自然限制，且不涉及$\theta$，可以选择先对$\theta$、再对$r$、最后对$z$积分。此时，我们需要将集合描述成
            $$\alpha\le z\le\beta,\quad \phi_1(z)\le r\le\phi_2(z),\quad\psi_1(r,z)\le\theta\le\psi_2(r,z)$$
            当$r$、$\theta$可任取时($\theta$实际无影响)，$z$的范围需要满足$z-z^2\ge0$、$z\ge0$，从而第一个不等式为
            $$0\le z\le1$$
            当$z$给定、$\theta$任取时，$r$的范围限制可写为
            $$r^2\le z-z^2,\quad r\le\frac{\sqrt3}{3}z$$
            结合自然范围即得到
            $$0\le r\le\min\bigg\{\sqrt{z-z^2},\frac{\sqrt3}{3}z\bigg\}$$
            最后，由于$\theta$不受到约束，其范围即为自然范围$[0,2\pi]$。

            综合以上，可最终写出积分
            $$\int_0^1\dr z\int_0^{\min\{\sqrt{z-z^2},z/\sqrt3\}}\dr r\int_0^{2\pi}f(\sqrt{r^2+z^2})r\dr\theta$$
            或对$\theta$积分后化为
            $$2\pi\int_0^1\dr z\int_0^{\min\{\sqrt{z-z^2},z/\sqrt3\}}f(\sqrt{r^2+z^2})r\dr r$$
            
            \item \textbf{球坐标换元}
            
            考虑球坐标换元
            $$x=\rho\sin\varphi\cos\theta,\quad y=\rho\sin\varphi\sin\theta,\quad z=\rho\cos\varphi$$
            这样约束条件即化为了
            $$\rho^2\le\rho\cos\varphi,\quad \rho\cos\varphi\le\sqrt{3\rho^2\sin^2\varphi}$$
            此时被积函数可写为
            $$f(\rho)$$
            且直接计算Jacobi行列式可知
            $$\dr x\dr y\dr z=\rho^2\sin\varphi\dr\rho\dr\varphi\dr\theta$$
            我们将约束条件与球坐标自然范围确定的$(\rho,\varphi,\theta)$区域记为$\Omega_0$，原积分化为
            $$\iiint_{\Omega_0}f(\rho)\rho^2\sin\varphi\dr\rho\dr\varphi\dr\theta$$
            
            \item \textbf{球坐标累次积分化}
            
            利用球坐标的自然范围，忽略$\rho=0$的特殊点可将两个约束化简为
            $$\rho\le\cos\varphi,\quad\cos\varphi\le\sqrt{3}\sin\varphi$$
            
            虽然约束中$\rho$被$\varphi$自然限制，且不涉及$\theta$，但我们希望能将难以处理的$f(\rho)$最后积分，因此先对$\theta$、再对$\varphi$、最后对$\rho$积分。
            
            此时，我们需要将集合描述成
            $$\alpha\le\rho\le\beta,\quad \phi_1(\rho)\le\varphi\le\phi_2(\rho),\quad\psi_1(\rho,\varphi)\le\theta\le\psi_2(\rho,\varphi)$$
            从第一个不等式可看出$\cos\varphi\ge0$，结合球坐标自然范围可以先得到$\varphi\in[0,\frac{\pi}{2}]$，此时第二个不等式即能通过$\tan\varphi$的单调性得到等价于
            $$0\le\varphi\le\frac{\pi}{6}$$
            我们在此基础上进一步分析。当$\varphi$、$\theta$可任取时($\theta$实际无影响)，$\rho$最大值被$\cos\varphi$的最大值($\varphi=0$时)限制，结合自然范围即得第一个不等式为
            $$0\le\rho\le1$$
            当$\rho$给定、$\theta$任取时，利用$\cos\varphi$的单调性，$\varphi$的范围同时被$\arccos\rho$与$\frac{\pi}{6}$限制，结合自然范围也即
            $$0\le\varphi\le\min\bigg\{\frac{\pi}{6},\arccos\rho\bigg\}$$
            最后，由于$\theta$不受到约束，其范围即为自然范围$[0,2\pi]$。

            综合以上，可最终写出积分
            $$\int_0^1\dr\rho\int_0^{\min\{\pi/6,\arccos\rho\}}\dr\varphi\int_0^{2\pi}f(\rho)\rho^2\sin\varphi\dr\theta$$

            \item \textbf{球坐标进一步化简}

            我们最后将上述的累次积分化简为定积分。对$\theta$积分后化为
            $$2\pi\int_0^1\dr\rho\int_0^{\min\{\pi/6,\arccos\rho\}}f(\rho)\rho^2\sin\varphi\dr\varphi$$
            再对$\varphi$积分即得
            $$2\pi\int_0^1\bigg(1-\cos\min\bigg\{\frac{\pi}{6},\arccos\rho\bigg\}\bigg)f(\rho)\rho^2\dr\rho$$
            由于范围内$\cos\frac{\pi}{6}=\frac{\sqrt3}{2}$、$\cos\arccos\rho=\rho$，拆分区域可最终将原积分写为
            $$\frac{6\pi-\pi^2}{3}\int_0^{\sqrt{3}/2}f(\rho)\rho^2\dr\rho+2\pi\int_{\sqrt{3}/2}^1(1-\rho)f(\rho)\rho^2\dr\rho$$
        \end{itemize}

        \note 由于原题只要求化为累次积分，对球坐标可以考虑化为先对$\theta$、再对$\rho$、最后对$\varphi$的积分，这样将得到更加简单的结果
        $$\int_0^{\pi/6}\dr\varphi\int_0^{\cos\varphi}\dr\rho\int_0^{2\pi}f(\rho)\rho^2\sin\varphi\dr\theta$$
        不过，对于抽象函数$f$，这样的形式无法进一步计算化简为定积分。
    }

    \item (总练习题7.14)用直角坐标写出柱坐标下的累次积分
    $$\int_0^{2\pi}\dr\theta\int_0^1\dr r\int_r^{\sqrt{2-r^2}}r\dr z$$
    对应的三重积分的积分区域边界曲面方程。

    \sol{
        \begin{itemize}
            \item \textbf{区域描述}
                
            根据条件，区域可描述为满足
            $$0\le\theta\le2\pi,\quad 0\le r\le1,\quad r\le z\le\sqrt{2-r^2}$$
            的$(r,\theta,z)$集合。

            \item \textbf{换元回直角坐标}
                
            转换为直角坐标后，利用$r=\sqrt{x^2+y^2}$，结合自然约束$r\ge0$、$\theta\in[0,2\pi]$可删去，三个约束条件即化为了
            $$\sqrt{x^2+y^2}\le1,\quad\sqrt{x^2+y^2}\le z\le\sqrt{2-x^2-y^2}$$

            将上述条件确定的$(x,y,z)$区域记为$\Omega_0$，进一步研究。

            由于已有$x^2+y^2\ge0$，可发现第二个条件满足时第一个条件自动满足，从而可以拆分为
            $$z\ge\sqrt{x^2+y^2},\quad z\le\sqrt{2-x^2-y^2}$$
            第一式已经保证了$z\ge0$，从而第二式可以同平方化为更简单的球面方程形式，最终即得到
            $$z\ge\sqrt{x^2+y^2},\quad x^2+y^2+z^2\le2$$
            将两个不等改写为等号就得到了边界曲面方程。
        \end{itemize}
    }

    \item (总练习题7.17)将三重积分
    $$\iiint_\Omega f(x,y,z)\dr V$$
    化为球坐标下的累次积分，其中$\Omega$为球坐标下的曲面$\rho=2$、$\rho=\cos\varphi$与$Oxy$平面围成的区域。

    \sol{
        分为如下步骤：
        \begin{itemize}
            \item \textbf{区域描述}
            
            $Oxy$平面的方程即即$z=0$。根据围成的区域必然以每个曲面为边界，由$\cos\varphi<2$恒成立，前两个约束必然为
            $$\cos\varphi\le\rho\le2$$
            $z=0$在球坐标下即对应$\rho\cos\varphi=0$，忽略$\rho=0$的特殊点也即$\cos\varphi=0$。由于$\cos\varphi\le0$时$\rho\ge\cos\varphi$不再成为边界(会被自然范围$\rho\ge0$限制)，它对应的约束必然为$\cos\varphi\ge0$，综合得到区域描述
            $$0\le\cos\varphi\le\rho\le2$$

            \item \textbf{旋转对称性换元}
            
            考虑球坐标换元
            $$x=\rho\sin\varphi\cos\theta,\quad y=\rho\sin\varphi\sin\theta,\quad z=\rho\cos\varphi$$
            此时被积函数可写为
            $$f(\rho\sin\varphi\cos\theta,\rho\sin\varphi\sin\theta,\rho\cos\varphi)$$
            且直接计算Jacobi行列式可知
            $$\dr x\dr y\dr z=\rho^2\sin\varphi\dr\rho\dr\varphi\dr\theta$$
            我们将上一部分讨论的约束条件与球坐标自然范围确定的$(\rho,\varphi,\theta)$区域记为$\Omega_0$，原积分化为
            $$\iiint_{\Omega_0}f(\rho\sin\varphi\cos\theta,\rho\sin\varphi\sin\theta,\rho\cos\varphi)\rho^2\sin\varphi\dr\rho\dr\varphi\dr\theta$$

            \item \textbf{累次积分化}
            
            由于约束中$\rho$被$\varphi$自然限制，且不涉及$\theta$，可以选择先对$\theta$、再对$\rho$、最后对$\varphi$积分。此时，我们需要将集合描述成
            $$\alpha\le\varphi\le\beta,\quad\phi_1(\varphi)\le\rho\le\phi_2(\varphi),\quad\psi_1(\rho,\varphi)\le\theta\le\psi_2(\rho,\varphi)$$
            对于$\varphi$的约束实际上只有，$\cos\varphi\ge0$，结合球坐标自然范围可以先得到$\varphi\in[0,\frac{\pi}{2}]$，
            当$\varphi$给定、$\theta$任取时，$\rho$的范围即为之前的
            $$\cos\varphi\le\rho\le2$$
            由于$\cos\varphi$已经大于等于0，它已经在自然范围中。最后，由于$\theta$不受到约束，其范围即为自然范围$[0,2\pi]$。

            综合以上，可最终写出积分
            $$\int_0^{\pi/2}\dr\varphi\int_{\cos\varphi}^2\dr\rho\int_0^{2\pi}f(\rho\sin\varphi\cos\theta,\rho\sin\varphi\sin\theta,\rho\cos\varphi)\rho^2\sin\varphi\dr\theta$$
        \end{itemize}
    }
\end{enumerate}

\subsection{重积分与定向}
\subsubsection{考虑定向的定积分}
上下限交换代表反定向

\subsubsection{固定方向的重积分}
重积分的本质是保持定向恒为$z$轴正方向的第二型曲面积分

考虑方向则换元无需Jacobi绝对值

\subsubsection{一维与二维换元}
一维换元并不要求一一对应是由于定向

不带Jacobi绝对值 二维也可以并非一一对应

\subsection{区域的性质}
\subsubsection{描述顺序}
何为``围成''

\begin{framed}
\begin{exercise}\label{get ineq}
    考虑曲面$y=0$、$z=0$、$x+z=\frac{\pi}{2}$与$y=\sqrt{x}$围成的区域$\Omega$，用重积分计算其体积。
\end{exercise}
\sol{
    根据重积分几何意义，其体积即为
    $$\iiint_\Omega1\dr V$$
    计算分为如下步骤：
    \begin{itemize}
        \item \textbf{区域描述}
        
        根据围成的区域必然以每个曲面为边界，由$0<\sqrt{x}$恒成立，$y=0$与$y=\sqrt{x}$对应的约束必然为
        $$0\le y\le\sqrt{x}$$
        否则，若$y=0$对应的约束为$y\le0$，$y=\sqrt{x}$将不可能成为区域边界；若$y=\sqrt{x}$对应的约束为$y\ge\sqrt{x}$，$y=0$将不可能成为区域边界。

        由于约束中已经出现了$\sqrt{x}$，使它有意义需要$x\ge0$，下面考虑$z=0$与$x+z=\frac{\pi}{2}$对应的约束。

        若$z=0$对应的约束为$z\le0$，$x$、$z$符号相反，无论$x+z\ge\frac{\pi}{2}$还是$x+z\le\frac{\pi}{2}$，均有趋于无穷的$x$、$z$满足条件，这与``围成''的有界性要求矛盾。由此，必然对应$z\ge0$。

        最后，若$x+z\ge\frac{\pi}{2}$，$x$、$z$均可趋于无穷，仍与有界性矛盾，由此得到最终的约束形式
        $$0\le y\le\sqrt{x},\quad x\ge0,\quad z\ge0,\quad x+z\le\frac{\pi}{2}$$
        这就是题中区域的代数描述。

        \item \textbf{累次积分化}
            
        由于约束中$y$被$x$自然限制，我们考虑先对$y$、再对$x$、最后对$z$积分($x$与$z$的顺序也可更换)。此时，我们需要将集合描述成
        $$\alpha\le z\le\beta,\quad\phi_1(x)\le z\le\phi_2(x),\quad\psi_1(x,z)\le y\le\psi_2(x,z)$$

        当$x$、$y$可任取时，$z$最小一定能取到0；$z$最大对应$x=0$时，为$\frac{\pi}{2}$。综合得第一个不等式为
        $$0\le z\le\frac{\pi}{2}$$
        当$z$给定、$y$可任取时，$x$受到的约束即
        $$0\le x\le\frac{\pi}{2}-z$$
        当$z$、$x$给定时，关于$y$的不等式即
        $$0\le y\le\sqrt{x}$$
        最终写出积分
        $$\int_0^{\pi/2}\dr z\int_0^{\pi/2-z}\dr x\int_0^{\sqrt{x}}1\dr y$$

        \item \textbf{积分计算}
        
        直接对$y$积分可将累次积分化为
        $$\int_0^{\pi/2}\dr z\int_0^{\pi/2-z}\sqrt{x}\dr x=\frac{2}{3}\int_0^{\pi/2}\bigg(\frac{\pi}{2}-z\bigg)^{3/2}\dr z$$
        由于被积函数是$z$的幂函数复合一次函数，可以直接得到一个原函数为$-\frac{2}{5}\big(\frac{\pi}{2}-z\big)^{5/2}$，从而即能计算得结果为
        $$\frac{4}{15}\frac{\pi^{5/2}}{2^{5/2}}=\frac{\pi^{5/2}\sqrt2}{30}$$
    \end{itemize}
}
\end{framed}

\note 两类容易确定的情况：若同时有$f(x,y,z)=g(x,y,z)$、$f(x,y,z)=h(x,y,z)$，且$g(x,y,z)\le h(x,y,z)$恒成立，则必然$g(x,y,z)\le f(x,y,z)\le h(x,y,z)$；若同时有$a_1x+b_1y=c_1$、$a_2x+b_2y=c_2$、$a_3x+b_3y=c_3$，且三个平面不平行，则必然选择了三条直线相交中间的三角形(可通过简单的几何确定)。

\begin{framed}
\begin{exercise}
    将累次积分
    $$\int_0^a\dr x\int_0^{x^2}\dr y\int_0^{x+y}f(z)\dr z$$
    化为定积分，这里参数$a>0$。
\end{exercise}
\sol{
    分为如下步骤：
    \begin{itemize}
        \item \textbf{区域描述}
            
        根据条件，区域可描述为满足
        $$0\le x\le a,\quad0\le y\le x^2,\quad 0\le z\le x+y$$
        的$(x,y,z)$集合。

        \item \textbf{累次积分化}
            
        为了将原题化为定积分，$f(z)$必须放在最后进行积分，因此需要考虑$x$、$y$、$z$的积分顺序(由于原区域$x$约束$y$，$y$、$x$、$z$的顺序比$x$、$y$、$z$更加自然)。

        此时，我们需要将集合描述成
        $$\alpha\le z\le\beta,\quad\phi_1(z)\le x\le\phi_2(z),\quad\psi_1(x,z)\le y\le\psi_2(x,z)$$
        当$x$、$y$可任取时，$z$最小一定能取到0；$z$最大对应$x=a$、$y=a^2$时，为$a^2+a$。综合得第一个不等式为
        $$0\le z\le a^2+a$$
        当$z$给定、$y$可任取时，若$x$太小、$x^2+x<z$，则即使$y$取到最大值$x^2$也无法保证$x+y\ge z$，因此，$x$实际的最小值并不是0，而是使得$x^2+x\ge z$的最小正数$x$，直接利用二次方程求根公式可得其为$\frac{-1+\sqrt{1+4z}}{2}$。$x$的最大值并无其他限制，一定可取到$a$，由此第二个不等式为
        $$\frac{-1+\sqrt{1+4z}}{2}\le x\le a$$
        当$z$、$x$给定时，关于$y$的不等式可改写为
        $$0\le y\le x^2,\quad y\ge z-x$$
        $y$必须大于等于两个下限中最大的，因此范围即
        $$\max\{0,z-x\}\le y\le x^2$$
        最终写出积分
        $$\int_0^{a^2+a}\dr z\int_{(-1+\sqrt{1+4z})/2}^a\dr x\int_{\max\{0,z-x\}}^{x^2}f(z)\dr y$$

        \item \textbf{积分计算}
            
        由于$f(z)$对$y$为常数，第一步积分可直接算得为
        $$\int_0^{a^2+a}\dr z\int_{(-1+\sqrt{1+4z})/2}^a(x^2-\max\{0,z-x\})f(z)\dr x$$
        为了计算此积分，我们必须讨论何时$z>x$。

        记$\varphi(z)=\frac{-1+\sqrt{1+4z}}{2}$。首先，由于给定$z$时$x$的最小值是通过$x^2+x=z$确定的，必然有
        $$\varphi(z)<z$$
        当$z\in(a,a^2+a)$时，$z>x$恒成立，最大值恒取$z-x$，否则积分将分为两段，由此上述积分分为三个部分求和：
        $$\int_a^{a^2+a}\dr z\int_{\varphi(z)}^a(x^2-z+x)f(z)\dr x$$
        $$\int_0^a\dr z\int_{\varphi(z)}^z(x^2-z+x)f(z)\dr x$$
        $$\int_0^a\dr z\int_z^ax^2f(z)\dr x$$

        \note 这些处理本质是\textbf{定积分计算}的知识，可参考上学期讲义7.3。

        三个部分均为对$x$的多项式积分，直接计算可得到积分结果分别为
        $$\int_a^{a^2+a}f(z)\bigg(\frac{1}{3}(a^3-\varphi^3(z))-z(a-\varphi(z))+\frac{1}{2}(a^2-\varphi^2(z))\bigg)\dr z$$
        $$\int_0^af(z)\bigg(\frac{1}{3}(z^3-\varphi^3(z))-z(z-\varphi(z))+\frac{1}{2}(z^2-\varphi^2(z))\bigg)\dr z$$
        $$\int_0^af(z)\frac{1}{3}(a^3-z^3)\dr z$$
        求和即得到最终的定积分。

        \note 可利用$\varphi^2(z)+\varphi(z)=z$进一步化简形式，不过结果并不能本质上简化。
    \end{itemize}
}
\end{framed}
有时化处的上下限中会存在$\max$与$\min$，可通过类似上方的讨论技巧计算出定积分结果。

\subsubsection{反射对称性}
为了介绍如何利用反射对称性化简重积分，我们先来看两种常见的关于平面反射对称的情况：
\begin{framed}
\begin{exercise}
    证明$(x_0,y_0,z_0)$与$(-x_0,y_0,z_0)$关于平面$x=0$对称，$(x_0,y_0,z)$与$(y_0,x_0,z_0)$关于平面$x=y$对称。
\end{exercise}
\sol{
    验证两点关于平面对称也即验证两点到平面距离相同，且连线与平面垂直：
    \begin{itemize}
        \item 利用点到平面的距离公式，$(x_0,y_0,z_0)$与$(-x_0,y_0,z_0)$到$x=0$的距离均为$|x_0|$，且连接两者的方向向量为$(2x_0,0,0)$，与平面的法向量$(1,0,0)$平行，从而对称。
        \item 将$y=x$写为$x-y=0$。利用点到平面的距离公式，$(x_0,y_0,z_0)$与$(y_0,x_0,z_0)$到$x=y$的距离均为为
        $$\frac{|x_0-y_0|}{\sqrt{1^2+1^2+0^2}}=\frac{|x_0-x_0|}{\sqrt2}$$
        且连接两者的方向向量为$(x_0-y_0,y_0-x_0,0)$，与平面的法向量$(1,-1,0)$平行，从而对称。
        
        \note 可以在$z=0$平面上画出$(x_0,y_0,0)$与$(y_0,x_0,0)$以感受几何直观。
    \end{itemize}
}
\end{framed}

对第一种情况，利用换元可以给出如下结论：
\begin{framed}
\begin{exercise}
    对空间区域$\Omega$，若$(x,y,z)\in\Omega$当且仅当$(-x,y,z)\in\Omega$，证明
    $$\iiint_\Omega f(x,y,z)\dr V=\iiint_\Omega f(-x,y,z)\dr V=\iiint_\Omega\frac{1}{2}(f(x,y,z)+f(-x,y,z))\dr V$$
    $$\iiint_\Omega f(x,y,z)\dr V=\iiint_{\Omega_0}(f(x,y,z)+f(-x,y,z))\dr V$$
    其中
    $$\Omega_0=\{(x,y,z)\in\Omega\mid x\ge0\}$$
\end{exercise}
\sol{
    我们依次证明两个等式：
    \begin{itemize}
        \item 考虑换元$x=-u$、$y=v$、$z=w$，直接计算可发现此为一一对应，换元后的被积函数为$f(-u,v,w)$，且计算Jacobi行列式(注意需要添加绝对值)得
        $$\dr x\dr y\dr z=\dr u\dr v\dr w$$
        设换元后的区域为$\tilde\Omega$，它应是所有使得$(x,y,z)\in\Omega$的$(u,v,w)$构成的集合，而又由于$(u,v,w)=(-x,y,z)$，利用条件$(x,y,z)\in\Omega$当且仅当$(u,v,w)\in\Omega$，这就得到了$\tilde\Omega=\Omega$，从而有
        $$\iiint_\Omega f(x,y,z)\dr x\dr y\dr z=\iiint_\Omega f(-u,v,w)\dr u\dr v\dr w$$
        类似定积分中积分变量可任意更改，我们将$u$、$v$、$w$重新写为$x$、$y$、$z$，即得到了
        $$\iiint_\Omega f(x,y,z)\dr V=\iiint_\Omega f(-x,y,z)\dr V$$
        这就是第一式的第一个等号。既然上述两个积分已经相等，第一个积分应当等于两者的平均值，这就是第一式的第二个等号。
        \item 设
        $$\Omega_1=\{(x,y,z)\in\Omega\mid x\le0\}$$
        上一章中我们已经讨论了平面上的三重积分为0，从而$x=0$时对积分无影响，这就可以得到
        $$\iiint_\Omega f(x,y,z)\dr V=\iiint_{\Omega_0}f(x,y,z)\dr V+\iiint_{\Omega_1}f(x,y,z)\dr V$$
        对比第二式可发现只需证明
        $$\iiint_{\Omega_0}f(-x,y,z)\dr V=\iiint_{\Omega_1}f(x,y,z)\dr V$$

        对上式左侧考虑换元$x=-u$、$y=v$、$z=w$，直接计算可发现此为一一对应，换元后的被积函数为$f(u,v,w)$，且计算Jacobi行列式(注意需要添加绝对值)得
        $$\dr x\dr y\dr z=\dr u\dr v\dr w$$
        设换元后的区域为$\tilde\Omega_0$，它应是所有使得$(x,y,z)\in\Omega_0$的$(u,v,w)$构成的集合，而又由于$(u,v,w)=(-x,y,z)$，利用条件$(x,y,z)\in\Omega$当且仅当$(-x,y,z)\in\Omega$，因此$(x,y,z)\in\Omega$且$x\ge0$当且仅当$(u,v,w)\in\Omega$且$u\le0$。这就得到了$\tilde\Omega_0=\Omega_1$，从而有
        $$\iiint_{\Omega_0}f(-x,y,z)\dr x\dr y\dr z=\iiint_{\Omega_1}f(u,v,w)\dr u\dr v\dr w$$

        类似定积分中积分变量可任意更改，我们将$u$、$v$、$w$重新写为$x$、$y$、$z$，即得结论。
    \end{itemize}
}
\end{framed}

对第二种情况则有如下结论：
\begin{framed}
\begin{exercise}\label{swap}
    对空间区域$\Omega$，若$(x,y,z)\in\Omega$当且仅当$(y,x,z)\in\Omega$，证明
    $$\iiint_\Omega f(x,y,z)\dr V=\iiint_\Omega f(y,x,z)\dr V=\iiint_\Omega\frac{1}{2}(f(x,y,z)+f(y,x,z))\dr V$$
    $$\iiint_\Omega f(x,y,z)\dr V=\iiint_{\Omega_0}(f(x,y,z)+f(y,x,z))\dr V$$
    其中
    $$\Omega_0=\{(x,y,z)\in\Omega\mid y\ge x\}$$
\end{exercise}
\sol{
    我们依次证明两个等式：
    \begin{itemize}
        \item 考虑换元$x=v$、$y=u$、$z=w$，直接计算可发现此为一一对应，换元后的被积函数为$f(v,u,w)$，且计算Jacobi行列式(注意需要添加绝对值)得
        $$\dr x\dr y\dr z=\dr u\dr v\dr w$$
        设换元后的区域为$\tilde\Omega$，它应是所有使得$(x,y,z)\in\Omega$的$(u,v,w)$构成的集合，而又由于$(u,v,w)=(y,x,z)$，利用条件$(x,y,z)\in\Omega$当且仅当$(u,v,w)\in\Omega$，这就得到了$\tilde\Omega=\Omega$，从而有
        $$\iiint_\Omega f(x,y,z)\dr x\dr y\dr z=\iiint_\Omega f(v,u,w)\dr u\dr v\dr w$$
        类似定积分中积分变量可任意更改，我们将$u$、$v$、$w$重新写为$x$、$y$、$z$，即得到了
        $$\iiint_\Omega f(x,y,z)\dr V=\iiint_\Omega f(y,x,z)\dr V$$
        这就是第一式的第一个等号。既然上述两个积分已经相等，第一个积分应当等于两者的平均值，这就是第一式的第二个等号。
        \item 设
        $$\Omega_1=\{(x,y,z)\in\Omega\mid x\ge y\}$$
        上一章中我们已经讨论了平面上的三重积分为0，从而$x=y$时对积分无影响，这就可以得到
        $$\iiint_\Omega f(x,y,z)\dr V=\iiint_{\Omega_0}f(x,y,z)\dr V+\iiint_{\Omega_1}f(x,y,z)\dr V$$
        对比第二式可发现只需证明
        $$\iiint_{\Omega_0}f(y,x,z)\dr V=\iiint_{\Omega_1}f(x,y,z)\dr V$$
        对上式左侧考虑换元$x=v$、$y=u$、$z=w$，直接计算可发现此为一一对应，换元后的被积函数为$f(u,v,w)$，且计算Jacobi行列式(注意需要添加绝对值)得
        $$\dr x\dr y\dr z=\dr u\dr v\dr w$$
        设换元后的区域为$\tilde\Omega_0$，它应是所有使得$(x,y,z)\in\Omega_0$的$(u,v,w)$构成的集合，而又由于$(u,v,w)=(y,x,z)$，利用条件$(x,y,z)\in\Omega$当且仅当$(y,x,z)\in\Omega$，因此$(x,y,z)\in\Omega$且$y\ge x$当且仅当$(u,v,w)\in\Omega$且$u\ge v$。这就得到了$\tilde\Omega_0=\Omega_1$，从而有
        $$\iiint_{\Omega_0}f(y,x,z)\dr x\dr y\dr z=\iiint_{\Omega_1}f(u,v,w)\dr u\dr v\dr w$$

        类似定积分中积分变量可任意更改，我们将$u$、$v$、$w$重新写为$x$、$y$、$z$，即得结论。
    \end{itemize}
}
\end{framed}

我们可以将上方两个例题中的第一个等式称为\textbf{函数折半}，第二个等式则是\textbf{区域折半}，它们都是利用反射对称化简的常用结论。例如，若区域关于$x=0$对称，且被积函数$f(x,y,z)=-f(-x,y,z)$，利用函数折半即可以直接得到积分为0。直观来看，这是由于$f$与区域都关于平面$x=0$对称，因此两侧的积分可以互相\textbf{抵消}。

\note 这实际上与定积分中对称性的应用完全相同(可见上学期讲义7.3.3)，也可以推广到区域对于其他平面对称的情况。

\begin{framed}
\begin{exercise}\label{three column v}
    将三个圆柱面$x^2+y^2=R^2$、$y^2+z^2=R^2$、$x^2+z^2=R^2$围成的区域$\Omega$中的重积分表示为累次积分，并计算其体积，这里参数$R>0$。
\end{exercise}
\sol{
    根据重积分几何意义，其体积即为
    $$\iiint_\Omega1\dr V$$
    计算分为如下步骤：
    \begin{itemize}
        \item \textbf{区域描述}
        
        根据围成的区域必然有界，可发现最终选择的区域必然在三个圆柱面内部，也即满足
        $$x^2+y^2\le R^2,\quad y^2+z^2\le R^2,\quad x^2+z^2\le R^2$$
        的$(x,y,z)$集合。

        \item \textbf{反射对称化简}
        
        由于直接计算累次积分将因三个约束导致过程复杂，且此区域\textbf{并不具有}旋转对称性(详见讲义下一部分分析)，需要考虑先行化简区域。

        我们进行四次反射对称以化简：
        \begin{enumerate}
            \item 由于$(x,y,z)\in\Omega$当且仅当$(-x,y,z)\in\Omega$，且被积函数$f(x,y,z)=1$也满足$f(x,y,z)=f(-x,y,z)$，考虑区域折半，记
            $$\Omega_1=\{(x,y,z)\in\Omega\mid x\ge0\}$$
            则原积分化为
            $$2\iiint_{\Omega_1}1\dr V$$
            \item 由于$(x,y,z)\in\Omega_1$当且仅当$(x,-y,z)\in\Omega_1$，且被积函数$f(x,y,z)=1$也满足$f(x,y,z)=f(x,-y,z)$，考虑区域折半，记
            $$\Omega_2=\{(x,y,z)\in\Omega_1\mid y\ge0\}$$
            则原积分化为
            $$4\iiint_{\Omega_2}1\dr V$$
            \item 由于$(x,y,z)\in\Omega_2$当且仅当$(x,y,-z)\in\Omega_2$，且被积函数$f(x,y,z)=1$也满足$f(x,y,z)=f(x,y,-z)$，考虑区域折半，记
            $$\Omega_3=\{(x,y,z)\in\Omega_2\mid z\ge0\}$$
            则原积分化为
            $$8\iiint_{\Omega_3}1\dr V$$
            \item 由于$(x,y,z)\in\Omega_3$当且仅当$(y,x,z)\in\Omega_3$，且被积函数$f(x,y,z)=1$也满足$f(x,y,z)=f(y,x,z)$，考虑区域折半，记
            $$\Omega_4=\{(x,y,z)\in\Omega_3\mid y\ge x\}$$
            则原积分化为
            $$16\iiint_{\Omega_4}1\dr V$$
        \end{enumerate}

        \note 注意反射对称必须\textbf{一步步}进行，因为反射后区域的对称性未必与原区域相同。例如，$\Omega$满足$(x,y,z)\in\Omega$当且仅当$(y,x,z)\in\Omega$，但$\Omega_1$、$\Omega_4$不满足，$\Omega_2$、$\Omega_3$满足。

        \item \textbf{累次积分化}
        
        根据之前的讨论，$\Omega_4$实际上是满足
        $$x^2+y^2\le R^2,\quad y^2+z^2\le R^2,\quad x^2+z^2\le R^2,\quad y\ge x\ge0,\quad z\ge0$$
        的$(x,y,z)$集合。

        可以发现，若$y^2+z^2\le R^2$，由$y\ge x\ge0$可得$x^2+z^2\le R^2$，从而约束可以减少为
        $$x^2+y^2\le R^2,\quad y^2+z^2\le R^2,\quad y\ge x\ge0,\quad z\ge0$$
        由于$x$、$z$都被$y$自然限制，考虑先对$x$、再对$z$、最后对$y$的积分顺序。
        
        \note 这样对$x$与对$z$的积分\textbf{不会互相影响}，之后的积分计算中将看到。若先对$x$、再对$y$、最后对$z$，积分将变得非常复杂，读者有兴趣可以自行尝试。
        
        此时，我们需要将集合描述成
        $$\alpha\le y\le\beta,\quad\phi_1(y)\le z\le\phi_2(y),\quad\psi_1(y,z)\le x\le\psi_2(y,z)$$
        当$x$、$z$可任取时，$y$最小一定能取到0；$y$最大对应$x=z=0$时，为$R$。综合得第一个不等式为
        $$0\le y\le R$$
        当$y$给定、$x$可任取时，$z$最小一定可取到0，最大受约束$y^2+z^2\le R^2$限制，不会超过$\sqrt{R^2-y^2}$，进一步分析可发现此时取$x=0$即符合要求，从而第二个不等式为
        $$0\le z\le\sqrt{R^2-y^2}$$
        当$z$、$y$给定时，由$x\ge0$，关于$x$的不等式可改写为
        $$0\le x\le y,\quad x\ge\sqrt{R^2-y^2}$$
        $x$必须小于等于两个上限中最小的，因此范围即
        $$0\le x\le\min\{y,\sqrt{R^2-y^2}\}$$
        最终写出积分
        $$16\int_0^R\dr y\int_0^{\sqrt{R^2-y^2}}\dr z\int_0^{\min\{y,\sqrt{R^2-y^2}\}}1\dr x$$

        \note 即使不进行化简，也理应可以写出累次积分进行计算，只是其中的$\max$、$\min$会相对复杂。

        \item \textbf{积分计算}
        
        第一步积分直接计算即得
        $$16\int_0^R\dr y\int_0^{\sqrt{R^2-y^2}}\min\{y,\sqrt{R^2-y^2}\}\dr z$$
        由于$\min$中内容对$z$为常数，可以直接进行第二步积分得到
        $$16\int_0^R\sqrt{R^2-y^2}\min\{y,\sqrt{R^2-y^2}\}\dr y$$
        直接计算可以发现，范围中$y>\sqrt{R^2-y^2}$当且仅当$y>\frac{\sqrt2}{2}R$，从而积分可写为
        $$16\int_0^{\sqrt2R/2}y\sqrt{R^2-y^2}\dr y+16\int_{\sqrt2R/2}^R(R^2-y^2)\dr y$$
        换元$t=y^2$即得
        $$16\int_0^{\sqrt2R/2}y\sqrt{R^2-y^2}\dr y=8\int_0^{R^2/2}\sqrt{R^2-t}\dr t$$
        由于这是$t$的幂函数复合一次函数，可直接写出被积函数的原函数
        $$-\frac{2}{3}(R^2-t)^{3/2}+C$$
        从而计算得
        $$16\int_0^{\sqrt2R/2}y\sqrt{R^2-y^2}\dr y=\bigg(\frac{16}{3}-\frac{4\sqrt{2}}{3}\bigg)R^3$$
        另一方面直接计算有
        $$16\int_{\sqrt2R/2}^R(R^2-y^2)\dr y=16R^3-8\sqrt2R^3-\frac{16}{3}R^3+\frac{4\sqrt{2}}{3}R^3$$
        求和两部分最终得到结果为
        $$(16-8\sqrt{2})R^3$$
    \end{itemize}
}
\end{framed}

\begin{framed}
\begin{exercise}
    计算四重积分
    $$\iiiint_\Omega(x+y+z+w)^2\dr x\dr y\dr z\dr w$$
    其中$\Omega$是曲面$|x|+|y|+|z|+|w|=1$围成的部分。
\end{exercise}
\sol{
    虽然我们并未学过四重积分，实际的计算过程是完全类似三重的。计算分为如下步骤：
    \begin{itemize}
        \item \textbf{区域描述}
        
        根据围成的区域必然有界，可发现最终选择的区域必然在曲面内部，也即满足
        $$|x|+|y|+|z|+|w|\le1$$
        的$(x,y,z,w)$集合。

        \item \textbf{反射对称化简}
        
        由于$\Omega$满足$(x,y,z,w)\in\Omega$当且仅当$(-x,y,z,w)\in\Omega$，而$f(x,y,z,w)=xy$满足$f(x,y,z,w)=-f(-x,y,z,w)$，考虑区域折半可发现
        $$\iiiint_\Omega xy\dr x\dr y\dr z\dr w=0$$
        同理可知$xz$、$xw$、$yz$、$yw$、$zw$的积分均为0，展开平方得原积分化为
        $$\iiiint_\Omega(x^2+y^2+z^2+w^2)\dr x\dr y\dr z\dr w$$
        与\exref{three column v}类似，依次考虑对$x=0$、$y=0$、$z=0$、$w=0$的反射对称，可将原积分化为
        $$16\iiiint_{\Omega_0}(x^2+y^2+z^2+w^2)\dr x\dr y\dr z\dr w$$
        这里
        $$\Omega_0=\{(x,y,z,w)\in\Omega\mid x\ge0,\ y\ge0,\ z\ge0,\ w\ge0\}$$
        此时可发现，由于四者都非负，可去除约束中的绝对值，最终写为
        $$x\ge0,\quad y\ge0,\quad z\ge0,\quad w\ge0,\quad x+y+z+w\le1$$
        更进一步地，分析可发现$(x,y,z,w)\in\Omega_0$当且仅当$(y,x,z,w)\in\Omega_0$，从而换元可发现
        $$\iiiint_{\Omega_0}x^2\dr x\dr y\dr z\dr w=\iiiint_{\Omega_0}y^2\dr x\dr y\dr z\dr w$$
        同理，$z^2$、$w^2$在$\Omega_0$的积分也与$x^2$相等，积分最终化为
        $$64\iiiint_{\Omega_0}x^2\dr x\dr y\dr z\dr w$$

        \item \textbf{累次积分化}
        
        由于积分对$x$的形式最复杂，我们最后对$x$积分，之前的顺序可任意选择，考虑$w$、$z$、$y$、$x$的顺序。此处我们直接写出累次积分的四个式子：
        \begin{enumerate}
            \item 当$w$、$z$、$y$可任取时，$x$最小值一定可取到0，最大值当$w=z=y=0$时取到，为1，由此范围是
            $$0\le x\le1$$
            \item 当$x$给定，$w$、$z$可任取时，$y$最小值一定可取到0，最大值当$w=z=0$时取到，为$1-x$，由此范围是
            $$0\le y\le1-x$$
            \item 当$x$、$y$给定，$w$可任取时，$z$最小值一定可取到0，最大值当$w=0$时取到，为$1-x-y$，由此范围是
            $$0\le z\le1-x-y$$
            \item 当$x$、$y$、$z$给定时，$w$的范围可直接化简得到
            $$0\le w\le1-x-y-z$$
        \end{enumerate}
        综合以上讨论，可得到累次积分
        $$\int_0^1\dr x\int_0^{1-x}\dr y\int_0^{1-x-y}\dr z\int_0^{1-x-y-z}x^2\dr w$$

        \item \textbf{积分计算}
        
        由于只涉及多项式积分，可直接计算逐步化简，对$w$积分得
        $$64\int_0^1\dr x\int_0^{1-x}\dr y\int_0^{1-x-y}(1-x-y-z)x^2\dr z$$
        对$z$积分得(换元$t=1-x-y-z$可更简便地计算)
        $$64\int_0^1\dr x\int_0^{1-x}\frac{1}{2}(1-x-y)^2x^2\dr y$$
        对$y$积分得(换元$t=1-x-y$可更简便地计算)
        $$64\int_0^1\frac{1}{6}(1-x)^3x^2\dr x$$
        对$x$积分得其为(换元$t=1-x$)
        $$\frac{32}{3}\int_0^1(t^5-2t^4+t^3)\dr t=\frac{8}{45}$$
    \end{itemize}
}
\end{framed}

\subsubsection{旋转对称性}
旋转对称的标识与换元

球坐标看成两次柱坐标，优先柱坐标

\begin{framed}
\begin{exercise}
    计算积分
    $$\iiint_\Omega(x^2+y^2+z^2)\dr V$$
    其中区域
    $$\Omega=\{(x,y,z)\mid x^2+y^2\le 2z\le1\}$$
\end{exercise}
\sol{
    分为如下步骤：
    \begin{itemize}
        \item \textbf{换元与累次积分化}
        
        由于此区域描述中只包含$x^2+y^2$而没有独立的$x$、$y$，关于$z$轴旋转对称，考虑柱坐标换元
        $$x=r\cos\theta,\quad y=r\sin\theta,\quad z=z$$
        此时被积函数为$r^2+z^2$，且直接计算Jacobi行列式可知
        $$\dr x\dr y\dr z=r\dr r\dr\theta\dr z$$
        换元后约束条件即化为了
        $$r^2\le 2z\le1$$
        由于$z$自然被$r$约束，而被积函数对$\theta$为常数，采取$\theta$、$z$、$r$的积分顺序。当$z$、$\theta$可任取时，$r$的范围是$[0,1]$，结合柱坐标的自然范围可进一步得到得
        $$0\le r\le1,\quad\frac{r^2}{2}\le z\le\frac{1}{2},\quad 0\le\theta\le 2\pi$$
        这已经是符合累次积分要求的范围，从而可最终写出积分
        $$\int_0^1\dr r\int_{r^2/2}^{1/2}\dr z\int_0^{2\pi}r(r^2+z^2)\dr\theta$$

        \item \textbf{积分计算}
        
        直接对$\theta$积分可将累次积分化为
        $$2\pi\int_0^1\dr r\int_{r^2/2}^{1/2}r(r^2+z^2)\dr z=2\pi\int_0^1\bigg(\frac{1-r^2}{2}r^3+\frac{1-r^6}{24}r\bigg)\dr r$$
        直接计算多项式积分可得结果
        $$2\pi\bigg(\frac{1}{8}-\frac{1}{12}+\frac{1}{48}-\frac{1}{192}\bigg)=\frac{11}{96}\pi$$
    \end{itemize}
}
\end{framed}

\begin{framed}
\begin{exercise}
    计算积分
    $$\iiint_\Omega(x^2+y^2+z^2)\dr V$$
    其中区域
    $$\Omega=\{(x,y,z)\mid x^2+y^2+z^2\le 2z\}$$
\end{exercise}
\sol{
    本题我们将介绍三种方法写出累次积分，柱坐标换元、直接球坐标换元与平移的球坐标换元：
    \begin{itemize}
        \item \textbf{柱坐标换元}
        
        由于此区域描述中只包含$x^2+y^2$而没有独立的$x$、$y$，关于$z$轴旋转对称，考虑柱坐标换元
        $$x=r\cos\theta,\quad y=r\sin\theta,\quad z=z$$
        此时被积函数为$r^2+z^2$，且直接计算Jacobi行列式可知
        $$\dr x\dr y\dr z=r\dr r\dr\theta\dr z$$
        换元后约束条件即化为了
        $$r^2\le 2z-z^2$$
        由于$r$自然被$z$约束，而被积函数对$\theta$为常数，采取$\theta$、$r$、$z$的积分顺序。当$r$、$\theta$可任取时，由于需要满足自然范围$r\ge0$，$z$的范围是$[0,2]$，结合柱坐标的自然范围可进一步得到
        $$0\le z\le2,\quad0\le r\le\sqrt{2z-z^2},\quad 0\le\theta\le 2\pi$$
        这已经是符合累次积分要求的范围，从而可最终写出积分
        $$\int_0^2\dr z\int_0^{\sqrt{2z-z^2}}\dr r\int_0^{2\pi}r(r^2+z^2)\dr\theta$$
        被积函数中包含$r$的项只有$r$与$r^3$，因此它对$r$积分后实际上是关于$z$的多项式，可以直接算出结果，由于过程相对复杂，我们在之后的方法中再进行计算。

        \item \textbf{直接球坐标换元}
        
        出于区域描述与被积函数均含有$x^2+y^2+z^2$，球坐标换元可能使得计算更加简单。考虑球坐标换元
        $$x=\rho\sin\varphi\cos\theta,\quad y=\rho\sin\varphi\sin\theta,\quad z=\rho\cos\varphi$$
        此时被积函数为$\rho^2$，且直接计算Jacobi行列式可知
        $$\dr x\dr y\dr z=\rho^2\sin\varphi\dr\rho\dr\varphi\dr\theta$$
        换元后约束条件即化为了
        $$\rho^2\le2\rho\cos\varphi$$
        忽略$\rho=0$的特殊点，由$\rho\ge0$化简得到
        $$\rho\le2\cos\varphi$$
        由于需要满足自然范围$\rho\ge0$，结合$\varphi$的自然范围，其取值范围是$[0,\frac{\pi}{2}]$，结合球坐标的自然范围可进一步得到
        $$0\le\varphi\le\frac{\pi}{2},\quad 0\le\rho\le2\cos\varphi,\quad 0\le\theta\le2\pi$$
        这已经是符合累次积分要求的范围，从而可最终写出积分
        $$\int_0^{\pi/2}\dr\varphi\int_0^{2\cos\varphi}\dr\rho\int_0^{2\pi}\rho^4\sin\varphi\dr\theta$$
        逐步计算得其为
        $$2\pi\int_0^{\pi/2}\dr\varphi\int_0^{2\cos\varphi}\rho^4\sin\varphi\dr\rho=\frac{64\pi}{5}\int_0^{\pi/2}\cos^5\varphi\sin\varphi\dr\varphi$$
        换元$t=\cos\varphi$即得到其为
        $$\frac{64\pi}{5}\int_0^1t^5\dr t=\frac{32\pi}{15}$$

        \item \textbf{平移的球坐标换元}
        
        可以发现区域实际上为球体$x^2+y^2+(z-1)^2\le1$，由此考虑平移后的球坐标换元(也即不以原点为球坐标中心)
        $$x=\rho\sin\varphi\cos\theta,\quad y=\rho\sin\varphi\sin\theta,\quad z=\rho\cos\varphi+1$$
        此时被积函数为$\rho^2+2\rho\cos\varphi+1$，且直接计算Jacobi行列式可知(由于加常数不影响求导，Jacobi行列式应与球坐标完全相同)
        $$\dr x\dr y\dr z=\rho^2\sin\varphi\dr\rho\dr\varphi\dr\theta$$
        换元后约束条件即化为了
        $$\rho^2\le1$$
        结合球坐标的自然范围可进一步得到
        $$0\le\rho\le1,\quad 0\le\varphi\le\pi,\quad 0\le\theta\le2\pi$$
        这已经是符合累次积分要求的范围，从而可最终写出积分
        $$\int_0^1\dr\rho\int_0^\pi\dr\varphi\int_0^{2\pi}\rho^2\sin\varphi(\rho^2+2\rho\cos\varphi+1)\dr\theta$$
        逐步计算得其为($\sin\varphi\cos\varphi$在$[0,\pi]$的积分可由对称性知为0)
        $$2\pi\int_0^1\dr\rho\int_0^\pi\rho^2\sin\varphi(\rho^2+2\rho\cos\varphi+1)\dr\varphi=2\pi\int_0^12(\rho^4+\rho^2)=\frac{32\pi}{15}$$
    \end{itemize}
}
\end{framed}
从以上三种方法的比较中可以看出，对于实际问题，采取不同的换元方式可能会有不同的计算效果。本质上，本题由于区域具有良好的球对称性，应当优先考虑\textbf{平移的球坐标换元}。

\begin{framed}
\begin{exercise}
    计算积分
    $$\iiint_\Omega(x^2+y^2+z^2)\dr V$$
    其中$\Omega$是曲面$x^2+y^2+z^2=R^2$、$x^2+y^2-z^2=0$在$z\ge0$部分围成的区域。
\end{exercise}
\sol{
    我们这里介绍柱坐标换元的方法。计算分为如下步骤：
    \begin{itemize}
        \item \textbf{区域描述}
        
        ``在$z\ge0$部分''应当理解为，$z=0$不为区域的边界，但区域恒满足$z\ge0$。由此，讨论$z$与$\pm\sqrt{x^2+y^2}$的大小关系可发现第二个约束只能化为
        $$z\ge\sqrt{x^2+y^2}$$ 
        根据``围成''的有界性要求，即得第一个约束必然为
        $$x^2+y^2+z^2\le R^2$$
        区域最终表述为满足以上条件的$(x,y,z)$集合。

        \item \textbf{换元与累次积分化}
        
        由于此区域描述中只包含$x^2+y^2$而没有独立的$x$、$y$，关于$z$轴旋转对称，考虑柱坐标换元
        $$x=r\cos\theta,\quad y=r\sin\theta,\quad z=z$$
        此时被积函数为$r^2+z^2$，且直接计算Jacobi行列式可知
        $$\dr x\dr y\dr z=r\dr r\dr\theta\dr z$$
        换元后约束条件即化为了(结合柱坐标的自然范围$r\ge0$)
        $$r^2+z^2\le R^2,\quad z\ge r$$
        由于$z\ge0$，第一个约束也即$z\le\sqrt{R^2-r^2}$，我们可以得到
        $$r\le z\le\sqrt{R^2-r^2}$$
        而被积函数对$\theta$为常数，采取$\theta$、$z$、$r$的积分顺序。当$z$、$\theta$可任取时，由于$r$需要满足$r\le\sqrt{R^2-r^2}$，实际取值范围是$[0,\frac{\sqrt2}{2}R]$，结合柱坐标的自然范围可进一步得到得
        $$0\le r\le\frac{\sqrt2}{2}R,\quad r\le z\le\sqrt{R^2-r^2},\quad 0\le\theta\le 2\pi$$
        这已经是符合累次积分要求的范围，从而可最终写出积分
        $$\int_0^{\sqrt2R/2}\dr r\int_r^{\sqrt{R^2-r^2}}\dr z\int_0^{2\pi}r(r^2+z^2)\dr\theta$$

        \item \textbf{积分计算}
        
        直接对$\theta$积分可将累次积分化为
        $$2\pi\int_0^{\sqrt2R/2}\dr r\int_r^{\sqrt{R^2-r^2}}r(r^2+z^2)\dr z$$
        对$z$计算多项式积分可进一步化为
        $$2\pi\int_0^{\sqrt2R/2}\bigg(r^3(\sqrt{R^2-r^2}-r)+\frac{r}{3}((R^2-r^2)\sqrt{R^2-r^2}-r^3)\bigg)\dr r$$
        整理得
        $$2\pi\int_0^{\sqrt2R/2}\bigg(\frac{R^2}{3}r\sqrt{R^2-r^2}+\frac{2}{3}r^3\sqrt{R^2-r^2}-\frac{4}{3}r^4\bigg)\dr r$$
        我们分别计算三部分：换元$s=r^2$可知(可直接找到幂函数复合一次函数的原函数)
        $$\int_0^{\sqrt2R/2}r\sqrt{R^2-r^2}\dr r=\frac{1}{2}\int_0^{R^2/2}\sqrt{R^2-s}\dr s=\frac{1}{3}\bigg(1-\frac{\sqrt2}{4}\bigg)R^3$$
        换元$r=R\sin s$可知
        $$\int_0^{\sqrt2R/2}r^3\sqrt{R^2-r^2}=R^5\int_0^{\pi/4}\sin^3s\cos^2s\dr s$$
        利用$\sin^2s=1-\cos^2s$，进一步换元$t=\cos s$可知
        $$\int_0^{\sqrt2R/2}r^3\sqrt{R^2-r^2}=R^5\int_{\sqrt2/2}^1(1-t^2)t^2\dr t=R^5\frac{4-\sqrt{2}}{30}$$
        直接计算可知
        $$\int_0^{\sqrt2/2}r^4\dr r=\frac{\sqrt2}{40}R^5$$
        综合可得结果为
        $$2\pi R^5\bigg(\frac{1}{9}\bigg(1-\frac{\sqrt2}{4}\bigg)+\frac{4-\sqrt{2}}{45}-\frac{\sqrt2}{30}\bigg)=\frac{2-\sqrt2}{5}\pi R^5$$
    \end{itemize}

    事实上，本题\textbf{球坐标换元}会导致出现$\sin^2\varphi\le\cos^2\varphi$这样的约束，更难正确分析区域，但最终区域事实上为长方体
    $$0\le\rho\le R,\quad0\le\varphi\le\frac{\pi}{4},\quad0\le\theta\le2\pi$$
    积分计算过程将远比柱坐标简单，大家可以自己尝试采用球坐标的计算。
}
\end{framed}

总的来说，当区域具有一定的旋转对称性但并非椭球体时，用柱坐标往往比球坐标更直接，因为柱坐标的约束往往是幂函数形式，而球坐标下容易出现三角函数的不等式。不过，若三角函数的不等式是容易确定范围的，球坐标有时会更加简单。

本质上来说，球坐标也可以看作先对$(x,y,z)$中的$x$、$y$采用柱坐标换元，再对$(r,\theta,z)$中的$r$、$z$采用柱坐标换元。因此，对于区域非椭球体的问题仍然建议\textbf{先采用柱坐标换元}，若难以计算再尝试球坐标。

即使区域不具有完全的旋转对称性，使用柱坐标或球坐标换元也可能可以方便计算：
\begin{framed}
\begin{exercise}
    将曲面
    $$x^2=a^2-az,\quad x^2+y^2=\frac{a^2}{4},\quad z=0$$
    围成的区域$\Omega$中的重积分表示为累次积分，并计算其体积，这里参数$a>0$。
\end{exercise}
\sol{
    根据重积分几何意义，其体积即为
    $$\iiint_\Omega1\dr V$$
    计算分为如下步骤：
    \begin{itemize}
        \item \textbf{区域描述}
        
        根据围成的有界性，第二个曲面必然对应约束
        $$x^2+y^2\le\frac{a^2}{4}$$
        否则$y$将可趋于无穷。
        
        将第一个曲面的表达式化为
        $$z=a-\frac{x^2}{a}$$
        由于第二个曲面对应的约束保证了$x^2\le\frac{a^2}{4}$，必然有$a-\frac{x^2}{a}>0$，由此，根据围成的区域必然以每个曲面为边界，只能为
        $$0\le z\le a-\frac{x^2}{a}$$
        综合以上，区域可描述为满足
        $$x^2+y^2\le\frac{a^2}{4},\quad0\le z\le a-\frac{x^2}{a}$$
        的$(x,y,z)$集合。

        \item \textbf{换元与累次积分化}
        
        虽然此区域并不旋转对称，若以柱坐标表述约束，区域描述将变得简单，由此考虑柱坐标换元
        $$x=r\cos\theta,\quad y=r\sin\theta,\quad z=z$$
        此时被积函数仍为1，且直接计算Jacobi行列式可知
        $$\dr x\dr y\dr z=r\dr r\dr\theta\dr z$$
        换元后约束条件即化为了
        $$r^2\le\frac{a^2}{4},\quad 0\le z\le a-\frac{r^2\cos^2\theta}{a}$$
        结合柱坐标的自然范围可得
        $$0\le r\le\frac{a}{2},\quad0\le\theta\le 2\pi,\quad 0\le z\le a-\frac{r^2\cos^2\theta}{a}$$
        这已经是符合累次积分要求的范围，从而可最终写出积分
        $$\int_0^{a/2}\dr r\int_0^{2\pi}\dr\theta\int_0^{a-r^2\cos^2\theta/a}r\dr z$$        

        \item \textbf{积分计算}
        
        直接对$z$积分可将累次积分化为
        $$\int_0^{a/2}\dr r\int_0^{2\pi}r\bigg(a-r^2\frac{\cos^2\theta}{a}\bigg)\dr\theta$$
        类似本讲义19.1.1第4题的对称性处理与计算可知
        $$\int_0^{2\pi}\cos^2\theta\dr\theta=\pi$$
        于是进一步计算得结果为
        $$\int_0^{a/2}\bigg(2\pi ra-\frac{\pi r^3}{a}\bigg)\dr r=\frac{15\pi}{64}a^3$$
    \end{itemize}
}
\end{framed}

\subsection{积分计算}
\subsubsection{对称性与换元}
\begin{framed}
\begin{exercise}\label{three column s}
    计算三个圆柱面$x^2+y^2=R^2$、$y^2+z^2=R^2$、$x^2+z^2=R^2$围成的区域$\Omega$的表面积。
\end{exercise}
\sol{
    分为如下步骤：
    \begin{itemize}
        \item \textbf{区域描述与对称化简}
        
        根据公式，我们可以知道$z=f(x,y)$在$(x,y)$属于某区域时的表面积计算方式(更换$x$、$y$、$z$字母顺序也可类似计算)。为了应用公式，我们需要先用这样的形式刻画表面。

        由于围成的区域以每个曲面为表面，我们可以将每个约束分别取到等号，其他约束取不等号(不等号方向的确定见\exref{three column v})，这样就把$\Omega$的表面分成了三个部分。三个部分的相交处会有两个约束取等号，实际是一条曲线，面积为0，因此$\Omega$表面积就是$S_1$、$S_2$、$S_3$面积之和。
        $$S_1:\quad x^2+y^2=R^2,\quad y^2+z^2\le R^2,\quad x^2+z^2\le R^2$$
        $$S_2:\quad y^2+z^2=R^2,\quad x^2+y^2\le R^2,\quad x^2+z^2\le R^2$$
        $$S_3:\quad x^2+z^2=R^2,\quad x^2+y^2\le R^2,\quad y^2+z^2\le R^2$$
        与\exref{three column v}中区域$\Omega$的对称性判断类似，直接计算可以发现可以发现
        $$(x,y,z)\in S_1\ \Longleftrightarrow\ (z,y,x)\in S_2\ \Longleftrightarrow\ (x,z,y)\in S_3$$
        也即$S_1$与$S_2$关于$x=z$对称、$S_1$与$S_3$关于$y=z$对称，三者面积应当相等，我们只需计算$S_1$的面积再乘3即可。

        进一步地，根据$x^2+y^2=R^2$，我们又可以将$S_1$分为两部分：
        $$S_1^+:\quad y=\sqrt{R^2-x^2},\quad y^2+z^2\le R^2,\quad x^2+z^2\le R^2$$
        $$S_1^-:\quad y=-\sqrt{R^2-x^2},\quad y^2+z^2\le R^2,\quad x^2+z^2\le R^2$$
        可以发现$(x,y,z)\in S_1^+$当且仅当$(x,-y,z)\in S_1^-$，因此两部分的面积相同，只需要计算$S_1^+$的面积再乘6即可。

        $S_1^+$中可以将$y$看作$f(x,z)$\ (对$z$偏导为0)，由此直接代入公式可知面积为
        $$\iint_D\sqrt{1+\bigg(\frac{\dr\sqrt{R^2-x^2}}{\dr x}\bigg)^2}\dr x\dr z$$
        这里$D$为(已经代入$y=\sqrt{R^2-x^2}$)
        $$\{(x,z)\mid R^2-x^2+z^2\le R^2,\ x^2+z^2\le R^2\}$$
        化简可知积分即
        $$\iint_D\frac{R}{\sqrt{R^2-x^2}}\dr x\dr z$$
        而$D$即
        $$\{(x,z)\mid z^2\le x^2,\ x^2+z^2\le R^2\}$$

        \item \textbf{进一步化简}
        
        考虑到$D$的对称性，我们事实上还可以作进一步的化简：由于$D$关于$x$的约束只涉及$x^2$，$(x,z)\in D$当且仅当$(-x,z)\in D$，记
        $$D_1=\{(x,z)\in D\mid x\ge0\}$$
        利用\textbf{区域折半}可发现积分化为
        $$2\iint_{D_1}\frac{R}{\sqrt{R^2-x^2}}\dr x\dr z$$
        接下来，由于$D_1$关于$z$的约束只涉及$z^2$，$(x,z)\in D_1$当且仅当$(x,-z)\in D$，记
        $$D_2=\{(x,z)\in D\mid z\ge0\}$$
        再次利用区域折半可发现积分化为
        $$4\iint_{D_2}\frac{R}{\sqrt{R^2-x^2}}\dr x\dr z$$
        由于已经保证了$x$、$z$均非负，$D_2$的约束可化简为
        $$0\le z\le x,\quad x^2+z^2\le R^2$$

        \item \textbf{累次积分化}
        
        由于被积函数对$z$为常数，考虑先$z$后$x$的积分顺序，我们需要将集合描述成
        $$\alpha\le x\le\beta,\quad \phi_1(x)\le z\le\phi_2(x)$$
        根据约束，$z$任取时$x$最小值最大值均在$z=0$时可取到，为
        $$0\le x\le R$$
        当$x$给定时，$z$最小一定可取到0，最大受$x$与$\sqrt{R^2-x^2}$控制，从而范围是
        $$0\le z\le\min\{x,\sqrt{R^2-x^2}\}$$

        这已经是符合累次积分要求的范围，从而可以写出累次积分
        $$4\int_0^R\dr x\int_0^{\min\{x,\sqrt{R^2-x^2}\}}\frac{R}{\sqrt{R^2-x^2}}\dr z$$
    
        \item \textbf{积分计算}
        
        对$z$积分可将累次积分化为
        $$4\int_0^R\frac{\min\{x,\sqrt{R^2-x^2}\}R}{\sqrt{R^2-x^2}}\dr x$$
        讨论$\sqrt{R^2-x^2}$与$x$何者更小可将区域分为两段，积分化为
        $$4\int_0^{\sqrt2R/2}\frac{xR\dr x}{\sqrt{R^2-x^2}}+4\int_{\sqrt2R/2}^RR\dr x$$
        对第一部分，换元$t=x^2$可得(可直接找到幂函数复合一次函数的原函数)
        $$4\int_0^{\sqrt2R/2}\frac{xR\dr x}{\sqrt{R^2-x^2}}=2\int_0^{R^2/2}\frac{R\dr t}{\sqrt{R^2-t}}=(4-2\sqrt2)R^2$$
        第二部分直接计算可得结果为
        $$(4-2\sqrt2)R^2$$
        两部分求和即得结果为
        $$(8-4\sqrt2)R^2$$
        乘6得到最终的表面积
        $$(48-24\sqrt2)R^2$$
    \end{itemize}
}
\end{framed}

\begin{framed}
\begin{exercise}
    将累次积分
    $$\int_0^a\dr x\int_0^x\dr y\int_0^yf(x+z)\dr z$$
    化为定积分，这里参数$a>0$。
\end{exercise}
\sol{
    分为如下步骤：
    \begin{itemize}
        \item \textbf{区域描述}
        
        根据条件，区域可描述为满足
        $$0\le x\le a,\quad0\le y\le x,\quad 0\le z\le y$$
        的$(x,y,z)$集合。

        \item \textbf{换元}
        
        为了能将其化为定积分，必须让$f$中只有一个参数，也即必须换元$u=x+z$，且最后对$u$积分。为了形式简单，我们换元
        $$u=x+z,\quad v=y,\quad w=z$$
        直接计算Jacobi行列式可得
        $$\dr x\dr y\dr z=\dr u\dr v\dr w$$
        且被积函数即化为了$f(u)$，约束为
        $$0\le u-w\le a,\quad 0\le v\le u-w,\quad 0\le w\le v$$
        将以上约束确定的$(u,v,w)$区域记为$\Omega_0$，原积分化为
        $$\iiint_{\Omega_0}f(u)\dr u\dr v\dr w$$

        \item \textbf{累次积分化}
        
        由于约束条件中$v$自然被$u$、$w$限制，而根据化为定积分的需要应最后对$u$积分，采用先对$v$、再对$w$、最后对$u$的积分顺序。此时，我们需要将集合描述成
        $$\alpha\le u\le\beta,\quad\phi_1(u)\le w\le\phi_2(u),\quad\psi_1(u,w)\le v\le\psi_2(u,w)$$
        当$v$可任取时，关于$u$的约束实际上为$0\le u\le w+a$与$0\le w\le u-w$。进一步任取$w$，可以发现$u$最小一定可取到0，当$w=u-w=a$时$u$最大，为$2a$\ (也可直接通过$x$与$z$最大均为$a$、$x+z$最大为$2a$得到)，综合得第一个不等式为
        $$0\le u\le 2a$$
        当$u$给定、$v$可任取时，可化简关于$w$的约束为
        $$u\le w+a,\quad w\ge0,\quad w\le\frac{u}{2}$$
        将两个下界取最大值即得到第二个不等式
        $$\max\{0,u-a\}\le w\le\frac{u}{2}$$
        最后，关于$v$的不等式直接为
        $$w\le v\le u-w$$
        综合以上讨论，我们可以写出累次积分
        $$\int_0^{2a}\dr u\int_{\max\{0,u-a\}}^{u/2}\dr w\int_w^{u-w}f(u)\dr v$$

        \item \textbf{积分计算}

        直接对$v$积分可将累次积分化为
        $$\int_0^{2a}\dr u\int_{\max\{0,u-a\}}^{u/2}(u-2w)f(u)\dr w$$
        再对$w$积分(这是简单的多项式积分)即得其为
        $$\int_0^{2a}\bigg(\frac{u^2}{4}-u\max\{0,u-a\}+(\max\{0,u-a\})^2\bigg)f(u)$$
        也可进一步按照0到$a$与$a$到$2a$拆分为
        $$\int_0^a\frac{u^2}{4}f(u)\dr u+\int_a^{2a}\bigg(\frac{u^2}{4}-au+a^2\bigg)f(u)\dr u$$
    \end{itemize}
}
\end{framed}

对更一般的问题，仍然可以遵循``先处理可以处理的部分''的技巧：
\begin{framed}
\begin{exercise}
    考虑参数方程
    $$\begin{cases}x=a(t-\sin t)\\y=a(1-\cos t)\end{cases}(0\le t\le2\pi)$$
    将此曲线与$x$轴围成的区域记作$D$，将
    $$\iint_Df(x)y\dr x\dr y$$
    化为定积分，这里$f$是$\rb$上的连续函数。
\end{exercise}
\sol{
    分为如下步骤：
    \begin{itemize}
        \item \textbf{区域描述与累次积分化}
        
        计算得$t=0$时$x=0$、$t=2\pi$时$x=2\pi a$，且求导可知$x$随$t$增加单调增。此外，进一步计算得$y$恒非负，从而区域可以描述为
        $$0\le x\le2\pi a,\quad0\le y\le\phi(x)$$
        这里$\phi(x)$是由参数方程确定的$y$关于$x$的函数。这已经是符合累次积分要求的范围，从而可以写出累次积分
        $$\int_0^{2\pi a}\dr x\int_0^{\phi(x)}f(x)y\dr y$$

        \item \textbf{积分计算}
        
        直接对$y$积分可将累次积分化为
        $$\int_0^{2\pi a}\frac{1}{2}f(x)\phi^2(x)\dr x$$
        不过，计算到此处并不是结束，因为$\phi(x)$仍是未知的函数。

        利用参数方程的定义，我们可以在$y=\phi(x)$中代入参数方程，得到
        $$\phi(a(t-\sin t))=a(1-\cos t),\quad t\in[0,2\pi]$$
        由此，我们将$x$直接换元为$a(t-\sin t)$，利用$t=0$时$x=0$、$t=2\pi$时$x=2\pi a$，即得积分化为
        $$\frac{1}{2}\int_0^{2\pi a}f(a(t-\sin t))a^2(1-\cos t)^2\dr(a(t-\sin t))=\frac{a^3}{2}\int_0^{2\pi}f(a(t-\sin t))(1-\cos t)^3\dr t$$
        若给定$f$的形式，上述积分已经可以计算，从而最终结果即
        $$\frac{a^3}{2}\int_0^{2\pi}f(a(t-\sin t))(1-\cos t)^3\dr t$$
    \end{itemize}
}
\end{framed}

此外，换元为长方体以化简计算区域也是常用的技巧。不过，换元并不是一定可以简化计算，例如如下的问题：
\begin{framed}
\begin{exercise}
    计算积分
    $$\iiint_\Omega y^4\dr V$$
    其中区域$\Omega$由曲面$x=az^2$、$x=bz^2$、$x=\alpha y$、$x=\beta y$与$x=h$围成，这里参数$b>a>0$、$\beta>\alpha>0$、$h>0$。
\end{exercise}
\sol{
    我们在区域描述后介绍换元与不换元两种计算方法：
    \begin{itemize}
        \item \textbf{区域描述}
        
        根据围成的区域必然以每个曲面为边界，从$az^2\le bz^2$可直接得到前两个曲面对应的约束为
        $$az^2\le x\le bz^2$$
        这已经带有要求$x\ge0$，从而由于$x=\alpha y$、$x=\beta y$为边界，必然也有$y\ge0$恒成立，于是$ay\le by$，$x=\alpha y$、$x=\beta y$两平面对应的约束为
        $$\alpha y\le x\le\beta y$$
        再结合有界性可知
        $$x\le h$$
        区域可描述为满足以上五个约束的$(x,y,z)$集合。
        
        \item \textbf{直接计算}
        
        可以发现，当$y$、$z$可任取时$x$的范围即
        $$0\le x\le h$$
        而前四个约束都可以化为$y$、$z$被$x$限制：
        $$\frac{x}{\beta}\le y\le\frac{x}{\alpha},\quad\sqrt{\frac{x}{b}}\le z\le\sqrt{\frac{x}{a}}$$
        这已经是符合累次积分要求的范围，从而可以写出累次积分
        $$\int_0^h\dr x\int_{x/\beta}^{x/\alpha}\dr y\int_{\sqrt{x/b}}^{\sqrt{x/a}}y^4\dr z$$
        对$z$计算积分得
        $$\int_0^h\dr x\int_{x/\beta}^{x/\alpha}y^4\sqrt{x}\bigg(\frac{1}{\sqrt a}-\frac{1}{\sqrt b}\bigg)\dr y$$
        再对$y$积分即得到
        $$\frac{1}{5}\int_0^hx^{11/2}\bigg(\frac{1}{\sqrt a}-\frac{1}{\sqrt b}\bigg)\bigg(\frac{1}{\alpha^5}-\frac{1}{\beta^5}\bigg)\dr x$$
        这就可以直接算出结果
        $$\frac{2}{65}h^{13/2}\bigg(\frac{1}{\sqrt a}-\frac{1}{\sqrt b}\bigg)\bigg(\frac{1}{\alpha^5}-\frac{1}{\beta^5}\bigg)$$

        \item \textbf{换元计算}
        
        由于排除$x=0$的特殊点(可排除的理由类似球坐标，详见本讲义19.3.2)时约束可写为
        $$0\le x\le h,\quad a\le\frac{x}{z^2}\le b,\quad\alpha\le\frac{x}{y}\le\beta$$
        可以想到换元
        $$u=x,\quad v=\frac{x}{z^2},\quad w=\frac{x}{y}$$
        这时被积函数为$u^4/w^4$约束可化为
        $$0\le u\le h,\quad a\le v\le b,\quad\alpha\le w\le\beta$$
        将此三约束确定的$(u,v,w)$区域记为$\Omega_0$。

        由于$\Omega_0$中除$w=0$的特殊点外可反解出
        $$x=u,\quad y=\frac{u}{w},\quad z=\sqrt{\frac{v}{w}}$$
        直接计算Jacobi行列式可知
        $$\begin{vmatrix}1&0&0\\1/w&0&-u/w^2\\1/(2\sqrt{uv})&-\sqrt{u}/(2v^{3/2})&0\end{vmatrix}=-\frac{1}{2}u^{3/2}v^{-3/2}w^{-2}$$
        其在$\Omega_0$中$u=0$、$v=0$、$w=0$的特殊点外恒负，换元是合理的，且
        $$\dr x\dr y\dr z=\frac{1}{2}u^{3/2}v^{-3/2}w^{-2}\dr u\dr v\dr w$$
        综合以上讨论可写出累次积分
        $$\frac{1}{2}\int_0^h\dr u\int_a^b\dr v\int_\alpha^\beta u^{11/2}v^{-3/2}w^{-6}\dr w$$
        进一步计算可得与前一种方法相同的结果。
    \end{itemize}

    \note 对比本题的两种方法，虽然换元后区域更简单，实际的计算量其实更大了，也更容易出错。由此，容易直接描述区域时建议\textbf{不要优先考虑换元}。
}
\end{framed}

\subsubsection{综合练习}
首先，由于描述累次积分时上下限常出现$\max$或$\min$，需要熟悉如何讨论计算带有分段函数的累次积分：
\begin{framed}
\begin{exercise}
    计算积分
    $$\iint_D\max\{xy,x^3\}\dr x\dr y$$
    其中
    $$D=\{(x,y)\mid x\in[-1,1],\ y\in[0,1]\}$$
\end{exercise}
\sol{
    \note 注意虽然$xy$与$x^3$均对直线$x=0$反对称，取$\max$后\textbf{不再具有}，因此无法直接得到积分为0。

    由于此区域的描述是可以直接化为累次积分的，且被积函数的形式对$y$相对简单，我们可以直接写出累次积分
    $$\int_{-1}^1\dr x\int_0^1\max\{xy,x^3\}\dr y$$
    为了进一步计算，需要讨论$xy$与$x^3$何者更大。当$x>0$时，$xy\ge x^3$当且仅当$y\ge x^2$，否则当且仅当$y\le x^2$，由此，可将其先拆分为
    $$\int_{-1}^0\dr x\int_0^1\max\{xy,x^3\}\dr y+\int_0^1\dr x\int_0^1\max\{xy,x^3\}\dr y$$
    再根据$\max$的取值进一步拆分为
    $$\int_{-1}^0\dr x\int_0^{x^2}xy\dr y+\int_{-1}^0\dr x\int_{x^2}^1x^3\dr y+\int_0^1\dr x\int_0^{x^2}x^3\dr y+\int_0^1\dr x\int_{x^2}^1xy\dr y$$
    分别计算每个部分
    $$\int_{-1}^0\dr x\int_0^{x^2}xy\dr y=\frac{1}{2}\int_{-1}^0x^5=-\frac{1}{12}$$
    $$\int_{-1}^0\dr x\int_{x^2}^1x^3\dr y=\int_{-1}^0(x^3-x^5)\dr x=-\frac{1}{12}$$
    $$\int_0^1\dr x\int_0^{x^2}x^3\dr y=\int_0^1x^5\dr x=\frac{1}{6}$$
    $$\int_0^1\dr x\int_{x^2}^1xy\dr y=\int_0^1\frac{x-x^5}{2}\dr x=\frac{1}{6}$$
    全部求和得到最终结果为$\frac{1}{6}$。
}
\end{framed}

\begin{framed}
\begin{exercise}
    计算积分
    $$\iiint_\Omega\frac{2xy+1}{x^2+y^2+z^2}\dr x\dr y\dr z$$
    其中$\Omega$曲面$x^2+y^2+z^2=2a^2$与$az=x^2+y^2$在$z\ge0$部分围成的区域。

    \note 本题中的函数在原点无定义，而区域包含原点，实际上是一个\textbf{广义积分}。由于期中范围尚未学到如何严谨处理，我们只以之前的方法进行形式上的计算并得出结果，不论证合理性。
\end{exercise}
\sol{
    分为如下步骤：
    \begin{itemize}
        \item \textbf{区域描述}

        根据``围成''的有界性要求，球面必然对应内部
        $$x^2+y^2+z^2\le 2a^2$$
        结合条件在$z\ge0$部分，可知必然(``在$z\ge0$部分''应当理解为，$z=0$不为区域的边界，但区域恒满足$z\ge0$)
        $$az\ge x^2+y^2$$
        由此$\Omega$的代数描述最终为满足如上两个不等式的$(x,y,z)$集合。

        \item \textbf{反射对称化简}
        
        由于区域$\Omega$满足$(x,y,z)\in\Omega$当且仅当$(-x,y,z)\in\Omega$，而$\frac{xy}{x^2+y^2+z^2}=-\frac{(-x)y}{(-x)^2+y^2+z^2}$，利用函数折半可得
        $$\iiint_\Omega\frac{2xy}{x^2+y^2+z^2}\dr x\dr y\dr z=0$$
        于是原积分化为
        $$\iiint_\Omega\frac{1}{x^2+y^2+z^2}\dr x\dr y\dr z$$

        \item \textbf{旋转对称性换元}
        
        \note 这里介绍球坐标换元的方法，\textbf{柱坐标换元}也可计算，可以作为积分计算的练习。柱坐标换元的积分区域分析将更加简单，不过积分计算相对更复杂。
        
        根据区域的旋转对称性，可以考虑球坐标换元
        $$x=\rho\sin\varphi\cos\theta,\quad y=\rho\sin\varphi\sin\theta,\quad z=\rho\cos\varphi$$
        这样约束条件即成为了
        $$\rho^2\le 2a^2,\quad a\rho\cos\varphi\ge\rho^2\sin^2\varphi$$
        结合球坐标自然范围，忽略$\rho=0$的特殊点可得
        $$0\le\rho\le\sqrt2a,\quad a\cos\varphi\ge\rho\sin^2\varphi$$        
        将上述条件与球坐标自然范围确定的$(\rho,\varphi,\theta)$区域记为$\Omega_0$。此时被积函数成为$\frac{1}{\rho^2}$，直接计算Jacobi行列式可知
        $$\dr x\dr y\dr z=\rho^2\sin\varphi\dr\rho\dr\varphi\dr\theta$$
        原积分化为
        $$\iiint_{\Omega_0}\sin\varphi\dr\rho\dr\varphi\dr\theta$$

        \item \textbf{累次积分化}
        
        由于$\rho$自然被$\varphi$约束，而区域与被积函数均和$\theta$无关，采取$\theta$、$\rho$、$\varphi$的积分顺序。此时，我们需要将集合描述成
        $$\alpha\le\varphi\le\beta,\quad \phi_1(\varphi)\le\rho\le\phi_2(\varphi),\quad\psi_1(\rho,\varphi)\le\theta\le\psi_2(\rho,\varphi)$$

        根据$\rho\ge0$的自然范围，存在合理的$\rho$至少需要$\cos\varphi\ge0$，结合球坐标自然范围可知第一个不等式为
        $$0\le\varphi\le\frac{\pi}{2}$$
        $\rho$最小一定可取到0，最大同时存在两个约束，由此可知$\varphi\ne0$时
        $$0\le\rho\le\min\bigg\{\sqrt2a,\frac{a\cos\varphi}{\sin^2\varphi}\bigg\}$$
        当$\varphi=0$时，分析可发现$a\cos\varphi\ge\rho\sin^2\varphi$恒成立，由此相当于认为$\min$中的第二项为无穷，取第一个约束$\sqrt2a$。

        最后，$\theta$的范围即为球坐标自然范围
        $$0\le\theta\le2\pi$$
        综合以上讨论可写出累次积分
        $$\int_0^{\pi/2}\dr\varphi\int_0^{\min\{\sqrt2a,a\cos\varphi/\sin^2\varphi\}}\dr\rho\int_0^{2\pi}\sin\varphi\dr\theta$$

        \item \textbf{积分计算}

        直接对$\theta$积分可将积分化为
        $$2\pi\int_0^{\pi/2}\dr\varphi\int_0^{\min\{\sqrt2a,a\cos\varphi/\sin^2\varphi\}}\sin\varphi\dr\rho$$
        再对$\rho$积分即
        $$2\pi\int_0^{\pi/2}\min\bigg\{\sqrt2a,\frac{a\cos\varphi}{\sin^2\varphi}\bigg\}\sin\varphi\dr\varphi$$
        利用三角函数的单调性可以发现，在$[0,\frac{\pi}{4}]$里$\min$中第一项更小，在$[\frac{\pi}{4},\frac{\pi}{2}]$里$\min$中第二项更小，因此积分化为
        $$2\pi a\bigg(\int_0^{\pi/4}\sqrt2\sin\varphi\dr\varphi+\int_{\pi/4}^{\pi/2}\frac{\cos\varphi}{\sin\varphi}\dr\varphi\bigg)$$
        两个被积函数的一个原函数(第二部分将$\cos\varphi$放入$\dr\varphi$中)分别为$-\sqrt2\cos\varphi$与$\ln\sin\varphi$，由此结果为(利用对数的性质进行了化简)
        $$2\pi a\bigg(\sqrt2-1+\frac{1}{2}\ln2\bigg)$$
    \end{itemize}
}
\end{framed}

\begin{framed}
\begin{exercise}
    计算积分
    $$\iiint_\Omega\frac{(x+y+z)^2\sqrt{1+x^2+y^2}}{(x^2+y^2+z^2)(1+x^2+y^2+z^2)}\dr V$$
    其中区域$\Omega$由曲面$x^2+y^2+1=z^2$、$3x^2+3y^2+3=z^2$、$x^2+y^2=1$在$z\ge0$部分围成。
\end{exercise}
\sol{
    分为如下步骤：
    \begin{itemize}
        \item \textbf{区域描述}

        根据``围成''的有界性要求，圆柱必然对应内部
        $$x^2+y^2\le1$$
        又由于围成需要以每个曲面为边界，根据$x^2+y^2+1<3x^2+3y^2+3$可知对$z$的约束必然为
        $$x^2+y^2+1\le z^2\le 3x^2+3y^2+3$$
        结合$z\ge0$的条件，可同开根号化为
        $$\sqrt{x^2+y^2+1}\le z\le\sqrt{3x^2+3y^2+3}$$

        由此$\Omega$的代数描述最终为满足如上三个不等式的$(x,y,z)$集合。

        \item \textbf{反射对称化简}
        
        由于区域$\Omega$满足$(x,y,z)\in\Omega$当且仅当$(-x,y,z)\in\Omega$，而
        $$\frac{xy\sqrt{1+x^2+y^2}}{(x^2+y^2+z^2)(1+x^2+y^2+z^2)}=-\frac{(-x)y\sqrt{1+(-x)^2+y^2}}{((-x)^2+y^2+z^2)(1+(-x)^2+y^2+z^2)}$$
        利用函数折半可得
        $$\iiint_\Omega\frac{xy\sqrt{1+x^2+y^2}}{(x^2+y^2+z^2)(1+x^2+y^2+z^2)}\dr x\dr y\dr z=0$$
        将分子中的$(x+y+z)^2$展开，类似可消去含$yz$、$zx$的项，于是原积分化为
        $$\iiint_\Omega\frac{(x^2+y^2+z^2)\sqrt{1+x^2+y^2}}{(x^2+y^2+z^2)(1+x^2+y^2+z^2)}\dr x\dr y\dr z=\iiint_\Omega\frac{\sqrt{1+x^2+y^2}}{1+x^2+y^2+z^2}\dr x\dr y\dr z$$


        \note 一般建议将分子上平方和等部分\textbf{完全展开}后观察可以用对称性消去的部分。

        \item \textbf{换元与累次积分化}
        
        根据区域的旋转对称性，可以考虑柱坐标换元
        $$x=r\cos\theta,\quad y=r\sin\theta,\quad z=z$$
        这样约束条件即成为了
        $$\sqrt{r^2+1}\le z\le\sqrt{3r^2+3},\quad r^2\le1$$
        结合球坐标自然范围可得
        $$0\le r\le1,\quad\sqrt{r^2+1}\le z\le\sqrt{3r^2+3},\quad0\le\theta\le2\pi$$        
        此时被积函数成为$\frac{\sqrt{1+r^2}}{(1+r^2+z^2)}$，直接计算Jacobi行列式可知
        $$\dr x\dr y\dr z=r\dr r\dr\theta\dr z$$
        由于我们得到的换元后区域描述已经符合累次积分的要求，可以直接写出累次积分
        $$\int_0^1\dr r\int_{\sqrt{r^2+1}}^{\sqrt{3r^2+3}}\dr z\int_0^{2\pi}\frac{r\sqrt{1+r^2}}{1+r^2+z^2}\dr\theta$$

        \item \textbf{积分计算}

        直接对$\theta$积分可将积分化为
        $$2\pi\int_0^1\dr r\int_{\sqrt{r^2+1}}^{\sqrt{3r^2+3}}\frac{r\sqrt{1+r^2}}{1+r^2+z^2}\dr z$$
        利用上册教材3.2节的积分公式，$a>0$时
        $$\int\frac{1}{a^2+z^2}\dr z=\frac{1}{a}\arctan\frac{z}{a}+C$$
        从而可直接得到被积函数对$z$的一个原函数为
        $$r\arctan\frac{z}{\sqrt{1+r^2}}$$
        而$\arctan\sqrt3=\frac{\pi}{3}$、$\arctan1=\frac{\pi}{4}$，从而直接得到对$z$积分后成为
        $$2\pi\int_0^1r\bigg(\frac{\pi}{3}-\frac{\pi}{4}\bigg)\dr r=\frac{\pi^2}{12}$$
    \end{itemize}

    \note 上学期常用的积分公式也务必这学期\textbf{保持记忆}。
}
\end{framed}

本质上，无论是换元还是选择累次积分顺序都是在找一种易于计算的描述方法，而涉及无穷时，描述方法的浮动空间更大：
\begin{framed}
\begin{exercise}\label{gauss}
    对参数$t>0$，设
    $$I(t)=\int_0^t\er^{-x^2}\dr x$$
    证明
    $$I^2(t)=\iint_{D_t}\er^{-(x^2+y^2)}\dr x\dr y$$
    其中
    $$D_t=\{(x,y)\in\rb^2\mid0\le x\le t,\ 0\le y\le t\}$$
    利用上式计算极限
    $$\lim_{t\to+\infty}I(t)$$
\end{exercise}
\sol{
    我们将证明分为三个部分：
    \begin{itemize}
        \item \textbf{化为二重积分}
        
        由于$\er^{-(x^2+y^2)}=\er^{-x^2}\er^{-y^2}$，利用\exref{separation}的结论(此结论基本可直接使用)可知
        $$\iint_{D_t}\er^{-(x^2+y^2)}\dr x\dr y=\int_0^t\er^{-x^2}\dr x\int_0^t\er^{-y^2}\dr y$$
        由于定积分更换变量无影响，可将所有$y$也换为$x$，这就得到了结论。

        \item \textbf{$J(t)$的构造与估算}

        本题最困难的点在于，$\er^{-x^2}$不存在初等原函数，$I(t)$是无法直接计算的，极限也难以算得。不过，将$I(t)$用二重积分表示后，若区域的形状改编，$\er^{-(x^2+y^2)}$的积分可能是可以计算的。可以想象，无论区域是何种形状，当放大到趋于整个第一象限时，极限应当是恒定的。由此，我们构造
        $$J(t)=\int_{C_t}\er^{-(x^2+y^2)}\dr x\dr y$$
        其中
        $$C_t=\{(x,y)\in\rb^2\mid x^2+y^2\le t^2,\ x\ge0,\ y\ge0\}$$
        我们先证明对任何$t>0$有
        $$J(t)\le I^2(t)\le J(2t)$$

        \note 从之后的过程中可以看到，这是为了配凑出\textbf{夹逼定理}的形式进行估算。

        首先，无论是$J(t)$、$I^2(t)$还是$J(2t)$，被积函数都是$\er^{-(x^2+y^2)}$，只是积分区域分别为
        $$C_t=\{(x,y)\in\rb^2\mid x^2+y^2\le t^2,\ x\ge0,\ y\ge0\}$$
        $$D_t=\{(x,y)\in\rb^2\mid0\le x\le t,\ 0\le y\le t\}$$
        $$C_{2t}=\{(x,y)\in\rb^2\mid x^2+y^2\le 4t^2,\ x\ge0,\ y\ge0\}$$
        直接通过代数方式可验证$C_t\subset D_t\subset C_{2t}$。

        由于被积函数$\er^{-(x^2+y^2)}$在$\rb^2$恒正，它在$D_t$内$C_t$外的部分(记为$D_t\backslash C_t$)积分应当非负，在$C_{2t}$内$D_t$外的部分(记为$C_{2t}\backslash D_t$)积分也应当非负，由此利用重积分的性质即得
        $$\iint_{D_t}\er^{-(x^2+y^2)}\dr x\dr y=\int_{C_t}\er^{-(x^2+y^2)}\dr x\dr y+\iint_{D_t\backslash C_t}\er^{-(x^2+y^2)}\dr x\dr y\ge\int_{C_t}\er^{-(x^2+y^2)}\dr x\dr y$$
        $$\int_{C_{2t}}\er^{-(x^2+y^2)}\dr x\dr y=\iint_{D_t}\er^{-(x^2+y^2)}\dr x\dr y+\int_{C_{2t}\backslash D_t}\er^{-(x^2+y^2)}\dr x\dr y\ge\iint_{D_t}\er^{-(x^2+y^2)}\dr x\dr y$$
        这就是得到了证明。

        \item \textbf{结论计算}

        由于$J(t)$的旋转对称性，采用\textbf{极坐标换元}进行计算：设
        $$x=r\cos\theta,\quad y=r\sin\theta$$
        结合极坐标自然范围可发现约束变为($x\ge0$、$y\ge0$即对应第一象限部分，由此可直接写出$\theta$的范围，这样结合几何意义的写法往往比代数分析更加简单)
        $$0\le r\le t,\quad0\le\theta\le\frac{\pi}{2}$$
        直接计算Jacobi行列式可得
        $$\dr x\dr y=r\dr r\dr\theta$$
        且被积函数成为$\er^{-r^2}$。换元后的区域描述已经符合累次积分的要求，从而可以写出累次积分
        $$J(t)=\int_0^t\dr r\int_0^{\pi/2}\er^{-r^2}r\dr\theta=\frac{\pi}{2}\int_0^t\er^{-r^2}r\dr r$$
        换元$s=r^2$可得
        $$J(t)=\frac{\pi}{4}\int_0^s\er^{-s}\dr s=\frac{\pi}{4}(1-\er^{-t})$$
        由此代入上一部分的不等式知
        $$\frac{\pi}{4}(1-\er^{-t})\le I^2(t)\le\frac{\pi}{4}(1-\er^{-2t})$$
        在不等式两侧同取极限，利用\textbf{夹逼定理}即得
        $$\lim_{t\to+\infty}I^2(t)=\frac{\pi}{4}$$
        由条件可知$I(t)$为正数，从而$I(t)=\sqrt{I^2(t)}$，利用复合函数极限性质得最终结论
        $$\lim_{t\to+\infty}I(t)=\sqrt{\frac{\pi}{4}}=\frac{\sqrt\pi}{2}$$
    \end{itemize}
}
\end{framed}

\subsubsection{重积分计算总结}

\newpage
\section{曲线与曲面积分}
\subsection{习题解答}
\subsubsection{作业解答}
\begin{enumerate}
    \item (习题8.1.4)计算
    $$\int_L\frac{1}{x^2+y^2+z^2}\dr s$$
    其中$L$为螺旋线段
    $$\begin{cases}x=a\cos t\\y=a\sin t\\z=bt\end{cases}(0\le t\le2\pi)$$
    这里参数$a>0$、$b>0$。

    \sol{
        分为如下步骤：
        \begin{itemize}
            \item \textbf{化为定积分}

            用下标$t$表示对$t$求导，直接计算有
            $$x_t=-a\sin t,\quad y_t=a\cos t,\quad z_t=b$$
            从而
            $$\dr s=\sqrt{x_t^2+y_t^2+z_t^2}\dr t=\sqrt{a^2+b^2}\dr t$$
            且进一步计算可发现被积函数为$a^2+b^2t^2$，由此原积分为
            $$\int_0^{2\pi}\frac{1}{a^2+b^2t^2}\sqrt{a^2+b^2}\dr t$$
            
            \item \textbf{积分计算}
            
            利用上册教材3.2节的积分公式，$c>0$时
            $$\int\frac{1}{c^2+z^2}\dr z=\frac{1}{c}\arctan\frac{z}{c}+C$$
            我们提出常数部分$\frac{1}{b^2}\sqrt{a^2+b^2}$再如此操作，可以得到被积函数的一个原函数(由$a>0$、$b>0$)
            $$\frac{\sqrt{a^2+b^2}}{ab}\arctan\frac{bt}{a}$$
            由此最终得到结果为
            $$\frac{\sqrt{a^2+b^2}}{ab}\arctan\frac{2\pi b}{a}$$
        \end{itemize}
    }

    \item (习题8.1.5)计算
    $$\oint_C(x+y)\dr s$$
    其中$C$为极坐标下的双纽线$r^2=a^2\cos(2\theta)$在$x\ge0$的部分，这里参数$a>0$。

    \sol{
        分为如下步骤：
        \begin{itemize}
            \item \textbf{参数化}
            
            由直角坐标与极坐标的关系
            $$\begin{cases}x=r\cos\theta\\y=r\sin\theta\end{cases}$$
            代入极坐标下曲线方程即得参数方程
            $$\begin{cases}x=a\sqrt{\cos(2\theta)}\cos\theta\\y=a\sqrt{\cos(2\theta)}\sin\theta\end{cases}$$
            由极坐标自然范围$\theta\in[0,2\pi]$，由$r$需非负可知$\cos(2\theta)\ge0$，再由$x\ge0$可知$\cos\theta\ge0$。综合这些条件可以得到参数的最终范围
            $$\theta\in\bigg[0,\frac{\pi}{4}\bigg]\cup\bigg[\frac{7\pi}{4},2\pi\bigg]$$

            \note 注意虽然参数范围并不连续，但曲线仍然是连续的，这是由于$\theta=0$和$\theta=2\pi$对应同一个点。

            \item \textbf{化为定积分与计算}     

            用下标$\theta$表示对$\theta$求导，直接计算有
            $$x_\theta=a\bigg(-\frac{\sin(2\theta)}{\sqrt{\cos(2\theta)}}\cos\theta-\sqrt{\cos(2\theta)}\sin\theta\bigg),\quad y_\theta=a\bigg(-\frac{\sin(2\theta)}{\sqrt{\cos(2\theta)}}\sin\theta+\sqrt{\cos(2\theta)}\cos\theta\bigg)$$
            从而
            $$\dr s=\sqrt{x_\theta^2+y_\theta^2}\dr\theta=a\sqrt{\frac{\sin^2(2\theta)}{\cos(2\theta)}+\cos(2\theta)}\dr\theta=\frac{a}{\sqrt{\cos(2\theta)}}\dr\theta$$

            \note 注意在$\cos(2\theta)=0$处$\dr s$不存在，形式上我们忽略这样的特殊点可得到正常的结果，这样处理的合理性详见本讲义22.3.3。

            进一步计算可发现被积函数为$a(\cos\theta+\sin\theta)\sqrt{\cos(2\theta)}$，由此原积分为
            $$\int_0^{\pi/4}a^2(\cos\theta+\sin\theta)\dr\theta+\int_{7\pi/4}^{2\pi}a^2(\cos\theta+\sin\theta)\dr\theta=\sqrt2a^2$$          
        \end{itemize}

        \note 事实上，利用区域的对称性可知$y$在曲线上积分为0，但由于我们\textbf{并未学过曲线积分如何换元}(重参数化)，建议大家化为定积分后再进行对称性处理，如本题中$\sin\theta$在两部分的积分为相反数，可抵消。 
    }

    \item (习题8.1.7)计算
    $$\int_L\sqrt{x^2+y^2}\dr s$$
    其中$L$为曲线段
    $$\begin{cases}x=a(\cos t+t\sin t)\\y=a(\sin t-t\cos t)\end{cases}(0\le t\le2\pi)$$
    这里参数$a$为实数。

    \sol{
        分为如下步骤：
        \begin{itemize}
            \item \textbf{化为定积分}
            
            用下标$t$表示对$t$求导，直接计算有
            $$x_t=at\cos t,\quad y_t=at\sin t$$
            从而(由范围内$t\ge0$)
            $$\dr s=\sqrt{x_t^2+y_t^2}\dr t=at\dr t$$

            进一步计算可发现被积函数为$a\sqrt{1+t^2}$，由此原积分为
            $$a^2\int_0^{2\pi}\sqrt{1+t^2}t\dr t$$

            \item \textbf{积分计算}
            
            换元$s=t^2$可将原积分化为
            $$\frac{a^2}{2}\int_0^{4\pi^2}\sqrt{1+s}\dr s$$
            被积函数已经为幂函数复合一次函数，直接利用原函数可得到积分为
            $$\frac{a^2}{3}\big((4\pi^2+1)^{3/2}-1\big)$$
        \end{itemize}
    }

    \item (习题8.1.11)计算
    $$\int_Lx^2\dr s$$
    其中$L$为圆周
    $$\begin{cases}x^2+y^2+z^2=a^2\\x+y+z=0\end{cases}$$
    这里参数$a>0$。

    \sol{
        分为如下步骤：
        \begin{itemize}
            \item \textbf{参数化}

            首先，从$x+y+z=0$可以直接解出一个变量，这里选择解出$z=-x-y$。

            代入第一个方程可得
            $$x^2+y^2+(x+y)^2=a^2$$
            现在，我们只需要给上面的$x$、$y$用参数方程表示，再代入$z=-x-y$即可得到$z$的参数。

            这是一个关于$x$、$y$的二次函数，对这类问题应采取\textbf{逐步配方}的方式。我们先对$y$后对$x$配方(这是为了让之后的计算更加简单，大家可以自行尝试先对$x$后对$y$)，按照$y$整理得到
            $$2y^2+2xy+2x^2-a^2=0$$
            对$y$配方得到
            $$2\bigg(y+\frac{x}{2}\bigg)^2+\frac{3}{2}x^2-a^2=0$$
            此时剩下的项已经是平方，配方结束，我们最终得到了
            $$2\bigg(y+\frac{x}{2}\bigg)^2+\frac{3}{2}x^2=a^2$$
            这已经与椭圆方程的形式类似(事实上也的确代表$x$、$y$构成椭圆)，从而令
            $$\sqrt2\bigg(y+\frac{x}{2}\bigg)=a\cos\theta,\quad\sqrt{\frac{3}{2}}x=a\sin\theta,\quad\theta\in[0,2\pi]$$
            反解出$x$、$y$就得到了
            $$x=a\frac{\sqrt2}{\sqrt3}\sin\theta,\quad y=a\bigg(\frac{1}{\sqrt2}\cos\theta-\frac{1}{\sqrt6}\sin\theta\bigg)$$
            最终代入$z$得到
            $$z=a\bigg(-\frac{1}{\sqrt2}\cos\theta-\frac{1}{\sqrt6}\sin\theta\bigg)$$

            \item \textbf{化为定积分与计算}
            
            用下标$\theta$表示对$\theta$求导，直接计算有
            $$x_\theta=a\frac{\sqrt2}{\sqrt3}\cos\theta,\quad y_\theta=a\bigg(-\frac{1}{\sqrt2}\sin\theta-\frac{1}{\sqrt6}\cos\theta\bigg),\quad z_\theta=a\bigg(\frac{1}{\sqrt2}\sin\theta-\frac{1}{\sqrt6}\cos\theta\bigg)$$
            从而
            $$\dr s=\sqrt{x_\theta^2+y_\theta^2+z_\theta^2}\dr\theta=a\dr\theta$$
            进一步计算可发现被积函数为$\frac{2}{3}a\cos^2\theta$，由此原积分为
            $$\int_0^{2\pi}\frac{2}{3}a^3\cos^2\theta\dr\theta$$
            类似本讲义19.1.1第4题进行区间拆分，再类似上册教材3.4节例2即可计算出上式为
            $$\frac{8}{3}a^3\int_0^{\pi/2}\cos^2\theta\dr\theta=\frac{2\pi a^3}{3}$$
        \end{itemize}
    }

    \item (习题8.4.2)计算
    $$\iint_S(y+z)\dr S$$
    其中$S$为平面$y+z=1$、$x=2$与三个坐标平面围成的立体的表面。

    \sol{
        分为如下步骤：
        \begin{itemize}
            \item \textbf{区域描述}
            
            我们先分析区域，根据``围成''的有界性要求与围成的区域必然以每个平面为边界，我们可以得到五个约束(分析类似\exref{get ineq})
            $$x\ge0,\quad x\le2,\quad y\ge0,\quad z\ge0,\quad y+z\le1$$
            由于积分区域为表面，需要将每个约束取等，其他取不等号。类似\exref{three column s}，我们将曲面积分分为了如下五个部分的积分之和：
            $$S_1:\quad x=0,\quad x\le2,\quad y\ge0,\quad z\ge0,\quad y+z\le1$$
            $$S_2:\quad x=2,\quad x\ge0,\quad y\ge0,\quad z\ge0,\quad y+z\le1$$
            $$S_3:\quad y=0,\quad x\ge0,\quad x\le2,\quad z\ge0,\quad y+z\le1$$
            $$S_4:\quad z=0,\quad x\ge0,\quad x\le2,\quad y\ge0,\quad y+z\le1$$
            $$S_5:\quad y+z=1,\quad x\ge0,\quad x\le2,\quad y\ge0,\quad z\ge0$$
            代入每个部分满足的条件课化简约束，得到每个部分的化简表述：
            $$S_1:\quad x=0,\quad y\ge0,\quad z\ge0,\quad y+z\le1$$
            $$S_2:\quad x=2,\quad y\ge0,\quad z\ge0,\quad y+z\le1$$
            $$S_3:\quad y=0,\quad 0\le x\le 2,\quad 0\le z\le1$$
            $$S_4:\quad z=0,\quad 0\le x\le 2,\quad 0\le y\le1$$
            $$S_5:\quad y+z=1,\quad 0\le x\le 2,\quad y\ge0,\quad z\ge0$$
            下面分别计算每个部分$y+z$的积分，再求和。

            \item \textbf{参数化与化为重积分1}
            
            由于$x$恒定，可看作$y$、$z$的函数，直接参数化为
            $$\rv=(x,y,z)=(0,u,v)$$
            用下标$u$、$v$表示对$u$、$v$求偏导并保持另一个不变，直接计算有
            $$\rv_u=(x_u,y_u,z_u)=(0,1,0),\quad\rv_v=(x_v,y_v,z_v)=(0,0,1)$$
            代入公式可知
            $$\dr S=\sqrt{|\rv_u|^2|\rv_v|^2-(\rv_u\cdot\rv_v)^2}\dr u\dr v=\dr u\dr v$$
            进一步计算可发现被积函数为$u+v$，约束成为
            $$u\ge0,\quad v\ge0,\quad u+v\le1$$
            由此原积分为
            $$\iint_D(u+v)\dr u\dr v$$
            其中$D$为上方三个约束确定的区域。
            
            \item \textbf{积分计算1}
            
            当$v$任取时，$u$的范围是$[0,1]$，而$u$给定时有$v\in[0,1-u]$，从而可将积分化为累次积分
            $$\int_0^1\dr u\int_0^{1-u}(u+v)\dr v$$
            进一步计算得其为
            $$\int_0^1\bigg(u(1-u)+\frac{1}{2}(1-u)^2\bigg)\dr u=\frac{1}{3}$$
            
            \item \textbf{参数化与化为重积分2}
            
            由于$x$恒定，可看作$y$、$z$的函数，直接参数化为
            $$\rv=(x,y,z)=(2,u,v)$$
            用下标$u$、$v$表示对$u$、$v$求偏导并保持另一个不变，直接计算有
            $$\rv_u=(x_u,y_u,z_u)=(0,1,0),\quad\rv_v=(x_v,y_v,z_v)=(0,0,1)$$
            代入公式可知
            $$\dr S=\sqrt{|\rv_u|^2|\rv_v|^2-(\rv_u\cdot\rv_v)^2}\dr u\dr v=\dr u\dr v$$
            进一步计算可发现被积函数为$u+v$，约束成为
            $$u\ge0,\quad v\ge0,\quad u+v\le1$$
            由此原积分为
            $$\iint_D(u+v)\dr u\dr v$$
            其中$D$为上方三个约束确定的区域。
            
            \item \textbf{积分计算2}
            
            此积分与$S_1$上完全相同，从而结果仍为$\frac{1}{3}$。
            
            \item \textbf{参数化与化为重积分3}     
            
            由于$y$恒定，可看作$x$、$z$的函数，直接参数化为
            $$\rv=(x,y,z)=(u,0,v)$$
            用下标$u$、$v$表示对$u$、$v$求偏导并保持另一个不变，直接计算有
            $$\rv_u=(x_u,y_u,z_u)=(1,0,0),\quad\rv_v=(x_v,y_v,z_v)=(0,0,1)$$
            代入公式可知
            $$\dr S=\sqrt{|\rv_u|^2|\rv_v|^2-(\rv_u\cdot\rv_v)^2}\dr u\dr v=\dr u\dr v$$
            进一步计算可发现被积函数为$v$，约束成为
            $$0\le u\le 2,\quad 0\le v\le1$$
            由此原积分为
            $$\iint_Dv\dr u\dr v$$
            其中$D$为上方两个约束确定的区域。
            
            \item \textbf{积分计算3}
            
            此积分可直接写为累次积分
            $$\int_0^2\dr u\int_0^1v\dr v=\int_0^2\frac{1}{2}\dr u=1$$
            
            \item \textbf{参数化与化为重积分4}
            
            由于$z$恒定，可看作$x$、$y$的函数，直接参数化为
            $$\rv=(x,y,z)=(u,v,0)$$
            用下标$u$、$v$表示对$u$、$v$求偏导并保持另一个不变，直接计算有
            $$\rv_u=(x_u,y_u,z_u)=(1,0,0),\quad\rv_v=(x_v,y_v,z_v)=(0,1,0)$$
            代入公式可知
            $$\dr S=\sqrt{|\rv_u|^2|\rv_v|^2-(\rv_u\cdot\rv_v)^2}\dr u\dr v=\dr u\dr v$$
            进一步计算可发现被积函数为$v$，约束成为
            $$0\le u\le 2,\quad 0\le v\le1$$
            由此原积分为
            $$\iint_Dv\dr u\dr v$$
            其中$D$为上方两个约束确定的区域。
            
            \item \textbf{积分计算4}
            
            此积分与$S_3$上完全相同，从而结果仍为1。
            
            \item \textbf{参数化与化为重积分5}
            
            由于$y+z=1$，可将$z$看作$1-y$成为$x$、$y$的函数，直接参数化为
            $$\rv=(x,y,z)=(u,v,1-v)$$
            用下标$u$、$v$表示对$u$、$v$求偏导并保持另一个不变，直接计算有
            $$\rv_u=(x_u,y_u,z_u)=(1,0,0),\quad\rv_v=(x_v,y_v,z_v)=(0,1,-1)$$
            代入公式可知
            $$\dr S=\sqrt{|\rv_u|^2|\rv_v|^2-(\rv_u\cdot\rv_v)^2}\dr u\dr v=\sqrt2\dr u\dr v$$
            进一步计算可发现被积函数为1，约束成为
            $$0\le u\le 2,\quad v\ge0,\quad 1-v\ge0$$
            由此原积分为
            $$\iint_D\sqrt2\dr u\dr v$$
            其中$D$为上方三个约束确定的区域。
            
            \item \textbf{积分计算5}
            
            此积分可直接写为累次积分
            $$\int_0^2\dr u\int_0^1\sqrt2\dr v=\int_0^2\sqrt2\dr u=2\sqrt2$$
            
            \item \textbf{综合结果}
            
            将五个部分求和可知最终的积分结果
            $$\frac{8}{3}+2\sqrt2$$
        \end{itemize}
    }

    \item (习题8.4.6)计算
    $$\iint_S\big(x(y+z)+z(x+y)\big)\dr S$$
    其中$S$为圆锥面$z=\sqrt{x^2+y^2}$被曲面$x^2+y^2=2ax$所割下的部分，这里参数$a>0$。

    \sol{
        分为如下步骤：
        \begin{itemize}
            \item \textbf{区域描述}
            
            根据``割下''的有界性要求，区域应为曲面$z=\sqrt{x^2+y^2}$中满足$x^2+y^2\le 2ax$的部分(大于等于的部分无界)。

            \item \textbf{参数化}
            
            虽然$z$是$x$、$y$的函数可以提供自然的参数化
            $$\rv=(x,y,z)=(u,v,\sqrt{u^2+v^2})$$
            由于曲面的旋转对称性，为了让计算更简洁可以考虑\textbf{极坐标}参数化
            $$\rv=(x,y,z)=(u\cos v,u\sin v,u)$$
            代入约束可得到参数范围
            $$u^2\le 2au\cos v$$
            将上述条件与极坐标自然范围确定的$(u,v)$区域记为$D$。

            \item \textbf{化为重积分}
            
            用下标$u$、$v$表示对$u$、$v$求偏导并保持另一个不变，直接计算有
            $$\rv_u=(x_u,y_u,z_u)=(\cos v,\sin v,1),\quad\rv_v=(x_v,y_v,z_v)=(-u\sin v,u\cos v,0)$$
            代入公式可知(利用极坐标自然范围可知$u\ge0$，可消除绝对值)
            $$\dr S=\sqrt{|\rv_u|^2|\rv_v|^2-(\rv_u\cdot\rv_v)^2}\dr u\dr v=\sqrt2|u|\dr u\dr v=\sqrt2u\dr u\dr v$$
            进一步计算可发现被积函数为
            $$u^2(\cos v\sin v+2\cos v+\sin v)$$
            由此原积分为
            $$\iint_D\sqrt2u^3(\cos v\sin v+2\cos v+\sin v)\dr u\dr v$$

            \item \textbf{积分计算}
            
            利用$u$的自然范围$u\ge0$，排除$u=0$的特殊点可将对$u$的约束化为
            $$0\le u\le 2a\cos v$$
            $u$可任取时，$v$需满足$\cos v\ge0$，结合自然范围可得到
            $$v\in\bigg[0,\frac{\pi}{2}\bigg]\cup\bigg[\frac{3\pi}{2},2\pi\bigg]$$
            从而可将积分化为累次积分
            $$\begin{aligned}&\int_0^{\pi/2}\dr v\int_0^{2a\cos v}\sqrt2u^3(\cos v\sin v+2\cos v+\sin v)\dr u\\+&\int_{3\pi/2}^{2\pi}\dr v\int_0^{2a\cos v}\sqrt2u^3(\cos v\sin v+2\cos v+\sin v)\dr u\end{aligned}$$
            对$u$积分得
            $$\begin{aligned}&4\sqrt2a^4\int_0^{\pi/2}\cos^4v(\cos v\sin v+2\cos v+\sin v)\dr v\\+&4\sqrt2a^4\int_{3\pi/2}^{2\pi}\cos^4v(\cos v\sin v+2\cos v+\sin v)\dr v\end{aligned}$$
            先利用对称性化简，在第二个部分中将$v$换元为$2\pi-v$可计算得到上式为
            $$\begin{aligned}&4\sqrt2a^4\int_0^{\pi/2}\cos^4v(\cos v\sin v+2\cos v+\sin v)\dr v\\+&4\sqrt2a^4\int_0^{\pi/2}\cos^4v(-\cos v\sin v+2\cos v-\sin v)\dr v\end{aligned}$$
            从而求和得其为
            $$16\sqrt2a^4\int_0^{\pi/2}\cos^5v\dr v$$
            将$v$换元为$\frac{\pi}{2}-v$可知被积函数换为$\sin^5v$时结果相同，由上册教材3.4节例2可得最终结果为
            $$16\sqrt2a^4\cdot\frac{4}{5}\cdot\frac{2}{3}=\frac{128\sqrt2a^4}{15}$$
        \end{itemize}

        \note 本题也可以第一种方式进行参数化，计算重积分时再换元，效果类似下一题的过程。
    }

    \item (习题8.4.9)求抛物面壳
    $$z=\frac{1}{2}(x^2+y^2),\quad z\in[0,1]$$
    的质量，其在$(x,y,z)$的面密度为$x+y+z$。

    \sol{
        根据曲面积分几何意义，质量即为
        $$\iint_S(x+y+z)\dr S$$
        其中$S$表示条件中的抛物面壳。计算分为如下步骤：
        \begin{itemize}
            \item \textbf{参数化}
            
            $z$是$x$、$y$的函数可以提供自然的参数化
            $$\rv=(x,y,z)=\bigg(u,v,\frac{1}{2}(u^2+v^2)\bigg)$$
            代入约束可得到参数范围
            $$0\le u^2+v^2\le 2$$
            将上述条件确定的$(u,v)$区域记为$D$。

            \item \textbf{化为重积分}
            
            用下标$u$、$v$表示对$u$、$v$求偏导并保持另一个不变，直接计算有
            $$\rv_u=(x_u,y_u,z_u)=(1,0,u),\quad\rv_v=(x_v,y_v,z_v)=(0,1,v)$$
            代入公式可知
            $$\dr S=\sqrt{|\rv_u|^2|\rv_v|^2-(\rv_u\cdot\rv_v)^2}\dr u\dr v=\sqrt{1+u^2+v^2}\dr u\dr v$$
            进一步计算可发现被积函数为
            $$u+v+\frac{1}{2}(u^2+v^2)$$
            由此原积分为
            $$\iint_D\bigg(u+v+\frac{1}{2}(u^2+v^2)\bigg)\sqrt{1+u^2+v^2}\dr u\dr v$$

            \item \textbf{积分计算}
            
            先进行反射对称性化简。由于$(u,v)\in D$当且仅当$(-u,v)\in D$，考虑函数折半可发现
            $$\iint_Du\sqrt{1+u^2+v^2}\dr u\dr v=0$$
            由于$(u,v)\in D$当且仅当$(u,-v)\in D$，考虑函数折半可发现
            $$\iint_Dv\sqrt{1+u^2+v^2}\dr u\dr v=0$$
            这就将原积分化为了
            $$\iint_D\frac{1}{2}(u^2+v^2)\sqrt{1+u^2+v^2}\dr u\dr v$$
            考虑到区域与函数的旋转对称性，采用极坐标换元
            $$u=r\cos\theta,\quad v=r\sin\theta$$
            此时有
            $$\dr u\dr v=r\dr r\dr\theta$$
            约束化简为
            $$0\le r\le\sqrt2$$
            被积函数为$\frac{1}{2}r^2\sqrt{1+r^2}$，再结合极坐标自然范围可将积分化为
            $$\int_0^{\sqrt2}\dr r\int_0^{2\pi}\frac{1}{2}r^3\sqrt{1+r^2}\dr\theta$$
            对$\theta$积分得到
            $$\pi\int_0^{\sqrt2}r^3\sqrt{1+r^2}\dr r$$
            换元$s=r^2$即得
            $$\frac{\pi}{2}\int_0^2s\sqrt{1+s}\dr s$$
            为了处理难于处理的部分$\sqrt{1+s}$，进一步换元$t=s+1$可将积分化为
            $$\frac{\pi}{2}\int_1^3(t\sqrt{t}-\sqrt{t})\dr t=\frac{\pi}{2}\bigg(\frac{2}{5}\cdot 3^{5/2}-\frac{2}{5}-\frac{2}{3}\cdot3^{3/2}+\frac{2}{3}\bigg)=\bigg(\frac{4\sqrt3}{5}+\frac{2}{15}\bigg)\pi$$

        \end{itemize}
    }

    \item (习题8.2.4)计算
    $$\int_L(x^2-2xy)\dr x+(y^2-2xy)\dr y$$
    其中$L$为
    \begin{enumerate}[(1)]
        \item 曲线$y=x^4$上点$(-1,1)$到点$(1,1)$的一段。
        
        \sol{
            分为如下步骤：
            \begin{itemize}
                \item \textbf{参数化}
                
                根据函数式可直接得到参数方程
                $$\begin{cases}x=t\\y=t^4\end{cases}(-1\le t\le1)$$

                \item \textbf{化为定积分与计算}
                
                用下标$t$表示对$t$求导，直接计算有
                $$x_t=1,\quad y_t=4t^3$$
                代入$\dr x=x_t\dr t$、$\dr y=y_t\dr t$与$x$、$y$的参数方程可得原积分为
                $$\int_{-1}^1(t^2-2t^5+4t^{11}-8t^8)\dr t=\frac{2}{3}-\frac{16}{9}=-\frac{10}{9}$$
            \end{itemize}
        }

        \item 点$(-1,1)$到点$(1,1)$的直线段。
        
        \sol{
            分为如下步骤：
            \begin{itemize}
                \item \textbf{参数化}
                
                根据直线方程可直接得到参数方程
                $$\begin{cases}x=t\\y=1\end{cases}(-1\le t\le1)$$

                \item \textbf{化为定积分与计算}
                
                用下标$t$表示对$t$求导，直接计算有
                $$x_t=1,\quad y_t=0$$
                代入$\dr x=x_t\dr t$、$\dr y=y_t\dr t$与$x$、$y$的参数方程可得原积分为
                $$\int_{-1}^1(t^2-2t)\dr t=\frac{2}{3}$$
            \end{itemize}
        }
    \end{enumerate}

    \item (习题8.2.8)计算
    $$\int_L(x^4-z^2)\dr x+2xy^2\dr y-y\dr z$$
    其中$L$为(沿$t$增加方向)
    $$\begin{cases}x=t\\y=t^2\\z=t^3\end{cases}(0\le t\le1)$$

    \sol{
        用下标$t$表示对$t$求导，直接计算有
        $$x_t=1,\quad y_t=2t,\quad z_t=3t^2$$
        代入$\dr x=x_t\dr t$、$\dr y=y_t\dr t$、$\dr z=z_t\dr t$与$x$、$y$、$z$的参数方程可得原积分为
        $$\int_0^1(t^4-t^6+4t^6-3t^4)\dr t=-\frac{2}{5}+\frac{3}{7}=\frac{1}{35}$$
    }

    \item (习题8.2.10)计算
    $$\oint_Lx\dr x+z\dr y+y\dr z$$
    其中闭曲线L由下列三曲线段(沿$t$增加方向)围成：
    $$\begin{cases}x=\cos t\\y=\sin t\\z=t\end{cases}\bigg(0\le t\le\frac{\pi}{2}\bigg),\quad\begin{cases}x=0\\y=1\\z=\frac{\pi}{2}(1-t)\end{cases}(0\le t\le1),\quad\begin{cases}x=t\\y=1-t\\z=0\end{cases}(0\le t\le1)$$

    \sol{
        分为如下步骤：
        \begin{itemize}
            \item \textbf{第一段化为定积分与计算}
            
            用下标$t$表示对$t$求导，直接计算有
            $$x_t=-\sin t,\quad y_t=\cos t,\quad z_t=1$$
            代入$\dr x=x_t\dr t$、$\dr y=y_t\dr t$、$\dr z=z_t\dr t$与$x$、$y$、$z$的参数方程可得原积分为
            $$\int_0^{\pi/2}(-\cos t\sin t+t\cos t+\sin t)\dr t$$
            直接计算有(第二部分利用了分部积分)
            $$\int_0^{\pi/2}-\cos t\sin t\dr t=-\frac{1}{2}\int_0^{\pi/2}\sin(2t)\dr t=-\frac{1}{2}$$
            $$\int_0^{\pi/2}t\cos t\dr t=\int_0^{\pi/2}t\dr\sin t=\frac{\pi}{2}-0-\int_0^{\pi/2}\sin t=\frac{\pi}{2}-1$$
            $$\int_0^{\pi/2}\sin t\dr t=1$$
            从而求和得第一段的积分结果为
            $$\frac{\pi}{2}-\frac{1}{2}$$
            
            \item \textbf{第二段化为定积分与计算}
            
            用下标$t$表示对$t$求导，直接计算有
            $$x_t=0,\quad y_t=0,\quad z_t=-\frac{\pi}{2}$$
            代入$\dr x=x_t\dr t$、$\dr y=y_t\dr t$、$\dr z=z_t\dr t$与$x$、$y$、$z$的参数方程可得原积分为
            $$\int_0^1-\frac{\pi}{2}\dr t$$
            从而第二段的积分结果为
            $$-\frac{\pi}{2}$$
            
            \item \textbf{第三段化为定积分与计算}
            
            用下标$t$表示对$t$求导，直接计算有
            $$x_t=1,\quad y_t=-1,\quad z_t=0$$
            代入$\dr x=x_t\dr t$、$\dr y=y_t\dr t$、$\dr z=z_t\dr t$与$x$、$y$、$z$的参数方程可得原积分为
            $$\int_0^1t\dr t$$
            从而第三段的积分结果为
            $$\frac{1}{2}$$
            
            \item \textbf{综合结果}
            
            将三个部分求和可知最终的积分结果为0。
        \end{itemize}
    }

    \item (习题8.2.13)设力场$\Fv(x,y,z)=(y,-x,x+y+z)$，求：
    \begin{enumerate}[(1)]
        \item 质点沿螺旋线$L_1$
        $$\begin{cases}x=a\cos t\\y=a\sin t\\z=\frac{c}{2\pi}t\end{cases}(0\le t\le2\pi)$$
        从$A(a,0,0)$运动到$B(a,0,c)$时力场做的功。

        \sol{
            根据第二型曲线积分的物理意义，做功即为
            $$\int_{L_1}\Fv(x,y,z)\cdot\dr\rv$$
            计算分为如下步骤：
            \begin{itemize}
                \item \textbf{化为定积分}
                
                用下标$t$表示对$t$求导，直接计算有
                $$x_t=-a\sin t,\quad y_t=a\cos t,\quad z_t=\frac{c}{2\pi}$$
                代入$\dr x=x_t\dr t$、$\dr y=y_t\dr t$、$\dr z=z_t\dr t$与$x$、$y$、$z$的参数方程可得原积分为
                $$\int_0^{2\pi}\bigg(-a^2+\frac{c}{2\pi}\bigg(a\cos t+a\sin t+\frac{c}{2\pi}t\bigg)\bigg)\dr t$$

                \item \textbf{积分计算}
                
                利用对称性可知$\cos t$与$\sin t$的积分均为0，由此积分化为
                $$\int_0^{2\pi}\bigg(-a^2+\frac{c^2}{4\pi^2}t\bigg)\dr t=-2\pi a^2+\frac{c^2}{2}$$
            \end{itemize}
        }

        \item 质点按(1)中的线段$AB$从$A$运动到$B$时力场做的功。
        
        \sol{
            根据第二型曲线积分的物理意义，做功即为
            $$\int_{L_2}\Fv(x,y,z)\cdot\dr\rv$$
            $L_2$代表从$A$到$B$的线段$AB$。计算分为如下步骤：
            \begin{itemize}
                \item \textbf{参数化}
                
                根据直线方程可直接得到参数方程
                $$\begin{cases}x=a\\y=0\\z=t\end{cases}(0\le t\le c)$$

                \item \textbf{化为定积分与计算}
                
                用下标$t$表示对$t$求导，直接计算有
                $$x_t=0,\quad y_t=0,\quad z_t=1$$
                代入$\dr x=x_t\dr t$、$\dr y=y_t\dr t$、$\dr z=z_t\dr t$与$x$、$y$、$z$的参数方程可得原积分为
                $$\int_0^c(a+t)\dr t=ac+\frac{1}{2}c^2$$
            \end{itemize}
        }
    \end{enumerate}

    \item (总练习题8.9)计算
    $$\oint_{L^+}\sqrt{x^2+y^2}\dr x+y\big(xy+\ln(x+\sqrt{x^2+y^2})\big)\dr y$$
    其中$L$为圆周$x^2+y^2=1$。

    \note 本题实际上是一个$(-1,0)$点为无穷的\textbf{广义积分}，我们只形式计算结果，不考虑敛散性问题。
    
    \sol{
        这里我们先介绍直接计算的方案，之后将介绍如何化简计算。分为如下步骤：
        \begin{itemize}
            \item \textbf{参数化与定向}
            
            利用极坐标可以得到圆的参数方程
            $$\begin{cases}x=\cos\theta\\y=\sin\theta\end{cases}(0\le\theta\le2\pi)$$
            为了确定定向，用下标$\theta$代表对$\theta$求导，直接计算有
            $$x_\theta=-\sin\theta,\quad y_\theta=\cos\theta$$
            我们\textbf{任取}一个使得$x_\theta$、$y_\theta$\textbf{不全为0}的点，如取$\theta=\pi$，此处有$(x,y)=(-1,0)$、$(x_\theta,y_\theta)=(0,-1)$，在曲线上的点$(x,y)$处画出向量$(x_\theta,y_\theta)$，可以发现它指向逆时针方向，与要求的定向相同，从而化为定积分无需加负号。

            \note 直观来说，不全为0是为了画出的$(x_\theta,y_\theta)$不是零向量，能看出方向。

            \note 这段推导的严谨逻辑见本讲义21.2.3。事实上，这里的``画出''也可以用更加代数的处理：考虑取值范围可知$(-1,0)$是曲线左侧端点，而$(0,-1)$代表向下，在左侧端点向下走必然为逆时针。

            \item \textbf{化为定积分与计算}
                
            代入$\dr x=x_\theta\dr\theta$、$\dr y=y_\theta\dr\theta$与$x$、$y$的参数方程可得原积分为
            $$\int_0^{2\pi}(-\sin\theta+\sin^2\theta\cos^2\theta+\sin\theta\cos\theta\ln(1+\cos\theta))\dr\theta$$
            先利用对称性化简。将$\theta$换元为$2\pi-\theta$可将上式化为
            $$\int_0^{2\pi}(\sin\theta+\sin^2\theta\cos^2\theta-\sin\theta\cos\theta\ln(1+\cos\theta))\dr\theta$$
            两式求和并平均即将积分化简为了(换元$t=2\theta$)
            $$\int_0^{2\pi}\sin^2\theta\cos^2\theta\dr\theta=\frac{1}{4}\int_0^{2\pi}\sin^2(2\theta)\dr\theta=\frac{1}{8}\int_0^{4\pi}\sin^2t\dr t$$
            与本讲义20.1.2第7题相同计算可知结果为
            $$\int_0^{\pi/2}\sin^2t\dr t=\frac{\pi}{4}$$
        \end{itemize}

        \note 极坐标与广义极坐标参数化后的曲线是\textbf{逆时针定向}的，建议熟悉此性质。
    }

    \item (总练习题8.11)求两积分
    $$I_1=\iint_{S_1}(x^2+y^2+z^2)\dr S,\quad I_2=\iint_{S_2}(x^2+y^2+z^2)\dr S$$
    之差，其中$S_1$为球面$x^2+y^2+z^2=a^2$，$S_2$为内接于此球的八面体$|x|+|y|+|z|=a$，这里参数$a>0$。
    
    \sol{
        分为如下步骤：
        \begin{itemize}
            \item \textbf{$I_1$参数化}
            
            利用球坐标可得到球面的一个参数化
            $$\rv=(x,y,z)=(a\cos u\sin v,a\sin u\sin v,a\cos v)$$
            球坐标自然范围即为参数范围
            $$0\le u\le2\pi,\quad0\le v\le\pi$$
            将上述条件确定的$(u,v)$区域记为$D$。

            \item \textbf{$I_1$化为重积分}
            
            用下标$u$、$v$表示对$u$、$v$求偏导并保持另一个不变，直接计算有
            $$\rv_u=(x_u,y_u,z_u)=(-a\sin u\sin v,a\cos u\sin v,0)$$
            $$\rv_v=(x_v,y_v,z_v)=(a\cos u\cos v,a\sin u\cos v,-a\sin v)$$
            代入公式可知(由$v$范围可消除$\sin v$绝对值)
            $$\dr S=\sqrt{|\rv_u|^2|\rv_v|^2-(\rv_u\cdot\rv_v)^2}\dr u\dr v=a^2\sin v\dr u\dr v$$
            进一步计算可发现被积函数为
            $$a^2$$
            由此原积分为
            $$\iint_Da^4\sin v\dr u\dr v$$
            
            \item \textbf{$I_1$计算}
            
            此积分可直接写为累次积分
            $$\int_0^{2\pi}\dr u\int_0^\pi a^4\sin v\dr v=2a^4\int_0^{2\pi}\dr u=4\pi a^4$$

            \item \textbf{$I_2$区域描述与参数化}
            
            根据绝对值的不同符号，$S_2$可以在八个卦限中分为八个部分。我们接下来先计算$x\ge0$、$y\ge0$、$z\ge0$部分的结果，此时积分的曲面为
            $$S_2^{+++}:\quad x+y+z=a,\quad x\ge0,\quad y\ge0,\quad z\ge0$$

            由于可将$z$看作$x$、$y$的函数，存在自然的参数化
            $$\vec{r}=(x,y,z)=(u,v,a-u-v)$$
            代入约束可得到参数范围
            $$u\ge0,\quad v\ge0,\quad u+v\le a$$
            将上述条件确定的$(u,v)$区域记为$D$。
            
            \item \textbf{$I_2$化为重积分}
            
            用下标$u$、$v$表示对$u$、$v$求偏导并保持另一个不变，直接计算有
            $$\rv_u=(x_u,y_u,z_u)=(1,0,-1)$$
            $$\rv_v=(x_v,y_v,z_v)=(0,1,-1)$$
            代入公式可知
            $$\dr S=\sqrt{|\rv_u|^2|\rv_v|^2-(\rv_u\cdot\rv_v)^2}\dr u\dr v=\sqrt3\dr u\dr v$$
            进一步计算可发现被积函数为
            $$u^2+v^2+(a-u-v)^2$$
            由此$S_2^{+++}$上的积分为
            $$\iint_D\sqrt3(u^2+v^2+(a-u-v)^2)\dr u\dr v$$
            
            \item \textbf{$I_2$计算}
            
            当$v$任取时，$u$的范围是$[0,a]$，而$u$给定时有$v\in[0,a-u]$，从而可将积分化为累次积分
            $$\int_0^a\dr u\int_0^{a-u}\sqrt3(u^2+v^2+(a-u-v)^2)\dr v$$
            对$v$积分得到
            $$\sqrt3\int_0^a\bigg(u^2(a-u)+\frac{2}{3}(a-u)^3\bigg)\dr u$$
            进一步计算得其为
            $$\sqrt3\bigg(\frac{1}{3}a^4-\frac{1}{4}a^4+\frac{1}{6}a^4\bigg)=\frac{\sqrt3}{4}a^4$$
            对$I_2$的剩余7个卦限的部分讨论可以发现能化为形式完全相同的重积分，从而积分结果最终为
            $$8\cdot\frac{\sqrt3}{4}a^4=2\sqrt3a^4$$

            \note 这部分说明也可以利用对称性进行，分析过程与三重积分的对称性完全类似，详见本讲义21.3.3。
            
            \item \textbf{综合结果}
            
            直接作差可知
            $$I_1-I_2=(4\pi-2\sqrt3)a^4$$
        \end{itemize}
    }
\end{enumerate}
\subsubsection{补充题解答}
\begin{enumerate}
    \item (习题7.4.2)求锥面$z=\sqrt{x^2+y^2}$被柱面$z^2=2x$割下部分的曲面面积。
    
    \sol{
        分为如下步骤：
        \begin{itemize}
            \item \textbf{区域描述}
            
            根据``割下''的有界性要求，区域应为曲面$z=\sqrt{x^2+y^2}$中满足$z^2\le 2x$的部分(大于等于的部分无界)，代入曲面的表达式可得
            $$x^2+y^2\le2x$$

            将满足以上条件的$(x,y)$区域记为$D$，由于$z$可看作$x$、$y$的函数，直接代入公式可知面积为
            $$\iint_D\sqrt{1+\bigg(\frac{\partial\sqrt{x^2+y^2}}{\partial x}\bigg)^2+\bigg(\frac{\partial\sqrt{x^2+y^2}}{\partial y}\bigg)^2}\dr x\dr y$$
            化简可知积分即
            $$\iint_D\sqrt2\dr x\dr y$$

            \item \textbf{积分计算}
            
            \note 本题也可以用之前的方法进行极坐标换元$x=1+r\cos\theta$、$y=r\sin\theta$后逐步计算，这里提供一个更简单的方法。
            
            利用重积分的几何意义，$\iint_D\dr x\dr y$即为区域的面积，而区域$(x-1)^2+y^2\le1$是一个半径为1的圆，面积为$\pi$，从而积分结果是
            $$\sqrt2\iint_D\dr x\dr y=\sqrt2\pi$$
        \end{itemize}
    }
    
    \item (习题7.4.6)给定参数$R>0$，设球体$x^2+y^2+z^2\le 2Rz$上任一点$(x,y,z)$处的体密度等于$x^2+y^2+z^2$，求该球体的质心坐标。
    
    \sol{
        利用质心坐标公式，设质心坐标为$(x_0,y_0,z_0)=\frac{1}{m}(x_1,y_1,z_1)$，应满足
        $$x_1=\iiint_\Omega x(x^2+y^2+z^2)\dr V,\quad y_1=\iiint_\Omega y(x^2+y^2+z^2)\dr V,\quad z_1=\iiint_\Omega z(x^2+y^2+z^2)\dr V$$
        $$m=\iiint_\Omega(x^2+y^2+z^2)\dr V$$

        分为如下步骤：
        \begin{itemize}
            \item \textbf{反射对称化简}
            
            由于区域$\Omega$满足$(x,y,z)\in\Omega$当且仅当$(-x,y,z)\in\Omega$，且
            $$x(x^2+y^2+z^2)=-(-x)((-x)^2+y^2+z^2)$$
            于是利用函数折半可发现
            $$x_1=0$$
            同理可得$y_1=0$，只需计算$z_1$与$m$。
            
            \item \textbf{换元与累次积分化}
            
            由于区域为球体
            $$x^2+y^2+(z-R)^2\le R^2$$
            由此考虑平移后的球坐标换元(以$(0,0,R)$为球坐标中心)
            $$x=\rho\sin\varphi\cos\theta,\quad y=\rho\sin\varphi\sin\theta,\quad z=\rho\cos\varphi+R$$
            此时$m$的被积函数为
            $$\rho^2+2R\rho\cos\varphi+R^2$$
            $z_1$的被积函数为
            $$(\rho^2+2R\rho\cos\varphi+R^2)(\rho\cos\varphi+R)$$
            且直接计算Jacobi行列式可知(由于加常数不影响求导，Jacobi行列式应与球坐标完全相同)
            $$\dr x\dr y\dr z=\rho^2\sin\varphi\dr\rho\dr\varphi\dr\theta$$
            换元后约束条件即化为了
            $$\rho^2\le R^2$$
            结合球坐标自然范围可知
            $$0\le\rho\le R,\quad0\le\varphi\le\pi,\quad0\le\theta\le2\pi$$
            这已经是符合累次积分要求的范围，从而可写出累次积分
            $$m=\int_0^R\dr\rho\int_0^\pi\dr\varphi\int_0^{2\pi}(\rho^2+2R\rho\cos\varphi+R^2)\rho^2\sin\varphi\dr\theta$$
            $$z_1=\int_0^R\dr\rho\int_0^\pi\dr\varphi\int_0^{2\pi}(\rho^2+2R\rho\cos\varphi+R^2)\rho^2\sin\varphi(\rho\cos\varphi+R)\dr\theta$$
            
            \item \textbf{积分计算}
            
            对$\theta$积分可得
            $$m=2\pi\int_0^R\dr\rho\int_0^\pi(\rho^2+2R\rho\cos\varphi+R^2)\rho^2\sin\varphi\dr\varphi$$
            $$z_1=2\pi\int_0^R\dr\rho\int_0^\pi(\rho^2+2R\rho\cos\varphi+R^2)\rho^2\sin\varphi(\rho\cos\varphi+R)\dr\varphi$$
            利用对称性，在上下的定积分中换元$\pi-\theta$，则$\sin$不变，$\cos$增加负号，可发现
            $$m=2\pi\int_0^R\dr\rho\int_0^\pi(\rho^2-2R\rho\cos\varphi+R^2)\rho^2\sin\varphi\dr\varphi$$
            $$z_1=2\pi\int_0^R\dr\rho\int_0^\pi(\rho^2-2R\rho\cos\varphi+R^2)\rho^2\sin\varphi(-\rho\cos\varphi+R)\dr\varphi$$
            相加并平均可得
            $$m=2\pi\int_0^R\dr\rho\int_0^\pi(\rho^2+R^2)\rho^2\sin\varphi\dr\varphi$$
            $$z_1=2\pi\int_0^R\dr\rho\int_0^\pi(2R\rho^4\sin\varphi\cos^2\varphi+(\rho^2+R^2)R\rho^2\sin\varphi)\dr\varphi$$
            由此可直接算出
            $$m=4\pi\int_0^R(\rho^2+R^2)\rho^2\dr\rho=\frac{32}{15}\pi R^5$$
            而根据
            $$\int_0^\pi\sin\varphi\cos^2\varphi\dr\varphi=\int_0^\pi\cos^2\varphi\dr\cos\varphi=\frac{2}{3}$$
            可知
            $$z_1=\frac{4\pi}{3}\int_0^R2R\rho^4\dr\rho+4\pi\int_0^R(\rho^2+R^2)R\rho^2\dr\rho=\frac{40}{15}\pi R^6$$
            
            \item \textbf{综合结果}
            
            根据以上的计算结果可知质心坐标
            $$\bigg(0,0,\frac{5}{4}R\bigg)$$
        \end{itemize}
    }
    
    \item (习题7.4.12)设引力常数为$k$，求高为$h$、顶角为$2\alpha$、密度为$\rho$的均匀圆锥体对位于其顶点的单位质点的引力。这里参数$0<\alpha<\frac{\pi}{2}$、$h>0$。
    
    \sol{
        设锥体顶点为坐标原点$O$，与地面圆心$A$的连线方向$\overrightarrow{OA}$为$z$轴正方向，$x$、$y$轴方向任取。将圆锥体记为$\Omega$，利用几何关系可以发现其满足
        $$\Omega=\{(x,y,z)\mid\sqrt{x^2+y^2}\le z\tan\alpha,\ 0\le z\le h\}$$

        利用引力公式，设引力为$(F_x,F_y,F_z)$，应满足
        $$F_x=k\iiint_\Omega\frac{x}{(x^2+y^2+z^2)^{3/2}}\rho\dr V$$
        $$F_y=k\iiint_\Omega\frac{y}{(x^2+y^2+z^2)^{3/2}}\rho\dr V$$
        $$F_z=k\iiint_\Omega\frac{z}{(x^2+y^2+z^2)^{3/2}}\rho\dr V$$
        由于$(x,y,z)\in\Omega$当且仅当$(-x,y,z)\in\Omega$，利用函数折半可得$F_x=0$，同理$F_y=0$，只需计算$F_z$。

        计算分为如下步骤：
        \begin{itemize}
            \item \textbf{换元与累次积分化}
            
            根据区域的旋转对称性，我们考虑柱坐标换元
            $$x=r\cos\theta,\quad y=r\sin\theta,\quad z=z$$
            这样约束条件即化为了
            $$r\le z\tan\alpha,\quad 0\le z\le h$$
            结合柱坐标的自然范围可得
            $$0\le z\le h,\quad0\le r\le z\tan\alpha,\quad 0\le\theta\le2\pi$$
            将上述条件确定的$(r,\theta,z)$区域记为$\Omega_0$，被积函数可写为
            $$\frac{z}{(r^2+z^2)^{3/2}}$$
            直接计算Jacobi行列式可知
            $$\dr x\dr y\dr z=r\dr r\dr\theta\dr z$$
            由此可写出累次积分
            $$F_z=k\rho\int_0^h\dr z\int_0^{z\tan\alpha}\dr r\int_0^{2\pi}\frac{zr}{(r^2+z^2)^{3/2}}\dr\theta$$
            
            \item \textbf{积分计算}
            
            对$\theta$积分可知
            $$F_z=2\pi k\rho\int_0^h\dr z\int_0^{z\tan\alpha}\frac{zr}{(r^2+z^2)^{3/2}}\dr r$$
            由于分母形式，考虑三角换元$r=z\tan\phi$，由$\alpha$的范围可知换元后$\phi\in[0,\alpha]$，可得
            $$F_z=2\pi k\rho\int_0^h\dr z\int_0^\alpha\sin\phi\dr\phi$$
            进一步计算即得
            $$F_z=2\pi k\rho\int_0^h(1-\cos\alpha)\dr z=2\pi k\rho h(1-\cos\alpha)$$
        \end{itemize}
        综合$F_x$、$F_y$、$F_z$可得结果。

        \note 本题结果与教材答案不同，是因为这里的建系$z$方向与教材相反。
    }
    
    \item (总练习题7.10)求极坐标下的心脏线$r=1+\sin\theta$所围区域中位于第一象限部分的面积。
    
    \sol{
        设区域为$D$根据重积分的几何意义知面积为
        $$\iint_D1\dr x\dr y$$
        计算分为如下步骤：
        \begin{itemize}
            \item \textbf{区域描述与换元}
            
            我们将条件确定的$(r,\theta)$区域记为$D_0$\ (注意它与积分区域$D$不同，因为$D$是针对$x$、$y$的)，根据``围成''的有界性要求与第一象限对应$\theta\in[0,\frac{\pi}{2}]$，它是满足
            $$0\le\theta\le\frac{\pi}{2},\quad0\le r\le1+\sin\theta$$
            的$(r,\theta)$集合。

            计算极坐标Jacobi行列式得$\dr x\dr y=r\dr r\dr\theta$，原积分化为
            $$\iint_{D_0}r\dr r\dr\theta$$

            \item \textbf{累次积分化}
            
            由于区域描述已经符合累次积分的要求，可以直接写出累次积分
            $$\int_0^{\pi/2}\dr\theta\int_0^{1+\sin\theta}r\dr r$$
            
            \item \textbf{积分计算}
            
            对$r$积分可得到积分化为
            $$\frac{1}{2}\int_0^{\pi/2}(1+\sin\theta)^2\dr\theta$$
            利用上册教材3.4节例2可得结果为
            $$\frac{1}{2}\bigg(\frac{\pi}{2}+2+\frac{\pi}{4}\bigg)=\frac{3}{8}\pi+1$$
        \end{itemize}
    }
    
    \item (习题8.1.6)计算
    $$\int_Lxy\dr s$$
    其中$L$是椭圆周$\frac{x^2}{a^2}+\frac{y^2}{b^2}=1$位于第一象限中的部分，这里参数$a>0$、$b>0$。

    \sol{
        分为如下步骤：
        \begin{itemize}
            \item \textbf{参数化}
            
            由其为椭圆，利用\textbf{广义极坐标}可参数化为
            $$x=a\cos\theta,\quad y=b\sin\theta$$
            由第一象限可知范围$0\le\theta\le\frac{\pi}{2}$。

            \item \textbf{化为定积分}

            用下标$\theta$表示对$\theta$求导，直接计算有
            $$x_\theta=-a\sin\theta,\quad y_\theta=b\cos\theta$$
            从而
            $$\dr s=\sqrt{x_\theta^2+y_\theta^2}\dr\theta=\sqrt{a^2\sin^2\theta+b^2\cos^2\theta}\dr\theta$$
            进一步计算可发现被积函数为$ab\sin\theta\cos\theta$，由此原积分为
            $$ab\int_0^{\pi/2}\sin\theta\cos\theta\sqrt{a^2\sin^2\theta+b^2\cos^2\theta}\dr\theta$$

            \item \textbf{积分计算}
            
            由于$\cos\theta\dr\theta=\dr\sin\theta$，且根号中的$\cos^2\theta$可以写为$1-\sin^2\theta$，换元$t=\sin\theta$可将积分化为
            $$ab\int_0^1t\sqrt{b^2+(a^2-b^2)t^2}\dr t$$
            再换元$s=t^2$即得其为
            $$\frac{ab}{2}\int_0^1\sqrt{b^2+(a^2-b^2)s}\dr s$$
            这对$s$是幂函数复合一次函数，可直接写出原函数得到积分，当$a^2-b^2=0$时积分为$\frac{ab^2}{2}$，当$a^2-b^2\ne0$时积分为
            $$\frac{ab}{2}\cdot\frac{2}{3}\cdot\frac{1}{a^2-b^2}(a^3-b^3)$$
            化简可得到两种情况统一为
            $$\frac{ab(a^2+ab+b^2)}{3(a+b)}$$
        \end{itemize}
    }

    \item (习题8.1.9)若椭圆周$\frac{x^2}{a^2}+\frac{y^2}{b^2}=1$上任何一点$(x,y)$处的线密度为$|y|$，求其质量，这里参数$0<b<a$。
    
    \sol{
        分为如下步骤：
        \begin{itemize}
            \item \textbf{参数化}
            
            由其为椭圆，利用\textbf{广义极坐标}可参数化为
            $$\begin{cases}x=a\cos\theta\\y=b\sin\theta\end{cases}(0\le\theta\le2\pi)$$

            \item \textbf{化为定积分}

            用下标$\theta$表示对$\theta$求导，直接计算有
            $$x_\theta=-a\sin\theta,\quad y_\theta=b\cos\theta$$
            从而
            $$\dr s=\sqrt{x_\theta^2+y_\theta^2}\dr\theta=\sqrt{a^2\sin^2\theta+b^2\cos^2\theta}\dr\theta$$
            进一步计算可发现被积函数为$b|\sin\theta|$，由此原积分为
            $$b\int_0^{2\pi}|\sin\theta|\sqrt{a^2\sin^2\theta+b^2\cos^2\theta}\dr\theta$$

            \item \textbf{积分计算}
            
            将$\theta$换元为$2\pi-\theta$可计算得到
            $$b\int_0^\pi|\sin\theta|\sqrt{a^2\sin^2\theta+b^2\cos^2\theta}\dr\theta=b\int_\pi^{2\pi}|\sin\theta|\sqrt{a^2\sin^2\theta+b^2\cos^2\theta}\dr\theta$$
            从而可将原积分化为
            $$2b\int_0^\pi|\sin\theta|\sqrt{a^2\sin^2\theta+b^2\cos^2\theta}\dr\theta$$
            同理，再将$\theta$换元为$\pi-\theta$可进一步化为
            $$4b\int_0^{\pi/2}|\sin\theta|\sqrt{a^2\sin^2\theta+b^2\cos^2\theta}\dr\theta$$

            \note 这些对称性技巧详见上学期讲义7.3.3，与重积分类似。

            此时$\sin\theta$已经恒正，可去除绝对值，从而类似上题换元$t=\sin\theta$可得积分为
            $$4b\int_0^1\sqrt{a^2-(a^2-b^2)t^2}\dr t$$
            根据条件$a^2-b^2>0$，由此我们考虑三角换元
            $$t=\frac{a}{\sqrt{a^2-b^2}}\sin\phi$$
            这时利用单调性可知$\phi$的范围是$\big[0,\arcsin\frac{\sqrt{a^2-b^2}}{a}\big]$，积分化为
            $$4b\int_0^{\arcsin(\sqrt{a^2-b^2}/a)}\frac{a^2}{\sqrt{a^2-b^2}}\cos^2\phi\dr\phi$$
            由$\cos^2\phi=\frac{1}{2}(\cos(2\phi)+1)$可知其一个原函数为$\frac{1}{2}\sin\phi\cos\phi+\frac{1}{2}\phi$，从而可写出积分结果
            $$\frac{2a^2b}{\sqrt{a^2-b^2}}\bigg(\frac{\sqrt{a^2-b^2}}{a}\cos\arcsin\frac{\sqrt{a^2-b^2}}{a}+\arcsin\frac{\sqrt{a^2-b^2}}{a}\bigg)$$
            利用三角函数的关系可直接算出
            $$\cos\arcsin\frac{\sqrt{a^2-b^2}}{a}=\sqrt{1-\frac{(\sqrt{a^2-b^2})^2}{a^2}}=\frac{b}{a}$$
            由此最终化简得到结果为
            $$2b^2+\frac{2a^2b}{\sqrt{a^2-b^2}}\arcsin\frac{\sqrt{a^2-b^2}}{a}$$
        \end{itemize}
    }
    
    \item (习题8.4.1)计算
    $$\iint_S\bigg(2x+\frac{4}{3}y+z\bigg)\dr S$$
    其中$S$为平面$\frac{x}{2}+\frac{y}{3}+\frac{z}{4}=1$在第一卦限中的部分。

    \sol{
        分为如下步骤：
        \begin{itemize}
            \item \textbf{参数化}
            
            将$z$看作$x$、$y$的函数可以得到参数化
            $$\rv=(x,y,z)=\bigg(u,v,4-2u-\frac{4}{3}v\bigg)$$
            代入第一卦限要求的$x\ge0$、$y\ge0$、$z\ge0$可得到参数范围
            $$u\ge0,\quad v\ge0,\quad u+\frac{2}{3}v\le2$$
            将上述条件确定的$(u,v)$区域记为$D$。

            \item \textbf{化为重积分}
            
            用下标$u$、$v$表示对$u$、$v$求偏导并保持另一个不变，直接计算有
            $$\rv_u=(x_u,y_u,z_u)=(1,0,-2),\quad\rv_v=(x_v,y_v,z_v)=\bigg(0,1,-\frac{4}{3}\bigg)$$
            代入公式可知
            $$\dr S=\sqrt{|\rv_u|^2|\rv_v|^2-(\rv_u\cdot\rv_v)^2}\dr u\dr v=\frac{\sqrt{61}}{3}\dr u\dr v$$
            进一步计算可发现被积函数为4，由此原积分为
            $$\frac{4\sqrt{61}}{3}\iint_D\dr u\dr v$$

            \item \textbf{积分计算}
            
            利用重积分的几何意义，$\iint_D\dr u\dr v$即为区域的面积，而画图可发现积分区域实际上是$(0,0)$、$(2,0)$与$(0,3)$连接而成的直角三角形，面积为3，从而积分结果是
            $$4\sqrt{61}$$
        \end{itemize}
    }

    \item (习题8.4.7)计算
    $$\iint_Syz\dr S$$
    其中$S$为螺旋曲面
    $$\begin{cases}x=u\cos v\\y=u\sin v\\z=v\end{cases}(0\le u\le a,\ 0\le v\le2\pi)$$
    这里参数$a>0$。

    \sol{
        分为如下步骤：
        \begin{itemize}
            \item \textbf{化为重积分}
            
            用下标$u$、$v$表示对$u$、$v$求偏导并保持另一个不变，直接计算有
            $$\rv_u=(x_u,y_u,z_u)=(\cos v,\sin v,0),\quad\rv_v=(x_v,y_v,z_v)=\bigg(-u\sin v,u\cos v,1\bigg)$$
            代入公式可知
            $$\dr S=\sqrt{|\rv_u|^2|\rv_v|^2-(\rv_u\cdot\rv_v)^2}\dr u\dr v=\sqrt{u^2+1}\dr u\dr v$$
            进一步计算可发现被积函数为$uv\sin v$，由此原积分为
            $$\iint_D\sqrt{u^2+1}uv\sin v\dr u\dr v$$

            \item \textbf{积分计算}
            
            由于区域描述已经符合累次积分的要求，可以直接写出累次积分
            $$\int_0^a\dr u\int_0^{2\pi}\sqrt{u^2+1}uv\sin v\dr u\dr v$$
            利用\exref{separation}的结论可写为
            $$\int_0^a\sqrt{u^2+1}u\dr u\int_0^{2\pi}v\sin v\dr v$$
            第一部分积分换元$t=u^2$可得其为(利用幂函数复合一次函数的原函数)
            $$\frac{1}{2}\int_0^{a^2}\sqrt{t+1}\dr t=\frac{1}{2}\cdot\frac{2}{3}\big((a^2+1)^{3/2}-1\big)=\frac{1}{3}\big((a^2+1)^{3/2}-1\big)$$
            第二部分进行分部积分可得其为
            $$-\int_0^{2\pi}v\dr\cos v=-2\pi\cos(2\pi)+\int_0^{2\pi}\cos v\dr v=-2\pi$$
            两部分乘积最终得到积分为
            $$-\frac{2}{3}\pi\big((a^2+1)^{3/2}-1\big)$$
        \end{itemize}
    }

    \item (习题8.4.12)求位于第一卦限中的球面$x^2+y^2+z^2=a^2$的质心坐标，设该球面的面密度为常数，这里参数$a>0$。
    
    \sol{
        记第一卦限中的球面为$S$，并设其面密度为$\rho$。利用质心坐标公式，设质心坐标为$(x_0,y_0,z_0)=\frac{1}{m}(x_1,y_1,z_1)$，应满足
        $$x_1=\rho\iint_Sx\dr S,\quad y_1=\rho\iint_Sy\dr S,\quad z_1=\rho\iint_Sz\dr S,\quad m=\rho\iint_S\dr S$$

        分为如下步骤：
        \begin{itemize}
            \item \textbf{对称性分析}
            
            利用第一型曲面积分的几何意义，$\iint_S\dr S$即为曲面面积，由对称性可知第一卦限球面的表面积应为全表面积的$\frac{1}{8}$，也即
            $$m=\rho\cdot\frac{1}{8}4\pi a^2=\frac{\rho}{2}\pi a^2$$
            
            由于$S$满足$(x,y,z)\in S$当且仅当$(y,x,z)\in S$，与三重积分反射对称性完全类似分析可得(详见本讲义21.3.3)
            $$\iint_Sx\dr S=\iint_Sy\dr S$$
            也即$x_1=y_1$，同理$y_1=z_1$，从而只需计算$z_1$\ (选择$z_1$计算是因为球坐标中$z$的形式更加简单)。

            \item \textbf{参数化}
            
            利用球坐标可得到球面的一个参数化
            $$\rv=(x,y,z)=(a\cos u\sin v,a\sin u\sin v,a\cos v)$$
            根据球坐标的几何意义可得到参数范围
            $$0\le u\le\frac{\pi}{2},\quad0\le v\le\frac{\pi}{2}$$
            将上述条件确定的$(u,v)$区域记为$D$。

            \note 也可通过代数方法从$x\ge0$、$y\ge0$、$z\ge0$分析范围，事实上$z\ge0$对应$v$的范围，$x\ge0$、$y\ge0$对应$u$的范围。

            \item \textbf{化为重积分}
            
            用下标$u$、$v$表示对$u$、$v$求偏导并保持另一个不变，直接计算有
            $$\rv_u=(x_u,y_u,z_u)=(-a\sin u\sin v,a\cos u\sin v,0)$$
            $$\rv_v=(x_v,y_v,z_v)=(a\cos u\cos v,a\sin u\cos v,-a\sin v)$$
            代入公式可知(由$v$范围可消除$\sin v$绝对值)
            $$\dr S=\sqrt{|\rv_u|^2|\rv_v|^2-(\rv_u\cdot\rv_v)^2}\dr u\dr v=a^2\sin v\dr u\dr v$$
            进一步计算可发现被积函数为
            $$a\cos v$$
            由此原积分为
            $$\rho\iint_Da^3\sin v\cos v\dr u\dr v$$

            \item \textbf{积分计算}
            
            由于区域描述已经符合累次积分的要求，可以直接写出累次积分
            $$\rho a^3\int_0^{\pi/2}\dr u\int_0^{\pi/2}\sin v\cos v\dr v$$
            换元$t=\sin v$即得其为
            $$\rho a^3\int_0^{\pi/2}\dr u\int_0^1t\dr t=\frac{1}{2}\rho a^3\int_0^{\pi/2}\dr u=\frac{\pi}{4}\rho a^3$$

            \item \textbf{综合结果}
            
            根据以上的计算结果可知质心坐标
            $$\bigg(\frac{a}{2},\frac{a}{2},\frac{a}{2}\bigg)$$
        \end{itemize}
    }
\end{enumerate}

\subsection{曲线积分}
\subsubsection{曲线的参数化}
曲线参数化技巧

弧长微分

正则性即参数前进时函数前进

\subsubsection{第一型曲线积分}
弧长$s$的含义

线段弯折不影响积分

重参数化不变性

反向参数化不变性

\subsubsection{第二型曲线积分}
定义

反向参数化符号相反

确定参数化的定向

\subsection{曲面积分}
\subsubsection{曲面的参数化}
曲面参数化技巧

面积微元

正则性

\subsubsection{再谈定向}
无法正则参数化的曲面

曲面定向(可用边界点)

重新参数化与Jacobi行列式(保定向/反定向)

\subsubsection{曲面积分}
面密度累计的定义

通量累计的定义

重参数化不变性

反向参数化不变/方向相反


\subsubsection{对称性}
第一型曲线/曲面与重积分完全类似使用对称性

第二型慎用对称性

\newpage
\section{多元微分与积分}
\subsection{习题解答}
\subsubsection{作业解答}
\begin{enumerate}
    \item (习题8.3.1)利用格林公式计算曲线积分：
    \begin{enumerate}
        \item[(3)]
        $$\oint_{L^+}\sqrt{x^2+y^2}\dr x+y(xy+\ln(x+\sqrt{x^2+y^2}))\dr y$$
        其中$L$是以点$A(1,1)$、$B(2,2)$、$C(1,3)$为顶点的三角形的边界线。

        \sol{
            分为如下步骤：
            \begin{itemize}
                \item \textbf{格林公式}
                
                将被积函数看作$P\dr x+Q\dr y$，由初等函数知$P$、$Q$在$L$围成的单连通区域$D$内可导且导函数连续，直接计算得
                $$Q_x-P_y=y^2+\frac{y}{x+\sqrt{x^2+y^2}}\bigg(1+\frac{x}{\sqrt{x^2+y^2}}\bigg)-\frac{y}{\sqrt{x^2+y^2}}=y^2$$

                由于三角形三边方程分别为
                $$AB\colon y=x,\quad BC\colon x+y=4,\quad AC\colon x=1$$
                利用代数或几何可确定$D$的范围
                $$y\ge x,\quad x+y\le 4,\quad x\ge1$$
                积分最终化为
                $$\iint_Dy^2\dr x\dr y$$
                
                \item \textbf{积分计算}
                
                将约束化为
                $$x\le y\le 4-x,\quad x\ge1$$
                可发现为使第一个不等式可以成立，$x$最大为2，也即$x\in[1,2]$，由此可写出累次积分
                $$\int_1^2\dr x\int_x^{4-x}y^2\dr y$$
                对$y$积分得
                $$\frac{1}{3}\int_1^2((4-x)^3-x^3)\dr x$$
                进一步计算可知结果为(计算第一部分时可换元$t=4-x$化简)
                $$\frac{1}{12}(3^4-2^4-2^4+1^4)=\frac{25}{6}$$
            \end{itemize}
        }

        \item[(4)]
        $$\oint(x+\er^x\sin y)\dr x+(x+\er^x\cos y)\dr y$$
        其中$L$是极坐标下的双扭线$r^2=\cos(2\theta)$在$x\ge0$的部分。

        \sol{
            分为如下步骤：
            \begin{itemize}
                \item \textbf{格林公式}
                
                将被积函数看作$P\dr x+Q\dr y$，由初等函数知$P$、$Q$在$L$围成的单连通区域$D$上可导且导函数连续，直接计算得
                $$Q_x-P_y=1+\er^x\cos y-\er^x\cos y=1$$

                由此积分最终化为
                $$\iint_D1\dr x\dr y$$
                这里$D$的约束即极坐标下的$r^2\le\cos(2\theta)$与$x\ge0$。
                
                \item \textbf{积分计算}
                
                换元为极坐标，为使约束$r^2\le\cos(2\theta)$合理需$\cos(2\theta)\ge0$，再由$x\ge0$可知(排除$r=0$的特殊点)\ $\cos\theta\ge0$，综合自然范围得到$\theta$的范围
                $$\theta\in\bigg[0,\frac{\pi}{4}\bigg]\cup\bigg[\frac{7\pi}{4},2\pi\bigg]$$
                而$r$的范围即
                $$0\le r\le\sqrt{\cos(2\theta)}$$
                进一步计算Jacobi行列式可知
                $$\dr x\dr y=r\dr r\dr\theta$$
                从而最终写成极坐标下的累次积分
                $$\int_0^{\pi/4}\dr\theta\int_0^{\sqrt{\cos(2\theta)}}r\dr r+\int_{7\pi/4}^{2\pi}\dr\theta\int_0^{\sqrt{\cos(2\theta)}}r\dr r$$
                对$r$积分得到其为
                $$\frac{1}{2}\bigg(\int_0^{\pi/4}\cos(2\theta)\dr\theta+\int_{7\pi/4}^{2\pi}\cos(2\theta)\dr\theta\bigg)$$
                在后一项将$\theta$换为$2\pi-\theta$可知与前一项相同，最终得到结果
                $$\int_0^{\pi/4}\cos(2\theta)\dr\theta=\frac{1}{2}\sin\frac{\pi}{2}=\frac{1}{2}$$
            \end{itemize}
        }
    \end{enumerate}

    \item (习题8.3.2(1))计算星形线
    $$\begin{cases}x=a\cos^3t\\y=a\sin^3t\end{cases}(0\le t\le 2\pi)$$
    围成的面积，这里参数$a>0$。

    \sol{
        分为如下步骤：
        \begin{itemize}
            \item \textbf{格林公式}
            
            设围成的区域为$D$，则根据重积分几何意义知面积为
            $$\iint_D\dr x\dr y$$
            利用格林公式，构造$P=0$、$Q=x$，则$Q_x-P_y=1$，设区域边界为$L$，面积应可以写为
            $$\oint_{L^+}x\dr y$$

            \item \textbf{定向与化为定积分}
            
            为了确定定向，用下标$t$代表对$t$求导，直接计算有
            $$x_t=-3a\cos^2t\sin t,\quad y_t=3a\sin^2t\cos t$$
            我们\textbf{任取}一个使得$x_t$、$y_t$\textbf{不全为0}的点，如取$t=\frac{\pi}{4}$，此处有$(x,y)=\frac{\sqrt2}{2}a(1,1)$、$(x_t,y_t)=\frac{3\sqrt2}{4}a(-1,1)$，在曲线上的点$(x,y)$处画出向量$(x_t,y_t)$，可以发现它指向逆时针方向，与要求的定向相同，从而化为定积分无需加负号。

            \note 无需画出整个曲线，只要看出这是在曲线右上角的点指向左上即知是逆时针。

            \note 事实上本题也可以不用提前确定定向，因为\textbf{面积一定为正}，若算出负数再添加负号即可。

            代入$\dr x=x_t\dr t$、$\dr y=y_t\dr t$与$x$、$y$的参数方程可得原积分为
            $$\int_0^{2\pi}3a^2\sin^2t\cos^4t\dr t$$
            
            \item \textbf{积分计算}
            
            将$\sin^2t$看作$1-\cos^2t$，原积分即化为了
            $$3a^2\int_0^{2\pi}(\cos^4t-\cos^6t)\dr t$$
            为了能使用上册教材3.4节例2的结论，不断进行对称性化简。将$t$换元为$2\pi-t$并折半区域可得积分为
            $$6a^2\int_0^\pi(\cos^4t-\cos^6t)\dr t$$
            将$t$换元为$\pi-t$并折半区域可得积分为
            $$12a^2\int_0^{\pi/2}(\cos^4t-\cos^6t)\dr t$$
            最后将$t$换元为$\frac{\pi}{2}-t$可得积分为
            $$12a^2\int_0^{\pi/2}(\sin^4t-\sin^6t)\dr t$$
            代入上册教材3.4节例2的结论得结果
            $$12a^2\bigg(\frac{1}{2}\cdot\frac{3}{4}\cdot\frac{2}{\pi}-\frac{1}{2}\cdot\frac{3}{4}\cdot\frac{5}{6}\cdot\frac{\pi}{2}\bigg)=\frac{3}{8}\pi a^2$$
        \end{itemize}
    }

    \item (习题8.3.4(1))证明如下积分与路径无关并求积分：
    $$\int_{(0,0)}^{(1,1)}(x+y)\dr x+(x-y)\dr y$$

    \sol{
        分为如下步骤：
        \begin{itemize}
            \item \textbf{格林公式}
            
            将被积函数看作$P\dr x+Q\dr y$，由初等函数知$P$、$Q$在$\rb^2$上可导且导函数连续，进一步计算验证可知
            $$Q_x-P_y=1-1=0$$
            由$\rb^2$单连通得积分与路径无关，选择路径为连接两点的线段。

            \item \textbf{参数化}
            
            根据直线方程可直接得到参数方程
            $$\begin{cases}x=t\\y=t\end{cases}(0\le t\le1)$$

            \item \textbf{化为定积分与计算}
            
            用下标$t$表示对$t$求导，直接计算有$x_t=y_t=1$，代入$\dr x=x_t\dr t$、$\dr y=y_t\dr t$与$x$、$y$的参数方程可知原积分为
            $$\int_0^12t\dr t=1$$
        \end{itemize}
    }

    \item (习题8.3.5)计算
    $$\int_L(x^4+4xy^3)\dr x+(6x^2y^2-5y^4)\dr y$$
    这里$L$是从$A(-2,-1)$到$B(3,0)$的任意路径。

    \sol{
        分为如下步骤：
        \begin{itemize}
            \item \textbf{格林公式}
            
            将被积函数看作$P\dr x+Q\dr y$，由初等函数知$P$、$Q$在$\rb^2$上可导且导函数连续，进一步计算验证可知
            $$Q_x-P_y=12xy^2-12xy^2=0$$
            由$\rb^2$单连通得积分与路径无关，选择路径为连接两点的线段。

            \item \textbf{参数化}
            
            根据直线方程可直接得到参数方程
            $$\begin{cases}x=-2+5t\\y=-1+t\end{cases}(0\le t\le1)$$

            \item \textbf{化为定积分与计算}
            
            用下标$t$表示对$t$求导，直接计算有
            $$x_t=5,\quad y_t=1$$
            代入$\dr x=x_t\dr t$、$\dr y=y_t\dr t$与$x$、$y$的参数方程可知原积分为
            $$\int_0^1\big(5(5t-2)^4+20(5t-2)(t-1)^3+6(5t-2)^2(t-1)^2-5(t-1)^4\big)\dr t$$
            换元$s=t-1$可将其化为稍简单的形式
            $$\int_{-1}^0\big(5(5s+3)^4+20(5s+3)s^3+6(5s+3)^2s^2-5s^4\big)\dr s$$
            在计算第一部分时换元$c=5s$可进一步化简，从而得到结果为四部分求和
            $$55+5+3-1=62$$
        \end{itemize}
    }

    \item (习题8.3.8)计算
    $$\int_{(0,1)}^{(1,1)}\bigg(\frac{x}{\sqrt{x^2+y^2}}+y\bigg)\dr x+\bigg(\frac{y}{\sqrt{x^2+y^2}}+x\bigg)\dr y$$

    \sol{
        分为如下步骤：
        \begin{itemize}
            \item \textbf{格林公式}
            
            将被积函数看作$P\dr x+Q\dr y$，由初等函数知$P$、$Q$除原点外可导且导函数连续，进一步计算验证可知
            $$Q_x-P_y=1-\frac{xy}{(x^2+y^2)^{3/2}}-1+\frac{xy}{(x^2+y^2)^{3/2}}=0$$
            不过，这只说明了区域不包含原点时$Q_x-P_y$在区域上的积分为0，利用格林公式，若两条路径围成的区域不包含原点，两条路径的积分应当相等。为了说明围成的区域包含原点时积分仍然相等，我们还需要说明\textbf{绕原点一周的积分为0}。

            对任何一个内部有原点的简单光滑闭曲线$L$，考虑绕原点的半径为$r$的圆$L_r$，当$r$充分小时，根据区域的定义可使得$L_r$在$L$内部。利用多连通区域的格林公式，由$Q_x-P_y$恒为0应有
            $$\int_{L^+}P\dr x+Q\dr y=\int_{L_r^+}P\dr x+Q\dr y$$
            对等号右侧，由于小圆上$\sqrt{x^2+y^2}$恒为$r$，积分化为了
            $$\int_{L_r^+}\bigg(\frac{x}{r}+y\bigg)\dr x+\bigg(\frac{y}{r}+x\bigg)\dr y$$
            将其看作$P_0\dr x+Q_0\dr y$，再次利用格林公式可计算得到$(Q_0)_x-(P_0)_y=1-1=0$，从而积分为0，因此$P\dr x+Q\dr y$在$L^+$上的积分也为0。

            综合以上讨论，原积分的确与路径无关，可取路径为连接两点的线段。

            \note 另一种看法是直接看出$P\dr x+Q\dr y=\dr(\sqrt{x^2+y^2}+xy)$，而$u=\sqrt{x^2+y^2}+xy$可以在$\rb^2$上定义，通过参数化可进一步证明任何光滑曲线上的积分都是$u$在两点处取值的差，详见本讲义22.2.2。

            \item \textbf{参数化}
            
            根据直线方程可直接得到参数方程
            $$\begin{cases}x=t\\y=1\end{cases}(0\le t\le1)$$

            \item \textbf{化为定积分与计算}
            
            用下标$t$表示对$t$求导，直接计算有
            $$x_t=1,\quad y_t=0$$
            代入$\dr x=x_t\dr t$、$\dr y=y_t\dr t$与$x$、$y$的参数方程可知原积分为
            $$\int_0^1\bigg(\frac{t}{\sqrt{t^2+1}}+1\bigg)\dr t$$
            第二部分积分即为1，第一部分中换元$s=t^2$可化为
            $$\frac{1}{2}\int_0^1\frac{1}{\sqrt{s+1}}\dr s$$
            这已经是幂函数复合一次函数的形式，可写出原函数以得到积分$\sqrt{2}-1$，综合可得最终结果为
            $$(\sqrt{2}-1)+1=\sqrt2$$
        \end{itemize}
    }

    \item (习题8.3.9)计算
    $$\int_L(x^2+y)\dr x+(x-y^2)\dr y$$
    这里$L$是从$A(0,0)$到$B(1,1)$沿曲线$y^3=x^2$的曲线段。

    \sol{
        分为如下步骤：
        \begin{itemize}
            \item \textbf{格林公式}
            
            将被积函数看作$P\dr x+Q\dr y$，由初等函数知$P$、$Q$在$\rb^2$上可导且导函数连续，进一步计算验证可知
            $$Q_x-P_y=1-1=0$$
            由$\rb^2$单连通得积分与路径无关，选择路径为连接两点的线段。

            \item \textbf{参数化}
            
            根据直线方程可直接得到参数方程
            $$\begin{cases}x=t\\y=t\end{cases}(0\le t\le1)$$

            \item \textbf{化为定积分与计算}
            
            用下标$t$表示对$t$求导，直接计算有$x_t=y_t=1$，代入$\dr x=x_t\dr t$、$\dr y=y_t\dr t$与$x$、$y$的参数方程可知原积分为
            $$\int_0^12t\dr t=1$$
        \end{itemize}
    }

    \item (习题8.3.12)设函数$u(x,y)$、$v(x,y)$在有界闭区域$D$上有连续的二阶偏导数，$D$有分段光滑边界$L$，证明(以下省略函数的参数$(x,y)$)：
    \begin{enumerate}[(1)]
        \item $$\iint_Dv\triangle u\dr\sigma=\oint_Lv\frac{\partial u}{\partial\nv}\dr s-\iint_D(\partial_1u\partial_1v+\partial_2u\partial_2v)\dr\sigma$$
        其中$\frac{\partial u}{\partial\nv}$表示$u$沿边界$L$外法向的方向导数。

        \sol{
            我们先分析再证明：
            \begin{itemize}
                \item \textbf{问题分析}
                
                为了将曲线上的积分和区域内的积分产生关联，我们必须使用\textbf{格林公式}，而原本右侧为第一型曲线积分，必须化为\textbf{第二型曲线积分}才能使用格林公式。

                设$\vec\tau=(\tau_1,\tau_2)$代表曲线逆时针方向的单位切向量，利用第一型曲线积分与第二型曲线积分的关系可知(这里曲线以逆时针为正定向)
                $$(\tau_1\dr s,\tau_2\dr s)=(\dr x,\dr y)$$

                \note 也可写为向量形式$\vec\tau\dr s=\dr\rv$。

                不过，题中为\textbf{法向量}的相关式子，而非切向量。利用几何关系可知，$L$的逆时针方向单位切向量为单位外法向量向逆时针方向旋转90度的结果。设单位外法向量$\nv=(n_1,n_2)$，计算可发现
                $$n_1=\tau_2,\quad n_2=-\tau_1$$
                由此有
                $$n_1\dr s=\dr y,\quad n_2\dr s=-\dr x$$
                我们接下来的目标即是利用这些式子将原问题转化为第二型曲线积分后再计算。

                \item \textbf{化为第二型曲线积分}
                
                由$u$可微，方向导数可以写为(可见上学期讲义14.4.2)
                $$n_1\partial_1u+n_2\partial_2u$$
                
                \note 也可写为向量形式$\nabla u\cdot\vec{n}$。

                由此可以将
                $$\oint_Lv\frac{\partial u}{\partial\nv}\dr s$$
                展开写为
                $$\oint_Lv\partial_1un_1\dr s+v\partial_2u\dr y$$
                从而其等于
                $$\oint_{L^+}v\partial_1u\dr y-v\partial_2u\dr x$$

                \item \textbf{格林公式}
                
                由$u$、$v$有连续的二阶偏导数，符合使用条件，从而上式可利用格林公式化为
                $$\iint_D(\partial_1(v\partial_1u)+\partial_2(v\partial_2u))\dr\sigma$$
                利用乘积求导公式得被积函数为$\partial_1v\partial_1u+\partial_2v\partial_2u+v\partial_1^2u+v\partial_2^2u$，再根据$\triangle u=\partial_1^2u+\partial_2^2u$的定义可最终计算出
                $$\oint_Lv\frac{\partial u}{\partial\nv}\dr s=\iint_D(\partial_1u\partial_1v+\partial_2u\partial_2v+v\triangle u)\dr\sigma$$
                移项得最终结论。

                \note 也可将$\partial_1u\partial_1v+\partial_2u\partial_2v$写为向量形式$\nabla u\cdot\nabla v$。
            \end{itemize}
        }

        \item $$\iint_D(u\triangle v-v\triangle u)\dr\sigma=\int_L\bigg(u\frac{\partial v}{\partial\nv}-v\frac{\partial u}{\partial\nv}\bigg)\dr s$$
        
        \sol{
            利用前一问结论有
            $$\iint_Du\triangle v\dr\sigma=\oint_Lu\frac{\partial v}{\partial\nv}\dr s-\iint_D(\partial_1v\partial_1u+\partial_2v\partial_2u)\dr\sigma$$
            $$\iint_Dv\triangle u\dr\sigma=\oint_Lv\frac{\partial u}{\partial\nv}\dr s-\iint_D(\partial_1u\partial_1v+\partial_2u\partial_2v)\dr\sigma$$
            两式作差，利用乘法交换律可消去右侧第二项，即得最终结论。
        }
    \end{enumerate}
    
    \item (习题8.3.13)若$u(x,y)$在有界闭区域$D$上有连续的二阶偏导数，$D$有分段光滑边界$L$，且
    $$\forall(x,y)\in D,\quad\triangle u(x,y)=0$$
    证明(以下省略函数的参数$(x,y)$)：
    \begin{enumerate}[(1)]
        \item $$\oint_Lu\frac{\partial u}{\partial\nv}\dr s=\iint_D\big((\partial_1u)^2+(\partial_2u)^2\big)\dr\sigma$$
        其中$\frac{\partial u}{\partial\nv}$表示$u$沿边界$L$外法向的方向导数。

        \sol{
            利用前一题结论可知
            $$\iint_Du\triangle u\dr\sigma=\oint_Lu\frac{\partial u}{\partial\nv}\dr s-\iint_D((\partial_1u)^2+(\partial_2u)^2)\dr\sigma$$
            由于$\triangle u$恒为0，第一项积分为0，从而有
            $$\oint_Lu\frac{\partial u}{\partial\nv}\dr s=\iint_D((\partial_1u)^2+(\partial_2u)^2)\dr\sigma$$
        }

        \item 若$u(x,y)$在$L$上恒为0，则$D$上也恒为0。
        
        \sol{
            由于$L$上$u$为0，前一问中的
            $$\oint_Lu\frac{\partial u}{\partial\nv}\dr s=\oint_L0\dr s=0$$
            由此可知
            $$\iint_D\big((\partial_1u)^2+(\partial_2u)^2\big)\dr\sigma=0$$
            由于$u$的偏导连续，且$(\partial_1u)^2+(\partial_2u)^2\ge0$恒成立，利用本讲义19.1.2第1题的结论(这里仅针对区域$D$内部，记为$D^\circ$)有
            $$\forall (x,y)\in D^\circ,\quad(\partial_1u(x,y))^2+(\partial_2u(x,y))^2\ge0$$
            也即$\partial_1u(x,y)=\partial_2u(x,y)=0$。

            通过拉格朗日中值定理可证明$u$在$D^\circ$中为常数(详细证明可见上学期讲义15.1第10题)，再由连续性可知$u$在$D$上为常数，根据边界为0可得恒为0，得证。

            \note 最后这段说明的严谨论证相对复杂，不过直观来看是自然的，一般写出核心说明步骤即可。
        }
    \end{enumerate}

    \item (总练习题8.3)计算
    $$\int_{(x_1,y_1)}^{(x_2,y_2)}f(x)\dr x+g(y)\dr y$$
    其中$f(x)$、$g(y)$为$\rb$上的连续函数。

    \sol{
        \note 注意本题并未给出可微条件，不符合教材中格林公式的使用要求，无法直接使用格林公式。

        设$(x_1,y_1)$到$(x_2,y_2)$的光滑曲线$L$参数化为
        $$\begin{cases}x=\phi(t)\\y=\psi(t)\end{cases}(a\le t\le b)$$
        代入$\dr x=x_t\dr t$、$\dr y=y_t\dr t$可发现积分化为
        $$\int_a^b(f(\phi(t))\phi'(t)+g(\psi(t))\psi'(t))\dr t$$
        利用定积分换元公式，计算第一部分积分时换元$x=\phi(t)$、计算第二部分积分时换元$y=\psi(t)$，可得积分化为
        $$\int_{x_1}^{x_2}f(x)\dr x+\int_{y_1}^{y_2}g(y)\dr y$$
        这就是最终结果。

        \note 注意定积分换元公式无需$\phi(t)$或$\psi(t)$是单调函数。

        \note 仍可以计算出$Q_x-P_y=0$，从而可以考虑配凑为全微分$f(x)\dr x+g(y)\dr y=\dr(F(x)+G(y))$，以此也能得到正确的结果。
    }
    
    \item (总练习题8.4)计算
    $$\int_{(1,2)}^{(5,6)}\frac{x\dr y-y\dr x}{(x-y)^2}$$
    其中积分路径不与直线$x=y$相交。
    
    \sol{
        分为如下步骤：
        \begin{itemize}
            \item \textbf{格林公式}
            
            由于$(1,2)$、$(5,6)$都满足$y>x$，将满足$y>x$的点记为区域$D$，可发现其单连通。

            将被积函数看作$P\dr x+Q\dr y$，由初等函数知$P$、$Q$在$D$中可导且导函数连续，进一步计算验证可知
            $$Q_x-P_y=\frac{1}{(x-y)^2}+\frac{2x}{(x-y)^3}+\frac{1}{(x-y)^2}-\frac{2y}{(x-y)^3}=0$$
            由$D$单连通得积分与$D$中路径无关，选择路径为连接两点的线段。
            
            \item \textbf{参数化}
            
            根据直线方程可直接得到参数方程
            $$\begin{cases}x=1+t\\y=2+t\end{cases}(0\le t\le 4)$$

            \item \textbf{化为定积分与计算}
            
            用下标$t$表示对$t$求导，直接计算有$x_t=y_t=1$，代入$\dr x=x_t\dr t$、$\dr y=y_t\dr t$与$x$、$y$的参数方程可知原积分为
            $$\int_0^4(-1)\dr t=-4$$
        \end{itemize}
    }
    
    \item (总练习题8.6)设$f(t)$在$\rb$上连续，$L$为给定定向的分段光滑的闭曲线，证明
    $$\oint_Lf(xy)(y\dr x+x\dr y)=0$$

    \sol{
        \note 与本部分第9题相同，本题也不能直接使用格林公式。不过，若$f$可导仍可以计算出$Q_x-P_y=0$，从而可以考虑配凑为全微分。

        由于$y\dr x+x\dr y=\dr(xy)$，设$F(u)$是$f(u)$的一个原函数，直接计算可发现
        $$\dr F(xy)=f(xy)(y\dr x+x\dr y)$$
        设$L$参数化为(这里$\phi(a)=\phi(b)$、$\psi(a)=\psi(b)$为曲线上某一点，且定向与要求一致)
        $$\begin{cases}x=\phi(t)\\y=\psi(t)\end{cases}(a\le t\le b)$$
        代入$\dr x=x_t\dr t$、$\dr y=y_t\dr t$可发现积分化为(除了有限个不光滑点外)
        $$\int_a^bf(\phi(t)\psi(t))(\phi'(t)\psi(t)+\psi'(t)\phi(t))\dr t$$
        利用配凑出的全微分，我们可以发现被积函数实际上就是
        $$\frac{\dr}{\dr t}F(\phi(t)\psi(t))$$
        从而根据$\phi$、$\psi$的分段光滑性利用微积分基本定理可得结果为
        $$F(\phi(b)\psi(b))-F(\phi(a)\psi(a))=0$$
        从而得证。

        \note 这里对分段光滑情况讨论的严谨性可见本讲义22.3.3。
    }

    \item (总练习题8.9)计算
    $$\oint_{L^+}\sqrt{x^2+y^2}\dr x+y\big(xy+\ln(x+\sqrt{x^2+y^2})\big)\dr y$$
    其中$L$为圆周$x^2+y^2=1$。

    \note 本题实际上是一个$(-1,0)$点为无穷的\textbf{广义积分}，我们只形式计算结果，不考虑敛散性问题。
    
    \sol{
        我们这里采用先化简后用格林公式的方法进行更简便的计算。分为如下步骤：
        \begin{itemize}
            \item \textbf{积分化简}
            
            由于$L$上满足$x^2+y^2=1$，我们可以将$x^2+y^2$改写为1以化简积分成
            $$\oint_{L^+}\dr x+y(xy+\ln(x+1))\dr y$$
            
            \item \textbf{格林公式}
            
            将被积函数看作$P\dr x+Q\dr y$，由初等函数知$P$、$Q$在$L$围成的单连通区域$D$\ (不含边界)内可导且导函数连续，进一步计算验证可知
            $$Q_x-P_y=y^2+\frac{y}{x+1}-0=y^2+\frac{y}{x+1}$$
            积分最终化为
            $$\iint_D\bigg(y^2+\frac{y}{\ln(1+x)}\bigg)\dr x\dr y$$

            \item \textbf{积分计算}
            
            由于$D$满足$(x,y)\in D$当且仅当$(x,-y)\in D$，利用函数折半可发现
            $$\iint_D\frac{y}{\ln(1+x)}\dr x\dr y=0$$
            从而积分可写为
            $$\iint_Dy^2\dr x\dr y$$
            换元为极坐标，约束变为$r^2\le1$，结合自然范围得到$r$、$\theta$的范围
            $$0\le r\le1,\quad 0\le\theta\le2\pi$$
            进一步计算Jacobi行列式可知
            $$\dr x\dr y=r\dr r\dr\theta$$
            从而最终写成极坐标下的累次积分
            $$\int_0^1\dr r\int_0^{2\pi}r^3\sin^2\theta\dr\theta$$
            类似本讲义19.1.1第10题对$\theta$积分可得到其为
            $$\pi\int_0^1r^3\dr r=\frac{\pi}{4}$$
        \end{itemize}
    }

    \item (总练习题8.12)计算积分
    $$F(a)=\iint_Sf(x,y,z)\dr S$$
    其中$S$为球面$x^2+y^2+z^2=a^2$，这里参数$a>0$，且
    $$f(x,y,z)=\begin{cases}x^2+y^2&z\ge\sqrt{x^2+y^2}\\0&z<\sqrt{x^2+y^2}\end{cases}$$

    \sol{
        分为如下步骤：
        \begin{itemize}
            \item \textbf{参数化}
            
            利用球坐标可得到球面的一个参数化
            $$\rv=(x,y,z)=(a\cos u\sin v,a\sin u\sin v,a\cos v)$$
            球坐标自然范围即为参数范围
            $$0\le u\le2\pi,\quad0\le v\le\pi$$
            将上述条件确定的$(u,v)$区域记为$D$。
            
            \item \textbf{化为重积分}
            
            用下标$u$、$v$表示对$u$、$v$求偏导并保持另一个不变，直接计算有
            $$\rv_u=(x_u,y_u,z_u)=(-a\sin u\sin v,a\cos u\sin v,0)$$
            $$\rv_v=(x_v,y_v,z_v)=(a\cos u\cos v,a\sin u\cos v,-a\sin v)$$
            代入公式可知(由$v$范围可消除$\sin v$绝对值)
            $$\dr S=\sqrt{|\rv_u|^2|\rv_v|^2-(\rv_u\cdot\rv_v)^2}\dr u\dr v=a^2\sin v\dr u\dr v$$
            进一步计算可发现被积函数为
            $$g(u,v)=\begin{cases}a^2\sin^2v&\cos v\ge\sin v\\0&\cos v<\sin v\end{cases}$$
            由此原积分为
            $$\iint_Da^2g(u,v)\sin v\dr u\dr v$$

            \item \textbf{积分计算}
            
            由于区域描述已经符合累次积分的要求，可以直接写出累次积分
            $$\int_0^{2\pi}\dr u\int_0^\pi g(u,v)\dr v$$

            根据$v$的范围$[0,\pi]$可以发现，其中满足$\cos v\ge\sin v$的范围实际上只有$[0,\frac{\pi}{4}]$，由此累次积分改写为
            $$\int_0^{2\pi}\dr u\int_0^{\pi/4}a^4\sin^3v\dr v$$
            换元$t=\cos v$，则$\sin^2v=1-t^2$，由此可将积分化为
            $$\int_0^{2\pi}\dr u\int_{\sqrt2/2}^1a^4(1-t^2)\dr t=\int_0^{2\pi}a^4\bigg(1-\frac{\sqrt2}{2}-\frac{1}{3}+\frac{1}{3}\cdot\frac{\sqrt2}{4}\bigg)\dr u=\bigg(\frac{4}{3}-\frac{5}{6}\sqrt2\bigg)\pi a^4$$
        \end{itemize}
    }
\end{enumerate}

\subsubsection{补充题解答}
\begin{enumerate}
    \item (习题8.2.1)计算
    $$\int_L2xy\dr x-x^2\dr y$$
    其中$L$为$O(0,0)$到$A(2,1)$的如下曲线段：
    \begin{enumerate}[(1)]
        \item 直线段。
        
        \sol{
            分为如下步骤：
            \begin{itemize}
                \item \textbf{参数化}
                
                根据直线方程可直接得到参数方程
                $$\begin{cases}x=2t\\y=t\end{cases}(0\le t\le1)$$

                \item \textbf{化为定积分与计算}
                
                用下标$t$表示对$t$求导，直接计算有
                $$x_t=2,\quad y_t=1$$
                代入$\dr x=x_t\dr t$、$\dr y=y_t\dr t$与$x$、$y$的参数方程可得原积分为
                $$\int_0^14t^2\dr t=\frac{4}{3}$$
            \end{itemize}
        }

        \item 以$y$轴为对称轴的抛物线段。
        
        \sol{
            分为如下步骤：
            \begin{itemize}
                \item \textbf{参数化}
                
                由于以$y$轴为对称轴的抛物线即$y=ax^2+b$，待定系数可解得$y=\frac{1}{4}x^2$，根据抛物线方程可直接得到参数方程
                $$\begin{cases}x=t\\y=\frac{1}{4}t^2\end{cases}(0\le t\le1)$$

                \item \textbf{化为定积分与计算}
                
                用下标$t$表示对$t$求导，直接计算有
                $$x_t=1,\quad y_t=\frac{1}{2}t$$
                代入$\dr x=x_t\dr t$、$\dr y=y_t\dr t$与$x$、$y$的参数方程可得原积分为
                $$\int_0^1\frac{t^3-t^3}{2}\dr t=0$$
            \end{itemize}
        }

        \item 折线段$OBA$，这里$B$为$(2,0)$。
        
        \sol{
            分为如下步骤：
            \begin{itemize}
                \item \textbf{参数化}
                
                根据直线方程可直接得到两部分的参数方程
                $$\begin{cases}x=2t\\y=0\end{cases}(0\le t\le1),\quad\begin{cases}x=2\\y=t\end{cases}(0\le t\le1)$$

                \item \textbf{化为定积分与计算}
                
                用下标$t$表示对$t$求导，直接计算有第一段中$x_t=2$、$y_t=0$，第二段中$x_t=0$、$y_t=1$，代入$\dr x=x_t\dr t$、$\dr y=y_t\dr t$与$x$、$y$的参数方程可得原积分为
                $$\int_0^10\dr t+\int_0^1(-2^2)\dr t=-4$$
            \end{itemize}
        }
        
        \item 折线段$OCA$，这里$C$为$(0,1)$。
        
        \sol{
            分为如下步骤：
            \begin{itemize}
                \item \textbf{参数化}
                
                根据直线方程可直接得到两部分的参数方程
                $$\begin{cases}x=0\\y=t\end{cases}(0\le t\le1),\quad\begin{cases}x=2t\\y=1\end{cases}(0\le t\le1)$$

                \item \textbf{化为定积分与计算}
                
                用下标$t$表示对$t$求导，直接计算有第一段中$x_t=0$、$y_t=1$，第二段中$x_t=2$、$y_t=0$，代入$\dr x=x_t\dr t$、$\dr y=y_t\dr t$与$x$、$y$的参数方程可得原积分为
                $$\int_0^10\dr t+\int_0^12\cdot(2t)\cdot2\dr t=4$$
            \end{itemize}
        }
    \end{enumerate}

    \item (习题8.2.3)计算
    $$\oint_{L^+}x^2\dr y-y^2\dr x$$
    其中$L$为椭圆周$\frac{x^2}{a^2}+\frac{y^2}{b^2}=1$，这里$a>0$、$b>0$为参数。

    \sol{
        分为如下步骤：
        \begin{itemize}
            \item \textbf{参数化}
                
            由其为椭圆，利用\textbf{广义极坐标}可参数化为
            $$\begin{cases}x=a\cos\theta\\y=b\sin\theta\end{cases}(0\le\theta\le2\pi)$$

            \item \textbf{化为定积分}
                
            用下标$\theta$表示对$\theta$求导，直接计算有
            $$x_\theta=-a\sin\theta,\quad y_\theta=b\cos\theta$$
            代入$\dr x=x_\theta\dr\theta$、$\dr y=y_\theta\dr\theta$与$x$、$y$的参数方程可得原积分为
            $$\int_0^{2\pi}(a^2b\cos^3\theta+ab^2\sin^3\theta)\dr\theta$$

            \item \textbf{积分计算}
            
            利用对称性，在$[0,\pi]$将$\theta$换元为$\pi+\theta$可得
            $$\int_0^\pi(a^2b\cos^3\theta+ab^2\sin^3\theta)\dr\theta=-\int_\pi^{2\pi}(a^2b\cos^3\theta+ab^2\sin^3\theta)\dr\theta$$
            移项即得到所求的积分为0。

            \note 本质上是由于$\sin^3\theta$、$\cos^3\theta$周期内积分，正负部分可相互抵消。
        \end{itemize}
    }

    \item (习题8.2.7)计算
    $$\oint_{L^+}\frac{\dr x+\dr y}{|x|+|y|}$$
    其中$L$为以$A(2,0)$、$B(0,2)$、$C(-2,0)$、$D(0,-2)$为顶点的正方形回路。

    \sol{
        分为如下步骤：
        \begin{itemize}
            \item \textbf{参数化}
                
            直接画图可发现逆时针定向即为$ABCDA$的顺序，根据直线方程可直接得到四条边$AB$、$BC$、$CD$、$DA$的参数方程
            $$\begin{cases}x=2-2t\\y=2t\end{cases}(0\le t\le1),\quad\begin{cases}x=-2t\\y=2-2t\end{cases}(0\le t\le1)$$
            $$\begin{cases}x=-2+2t\\y=-2t\end{cases}(0\le t\le1),\quad\begin{cases}x=2t\\y=-2+2t\end{cases}(0\le t\le1)$$

            \item \textbf{化为定积分与计算}
                
            用下标$t$表示对$t$求导，直接计算可知四段分别有
            $$x_t=-2,\quad y_t=2$$
            $$x_t=-2,\quad y_t=-2$$
            $$x_t=2,\quad y_t=-2$$
            $$x_t=2,\quad y_t=2$$
            且根据正负性可发现四段中$\frac{1}{|x|+|y|}$分别为
            $$\frac{1}{x+y},\quad\frac{1}{y-x},\quad\frac{1}{-x-y},\quad\frac{1}{x-y}$$
            
            代入$\dr x=x_t\dr t$、$\dr y=y_t\dr t$与$x$、$y$的参数方程可得原积分为
            $$\int_0^20\dr t+\int_0^2\frac{-4\dr t}{2}+\int_0^20\dr t+\int_0^2\frac{4\dr t}{2}=0$$
        \end{itemize}
    }

    \item (习题8.2.11)计算
    $$\int_L(y-z,z-x,x-y)\cdot\dr\rv$$
    其中$L$为圆周
    $$\begin{cases}x^2+y^2+z^2=a^2\\y=x\tan\beta\end{cases}$$
    且从$x$轴正向看去，$L$沿逆时针方向。这里参数$a>0$、$0<\beta<\frac{\pi}{2}$。

    \sol{
        分为如下步骤：
        \begin{itemize}
            \item \textbf{参数化}
                
            将$y=x\tan\beta$代入第一个方程得到
            $$\frac{x^2}{\cos^2\beta}+z^2=a^2$$
            这已经符合椭圆方程的形式，从而可利用\textbf{广义极坐标}参数化为
            $$\begin{cases}x=a\cos\beta\cos\theta\\z=a\sin\theta\end{cases}(0\le\theta\le2\pi)$$
            进一步代入即得此时$y=a\sin\beta\cos\theta$。

            \note 注意这里$\beta$为常数，并非真正的球坐标。

            \item \textbf{定向}

            为了确定定向，用下标$\theta$代表对$\theta$求导，直接计算有
            $$x_\theta=-a\cos\beta\sin\theta,\quad y_\theta=-a\sin\beta\sin\theta,\quad z_\theta=a\cos\theta$$
            我们\textbf{任取}一个使得$x_\theta$、$y_\theta$、$z_\theta$\textbf{不全为0}的点，如取$\theta=\frac{\pi}{2}$，此处有$(x,y,z)=(0,0,a)$、$(x_\theta,y_\theta,z_\theta)=(-a\cos\beta,-a\sin\beta,0)$，在曲线上的点$(x,y,z)$处画出向量$(x_\theta,y_\theta,z_\theta)$，可以发现它指向逆时针方向，与要求的定向相同，从而化为定积分无需添加负号。

            \note 事实上还有更简单的无需画出原曲线的判断方法：根据极坐标参数化的几何性质或在$Oxz$平面画图可以发现，从$y$轴\textbf{负向}看去，参数化后曲线在$Oxz$平面上的投影是逆时针方向的。而曲线在$x$增大时$y$值逐渐增大，几何上可以进一步观察到$x$轴正向看去相当于$y$轴负向看去，因此仍为逆时针。

            \item \textbf{化为定积分与计算}

            代入$\dr x=x_\theta\dr\theta$、$\dr y=y_\theta\dr\theta$、$\dr z=z_\theta\dr\theta$与$x$、$y$、$z$的参数方程可得原积分为(仍然注意$\beta$为常数)
            $$a^2\int_0^{2\pi}(\cos\beta-\sin\beta)\dr\theta=2\pi a^2(\cos\beta-\sin\beta)$$
        \end{itemize}
    }

    \item (总练习题8.7)设$P$、$Q$为二元连续函数，证明
    $$\bigg|\int_lP(x,y)\dr x+Q(x,y)\dr y\bigg|\le LM$$
    其中$L$为光滑曲线段$l$的弧长，$M$为$l$上$\sqrt{P^2(x,y)+Q^2(x,y)}$的最大值。

    \note 在$l$为光滑曲线段时，可以证明其为有界闭集，从而根据有界闭集上连续函数的最值定理(可见上学期讲义14.3.2)，$M$必然存在。

    \sol{
        由于对1的第一重曲线积分的\textbf{几何意义}即为弧长，有
        $$L=\int_l\dr s$$
        为了与此形式靠近，我们需要设法将其转化为第一型曲线积分。

        设$\vec\tau=(\tau_1,\tau_2)$代表曲线段对应定向的单位切向量，利用第一型曲线积分与第二型曲线积分的关系可知
        $$(\dr x,\dr y)=(\tau_1\dr s,\tau_2\dr s)$$
        也即(右侧省略$\tau$、$P$、$Q$的参数$(x,y)$)
        $$\int_lP(x,y)\dr x+Q(x,y)\dr y=\int_l(P\tau_1+Q\tau_2)\dr s$$
        由于$\tau$是单位向量，由\textbf{柯西不等式}可得
        $$|P\tau_1+Q\tau_2|\le\sqrt{P^2+Q^2}\sqrt{\tau_1^2+\tau_2^2}=\sqrt{P^2+Q^2}\le M$$
        由此利用第一型曲线积分的性质可知
        $$-ML=\int_l(-M)\dr s\le\int_lP(x,y)\dr x+Q(x,y)\dr y\le\int_lM\dr s=ML$$
        这就是要证的结论。
    }

    \item (习题8.3.1)利用格林公式计算曲线积分：
    \begin{enumerate}
        \item[(2)]
        $$\oint_{L^+}y^2\dr x+x^2\dr y$$
        其中$L$是以点$O(0,0)$、$B(1,0)$、$C(0,1)$为顶点的三角形的边界线。

        \sol{
            分为如下步骤：
            \begin{itemize}
                \item \textbf{格林公式}
                
                将被积函数看作$P\dr x+Q\dr y$，由初等函数知$P$、$Q$在$L$围成的单连通区域$D$内可导且导函数连续，直接计算得
                $$Q_x-P_y=2x-2y$$

                由于三角形三边方程分别为
                $$x=0,\quad y=0,\quad x+y=1$$
                利用代数或几何可确定$D$的范围
                $$x\ge0,\quad y\ge0,\quad x+y\le1$$
                积分最终化为
                $$\iint_D(2x-2y)\dr x\dr y$$
                
                \item \textbf{积分计算}
                
                由于区域$D$满足$(x,y)\in D$当且仅当$(y,x)\in D$，换元可发现
                $$\iint_Dx\dr x\dr y=\iint_Dy\dr x\dr y$$
                从而所求的积分为0。
            \end{itemize}
        }

        \item[(5)]
        $$\int_{L^+}(\er^x\sin y+\sin x-8y)\dr x+(\er^x\cos y-\sin y)\dr y$$
        其中$L$为上半圆$0\le y\le\sqrt{ax-x^2}$的边界，这里参数$a>0$。

        \sol{
            分为如下步骤：
            \begin{itemize}
                \item \textbf{格林公式}
                
                将被积函数看作$P\dr x+Q\dr y$，由初等函数知$P$、$Q$在题中所述的半圆$D$内可导且导函数连续，直接计算得
                $$Q_x-P_y=\er^x\cos y-\er^x\cos y+8=8$$
                由此可利用格林公式将积分化为
                $$\iint_D(Q_x-P_y)\dr x\dr y=8\iint_D\dr x\dr y$$
                
                \item \textbf{积分计算}
                
                利用重积分的几何意义，$\iint_D\dr x\dr y$即为$D$的面积，而$D$是半径为$\frac{a}{2}$的半圆，由此得所求的积分为
                $$8\cdot\frac{1}{2}\cdot\pi\bigg(\frac{a}{2}\bigg)^2=\pi a^2$$
            \end{itemize}
        }
    \end{enumerate}

    \item (习题8.3.3)设$f(t)$有连续的一阶导数，$L$为给定定向的分段光滑的闭曲线，证明
    $$\oint_Lf(xy)(y\dr x+x\dr y)=0$$

    \sol{
        将被积函数看作$P\dr x+Q\dr y$，由$f$有连续一阶导数知$P$、$Q$在$L$围成的区域$D$内可导且导函数连续，直接计算得
        $$Q_x-P_y=f(xy)+xyf'(xy)-(f(xy)+xyf'(xy))=0$$
        由此可利用格林公式知(这里正对应逆时针，负对应顺时针)
        $$\oint_Lf(xy)(y\dr x+x\dr y)=\pm\iint_D(Q_x-P_y)\dr x\dr y=0$$

        \note 将本题与本讲义22.1.1第11题进行对比可以发现格林公式的\textbf{连续导数}条件有时可以减弱。
    }

    \item (习题8.3.10)设$D$是有分段光滑边界$L$的有界闭区域，函数$P(x,y)$、$Q(x,y)$在$D$上有连续的一阶偏导数，证明(以下省略函数的参数$(x,y)$)
    $$\oint_L\big(P\cos(\nv,x)+Q\cos(\nv,y)\big)\dr s=\iint_D(\partial_1P+\partial_2Q)\dr\sigma$$
    其中$\nv$为$L$对应位置的单位外法向量，$(\nv,x)$、$(\nv,y)$为其与两个坐标轴正向的夹角。

    \sol{
        设$\nv=(n_1,n_2)$，利用内积性质可知
        $$\cos(\nv,x)=\frac{\nv\cdot(1,0)}{|nv||(1,0)|}=n_1$$
        同理$\cos(\nv,y)=n_2$，从而左侧的被积函数即为$(P,Q)\cdot\vec{n}$。接下来的的证明过程见\exref{2d gauss}。

        \note 注意这本质上是二维形式的\textbf{高斯定理}，之后将学到三维高斯定理的形式。
    }

    \item (习题8.3.11)考虑分段光滑的简单闭曲线$L$，计算
    $$\oint_L(x\cos(\nv,x)+y\cos(\nv,y))\dr s$$
    其中$\nv$为$L$对应位置的单位外法向量，$(\nv,x)$、$(\nv,y)$为其与两个坐标轴正向的夹角。

    \sol{
        利用上题结论，将被积函数看作$P\cos(\nv,x)+Q\cos(\nv,y)$，设$L$围成区域为$D$即得积分化为
        $$\iint_D(\partial_1P+\partial_2Q)\dr\sigma=\iint_D\bigg(\frac{\partial x}{\partial x}+\frac{\partial y}{\partial y}\bigg)\dr\sigma=2\iint_D\dr\sigma$$
        也即结果为$D$面积的两倍。
    }
    
    \item (总练习题8.10)设$u(x,y)$在$\rb^2$上有连续的二阶偏导数，证明以下两命题等价(以下省略函数的参数$(x,y)$)：
    \begin{itemize}
        \item $\triangle u=0$在$\rb^2$上处处成立；
        \item 对任何分段光滑的简单闭曲线$L$，有
        $$\oint_{L^+}\frac{\partial u}{\partial\nv}\dr s=0$$
        其中$\nv$为$L$对应位置的单位外法向量。
    \end{itemize}

    \sol{
        我们先化简积分后再进行证明：
        \begin{itemize}
            \item \textbf{化简积分}
            
            利用本讲义22.1.1第7题(1)的结论，取$v(x,y)=1$，则$v$的各偏导均为0，从而计算得到
            $$\oint_L\frac{\partial u}{\partial\vec{n}}\dr s=\iint_D\triangle u\dr\sigma$$

            \note 注意实际上\textbf{无需记忆}本讲义22.1.1第7题(1)的结论，只要掌握了分析方法，本部分结论是可以通过完全相同的方式计算出的。

            由此，第二个命题即等价于对任何边界分段光滑的有界区域$D$有
            $$\iint_D\triangle u\dr\sigma=0$$
            我们接下来记$z(x,y)=\triangle u(x,y)$，由$u$各二阶偏导连续可知$z$连续，此时第一个命题即$z$恒为0，第二个命题即$z$在任何边界分段光滑的有界区域积分为0。

            \item \textbf{上推下}
            
            若$z$恒为0，利用重积分定义即得
            $$\iint_Dz(x,y)\dr\sigma=0$$
            从而上推下得证。

            \item \textbf{下推上}
            
            利用\textbf{反证}法。若$z$在某点$(x_0,y_0)$非零，设
            $$c=z(x_0,y_0)$$
            考虑$c>0$的情况，利用\textbf{连续性}可知存在$\delta>0$使得对任何满足$\sqrt{(x-x_0)^2+(y-y_0)^2}<\delta$的$(x,y)$有
            $$|z(x,y)-z(x_0,y_0)|<\frac{c}{2}$$
            利用三角不等式得
            $$z(x,y)>\frac{c}{2}$$
            由此，设$D$为以$(x_0,y_0)$为圆心、$\frac{\delta}{2}$为半径的圆，利用重积分性质与几何意义(对1的重积分为面积)有
            $$\iint_Dz(x,y)\dr\sigma\ge\iint_D\frac{c}{2}\dr\sigma=\frac{c}{2}\iint_D\dr\sigma=\frac{c\delta^2\pi}{8}>0$$
            由于$D$是圆，边界是分段光滑的，从而$z$在$D$上积分应为0，这就得到了矛盾。若$c<0$可类似证明。

            \note 这类利用\textbf{连续性}估算的操作是高等数学中的一个基本技巧，本讲义19.1.2第1题也用到了类似的操作。
        \end{itemize}
    }
\end{enumerate}

\subsection{斯托克斯公式}
\subsubsection{格林公式}
格林公式的形式

基本区域推导

多连通区域结果

计算方式

\subsubsection{保守场}
保守场、有势场、无旋场等价

势函数唯一性

有势场等价于保守场的推广($u=\sqrt{x^2+y^2}$)

三维的保守场和有势场

三维旋度与无旋场

\subsubsection{斯托克斯公式}
利用旋度的转换

斯托克斯公式的形式

曲面的选择

\subsection{高斯公式与向量场}
\subsubsection{散度}
封闭区域第二型曲面积分

散度的形式

计算方式

\subsubsection{场的微分}
梯度、散度、旋度

简单的向量运算与操作

法向量$\nv$与切向量$\vec{\tau}$

$\nv\dr S$与平面上的第二型曲面积分

平面曲线的切向量法向量

\subsubsection{非正则情况}
改变特殊点与连续性质

非正则参数化下仍可曲线曲面积分

分段光滑情况的曲线曲面积分

非光滑情况的Green、Gauss、Stolkes


\newpage
\section{常微分方程的积分}
\subsection{习题解答}
\subsubsection{作业解答}
\begin{enumerate}
    \item (习题8.5.2)计算
    $$\iint_{S^+}xz\dr y\dr z+xy\dr z\dr x+yz\dr x\dr y$$
    其中$S$为圆柱面$x^2+y^2=R^2$与平面$x=0$、$y=0$、$z=0$、$z=h$在第一卦限中围成立体的表面，这里参数$R>0$、$a>0$。

    \sol{
        分为如下步骤：
        \begin{itemize}
            \item \textbf{区域描述}

            我们先分析区域。根据``围成''的有界性要求与围成的区域必然以每个平面为边界，结合第一卦限的限制，我们可以得到五个约束
            $$x\ge0,\quad y\ge0,\quad x^2+y^2\le R^2,\quad z\ge0,\quad z\le h$$
            由于积分区域为表面，需要将每个约束取等，其他取不等号。由此我们将曲面积分分为了如下五个部分的积分之和：
            $$S_1:\quad x=0,\quad y\ge0,\quad x^2+y^2\le R^2,\quad z\ge0,\quad z\le h$$
            $$S_2:\quad y=0,\quad x\ge0,\quad x^2+y^2\le R^2,\quad z\ge0,\quad z\le h$$
            $$S_3:\quad x^2+y^2=R^2,\quad x\ge0,\quad y\ge0,\quad z\ge0,\quad z\le h$$
            $$S_4:\quad z=0,\quad x\ge0,\quad y\ge0,\quad x^2+y^2\le R^2,\quad z\le h$$
            $$S_5:\quad z=h,\quad x\ge0,\quad y\ge0,\quad x^2+y^2\le R^2,\quad z\ge0$$

            代入每个部分满足的条件可得到每个部分的化简表述：
            $$S_1:\quad x=0,\quad 0\le y\le R,\quad 0\le z\le h$$
            $$S_2:\quad y=0,\quad 0\le x\le R,\quad 0\le z\le h$$
            $$S_3:\quad x^2+y^2=R^2,\quad x\ge0,\quad y\ge0,\quad 0\le z\le h$$
            $$S_4:\quad z=0,\quad x\ge0,\quad y\ge0,\quad x^2+y^2\le R^2$$
            $$S_5:\quad z=h,\quad x\ge0,\quad y\ge0,\quad x^2+y^2\le R^2$$
            下面分别计算每个部分的积分，再求和。

            \item \textbf{参数化与定向1}
            
            $x$可以看作$y$、$z$的函数，从而有自然的参数化
            $$\vec{r}=(x,y,z)=(0,u,v)$$
            代入约束可以得到参数范围
            $$0\le u\le R,\quad 0\le v\le h$$
            将上述条件确定的$(u,v)$区域记为$D$。

            为了确定定向，用下标$u$、$v$表示对$u$、$v$求偏导并保持另一个不变，直接计算有
            $$\rv_u=(x_u,y_u,z_u)=(0,1,0),\quad\rv_v=(x_v,y_v,z_v)=(0,0,1)$$
            我们\textbf{任取}一个使得$\rv_u$、$\rv_v$\textbf{不平行}(也可表述为线性无关/叉乘非零)的$D$中点，如取$(u,v)=(0,0)$，此处有$\rv=(0,0,0)$、$\rv_u=(0,1,0)$、$\rv_v=(0,0,1)$，计算得$\rv_u\times\rv_v=(1,0,0)$，即指向$x$轴正方向，而第一卦限的曲面$x=0$处外侧应为$x$轴负方向，两者相反，因此\textbf{需要添加负号}。

            \note 只要曲面边界处$\rv_u$、$\rv_v$存在，用边界点判断定向可以得到正确的结果，详见本讲义21.3.2。
            
            \item \textbf{化为重积分与计算1}
            
            代入
            $$\dr y\dr z=\frac{D(y,z)}{D(u,v)}\dr u\dr v,\quad\dr z\dr x=\frac{D(z,x)}{D(u,v)}\dr u\dr v,\quad\dr x\dr y=\frac{D(x,y)}{D(u,v)}\dr u\dr v$$
            与$x$、$y$、$z$的参数方程可得原积分为(最左侧的负号来自定向)
            $$-\iint_D0\dr u\dr v$$
            由此为0。
            
            \item \textbf{参数化与定向2}
            
            $y$可以看作$x$、$z$的函数，从而有自然的参数化
            $$\vec{r}=(x,y,z)=(u,0,v)$$
            代入约束可以得到参数范围
            $$0\le u\le R,\quad 0\le v\le h$$
            将上述条件确定的$(u,v)$区域记为$D$。

            为了确定定向，用下标$u$、$v$表示对$u$、$v$求偏导并保持另一个不变，直接计算有
            $$\rv_u=(x_u,y_u,z_u)=(1,0,0),\quad\rv_v=(x_v,y_v,z_v)=(0,0,1)$$
            我们\textbf{任取}一个使得$\rv_u$、$\rv_v$\textbf{不平行}(也可表述为线性无关/叉乘非零)的$D$中点，如取$(u,v)=(0,0)$，此处有$\rv=(0,0,0)$、$\rv_u=(1,0,0)$、$\rv_v=(0,0,1)$，计算得$\rv_u\times\rv_v=(0,-1,0)$，即指向$y$轴负方向，而第一卦限的曲面$y=0$处外侧应为$y$轴负方向，两者相同，无需添加负号。

            \item \textbf{化为重积分与计算2}
            
            代入
            $$\dr y\dr z=\frac{D(y,z)}{D(u,v)}\dr u\dr v,\quad\dr z\dr x=\frac{D(z,x)}{D(u,v)}\dr u\dr v,\quad\dr x\dr y=\frac{D(x,y)}{D(u,v)}\dr u\dr v$$
            与$x$、$y$、$z$的参数方程可得原积分为
            $$\iint_D0\dr u\dr v$$
            由此为0。
            
            \item \textbf{参数化与定向3}
            
            由于$x^2+y^2=R^2$是圆周的一部分，而$z$无特殊限制，从而可以利用柱坐标参数化
            $$\vec{r}=(x,y,z)=(R\cos u,R\sin u,v)$$
            代入约束并结合柱坐标自然范围可以得到参数范围
            $$0\le u\le\frac{\pi}{2},\quad 0\le v\le h$$
            将上述条件确定的$(u,v)$区域记为$D$。

            为了确定定向，用下标$u$、$v$表示对$u$、$v$求偏导并保持另一个不变，直接计算有
            $$\rv_u=(x_u,y_u,z_u)=(-R\sin u,R\cos u,0),\quad\rv_v=(x_v,y_v,z_v)=(0,0,1)$$
            我们\textbf{任取}一个使得$\rv_u$、$\rv_v$\textbf{不平行}(也可表述为线性无关/叉乘非零)的$D$中点，如取$(u,v)=(0,0)$，此处有$\rv=(R,0,0)$、$\rv_u=(0,R,0)$、$\rv_v=(0,0,1)$，计算得$\rv_u\times\rv_v=(R,0,0)$，即指向$x$轴正方向。根据约束此点处$x$取到最大值，因此$x$轴正方向的确指向外侧，无需添加负号。

            \item \textbf{化为重积分与计算3}
            
            代入
            $$\dr y\dr z=\frac{D(y,z)}{D(u,v)}\dr u\dr v,\quad\dr z\dr x=\frac{D(z,x)}{D(u,v)}\dr u\dr v,\quad\dr x\dr y=\frac{D(x,y)}{D(u,v)}\dr u\dr v$$
            与$x$、$y$、$z$的参数方程可得原积分为
            $$\iint_D(vR\cos u\cdot R\cos u+R^2\cos u\sin u\cdot R\sin u)\dr u\dr v$$
            由区域描述可直接写出累次积分
            $$\int_0^{\pi/2}\dr u\int_0^h(R^2v\cos^2u+R^3\cos u\sin^2u)\dr v=\int_0^{\pi/2}\bigg(\frac{1}{2}h^2R^2\cos^2u+hR^3\cos u\sin^2u\dr u\bigg)$$
            换元$u=\frac{\pi}{2}-t$并利用上册教材3.4节例2结论可得
            $$\int_0^{\pi/2}\cos^2u\dr u=\int_0^{\pi/2}\sin^2t\dr t=\frac{\pi}{4}$$
            换元$t=\sin u$可得
            $$\int_0^{\pi/2}\cos u\sin^2u\dr u=\int_0^1t^2\dr t=\frac{1}{3}$$
            综合得结果为
            $$\frac{\pi}{8}h^2R^2+\frac{1}{3}hR^3$$
            
            \item \textbf{参数化与定向4}
            
            $z$可以看作$x$、$y$的函数，从而有自然的参数化
            $$\vec{r}=(x,y,z)=(u,v,0)$$
            代入约束可以得到参数范围
            $$u\ge0,\quad v\ge0,\quad u^2+v^2\le R^2$$
            将上述条件确定的$(u,v)$区域记为$D$。

            为了确定定向，用下标$u$、$v$表示对$u$、$v$求偏导并保持另一个不变，直接计算有
            $$\rv_u=(x_u,y_u,z_u)=(1,0,0),\quad\rv_v=(x_v,y_v,z_v)=(0,1,0)$$
            我们\textbf{任取}一个使得$\rv_u$、$\rv_v$\textbf{不平行}(也可表述为线性无关/叉乘非零)的$D$中点，如取$(u,v)=(0,0)$，此处有$\rv=(0,0,0)$、$\rv_u=(1,0,0)$、$\rv_v=(0,1,0)$，计算得$\rv_u\times\rv_v=(0,0,1)$，即指向$z$轴正方向，而第一卦限的曲面$z=0$处外侧应为$z$轴负方向，两者相反，因此\textbf{需要添加负号}。

            \item \textbf{化为重积分与计算4}
            
            代入
            $$\dr y\dr z=\frac{D(y,z)}{D(u,v)}\dr u\dr v,\quad\dr z\dr x=\frac{D(z,x)}{D(u,v)}\dr u\dr v,\quad\dr x\dr y=\frac{D(x,y)}{D(u,v)}\dr u\dr v$$
            与$x$、$y$、$z$的参数方程可得原积分为(最左侧的负号来自定向)
            $$-\iint_D0\dr u\dr v$$
            由此为0。
            
            \item \textbf{参数化与定向5}
            
            $z$可以看作$x$、$y$的函数，从而有自然的参数化
            $$\vec{r}=(x,y,z)=(u,v,h)$$
            代入约束可以得到参数范围
            $$u\ge0,\quad v\ge0,\quad u^2+v^2\le R^2$$
            将上述条件确定的$(u,v)$区域记为$D$。

            为了确定定向，用下标$u$、$v$表示对$u$、$v$求偏导并保持另一个不变，直接计算有
            $$\rv_u=(x_u,y_u,z_u)=(1,0,0),\quad\rv_v=(x_v,y_v,z_v)=(0,1,0)$$
            我们\textbf{任取}一个使得$\rv_u$、$\rv_v$\textbf{不平行}(也可表述为线性无关/叉乘非零)的$D$中点，如取$(u,v)=(0,0)$，此处有$\rv=(0,0,h)$、$\rv_u=(1,0,0)$、$\rv_v=(0,1,0)$，计算得$\rv_u\times\rv_v=(0,0,1)$，即指向$z$轴正方向，而第一卦限的曲面$z=h$处外侧应为$z$轴正方向，两者相同，无需添加负号。

            \item \textbf{化为重积分与计算5}
            
            代入
            $$\dr y\dr z=\frac{D(y,z)}{D(u,v)}\dr u\dr v,\quad\dr z\dr x=\frac{D(z,x)}{D(u,v)}\dr u\dr v,\quad\dr x\dr y=\frac{D(x,y)}{D(u,v)}\dr u\dr v$$
            与$x$、$y$、$z$的参数方程可得原积分为
            $$\iint_Dvh\dr u\dr v$$
            由于$D$的旋转对称性，考虑\textbf{极坐标换元}
            $$u=r\cos\theta,\quad v=r\sin\theta$$
            此时有
            $$\dr u\dr v=r\dr r\dr\theta$$
            约束条件即第一象限的扇形内，可化简为
            $$0\le r\le R,\quad 0\le\theta\le\frac{\pi}{2}$$
            被积函数为$hr\sin\theta$，再结合极坐标自然范围可将积分化为
            $$\int_0^{R}\dr r\int_0^{\pi/2}hr^2\sin\theta\dr\theta$$
            对$\theta$积分得到结果为
            $$\int_0^Rhr^2\dr r=\frac{1}{3}hR^3$$

            \item \textbf{综合结果}
            
            将五个部分求和可知最终的积分结果
            $$\frac{\pi}{8}h^2R^2+\frac{3}{2}hR^3$$
        \end{itemize}
    }

    \item (习题8.5.4)计算
    $$\iint_{S^+}z^2\dr y\dr z+x\dr z\dr x-3z\dr x\dr y$$
    其中$S$为曲面$z=4-y^2$、平面$x=0$、$x=1$、$z=0$围成立体的表面。

    \sol{
        分为如下步骤：
        \begin{itemize}
            \item \textbf{区域描述}
            
            我们先分析区域。根据``围成''的有界性要求与围成的区域必然以每个平面为边界，我们可以得到四个约束(这里$0\le z\le 4-y^2$是由于否则$y$将无界)
            $$x\ge0,\quad x\le1,\quad z\ge0,\quad z\le 4-y^2$$
            由于积分区域为表面，需要将每个约束取等，其他取不等号。由此我们将曲面积分分为了如下四个部分的积分之和：
            $$S_1:\quad x=0,\quad x\le1,\quad z\ge0,\quad z\le 4-y^2$$
            $$S_2:\quad x=1,\quad x\ge0,\quad z\ge0,\quad z\le 4-y^2$$
            $$S_3:\quad z=0,\quad x\ge0,\quad x\le1,\quad z\le 4-y^2$$
            $$S_4:\quad z=4-y^2,\quad x\ge0,\quad x\le1,\quad z\ge0$$

            代入每个部分满足的条件可得到每个部分的化简表述：
            $$S_1:\quad x=0,\quad 0\le z\le 4-y^2$$
            $$S_2:\quad x=1,\quad 0\le z\le 4-y^2$$
            $$S_3:\quad z=0,\quad 0\le x\le1,\quad y^2\le4$$
            $$S_4:\quad z=4-y^2,\quad 0\le x\le1,\quad y^2\le4$$
            下面分别计算每个部分的积分，再求和。
            
            \item \textbf{参数化与定向1}
            
            $x$可以看作$y$、$z$的函数，从而有自然的参数化
            $$\vec{r}=(x,y,z)=(0,u,v)$$
            代入约束可以得到参数范围
            $$0\le v\le 4-u^2$$
            将上述条件确定的$(u,v)$区域记为$D$。

            为了确定定向，用下标$u$、$v$表示对$u$、$v$求偏导并保持另一个不变，直接计算有
            $$\rv_u=(x_u,y_u,z_u)=(0,1,0),\quad\rv_v=(x_v,y_v,z_v)=(0,0,1)$$
            我们\textbf{任取}一个使得$\rv_u$、$\rv_v$\textbf{不平行}(也可表述为线性无关/叉乘非零)的$D$中点，如取$(u,v)=(0,0)$，此处有$\rv=(0,0,0)$、$\rv_u=(0,1,0)$、$\rv_v=(0,0,1)$，计算得$\rv_u\times\rv_v=(1,0,0)$，即指向$x$轴正方向，而由于$x$的范围是$0\le x\le1$，此处为$x$最小处，外侧应为$x$轴负方向，两者相反，因此\textbf{需要添加负号}。
            
            \item \textbf{化为重积分与计算1}
            
            代入
            $$\dr y\dr z=\frac{D(y,z)}{D(u,v)}\dr u\dr v,\quad\dr z\dr x=\frac{D(z,x)}{D(u,v)}\dr u\dr v,\quad\dr x\dr y=\frac{D(x,y)}{D(u,v)}\dr u\dr v$$
            与$x$、$y$、$z$的参数方程可得原积分为(最左侧的负号来自定向)
            $$-\iint_Dv^2\dr u\dr v$$
            这里可以直接进行计算，不过如果可以\textbf{观察}到本题的某种对称性，事实上无需计算这个积分。我们接下来对$S_2$进行处理以说明为何无需计算积分：

            \item \textbf{参数化与定向2}
            
            $x$可以看作$y$、$z$的函数，从而有自然的参数化
            $$\vec{r}=(x,y,z)=(1,u,v)$$
            代入约束可以得到参数范围
            $$0\le v\le 4-u^2$$
            将上述条件确定的$(u,v)$区域记为$D$。

            为了确定定向，用下标$u$、$v$表示对$u$、$v$求偏导并保持另一个不变，直接计算有
            $$\rv_u=(x_u,y_u,z_u)=(0,1,0),\quad\rv_v=(x_v,y_v,z_v)=(0,0,1)$$
            我们\textbf{任取}一个使得$\rv_u$、$\rv_v$\textbf{不平行}(也可表述为线性无关/叉乘非零)的$D$中点，如取$(u,v)=(0,0)$，此处有$\rv=(1,0,0)$、$\rv_u=(0,1,0)$、$\rv_v=(0,0,1)$，计算得$\rv_u\times\rv_v=(1,0,0)$，即指向$x$轴正方向，而由于$x$的范围是$0\le x\le1$，此处为$x$最大处，外侧应为$x$轴正方向，两者相同，无需添加负号。
            
            \item \textbf{化为重积分与计算2}
            
            代入
            $$\dr y\dr z=\frac{D(y,z)}{D(u,v)}\dr u\dr v,\quad\dr z\dr x=\frac{D(z,x)}{D(u,v)}\dr u\dr v,\quad\dr x\dr y=\frac{D(x,y)}{D(u,v)}\dr u\dr v$$
            与$x$、$y$、$z$的参数方程可得原积分为
            $$\iint_Dv^2\dr u\dr v$$
            可以发现，此积分与分析$S_1$得到的积分形式相同，区域描述也完全相同，但并没有添加定向带来的负号。由此，在$S_1$、$S_2$上的积分互为相反数，求和为0。

            \item \textbf{参数化与定向3}
            
            $z$可以看作$x$、$y$的函数，从而有自然的参数化
            $$\vec{r}=(x,y,z)=(u,v,0)$$
            代入约束可以得到参数范围
            $$0\le u\le 1,\quad -2\le v\le 2$$
            将上述条件确定的$(u,v)$区域记为$D$。

            为了确定定向，用下标$u$、$v$表示对$u$、$v$求偏导并保持另一个不变，直接计算有
            $$\rv_u=(x_u,y_u,z_u)=(1,0,0),\quad\rv_v=(x_v,y_v,z_v)=(0,1,0)$$
            我们\textbf{任取}一个使得$\rv_u$、$\rv_v$\textbf{不平行}(也可表述为线性无关/叉乘非零)的$D$中点，如取$(u,v)=(0,0)$，此处有$\rv=(0,0,0)$、$\rv_u=(1,0,0)$、$\rv_v=(0,1,0)$，计算得$\rv_u\times\rv_v=(0,0,1)$，即指向$z$轴正方向，而由于$z$的范围是$0\le z\le4-y^2$，此处为$z$最小处，外侧应为$z$轴负方向，两者相反，因此\textbf{需要添加负号}。

            \item \textbf{化为重积分与计算3}
            
            代入
            $$\dr y\dr z=\frac{D(y,z)}{D(u,v)}\dr u\dr v,\quad\dr z\dr x=\frac{D(z,x)}{D(u,v)}\dr u\dr v,\quad\dr x\dr y=\frac{D(x,y)}{D(u,v)}\dr u\dr v$$
            与$x$、$y$、$z$的参数方程可得原积分为(最左侧的负号来自定向)
            $$-\iint_D0\dr u\dr v$$
            由此为0。

            \item \textbf{参数化与定向4}
            
            $z$可以看作$x$、$y$的函数，从而有自然的参数化
            $$\vec{r}=(x,y,z)=(u,v,4-v^2)$$
            代入约束可以得到参数范围
            $$0\le u\le 1,\quad -2\le v\le 2$$
            将上述条件确定的$(u,v)$区域记为$D$。

            为了确定定向，用下标$u$、$v$表示对$u$、$v$求偏导并保持另一个不变，直接计算有
            $$\rv_u=(x_u,y_u,z_u)=(1,0,0),\quad\rv_v=(x_v,y_v,z_v)=(0,1,-2v)$$
            我们\textbf{任取}一个使得$\rv_u$、$\rv_v$\textbf{不平行}(也可表述为线性无关/叉乘非零)的$D$中点，如取$(u,v)=(0,0)$，此处有$\rv=(0,0,4)$、$\rv_u=(1,0,0)$、$\rv_v=(0,1,0)$，计算得$\rv_u\times\rv_v=(0,0,1)$，即指向$z$轴正方向，而而由于$z$的范围是$0\le z\le4-y^2$，此处为$z$最大处，外侧应为$z$轴正方向，两者相同，无需添加负号。

            \item \textbf{化为重积分与计算4}
            
            代入
            $$\dr y\dr z=\frac{D(y,z)}{D(u,v)}\dr u\dr v,\quad\dr z\dr x=\frac{D(z,x)}{D(u,v)}\dr u\dr v,\quad\dr x\dr y=\frac{D(x,y)}{D(u,v)}\dr u\dr v$$
            与$x$、$y$、$z$的参数方程可得原积分为
            $$\iint_D(2uv-3(4-v^2))\dr u\dr v$$
            由区域描述可直接写出累次积分并计算
            $$\int_0^1\dr u\int_{-2}^2(2uv-12+3v^2)\dr v=\int_0^1(-48+16)\dr u=-32$$

            \item \textbf{综合结果}
            
            将四个部分求和可知最终的积分结果$-32$。
        \end{itemize}
    }

    \item (习题8.5.6)计算
    $$\iint_S(y-z)\dr y\dr z+(z-x)\dr z\dr x+(x-y)\dr x\dr y$$
    其中$S$为球面$x^2+y^2+z^2=2Rx$在$z\ge0$的部分被柱面$x^2+y^2=2rx$所截部分的上侧，这里参数$R>r>0$。

    \sol{
        分为如下步骤：
        \begin{itemize}
            \item \textbf{区域描述}
            
            若$x^2+y^2=2rx$对应的约束为$x^2+y^2\ge 2rx$，利用几何关系或代数计算可以发现此时确定的区域将穿过$z=0$。由于``在$z\ge0$部分''应当理解为，$z=0$不为区域的边界，但区域恒满足$z\ge0$，可知矛盾，从而实际约束为
            $$x^2+y^2+z^2=2Rx,\quad x^2+y^2\le 2rx,\quad z\ge0$$

            \item \textbf{参数化}
            
            \note 本题看似可以进行\textbf{平移的球坐标}参数化，但之后会发现区域描述非常复杂，由此直接计算或使用柱坐标是更好的。
            
            $z$可以看作$x$、$y$的函数，从而有自然的参数化
            $$\vec{r}=(x,y,z)=(u,v,\sqrt{2Ru-u^2-v^2})$$
            代入约束可以得到参数范围
            $$u^2+v^2\le 2ru$$
            将上述条件确定的$(u,v)$区域记为$D$。

            \item \textbf{定向}
            
            用下标$u$、$v$表示对$u$、$v$求偏导并保持另一个不变，直接计算有
            $$\rv_u=(x_u,y_u,z_u)=\bigg(1,0,\frac{R-u}{\sqrt{2Ru-u^2-v^2}}\bigg)$$
            $$\rv_v=(x_v,y_v,z_v)=(0,1,-\frac{v}{\sqrt{2Ru-u^2-v^2}})$$
            我们\textbf{任取}一个使得$\rv_u$、$\rv_v$\textbf{不平行}(也可表述为线性无关/叉乘非零)的$D$中点，如取$(u,v)=(r,0)$，此处有$\rv=(r,0,\sqrt{2Rr-r^2})$、$\rv_u=\big(1,0,\sqrt{\frac{R-r}{2R}}\big)$、$\rv_v=(0,1,0)$。在曲面上的点$\rv$处画出向量$\rv_u$与$\rv_v$，通过右手定则可发现$\rv_u\times\rv_v$指向上方，与要求的定向相同，无需添加负号。

            \note 事实上还有更简单的无需画图的判断方法：上侧代表取$z$分量大于0的法向量，直接计算可发现$\rv_u\times\rv_v$的$z$分量是1，大于0，因此与要求的定向相同。

            \item \textbf{化为重积分}
            
            代入
            $$\dr y\dr z=\frac{D(y,z)}{D(u,v)}\dr u\dr v,\quad\dr z\dr x=\frac{D(z,x)}{D(u,v)}\dr u\dr v,\quad\dr x\dr y=\frac{D(x,y)}{D(u,v)}\dr u\dr v$$
            与$x$、$y$、$z$的参数方程可得原积分为(这里将$\sqrt{2Ru-u^2-v^2}$写为$z$以简化书写)
            $$\iint_D\bigg((v-z)\frac{u-R}{z}+(z-u)\frac{v}{z}+(u-v)\bigg)\dr u\dr v$$
            化简得
            $$\iint_DR\dr u\dr v$$

            \item \textbf{积分计算}
            
            利用重积分的几何意义，$\iint_D\dr u\dr v$即为区域的面积，而此积分区域是圆$(u-r)^2+v^2\le r^2$，从而面积为$\pi r^2$，最终得到积分结果
            $$\pi r^2R$$
        \end{itemize}

        \note 本题的五个步骤是第二型曲线或曲面积分的一个\textbf{标准}的处理流程，而计算第一型曲线或曲面积分只需去除``定向''步骤。
    }

    \item (习题8.5.7)计算
    $$\iint_{S^+}(x+a)\dr y\dr z+(y+b)\dr z\dr x+(z+c)\dr x\dr y$$
    其中$S$为球面$x^2+y^2+z^2=R^2$，这里参数$R>0$，$a,b,c\in\rb$。

    \sol{
        我们介绍两种方法。首先是直接计算，分为如下步骤：
        \begin{itemize}
            \item \textbf{参数化}
            
            利用球坐标可得到球面的一个参数化
            $$\rv=(x,y,z)=(R\cos u\sin v,R\sin u\sin v,R\cos v)$$
            球坐标自然范围即为参数范围
            $$0\le u\le2\pi,\quad0\le v\le\pi$$
            将上述条件确定的$(u,v)$区域记为$D$。

            \item \textbf{定向}
            
            用下标$u$、$v$表示对$u$、$v$求偏导并保持另一个不变，直接计算有
            $$\rv_u=(x_u,y_u,z_u)=(-R\sin u\sin v,R\cos u\sin v,0)$$
            $$\rv_v=(x_v,y_v,z_v)=(R\cos u\cos v,R\sin u\cos v,-R\sin v)$$
            我们\textbf{任取}一个使得$\rv_u$、$\rv_v$\textbf{不平行}(也可表述为线性无关/叉乘非零)的$D$中点，如取$(u,v)=(0,\frac{\pi}{2})$，此处有$\rv=(R,0,0)$、$\rv_u=(0,R,0)$、$\rv_v=(0,0,-R)$，计算得$\rv_u\times\rv_v=(-R^2,0,0)$，即指向$x$轴负方向，而由于$x=R$处为$x$最大处，外侧应为$x$轴正方向，两者相反，因此\textbf{需要添加负号}。
            
            \item \textbf{化为重积分}
            
            代入
            $$\dr y\dr z=\frac{D(y,z)}{D(u,v)}\dr u\dr v,\quad\dr z\dr x=\frac{D(z,x)}{D(u,v)}\dr u\dr v,\quad\dr x\dr y=\frac{D(x,y)}{D(u,v)}\dr u\dr v$$
            与$x$、$y$、$z$的参数方程并计算化简可得原积分为(这里已经添加了定向带来的负号)
            $$\iint_D(R^2+aR\cos u\sin v+bR\sin u\sin v+cR\cos v)\sin v\dr u\dr v$$

            \item \textbf{积分计算}
            
            由区域描述可直接写出累次积分(先对$u$积分以消去更多项)
            $$\int_0^\pi\dr v\int_0^{2\pi}(R^2+aR\cos u\sin v+bR\sin u\sin v+cR\cos v)\sin v\dr u$$
            对$u$积分得到
            $$2\pi\int_0^{\pi}(R^2\sin v+cR\sin v\cos v)\dr v$$
            将$\sin v\cos v$写为$\frac{1}{2}\sin(2v)$，由此再对$v$积分即得结果为
            $$4\pi R^2$$
        \end{itemize}

        不过，本题还可以利用\textbf{第二型曲面积分与第一型的关系}简化计算，此时计算将更加简单：
        \begin{itemize}
            \item \textbf{化为第一型曲面积分}

            利用第一型曲面积分与第二型曲面积分的关系可知(这里$\vec{n}$代表第二型曲面积分对应定向的单位法向量)
            $$\vec{n}\dr S=(\dr y\dr z,\dr z\dr x,\dr x\dr y)$$
            利用球面的几何性质可知它$(x,y,z)$处的的一个外法向为$(x,y,z)$，从而指向外侧单位法向量为(利用球面上的点模长恒为$R$)
            $$\frac{1}{R}(x,y,z)$$
            由此即得原积分为
            $$\iint_S(x+a,y+b,z+c)\cdot\frac{1}{R}(x,y,z)\dr S=\frac{1}{R}\iint_S(R^2+ax+by+cz)\dr S$$

            \item \textbf{积分计算}
            
            由于$(x,y,z)\in S$当且仅当$(-x,y,z)\in S$，与三重积分反射对称性完全类似分析可得(详见本讲义21.3.3)
            $$\iint_Sx\dr S=\iint_S(-x)\dr S$$
            从而两侧必然为0，同理
            $$\iint_Sy\dr S=\iint_Sz\dr S=0$$
            由此积分化为
            $$\iint_SR\dr S$$
            利用第一型曲面积分的几何意义，$\iint_S\dr S$即为曲面面积，而$S$是半径为$R$的球面，从而面积为$4\pi R$，最终得到积分结果
            $$4\pi R^2$$
        \end{itemize}

        \note 对平面、球面、圆柱面的第二型曲面积分，由于法向量易于得到，往往可以考虑转化为第一型曲面积分计算。
    }

    \item (习题8.5.10)计算
    $$\iint_Sx\dr y\dr z+y\dr z\dr x+z\dr x\dr y$$
    其中$S$为螺旋面
    $$\begin{cases}x=u\cos v\\y=u\sin v\\z=cv\end{cases}(0\le a\le u\le b,\ 0\le v\le2\pi)$$
    的上侧，这里参数$c\in\rb$。
    
    \sol{
        分为如下步骤：
        \begin{itemize}
            \item \textbf{定向}
            
            用下标$u$、$v$表示对$u$、$v$求偏导并保持另一个不变，直接计算有
            $$\rv_u=(x_u,y_u,z_u)=(\cos v,\sin v,0)$$
            $$\rv_v=(x_v,y_v,z_v)=(-u\sin v,u\cos v,c)$$
            我们\textbf{任取}一个使得$\rv_u$、$\rv_v$\textbf{不平行}(也可表述为线性无关/叉乘非零)的$D$中点，如取$(u,v)=(a,0)$，此处有$\rv=(a,0,0)$、$\rv_u=(1,0,0)$、$\rv_v=(0,a,c)$，计算得$\rv_u\times\rv_v=(0,-c,a)$，由于上侧要求法向量的$z$分量大于0，而$a>0$已经满足要求，无需添加负号。
            
            \item \textbf{化为重积分}
            
            代入
            $$\dr y\dr z=\frac{D(y,z)}{D(u,v)}\dr u\dr v,\quad\dr z\dr x=\frac{D(z,x)}{D(u,v)}\dr u\dr v,\quad\dr x\dr y=\frac{D(x,y)}{D(u,v)}\dr u\dr v$$
            与$x$、$y$、$z$的参数方程并计算化简可得原积分为
            $$\iint_D(uc\cos v\sin v-uc\sin v\cos v+cuv)\dr u\dr v=\iint_Dcuv\dr u\dr v$$

            \item \textbf{积分计算}
            
            由区域描述可直接写出累次积分
            $$\int_a^b\dr u\int_0^{2\pi}cuv\dr v$$
            对$v$积分得到
            $$2\pi^2\int_a^bcu\dr u=\pi^2c(a^2-b^2)$$
        \end{itemize}
    }

    \item (习题8.6.1(1))利用高斯公式计算
    $$\iint_Sx^2\dr y\dr z+y^2\dr z\dr x+z^2\dr x\dr y$$
    其中$S$为不包含上下底的圆柱面
    $$x^2+y^2=a^2,\quad z\in[0,h]$$
    的外侧($x^2+y^2\ge a^2$侧)，这里参数$a>0$、$h>0$。
    
    \sol{
        分为如下步骤：
        \begin{itemize}
            \item \textbf{高斯公式}
            
            将被积函数看作$P\dr y\dr z+Q\dr z\dr x+R\dr x\dr y$，由初等函数知$P$、$Q$、$R$在$\rb^3$上可导且导函数连续，直接计算得
            $$P_x+Q_y+R_z=2(x+y+z)$$
            由于积分区域并非封闭曲面，我们补上顶部的
            $$S_1:\quad z=h,\quad x^2+y^2\le a^2$$
            与底部的
            $$S_2:\quad z=0,\quad x^2+y^2\le a^2$$
            并在顶部取上侧为正向、底部取下侧为正向。

            根据定义$S_1$在$S$的上方，$S_2$在$S$的下方，因此$S_1$的上侧、$S$的外侧、$S_2$的下侧构成三者围成的封闭曲面的外侧。设围成的区域为$\Omega$，所求的积分为$I$，并记
            $$I_1=\iint_{S_1}x^2\dr y\dr z+y^2\dr z\dr x+z^2\dr x\dr y$$
            $$I_2=\iint_{S_2}x^2\dr y\dr z+y^2\dr z\dr x+z^2\dr x\dr y$$
            $$I_3=\iiint_\Omega2(x+y+z)\dr V$$
            利用高斯公式有
            $$I_3=I+I_1+I_2$$
            从而$I=I_3-I_1-I_2$。

            \item \textbf{$I_1$计算}
            
            由于$I_1$是平面区域上的第二类曲面积分，化为第一类可方便计算。利用第一型曲面积分与第二型曲面积分的关系可知(这里$\nv$代表第二型曲面积分对应定向的单位法向量)
            $$\nv\dr S=(\dr y\dr z,\dr z\dr x,\dr x\dr y)$$   
            根据几何可知$z=h$平面的一个上侧的单位法向量是$(0,0,1)$，由此即得原积分为(已代入$z=h$)
            $$\iint_{S_1}(x^2,y^2,z^2)\cdot(\dr y\dr z,\dr z\dr x,\dr x\dr y)=\iint_{S_0}(x^2,y^2,h^2)\cdot(0,0,1)\dr S=h^2\iint_{S_1}\dr S$$
            利用第一型曲面积分的几何意义，$\iint_{S_1}\dr S$即为曲面面积，而$S_1$是圆柱的顶面，是一个半径为$a$的圆，从而面积为$\pi a^2$，最终得到
            $$I_1=\pi a^2h^2$$

            \item \textbf{$I_2$计算}
            
            由于$I_2$是平面区域上的第二类曲面积分，化为第一类可方便计算。利用第一型曲面积分与第二型曲面积分的关系可知(这里$\nv$代表第二型曲面积分对应定向的单位法向量)
            $$\nv\dr S=(\dr y\dr z,\dr z\dr x,\dr x\dr y)$$   
            根据几何可知$z=0$平面的一个下侧的单位法向量是$(0,0,-1)$，由此即得原积分为(已代入$z=0$)
            $$I_2=\iint_{S_2}(x^2,y^2,z^2)\cdot(\dr y\dr z,\dr z\dr x,\dr x\dr y)=\iint_{S_2}(x^2,y^2,0)\cdot(0,0,-1)\dr S=0$$

            \item \textbf{$I_3$计算}
            
            由于$\Omega$满足$(x,y,z)\in\Omega$当且仅当$(-x,y,z)\in\Omega$，利用区域折半可知
            $$\iiint_\Omega x\dr V=0$$
            同理
            $$\iiint_\Omega y\dr V=0$$
            由此可知
            $$I_3=2\iiint_\Omega z\dr x\dr y\dr z$$
            由于此被积函数对$x$、$y$为常数，先对$x$、$y$积分，再对$z$积分，将其写成
            $$I_3=2\int_0^h \dr z\iint_Dz\dr x\dr y$$
            $D$为$(x,y)$的区域$x^2+y^2\le a^2$，利用重积分的几何意义可知$\iint_D\dr x\dr y$为其面积$\pi a^2$，从而积分结果为
            $$I_3=2\pi a^2\int_0^hz\dr z=\pi a^2h^2$$

            \item \textbf{综合结果}
            
            将三个部分综合可知最终的积分结果
            $$I=I_3-I_1-I_2=\pi a^2h^2-0-\pi a^2h^2=0$$
        \end{itemize}
    }

    \item (习题8.6.4)计算
    $$\oiint_{S^+}x^3\dr y\dr z+y^3\dr z\dr x+z^3\dr x\dr y$$
    其中$S$为球面$x^2+y^2+z^2=a^2$，这里参数$a>0$。

    \sol{
        分为如下步骤：
        \begin{itemize}
            \item \textbf{高斯公式}
            
            将被积函数看作$P\dr y\dr z+Q\dr z\dr x+R\dr x\dr y$，由初等函数知$P$、$Q$、$R$在$\rb^3$上可导且导函数连续，直接计算得
            $$P_x+Q_y+R_z=3(x^2+y^2+z^2)$$

            由此设球体$x^2+y^2+z^2\le a^2$为区域$\Omega$，利用高斯公式可将积分化为
            $$3\iiint_\Omega(x^2+y^2+z^2)\dr V$$

            \item \textbf{积分计算}
            
            根据区域的旋转对称性，可以考虑球坐标换元
            $$x=\rho\sin\varphi\cos\theta,\quad y=\rho\sin\varphi\sin\theta,\quad z=\rho\cos\varphi$$
            
            将约束转化为$\rho^2\le a^2$后，结合球坐标自然范围可得最终范围
            $$0\le\rho\le a,\quad0\le\varphi\le\pi,\quad 0\le\theta\le2\pi$$        
            将上述条件确定的$(\rho,\varphi,\theta)$区域记为$\Omega_0$。此时被积函数成为$\rho^2$，直接计算Jacobi行列式可知
            $$\dr x\dr y\dr z=\rho^2\sin\varphi\dr\rho\dr\varphi\dr\theta$$
            原积分化为
            $$3\iiint_{\Omega_0}\rho^4\sin\varphi\dr\rho\dr\varphi\dr\theta$$
            利用区域描述可直接写出累次积分
            $$\int_0^a\dr\rho\int_0^\pi\dr\varphi\int_0^{2\pi}\rho^4\sin\varphi\dr\theta$$
            依次积分得到其为
            $$6\pi\int_0^a\dr\rho\int_0^\pi\rho^4\sin\varphi\dr\varphi=12\pi\int_0^a\rho^4\dr\rho=\frac{12}{5}\pi a^5$$
        \end{itemize}
    }

    \item (习题8.6.7)计算
    $$\oiint_{S^+}(xy^2+y+z)\dr y\dr z+(yz^2+xz)\dr z\dr x+(zx^2+5x^2y^2)\dr x\dr y$$
    其中$S$为椭球面$\frac{x^2}{a^2}+\frac{y^2}{b^2}+\frac{z^2}{c^2}=1$，这里参数$a>0$、$b>0$、$c>0$。

    \sol{
        分为如下步骤：
        \begin{itemize}
            \item \textbf{高斯公式}
            
            将被积函数看作$P\dr y\dr z+Q\dr z\dr x+R\dr x\dr y$，由初等函数知$P$、$Q$、$R$在$\rb^3$上可导且导函数连续，直接计算得
            $$P_x+Q_y+R_z=y^2+z^2+x^2$$

            由此设椭球体$\frac{x^2}{a^2}+\frac{y^2}{b^2}+\frac{z^2}{c^2}\le1$为区域$\Omega$，利用高斯公式可将积分化为
            $$\iiint_\Omega(x^2+y^2+z^2)\dr V$$

            \item \textbf{积分计算}
            
            由区域特性考虑\textbf{广义球坐标}换元
            $$x=a\rho\sin\varphi\cos\theta,\quad y=b\rho\sin\varphi\sin\theta,\quad z=c\rho\cos\varphi$$
            这样约束条件即化为了$\rho^2\le1$，结合球坐标自然范围有
            $$0\le\rho\le1,\quad0\le\varphi\le\pi,\quad0\le\theta\le2\pi$$
            将上述条件确定的$(\rho,\varphi,\theta)$区域记为$\Omega_0$，被积函数仍为1。 

            直接类似球坐标计算Jacobi行列式可知
            $$\dr x\dr y\dr z=abc\rho^2\sin\varphi\dr\rho\dr\varphi\dr\theta$$
            于是写出累次积分(这里积分顺序可任取)
            $$\int_0^\pi\dr\varphi\int_0^{2\pi}\dr\theta\int_0^1\rho^4abc(a^2\sin^2\varphi\cos^2\theta+b^2\sin^2\varphi\sin^2\theta+c^2\cos^2\varphi)\sin\varphi\dr\rho$$
            对$\rho$积分得到
            $$\frac{abc}{5}\int_0^\pi\dr\varphi\int_0^{2\pi}(a^2\sin^3\varphi\cos^2\theta+b^2\sin^3\varphi\sin^2\theta+c^2\cos^2\varphi\sin\varphi)\dr\theta$$
            类似\exref{sin 2 and 4 to 2pi}并利用3.4节例2可算出
            $$\int_0^{2\pi}\cos^2\theta\dr\theta=\int_0^{2\pi}\sin^2\theta\dr\theta=\pi$$
            $$\int_0^\pi\sin^3\varphi\dr\varphi=2\int_0^{\pi/2}\sin^3\varphi\dr\varphi=\frac{4}{3}$$
            由此对$\theta$的积分结果为
            $$\frac{abc\pi}{5}\int_0^\pi(a^2\sin^2\varphi+b^2\sin^2\varphi+2c^2\cos^2\varphi\sin\varphi)\dr\varphi$$
            进一步计算可将其化为
            $$\frac{abc\pi}{5}\bigg(\frac{4(a^2+b^2)}{3}+2c^2\int_0^\pi\cos^2\varphi\sin\varphi\dr\varphi\bigg)$$
            换元$t=\cos\varphi$可得
            $$\int_0^\pi\cos^2\varphi\sin\varphi\dr\varphi=\int_{-1}^1t^2\dr t=\frac{2}{3}$$
            代入最终得到结果为
            $$\frac{4\pi abc}{15}(a^2+b^2+c^2)$$
        \end{itemize}
    }

    \item (习题8.6.11)若$u(x,y,z)$在$\rb^3$上有连续的二阶偏导数，且
    $$\forall(x,y,z)\in\rb^3,\quad\triangle u(x,y,z)=0$$
    设$S$为光滑闭曲面，围成的区域为$\Omega$，证明(以下省略函数的参数$(x,y,z)$)：
    \begin{enumerate}[(1)]
        \item $$\oiint_Su\frac{\partial u}{\partial\nv}\dr S=\iiint_\Omega\big((\partial_1u)^2+(\partial_2u)^2+(\partial_3u)^3\big)\dr V$$
        其中$\frac{\partial u}{\partial\nv}$表示$u$沿$S$外法向的方向导数。

        \sol{
            我们先分析再证明：
            \begin{itemize}
                \item \textbf{问题分析}
                
                为了将曲面上的积分和区域内的积分产生关联，我们必须使用\textbf{高斯公式}，而原本右侧为第一型曲面积分，必须化为\textbf{第二型曲面积分}才能使用高斯公式。

                设$\nv=(n_1,n_2,n_3)$代表曲面外侧的单位法向量，利用第一型曲面积分与第二型曲面积分的关系可知(以外侧为正定向)
                $$(n_1\dr S,n_2\dr S,n_3\dr S)=(\dr y\dr z,\dr z\dr x,\dr x\dr y)$$

                \item \textbf{化为第二型曲面积分}
                
                由$u$可微，方向导数可以写为(可见上学期讲义14.4.2)
                $$n_1\partial_1u+n_2\partial_2u+n_3\partial_3u$$
                
                \note 也可写为向量形式$\nabla u\cdot\vec{n}$。

                由此可以将
                $$\oiint_Su\frac{\partial u}{\partial\nv}\dr S$$
                展开写为
                $$\oiint_S(u\partial_1un_1\dr S+u\partial_2un_2\dr S+u\partial_3un_3\dr S)$$
                从而其等于
                $$\oiint_S(u\partial_1u\dr y\dr z+u\partial_2u\dr z\dr x+u\partial_3u\dr x\dr y)$$

                \item \textbf{高斯公式}
                
                由$u$有连续的二阶偏导数，符合使用条件，从而上式可利用高斯公式化为
                $$\iint_\Omega(\partial_1(u\partial_1u)+\partial_2(u\partial_2u)+\partial_3(u\partial_3u))\dr V$$
                利用乘积求导公式得被积函数为
                $$(\partial_1u)^2+(\partial_2u)^2+(\partial_3u)^3+u(\partial_1^2u+\partial_2^2u+\partial_3^2u)$$
                再根据$\triangle u=\partial_1^2u+\partial_2^2u+\partial_3^2u=0$的条件可进一步化简为
                $$\iiint_\Omega\big((\partial_1u)^2+(\partial_2u)^2+(\partial_3u)^3\big)\dr V$$
                这就是要证的结论。
            \end{itemize}
        }

        \item 若$u(x,y,z)$在$S$上恒为0，则$\Omega$上也恒为0。
        
        \sol{
            由于$S$上$u$为0，前一问中的
            $$\oiint_Su\frac{\partial u}{\partial\nv}\dr S=\oiint_S0\dr S=0$$
            由此可知
            $$\iiint_\Omega\big((\partial_1u)^2+(\partial_2u)^2+(\partial_3u)^3\big)\dr V$$
            由于$u$的偏导连续，且$(\partial_1u)^2+(\partial_2u)^2+(\partial_3u)^2\ge0$恒成立，利用本讲义19.1.2第1题的结论(这里仅针对区域$\Omega$内部，记为$\Omega^\circ$)有
            $$\forall(x,y,z)\in\Omega^\circ,\quad(\partial_1u(x,y,z))^2+(\partial_2u(x,y,z))^2+(\partial_3u(x,y,z))^2\ge0$$
            也即$\partial_1u(x,y,z)=\partial_2u(x,y,z)=\partial_3u(x,y,z)=0$。

            通过拉格朗日中值定理可证明$u$在$\Omega^\circ$中为常数(详细证明可见上学期讲义15.1第10题)，再由连续性可知$u$在$\Omega$上为常数，根据边界为0可得恒为0，得证。

            \note 与本讲义22.1.1第8题(2)相同，最后这段说明的严谨论证相对复杂，不过直观来看是自然的，一般写出核心说明步骤即可。
        }
    \end{enumerate}

    \item (习题8.6.14)用斯托克斯公式计算
    $$\oint_L(y^2-z^2+x^2)\dr x+(z^2-x^2+y^2)\dr y+(x^2-y^2+z^2)\dr z$$
    其中$L$为平面$x+y+z=\frac{3}{2}R$与立方体
    $$\{(x,y,z)\mid x\in[0,R],\ y\in[0,R],\ z\in[0,R]\}$$
    的交线，且从$x$轴正向看去沿逆时针方向，这里参数$R>0$。

    \sol{
        分为如下步骤：
        \begin{itemize}
            \item \textbf{斯托克斯公式}
            
            将被积函数看作$P\dr x+Q\dr y+R\dr z$，由初等函数知$P$、$Q$、$R$在$\rb^3$上可导且导函数连续，直接计算得
            $$R_y-Q_z=-2y-2z,\quad P_z-R_x=-2z-2x,\quad Q_x-P_y=-2x-2y$$
            为了应用斯托克斯公式，考虑在平面$x+y+z=\frac{3}{2}R$在立方体里的部分$S$，则原积分可化为(这里利用了第二型曲面积分的性质提出$-2$以化简)
            $$-2\iint_S(y+z)\dr y\dr z+(z+x)\dr z\dr x+(x+y)\dr x\dr y$$
            利用右手定则，$S$应选择指向$x$轴正向的法向作为正向。

            \item \textbf{化为第一型曲面积分}
            
            利用第一型曲面积分与第二型曲面积分的关系可知(这里$\vec{n}$代表第二型曲面积分对应定向的单位法向量)
            $$\vec{n}\dr S=(\dr y\dr z,\dr z\dr x,\dr x\dr y)$$
            利用平面的几何性质可知它的一个法向为$(1,1,1)$，指向$x$轴正向即代表$x$分量为正，从而对应分量的单位法向量为
            $$\frac{1}{\sqrt3}(1,1,1)$$
            由此即得原积分为
            $$-2\iint_S(y+z,z+x,x+y)\cdot(\dr y\dr z,\dr z\dr x,\dr x\dr y)=-2\iint_S(y+z,z+x,x+y)\cdot\frac{1}{\sqrt3}(1,1,1)\dr S$$
            也即
            $$-\frac{4}{\sqrt{3}}\iint_S(x+y+z)\dr S$$

            \item \textbf{积分计算}
            
            由于平面上的点满足$x+y+z=\frac{3}{2}R$，积分可以化为
            $$-2\sqrt3 R\iint_S\dr S$$

            利用第一型曲面积分的几何意义，$\iint_S\dr S$即为曲面面积，而通过空间几何观察与计算可以发现$S$是一个以$\frac{\sqrt2}{2}R$为边长的正六边形，从而利用等边三角形面积公式可得到其面积为
            $$6\cdot\frac{\sqrt3}{4}\bigg(\frac{\sqrt2}{2}R\bigg)^2=\frac{3\sqrt3}{4}R^2$$
            从而最终结果为
            $$-2\sqrt3 R\cdot\frac{3\sqrt3}{4}R^2=-\frac{9}{2}R^3$$
        \end{itemize}
    }

    \item (习题8.6.16)用斯托克斯公式计算
    $$\oint_Lxz\dr x+xy\dr y+3xz\dr z$$
    其中$L$为平面$2x+y+z=2$在第一卦限部分的边界，且从$x$轴正向看去沿逆时针方向。

    \sol{
        分为如下步骤：
        \begin{itemize}
            \item \textbf{斯托克斯公式}
            
            将被积函数看作$P\dr x+Q\dr y+R\dr z$，由初等函数知$P$、$Q$、$R$在$\rb^3$上可导且导函数连续，直接计算得
            $$R_y-Q_z=0,\quad P_z-R_x=x-3z,\quad Q_x-P_y=y$$
            为了应用斯托克斯公式，考虑在平面$2x+y+z=2$在第一卦限的部分$S$，则原积分可化为
            $$\iint_S(x-3z)\dr z\dr x+y\dr x\dr y$$
            利用右手定则，$S$应选择指向$x$轴正向的法向作为正向。

            \item \textbf{参数化与定向}
            
            $z$可以看作$x$、$y$的函数，从而有自然的参数化
            $$\vec{r}=(x,y,z)=(u,v,2-2u-v)$$
            代入约束可以得到参数范围
            $$u\ge0,\quad v\ge0,\quad 2u+v\le 2$$
            将上述条件确定的$(u,v)$区域记为$D$。

            为了确定定向，用下标$u$、$v$表示对$u$、$v$求偏导并保持另一个不变，直接计算有
            $$\rv_u=(x_u,y_u,z_u)=(1,0,-2),\quad\rv_v=(x_v,y_v,z_v)=(0,1,-1)$$
            我们\textbf{任取}一个使得$\rv_u$、$\rv_v$\textbf{不平行}(也可表述为线性无关/叉乘非零)的$D$中点，如取$(u,v)=(0,0)$，此处有$\rv=(0,0,2)$、$\rv_u=(1,0,-2)$、$\rv_v=(0,1,-1)$，计算得$\rv_u\times\rv_v=(2,1,1)$，即指向$x$轴正方向，符合要求，无需添加负号。

            \item \textbf{化为重积分与计算}
            
            代入
            $$\dr y\dr z=\frac{D(y,z)}{D(u,v)}\dr u\dr v,\quad\dr z\dr x=\frac{D(z,x)}{D(u,v)}\dr u\dr v,\quad\dr x\dr y=\frac{D(x,y)}{D(u,v)}\dr u\dr v$$
            与$x$、$y$、$z$的参数方程并计算化简可得原积分为
            $$\iint_D(7u+4v-6)\dr u\dr v$$
            将区域描述为
            $$0\le u\le1,\quad 0\le v\le2-2u$$
            即可写出累次积分
            $$\int_0^1\dr u\int_0^{2-2u}(7u+4v-6)\dr v$$
            对$v$积分得到
            $$\int_0^1(7(2-2u)u-6(2-2u)+2(2-2u)^2)\dr u$$
            进一步计算可得到结果为
            $$\frac{7}{3}-6+\frac{8}{3}=-1$$
        \end{itemize}
    }

    \item (总练习题8.1)设一细金属线的形状可由参数方程
    $$\begin{cases}x=t\\y=\frac{2\sqrt2}{3}t^{3/2}\\z=\frac{t^2}{2}\end{cases}(0\le t\le2)$$
    表示，且在$t$对应点处密度为$\frac{1}{1+t}$，求质心坐标。

    \sol{
        设此金属线为$L$。利用质心坐标公式，设质心坐标为$(x_0,y_0,z_0)=\frac{1}{m}(x_1,y_1,z_1)$，应满足(这里$t$代表曲线在题干参数化下的参数$t$)
        $$x_1=\int_S\frac{x}{1+t}\dr s,\quad y_1=\int_S\frac{y}{1+t}\dr s,\quad z_1=\int_S\frac{z}{1+t}\dr s,\quad m=\int_S\frac{1}{1+t}\dr s$$
        计算分为如下步骤：
        \begin{itemize}
            \item \textbf{化为定积分}
            
            用下标$t$表示对$t$求导，直接计算有
            $$x_t=1,\quad y_t=\sqrt2\sqrt{t},\quad z_t=t$$
            从而(利用$t\in[0,2]$可去除绝对值)
            $$\dr s=\sqrt{x_t^2+y_t^2+z_t^2}\dr t=\sqrt{1+2t+t^2}\dr t=(t+1)\dr t$$
            由此
            $$x_1=\int_0^2t\dr t,\quad y_1=\int_0^2\frac{2\sqrt2}{3}t^{3/2}\dr t,\quad z_1=\int_0^2\frac{t^2}{2}\dr t,\quad m=\int_0^21\dr t$$
            
            \item \textbf{计算与综合结果}
            
            直接计算即可得到积分结果
            $$x_1=2,\quad y_1=\frac{32}{15},\quad z_1=\frac{4}{3},\quad m=2$$
            从而质心坐标为
            $$\bigg(1,\frac{16}{15},\frac{2}{3}\bigg)$$
        \end{itemize}
    }

    \item (总练习题8.8)利用
    $$x^2+xy+y^2=\frac{1}{2}(x^2+y^2)+\frac{1}{2}(x+y)^2$$
    估计积分
    $$I_a=\oint_{L_a^+}\frac{y\dr x-x\dr y}{(x^2+xy+y^2)^2}$$
    的范围，其中$L_a$为圆周$x^2+y^2=a^2$，这里参数$a>0$。

    \note 本题题目表述并不清晰，大家阅读答案理解估算方式即可。

    \sol{
        \note 注意第二型曲线/曲面积分的被积函数大小关系\textbf{无法推出积分大小关系}，必须\textbf{化为第一型或定积分/重积分}后估算。

        利用第一型曲线积分与第二型曲线积分的关系可知
        $$\vec\tau\dr s=(\dr x,\dr y)$$
        这里$\vec\tau=(\tau_1,\tau_2)$是曲线逆时针方向的单位切向量。对于圆$x^2+y^2=a^2$，利用几何知它在每点$(x,y)$处的一个外法向量是$(x,y)$，而逆时针方向切向量为其向左旋转$\frac{\pi}{2}$，由几何关系知为$(-y,x)$，从而单位切向量是
        $$\frac{1}{\sqrt{x^2+y^2}}(-y,x)$$
        代入可发现
        $$\begin{aligned}I_a=&\oint_{L_a^+}\frac{1}{(x^2+xy+y^2)^2}(-y,x)\cdot(\dr x,\dr y)\\=&\oint_{L_a}\frac{1}{(x^2+xy+y^2)^2}(-y,x)\cdot\frac{1}{\sqrt{x^2+y^2}}(-y,x)\dr s\\=&\oint_{L_a}\frac{\sqrt{x^2+y^2}}{(x^2+xy+y^2)^2}\dr s\end{aligned}$$
        此时被积函数非负，从而可知$I(a)\ge0$。计算内积并利用曲线上$x^2+y^2=a^2$可化简得到$a>0$时积分为
        $$I(a)=a\oint_{L_a}\frac{1}{(x^2+xy+y^2)^2}\dr s$$
        接下来需要估算，利用基本不等式可知
        $$\frac{1}{2}(x^2+y^2)\le x^2+xy+y^2\le\frac{3}{2}(x^2+y^2)$$
        也即$L_a$上有
        $$\frac{1}{2}a^2\le x^2+xy+y^2\le\frac{3}{2}a^2$$
        于是可将被积函数放缩得到
        $$a\oint_{L_a}\bigg(\frac{3a^2}{2}\bigg)^{-2}\dr s\le I(a)\le a\oint_{L_a}\bigg(\frac{a^2}{2}\bigg)^{-2}\dr s$$
        根据第一型曲线积分的几何意义，$\oint_{L_a}\dr s$为圆$x^2+y^2=a^2$的周长$2\pi a$，最终得到
        $$\frac{8\pi}{9a^2}\le I(a)\le\frac{8\pi}{a^2}$$
    }

    \item (总练习题8.13)设$L$是平面$2x+2y+z=2$上的一条光滑的简单闭曲线，且从$x$轴正向看去沿逆时针方向。证明曲线积分
    $$\oint_L2y\dr x+3z\dr y-x\dr z$$
    只与$L$所围区城的面积有关，而与$L$的形状及位置无关。

    \sol{
        由于需要将曲线的积分转化为曲面的积分，需要应用\textbf{斯托克斯公式}。分为如下步骤：
        \begin{itemize}
            \item \textbf{斯托克斯公式}
            
            将被积函数看作$P\dr x+Q\dr y+R\dr z$，由初等函数知$P$、$Q$、$R$在$\rb^3$上可导且导函数连续，直接计算得
            $$R_y-Q_z=-3,\quad P_z-R_x=1,\quad Q_x-P_y=-2$$
            为了应用斯托克斯公式，考虑$L$在平面上围成的部分$S$，则原积分可化为
            $$\iint_S-3\dr y\dr z+\dr z\dr x-2\dr x\dr y$$
            利用右手定则，$S$应选择指向$x$轴正向的法向作为正向。

            \item \textbf{化为第一型曲面积分}
            
            由于曲面的\textbf{面积}与第一型曲面积分相关，需要进一步化为第一型曲面积分。

            利用第一型曲面积分与第二型曲面积分的关系可知(这里$\vec{n}$代表第二型曲面积分对应定向的单位法向量)
            $$\vec{n}\dr S=(\dr y\dr z,\dr z\dr x,\dr x\dr y)$$
            利用平面的几何性质可知它的一个法向为$(2,2,1)$，指向$x$轴正向即代表$x$分量为正，从而对应分量的单位法向量为
            $$\frac{1}{3}(2,2,1)$$
            由此即得原积分为
            $$\iint_S(-3,1,-2)\cdot(\dr y\dr z,\dr z\dr x,\dr x\dr y)=\iint_S(-3,1,-2)\cdot\frac{1}{3}(2,2,1)\dr S=-2\iint_S\dr S$$
            也即最终的积分结果为面积的$-2$倍，这就得到了证明。
        \end{itemize}

        \note 类似可发现在\textbf{平面上的曲线}对\textbf{一次函数}的第二型曲线积分只与围成的面积有关。
    }

    \item (总练习题8.14)考虑向量值函数
    $$\Fv(x,y,z)=\bigg(-\frac{y}{x^2+y^2},\frac{x}{x^2+y^2},z\bigg)$$
    \begin{enumerate}[(1)]
        \item 证明$\rot\Fv=\vec{0}$处处成立。
        
        \sol{
            设$\Fv=(P,Q,R)$，直接计算可发现
            $$R_y=Q_z=0,\quad P_z=R_x=0,\quad Q_x=P_y=\frac{1}{x^2+y^2}-\frac{2x^2}{(x^2+y^2)^2}$$
            从而由定义有$\rot\Fv=\vec{0}$。
        }

        \item 当$L$为平面上的圆周$x^2+y^2=1$时证明
        $$\oint_{L^+}\Fv(x,y,z)\cdot\dr\rv\ne0$$
        这与斯托克斯公式结果矛盾吗？

        \sol{
            分为如下步骤：
            \begin{itemize}
                \item \textbf{参数化与定向}
                
                利用极坐标可以得到圆的参数方程
                $$\begin{cases}x=\cos\theta\\y=\sin\theta\end{cases}(0\le\theta\le2\pi)$$
                为了确定定向，用下标$\theta$代表对$\theta$求导，直接计算有
                $$x_\theta=-\sin\theta,\quad y_\theta=\cos\theta$$
                我们\textbf{任取}一个使得$x_\theta$、$y_\theta$\textbf{不全为0}的点，如取$\theta=\pi$，此处有$(x,y)=(-1,0)$、$(x_\theta,y_\theta)=(0,-1)$，在曲线上的点$(x,y)$处画出向量$(x_\theta,y_\theta)$，可以发现它指向逆时针方向，与要求的定向相同，从而化为定积分无需加负号。

                \item \textbf{化为定积分与计算}
                
                代入$\dr x=x_\theta\dr\theta$、$\dr y=y_\theta\dr\theta$=与$x$、$y$的参数方程可得原积分为
                $$\int_0^{2\pi}\dr\theta=2\pi$$

                \item \textbf{分析矛盾}
                
                与斯托克斯公式\textbf{并不矛盾}。这是由于任取一个以$L$为边界的曲面，都将经过某个满足$x=y=0$的点(这在几何上是直观的，详细证明相对复杂)，而这点处$\Fv$没有定义，因此无法直接使用斯托克斯公式转化为曲面积分以得到积分为0。
            \end{itemize}
        }
    \end{enumerate}

    \item (总练习题8.17)给定点$(x_0,y_0,z_0)\in\rb^3$，考虑函数
    $$u(x,y,z)=\ln\frac{1}{\sqrt{(x-x_0)^2+(y-y_0)^2+(z-z_0)^2}}$$
    求所有使得$u$的梯度模长为1的点。

    \sol{
        将$u(x,y,z)$的表达式化简为
        $$-\frac{1}{2}\ln\big((x-x_0)^2+(y-y_0)^2+(z-z_0)^2\big)$$
        直接计算可知
        $$\frac{\partial u(x,y,z)}{\partial x}=-\frac{x-x_0}{(x-x_0)^2+(y-y_0)^2+(z-z_0)^2}$$
        $$\frac{\partial u(x,y,z)}{\partial y}=-\frac{y-y_0}{(x-x_0)^2+(y-y_0)^2+(z-z_0)^2}$$
        $$\frac{\partial u(x,y,z)}{\partial z}=-\frac{z-z_0}{(x-x_0)^2+(y-y_0)^2+(z-z_0)^2}$$
        由此直接计算可知梯度模长为
        $$\sqrt{\bigg(\frac{\partial u(x,y,z)}{\partial x}\bigg)^2+\bigg(\frac{\partial u(x,y,z)}{\partial y}\bigg)^2+\bigg(\frac{\partial u(x,y,z)}{\partial z}\bigg)^2}=\frac{1}{\sqrt{(x-x_0)^2+(y-y_0)^2+(z-z_0)^2}}$$
        由此其为1当且仅当
        $$(x-x_0)^2+(y-y_0)^2+(z-z_0)^2=1$$
    }
\end{enumerate}

\subsubsection{补充题解答}
\begin{enumerate}
    \item (习题8.3.2(2))计算心脏线
    $$\begin{cases}x=a(1-\cos t)\cos t\\y=a(1-\cos t)\sin t\end{cases}(0\le t\le 2\pi)$$
    围成的面积，这里参数$a>0$。

    \sol{
        分为如下步骤：
        \begin{itemize}
            \item \textbf{格林公式}
            
            设围成的区域为$D$，则根据重积分几何意义知面积为
            $$\iint_D\dr x\dr y$$
            利用格林公式，构造$P=0$、$Q=x$，则$Q_x-P_y=1$，设区域边界为$L$，面积应可以写为
            $$\oint_{L^+}x\dr y$$

            \item \textbf{定向与化为定积分}
            
            为了确定定向，用下标$t$代表对$t$求导，直接计算有
            $$x_t=-a(1-\cos t)\sin t+a\sin t\cos t,\quad y_t=a(1-\cos t)\cos t+a\sin^2t$$
            我们\textbf{任取}一个使得$x_t$、$y_t$\textbf{不全为0}的点，如取$t=\pi$，此处有$(x,y)=(-2a,0)$、$(x_t,y_t)=(0,-2a)$，在曲线上的点$(x,y)$处画出向量$(x_t,y_t)$，可以发现它指向逆时针方向，与要求的定向相同，从而化为定积分无需加负号。

            \note 事实上本题也可以不用提前确定定向，因为\textbf{面积一定为正}，若算出负数再添加负号即可。

            代入$\dr x=x_t\dr t$、$\dr y=y_t\dr t$与$x$、$y$的参数方程可得原积分为
            $$a^2\int_0^{2\pi}\big((1-\cos t)^2\cos^2t+(1-\cos t)\cos t\sin^2t\big)\dr t$$
            
            \item \textbf{积分计算}
            
            将积分展开为
            $$a^2\int_0^{2\pi}(\cos^2t-2\cos^3t+\cos^4t+\cos t\sin^2t-\cos^2t\sin^2t)\dr t$$
            化简即得
            $$a^2\int_0^{2\pi}(\cos^4t+\sin^4t-2\cos^3t+\cos t\sin^2t)\dr t$$
            将$t$换元为$\frac{\pi}{2}-t$并利用周期函数在任何周期积分相同将积分区域移动回$[0,2\pi]$可发现
            $$\int_0^{2\pi}\cos^4t\dr t=\int_0^{2\pi}\sin^4t\dr t$$
            将$t$换元为$\pi-t$并利用周期函数在任何周期积分相同将积分区域移动回$[0,2\pi]$可发现
            $$\int_0^{2\pi}\cos^3t\dr t=-\int_0^{2\pi}\cos^3t\dr t,\quad\int_0^{2\pi}\cos t\sin^2t\dr t=-\int_0^{2\pi}\cos t\sin^2t\dr t$$
            由此上方两积分均为0，所求的积分最终化为
            $$2a^2\int_0^{2\pi}\sin^4t\dr t$$
            类似\exref{sin 2 and 4 to 2pi}的对称分析即得积分结果为
            $$\frac{3}{2}\pi a^2$$
        \end{itemize}
    }

    \item (习题8.3.4)证明下列积分与路径无关并求积分：
    \begin{enumerate}
        \item[(2)]$$\int_{(a_1,b_1)}^{(a_2,b_2)}xy(1+y)\dr x+x^2\bigg(\frac{1}{2}+y\bigg)\dr y$$
        
        \sol{
            分为如下步骤：
            \begin{itemize}
                \item \textbf{格林公式}
                
                将被积函数看作$P\dr x+Q\dr y$，由初等函数知$P$、$Q$在$\rb^2$上可导且导函数连续，进一步计算验证可知
                $$Q_x-P_y=x+2xy-x-xy-xy=0$$
                由$\rb^2$单连通得积分与路径无关。

                此时已经可以直接考虑两点的线段进行参数化计算，但由于两点均未给定，计算过程相对复杂，我们这里介绍\textbf{配凑全微分}的做法。

                \item \textbf{配凑全微分}
                
                设关于$x$、$y$的表达式$u$满足$P=u_x$、$Q=u_y$，将$P$对$x$进行不定积分可得存在函数$\varphi$使得
                $$u=\frac{1}{2}x^2y(1+y)+\varphi(y)$$
                而又由于
                $$Q=u_y=\frac{1}{2}x^2(2y+1)+\varphi'(y)$$
                代入可发现$\varphi'(y)=0$，从而
                $$u=\frac{1}{2}x^2(2y+1)+C$$

                \item \textbf{结果计算}
                
                利用教材中的定理可得曲线积分结果为$u$在$(a_2,b_2)$处的值减去$u$在$(a_1,b_1)$处的值，即
                $$\frac{1}{2}a_2^2(2b_2+1)-\frac{1}{2}a_1^2(2b_1+1)$$
            \end{itemize}
        }

        \item[(3)]$$\int_{(0,0)}^{(a,b)}\er^x\cos y\dr x-\er^x\sin y\dr y$$
        
        \sol{
            分为如下步骤：
            \begin{itemize}
                \item \textbf{格林公式}
                
                将被积函数看作$P\dr x+Q\dr y$，由初等函数知$P$、$Q$在$\rb^2$上可导且导函数连续，进一步计算验证可知
                $$Q_x-P_y=-\er^x\cos y+\er^x\cos y=0$$
                由$\rb^2$单连通得积分与路径无关。

                此时已经可以直接考虑两点的线段进行参数化计算，但由于$\er^t\cos t$相关的积分计算过程相对复杂，我们这里介绍\textbf{配凑全微分}的做法。

                \item \textbf{配凑全微分}
                
                设关于$x$、$y$的表达式$u$满足$P=u_x$、$Q=u_y$，将$P$对$x$进行不定积分可得存在函数$\varphi$使得
                $$u=\er^x\cos y+\varphi(y)$$
                而又由于
                $$Q=u_y=-\er^x\sin y+\varphi'(y)$$
                代入可发现$\varphi'(y)=0$，从而
                $$u=\er^x\cos y+C$$

                \item \textbf{结果计算}
                
                利用教材中的定理可得曲线积分结果为$u$在$(a,b)$处的值减去$u$在$(0,0)$处的值，即
                $$\er^a\cos b-1$$
            \end{itemize}
        }
    \end{enumerate}

    \item (习题8.3.6)求满足下列等式的函数$u(x,y)$：
    \begin{enumerate}[(1)]
        \item $\dr u=(x^2+2xy-y^2)\dr x+(x^2-2xy-y^2)\dr y$
        
        \sol{
            由于$u$满足$P=\partial_1u(x,y)$、$Q=\partial_2u(x,y)$，将$P$对$x$进行不定积分可得存在函数$\varphi$使得
            $$u(x,y)=\frac{1}{3}x^3+x^2y-y^2x+\varphi(y)$$
            而又由于
            $$Q=\partial_2u(x,y)=x^2-2yx+\varphi'(y)$$
            代入可发现$\varphi'(y)=-y^2$，从而
            $$u(x,y)=\frac{1}{3}x^3+x^2y-y^2x-\frac{1}{3}y^3+C$$
        }

        \item $\dr u=(2x\cos y-y^2\sin x)dx+(2y\cos x-x^2\sin y)\dr y$
        
        \sol{
            由于$u$满足$P=\partial_1u(x,y)$、$Q=\partial_2u(x,y)$，将$P$对$x$进行不定积分可得存在函数$\varphi$使得
            $$u(x,y)=x^2\cos y+y^2\cos x+\varphi(y)$$
            而又由于
            $$Q=\partial_2u(x,y)=-x^2\sin y+2y\cos x+\varphi'(y)$$
            代入可发现$\varphi'(y)=0$，从而
            $$u(x,y)=x^2\cos y+y^2\cos x+C$$
        }
    \end{enumerate}

    \item (习题8.3.7)求实数$a$、$b$使得
    $$\frac{(y^2+2xy+ax^2)\dr x-(x^2+2xy+by^2)\dr y}{(x^2+y^2)^2}$$
    是某个函数$u(x,y)$的全微分，并求$u(x,y)$。

    \sol{
        \note 本题的直接计算方法需要较大的计算量，这里介绍一种观察结构实现的待定系数解法。

        由于对$x$、$y$求偏导均为分母$(x^2+y^2)^2$、分子齐次二次多项式，可以假设$u$分母为$\frac{1}{x^2+y^2}$，分子考虑次数发现应为一次，从而设其为
        $$u(x,y)=\frac{\alpha x+\beta y}{x^2+y^2}$$

        \note 这样设法本质上\textbf{难以保证严谨性}，但先如此假设可以很大程度方便计算，若假设失败再尝试硬算。

        直接计算可知
        $$\dr u=\frac{(\alpha(x^2+y^2)-2x(\alpha x+\beta y))\dr x+(\beta(x^2+y^2)-2y(\alpha x+\beta y))\dr y}{(x^2+y^2)^2}$$
        对比系数可得
        $$\begin{cases}-\alpha=a\\-2\beta=2\\\alpha=1\\\beta=-1\\-2\alpha=-2\\-\beta=-b\end{cases}$$
        由此可发现仅当$a=b=1$时有解，此时
        $$u(x,y)=\frac{x-y}{x^2+y^2}$$
        利用全微分的\textbf{唯一性}(详见本讲义22.2.2)可知所有满足条件的$u(x,y)$为
        $$\frac{x-y}{x^2+y^2}+C$$        
    }

    \item (总练习题8.2)计算
    $$\int_{(3,4)}^{(5,12)}\frac{x\dr x+y\dr y}{\sqrt{x^2+y^2}}$$
    其中积分路径不过坐标原点。

    \sol{
        分为如下步骤：
        \begin{itemize}
            \item \textbf{格林公式}
            
            本讲义22.1.1第5题中，我们证明了
            $$\bigg(\frac{x}{\sqrt{x^2+y^2}}+y\bigg)\dr x+\bigg(\frac{y}{\sqrt{x^2+y^2}}+x\bigg)\dr y$$
            的积分与路径无关，而这是本题的被积函数加上$y\dr x+x\dr y=\dr(xy)$，由此可用完全相同的方式从本题$Q_x-P_y=0$与绕原点一周的积分为0说明本题的积分与路径无关。

            \note 同样，在看出$P\dr x+Q\dr y=\dr\sqrt{x^2+y^2}$后也能证明结论，但相应定理在教材中并未证明，因此\textbf{不建议直接使用}。

            \item \textbf{配凑全微分}
            
            本题我们介绍\textbf{直接配凑全微分}的做法。

            由于分子有
            $$x\dr x=\frac{1}{2}x^2,\quad y\dr y=\frac{1}{2}y^2$$
            设$v=x^2+y^2$，则有
            $$\frac{x\dr x+y\dr y}{\sqrt{x^2+y^2}}=\frac{\dr v}{2\sqrt{v}}$$
            而直接不定积分即可得到取
            $$u=\sqrt{v}+C=\sqrt{x^2+y^2}+C$$
            时有
            $$P=u_x,\quad Q=u_y$$

            \item \textbf{结果计算}
                
            利用教材中的定理可得曲线积分结果为$u$在$(5,12)$处的值减去$u$在$(3,4)$处的值，即
            $$13-5=8$$
        \end{itemize}
    }

    \item (总练习题8.5)计算
    $$\int_{(1,2\pi)}^{(2,\pi)}\bigg(1-\frac{y^2}{x^2}\cos\frac{y}{x}\bigg)\dr x+\bigg(\sin\frac{y}{x}+\frac{y}{x}\cos\frac{y}{x}\bigg)\dr y$$
    其中积分路径不过$y$轴。

    \sol{
        分为如下步骤：
        \begin{itemize}
            \item \textbf{格林公式}
            
            将被积函数看作$P\dr x+Q\dr y$，由初等函数知$P$、$Q$在$x>0$上可导且导函数连续，进一步计算验证可知
            $$Q_x-P_y=-\frac{y}{x^2}\bigg(2\cos\frac{y}{x}-\frac{y}{x}\sin\frac{y}{x}\bigg)+\frac{1}{x}\bigg(\frac{2y}{x}\cos\frac{y}{x}-\frac{y^2}{x^2}\sin\frac{y}{x}\bigg)=0$$
            由$x>0$单连通得积分与路径无关。

            \item \textbf{配凑全微分}
            
            本题我们仍然介绍\textbf{直接配凑全微分}的做法。

            由于被积函数均与$\frac{y}{x}$有关，直接计算有
            $$\dr\frac{y}{x}=-\frac{y}{x^2}\dr x+\frac{1}{x}\dr y$$
            由此设$v=\frac{y}{x}$，有
            $$\dr v=\frac{-v\dr x+\dr y}{x}$$
            也即
            $$\dr y=x\dr v+v\dr x$$
            而$P\dr x+Q\dr y$为
            $$(1-v^2\cos v)\dr x+(\sin v+v\cos v)\dr y$$
            代入$\dr y$的表达式并化简得其为
            $$(1+v\sin v)\dr x+(x\sin v+xv\cos v)\dr v$$
            此时根据形式已经可以看出
            $$v\sin v=\frac{\partial(xv\sin v)}{\partial x},\quad x\sin v+xv\cos v=\frac{\partial(xv\sin v)}{\partial v}$$
            由此取
            $$u=x+xv\sin v+C=x+y\sin\frac{y}{x}+C$$
            时即有
            $$P=u_x,\quad Q=u_y$$

            \note 化为$x$、$v$相关的全微分后也可先对$x$不定积分以确定$u$。

            \item \textbf{结果计算}
                
            利用教材中的定理可得曲线积分结果为$u$在$(2,\pi)$处的值减去$u$在$(1,2\pi)$处的值，即
            $$(2+\pi)-1=\pi+1$$
        \end{itemize}
    }

    \item (习题8.5.1)计算
    $$\iint_{S^+}yz\dr y\dr z+xz\dr z\dr x+xy\dr x\dr y$$
    其中$S$为平面$x=0$、$y=0$、$z=0$、$x+y+z=a$围成四面体的表面，这里参数$a>0$。

    \sol{
        分为如下步骤：
        \begin{itemize}
            \item \textbf{区域描述}
            
            我们先分析区域。根据``围成''的有界性要求与围成的区域必然以每个平面为边界，我们可以得到四个约束(类似对平面三角形区域的分析)
            $$x\ge0,\quad y\le0,\quad z\ge0,\quad x+y+z\le a$$
            由于积分区域为表面，需要将每个约束取等，其他取不等号。由此我们将曲面积分分为了如下四个部分的积分之和：
            $$S_1:\quad x=0,\quad y\ge0,\quad z\ge0,\quad x+y+z\le1$$
            $$S_2:\quad y=0,\quad x\ge0,\quad z\ge0,\quad x+y+z\le1$$
            $$S_3:\quad z=0,\quad x\ge0,\quad y\ge0,\quad x+y+z\le1$$
            $$S_4:\quad x+y+z=1,\quad x\ge0,\quad y\ge0,\quad z\ge0$$
            下面分别计算每个部分的积分，再求和。
            
            \item \textbf{参数化与定向1}
            
            $x$可以看作$y$、$z$的函数，从而有自然的参数化
            $$\vec{r}=(x,y,z)=(0,u,v)$$
            代入约束可以得到参数范围
            $$u\ge0,\quad v\ge0,\quad u+v\le1$$
            将上述条件确定的$(u,v)$区域记为$D$。

            为了确定定向，用下标$u$、$v$表示对$u$、$v$求偏导并保持另一个不变，直接计算有
            $$\rv_u=(x_u,y_u,z_u)=(0,1,0),\quad\rv_v=(x_v,y_v,z_v)=(0,0,1)$$
            我们\textbf{任取}一个使得$\rv_u$、$\rv_v$\textbf{不平行}(也可表述为线性无关/叉乘非零)的$D$中点，如取$(u,v)=(0,0)$，此处有$\rv=(0,0,0)$、$\rv_u=(0,1,0)$、$\rv_v=(0,0,1)$，计算得$\rv_u\times\rv_v=(1,0,0)$，即指向$x$轴正方向，而由于$x$的约束有$x\ge0$，此处为$x$最小处，外侧应为$x$轴负方向，两者相反，因此\textbf{需要添加负号}。
            
            \item \textbf{化为重积分与计算1}
            
            代入
            $$\dr y\dr z=\frac{D(y,z)}{D(u,v)}\dr u\dr v,\quad\dr z\dr x=\frac{D(z,x)}{D(u,v)}\dr u\dr v,\quad\dr x\dr y=\frac{D(x,y)}{D(u,v)}\dr u\dr v$$
            与$x$、$y$、$z$的参数方程可得原积分为(最左侧的负号来自定向)
            $$-\iint_Duv\dr u\dr v$$
            将区域描述为
            $$0\le u\le1,\quad 0\le v\le1-u$$
            即可写出累次积分并对$v$积分得到
            $$-\int_0^1\dr u\int_0^{1-u}uv\dr v=-\frac{1}{2}\int_0^1(1-u)^2u\dr v$$
            直接展开计算或换元$t=1-u$即可算出积分结果
            $$-\frac{1}{24}$$

            \item \textbf{参数化与定向2}
            
            $y$可以看作$x$、$z$的函数，从而有自然的参数化
            $$\vec{r}=(x,y,z)=(u,0,v)$$
            代入约束可以得到参数范围
            $$u\ge0,\quad v\ge0,\quad u+v\le1$$
            将上述条件确定的$(u,v)$区域记为$D$。

            为了确定定向，用下标$u$、$v$表示对$u$、$v$求偏导并保持另一个不变，直接计算有
            $$\rv_u=(x_u,y_u,z_u)=(1,0,0),\quad\rv_v=(x_v,y_v,z_v)=(0,0,1)$$
            我们\textbf{任取}一个使得$\rv_u$、$\rv_v$\textbf{不平行}(也可表述为线性无关/叉乘非零)的$D$中点，如取$(u,v)=(0,0)$，此处有$\rv=(0,0,0)$、$\rv_u=(1,0,0)$、$\rv_v=(0,0,1)$，计算得$\rv_u\times\rv_v=(0,-1,0)$，即指向$y$轴负方向，而由于$y$的约束有$y\ge0$，此处为$x$最小处，外侧应为$y$轴负方向，两者相同，无需添加负号。
            
            \item \textbf{化为重积分与计算2}
            
            代入
            $$\dr y\dr z=\frac{D(y,z)}{D(u,v)}\dr u\dr v,\quad\dr z\dr x=\frac{D(z,x)}{D(u,v)}\dr u\dr v,\quad\dr x\dr y=\frac{D(x,y)}{D(u,v)}\dr u\dr v$$
            与$x$、$y$、$z$的参数方程可得原积分为
            $$\iint_D(-uv)\dr u\dr v$$
            可以发现，此积分与分析$S_1$得到的积分形式相同，区域描述也完全相同，由此结果也为$-\frac{1}{24}$。

            \item \textbf{参数化与定向3}
            
            $z$可以看作$x$、$y$的函数，从而有自然的参数化
            $$\vec{r}=(x,y,z)=(0,u,v)$$
            代入约束可以得到参数范围
            $$u\ge0,\quad v\ge0,\quad u+v\le1$$
            将上述条件确定的$(u,v)$区域记为$D$。

            为了确定定向，用下标$u$、$v$表示对$u$、$v$求偏导并保持另一个不变，直接计算有
            $$\rv_u=(x_u,y_u,z_u)=(0,1,0),\quad\rv_v=(x_v,y_v,z_v)=(0,0,1)$$
            我们\textbf{任取}一个使得$\rv_u$、$\rv_v$\textbf{不平行}(也可表述为线性无关/叉乘非零)的$D$中点，如取$(u,v)=(0,0)$，此处有$\rv=(0,0,0)$、$\rv_u=(0,1,0)$、$\rv_v=(0,0,1)$，计算得$\rv_u\times\rv_v=(0,0,1)$，即指向$z$轴正方向，而由于$z$的约束有$z\ge0$，此处为$z$最小处，外侧应为$z$轴负方向，两者相反，因此\textbf{需要添加负号}。
            
            \item \textbf{化为重积分与计算3}
            
            代入
            $$\dr y\dr z=\frac{D(y,z)}{D(u,v)}\dr u\dr v,\quad\dr z\dr x=\frac{D(z,x)}{D(u,v)}\dr u\dr v,\quad\dr x\dr y=\frac{D(x,y)}{D(u,v)}\dr u\dr v$$
            与$x$、$y$、$z$的参数方程可得原积分为
            $$-\iint_Duv\dr u\dr v$$
            可以发现，此积分与分析$S_1$得到的积分形式相同，区域描述也完全相同，由此结果也为$-\frac{1}{24}$。

            \item \textbf{参数化与定向4}
            
            $z$可以看作$x$、$y$的函数，从而有自然的参数化
            $$\vec{r}=(x,y,z)=(u,v,1-u-v)$$
            代入约束可以得到参数范围
            $$u\ge0,\quad v\ge0,\quad u+v\le1$$
            将上述条件确定的$(u,v)$区域记为$D$。

            为了确定定向，用下标$u$、$v$表示对$u$、$v$求偏导并保持另一个不变，直接计算有
            $$\rv_u=(x_u,y_u,z_u)=(1,0,-1),\quad\rv_v=(x_v,y_v,z_v)=(0,1,-1)$$
            我们\textbf{任取}一个使得$\rv_u$、$\rv_v$\textbf{不平行}(也可表述为线性无关/叉乘非零)的$D$中点，如取$(u,v)=(0,0)$，此处有$\rv=(0,0,1)$、$\rv_u=(1,0,-1)$、$\rv_v=(0,1,-1)$，计算得$\rv_u\times\rv_v$的$z$分量为1，即指向$z$轴正方向，而而由于$z$的范围是$0\le z\le1-x-y$，此处为$z$最大处，外侧应为$z$轴正方向，两者相同，无需添加负号。

            \item \textbf{化为重积分与计算4}
            
            代入
            $$\dr y\dr z=\frac{D(y,z)}{D(u,v)}\dr u\dr v,\quad\dr z\dr x=\frac{D(z,x)}{D(u,v)}\dr u\dr v,\quad\dr x\dr y=\frac{D(x,y)}{D(u,v)}\dr u\dr v$$
            与$x$、$y$、$z$的参数方程可得原积分为
            $$\iint_D(v(1-u-v)+u(1-u-v)+uv)\dr u\dr v=\iint_D(u+v-u^2-v^2-uv)\dr u\dr v$$
            将区域描述为
            $$0\le u\le1,\quad 0\le v\le1-u$$
            即可写出累次积分
            $$\int_0^1\dr u\int_0^{1-u}(u+v-u^2-v^2-uv)\dr v$$
            对$v$积分得到
            $$\int_0^1\dr u\bigg((1-u)u+\frac{(1-u)^2}{2}-(1-u)u^2-\frac{(1-u)^3}{3}-\frac{u(1-u)^2}{2}\bigg)\dr u$$
            进一步计算可得到结果为(这里换元$t=1-u$可以方便部分项的计算)
            $$\frac{1}{6}+\frac{1}{6}-\frac{1}{12}-\frac{1}{12}-\frac{1}{24}=\frac{1}{8}$$

            \item \textbf{综合结果}
            
            将四个部分求和可知最终的积分结果为0。
        \end{itemize}
    }

    \item (习题8.5.3)计算
    $$\iint_Sx^2y^2z\dr x\dr y$$
    其中$S$为球面$x^2+y^2+z^2=R^2$在$z\le0$的部分的下侧，这里参数$R>0$。

    \sol{
        分为如下步骤：
        \begin{itemize}
            \item \textbf{参数化}
            
            利用球坐标可得到球面的一个参数化
            $$\rv=(x,y,z)=(R\cos u\sin v,R\sin u\sin v,R\cos v)$$
            根据球坐标的几何意义，$z\le0$对应$v\in[\frac{\pi}{2},\pi]$，结合自然范围得到参数范围
            $$0\le u\le2\pi,\quad\frac{\pi}{2}\le v\le\pi$$
            将上述条件确定的$(u,v)$区域记为$D$。

            \item \textbf{定向}
            
            用下标$u$、$v$表示对$u$、$v$求偏导并保持另一个不变，直接计算有
            $$\rv_u=(x_u,y_u,z_u)=(-R\sin u\sin v,R\cos u\sin v,0)$$
            $$\rv_v=(x_v,y_v,z_v)=(R\cos u\cos v,R\sin u\cos v,-R\sin v)$$
            我们\textbf{任取}一个使得$\rv_u$、$\rv_v$\textbf{不平行}(也可表述为线性无关/叉乘非零)的$D$中点，如取$(u,v)=(0,\frac{\pi}{2})$，此处有$\rv=(R,0,0)$、$\rv_u=(0,R,0)$、$\rv_v=(0,0,-R)$，计算得$\rv_u\times\rv_v=(-R^2,0,0)$，即指向$x$轴负方向。根据几何可发现$z\le0$部分的下侧即为球面外侧，而由于$x=R$处为$x$最大处，外侧应为$x$轴正方向，两者相反，因此\textbf{需要添加负号}。
            
            \item \textbf{化为重积分}
            
            代入
            $$\dr y\dr z=\frac{D(y,z)}{D(u,v)}\dr u\dr v,\quad\dr z\dr x=\frac{D(z,x)}{D(u,v)}\dr u\dr v,\quad\dr x\dr y=\frac{D(x,y)}{D(u,v)}\dr u\dr v$$
            与$x$、$y$、$z$的参数方程并计算化简可得原积分为(这里已经添加了定向带来的负号)
            $$\iint_D(R^5\cos^2u\sin^2v\sin^2u\sin^2v\cos v)R^2\sin v\cos v\dr u\dr v$$
            进一步简化得到
            $$R^7\iint_D\cos^2u\sin^2u\sin^5v\cos^2v\dr u\dr v$$

            \item \textbf{积分计算}
            
            由区域描述可直接写出累次积分
            $$R^7\int_0^{2\pi}\dr u\int_{\pi/2}^{\pi}\cos^2u\sin^2u\sin^5v\cos^2v\dr v$$
            利用\exref{separation}可将它写成
            $$R^7\int_0^{2\pi}\cos^2u\sin^2u\dr u\int_{\pi/2}^\pi\sin^5v\cos^2v\dr v$$
            与\exref{sin 2 and 4 to 2pi}相同可计算得到
            $$\int_0^{2\pi}\cos^2u\sin^2u\dr u=\frac{\pi}{4}$$
            而换元$t=\pi-v$可得
            $$\int_{\pi/2}^\pi\sin^5v\cos^2v\dr v=\int_0^{\pi/2}\sin^5t\cos^2t\dr t$$
            被积函数可以进一步改写为$\sin^5t-\sin^7t$，从而利用上册教材3.4节例2知积分结果为
            $$\frac{2}{3}\cdot\frac{4}{5}-\frac{2}{3}\cdot\frac{4}{5}\cdot\frac{6}{7}=\frac{8}{105}$$
            综合得到最终结果
            $$R^7\cdot\frac{\pi}{4}\cdot\frac{8}{105}=\frac{2\pi}{105}$$

            \note 由此可见，熟悉$\sin^nt$积分结论可以大幅方便一些计算。
        \end{itemize}
    }

    \item (习题8.5.5)计算
    $$\iint_{S^+}xy^2\dr y\dr z+yz^2\dr z\dr x+zx^2\dr x\dr y$$
    其中$S$为椭球面$\frac{x^2}{a^2}+\frac{y^2}{b^2}+\frac{z^2}{c^2}=1$，这里参数$a>0$、$b>0$、$c>0$。

    \sol{
        分为如下步骤：
        \begin{itemize}
            \item \textbf{参数化}
            
            利用\textbf{广义球坐标}可得到椭球面的一个参数化
            $$\rv=(x,y,z)=(a\cos u\sin v,b\sin u\sin v,c\cos v)$$
            广义球坐标自然范围即为参数范围
            $$0\le u\le2\pi,\quad0\le v\le\pi$$
            将上述条件确定的$(u,v)$区域记为$D$。

            \item \textbf{定向}
            
            用下标$u$、$v$表示对$u$、$v$求偏导并保持另一个不变，直接计算有
            $$\rv_u=(x_u,y_u,z_u)=(-a\sin u\sin v,b\cos u\sin v,0)$$
            $$\rv_v=(x_v,y_v,z_v)=(a\cos u\cos v,b\sin u\cos v,-c\sin v)$$
            我们\textbf{任取}一个使得$\rv_u$、$\rv_v$\textbf{不平行}(也可表述为线性无关/叉乘非零)的$D$中点，如取$(u,v)=(0,\frac{\pi}{2})$，此处有$\rv=(a,0,0)$、$\rv_u=(0,b,0)$、$\rv_v=(0,0,-c)$，计算得$\rv_u\times\rv_v=(-bc,0,0)$，即指向$x$轴负方向。由于椭球面$x=a$处为$x$最大处，外侧应为$x$轴正方向，两者相反，因此\textbf{需要添加负号}。
            
            \item \textbf{化为重积分}
            
            代入
            $$\dr y\dr z=\frac{D(y,z)}{D(u,v)}\dr u\dr v,\quad\dr z\dr x=\frac{D(z,x)}{D(u,v)}\dr u\dr v,\quad\dr x\dr y=\frac{D(x,y)}{D(u,v)}\dr u\dr v$$
            与$x$、$y$、$z$的参数方程并计算化简可得原积分为(这里已经添加了定向带来的负号)
            $$abc\iint_D(b^2\cos^2u\sin^2u\sin^5v+c^2\sin^2u\sin^3v\cos^2v+a^2\cos^2u\cos^2v\sin^3v)\dr u\dr v$$

            \item \textbf{积分计算}
            
            由区域描述可写出累次积分，与上题类似通过\exref{separation}将其化为
            $$\begin{aligned}&ab^3c\int_0^{2\pi}\cos^2u\sin^2u\dr u\int_0^\pi\sin^5v\dr v\\+&abc^3\int_0^{2\pi}\sin^2u\dr u\int_0^\pi\sin^3v\cos^2v\dr v\\+&a^3bc\int_0^{2\pi}\cos^2u\dr u\int_0^\pi\sin^3v\cos^2v\dr v\end{aligned}$$
            
            与\exref{sin 2 and 4 to 2pi}类似利用对称性与上册教材3.4节例2可计算得到
            $$\int_0^{2\pi}\cos^2u\sin^2u\dr u=\frac{\pi}{4}$$
            $$\int_0^{2\pi}\sin^2u\dr u=\int_0^{2\pi}\cos^2u\dr u=\pi$$
            $$\int_0^\pi\sin^5v\dr v=2\int_0^{\pi/2}\sin^5v\dr v=\frac{16}{15}$$
            与上题类似并利用对称性可得
            $$\int_0^\pi\sin^3v\cos^2v\dr v=2\int_0^{\pi/2}(\sin^3v-\sin^5v)\dr v=\frac{4}{15}$$
            综合算出结果为
            $$abc\bigg(b^2\frac{\pi}{4}\cdot\frac{16}{15}+c^2\frac{4\pi}{15}+a^2\frac{4\pi}{15}\bigg)=\frac{4\pi}{15}(a^2+b^2+c^2)abc$$
        \end{itemize}
    }

    \item (习题8.5.8)计算
    $$\iint_{S^+}x\dr y\dr z+y\dr z\dr x+z\dr x\dr y$$
    其中$S$为椭球面$\frac{x^2}{a^2}+\frac{y^2}{b^2}+\frac{z^2}{c^2}=1$，这里参数$a>0$、$b>0$、$c>0$。

    \sol{
        分为如下步骤：
        \begin{itemize}
            \item \textbf{参数化}
            
            利用\textbf{广义球坐标}可得到椭球面的一个参数化
            $$\rv=(x,y,z)=(a\cos u\sin v,b\sin u\sin v,c\cos v)$$
            广义球坐标自然范围即为参数范围
            $$0\le u\le2\pi,\quad0\le v\le\pi$$
            将上述条件确定的$(u,v)$区域记为$D$。

            \item \textbf{定向}
            
            用下标$u$、$v$表示对$u$、$v$求偏导并保持另一个不变，直接计算有
            $$\rv_u=(x_u,y_u,z_u)=(-a\sin u\sin v,b\cos u\sin v,0)$$
            $$\rv_v=(x_v,y_v,z_v)=(a\cos u\cos v,b\sin u\cos v,-c\sin v)$$
            我们\textbf{任取}一个使得$\rv_u$、$\rv_v$\textbf{不平行}(也可表述为线性无关/叉乘非零)的$D$中点，如取$(u,v)=(0,\frac{\pi}{2})$，此处有$\rv=(a,0,0)$、$\rv_u=(0,b,0)$、$\rv_v=(0,0,-c)$，计算得$\rv_u\times\rv_v=(-bc,0,0)$，即指向$x$轴负方向。由于椭球面$x=a$处为$x$最大处，外侧应为$x$轴正方向，两者相反，因此\textbf{需要添加负号}。
            
            \item \textbf{化为重积分}
            
            代入
            $$\dr y\dr z=\frac{D(y,z)}{D(u,v)}\dr u\dr v,\quad\dr z\dr x=\frac{D(z,x)}{D(u,v)}\dr u\dr v,\quad\dr x\dr y=\frac{D(x,y)}{D(u,v)}\dr u\dr v$$
            与$x$、$y$、$z$的参数方程并计算化简可得原积分为(这里已经添加了定向带来的负号)
            $$abc\iint_D(\cos^2u\sin^3v+\sin^2u\sin^3v+\sin v\cos^2v)\dr u\dr v$$

            \item \textbf{积分计算}
            
            由区域描述可写出累次积分，与上题类似通过\exref{separation}将其化为
            $$\begin{aligned}&abc\int_0^{2\pi}\cos^2u\dr u\int_0^\pi\sin^3v\dr v\\+&abc\int_0^{2\pi}\sin^2u\dr u\int_0^\pi\sin^3v\dr v\\+&abc\int_0^{2\pi}\dr u\int_0^\pi\sin^2v\cos v\dr v\end{aligned}$$
            并计算得
            $$\int_0^{2\pi}\sin^2u\dr u=\int_0^{2\pi}\cos^2u\dr u=\pi$$
            $$\int_0^{2\pi}\dr u=2\pi$$
            $$\int_0^\pi\sin^3v\dr v=2\int_0^{\pi/2}\sin^3v\dr v=\frac{4}{3}$$
            $$\int_0^\pi\sin v\cos^2v\dr v=2\int_0^{\pi/2}(\sin v-\sin^3v)\dr v=\frac{2}{3}$$
            综合算出结果为
            $$4\pi abc$$
        \end{itemize}
    }

    \item (习题8.5.9)计算
    $$\iint_S-2\dr y\dr z+2y\dr z\dr x+\er^z\sin(x+2y)\dr x\dr y$$
    其中$S$为曲面
    $$y=\er^x,\quad y\in[1,2],\quad z\in[0,2]$$
    取$x$分量为正的法向量作为正向。

    \sol{
        分为如下步骤：
        \begin{itemize}
            \item \textbf{参数化与定向}
            
            $y$可以看作$x$、$z$的函数，从而有自然的参数化
            $$\vec{r}=(x,y,z)=(u,\er^u,v)$$
            代入约束可以得到参数范围
            $$0\le u\le\ln2,\quad 0\le v\le 2$$
            将上述条件确定的$(u,v)$区域记为$D$。

            为了确定定向，用下标$u$、$v$表示对$u$、$v$求偏导并保持另一个不变，直接计算有
            $$\rv_u=(x_u,y_u,z_u)=(1,\er^u,0),\quad\rv_v=(x_v,y_v,z_v)=(0,0,1)$$
            我们\textbf{任取}一个使得$\rv_u$、$\rv_v$\textbf{不平行}(也可表述为线性无关/叉乘非零)的$D$中点，如取$(u,v)=(0,0)$，此处有$\rv=(0,1,0)$、$\rv_u=(1,1,0)$、$\rv_v=(0,0,1)$，计算得$\rv_u\times\rv_v=(1,-1,0)$。由于它的$x$分量为正，与要求定向相同，无需添加负号。

            \item \textbf{化为重积分}
            
            代入
            $$\dr y\dr z=\frac{D(y,z)}{D(u,v)}\dr u\dr v,\quad\dr z\dr x=\frac{D(z,x)}{D(u,v)}\dr u\dr v,\quad\dr x\dr y=\frac{D(x,y)}{D(u,v)}\dr u\dr v$$
            与$x$、$y$、$z$的参数方程可得原积分为
            $$\iint_D(-2\er^u-2\er^u)\dr u\dr v=-4\iint_D\er^u\dr u\dr v$$

            \item \textbf{积分计算}
            
            由区域描述可写出累次积分
            $$-4\int_0^{\ln 2}\dr u\int_0^2\er^u\dr v$$
            对$v$积分可得结果为
            $$-8\int_0^{\ln 2}\er^u\dr u=-8$$
        \end{itemize}
    }
\end{enumerate}

\subsection{常微分方程与通解}
\subsubsection{常微分方程简介}
常微分方程的定义

解与区间上解的定义

初值问题

如何判断阶数与求导？

\subsubsection{通解}
$n$阶方程$n$个任意常数的直观理解

推导两个任意常数的独立性公式

多个任意常数的情况

\subsubsection{动力系统}
通解并非全部解

局部与整体($x^2y''=(y')^2$)

向量场的积分曲线

存在性与唯一性简述

\subsection{基本求解方式}
\subsubsection{换元与配凑}
一阶常系数线性微分方程

通过合适换元化为简单情况

几种换元方式与积分

\subsubsection{恰当方程与积分因子}
全微分、恰当方程定义

积分因子定义与表达式

只与$x$/$y$有关的积分因子

更一般的积分因子法

\newpage
\section{线性微分方程组}
\subsection{习题解答}
\subsubsection{作业解答}
\note 教材第九章的解答只进行\textbf{通解与特解的获取}，不研究这是否构成全部解等严谨性问题，一些进一步的分析可见本讲义第二十三章或25.2.5。注意理论来说求出的通解可以\textbf{完全不进行化简与合并}，实际几乎只取决于改卷标准，解答中我们尽量化到相对简洁的形式。

\begin{enumerate}
    \item (习题8.7.1)求下列向量场在指定点的散度：
    \begin{enumerate}[(1)]
        \item 在$(6,1,3)$处的
        $$\Fv=\bigg(5xy,\frac{xy}{z},3\ln(xz)\bigg)$$

        \sol{
            设$\Fv=(P,Q,R)$，根据散度定义有
            $$\nabla\cdot\Fv=P_x+Q_y+R_z=5y+\frac{x}{z}+\frac{3}{z}$$
            从而代入可得值为
            $$5+2+1=8$$
        }

        \item 在$(3,4,3)$处的
        $$\Fv=\sqrt{x^2+y^2}\vec{i}+\ln(y+\sqrt{y^2+z^2})\vec{j}+\sqrt{y^2+z^2}\vec{k}$$

        \sol{
            设$\Fv=(P,Q,R)$，根据散度定义有
            $$\nabla\cdot\Fv=P_x+Q_y+R_z=\frac{x}{\sqrt{x^2+y^2}}+\frac{1+\frac{y}{\sqrt{y^2+z^2}}}{y+\sqrt{y^2+z^2}}+\frac{z}{\sqrt{y^2+z^2}}$$
            从而代入可得值为
            $$\frac{3}{5}+\frac{1}{9}\cdot\frac{9}{5}+\frac{3}{5}=\frac{7}{5}$$
        }
    \end{enumerate}

    \item (习题8.7.2)设$\vec{a}$为常向量，$\rv=(x,y,z)$，$r=|\rv|$，$f(r)$为$r$的可微函数，求下列各式的值(以下的$\nabla$表示$\nabla_{x,y,z}$，即对$x,y,z$的梯度或散度)：
    \begin{enumerate}[(1)]
        \item $\nabla\cdot(r\vec{a})$
        
        \sol{
            设$\vec{a}=(a_1,a_2,a_3)$，根据散度定义知其为
            $$\frac{\partial a_1r}{\partial x}+\frac{\partial a_2r}{\partial y}+\frac{\partial a_3r}{\partial z}$$
            直接计算有
            $$\frac{\partial r}{\partial x}=\frac{x}{r},\quad \frac{\partial r}{\partial y}=\frac{y}{r},\quad \frac{\partial r}{\partial z}=\frac{z}{r}$$
            从而结果为
            $$\frac{a_1x+a_2y+a_3z}{r}=\frac{\vec{a}\cdot\rv}{r}$$
        }

        \item $\nabla\cdot(r^2\vec{a})$
        
        \sol{
            设$\vec{a}=(a_1,a_2,a_3)$，根据散度定义知其为
            $$\frac{\partial a_1r^2}{\partial x}+\frac{\partial a_2r^2}{\partial y}+\frac{\partial a_3r^2}{\partial z}$$
            直接计算有
            $$\frac{\partial r^2}{\partial x}=2x,\quad \frac{\partial r^2}{\partial y}=2y,\quad \frac{\partial r^2}{\partial z}=2z$$
            从而结果为
            $$2a_1x+2a_2y+2a_3z=2\vec{a}\cdot\rv$$
        }

        \item $\nabla\cdot(r^n\vec{a})$，其中$n$为整数。
        
        \sol{
            设$\vec{a}=(a_1,a_2,a_3)$，根据散度定义知其为
            $$\frac{\partial a_1r^n}{\partial x}+\frac{\partial a_2r^n}{\partial y}+\frac{\partial a_3r^n}{\partial z}$$
            直接通过复合函数求导公式计算有
            $$\frac{\partial r^n}{\partial x}=nr^{n-2}x,\quad \frac{\partial r^n}{\partial y}=nr^{n-2}y,\quad \frac{\partial r^n}{\partial z}=nr^{n-2}z$$
            从而结果为
            $$nr^{n-2}a_1x+nr^{n-2}a_2y+nr^{n-2}a_3z=nr^{n-2}\vec{a}\cdot\rv$$
        }

        \item $\nabla\cdot\frac{\rv}{r^3}$
        
        \sol{
            根据散度定义知其为
            $$\frac{\partial xr^{-3}}{\partial x}+\frac{\partial yr^{-3}}{\partial y}+\frac{\partial zr^{-3}}{\partial z}$$
            利用前一问的计算结果与乘积求导公式有
            $$\frac{\partial xr^{-3}}{\partial x}=r^{-3}-3r^{-5}x^2,\quad \frac{\partial yr^{-3}}{\partial y}=r^{-3}-3r^{-5}y^2,\quad \frac{\partial zr^{-3}}{\partial z}=r^{-3}-3r^{-5}z^2$$
            从而结果为(利用$x^2+y^2+z^2=r^2$)
            $$3r^{-3}-3r^{-5}(x^2+y^2+z^2)=0$$
        }

        \item $\nabla\cdot(f(r)\vec{a})$
        
        \sol{
            设$\vec{a}=(a_1,a_2,a_3)$，根据散度定义知其为
            $$\frac{\partial a_1f(r)}{\partial x}+\frac{\partial a_2f(r)}{\partial y}+\frac{\partial a_3f(r)}{\partial z}$$
            直接通过复合函数求导公式计算有
            $$\frac{\partial f(r)}{\partial x}=f'(r)\frac{x}{r},\quad\frac{\partial f(r)}{\partial y}=f'(r)\frac{y}{r},\quad\frac{\partial f(r)}{\partial z}=f'(r)\frac{z}{r}$$
            从而结果为
            $$f'(r)\frac{xa_1+ya_2+za_3}{r}=\frac{f'(r)}{r}\vec{a}\cdot\rv$$
        }

        \item $\nabla\cdot\nabla(f(r))$
        
        \sol{
            根据梯度定义与上一问计算可知
            $$\nabla(f(r))=\bigg(f'(r)\frac{x}{r},f'(r)\frac{y}{r},f'(r)\frac{z}{r}\bigg)$$
            设$\phi(r)=\frac{f'(r)}{r}$，根据散度定义知结果为
            $$\frac{\partial x\phi(r)}{\partial x}+\frac{\partial y\phi(r)}{\partial y}+\frac{\partial z\phi(r)}{\partial z}$$
            直接通过乘积求导与复合函数求导公式计算有
            $$\frac{\partial x\phi(r)}{\partial x}=\phi(r)+\phi'(r)\frac{x^2}{r},\quad\frac{\partial y\phi(r)}{\partial y}=\phi(r)+\phi'(r)\frac{y^2}{r},\quad\frac{\partial z\phi(r)}{\partial z}=\phi(r)+\phi'(r)\frac{z^2}{r}$$
            从而结果为(利用$x^2+y^2+z^2=r^2$)
            $$3\phi(r)+\phi'(r)\frac{x^2+y^2+z^2}{r}=3\phi(r)+r\phi'(r)$$
            进一步计算，由于
            $$\phi'(r)=\frac{f''(r)r-f'(r)}{r^2}$$
            代入可得最终结果
            $$\frac{2f'(r)}{r}+f''(r)$$
        }

        \item $\nabla\cdot(f(r)\rv)$
        
        \sol{
            根据散度定义知结果为
            $$\frac{\partial xf(r)}{\partial x}+\frac{\partial yf(r)}{\partial y}+\frac{\partial zf(r)}{\partial z}$$  
            从上一问关于$\phi(r)$的计算过程中可直接看出本题结果
            $$3f(r)+rf'(r)$$
        }
    \end{enumerate}

    \item (习题8.7.4)求下列向量场的旋度：
    \begin{enumerate}[(1)]
        \item $\Fv=(y^2+x^2,z^2+y^2,x^2+z^2)$

        \sol{
            设$\Fv=(P,Q,R)$，直接计算可发现
            $$R_y-Q_z=-2z,\quad P_z-R_x=-2x,\quad Q_x-P_y=-2y$$
            从而由定义有$\rot\Fv=(-2z,-2x,-2y)$。
        }

        \item $\Fv=(yz,zx,xy)$

        \sol{
            设$\Fv=(P,Q,R)$，直接计算可发现
            $$R_y-Q_z=x-x=0,\quad P_z-R_x=y-y=0,\quad Q_x-P_y=z-z=0$$
            从而由定义有$\rot\Fv=(0,0,0)$。
        }
    \end{enumerate}

    \item (习题8.7.7)证明下列向量场为有势场，并计算一个势函数：
    \begin{enumerate}[(1)]
        \item $\Fv=(-z\sin(xz)-y\cos(xy),-x\cos(xy),-x\sin(xz))$

        \sol{
            分为如下步骤：
            \begin{itemize}
                \item \textbf{有势场验证}
                
                设$\Fv=(P,Q,R)$，直接计算可发现
                $$R_y=Q_z=0,\quad P_z=R_x=-\sin(xz)-xz\cos(xz),\quad Q_x=P_y=-\cos(xy)+xy\sin(xy)$$
                从而由定义有$\rot\Fv=\vec{0}$，其为有势场。

                \item \textbf{势函数计算}
                
                与二维情况类似，设势函数为$u$，先由$\partial_1u(x,y,z)=-P$，计算$-P$对$x$的不定积分，得到
                $$u(x,y,z)=\sin(xy)-\cos(xz)+\phi(y,z)$$
                代入$\partial_2u(x,y,z)=-Q$、$\partial_3u(x,y,z)=-R$可发现
                $$\frac{\partial\phi(y,z)}{\partial y}=0,\quad\frac{\partial\phi(y,z)}{\partial z}=0$$
                由此其可取为0，一个势函数是
                $$u(x,y,z)=\sin(xy)-\cos(xz)$$
            \end{itemize}
        }

        \item $\Fv=\big(yz\ln(1+z^2),xz\ln(1+z^2),xy\ln(1+z^2)+\frac{2xyz^2}{1+z^2}\big)$

        \sol{
            分为如下步骤：
            \begin{itemize}
                \item \textbf{有势场验证}
                
                设$\Fv=(P,Q,R)$，直接计算可发现
                $$R_y=Q_z=x\ln(1+z^2)+\frac{2xz^2}{1+z^2}$$
                $$P_z=R_x=y\ln(1+z^2)+\frac{2yz^2}{1+z^2}$$
                $$Q_x=P_y=z\ln(1+z^2)$$
                从而由定义有$\rot\Fv=\vec{0}$，其为有势场。

                \item \textbf{势函数计算}
                
                与二维情况类似，设势函数为$u$，先由$\partial_1u(x,y,z)=-P$，计算$-P$对$x$的不定积分，得到
                $$u(x,y,z)=-xyz\ln(1+z^2)+\phi(y,z)$$
                代入$\partial_2u(x,y,z)=-Q$、$\partial_3u(x,y,z)=-R$可发现
                $$\frac{\partial\phi(y,z)}{\partial y}=0,\quad\frac{\partial\phi(y,z)}{\partial z}=0$$
                由此其可取为0，一个势函数是
                $$u(x,y,z)=-xyz\ln(1+z^2)$$
            \end{itemize}
        }
    \end{enumerate}

    \item (习题9.1.2)指出下列关于$x$的表达式中的任意常数$C_1$、$C_2$是否独立：
    \begin{enumerate}
        \item[(2)]$y=C_1(1-\cos(2x))+C_2\sin^2x$
        
        \sol{
            提供两种方法：
            \begin{itemize}
                \item \textbf{求导计算}
                
                直接利用倍角公式计算可知(求导代表对$x$)
                $$\frac{D(y,y')}{D(C_1,C_2)}=(1-\cos(2x))\cdot(2\sin x\cos x)-\sin^2x\cdot(2\sin(2x))=0$$
                从而$C_1$、$C_2$不独立。
                
                \item \textbf{定义验证}
                
                由$1-\cos(2x)=2\sin^2x$可知
                $$y=(2C_1+C_2)\sin^2x$$
                因此取一个任意常数$C=2C_1+C_2$可替代两个任意常数的作用，$C_1$、$C_2$不独立。
            \end{itemize}

            \note 独立性相关的详细讨论可见本讲义23.2.2。
        }

        \item[(4)]$y=C_1\er^{\lambda x}+C_2x\er^{\lambda x}$
        
        \sol{
            直接求导计算可知(求导代表对$x$)
            $$\frac{D(y,y')}{D(C_1,C_2)}=\er^{\lambda x}(\er^{\lambda x}+\lambda x\er^{\lambda x})-x\er^{\lambda x}\cdot\lambda\er^{\lambda x}=\er^{2\lambda x}$$
            由于此式恒非零，$C_1$、$C_2$必然独立。
        }
    \end{enumerate}

    \item (习题9.1.3)验证下列表达式是相应微分变量为$x$的微分方程的解，并指出是特解还是通解：
    \begin{enumerate}
        \item[(3)]表达式为
        $$y=C_1\er^{\lambda_1x}+C_2\er^{\lambda_2x}$$
        微分方程为
        $$y''-(\lambda_1+\lambda_2)y'+\lambda_1\lambda_2y=0$$
        这里参数$\lambda_1\ne\lambda_2$。

        \sol{
            分为如下步骤：
            \begin{itemize}
                \item \textbf{验证为解}
                
                直接计算可知
                $$y'=\lambda_1C_1\er^{\lambda_1x}+\lambda_2C_2\er^{\lambda_2x}$$
                $$y''=\lambda_1^2C_1\er^{\lambda_1x}+\lambda_2^2C_2\er^{\lambda_2x}$$
                从而代入可验证的确为解。

                \item \textbf{参数独立性}
                
                直接求导计算可知
                $$\frac{D(y,y')}{D(C_1,C_2)}=\er^{\lambda_1x}\lambda_2\er^{\lambda_2x}-\er^{\lambda_2x}\lambda_1\er^{\lambda_1x}=(\lambda_2-\lambda_1)\er^{(\lambda_1+\lambda_2)x}$$
                由于$\lambda_1\ne\lambda_2$，此式恒非零，$C_1$、$C_2$必然独立。由于这是二阶微分方程，有两个任意常数即为通解。
            \end{itemize}
        }

        \item[(6)]表达式为
        $$y=x\bigg(\int_1^x\frac{\er^t}{t}\dr t+C\bigg)$$
        微分方程为
        $$xy'-y=x\er^x$$

        \sol{
            分为如下步骤：
            \begin{itemize}
                \item \textbf{验证为解}
                
                利用乘积求导与变上限积分求导结论可知
                $$y'=\int_1^x\frac{\er^t}{t}\dr t+C+x\cdot\frac{\er^x}{x}=\er^x+\int_1^x\frac{\er^t}{t}\dr t+C$$
                从而代入可验证的确为解。

                \item \textbf{参数独立性}
                
                直接求导计算可知
                $$\frac{\partial y}{\partial C}=x$$
                它在$x\ne0$时其非零，从而在不包含$x=0$的区间中其为任意常数。由于这是一阶微分方程，有一个任意常数即为通解。
            \end{itemize}
        }
    \end{enumerate}

    \item (习题9.2.1)求下列微分方程的通积分：
    \begin{enumerate}
        \item[(3)] $\sqrt{1+x^2}\dr y-\sqrt{1-y^2}\dr x=0$
        
        \sol{
            分为如下步骤：
            \begin{itemize}
                \item \textbf{通解计算}
                
                由于此方程有分离变量形式，在$y\in(-1,1)$时，两边同乘$\frac{1}{\sqrt{(1-y^2)(1+x^2)}}$并移项得到
                $$\frac{\dr y}{\sqrt{1-y^2}}=\frac{\dr x}{\sqrt{1+x^2}}$$
                利用上册教材3.2节的积分公式，两侧同积分有($C$为任意常数)
                $$\arcsin y=C+\ln(x+\sqrt{x^2+1})$$

                \item \textbf{特解获取}
                
                由于过程中除以了可能为0的$\sqrt{1-y^2}$，额外讨论$y=\pm1$的情况可发现$y=\pm1$为特解。

                \note 理论来说题目表述为``求通积分''或``求通解''时无需讨论特解，这里为了完整性列出。
            \end{itemize}

            \note 由于$y=\pm1$与$y=\sin(C+\ln(x+\sqrt{x^2+1}))$有相交，本题的全部解情况较为复杂，存在解曲线的拼接。
        }
        
        \item[(7)] $(2x^2+y^2)\dr x+(2xy+3y^2)\dr y=0$
        
        \sol{
            将此方程看作$P\dr x+Q\dr y=0$，直接计算可验证
            $$Q_x=P_y=2y$$
            从而这是一个\textbf{全微分方程}。设$\dr u=P\dr x+Q\dr y$，先由$\partial_1u(x,y)=P$，计算$P$对$x$的不定积分，得到
            $$u(x,y)=\frac{2}{3}x^3+xy^2+\varphi(y)$$
            代入$\partial_2u(x,y)=Q$可发现
            $$2xy+\varphi'(y)=2xy+3y^2$$
            也即$\varphi'(y)=3y^2$，从而($C$为任意常数)
            $$u(x,y)=\frac{2}{3}x^3+y^2x+y^3+C$$
            最终通解为
            $$\frac{2}{3}x^3+y^2x+y^3+C=0$$

            \note 解决微分方程问题的基本方法是\textbf{先判断类型}(可以考虑逐一验证学过的类型)，再\textbf{根据类型选择求解方式}。
        }

        \item[(10)] $(2x-y+4)\dr y+(x-2y+5)\dr x=0$
        
        \sol{
            分为如下步骤：
            \begin{itemize}
                \item \textbf{初步化简}
                
                仿照教材中化为变量分离方程的思路进行操作。设$u=x-x_0$、$v=y-y_0$，有$\dr u=\dr x$、$\dr v=\dr y$。我们希望能够消去常数项，即
                $$0=(2x-y+4)\dr y+(x-2y+5)\dr x=(2u-v)\dr v+(u-2v)\dr u$$
                解得$x_0=-1$、$y_0=2$，从而$u=x+1$、$v=y-2$。
                
                \item \textbf{齐次方程}
                
                此方程已经是关于$u$、$v$的齐次方程，在$u\ne0$时，设$z=\frac{v}{u}$，则根据$v=zu$有
                $$\dr v=u\dr z+z\dr u$$
                代入后，由$u\ne0$可两端同除以$u$，得到
                $$(2u-zu)\dr z+(1-z^2)\dr u=0$$
                
                \item \textbf{最终求解}
                
                由于此方程有变量分离形式，当$u\ne0$、$z\ne\pm1$时，两边同乘$\frac{1}{(1-z^2)u}$并移项得到
                $$\frac{z-2}{1-z^2}\dr z=\frac{\dr u}{u}$$
                两侧同不定积分，通过待定系数法可分解计算出
                $$\int\frac{z-2}{1-z^2}\dr z=\int\bigg(\frac{1}{2}\cdot\frac{1}{z-1}-\frac{3}{2}\cdot\frac{1}{z+1}\bigg)\dr z=\frac{1}{2}\ln|z-1|+\frac{3}{2}\ln|z+1|+C$$
                从而可得到通解形式($C$为任意常数)
                $$\frac{1}{2}\ln|z-1|-\frac{3}{2}\ln|z+1|=\ln|u|+C$$
                两边同作$\er$的指数得
                $$|z-1|^{1/2}|z+1|^{-3/2}=\er^C|u|$$
                为了消除根号，同平方得
                $$|z-1||z+1|^{-3}=\er^{2C}u^2$$
                考虑正负可去除两侧绝对值并将恒正的$\er^{2C}$改写为新的任意常数$C$，且要求$C$\textbf{非零}。不过，进一步观察可发现$C=0$时恰好对应$z=1$或$z=-1$，可验证其为特解，综合可将通解改写为
                $$(z-1)(z+1)^{-3}=Cu^2$$
                代入也即
                $$\frac{y-x-3}{x+1}\bigg(\frac{y+x-1}{x+1}\bigg)^{-3}=C(x+1)^2$$
                当$x\ne1$、$y+x-1\ne0$时，上式可进一步化简为
                $$y-x+3=C(y+x-1)^3$$
                可验证$y+x-1=0$时其仍然为解，且对应的解在$x=1$时仍满足方程，从而这是合理的通解形式。
            \end{itemize}
        }
    \end{enumerate}

    \item (习题9.2.4(3))求以下微分变量为$x$的微分方程的通积分：
    $$nxy'+2y=xy^{n+1}$$
    这里参数$n$为正整数。

    \sol{
        分为如下步骤：
        \begin{itemize}
            \item \textbf{通解计算}
            
            这符合\textbf{伯努利方程}的形式，从而在$y\ne0$时，两边同乘$\frac{1}{y^{n+1}}$并换元$u=\frac{1}{y^n}$，则由复合函数求导公式可知
            $$u'=-\frac{ny'}{y^{n+1}}$$
            从而方程化为
            $$-xu'+2u=x$$
            这已经是一阶线性微分方程，可以直接求解，我们下面进行适当的配凑以方便计算。由于
            $$(x^ku)'=kx^{k-1}u+x^ku'$$
            取$k=-2$，在$x\ne0$时方程两边同乘$x^{-3}$可得
            $$(x^{-2}u)'=-x^{-2}$$
            从而同积分得到通解($C$为任意常数)
            $$\frac{u}{x^2}=\frac{1}{x}+C$$
            代入$u$的表达式并化简成
            $$1=y^n(Cx^2+x)$$
            这已经是合理的通解形式(注意此形式已经可以保证$y\ne0$、$x\ne0$)。

            \item \textbf{特解获取}
            
            由于过程中除以了$y^{n+1}$，额外讨论$y=0$的情况可发现$y=0$为一个特解。
        \end{itemize}
    }
    
    \item (习题9.2.5)将下列微分变量为$x$的微分方程化为线性微分方程：
    \begin{enumerate}
        \item[(3)]$3xy^2y'+y^3+x^3=0$
        
        \sol{
            利用复合函数求导公式可知$(y^3)'=3y^2y'$，从而根据形式可看出令$u=y^3$方程将化为
            $$xu'+u+x^3=0$$
            在$x\ne0$时其为
            $$u'+\frac{1}{x}u=-x^2$$
            这已经是线性微分方程。
        }

        \item[(4)]$y'=\sec y+x\tan y$
        
        \sol{
            为了进一步化简，用$\sin$、$\cos$写出右侧得到
            $$y'=\frac{1}{\cos y}+x\frac{\sin y}{\cos y}$$
            两侧同乘$\cos y$得
            $$y'\cos y=1+x\sin y$$
            同样利用复合函数求导公式，设$u=\sin y$即有
            $$u'=1+xu$$
            这已经是线性微分方程。
        }
    \end{enumerate}

    \item (习题9.2.8)镭的衰变速率与其现存最成正比。经测定，一块镭经过1600年后，只剩原始量$R_0$的一半，问1克镭经过一年后衰变多少毫克？
    
    \sol{
        设现存量的克数为$R(t)$，这里$t$为经过的年数，由条件可知存在$k>0$使得
        $$R'(t)=-kR(t),\quad R(0)=R_0,\quad R(1600)=\frac{1}{2}R_0$$
        从$R'(t)=-kR(t)$利用一阶线性微分方程知识可求解出存在$C$使得
        $$R(t)=C\er^{-kt}$$
        代入$R(0)=R_0$可知$C=R_0$，代入$R(1600)=\frac{1}{2}R_0$可知$k=\frac{\ln 2}{1600}$，从而化简得到表达式
        $$R(t)=R_02^{-\frac{t}{1600}}$$
        代入$R_0$为1、$t$为1可知最终克数为
        $$2^{-\frac{1}{1600}}$$
        从而衰变的毫克数为
        $$1000\big(1-2^{-\frac{1}{1600}}\big)\approx0.433$$
    }

    \item (总练习题9.3)假定降落伞在降落过程中所受的阻力与其下降速度成正比，比例系数为$k$。证明一个质量为$m$的跳伞运动员在空中运动的速度函数$v(t)$满足
    $$v'(t)+\frac{k}{m}v(t)=g$$
    并求出此方程的通解与$t$趋于正无穷时的极限。

    \sol{
        分为如下步骤：
        \begin{itemize}
            \item \textbf{方程构造}
            
            设向下的力为$F(t)$，向下的速度大小为$v(t)$，根据牛顿第二定律可知
            $$mv'(t)=F(t)$$
            再由题中条件有
            $$F(t)=mg-kv(t)$$
            代入化简得方程。

            \item \textbf{通解计算}
            
            这是一个一阶线性微分方程，直接计算可得
            $$\int\frac{k}{m}\dr t=\frac{kt}{m}+C$$
            由此可取一个积分因子
            $$\mu(t)=\er^{kt/m}$$
            方程两边同乘$\mu(t)$有
            $$\big(\er^{kt/m}v(t)\big)'=g\er^{kt/m}$$
            积分可得($C$为任意常数)
            $$\er^{kt/m}v(t)=\frac{mg}{k}\er^{kt/m}+C$$
            最终解出
            $$v(t)=\frac{mg}{k}+C\er^{-kt/m}$$

            \note \textbf{积分因子法}是最推荐掌握的一阶线性方程解法，因为它易于记忆，也相对灵活。

            \item \textbf{极限情况}
            
            根据物理意义可知$k>0$、$m>0$，由此
            $$\lim_{t\to+\infty}v(t)=\frac{mg}{k}$$
            也即最终速度的极限恰好使阻力与重力平衡。
        \end{itemize}
    }

    \item (总练习题9.4)设$p(t)$是时刻$t$时某一种群生物的数量，统计资料表明可以认为其满足
    $$p'(t)=(a-bp(t))p(t)$$
    这里参数$a>0$、$b>0$。设$p(t_0)=p_0$\ (这里$p_0>0$)，求函数$p(t)$在$t\ge t_0$时的表达式。

    \sol{
        这符合\textbf{伯努利方程}的形式，从而在$p(t)\ne0$时，两边同乘$\frac{1}{p^2(t)}$并换元$u(t)=\frac{1}{p(t)}$，可利用复合函数求导计算得方程化为
        $$-u'(t)=au(t)-b$$
        这是常系数一阶线性微分方程，与上题完全类似可解出存在$C$使得
        $$u(t)=\frac{b}{a}+C\er^{-at}$$
        从而
        $$p(t)=\frac{a}{b+Ca\er^{-at}}$$
        代入初始条件可解出
        $$C=\frac{a-bp_0}{ap_0}\er^{at_0}$$
        从而最终有
        $$p(t)=\frac{ap_0}{bp_0+(a-bp_0)\er^{-a(t-t_0)}}$$
        可以验证$t\ge t_0$时分母一定大于等于$a\er^{-a(t-t_0)}$，不会为0，且$p(t)\ne0$恒成立，从而此表达式合理。
    }
\end{enumerate}
\subsubsection{补充题解答}
\begin{enumerate}
    \item (习题8.6.2)计算
    $$\oiint_{S^+}xz^2\dr y\dr z+(x^2y-z^3)\dr z\dr x+(2xy+y^2z)\dr x\dr y$$
    其中$S$为$z=\sqrt{a^2-x^2-y^2}$与$z=0$围成的半球面，这里参数$a>0$。

    \sol{
        分为如下步骤：
        \begin{itemize}
            \item \textbf{高斯公式}
            
            将被积函数看作$P\dr y\dr z+Q\dr z\dr x+R\dr x\dr y$，由初等函数知$P$、$Q$、$R$在$\rb^3$上可导且导函数连续，直接计算得
            $$P_x+Q_y+R_z=z^2+x^2+y^2$$

            由此设半球$0\le z\le\sqrt{a^2-x^2-y^2}$为区域$\Omega$，利用高斯公式可将积分化为
            $$\iiint_\Omega(x^2+y^2+z^2)\dr V$$

            \item \textbf{积分计算}
            
            根据区域的旋转对称性，可以考虑球坐标换元
            $$x=\rho\sin\varphi\cos\theta,\quad y=\rho\sin\varphi\sin\theta,\quad z=\rho\cos\varphi$$
            
            利用球坐标的$\varphi$的几何意义可知$z\ge0$对应
            $$0\le\varphi\le\frac{\pi}{2}$$
            将另一个约束转化为$\rho^2\le a^2$后，结合球坐标自然范围可得最终范围
            $$0\le\rho\le a,\quad0\le\varphi\le\frac{\pi}{2},\quad 0\le\theta\le2\pi$$        
            将上述条件确定的$(\rho,\varphi,\theta)$区域记为$\Omega_0$。此时被积函数成为$\rho^2$，直接计算Jacobi行列式可知
            $$\dr x\dr y\dr z=\rho^2\sin\varphi\dr\rho\dr\varphi\dr\theta$$
            原积分化为
            $$\iiint_{\Omega_0}\rho^4\sin\varphi\dr\rho\dr\varphi\dr\theta$$
            利用区域描述可直接写出累次积分
            $$\int_0^a\dr\rho\int_0^{\pi/2}\dr\varphi\int_0^{2\pi}\rho^4\sin\varphi\dr\theta$$
            依次积分得到其为
            $$2\pi\int_0^a\dr\rho\int_0^{\pi/2}\rho^4\sin\varphi\dr\varphi=2\pi\int_0^a\rho^4\dr\rho=\frac{2}{5}\pi a^5$$
        \end{itemize}
    }

    \item (习题8.6.6)计算
    $$\iint_S\Fv\cdot\nv\dr S$$
    其中$\Fv=(x^4,y^4,z^4)$，$S$为锥面$x^2+y^2=z^2$在$z\in[0,h]$部分的外侧($x^2+y^2\ge z^2$侧)，$\nv$为$S$对应位置的单位外法向量，这里参数$h>0$。

    \sol{
        分为如下步骤：
        \begin{itemize}
            \item \textbf{化为第二型曲面积分}
            
            利用第一型曲面积分与第二型曲面积分的关系可知
            $$\nv\dr S=(\dr y\dr z,\dr z\dr x,\dr x\dr y)$$
            由此原积分即化为($S$以外侧为正定向)
            $$\iint_Sx^4\dr y\dr z+y^4\dr z\dr x+z^4\dr x\dr y$$

            \item \textbf{高斯公式}
            
            将被积函数看作$P\dr y\dr z+Q\dr z\dr x+R\dr x\dr y$，由初等函数知$P$、$Q$、$R$在$\rb^3$上可导且导函数连续，直接计算得
            $$P_x+Q_y+R_z=4(x^3+y^3+z^3)$$
            由于积分区域并非封闭曲面，我们补上顶部的
            $$S_0:\quad z=h,\quad x^2+y^2\le h^2$$
            并取上侧为正向。

            根据定义$S_0$在$S$的上方，因此$S_0$的上侧、$S$的下侧构成两者围成的封闭曲面的外侧。设围成的区域为$\Omega$，由高斯公式可知
            $$\iint_S(x^4\dr y\dr z+y^4\dr z\dr x+z^4\dr x\dr y)+\iint_{S_0}(x^4\dr y\dr z+y^4\dr z\dr x+z^4\dr x\dr y)=\iiint_\Omega4(x^3+y^3+z^3)\dr V$$
            由此设
            $$I_1=\iint_{S_0}(x^4\dr y\dr z+y^4\dr z\dr x+z^4\dr x\dr y),\quad I_2=\iiint_\Omega(x^3+y^3+z^3)\dr V$$
            结果即为$4I_2-I_1$。

            \item \textbf{$I_1$计算}
            
            由于$I_1$是平面区域上的第二类曲面积分，化为第一类可方便计算。利用第一型曲面积分与第二型曲面积分的关系可知(这里$\nv_0$代表第二型曲面积分对应定向的单位法向量)
            $$\nv_0\dr S=(\dr y\dr z,\dr z\dr x,\dr x\dr y)$$   
            根据几何可知$z=h$平面的一个上侧的单位法向量是$(0,0,1)$，由此即得原积分为(已代入$z=h$)
            $$\iint_{S_0}(x^4,y^4,z^4)\cdot(\dr y\dr z,\dr z\dr x,\dr x\dr y)=\iint_{S_0}(x^4,y^4,h^4)\cdot(0,0,1)\dr S=h^4\iint_{S_0}\dr S$$
            利用第一型曲面积分的几何意义，$\iint_{S_0}\dr S$即为曲面面积，而$S_0$的约束为$z=h$、$x^2+y^2\le h^2$，是一个半径为$h$的圆，从而面积为$\pi h^2$，最终得到
            $$I_1=\pi h^6$$

            \item \textbf{$I_2$计算}
            
            由于$\Omega$满足$(x,y,z)\in\Omega$当且仅当$(-x,y,z)\in\Omega$，利用区域折半可知
            $$\iiint_\Omega x^3\dr V=0$$
            同理
            $$\iiint_\Omega y^3\dr V=0$$
            由此可知
            $$I_2=\iiint_\Omega z^3\dr x\dr y\dr z$$
            由于此积分对$x$、$y$为常数，先对$x$、$y$积分，再对$z$积分，将其写成
            $$I_2=\int_0^h\dr z\iint_{D_z}z^3\dr x\dr y$$
            $D_z$为给定$z$时$(x,y)$的区域$x^2+y^2\le z^2$，利用重积分的几何意义可知$\iint_{D_z}\dr x\dr y$为其面积$\pi z^2$，从而积分结果为
            $$I_2=\int_0^hz^3\cdot\pi z^2\dr z=\frac{\pi h^6}{6}$$

            \item \textbf{综合结果}

            综合以上最终可知积分为
            $$4I_2-I_1=-\frac{\pi h^6}{3}$$
        \end{itemize}
    }

    \item (习题8.6.9)设$S$为封闭的简单曲面，证明
    $$\oiint_S\cos(\Fv,\nv)\dr S=0$$
    其中$\Fv=(\cos z,\sin z,1)$，$\nv$为$S$对应位置的单位外法向量，$(\Fv,\nv)$表示两者夹角。

    \sol{
        为了能够进行计算，分为如下步骤：
        \begin{itemize}
            \item \textbf{化为第二型曲面积分}
            
            由单位向量定义与内积性质可计算知
            $$\cos(\Fv,\nv)=\frac{\Fv\cdot\nv}{|\Fv||\nv|}=\frac{1}{\sqrt2}\Fv\cdot\nv$$
            原积分即
            $$\frac{1}{\sqrt2}\oiint_S\Fv\cdot\nv\dr S$$
            与上题相同可知这可以进一步化为
            $$\frac{1}{\sqrt2}\oiint_{S^+}\cos z\dr y\dr z+\sin z\dr z\dr x+\dr x\dr y$$

            \item \textbf{高斯公式}

            将被积函数看作$P\dr y\dr z+Q\dr z\dr x+R\dr x\dr y$，由初等函数知$P$、$Q$、$R$在$\rb^3$上可导且导函数连续，直接计算得
            $$P_x+Q_y+R_z=0+0+0=0$$
            由此设$S$围成了区域$\Omega$，利用高斯公式可将积分化为
            $$\iiint_\Omega0\dr V=0$$
            这就得到了证明。
        \end{itemize}
    }

    \item (习题8.6.13)计算
    $$\oint_L(y^2+z^2-x^2)\dr x+(z^2+x^2-y^2)\dr y+(x^2+y^2-z^2)\dr z$$
    其中$L$为球面$z=\sqrt{2Rx-x^2-y^2}$与柱面$x^2+y^2=2rx$的交线，且定向与球面上侧成右手系，这里参数$R>r>0$。

    \sol{
        分为如下步骤：
        \begin{itemize}
            \item \textbf{斯托克斯公式}
            
            将被积函数看作$P\dr x+Q\dr y+R\dr z$，由初等函数知$P$、$Q$、$R$在$\rb^3$上可导且导函数连续，直接计算得
            $$R_y-Q_z=2y-2z,\quad P_z-R_x=2z-2x,\quad Q_x-P_y=2x-2y$$
            为了应用斯托克斯公式，考虑$L$在球面上围成的部分$S$，则原积分可化为
            $$2\iint_S(y-z)\dr y\dr z+(z-x)\dr z\dr x+(x-y)\dr x\dr y$$
            利用右手定则，$S$应选择指向上侧的法向作为正向。

            \item \textbf{积分计算}
            
            由本讲义23.1.1第3题可直接得到
            $$\iint_S(y-z)\dr y\dr z+(z-x)\dr z\dr x+(x-y)\dr x\dr y=\pi r^2R$$
            从而本题结果为
            $$2\pi r^2R$$
        \end{itemize}
    }

    \item (总练习题8.15)设$L$为平面
    $$x\cos\alpha+y\cos\beta+z\cos\gamma=d$$
    上的一条光滑的简单闭曲线，且定向与平面的法向量$\nv=(\cos\alpha,\cos\beta,\cos\gamma)$成右手系。设$L$围成的面积为$|S|$，计算
    $$\oint_L\begin{vmatrix}\dr x&\dr y&\dr z\\\cos\alpha&\cos\beta&\cos\gamma\\x&y&z\end{vmatrix}$$
    这里参数$\cos^2\alpha+\cos^2\beta+\cos^2\gamma=1$。

    \sol{
        分为如下步骤：
        \begin{itemize}
            \item \textbf{斯托克斯公式}
            
            将被积函数看作$P\dr x+Q\dr y+R\dr z$，这里
            $$P=\cos\beta z-\cos\gamma y,\quad Q=\cos\gamma x-\cos\alpha z,\quad R=\cos\alpha y-\cos\beta x$$
            由初等函数知$P$、$Q$、$R$在$\rb^3$上可导且导函数连续，直接计算得
            $$R_y-Q_z=2\cos\alpha,\quad P_z-R_x=2\cos\beta,\quad Q_x-P_y=2\cos\gamma$$
            为了应用斯托克斯公式，考虑$L$在平面上围成的部分$S$，则原积分可化为
            $$2\iint_S\cos\alpha\dr y\dr z+cos\beta\dr z\dr x+\cos\gamma\dr x\dr y$$
            利用右手定则，$S$应选择法向$\nv=(\cos\alpha,\cos\beta,\cos\gamma)$作为正向。

            \item \textbf{化为第一型曲面积分}
            
            由于曲面的\textbf{面积}与第一型曲面积分相关，需要进一步化为第一型曲面积分。

            利用第一型曲面积分与第二型曲面积分的关系可知(这里$\vec{n}$代表第二型曲面积分对应定向的单位法向量)
            $$\vec{n}\dr S=(\dr y\dr z,\dr z\dr x,\dr x\dr y)$$
            由平面选取了单位法向量$\nv$可得到原积分为
            $$2\iint_S(\cos\alpha,\cos\beta,\cos\gamma)\cdot(\dr y\dr z,\dr z\dr x,\dr x\dr y)=2\iint_S\nv\cdot\nv\dr S$$
            利用条件中法向量的单位性即得此为$2\iint_S\dr S$，根据第一型曲面积分的几何意义知这就是$2|S|$。
        \end{itemize}
    }

    \item (总练习题8.16)给定点$(x_0,y_0,z_0)\in\rb^3$，设
    $$\rv=(x-x_0,y-y_0,z-z_0),\quad r=|\rv|$$
    证明
    $$\iiint_\Omega\frac{\dr x\dr y\dr z}{r}=\frac{1}{2}\oiint_S\cos(\rv,\nv)\dr S$$
    其中$S$光滑的简单闭曲面，$\Omega$为其围成的区域，且$(x_0,y_0,z_0)$不在其中。$\nv$为$S$对应位置的单位外法向量，$(\rv,\nv)$表示$\rv$与它的夹角。

    \sol{
        设右侧的曲面积分为$I$，也即要证明
        $$I=2\iiint_\Omega\frac{\dr x\dr y\dr z}{r}$$
        为了将曲面上的积分和区域内的积分产生关联，我们必须使用\textbf{高斯公式}，由此分为如下步骤：
        \begin{itemize}
            \item \textbf{化为第二型曲面积分}
            
            由单位向量定义与内积性质可计算知
            $$\cos(\rv,\nv)=\frac{\rv\cdot\nv}{|\rv||\nv|}=\frac{1}{r}\rv\cdot\nv$$
            原积分即
            $$\oiint_S\frac{\rv}{r}\cdot\nv\dr S$$
            利用第一型曲面积分与第二型曲面积分的关系可知
            $$\nv\dr S=(\dr y\dr z,\dr z\dr x,\dr x\dr y)$$
            由此即有(代入$\rv$的表达式)
            $$I=\oiint_{S^+}\frac{x}{r}\dr y\dr z+\frac{y}{r}\dr z\dr x+\frac{z}{r}\dr x\dr y$$

            \item \textbf{高斯公式}

            将被积函数看作$P\dr y\dr z+Q\dr z\dr x+R\dr x\dr y$，由条件可知$P$、$Q$、$R$在$\Omega$上可导且导函数连续，直接计算得(代入$x^2+y^2+z^2=r^2$)
            $$P_x+Q_y+R_z=\frac{1}{r}-\frac{x^2}{r^3}+\frac{1}{r}-\frac{y^2}{r^3}+\frac{1}{r}-\frac{z^2}{r^3}=\frac{2}{r}$$
            由此利用高斯公式可得到
            $$I=\iiint_\Omega\frac{2}{r}\dr V$$
            这就是要证的结论。
        \end{itemize}
    }

    \item (习题8.7.5)设数量场$u=f(x,y,z)$与向量场$\Fv=\vec{g}(x,y,z)$各分量都有连续的二阶偏导数，证明(以下的$\nabla$表示$\nabla_{x,y,z}$，即对$x,y,z$的梯度、散度或旋度)：
    \begin{enumerate}[(1)]
        \item $\nabla\times(\nabla u)=\vec{0}$
        
        \sol{
            由于
            $$\nabla u=(u_x,u_y,u_z)$$
            将其看作$(P,Q,R)$，由于$u$的二阶偏导连续，它们可交换，从而直接计算可发现
            $$R_y=Q_z=u_{yz},\quad P_z=R_x=u_{zx},\quad Q_x=P_y=u_{xy}$$
            由定义即有$\nabla\times(\nabla u)=\vec{0}$。
        }
        
        \item $\nabla\cdot(\nabla\times\Fv)=0$
        
        \sol{
            设$\Fv=(P,Q,R)$，由于
            $$\nabla\times\Fv=(R_y-Q_z,P_z-R_x,Q_x-P_y)$$
            由于$P$、$Q$、$R$的二阶偏导连续，它们可交换，从而直接计算可发现
            $$(R_y-Q_z)_x+(P_z-R_x)_y+(Q_x-P_y)_z=R_{yx}-R_{xy}-Q_{zx}+Q_{xz}+P_{zy}-P_{yz}=0$$
            由定义即有$\nabla\cdot(\nabla\times\Fv)=0$。
        }
        
        \item $\nabla\times(u\Fv)=u\nabla\times\Fv+(\nabla u)\times\Fv$
        
        \sol{
            设$\Fv=(P,Q,R)$，直接计算有
            $$\nabla\times(u\Fv)=((uR)_y-(uQ)_z,(uP)_z-(uR)_x,(uQ)_x-(uP)_y)$$
            利用乘积求导公式可将其展开为关于$u$导数与关于$\Fv$导数的两个向量和，即
            $$(u_yR-u_zQ,u_zP-u_xR,u_xQ-u_yP)+(uR_y-uQ_z,uP_z-uR_x,uQ_x-uP_y)$$
            进一步计算可发现第一部分即为$(\nabla u)\times\Fv$，，第二部分即为$u\nabla\times\Fv$，从而得证。
        }
    
    \end{enumerate}

    \item (习题8.7.8)设$\rv=(x,y,z)$，$r=|\rv|$，求定义在$(0,+\infty)$函数$f$使得(以下的$\nabla$表示$\nabla_{x,y,z}$，即对$x,y,z$的梯度或散度)：
    \begin{enumerate}[(1)]
        \item $\nabla\cdot(f(r)\rv)=0$
        
        \sol{
            利用本讲义24.1.1第2题(7)的结论可知其成立当且仅当
            $$\forall r>0,\quad 3f(r)+rf'(r)=0$$
            直接计算可发现$(rf(r))'=f(r)+rf'(r)$，由此可以启发我们，为了配凑出三倍，在两侧同乘$r^2$，这样即有
            $$0=3r^2f(r)+r^3f'(r)=(r^3f(r))'$$
            于是可得
            $$r^3f(r)=C$$
            也即满足条件的$f(r)$为所有$\frac{C}{r^3}$，这里$C$为任意实数。

            \note 更多这类问题的求解方法将在下一章\textbf{常微分方程}中介绍。
        }
        
        \item $\nabla\cdot\nabla(f(r))=0$
        
        \sol{
            利用本讲义24.1.1第2题(6)的结论可知其成立当且仅当
            $$\forall r>0,\quad \frac{2f'(r)}{r}+f''(r)=0$$
            与上一问完全类似，为了配凑出两倍，在两侧同乘$r^2$，这样即有
            $$0=2rf'(r)+r^2f''(r)=(r^2f'(r))'$$
            于是可得
            $$r^2f'(r)=C$$
            也即
            $$f'(r)=\frac{C}{r^2}$$
            进一步积分得到
            $$f(r)=-\frac{C}{r}+C_1$$
            这就是所有满足条件的$f(r)$，这里$C$、$C_1$为任意实数。

            \note 更多这类问题的求解方法将在下一章\textbf{常微分方程}中介绍。
        }
    \end{enumerate}
\end{enumerate}

\subsection{一阶线性方程}
\subsubsection{齐次与非齐次线性方程}
线性算子

线性齐次微分方程

线性非齐次微分方程

解的关系

\subsubsection{齐次方程的解}
齐次方程解的形式(线性无关性与解的线性无关性、W行列式)

证明通解是全部解

一点为0恒为0

\subsubsection{常数变易法}
常数变易法的意义

一般线性微分方程的常数变易法

\subsection{二阶线性方程}
\subsubsection{二阶线性齐次常系数方程}
二阶线性递推出发

配凑等比

特征方程

三种情况

\subsubsection{二阶线性方程}
通解一般难以确定

W行列式与通解确定

特殊性质

\subsubsection{高阶线性方程}
常系数情况的类似配凑

算子分解

非齐次情况的W行列式


\newpage
\section{期中复习}
本章为针对教材七到九三章的期中考试复习题，请大家\textbf{务必关注每题解答后的注释}，强调的重点往往在注释中。

\subsection{习题解答}
\note 教材第九章的解答只进行\textbf{通解与特解的获取}，不研究这是否构成全部解等严谨性问题，一些进一步的分析可见本讲义第二十三章或25.2.5。注意理论来说求出的通解可以\textbf{完全不进行化简与合并}，实际几乎只取决于改卷标准，解答中我们尽量化到相对简洁的形式。

\subsubsection{作业解答}
\begin{enumerate}
    \item (习题9.2.13(1))求解微分变量为$x$的微分方程
    $$x^2y''=(y')^2$$

    \sol{
        分为如下步骤：
        \begin{itemize}
            \item \textbf{化为一阶求解}
            
            设$z=y'$，则方程变为
            $$x^2z'=z^2$$
            这是一个变量分离的方程，当$x\ne0$、$z\ne0$时方程两边同除以$x^2z^2$可将其化为
            $$\frac{\dr z}{z^2}=\frac{\dr x}{x^2}$$
            两边同积分得到($C$为任意常数)
            $$-\frac{1}{z}=-\frac{1}{x}+C$$
            化简得到
            $$\frac{1}{z}=\frac{1}{x}-C$$
            将它改写为
            $$z=\frac{x}{1-Cx}$$
            计算可验证上述解在$x=0$时仍满足方程。此外，验证可发现$z=0$为此方程的特解。

            \item \textbf{进一步积分}
            
            代入$z=y'$，分类讨论：
            \begin{enumerate}
                \item 对特解$z=0$，积分结果为($C_0$为任意常数)
                $$y=C_0$$
                \item 对通解$z=\frac{x}{1-Cx}$，若$C=0$，积分结果为($C_0$为任意常数)
                $$y=\frac{1}{2}x^2+C_0$$
                \item 对通解$z=\frac{x}{1-Cx}$，若$C\ne0$，将其化为
                $$y=\int\frac{x}{1-Cx}\dr x=\int\bigg(\frac{1}{C}\cdot\frac{1}{1-Cx}-\frac{1}{C}\bigg)\dr x$$
                可得到积分结果为($C_0$为任意常数)
                $$y=-\frac{1}{C^2}\ln|1-Cx|-\frac{x}{C}+C_0$$
            \end{enumerate}
        \end{itemize}
        
        \note 本题通解与全部解的关系\textbf{非常复杂}，详见本讲义23.2.3的讨论。
    }
    
    \item (习题9.2.14)判断下列微分方程是否为全微分方程，若是则求出通积分：
    \begin{enumerate}
        \item[(3)]$\er^y\dr x+(x\er^y-2y)\dr y=0$
        
        \sol{
            将此方程看作$P\dr x+Q\dr y=0$，直接计算可验证
            $$Q_x=P_y=\er^y$$
            从而这是一个全微分方程。设$\dr u=P\dr x+Q\dr y$，先由$\partial_1u(x,y)=P$，计算$P$对$x$的不定积分，得到
            $$u(x,y)=x\er^y+\varphi(y)$$
            代入$\partial_2u(x,y)=Q$可发现
            $$x\er^y+\varphi'(y)=x\er^y-2y$$
            也即$\varphi'(y)=-y^2$，从而($C$为任意常数)
            $$u(x,y)=x\er^y-y^2+C$$
            最终通解为
            $$x\er^y-y^2+C=0$$
        }

        \item[(8)]$(ax^2+by^2)\dr x+lxy\dr y=0$，这里$a$、$b$、$l$为常数。
        
        \sol{
            将此方程看作$P\dr x+Q\dr y=0$，直接计算可验证
            $$Q_x-P_y=(l-2b)y$$
            从而这是全微分方程当且仅当$l=2b$。此时，设$\dr u=P\dr x+Q\dr y$，先由$\partial_1u(x,y)=P$，计算$P$对$x$的不定积分，得到
            $$u(x,y)=\frac{1}{3}ax^3+bxy^2$$
            代入$\partial_2u(x,y)=Q$可发现
            $$2bxy+\varphi'(y)=lxy$$
            由$l=2b$也即$\varphi'(y)=0$，从而($C$为任意常数)
            $$u(x,y)=\frac{1}{3}ax^3+bxy^2+C$$
            最终通解为
            $$\frac{1}{3}ax^3+bxy^2+C=0$$
        }
    \end{enumerate}

    \item (习题9.2.15(2))用观察法求以下微分方程的积分因子，并计算通积分：
    $$(1+x^2+y^2+x)\dr x+y\dr y=0$$

    \sol{
        由于出现了$x^2+y^2$作为整体，且用$x\dr x+y\dr y$可以进行配凑
        $$x\dr x=\frac{1}{2}x^2,\quad y\dr y=\frac{1}{2}y^2,\quad x\dr x+y\dr y=\frac{1}{2}\dr(x^2+y^2)$$
        设$v=x^2+y^2$，方程则化为
        $$(1+v)\dr x+\frac{1}{2}\dr v=0$$
        这已经是变量分离的方程，从而
        $$\mu(x,y)=\frac{1}{1+v}=\frac{1}{1+x^2+y^2}$$
        即是一个\textbf{积分因子}。乘积分因子后方程化为
        $$\dr x+\frac{\dr v}{2(1+v)}=0$$
        分别积分并求和也即(由$v\ge0$可去除$\ln$的绝对值)
        $$\dr\bigg(x+\frac{1}{2}\ln(1+v)\bigg)=0$$
        由此通解为($C$为任意常数)
        $$x+\frac{1}{2}\ln(1+v)=C$$
        代入得到
        $$x+\frac{1}{2}\ln(1+x^2+y^2)=C$$
    }

    \item (习题9.2.16(6))利用积分因子求解微分方程：
    $$\er^x\dr x+(\er^x\cot y+2y\cos y)\dr y=0$$

    \sol{
        由于$\cot y$难以处理，将它看作$\frac{\cos y}{\sin y}$，方程两边同乘$\sin y$可以得到
        $$\er^x\sin y\dr x+(\er^x\cos y+2y\sin y\cos y)\dr y=0$$
        将此方程看作$P\dr x+Q\dr y=0$，直接计算可验证
        $$Q_x-P_y=\er^x\cos y$$
        从而这已经是一个全微分方程，$\sin y$即是积分因子。

        \note 若没有观察出此性质，也可以按照教材方法验证其存在\textbf{只和$y$有关}的积分因子并计算。

        设$\dr u=P\dr x+Q\dr y$，先由$\partial_1u(x,y)=P$，计算$P$对$x$的不定积分，得到
        $$u(x,y)=\er^x\sin y+\varphi(y)$$
        代入$\partial_2u(x,y)=Q$可发现
        $$\er^x\cos y+\varphi'(y)=\er^x\cos y+2y\sin y\cos y$$
        也即$\varphi'(y)=2y\sin y\cos y=y\sin 2y$。利用分部积分计算：
        $$\int y\sin(2y)\dr y=-\frac{1}{2}\int y\dr(\cos(2y))=-\frac{1}{2}y\cos(2y)+\frac{1}{2}\int\cos(2y)\dr y$$
        于是有($C$为任意常数)
        $$\varphi(y)=-\frac{1}{2}y\cos(2y)+\frac{1}{4}\sin(2y)+C$$
        从而
        $$u(x,y)=\er^x\sin y-\frac{1}{2}y\cos(2y)+\frac{1}{4}\sin(2y)+C$$
        最终通解为
        $$\er^x\sin y-\frac{1}{2}y\cos(2y)+\frac{1}{4}\sin(2y)+C=0$$
    }

    \item (习题9.3.1)设$f(x,y)$的偏导数$\partial_2f(x,y)$在闭矩形$R$上有界，证明$f(x,y)$在$R$上对$y$满足Lipschitz条件。
    
    \sol{
        设$|\partial_2f(x,y)|\le M$在$R$上恒成立。对
        $$(x_0,y_1)\in R,\quad (x_0,y_2)\in R,\quad y_1\ne y_2$$
        将$f(x_0,y)$看作对$y$的函数，由偏导数存在有界性，利用\textbf{微分中值定理}可知存在$y_1$、$y_2$之间的$y_0$使得
        $$\partial_2f(x_0,y_0)=\frac{f(x_0,y_1)-f(x_0,y_2)}{y_1-y_2}$$
        移项取绝对值即
        $$|f(x_0,y_1)-f(x_0,y_2)|=|\partial_2f(x_0,y_0)||y_1-y_2|\le M|y_1-y_2|$$
        而当$y_1=y_2$时，取$y_0=y_1$，代入可发现上式同样成立，由此其满足Lipshitz常数为$M$的Lipschitz条件。

        \note 这是Lipschitz条件的\textbf{常用判定方式}。
    }

    \item (习题9.3.3)考虑函数$f(x,y)=x^2y+x$，初值$x_0=0$、$y_0=0$，区域为
    $$R=\{(x,y)\mid |x|\le1,\ |y|\le1\}$$
    应用本节的存在唯一性定理，估计Lipschitz常数$L$与$f(x,y)$在$R$上的最大值$M$，进一步估计$y'=f(x,y)$过点$(x_0,y_0)$的解的存在区间$2h$。

    \sol{
        分为如下步骤：
        \begin{itemize}
            \item \textbf{估计$L$}
            
            直接计算可发现
            $$|f(x_0,y_1)-f(x_0,y_2)|=|x_0^2y_1-x_0^2y_2|=x_0^2|y_1-y_2|$$
            由此根据Lipshitz条件的定义，Lipshitz常数必须大于等于$x_0^2$在$R$上的最大值(否则给定$x_0$使得$x_0^2$取到最大，取$y_1\ne y_2$即得矛盾)，也即$L$可取任何大于等于1的实数。
            
            \item \textbf{估计$M$}
            
            直接放缩有
            $$|f(x,y)|=|x^2y+x|\le x^2|y|+|x|=2$$
            且当$x=y=1$时可取到，从而$|f(x,y)|$在$R$上最大值为2，即$M=2$。

            \item \textbf{估计$h$}
            
            本题中$a=b=1$，因此
            $$h=\min\bigg\{1,\frac{1}{M}\bigg\}=\frac{1}{2}$$
            估计得到的$2h=1$。
        \end{itemize}
    }

    \item (习题9.3.4)计算初值问题
    $$\frac{\dr y}{\dr x}=x+y,\quad y(0)=-2$$
    的皮卡序列及其极限。

    \sol{
        \begin{itemize}
            \item \textbf{皮卡序列确定}
            
            根据定义有
            $$y_0(x)=-2$$
            尝试计算发现
            $$y_1(x)=-2+\int_0^x(t+y_0(t))\dr t=-2-2x+\frac{1}{2}x^2$$
            $$y_2(x)=-2+\int_0^x(t+y_1(t))\dr t=-2-2x-\frac{1}{2}x^2+\frac{1}{6}x^3$$
            $$y_3(x)=-2+\int_0^x(t+y_2(t))\dr t=-2-2x-\frac{1}{2}x^2-\frac{1}{6}x^3+\frac{1}{24}x^4$$
            以此(或多算几项后)可猜测$n\ge2$时
            $$y_n(x)=-2-2x-\sum_{k=2}^n\frac{x^k}{k!}+\frac{x^{n+1}}{(n+1)!}$$
            下面归纳证明。当$n=2$时已经满足，而若$n$时满足，$n+1$时有
            $$y_{n+1}(x)=-2+\int_0^x(t+y_n(t))\dr t$$
            由于
            $$t+y_n(t)=-2-t-\sum_{k=2}^n\frac{t^k}{k!}+\frac{t^{n+1}}{(n+1)!}$$
            作不定积分可得到
            $$-2t-\frac{1}{2}t^2-\sum_{k=2}^n\frac{t^{k+1}}{(k+1)!}+\frac{t^{n+2}}{(n+2)!}$$
            利用微积分基本定理，并将求和下标从$k$改为$p=k+1$可发现
            $$y_{n+1}(x)=-2-2x-\frac{1}{2}x^2-\sum_{p=3}^{n+1}\frac{x^p}{p!}+\frac{x^{n+2}}{(n+2)!}$$
            再将$-\frac{1}{2}x^2$合并到求和中即得结果。
            
            \item \textbf{极限计算}
            
            根据本节定理，皮卡序列的极限为此微分方程初值问题的解，由此只需求解初值问题即可。

            将方程看作$y'(x)-y(x)=x$，这是一个一阶线性微分方程，直接计算可得
            $$\int(-1)\dr x=-x+C$$
            由此可取一个积分因子
            $$\mu(x)=\er^{-x}$$
            方程两边同乘$\mu(x)$有
            $$\big(\er^{-x}y(x)\big)'=x\er^{-x}$$
            分部积分可得($C$为任意常数)
            $$\er^{-x}y(x)=\int x\er^{-x}\dr x=-\int x\dr\er^{-x}=-x\er^{-x}+\int\er^{-x}\dr x=(-x-1)\er^{-x}+C$$
            代入初值条件$y(0)=-2$即可解出
            $$y(x)=-x-1-\er^x$$
            
            \note 直接计算极限所需的结论事实上在教材10.5节例9(1)，\textbf{超出了当前知识范围}。
        \end{itemize}
    }

    \item (习题9.3.6(1))利用解的存在唯一性定理计算以下初值问题的解的存在区间：
    $$y'=\sqrt{|x-y|},\quad y(4)=1$$
    区域为
    $$R=\{(x,y)\mid|x-4|\le2,\ |y-1|\le1\}$$
    
    \sol{
        分为如下步骤：
        \begin{itemize}
            \item \textbf{初步分析}
            
            由于条件的区域为$|x-4|\le2$、$|y-1|\le1$，可发现$x\ge y$恒成立，从而可去除绝对值。设
            $$f(x,y)=\sqrt{x-y}$$
            直接计算可知
            $$\partial_2f(x,y)=-\frac{1}{2\sqrt{x-y}}$$
            若$\partial_2f(x,y)$有界，利用本部分第5题结论可知其对$y$满足Lipschitz条件，但在$x\to2^+$、$y\to2^-$时，$\partial_2f(x,y)$将趋于无穷，不满足。

            不过，只要略微缩小区间，给定$\varepsilon>0$，考虑新区域
            $$R_\varepsilon=\{(x,y)\mid |x-4|\le 2-\varepsilon,\ |y-1|\le1\}$$
            这之中有$x-y\ge\varepsilon$，从而
            $$|\partial_2f(x,y)|\le\frac{1}{2\sqrt\varepsilon}$$
            有界。由此，我们先尝试在$R_\varepsilon$中进行研究。

            \item \textbf{缩小区域研究}
            
            取定某$0<\varepsilon<2$，定义
            $$R_\varepsilon=\{(x,y)\mid |x-4|\le 2-\varepsilon,\ |y-1|\le1\}$$
            根据定义可发现区域中
            $$|\partial_2f(x,y)|=\frac{1}{2\sqrt{x-y}}\le\frac{1}{2\sqrt{4-(2-\varepsilon)-2}}=\frac{1}{2\sqrt\varepsilon}$$
            由此利用本部分第5题结论可知$f(x,y)$满足Lipschitz条件。

            此时直接估算可发现$R_\varepsilon$中
            $$|f(x,y)|\le\sqrt{(4+2-\varepsilon)-(1-1)}=\sqrt{6-\varepsilon}$$

            根据本节定理，这里$a=2-\varepsilon$、$b=1$，由此对应的最大存在区间
            $$h_\varepsilon=\min\bigg\{2-\varepsilon,\frac{1}{\sqrt{6-\varepsilon}}\bigg\}$$
            而解在$[4-h_\varepsilon,4+h_\varepsilon]$存在。

            \item \textbf{还原区域}
            
            为了得知原问题解的最大存在区间，我们需要使得$h_\varepsilon$尽量大。由于$0<\varepsilon<2$可以发现当$2-\varepsilon=\frac{1}{\sqrt{6-\varepsilon}}$时取到最大值。

            此时$\varepsilon$的取值对应一个三次方程，可以代入求根公式求解，但由于形式过于复杂，这里只给出数值近似解
            $$\varepsilon\approx1.5272,\quad h_\varepsilon\approx0.4728$$
            由此解至少保证在区间$[3.5272,4.4728]$存在。

            \note 教材答案中的$\frac{1}{\sqrt6}$相当于只考虑了$\varepsilon\to0$的情况，比此处答案的区间更小。
        \end{itemize}
        
    }

    \item (习题9.4.1)证明下列关于$x$的函数组线性无关：
    \begin{enumerate}[(1)]
        \item $\er^{\lambda x},x\er^{\lambda x}$
        
        \sol{
            \note 注意教材从Wronski行列式恒为0推线性相关事实上\textbf{并未用到}$\varphi_1(x)$、$\varphi_2(x)$是方程的解，由此只要\textbf{一点处}Wronski行列式非零，就必然线性无关。

            设$y_1=\er^{\lambda x}$、$y_2=x\er^{\lambda x}$，计算Wronski行列式可得(求导代表对$x$)
            $$y_1y_2'-y_2y_1'=\er^{\lambda x}(\er^{\lambda x}+\lambda x\er^{\lambda x})-\lambda\er^{\lambda x}x\er^{\lambda x}=\er^{2\lambda x}$$
            由于它恒非0，$\er^{\lambda x}$、$x\er^{\lambda x}$必然线性无关。
        }
        
        \item $\cos(\beta x),\sin(\beta x)$，这里参数$\beta\ne0$。
        
        \sol{
            设$y_1=\cos(\beta x)$、$y_2=\sin(\beta x)$，计算Wronski行列式可得(求导代表对$x$)
            $$y_1y_2'-y_2y_1'=\beta\cos^2(\beta x)+\beta\sin^2(\beta x)=\beta$$
            由于它恒非0，$\cos(\beta x)$、$\sin(\beta x)$必然线性无关。
        }

        \item $\er^{ax}\cos(\beta x),\er^{ax}\sin(\beta x)$，这里参数$\beta\ne0$。
        
        \sol{
            设$y_1=\er^{ax}\cos(\beta x)$、$y_2=\er^{ax}\sin(\beta x)$，计算Wronski行列式可得(求导代表对$x$)
            $$y_1y_2'-y_2y_1'=\er^{2ax}\big(\beta\cos^2(\beta x)+a\sin(\beta x)\cos(\beta x)+\beta\sin^2(\beta x)-a\sin(\beta x)\cos(\beta x)\big)=\beta\er^{2ax}$$
            由于它恒非0，$\er^{ax}\cos(\beta x)$、$\er^{ax}\sin(\beta x)$必然线性无关。
        }
    \end{enumerate}

    \item (习题9.4.3)设$y=\varphi(x)$是$(a,b)$上的线性齐次微分方程
    $$y''(x)+p(x)y'(x)+q(x)y(x)=0$$
    的一个不恒为0的解(这里$p(x)$、$q(x)$在$(a,b)$连续)，若有$x_0\in(a,b)$使得$\varphi(x_0)=0$，证明$\varphi'(x_0)\ne0$。

    \sol{
        考虑恒为0的函数$\psi(x)$，由定义可知$\psi(x_0)=\psi'(x_0)=0$，且$\psi(x)$为微分方程一个解。

        考虑\textbf{反证法}，若$\varphi'(x_0)=\varphi(x_0)=0$，且$\varphi(x)$为线性齐次微分方程的解，根据初值问题
        $$\begin{cases}y''(x)+p(x)y'(x)+q(x)y(x)=0\\y(x_0)=0,\quad y'(x_0)=0\end{cases}$$
        的解的\textbf{唯一性}(本节定理1)可知必然$\varphi(x)=\psi(x)$恒成立，从而$\varphi(x)$也是恒为0的解，矛盾。
    }

    \item (习题9.4.4)设$y=\varphi_1(x)$、$y=\varphi_2(x)$是$(a,b)$上的线性齐次微分方程
    $$y''(x)+p(x)y'(x)+q(x)y(x)=0$$
    的两个线性无关的解(这里$p(x)$、$q(x)$在$(a,b)$连续)，证明它们没有公共的零点。

    \sol{
        考虑\textbf{反证法}，若$\varphi_1(x_0)=\varphi_2(x_0)=0$，计算可知Wronski行列式在$x_0$处的值为
        $$\varphi_1(x_0)\varphi_2'(x_0)-\varphi_2(x_0)\varphi_1'(x_0)=0$$
        根据本节定理2推论可知$\varphi_1(x)$、$\varphi_2(x)$线性相关，矛盾。

    }

    \item (习题9.4.5)证明$(a,b)$上的线性齐次微分方程
    $$y''(x)+p(x)y'(x)+q(x)y(x)=0$$
    存在两个线性无关的解(这里$p(x)$、$q(x)$在$(a,b)$连续)。

    \sol{
        任取$x_0\in(a,b)$，考虑初值问题
        $$\begin{cases}y''(x)+p(x)y'(x)+q(x)y(x)=0\\y(x_0)=1,\quad y'(x_0)=0\end{cases}$$
        $$\begin{cases}y''(x)+p(x)y'(x)+q(x)y(x)=0\\y(x_0)=0,\quad y'(x_0)=1\end{cases}$$
        由于它们的解都\textbf{存在唯一}(本节定理1)，设第一个初值问题的解为$\varphi_1(x)$，第二个初值问题的解为$\varphi_2(x)$，计算可知Wronski行列式在$x_0$处的值为
        $$\varphi_1(x_0)\varphi_2'(x_0)-\varphi_2(x_0)\varphi_1'(x_0)=1$$
        根据本节定理2可知$\varphi_1(x)$、$\varphi_2(x)$线性无关，从而符合要求。

        \note 这几道习题中可以看出初值问题解的\textbf{存在唯一性}是常用来构造证明的。
    }

    \item (习题9.5.1)求下列微分变量为$x$的微分方程的通解：
    \begin{enumerate}
        \item[(3)]$y''+6y'+9y=0$
        
        \sol{
            考虑其特征方程$\lambda^2+6\lambda+9=0$，两根为$-3$、$-3$，由此通解为($C_1$、$C_2$为任意常数)
            $$(C_1x+C_2)\er^{-3x}$$
        }

        \item[(6)]$y'''+2y''-y'=0$
        
        \sol{
            设$z=y'$，对应的方程为$z''+2z'-z=0$，考虑其特征方程$\lambda^2+2\lambda-1=0$，两根为$-1\pm\sqrt2$，由此通解为($C_1$、$C_2$为任意常数)
            $$z=C_1\er^{(-1+\sqrt2)x}+C_2\er^{(-1-\sqrt2)x}$$
            再进行一次积分可得到原方程通解($C_3$为任意常数)
            $$y=\frac{C_1}{-1+\sqrt2}\er^{(-1+\sqrt2)x}+\frac{C_2}{-1-\sqrt2}\er^{(-1-\sqrt2)x}+C_3$$
            也可将$\frac{C_1}{-1+\sqrt2}$、$\frac{C_2}{-1-\sqrt2}$改写为新的任意常数$C_1$、$C_2$，简化为
            $$y=C_1\er^{(-1+\sqrt2)x}+C_2\er^{(-1-\sqrt2)x}+C_3$$
        }
    \end{enumerate}

    \item (习题9.5.3)用待定系数法求下列微分变量为$x$的微分方程的一个特解：
    \begin{enumerate}
        \item[(6)]$y''-9y=\er^{3x}\cos x$
        
        \sol{
            计算可知$y''-9y$的特征方程为$\lambda^2-9=0$，两根为$\pm3$。查表得可以待定系数
            $$y=\er^{3x}(\alpha\cos x+\beta\sin x)$$
            直接计算可发现
            $$\er^{3x}(-6\alpha\sin x+6\beta\cos x-\alpha\cos x-\beta\sin x)=\er^{3x}\cos x$$
            从而对比系数得到可取$\alpha=-\frac{1}{37}$、$\beta=\frac{6}{10}$，一个特解为
            $$y=\er^{3x}\bigg(\frac{6}{37}\sin x-\frac{1}{37}\cos x\bigg)$$
        }

        \item[(7)]$y''-y=2\er^x-x^2$
        
        \sol{
            计算可知$y''-y$的特征方程为$\lambda^2-1=0$，两根为$\pm1$。由于右端分为$2\er^x$与$-x^2$两部分，我们考虑
            $$y_1''-y_1=\er^x,\quad y_2''-y_2=-x^2$$
            对$y_1$，查表得可以待定系数
            $$y_1=\alpha x\er^x$$
            直接计算可发现
            $$2\alpha\er^x=2\er^x$$
            从而对比系数得到可取$\alpha=1$，$y_1$的一个特解为
            $$y_1=x\er^x$$
            对$y_2$，查表得可以待定系数
            $$y_2=\beta x^2+\gamma x+\tau$$
            直接计算可发现
            $$-\beta x^2-\gamma x+2\beta-\tau=-x^2$$
            从而对比系数得到可取$\beta=1$、$\gamma=0$、$\tau=2$，$y_2$的一个特解为
            $$y_2=x^2+2$$
            进一步计算可发现原方程的一个特解为
            $$y_1+y_2=x\er^x+x^2+2$$
        }
    \end{enumerate}

    \item (习题9.5.4(6))写出以下微分变量为$x$的微分方程待定系数的特解形式：
    $$y''-2y'+2y=\er^x+x\cos x$$

    \sol{
        计算可知$y''-2y'+2y$的特征方程为$\lambda^2-2\lambda+2=0$，两根为$1\pm\ir$。由于右端分为$\er^x$与$x\cos x$两部分，我们考虑
        $$y_1''-2y_1'+2y_1=\er^x,\quad y_2''-2y_2'+2y_2=x\cos x$$
        对$y_1$，查表得可以待定系数
        $$y_1=\alpha\er^x$$
        对$y_2$，查表得可以待定系数
        $$y_2=(\beta x+\gamma)\cos x+(\tau x+\xi)\sin x$$
        进一步计算发现可设原方程的一个特解为
        $$y_1+y_2=\alpha\er^x+(\beta x+\gamma)\cos x+(\tau x+\xi)\sin x$$
    }
\end{enumerate}

\subsubsection{补充题解答}
\begin{enumerate}
    \item (习题9.1.2)指出下列关于$x$的表达式中的任意常数$C$是否相互独立：
    \begin{enumerate}
        \item[(3)]$y=C_1\er^{\lambda_1x}+C_2\er^{\lambda_2x}$，这里参数$\lambda_1\ne\lambda_2$。
        
        \sol{
            直接求导计算可知(求导代表对$x$)
            $$\frac{D(y,y')}{D(C_1,C_2)}=\er^{\lambda_1x}\lambda_2\er^{\lambda_2x}-\lambda_1\er^{\lambda_1x}\er^{\lambda_2x}=(\lambda_1-\lambda_2)\er^{(\lambda_1+\lambda_2)x}$$
            由条件其恒非零，从而$C_1$、$C_2$独立。
        }
        
        \item[(5)]$y=C_1x+C_2x^2+C_3(x+x^2)$
        
        \sol{
            提供两种方法：
            \begin{itemize}
                \item \textbf{求导计算}
                
                直接求导计算可知(求导代表对$x$)
                $$\frac{D(y,y',y'')}{D(C_1,C_2,C_3)}=\begin{vmatrix}x&x^2&x+x^2\\1&2x&1+2x\\0&2&2\end{vmatrix}=0$$
                从而$C_1$、$C_2$、$C_3$不独立。
                
                \item \textbf{定义验证}
                
                将原式改写为
                $$y=(C_1+C_3)x+(C_2+C_3)x^2$$
                因此取两个任意常数$D_1=C_1+C_3$、$D_2=C_2+C_3$可替代三个任意常数的作用，$C_1$、$C_2$、$C_3$不独立。
            \end{itemize}
        }
    \end{enumerate}

    \item (习题9.1.5(3))求以下初值问题的解：
    $$y'''(x)=x,\quad y(0)=a_0,\quad y'(0)=a_1,\quad y''(0)=a_2$$

    \sol{
        作不定积分可知($C_1$、$C_2$、$C_3$为任意常数)
        $$y''(x)=\frac{1}{2}x^2+C_1$$
        $$y'(x)=\frac{1}{6}x^3+C_1x+C_2$$
        $$y(x)=\frac{1}{24}x^4+\frac{1}{2}C_1x^2+C_2x+C_3$$
        代入初值条件即得
        $$\begin{cases}C_3=a_0\\C_2=a_1\\C_1=a_2\end{cases}$$
        从而
        $$y(x)=\frac{1}{24}x^4+\frac{1}{2}a_2x^2+a_1x+a_0$$
    }

    \item (习题9.2.1(2))求以下微分变量为$x$的微分方程的通积分：
    $$a(xy'+2y)=xyy'$$
    这里参数$a>0$。

    \sol{
        分为如下步骤：
        \begin{itemize}
            \item \textbf{通解计算}
            
            将方程整理为
            $$x(y-a)y'=2ya$$
            这是一个变量分离的方程，当$x\ne0$、$y\ne0$时方程两边同除以$2xy$可将其化为
            $$\frac{y-a}{y}\dr y=\frac{2a}{x}\dr x$$
            由于$\frac{y-a}{y}=1-\frac{a}{y}$，两边同积分得到($C$为任意常数)
            $$y-a\ln|y|=2a\ln|x|+C$$
            由$a>0$，两边同除以$a$后作$\er$的指数得
            $$\frac{\er^{y/a}}{y}=\er^{C/a}x^2$$
            考虑正负可去除两侧绝对值并将恒正的$\er^{C/a}$改写为新的\textbf{非零}任意常数$C$，再同乘$y$化简得到
            $$\er^{y/a}=Cyx^2$$
            这已经是合理的通解形式(注意此形式已经可以保证$x\ne0$、$y\ne0$)。

            \item \textbf{特解获取}
            
            由于过程中除以了$y$，额外讨论$y=0$的情况可发现$y=0$为一个特解。

            \note 若将通解改写为$yx^2=C\er^{y/a}$的形式，$C$即可取到0，且$C$取0对应特解。
        \end{itemize}
    }
    
    \item (习题9.2.4(1))求以下微分变量为$x$的微分方程的通积分：
    $$y'+\frac{y}{x}=y^3$$

    \sol{
        分为如下步骤：
        \begin{itemize}
            \item \textbf{通解计算}
            
            这符合\textbf{伯努利方程}的形式，从而在$y\ne0$时，两边同乘$\frac{1}{y^3}$并换元$u=\frac{1}{y^2}$，则由复合函数求导公式可知
            $$u'=-\frac{2y'}{y^3}$$
            从而方程化为
            $$-\frac{1}{2}u'+\frac{u}{x}=1$$
            也即
            $$u'-\frac{2}{x}u=-2$$
            这已经是一阶线性微分方程，可以直接求解。直接计算可得
            $$\int-\frac{2}{x}\dr x=-2\ln|x|+C$$
            由此可取一个积分因子
            $$\mu(x)=\er^{-2\ln|x|}=x^{-2}$$
            方程两边同乘$\mu(x)$\ (由方程形式可以保证$x\ne0$)有
            $$\big(x^{-2}u\big)'=-2x^{-2}$$
            作不定积分即得($C$为任意常数)
            $$u=x^2\cdot\bigg(\frac{2}{x}+C\bigg)=2x+Cx^2$$
            代入$u$的表达式最终得到
            $$1=y^2(2x+Cx^2)$$
            这已经是合理的通解形式(注意此形式已经可以保证$x\ne0$、$y\ne0$)。

            \item \textbf{特解获取}
            
            由于过程中除以了$y^3$，额外讨论$y=0$的情况可发现$y=0$为一个特解。
        \end{itemize}
    }

    \item (习题9.2.5)将下列微分变量为$x$的微分方程化为线性微分方程：
    \begin{enumerate}[(1)]
        \item $y'=\frac{x^2+y^2}{2y}$
        
        \sol{
            将方程改写为
            $$2yy'=x^2+y^2$$
            利用复合函数求导公式可知$(y^2)'=2yy'$，从而根据形式可看出令$u=y^2$方程将化为
            $$u'=x^2+u$$
            这已经是线性微分方程。
        }
        
        \item $y'=\frac{y}{x+y^2}$
        
        \sol{
            由于此方程为
            $$\frac{\dr y}{\dr x}=\frac{y}{x+y^2}$$
            利用反函数求导公式可知
            $$\frac{\dr x}{\dr y}=\frac{x+y^2}{y}$$
            也即
            $$\frac{\dr x}{\dr y}-\frac{x}{y}=y$$
            这已经是关于$x$、微分变量为$y$的线性微分方程。
        }
    \end{enumerate}

    \item (习题9.2.9)求在$[0,+\infty)$上连续的函数$\varphi$使得
    $$\forall x\in[0,+\infty),\quad\int_0^1\varphi(tx)\dr t=a\varphi(x)$$
    这里参数$a$满足$0<a<1$。

    \note 这里与教材上的写法略有不同，是为了避免出现\textbf{瑕积分}相关的讨论。

    \sol{
        为了能对左侧的积分进行操作，我们希望被积函数为$\varphi(t)$，换元得到$x>0$时
        $$\int_0^1\varphi(tx)\dr t=\frac{1}{x}\int_0^1\varphi(tx)\dr(tx)=\frac{1}{x}\int_0^x\varphi(t)\dr t$$
        由此，设
        $$\Phi(x)=\int_0^x\varphi(t)\dr t$$
        根据变上限积分的求导公式可知$\Phi'(x)=\varphi(x)$，由此方程改写为
        $$\frac{1}{x}\Phi(x)=a\Phi'(x)$$
        这已经是变量分离的方程，在$\Phi(x)\ne0$时两侧同除以$a\Phi(x)$得到(除以$a$是为了方便确定$\Phi(x)$)
        $$\frac{1}{ax}=\frac{\Phi(x)}{\Phi'(x)}$$
        同积分，由$x>0$即得($C$为任意常数)
        $$\ln|\Phi(x)|=\frac{1}{a}\ln x+C$$
        两边同作$\er$的指数，并考虑正负可去除绝对值，将恒正的$\er^C$改写为新的任意常数$C$，且要求$C$非零，得到
        $$\Phi(x)=Cx^{\frac{1}{a}}$$
        不过，进一步观察可发现$C=0$时恰好对应$\Phi(x)=0$，可验证其为特解，综合可将通解改写为($x>0$时)
        $$\Phi(x)=Cx^{\frac{1}{a}}$$
        求导即得($x>0$时)
        $$\varphi(x)=Cx^{\frac{1}{a}-1}$$
        由于$0<a<1$，此函数在0处也有定义，且可验证它在$[0,+\infty)$连续，符合要求。
    }

    \item (习题9.2.14(7))判断下列微分方程是否为全微分方程，若是则求出通积分：
    $$(y\er^x+2\er^x+y^2)\dr x+(\er^x+2xy)\dr y=0$$

    \sol{
        将此方程看作$P\dr x+Q\dr y=0$，直接计算可验证
        $$Q_x=P_y=\er^x+2y$$
        从而这是一个全微分方程。设$\dr u=P\dr x+Q\dr y$，先由$\partial_1u(x,y)=P$，计算$P$对$x$的不定积分，得到
        $$u(x,y)=y\er^x+2\er^x+y^2x+\varphi(y)$$
        代入$\partial_2u(x,y)=Q$可发现
        $$\er^x+2yx+\varphi'(y)=\er^x+2xy$$
        也即$\varphi'(y)=0$，从而($C$为任意常数)
        $$u(x,y)=y\er^x+2\er^x+y^2x+C$$
        最终通解为
        $$y\er^x+2\er^x+y^2x+C=0$$
    }

    \item (习题9.2.16(1))利用积分因子求解微分方程：
    $$x\dr x=(2xy^3\dr x+3x^2y^2\dr y)\sqrt{1+x^2}$$

    \sol{
        从题目中配凑的形式可猜测$2xy^3\dr x+3x^2y^2\dr y$是一个全微分，尝试计算发现令$v=x^2y^3$，则
        $$\dr v=2xy^3\dr x+3x^2y^2\dr y$$
        此时方程化为
        $$x\dr x=\sqrt{1+x^2}\dr v$$
        这已经是变量分离的方程，从而
        $$\mu(x,y)=\frac{1}{\sqrt{1+x^2}}$$
        即是一个\textbf{积分因子}。乘积分因子后方程化为
        $$\frac{x}{\sqrt{1+x^2}}\dr x=\dr v$$
        直接计算可知(换元$t=x^2$，之后成为幂函数复合一次函数)
        $$\int\frac{x}{\sqrt{1+x^2}}\dr x=\frac{1}{2}\int\frac{\dr t}{\sqrt{1+t}}=\sqrt{1+t}+C=\sqrt{1+x^2}+C$$
        由此通解为($C$为任意常数)
        $$\sqrt{1+x^2}+C=v$$
        代入得到
        $$\sqrt{1+x^2}+C=x^2y^3$$
    }

    \item (总练习题9.2)给定重力加速度$g$，考虑摆长$l$、摆锤质量$m$的单摆，设它在$t$时刻与平衡位置夹角为$\theta(t)$\ (含正负)，证明它的运动满足方程
    $$\theta''(t)=-\frac{g}{l}\sin(\theta(t))$$
    当$\theta$很小时，$\sin(\theta(t))$近似为$\theta(t)$，求出此时微分方程的通解。

    \sol{
        分为如下步骤：
        \begin{itemize}
            \item \textbf{方程构造}
            
            设切向受到向平衡位置的力为$F(t)$，远离平衡位置的加速度为$a(t)$，根据牛顿第二定律可知
            $$-ma(t)=F(t)$$
            再根据圆周运动的公式可知圆周运动的切向加速度即为角加速度乘半径，也即
            $$a(t)=l\theta''(t)$$
            根据受力分解可知$F(t)$是重力在切向的分量，从而
            $$F(t)=mg\sin(\theta(t))$$
            综合以上三式即得结果。

            \item \textbf{通解计算}
            
            此时的方程为
            $$\theta''(t)=-\frac{g}{l}\theta(t)$$
            这是一个不含未知数$t$的二阶微分方程，设$p=\theta'(t)$，利用教材中的分析可知(这里右侧将$\theta$也看作表达式)
            $$\theta''(t)=p\frac{\dr p}{\dr\theta}$$
            由此有
            $$p\frac{\dr p}{\dr\theta}=-\frac{g}{l}\theta$$
            这已经是变量分离的方程，也即
            $$p\dr p=-\frac{g}{l}\theta\dr\theta$$
            同积分得到($C$为任意常数)
            $$\frac{1}{2}p^2=-\frac{g}{2l}\theta^2+C$$
            将$2C$改写为新的任意常数$C$，方程重新写为(由物理意义$\frac{g}{l}>0$，从而为了有合理的解，至少有$C\ge0$)
            $$(\theta'(t))^2=-\frac{g}{l}\theta^2(t)+C$$
            这事实上仍是一个变量分离的微分方程。在$C-\frac{g}{l}\theta^2(t)>0$时将其写为
            $$\theta'(t)=\pm\sqrt{C-\frac{g}{l}\theta^2(t)}$$
            也即
            $$\frac{\dr\theta}{\pm\sqrt{C-\frac{g}{l}\theta^2}}=\dr t$$
            将$C$改写为新的任意常数$\frac{g}{l}C$，即可提出根号中的部分化简成
            $$\frac{\dr\theta}{\pm\sqrt{C-\theta^2}}=\sqrt{\frac{g}{l}}\dr t$$
            由$C\ge0$，利用上册教材3.2节的积分公式可知($C_0$为任意常数)
            $$\pm\arcsin\frac{\theta}{\sqrt{C}}=\sqrt{\frac{g}{l}}t+C_0$$
            两侧同取$\sin$，并将$\sqrt{C}$改写为新的任意常数$C_1$，利用左侧的正负，$C_1$可取任何实数，由此最终得到
            $$\theta(t)=C_1\sin\bigg(\sqrt{\frac{g}{l}}t+C_0\bigg)$$
            
            \note 这里由于总练习题9.2对应教材9.2节的部分，没有引入更复杂的知识。事实上利用教材9.5节的特征方程方法可以直接求解，得到($C_1$、$C_2$为任意常数)
            $$\theta(t)=C_1\cos\sqrt{\frac{g}{l}}t+C_2\sin\sqrt{\frac{g}{l}}t$$
            可验证两结果等价。
        \end{itemize}
    }

    \item (习题9.3.2)判断下列函数在相应的闭区域上是否满足Lipschitz条件：
    \begin{enumerate}
        \item[(2)]函数为$f(x,y)=\sqrt{y}$，区域为
        $$R=\{(x,y)\mid |x|\le a,\ 0<\delta\le y\le b\}$$
        这里参数$a>0$、$b>\delta>0$。

        \sol{
            直接计算可知在$R$上有
            $$|\partial_2f(x,y)|=\frac{1}{2\sqrt y}\le\frac{1}{2\delta}$$
            根据本讲义25.1.1习题5可知此函数满足Lipschitz常数为$\frac{1}{2\delta}$的Lipschitz条件。
        }

        \item[(4)]函数为$f(x,y)=xy^2$，区域为
        $$D=\{(x,y)\mid |x|\le a,\ y\in\rb\}$$
        这里参数$a>0$。

        \sol{
            若其满足Lipschitz条件，设Lipschitz常数为$L$。

            对$x_0\in[-a,a]$与$y\in\rb$，直接计算可知
            $$|f(x_0,y_1)-f(x_0,y_2)|=x_0|y_1^2-y_2^2|=x_0|y_1+y_2||y_1-y_2|$$
            取$x_0=a$、$y_1=\frac{L+1}{a}$、$y_2=0$，即有
            $$|f(x_0,y_1)-f(x_0,y_2)|=(L+1)|y_1-y_2|>L|y_1-y_2|$$
            矛盾，从而其不满足Lipschitz条件。
        }
    \end{enumerate}

    \item (习题9.3.5)求初值问题
    $$\frac{\dr y}{\dr x}=x-y^2,\quad y(0)=0$$
    皮卡序列的前三项$y_0(x)$、$y_1(x)$、$y_2(x)$。

    \sol{
        根据定义可知
        $$y_0(x)=0$$
        $$y_1(x)=\int_0^x(t-y_0^2(t))\dr t=\frac{1}{2}x^2$$
        $$y_2(x)=\int_0^x(t-y_1^2(t))\dr t=\frac{1}{2}x^2-\frac{1}{20}x^5$$
    }

    \item (习题9.3.6(3))利用解的存在唯一性定理计算以下初值问题的解的存在区间：
    $$y'=x+\er^y,\quad y(1)=0$$
    区域为
    $$R=\{(x,y)\mid|x-1|\le1,\ |y|\le1\}$$
    
    \sol{
        分为如下步骤：
        \begin{itemize}
            \item \textbf{估计$L$}
            
            设$f(x,y)=x+\er^y$，直接计算可知在$R$上有
            $$|\partial_2f(x,y)|=\er^y\le\er$$
            根据本讲义25.1.1习题5可知此函数满足Lipschitz常数为$\er$的Lipschitz条件。
            
            \item \textbf{估计$M$}
            
            由于区域中$x\ge0$，直接放缩有
            $$|f(x,y)|=x+\er^y\le 2+\er$$
            且当$x=2$、$y=1$时可取到，从而$|f(x,y)|$在$R$上最大值为$2+\er$，即$M=2+\er$。

            \item \textbf{估计$h$}
            
            本题中$a=b=1$，因此
            $$h=\min\bigg\{1,\frac{1}{2+\er}\bigg\}=\frac{1}{2+\er}$$
            从而估计得到的存在区间为
            $$\bigg[1-\frac{1}{2+\er},1+\frac{1}{2+\er}\bigg]$$
        \end{itemize}
    }
\end{enumerate}

\subsection{复习题}
\subsubsection{重积分}
\noindent\textbf{二重积分}
\begin{framed}
\begin{exercise}\label{sin 2 and 4 to 2pi}
    计算
    $$\iint_D(|x|+y+xy)^2\dr x\dr y$$
    其中
    $$D=\{(x,y)\mid x^2+y^2\le a^2\}$$
    这里参数$a>0$。
\end{exercise}
\sol{
    分为如下步骤：
    \begin{itemize}
        \item \textbf{反射对称化简}
        
        将被积函数\textbf{展开}得到其为
        $$x^2+y^2+x^2y^2+2|x|y+2xy^2+2|x|xy$$
        由于区域$D$满足$(x,y)\in D$当且仅当$(-x,y)\in D$，且
        $$2(-x)y^2+2|-x|(-x)y=-(2xy^2+2|x|xy)$$
        利用函数折半可发现
        $$\iint_D(2xy^2+2|x|xy)\dr x\dr y=0$$
        同理，由$(x,y)\in D$当且仅当$(x,-y)\in D$可得
        $$\iint_D2|x|y\dr x\dr y=0$$
        综合以上可将积分化为
        $$\iint_D(x^2+y^2+x^2y^2)\dr x\dr y$$

        \item \textbf{换元与累次积分化}
        
        由区域的旋转对称性，考虑极坐标换元$x=r\cos\theta$、$y=r\sin\theta$。

        此时，原本的被积函数可写为$r^2+r^4\cos^2\theta\sin^2\theta$，且计算Jacobi行列式得$\dr x\dr y=r\dr r\dr\theta$。区域的条件可写为
        $$r^2\le a^2$$
        结合极坐标自然范围即得
        $$0\le r\le a,\quad 0\le\theta\le2\pi$$
        此范围已经可以直接写出累次积分，从而得到积分
        $$\int_0^a\dr r\int_0^{2\pi}(r^3+r^5\cos^2\theta\sin^2\theta)\dr\theta$$

        \item \textbf{积分计算}
        
        我们先来计算
        $$I_0=\int_0^{2\pi}\cos^2\theta\sin^2\theta\dr\theta$$
        为了利用上册教材3.4节例2的熟知结论，我们先将积分改写为
        $$I_0=\int_0^{2\pi}(1-\sin^2\theta)\sin^2\theta\dr\theta=\int_0^{2\pi}\sin^2\theta\dr\theta-\int_0^{2\pi}\sin^4\theta\dr\theta$$
        接下来可利用对称性化简，将$\theta$换元为$2\pi-\theta$，考虑函数折半可发现
        $$I_0=2\int_0^\pi\sin^2\theta\dr\theta-2\int_0^\pi\sin^4\theta\dr\theta$$
        将$\theta$换元为$\pi-\theta$再次折半，得到
        $$I_0=4\int_0^{\pi/2}\sin^2\theta\dr\theta-4\int_0^{\pi/2}\sin^4\theta\dr\theta$$
        这时就可以直接用结论得到
        $$I_0=4\cdot\frac{1}{2}\cdot\frac{\pi}{2}-4\cdot\frac{3}{4}\cdot\frac{1}{2}\cdot\frac{\pi}{2}=\frac{\pi}{4}$$
        由此可进一步计算出积分为
        $$\int_0^a\bigg(2\pi r^3+\frac{\pi}{4}r^5\bigg)\dr r=\frac{\pi}{2}a^4+\frac{\pi}{24}a^6$$
    \end{itemize}
}
\end{framed}

\note 反射对称化简的基本定理可见本讲义20.3.2\ (给出了三重积分版本，二重积分类似)，注意交换$x$、$y$也是一种反射对称。务必先\textbf{展开}被积函数再进行化简。

\note 反射对称性的观察基本来自将约束中的$(x,y,z)$\textbf{代换}为$(-x,y,z)$或$(y,x,z)$等，观察是否能得到相同的一组约束，能得到即意味着有对称性。

\note \textbf{非常建议记忆}上册教材3.4节例2的结论，很多化出的三角积分都可以利用对称性向其配凑。

\begin{framed}
\begin{exercise}
    计算
    $$\iint_D\bigg(\frac{3x^2\sin y}{y}+2\er^{x^2}\bigg)\dr x\dr y$$
    其中$D$是由曲线$y=x$、$y=x^3$在$x\ge0$部分围成的区域。
\end{exercise}
\sol{
    分为下步骤：
    \begin{itemize}
        \item \textbf{区域描述}
        
        利用``围成''需要区域有界，我们分别讨论每个曲线对应的约束(通常也即将曲线表达式里的等号变为大于等于或小于等于)。

        由于$y$必然需要从两边被约束(否则可以向一侧无限放大)，区域只可能为$x\le y\le x^3$或$x^3\le y\le x$\ (结合$x\ge0$的要求，事实上左边均需加上大于等于0)。对于第一种，$x\to+\infty$时符合要求，不满足有界性，因此最终得到的范围只能为
        $$0\le x^3\le y\le x$$
        区域$D$即为满足以上约束的$(x,y)$结合。

        \note 本题直接画出函数图像事实上更容易确定上述形式。

        \item \textbf{累次积分化}
        
        由于函数$\frac{3x^2\sin y}{y}$对$y$不存在初等原函数，必须对$x$先积分，而$\er^{x^2}$对$x$不存在初等原函数，必须对$y$积分。记
        $$I_1=\iint_D\frac{3x^2\sin y}{y}\dr x\dr y,\quad I_2=\iint_D2\er^{x^2}\dr x\dr y$$
        对$I_1$，为将其写为先$x$后$y$的累次积分应将区域写成
        $$\alpha\le y\le\beta,\quad\phi_1(y)\le x\le\phi_2(y)$$
        当$x$可任取时，$y$最小为0，在$x=0$时取到；最大必然在$x^3=x=y$时取到，可发现$x=1$时$y$有最大值，为1。给定$y$后，改写不等式可发现
        $$0\le y\le1,\quad y\le x\le\sqrt[3]{y}$$
        从而可写出累次积分
        $$I_1=\int_0^1\dr y\int_y^{\sqrt[3]{y}}\frac{3x^2\sin y}{y}\dr x$$
        对$I_1$，为将其写为先$y$后$x$的累次积分应将区域写成
        $$\alpha\le x\le\beta,\quad\phi_1(x)\le y\le\phi_2(x)$$
        当$y$可任取时，$x$最小为0，在$y=0$时取到；最大必然在$x^3=x$时取到，可发现为1。给定$x$后，添加对$y$的不等式即得范围
        $$0\le x\le1,\quad x^3\le y\le x$$
        从而可写出累次积分
        $$I_2=\int_0^1\dr x\int_{x^3}^x2\er^{x^2}\dr y$$

        \item \textbf{$I_1$积分计算}
        
        直接对$x$积分可知
        $$I_1=\int_0^1(y-y^3)\frac{\sin y}{y}=\int_0^1(\sin y-y^2\sin y)\dr y$$
        利用分部积分可知
        $$\int_0^1y^2\sin y\dr y=-\int_0^1y^2\dr\cos y=-\cos 1+\int_0^12y\cos y\dr y$$
        再利用分部积分得
        $$\int_0^12y\cos y\dr y=\int_0^12y\dr\sin y=2\sin 1-2\int_0^1\sin y\dr y$$
        综合可得结果
        $$I_1=\cos 1-2\sin 1+3\int_0^1\sin y\dr y=3-2\cos1-2\sin1$$

        \item \textbf{$I_2$积分计算}
        
        直接对$y$积分可知
        $$I_2=\int_0^12(x-x^3)\er^{x^2}\dr x$$
        换元$t=x^2$可得
        $$I_2=\int_0^1(1-t)\er^t\dr t$$
        利用分部积分可得
        $$\int_0^1t\er^t\dr t=\int_0^1t\dr\er^t=\er-\int_0^1\er^t\dr t$$
        由此即知
        $$I_2=2\int_0^1\er^t\dr t-\er=\er-2$$

        \item \textbf{综合结果}

        综合以上两个部分可知积分最终为
        $$1-2\sin1-2\cos1+\er$$
    \end{itemize}
}
\end{framed}

\note 累次积分往往应最优先计算\textbf{容易计算的部分}，其次是\textbf{容易确定范围的部分}，如$x^3\le y\le x$实际上自然有$y$被$x$约束，若被积函数无特殊性质一般考虑先对$y$积分再对$x$积分。

\note 为了确定哪些部分可以计算，往往需要知道一些\textbf{不存在初等原函数}的形式，如$\frac{\sin x}{x}$、$\cos x^2$、$\er^{x^2}$等。

\note 大家务必复习上学期的基本\textbf{定积分计算技巧}。

\begin{framed}
\begin{exercise}
    计算
    $$\iint_D\bigg(\frac{1}{2}x-y\bigg)\dr x\dr y$$
    其中$D$是由直线$y=0$、$y=2$、$y=x$、$y=x+2$围成的区域。
\end{exercise}
\sol{
    分为如下步骤：
    \begin{itemize}
        \item \textbf{区域描述}
        
        根据围成的区域必然以每条直线为边界，由于$0<2$恒成立，$y=0$与$y=2$对应的约束必然为(若$y\le0$则$y=2$不会为边界，若$y\ge2$则$y=0$不会为边界)
        $$0\le y\le2$$
        同理，由于$x<x+2$恒成立，两者对应的约束必然为
        $$x\le y\le x+2$$
        区域$D$即为满足以上约束的$(x,y)$结合。

        \item \textbf{累次积分化}
        
        由于区域描述中给出了$y$的最大范围，我们只需用$y$约束$x$即可将区域写为累次积分要求的形式。改写约束可发现
        $$y-2\le x\le y$$
        从而结合$0\le y\le2$可写出累次积分
        $$\int_0^2\dr y\int_{y-2}^y\bigg(\frac{1}{2}x-y\bigg)\dr x$$

        \item \textbf{积分计算}
        
        直接对$x$积分可将累次积分化为
        $$\int_0^2\bigg(\frac{1}{4}(y^2-(y-2)^2)-2y\bigg)\dr y=\int_0^2(-y-1)\dr y=-4$$
    \end{itemize}
}
\end{framed}

\note 区域由``围成''给出时往往可以通过\textbf{有界性}与\textbf{以每个曲线/曲面为边界}直接从\textbf{纯代数}方法确定出区域。

\note 不过，二维情况这样确定区域未必比直接画图方便，且曲线的情况可能更加复杂，仍然可以考虑\textbf{画图确定区域}。三维情况则\textbf{非常不推荐画图}。

\note 当原本的区域描述方便计算时，一般\textbf{不建议考虑换元}，因为Jacobi行列式未必容易计算。

\begin{framed}
\begin{exercise}
    计算
    $$\iint_D\sqrt{|y-x^2|}\dr x\dr y$$
    其中
    $$D=\{(x,y)\mid x\in[-1,1],\ y\in[0,2]\}$$
\end{exercise}
\sol{
    分为如下步骤：
    \begin{itemize}
        \item \textbf{累次积分化}
        
        由区域描述可直接写出对应的累次积分(这里对$y$形式稍简单，考虑先对$y$积分)
        $$\int_{-1}^1\dr x\int_0^2\sqrt{|y-x^2|}\dr y$$

        \item \textbf{积分计算}
        
        给定$x$后，由$y\in[0,2]$当$y\in[0,x^2]$时积分中为$\sqrt{x^2-y}$，否则为$\sqrt{y-x^2}$，从而积分可写为
        $$\int_{-1}^1\dr x\bigg(\int_0^{x^2}\sqrt{x^2-y}\dr y+\int_{x^2}^2\sqrt{y-x^2}\dr y\bigg)$$
        两个被积函数都是幂函数复合一次函数，可直接写出对$y$的原函数
        $$-\frac{2}{3}(x^2-y)^{3/2}+C,\quad\frac{2}{3}(y-x^2)^{3/2}+C$$
        由此可将积分化为(注意$(x^2)^{3/2}$是$|x|^3$)
        $$\frac{2}{3}\int_{-1}^1(|x|^3+(2-x^2)^{3/2})\dr x$$
        由于被积函数为偶函数，考虑$x$换元为$-x$可发现
        $$\int_0^1(|x|^3+(2-x^2)^{3/2})\dr x=\int_{-1}^0(|x|^3+(2-x^2)^{3/2})\dr x$$
        由此原积分进一步化为其在$[0,1]$积分的两倍，去除绝对值写为
        $$\frac{4}{3}\int_0^1(x^3+(2-x^2)^{3/2})\dr x=\frac{1}{3}+\frac{4}{3}\int_0^1(2-x^2)^{3/2}\dr x$$
        换元$x=\sqrt{2}\sin\theta$，由$\sin\frac{\pi}{4}=\frac{\sqrt2}{2}$可知
        $$\int_0^1(2-x^2)^{3/2}\dr x=4\int_0^{\pi/4}\cos^4\theta\dr\theta$$
        也即原积分为
        $$\frac{1}{3}+\frac{16}{3}\int_0^{\pi/4}\cos^4\theta\dr\theta$$
        利用$2\cos^2\theta=1+\cos(2\theta)$并换元$\phi=2\theta$可得其为
        $$\frac{1}{3}+\frac{4}{3}\int_0^{\pi/4}(1+\cos(2\theta))^2\dr\theta=\frac{1}{3}+\frac{2}{3}\int_0^{\pi/2}(1+\cos\phi)^2\dr\phi$$
        展开有(第三部分可换元$t=\frac{\pi}{2}-\theta$后使用上册教材3.4节例2)
        $$\int_0^{\pi/2}1\dr\phi=\frac{\pi}{2},\quad\int_0^{\pi/2}2\cos\phi\dr\phi=2,\quad\int_0^{\pi/2}\cos^2\phi\dr\phi=\frac{\pi}{4}$$
        从而综合得到结果为
        $$\frac{1}{2}\pi+\frac{5}{3}$$
    \end{itemize}
}
\end{framed}

\note 对于带有绝对值或$\max$、$\min$的积分，要学会写为定积分后再\textbf{讨论范围}积分。

\begin{framed}
\begin{exercise}
    计算
    $$\iint_Dr^2\sin\theta\sqrt{1-r^2\cos(2\theta)}\dr r\dr\theta$$
    其中
    $$D=\bigg\{(r,\theta)\ \bigg|\ 0\le r\le\sec\theta,\ 0\le\theta\le\frac{\pi}{4}\bigg\}$$
\end{exercise}
\sol{
    分为如下步骤：
    \begin{itemize}
        \item \textbf{换元回直角坐标}
        
        由于在极坐标下强行计算非常困难，我们尝试换元回直角坐标，也即令
        $$x=r\cos\theta,\quad y=r\sin\theta$$
        由于$D$对$r$、$\theta$的约束\textbf{包含在极坐标自然范围内}，换元是合理的。

        先处理区域。由$0\le\theta\le\frac{\pi}{4}$可知其在第一象限中，且
        $$\arctan\frac{y}{x}\le\frac{\pi}{4}$$
        利用单调性即得范围
        $$0\le y\le x$$
        此时，已经可以保证$\cos\theta>0$，从而第一个约束同乘$\cos\theta$得到$0\le r\cos\theta\le1$，也即最终得到范围
        $$0\le y\le x\le1$$
        将上述条件确定的$(x,y)$区域记为$D_0$。

        计算Jacobi行列式可得
        $$\dr x\dr y=r\dr r\dr\theta$$
        从而直角坐标下的被积函数应为
        $$r\sin\theta\sqrt{1-r^2\cos(2\theta)}$$
        为了配凑$r^2$，直接展开$\cos(2\theta)$为$\cos^2\theta-\sin^2\theta$，最终得到被积函数
        $$y\sqrt{1-x^2+y^2}$$

        \note 即使用其他方式展开$\cos(2\theta)$，被积函数仍然是相同的形式。

        积分最终化为
        $$\iint_{D_0}y\sqrt{1-x^2+y^2}\dr x\dr y$$

        \item \textbf{累次积分化}
        
        由于对$y$的形式相对简单，我们先对$y$积分，这时需将约束写为
        $$\alpha\le x\le\beta,\quad \phi_1(x)\le y\le\phi_2(x)$$
        根据约束的条件，$y$任取时$x$的范围是$[0,1]$，而给定$x$后$y$的范围即$[0,x]$，因此可写出累次积分
        $$\int_0^1\dr x\int_0^xy\sqrt{1-x^2+y^2}\dr y$$

        \item \textbf{积分计算}
        
        由于$y\ge0$，可换元$t=y^2$将积分化为
        $$\int_0^1\dr x\frac{1}{2}\int_0^{x^2}\sqrt{1-x^2+t}\dr t$$
        此对$t$为幂函数符合一次函数，可直接写出原函数计算得结果为
        $$\frac{1}{3}\int_0^1(1-(1-x^2)^{3/2})\dr x=\frac{1}{3}-\frac{1}{3}\int_0^1(1-x^2)^{3/2}\dr x$$
        换元$x=\cos t$可将上式写为
        $$\frac{1}{3}-\frac{1}{3}\int_0^{\pi/2}\sin^4t\dr t$$
        由此利用上册教材3.4节例2结论可知结果为
        $$\frac{1}{3}-\frac{1}{3}\cdot\frac{3}{4}\cdot\frac{1}{2}\cdot\frac{\pi}{2}=\frac{1}{3}-\frac{\pi}{16}$$

        \note 用$\cos t$而非$\sin t$换元就是为了使得之后能直接使用结论。
    \end{itemize}
}
\end{framed}

\note 直角坐标与极坐标、柱坐标、球坐标的\textbf{相互转化}需要熟悉，核心是配凑出相应的形式。

\note 本题即使最开始的字母不为$r$与$\theta$也应看出是极坐标，因为有一个变量作为角度出现，另一个变量则被角度限制。

\

\noindent\textbf{三重积分}

\begin{framed}
\begin{exercise}
    计算
    $$\int_{-1}^1\dr x\int_0^{\sqrt{1-x^2}}\dr y\int_1^{1+\sqrt{1-x^2-y^2}}\frac{\dr z}{\sqrt{x^2+y^2+z^2}}$$
\end{exercise}
\sol{
    分为如下步骤：
    \begin{itemize}
        \item \textbf{区域描述}
    
        根据累次积分，区域可描述为满足
        $$-1\le x\le1,\quad 0\le y\le\sqrt{1-x^2},\quad1\le z\le1+\sqrt{1-x^2-y^2}$$
        的$(x,y,z)$集合。

        第一个约束可保证$1-x^2\ge0$，而第二个约束包含$y\ge0$，由此可同平方可将前两个约束改写为
        $$x^2+y^2\le1,\quad y\ge0$$
        同理，已经保证$1-x^2-y^2\ge0$后第三个约束右侧可同平方改写为$x^2+y^2+(z-1)^2\le1$，综合以上得到了约束
        $$y\ge0,\quad z-1\ge0,\quad x^2+y^2+(z-1)^2\le1$$

        \item \textbf{换元与累次积分化}
        
        由于区域是以$(0,0,1)$为中心的球体的一部分，考虑平移的球坐标换元
        $$x=\rho\sin\varphi\cos\theta,\quad y=\rho\sin\varphi\sin\theta,\quad z=\rho\cos\varphi+1$$
        此时被积函数为$\frac{1}{\sqrt{\rho^2+2\rho\cos\varphi+1}}$，且直接计算Jacobi行列式可知(由于加常数不影响求导，Jacobi行列式应与球坐标完全相同)
        $$\dr x\dr y\dr z=\rho^2\sin\varphi\dr\rho\dr\varphi\dr\theta$$
        下面分析换元后的约束条件，这里采用\textbf{几何}的方式，也可从代数的分析得到相同的结果。

        首先，约束$x^2+y^2+(z-1)^2\le1$结合自然范围即$0\le\rho\le1$，对应球体内部。接下来，$z\ge1$对应球体的上半部分，也即与$z$轴正半轴夹角不超过$\frac{\pi}{2}$，$0\le\varphi\le\frac{\pi}{2}$。最后，$y\ge0$意味着$z$轴正方向看去，从$x$轴正向逆时针旋转并未超过半圈，也即$0\le\theta\le\pi$。

        综合以上，我们得到了区域的描述
        $$0\le\rho\le1,\quad0\le\varphi\le\frac{\pi}{2},\quad0\le\theta\le\pi$$
        由于被积函数与$\theta$无关，且对$\varphi$的形式更简单，考虑$\theta$、$\varphi$、$\rho$的积分顺序，根据此形式写出对应的累次积分
        $$\int_0^1\dr\rho\int_0^{\pi/2}\dr\varphi\int_0^\pi\frac{\rho^2\sin\varphi}{\sqrt{\rho^2+2\rho\cos\varphi+1}}\dr\theta$$

        \item \textbf{积分计算}
        
        对$\theta$积分可得
        $$\pi\int_0^1\dr\rho\int_0^{\pi/2}\frac{\rho^2\sin\varphi\dr\varphi}{\sqrt{\rho^2+2\rho\cos\varphi+1}}$$
        换元$t=\cos\varphi$得到
        $$\pi\int_0^1\dr\rho\int_0^1\frac{\rho^2\dr t}{\sqrt{\rho^2+2\rho t+1}}$$
        这已经是$t$的幂函数复合一次函数的形式，可直接写出原函数得到积分
        $$\pi\int_0^1\rho\big(\sqrt{\rho^2+2\rho+1}-\sqrt{\rho^2+1}\big)\dr\rho=\pi\int_0^1\rho\big(\rho+1-\sqrt{\rho^2+1}\big)\dr\rho$$
        可以先计算出$\rho$与$\rho^2$的积分将其化为
        $$\frac{5\pi}{6}-\pi\int_0^1\rho\sqrt{\rho^2+1}\dr\rho$$
        换元$s=\rho^2$即可得到其为
        $$\frac{5\pi}{6}-\frac{\pi}{2}\int_0^1\sqrt{s+1}\dr\rho=\frac{5\pi}{6}-\frac{\pi}{2}\cdot\frac{2}{3}(2^{3/2}-1)=\frac{7-4\sqrt2}{6}\pi$$
    \end{itemize}

    \note 本题直接进行球坐标换元也可以计算积分，大家可以自行尝试。不过，这种区域是非原点为中心的球面相关还是更建议大家用平移的球坐标。
}
\end{framed}

\note 换元为球坐标时\textbf{必须掌握自然范围}，也建议\textbf{熟悉参数的几何意义}以简化区域分析。

\begin{framed}
\begin{exercise}
    计算
    $$\iiint_\Omega(y^2+z^2)\dr V$$
    其中
    $$\Omega=\{(x,y,z)\mid0\le z\le x^2+y^2\le1\}$$
\end{exercise}
\sol{
    分为如下步骤：
    \begin{itemize}
        \item \textbf{换元与累次积分化}
        
        考虑到区域的旋转对称性，进行\textbf{柱坐标}换元
        $$x=r\cos\theta,\quad y=r\sin\theta,\quad z=z$$
        这样约束条件即化为了
        $$0\le z\le r^2\le1$$
        此范围中$z$自然被$r$限制，而$\theta$无限制，由此考虑先对$\theta$、再对$z$、最后对$r$积分。在$z$可任取时结合自然范围知$r$的范围
        $$0\le r\le1$$
        给定$r$、任取$\theta$时$z$的范围即
        $$0\le z\le r^2$$
        而$\theta$的范围即为自然范围$0\le\theta\le2\pi$。

        将上述条件确定的$(r,\theta,z)$区域记为$\Omega_0$，被积函数可写为
        $$z^2+r^2\sin^2\theta$$
        直接计算Jacobi行列式可知
        $$\dr x\dr y\dr z=r\dr r\dr\theta\dr z$$
        于是积分为
        $$\int_0^1\dr r\int_0^{r^2}\dr z\int_0^{2\pi}(z^2+r^2\sin^2\theta)r\dr\theta$$

        \note 也可利用反射对称性将被积函数先化为$\frac{1}{2}(x^2+y^2)+z^2$以使换元后的形式更加简单。

        \item \textbf{积分计算}
        
        类似\exref{sin 2 and 4 to 2pi}的处理可知$\sin^2\theta$在0到$2\pi$的积分为$\pi$，从而对$\theta$积分后为
        $$\int_0^1\dr r\int_0^{r^2}(2\pi z^2r+\pi r^3)\dr z$$
        对$z$积分即得到结果为
        $$\int_0^1\bigg(\frac{2\pi}{3}r^7+\pi r^5\bigg)\dr r=\frac{\pi}{4}$$
    \end{itemize}
}
\end{framed}

\note 若区域的描述中既有$x^2+y^2$相关又有$x^2+y^2+z^2$相关，用球坐标计算往往\textbf{区域确定更困难而计算更简单}，用柱坐标则往往\textbf{区域确定更简单而计算更困难}。若两者熟练度均较高，可以考虑\textbf{先使用球坐标}，若发现区域难以描述再改用柱坐标。

\begin{framed}
\begin{exercise}
    计算
    $$\iiint_\Omega\frac{1}{(x+y+z+1)^2}\dr V$$
    其中$\Omega$是由平面$x=0$、$y=0$、$z=0$、$x+y+z=1$围成的四面体。
\end{exercise}
\sol{
    我们介绍三种方法。首先是严谨的\textbf{直接代数计算方法}：
    \begin{itemize}
        \item \textbf{区域描述}
        
        根据``围成''的有界性要求讨论可以发现，四面体必然为满足
        $$x\ge0,\quad y\ge0,\quad z\ge0,\quad x+y+z\le1$$
        的$(x,y,z)$集合。

        \item \textbf{累次积分化}
        
        由于对称性，我们可以任取顺序，不妨考虑先$x$、后$y$、最后$z$，此时需要将集合描述成
        $$\alpha\le z\le\beta,\quad\phi_1(z)\le y\le\phi_2(z),\quad\psi_1(y,z)\le x\le\psi_2(y,z)$$
        当$x$、$y$可任取时，$z$的最大范围即$[0,1]$，最大值在$x=y=0$时取到；当$z$给定、$x$可任取时，$y$的最大范围即$[0,1-z]$，最大值在$x=0$时取到；当$y$、$z$给定时，$x$的范围即$[0,1-z-y]$。综合以上讨论可写出累次积分
        $$\int_0^1\dr z\int_0^{1-z}\dr y\int_0^{1-y-z}\frac{1}{(1+x+y+z)^2}\dr x$$

        \item \textbf{积分计算}
        
        直接写出原函数可知对$x$积分后为
        $$\int_0^1\dr z\int_0^{1-z}\bigg(\frac{1}{1+y+z}-\frac{1}{2}\bigg)\dr y$$
        同样直接写出原函数可知对$y$积分后为
        $$\int_0^1\bigg(\ln2-\ln(1+z)-\frac{1-z}{2}\bigg)\dr z$$
        这可以进一步化简为
        $$\ln 2-\frac{1}{4}-\int_0^1\ln(1+z)\dr z$$
        进一步计算，直接分部积分得到
        $$\int_0^1\ln(1+z)\dr z=\ln2-\int_0^1\frac{z\dr z}{1+z}=\ln2-1+\int_0^1\frac{\dr z}{1+z}=2\ln2-1$$
        从而最终计算结果为
        $$\frac{3}{4}-\ln2$$
    \end{itemize}
    
    第二种是\textbf{换元计算方法}：
    \begin{itemize}
        \item \textbf{换元}
        
        由于被积函数为$x+y+z$的函数，希望换元使得$x+y+z$看作整体。

        考虑$u=x+y+z$、$v=y$、$w=z$，直接计算Jacobi行列式可知
        $$\dr x\dr y\dr z=\dr u\dr v\dr w$$
        且被积函数成为$\frac{1}{(1+u)^2}$，约束化为
        $$u-v-w\ge0,\quad v\ge0,\quad w\ge0,\quad u\le1$$
        将上方约束确定的$(u,v,w)$区域记为$\Omega_0$，积分即化为
        $$\iiint_{\Omega_0}\frac{1}{(1+u)^2}\dr u\dr v\dr w$$
        
        \item \textbf{累次积分化}
        
        我们希望最后对$u$积分，先计算好处理的部分，因此选择$v$、$w$、$u$的积分顺序，此时需要将集合描述成
        $$\alpha\le u\le\beta,\quad\phi_1(u)\le w\le\phi_2(u),\quad\psi_1(u,w)\le v\le\psi_2(u,w)$$
        由于$v+w\le u\le1$，当$v$、$w$可任取时，$u$最小在$v=w=0$时取到，为0，最大为1；当$u$给定、$v$可任取时，$w$受到$0\le w\le u-v$约束，从而最小为0，最大在$v=0$时取到，为$u$；当$u$、$w$给定时，$v$的约束即$0\le v\le u-w$，最小为0、最大为$u-w$。
        
        综合以上讨论可写出累次积分
        $$\int_0^1\dr u\int_0^u\dr w\int_0^{u-w}\frac{1}{(1+u)^2}\dr v$$

        \item \textbf{积分计算}
        
        对$v$积分可以得到
        $$\int_0^1\dr u\int_0^u\frac{u-w}{(1+u)^2}\dr w$$
        再对$w$积分得到
        $$\frac{1}{2}\int_0^1\frac{u^2}{(1+u)^2}\dr u$$
        换元$t=u+1$得到其为
        $$\frac{1}{2}\int_1^2\frac{t^2-2t+1}{t^2}\dr t=\frac{1}{2}\bigg(1-2\ln2+1-\frac{1}{2}\bigg)=\frac{3}{4}-\ln2$$  
    \end{itemize}

    最后是方便计算的\textbf{几何方法}：
    \begin{itemize}
        \item \textbf{几何角度研究}
        
        由于被积函数当$x+y+z$恒定时恒定，我们考虑在垂直于$x+y+z$的方向累计。

        计算可发现原点与平面$x+y+z=t$的距离为$\frac{t}{\sqrt3}$，由此，给定到原点的距离$r$，考虑原区域与每个平面
        $$x+y+z=\sqrt3r$$
        相交的部分。这是由$(\sqrt3r,0,0)$、$(0,\sqrt3r,0)$、$(0,0,\sqrt3r)$连成的三角形，从而是边长$\sqrt6r$的等边三角形，面积为
        $$\frac{3\sqrt3}{2}r^2$$
        在此三角形上，函数值恒为
        $$\frac{1}{(1+\sqrt3r)^2}$$
        由此可得到积分
        $$\int_0^{1/\sqrt3}\frac{1}{(1+\sqrt3r)^2}\frac{3\sqrt3}{2}r^2\dr r$$
        
        \item \textbf{积分计算}
        
        换元$t=\sqrt3r$可以将积分写为
        $$\frac{1}{2}\int_0^1\frac{t^2}{(1+t)^2}\dr t$$
        这已经与上一种方法中的积分完全相同，可相同算出结果。
    \end{itemize}

    \note 几何方法用换元的严谨化事实上需要一定的\textbf{线性代数}知识，也即对区域进行正交变换，旋转$x+y+z=0$平面为$Oxy$平面。
}
\end{framed}

\note 需要熟练掌握通过\textbf{代数方法}确定区域的技巧，更多例子可见本讲义第二十章。是否需要换元以简化过程可以视实际复杂程度而定。

\

\noindent\textbf{间接计算}
\begin{framed}
\begin{exercise}
    计算曲线$x^3+y^3=xy$围成区域的面积。
\end{exercise}
\sol{
    根据重积分几何意义，设区域为$D$，其面积即为
    $$\iint_D1\dr\sigma$$
    计算分为如下步骤：
    \begin{itemize}
        \item \textbf{区域描述}
        
        本题\textbf{几乎无法不画图确定区域}：积分区域并不是仅仅$x^3+y^3\le xy$或$x^3+y^3\ge xy$，而是
        $$x^3+y^3\le xy,\quad x\ge0,\quad y\ge0$$
        本质上，这是由于这条曲线并不是简单闭曲线，它与自身有交点$(0,0)$，因此即使考虑$x^3+y^3\le xy$，在第一象限外的部分仍然是无界的，只在第一象限表示一个区域。

        \note 如果没有想到极坐标代换的方法，本题的作图也并不容易。

        \item \textbf{换元与累次积分化}
        
        考虑到直接计算无法解出$x$、$y$的关系，本题应考虑换元。由于方程两边分别是\textbf{齐次}三次、二次，若换元为极坐标可以将$r$用$\theta$表示，由此即想到换元极坐标
        $$x=r\cos\theta,\quad y=r\sin\theta$$
        此时原本的被积函数仍为1，且计算Jacobi行列式得$\dr x\dr y=r\dr r\dr\theta$。由第一象限可知
        $$0\le\theta\le\frac{\pi}{2}$$
        而曲线的约束即
        $$r^3(\cos^3\theta+\sin^3\theta)\le r^2\cos\theta\sin\theta$$
        由于第一象限中$\cos^3\theta+\sin^3\theta>0$，忽略$r=0$的特殊点即写出了约束
        $$0\le r\le\frac{\cos\theta\sin\theta}{\cos^3\theta+\sin^3\theta}$$
        这样关于$\theta$与$r$的约束已经可以写出累次积分
        $$\int_0^{\pi/2}\dr\theta\int_0^{\cos\theta\sin\theta/(\cos^3\theta+\sin^3\theta)}r\dr r$$

        \item \textbf{积分计算}
        
        对$r$积分可将其化为定积分
        $$\frac{1}{2}\int_0^{\pi/2}\frac{\cos^2\theta\sin^2\theta}{(\cos^3\theta+\sin^3\theta)^2}\dr\theta$$
        至此已经是三角函数的有理函数的形式，可以利用万能代换强行计算。不过，设被积函数为$R(\cos\theta,\sin\theta)$，我们可以发现$R(\cos\theta,\sin\theta)=R(-\cos\theta,-\sin\theta)$，此时\textbf{一定可以}用$t=\tan\theta$代换。具体来说，我们强行从前面提出$\cos^{-2}\theta$可以将积分写为
        $$\frac{1}{2}\int_0^{\pi/2}\frac{\cos^4\theta\sin^2\theta}{(\cos^3\theta+\sin^3\theta)^2}\dr\tan\theta$$
        分子分母同除以$\cos^6\theta$即将其写为了
        $$\frac{1}{2}\int_0^{\pi/2}\frac{\tan^2\theta}{(1+\tan^3\theta)^2}\dr\tan\theta$$
        由此可进一步写为
        $$\frac{1}{6}\int_0^{\pi/2}\frac{1}{(1+\tan^3\theta)^2}\dr\tan^3\theta$$
        我们试着换元$t=\tan^3\theta$，这时上限将变为$\tan\frac{\pi}{2}$，从极限的角度这应该是$+\infty$，暂且忽略严谨性问题，得到形式上的结果
        $$\frac{1}{6}\int_0^{+\infty}\frac{1}{(1+t)^2}\dr t$$
        此时，被积函数的一个原函数是$-\frac{1}{1+t}$，它在0处是$-1$，而无穷处的极限是0，可以形式上用微积分基本定理得到结果
        $$\frac{1}{6}(0-(-1))=\frac{1}{6}$$

        \item (\extra)\textbf{严谨性分析}

        我们下面论证上述计算过程的合理性。
        
        由于$\frac{\cos^4\theta\sin^2\theta}{(\cos^3\theta+\sin^3\theta)^2}$在$[0,\frac{\pi}{2}]$连续，它在0到$x$的变上限积分也对$x$连续。考虑$\varepsilon>0$且充分小的
        $$I(\varepsilon)=\int_0^{\frac{\pi}{2}-\varepsilon}\frac{\cos^4\theta\sin^2\theta}{(\cos^3\theta+\sin^3\theta)^2}\dr\theta$$

        由复合函数连续性可知
        $$\lim_{\varepsilon\to0^+}I(\varepsilon)=I(0)$$
        
        而换元可得到
        $$I(\varepsilon)=\frac{1}{6}\int_0^{\tan^3(\pi/2-\varepsilon)}\frac{1}{(1+t)^2}\dr t=\frac{1}{6}\int_0^{\tan^{-3}\varepsilon}\frac{1}{(1+t)^2}\dr t$$
        利用微积分基本定理得
        $$I(\varepsilon)=\frac{1}{6}-\frac{1}{6}\cdot\frac{1}{1+\tan^{-3}\varepsilon}$$
        令$\varepsilon\to0^+$，计算极限的确可以得到$I(0)=\frac{1}{6}$。

        \note 更一般的情况将在下半学期研究。
    \end{itemize}

    \note 本题由于方程两边分别的齐次性，也可以考虑先设$y=tx$，得到$x^3(1+t^3)=x^2t$，从而在第一象限可参数化出
    $$x=\frac{t}{1+t^3},\quad y=\frac{t^2}{1+t^3}$$
    再利用Green公式转化为第二型曲线积分算面积(类似本讲义22.1.1第2题)。不过，这个做法仍然中间会出现广义积分，因为参数的范围是$t\in[0,+\infty)$。
}
\end{framed}

\note 由此可见重积分问题的计算难度往往\textbf{无法目测确定}，在尝试入手失败后建议\textbf{先转去做其他题目}。个人对这道题的评价是见到的往年题中难度最高的具体函数计算题，但它仅仅出现在卷子的中间部分。

\begin{framed}
\begin{exercise}
    计算曲面$\big(x^2+\frac{y^2}{4}+\frac{z^2}{5}\big)^2=x$围成区域的体积。
\end{exercise}
\sol{
    根据重积分几何意义，设区域为$\Omega$，其体积即为
    $$\iint_\Omega1\dr V$$
    计算分为如下步骤：
    \begin{itemize}
        \item \textbf{区域描述}
        
        由``围成''的有界性要求可知区域为满足
        $$\bigg(x^2+\frac{y^2}{4}+\frac{z^2}{5}\bigg)^2\le x$$
        的$(x,y,z)$构成的集合。
        
        \item \textbf{旋转对称性换元}
        
        由于出现了$x^2+\frac{y^2}{4}+\frac{z^2}{5}$作为核心部分，可以想到\textbf{广义球坐标}。不过，右端为$x$，为了使$x$的形式更加简单，我们更改通常广义球坐标的顺序，换元成
        $$x=\rho\cos\varphi,\quad y=2\rho\sin\varphi\cos\theta,\quad z=\sqrt5\rho\sin\varphi\sin\theta$$
        换元后的约束即为
        $$\rho^4\le\rho\cos\varphi$$
        忽略$\rho=0$的特殊点得到
        $$\rho^3\le\cos\varphi$$
        将上述条件确定的$(\rho,\varphi,\theta)$区域记为$\Omega_0$，被积函数仍为1，直接计算Jacobi行列式可知
        $$\dr x\dr y\dr z=2\sqrt5\rho^2\sin\varphi\dr\rho\dr\varphi\dr\theta$$
        于是原积分即
        $$2\sqrt5\iiint_{\Omega_0}\rho^2\sin\varphi\dr\rho\dr\varphi\dr\theta$$

        \item \textbf{累次积分化}
        
        由于被积函数对$\theta$为常数，且$\rho$被$\varphi$约束，考虑先$\theta$、再$\rho$、最后$\varphi$的积分顺序。此时，我们需要将集合描述成
        $$\alpha\le\varphi\le\beta,\quad\phi_1(\varphi)\le\rho\le\phi_2(\varphi),\quad\psi_1(\rho,\varphi)\le\theta\le\psi_2(\rho,\varphi)$$
        由$\rho$的自然范围可知需$\cos\varphi\ge0$，结合$\varphi$的自然范围得到
        $$0\le\varphi\le\frac{\pi}{2}$$
        此时对$\rho$的约束可写为
        $$0\le\rho\le\sqrt[3]{\cos\varphi}$$
        而$\theta$的范围即为自然范围
        $$0\le\theta\le2\pi$$
        综合以上讨论可写出累次积分
        $$2\sqrt5\int_0^{\pi/2}\dr\varphi\int_0^{\sqrt[3]{\cos\varphi}}\dr\rho\int_0^{2\pi}\rho^2\sin\varphi\dr\theta$$

        \item \textbf{积分计算}
        
        对$\theta$积分可得到
        $$4\sqrt5\pi\int_0^{\pi/2}\dr\varphi\int_0^{\sqrt[3]{\cos\varphi}}\rho^2\sin\varphi\dr\rho$$
        再对$\rho$积分得
        $$\frac{4\sqrt5}{3}\pi\int_0^{\pi/2}\cos\varphi\sin\varphi\dr\varphi$$
        这即为
        $$\frac{4\sqrt5}{3}\pi\int_0^{\pi/2}\sin\varphi\dr\sin\varphi=\frac{2\sqrt5\pi}{3}$$
    \end{itemize}

    \note 本题直接用球坐标换元也可计算，不过区域描述与计算都将稍复杂。
}
\end{framed}

\note 注意通常球坐标换元$(x,y,z)=(\rho\sin\varphi\cos\theta,\rho\sin\varphi\sin\theta,\rho\cos\varphi)$后$z$的形式更加简单，而实际上为了使$x$或$y$的形式更加简单，也可以更改三个变量的\textbf{顺序}，得到的Jacobi行列式绝对值不会改变(相当于原行列式交换了行列)。

\

\noindent\textbf{抽象函数}
\begin{framed}
\begin{exercise}
    设$f(x)$为$\rb$上恒正的连续函数，对参数$t>0$，记
    $$F(t)=\frac{\iiint_{x^2+y^2+z^2\le t^2}f(x^2+y^2+z^2)\dr x\dr y\dr z}{\iint_{x^2+y^2\le t^2}(x^2+y^2)f(x^2+y^2)\dr x\dr y}$$
    证明$F(t)$在$t>0$时严格单调减。
\end{exercise}
\sol{
    分为如下步骤：
    \begin{itemize}
        \item \textbf{分母化简}
        
        考虑到区域和函数的旋转对称性，进行\textbf{极坐标}换元
        $$x=r\cos\theta,\quad y=r\sin\theta$$
        此时原本的被积函数为$r^2f(r^2)$，且计算Jacobi行列式得$\dr x\dr y=r\dr r\dr\theta$。此时区域要求仅$r^2\le t^2$，结合自然范围可知约束即化为
        $$0\le r\le t,\quad0\le\theta\le2\pi$$
        由此分母为累次积分(最后对$r$积分以保证含有未知函数$f$的部分最后积分)
        $$\int_0^t\dr r\int_0^{2\pi}r^3f(r^2)\dr\theta$$
        对$\theta$计算积分得结果
        $$2\pi\int_0^tr^3f(r^2)\dr r$$

        \item \textbf{分子化简}
        
        考虑到区域和函数的旋转对称性，进行\textbf{球坐标}换元
        $$x=\rho\sin\varphi\cos\theta,\quad y=\rho\sin\varphi\sin\theta,\quad z=\rho\cos\varphi$$
        此时原本的被积函数为$f(\rho^2)$，且计算Jacobi行列式得$\dr x\dr y\dr z=\rho^2\sin\varphi\dr\rho\dr\varphi\dr\theta$。此时区域要求仅$\rho^2\le t^2$，结合自然范围可知约束即化为
        $$0\le\rho\le t,\quad0\le\varphi\le\pi,\quad0\le\theta\le2\pi$$
        由此分子为累次积分(最后对$\rho$积分以保证含有未知函数$f$的部分最后积分)
        $$\int_0^t\dr\rho\int_0^\pi\dr\varphi\int_0^{2\pi}\rho^2f(\rho^2)\sin\varphi\dr\theta$$
        对$\theta$计算积分得结果
        $$2\pi\int_0^t\dr\rho\int_0^\pi\rho^2f(\rho^2)\sin\varphi\dr\varphi$$
        再对$\varphi$积分得
        $$4\pi\int_0^t\rho^2f(\rho^2)\dr\rho$$
        
        \item \textbf{导数计算}
        
        由于定积分更换变量无影响，我们将上下的变量都写为$r$，有
        $$F(t)=2\frac{\int_0^tr^2f(r^2)\dr r}{\int_0^tr^3f(r^2)\dr r}$$
        为了判断单调性，我们对$F(t)$求导，由变上限积分的导数结论可知
        $$\frac{\dr}{\dr t}\int_0^tr^2f(r^2)\dr r=t^2f(t^2),\quad\frac{\dr}{\dr t}\int_0^tr^3f(r^2)\dr r=t^3f(t^2)$$
        由此可知
        $$F'(t)=2\bigg(\int_0^tr^3f(r^2)\dr r\bigg)^{-2}\bigg(t^2f(t^2)\int_0^tr^3f(r^2)\dr r-t^3f(t^2)\int_0^tr^2f(r^2)\dr r\bigg)$$
        可直接提出相同部分得到
        $$F'(t)=2t^2f(t^2)\bigg(\int_0^tr^3f(r^2)\dr r\bigg)^{-2}\bigg(\int_0^tr^3f(r^2)\dr r-t\int_0^tr^2f(r^2)\dr r\bigg)$$
        也即
        $$F'(t)=2t^2f(t^2)\bigg(\int_0^tr^3f(r^2)\dr r\bigg)^{-2}\int_0^t(r-t)r^2f(r^2)\dr r$$
        由积分范围$r-t\le0$、且$r^2\ge0$，由条件$f(r^2)>0$，由此可知
        $$\forall r\in(0,t),\quad(r-t)r^2f(r^2)<0$$
        由于此函数连续，利用上学期知识能进一步得到(可见上学期讲义6.3.1)
        $$\int_0^t(r-t)r^2f(r^2)\dr r<0$$
        同理可知
        $$\int_0^tr^3f(r^2)\dr r>0$$
        且$t>0$时$t^2f(t^2)>0$，综合以上就得到$t>0$时
        $$F'(t)<0$$
        这就证明了严格单调减。
    \end{itemize}
}
\end{framed}

\note 对抽象函数的处理本质上和对具体函数并无区别，只是抽象函数往往需要视为\textbf{最难以计算的部分}最后处理。

\begin{framed}
\begin{exercise}
    设$f(x)$为$\rb$上的连续函数，且$f(0)=0$、$f'(0)=10$。对参数$t>0$，记
    $$V(t)=\bigg\{(x,y,z)\in\rb^3\ \bigg|\ x^2+16y^2+\frac{z^2}{25}\le t^2\bigg\}$$
    计算
    $$\lim_{t\to0^+}\frac{1}{t^5}\iiint_{V(t)}f\bigg(x^2+16y^2+\frac{z^2}{25}\bigg)\dr x\dr y\dr z$$
\end{exercise}
\sol{
    分为如下步骤：
    \begin{itemize}
        \item \textbf{积分化简}
        
        考虑到区域和函数的形式，进行\textbf{广义球坐标}换元
        $$x=\rho\sin\varphi\cos\theta,\quad y=\frac{1}{4}\rho\sin\varphi\sin\theta,\quad z=5\rho\cos\varphi$$
        此时原本的被积函数为$f(\rho^2)$，且计算Jacobi行列式得$\dr x\dr y\dr z=\frac{5}{4}\rho^2\sin\varphi\dr\rho\dr\varphi\dr\theta$。此时$V(t)$要求仅$\rho^2\le t^2$，结合自然范围可知约束即化为
        $$0\le\rho\le t,\quad0\le\varphi\le\pi,\quad0\le\theta\le2\pi$$
        由此积分可化为累次积分(最后对$\rho$积分以保证含有未知函数$f$的部分最后积分)
        $$\frac{5}{4}\int_0^t\dr\rho\int_0^\pi\dr\varphi\int_0^{2\pi}\rho^2f(\rho^2)\sin\varphi\dr\theta$$
        对$\theta$计算积分得结果
        $$\frac{5}{2}\pi\int_0^t\dr\rho\int_0^\pi\rho^2f(\rho^2)\sin\varphi\dr\varphi$$
        再对$\varphi$积分得
        $$5\pi\int_0^t\rho^2f(\rho^2)\dr\rho$$
        
        \item \textbf{极限计算}
        
        由上方计算知原极限即
        $$\lim_{t\to0^+}\frac{5\pi\int_0^t\rho^2f(\rho^2)\dr\rho}{t^5}$$
        此形式对$t$为负也可定义，我们考虑更一般的结果
        $$\lim_{t\to0}\frac{5\pi\int_0^t\rho^2f(\rho^2)\dr\rho}{t^5}$$

        \note 将右极限转化为极限是为了保证\textbf{洛必达使用的严谨性}，站在这学期的纯计算角度而言，也可以跳过这一步。

        验证条件并使用洛必达法则，由变上限积分的导数结论可知
        $$\frac{\dr}{\dr t}\int_0^t\rho^2f(\rho^2)\dr\rho=t^2f(t^2)$$
        从而只需计算
        $$\lim_{t\to0}\frac{\pi t^2f(t^2)}{t^4}=\lim_{t\to0}\frac{\pi f(t^2)}{t^2}$$
        利用$f(0)=0$、$f'(0)=10$可知
        $$\lim_{s\to0}\frac{f(s)-f(0)}{s}=\lim_{s\to0}\frac{f(s)}{s}=10$$
        利用复合函数极限性质，换元$s=t^2$最终得到原极限为$10\pi$。

        \note 最后$\frac{f(t^2)}{t^2}$的极限计算过程里，由于并不知道$f$在0以外的点的导数，无法再通过洛必达进行处理。不过，强行使用(相当于假设了$f$更强的光滑性)仍可以得到正确的结果。

        \note 也可以作泰勒展开$f(t^2)=f(0)+t^2f'(0)+o(t^2)$得到结论。
    \end{itemize}
}
\end{framed}

\note 注意熟悉上学期的一些\textbf{极限相关基本结论}，如基本等价无穷小、洛必达法则、变上限积分等。

\subsubsection{曲线积分}
\noindent\textbf{第一型曲线积分}
\begin{framed}
\begin{exercise}
    计算
    $$\int_L(|x|+x^3+x^5+|xy|)\dr s$$
    其中$L$为椭圆$4x^2+y^2=1$在$y\ge0$的部分。
\end{exercise}
\sol{
    \item \textbf{反射对称化简}
    
    由于曲线$L$满足$(x,y)\in L$当且仅当$(-x,y)\in L$，且
    $$x^3=-(-x)^3,\quad x^5=-(-x)^5$$
    利用函数折半可发现
    $$\int_Lx^3\dr s=\int_Lx^5\dr s=0$$
    从而积分可化为
    $$\int_L(|x|+|xy|)\dr s$$

    \item \textbf{参数化}
            
    由其为椭圆，利用\textbf{广义极坐标}可参数化为
    $$x=\frac{1}{2}\cos\theta,\quad y=\sin\theta$$
    由$y\ge0$，根据极坐标几何意义可知范围$0\le\theta\le\pi$。

    \item \textbf{化为定积分}

    用下标$\theta$表示对$\theta$求导，直接计算有
    $$x_\theta=-\frac{1}{2}\sin\theta,\quad y_\theta=\cos\theta$$
    从而
    $$\dr s=\sqrt{x_\theta^2+y_\theta^2}\dr\theta=\sqrt{\frac{1}{4}\sin^2\theta+\cos^2\theta}\dr\theta$$
    进一步计算可发现被积函数为$\frac{1}{2}(|\cos\theta|+|\sin\theta\cos\theta|)$，由此原积分为
    $$\int_0^\pi\frac{1}{2}(|\cos\theta|+|\sin\theta\cos\theta|)\sqrt{\frac{1}{4}\sin^2\theta+\cos^2\theta}\dr\theta$$

    \item \textbf{积分计算}
            
    为了消除绝对值，我们需要按照0到$\frac{\pi}{2}$与$\frac{\pi}{2}$到$\pi$两段，在第二段中换元$\theta$为$\pi-\theta$可发现
    $$\int_{\pi/2}^\pi(|\cos\theta|+|\sin\theta\cos\theta|)\sqrt{\frac{1}{4}\sin^2\theta+\cos^2\theta}\dr\theta=\int_0^{\pi/2}(|\cos\theta|+|\sin\theta\cos\theta|)\sqrt{\frac{1}{4}\sin^2\theta+\cos^2\theta}\dr\theta$$
    由此积分化为了
    $$\int_0^{\pi/2}(|\cos\theta|+|\sin\theta\cos\theta|)\sqrt{\frac{1}{4}\sin^2\theta+\cos^2\theta}\dr\theta$$
    此时可去掉绝对值得到
    $$\int_0^{\pi/2}(\cos\theta+\sin\theta\cos\theta)\sqrt{\frac{1}{4}\sin^2\theta+\cos^2\theta}\dr\theta$$

    \note 也可直接在曲线积分中利用对称性去除绝对值，将其写为第一象限中积分的两倍，具体方法的选择取决于对\textbf{定积分的对称性处理}是否熟悉。

    由于根号下的$\cos^2\theta$可以看作$1-\sin^2\theta$，我们可以将积分改写为
    $$\int_0^{\pi/2}(1+\sin\theta)\sqrt{1-\frac{3}{4}\sin^2\theta}\dr\sin\theta$$
    代入$t=\sin\theta$进一步化为
    $$\int_0^1(1+t)\sqrt{1-\frac{3}{4}t^2}\dr t$$
    我们拆成两部分计算，将积分写为
    $$I_1+I_2,\quad I_1=\int_0^1\sqrt{1-\frac{3}{4}t^2}\dr t,\quad I_2=\int_0^1t\sqrt{1-\frac{3}{4}t^2}\dr t$$
    换元$s=t^2$直接计算有
    $$I_2=\int_0^1\frac{1}{2}\sqrt{1-\frac{3}{4}s}\dr s$$
    被积函数已经是幂函数复合一次函数的形式，可直接写出一个原函数
    $$-\frac{4}{9}\bigg(1-\frac{3}{4}s\bigg)^{3/2}$$
    从而得积分结果
    $$I_2=\frac{4}{9}\bigg(1-\frac{1}{8}\bigg)=\frac{7}{18}$$

    对$I_1$，为了化简根号下的部分，重新设$t=\frac{2}{\sqrt3}\sin\phi$，则积分变为0到$\arcsin\frac{\sqrt3}{2}=\frac{\pi}{3}$，代入可发现原积分成为
    $$I_1=\frac{2}{\sqrt3}\int_0^{\pi/3}\cos^2\phi\dr\phi$$
    利用$\cos^2\phi=\frac{1}{2}(\cos(2\phi)+1)$可直接写出被积函数的一个原函数
    $$\frac{1}{4}\sin(2\phi)+\frac{1}{2}\phi$$
    从而可得到结果
    $$I_1=\frac{2}{\sqrt3}\bigg(\frac{1}{4}\sin\frac{2\pi}{3}+\frac{\pi}{6}\bigg)=\frac{1}{4}+\frac{\pi}{3\sqrt3}$$
    综合以上得结果为
    $$I_1+I_2=\frac{23}{36}+\frac{\sqrt3\pi}{9}$$
}
\end{framed}

\note 第一型曲线积分与第一型曲面积分有与重积分\textbf{完全类似}的对称性化简方式，但\textbf{第二型曲线/曲面积分不同}，因此不建议在第二型积分化成第一型积分或定积分前运用任何对称性。

\begin{framed}
\begin{exercise}
    计算
    $$\oint_C(x+y+1)\dr s$$
    其中$C$是以$O(0,0)$、$A(1,0)$、$B(0,1)$为顶点的三角形。
\end{exercise}
\sol{
    分为如下步骤：
    \begin{itemize}
        \item \textbf{参数化}
                
        根据直线方程可直接得到$OA$、$AB$、$BO$的参数方程
        $$\begin{cases}x=t\\y=0\end{cases}(0\le t\le1),\quad\begin{cases}x=1-t\\y=t\end{cases}(0\le t\le1),\quad\begin{cases}x=0\\y=1-t\end{cases}(0\le t\le1)$$

        \item \textbf{化为定积分与计算}
                
        用下标$t$表示对$t$求导，直接计算有第一段中
        $$x_t=1,\quad y_t=0,\quad\dr s=\sqrt{x_t^2+y_t^2}\dr t=\dr t$$
        第二段中
        $$x_t=-1,\quad y_t=1,\quad\dr s=\sqrt{x_t^2+y_t^2}\dr t=\sqrt2\dr t$$
        第三段中
        $$x_t=0,\quad y_t=-1,\quad\dr s=\sqrt{x_t^2+y_t^2}\dr t=\dr t$$
        再代入每段中$x$、$y$的表达式可得结果为
        $$\int_0^1(t+1)\dr t+\int_0^12\sqrt2\dr t+\int_0^1(2-t)\dr t=3+2\sqrt2$$
        \end{itemize}
}
\end{framed}

\note 注意分段曲线可能需要\textbf{分段进行参数化}，积分结果为每段相加。

\

\noindent\textbf{第二型曲线积分}
\begin{framed}
\begin{exercise}\label{green out}
    计算
    $$\oint_{L^+}\bigg(\frac{y^2+y+4x^2}{4x^2+y^2}+\sin x^2\bigg)\dr x+\bigg(\frac{4x^2-x+y^2}{4x^2+y^2}+\sin y^2\bigg)\dr y$$
    其中$L$为$x^2+y^2=9$的$y\ge0$部分与$\frac{x^2}{9}+\frac{y^2}{16}=1$的$y\le0$部分围成的闭曲线。
\end{exercise}
\sol{
    分为如下步骤：
    \begin{itemize}
        \item \textbf{格林公式}
        
        将被积函数看作$P\dr x+Q\dr y$，由初等函数知$P$、$Q$在$\rb^2$去除原点外的区域内可导且导函数连续，直接计算得
        $$Q_x-P_y=\frac{4x^2-y^2}{(4x^2+y^2)^2}-\frac{4x^2-y^2}{(4x^2+y^2)^2}=0$$
        不过，此区域并非单连通，无法得到积分为0。我们需要利用多连通区域的格林公式将积分曲线\textbf{转化为更好计算的曲线}。

        为了化简分母，考虑曲线
        $$L_0\colon4x^2+y^2=r^2$$
        我们取一个足够小的$r>0$使得$L_0$包含在$L$中，设$L_0$与$L$之间的区域为$D$，利用\textbf{多连通区域的格林公式}有
        $$\oint_{L^+}(P\dr x+Q\dr y)-\oint_{L_0^+}(P\dr x+Q\dr y)=\iint_D(Q_x-P_y)\dr x\dr y=0$$
        由此积分转化为
        $$\oint_{L_0^+}\bigg(\frac{y^2+y+4x^2}{4x^2+y^2}+\sin x^2\bigg)\dr x+\bigg(\frac{4x^2-x+y^2}{4x^2+y^2}+\sin y^2\bigg)\dr y$$

        \note 保留抽象的$r$往往比取出具体的$r$更易于计算，这样的$r$的存在性论证相对麻烦，不过基本可通过直观确定。

        \item \textbf{积分化简}
        
        由于$L_0$上满足$4x^2+y^2=r^2$，我们可以将$4x^2+y^2$改写为$r^2$以化简积分成
        $$\oint_{L_0^+}\bigg(\frac{y}{r^2}+1+\sin x^2\bigg)\dr x+\bigg(\frac{-x}{r^2}+1+\sin y^2\bigg)\dr y$$

        \item \textbf{格林公式}
        
        由于此积分仍然难以直接看出计算方式，考虑重新应用格林公式。将新的被积函数看作$P_0\dr x+Q_0\dr y$，由初等函数知$P_0$、$Q_0$在$L_0$围成的单连通区域$D_0$内可导且导函数连续，进一步计算验证可知
        $$(Q_0)_x-(P_0)_y=-\frac{1}{r^2}-\frac{1}{r^2}=-\frac{2}{r^2}$$
        积分最终化为
        $$-\iint_{D_0}\frac{2}{r^2}\dr x\dr y$$

        \item \textbf{积分计算}

        利用重积分的几何意义，$\iint_D\dr x\dr y$即为区域的面积，而此积分区域是椭圆$4x^2+y^2=r^2$，长轴长$r$、短轴长$\frac{r}{2}$，从而面积为
        $$\frac{r^2\pi}{2}$$
        最终得到积分结果$-\pi$。

        \note 结果中不应该有$r$的存在，因为由多连通区域的格林公式，对任何$r$的积分应相同。
    \end{itemize}

    \note 本题也可以先用格林公式说明$(1+\sin x^2)\dr x+(1+\sin y^2)\dr y$在封闭曲线的上的积分为0，再作差后处理剩余部分。
}
\end{framed}

\note 对于平面第二型曲线积分，无论曲线为何应\textbf{尝试计算}$Q_x-P_y$，若简单再考虑如何使用格林公式，常见如挖点(\exref{green out})、补线(\exref{green add})等。

\note 注意使用格林公式时若需要挖点，本质上是希望改变积分曲线以\textbf{尽量简化分母}。

\note 若简化完分母后在整个区域有定义，可以\textbf{再次使用格林公式}，由此可见，即使是相同的曲线积分(如在$x^2+y^2=1$上考虑$\dr x+\dr y$与$(x^2+y^2)(\dr x+\dr y)$的积分)，使用格林公式后也可能得到\textbf{不同形式的重积分}。对高斯公式、斯托克斯公式同理。

\note 重积分、第一型曲线积分、第一型曲面积分的\textbf{几何意义}往往可以用来快速计算一些简单的积分结果。

\begin{framed}
\begin{exercise}\label{green add}
    计算
    $$\int_L\frac{\big(x-\frac{1}{2}-y\big)\dr x+\big(x-\frac{1}{2}+y\big)\dr y}{\big(x-\frac{1}{2}\big)^2+y^2}$$
    其中$L$为$(0,-1)$沿着抛物线$y^2=1-x$到$(0,1)$的曲线段。
\end{exercise}
\sol{
    我们介绍两种方法，首先是格林公式分析后转化到另\textbf{一条曲线}计算：
    \begin{itemize}
        \item \textbf{格林公式}
        
        将被积函数看作$P\dr x+Q\dr y$，由初等函数知$P$、$Q$在$\rb^2$去除原点外的区域内可导且导函数连续，直接计算得
        $$Q_x-P_y=\frac{y^2-\big(x-\frac{1}{2}\big)^2-2\big(x-\frac{1}{2}\big)y}{\big(\big(x-\frac{1}{2}\big)^2+y^2\big)^2}-\frac{y^2-\big(x-\frac{1}{2}\big)^2-2\big(x-\frac{1}{2}\big)y}{\big(\big(x-\frac{1}{2}\big)^2+y^2\big)^2}=0$$

        由此，考虑另一条从$(0,-1)$出发到$(0,1)$的曲线$L_0$，若$L$与$L_0$围成的区域\textbf{不包含$(\frac{1}{2},0)$}，设围成的区域为$D$，可利用格林公式得到(这里有正负是因为无法保证将$L_0$反向拼接上$L$得到的曲线是顺时针还是逆时针)
        $$\int_L(P\dr x+Q\dr y)-\int_{L_0}(P\dr x+Q\dr y)=\pm\iint_D(Q_x-P_y)\dr x\dr y=0$$

        \note 由此，我们\textbf{不能}取$L_0$是连接两点的线段。

        为了让计算尽量简单，我们希望$L_0$是某个以$(\frac{1}{2},0)$为圆心的圆周的一部分，直接代入$(0,-1)$与$(0,1)$发现可取
        $$\bigg(x-\frac{1}{2}\bigg)^2+y^2=\frac{5}{4}$$
        且逆时针从$(0,-1)$走到$(0,1)$\ (可以直接画图确定此时$L_0$与$L$围成的区域不包含原点，若顺时针则会包含)。

        \note 若两点到$(\frac{1}{2},0)$的距离并不相同，则只能采用更复杂的计算方式，如通过绕开这点的水平、竖直线段计算积分。

        \item \textbf{参数化}
        
        利用极坐标可以得到整个圆周的参数方程
        $$\begin{cases}x=\frac{1}{2}+\frac{\sqrt5}{2}\cos\theta\\y=\frac{\sqrt5}{2}\sin\theta\end{cases}$$

        为了进一步计算，我们还需要确定参数的范围。直接计算可发现$(0,-1)$与$(0,1)$分别对应点(这里的$\theta$本质上都可以加$2\pi$的整数倍，如此写是为了方便描述范围)
        $$\theta=\arcsin\frac{2}{\sqrt5}-\pi,\quad\theta=\pi-\arcsin\frac{2}{\sqrt5}$$
        且通过极坐标几何意义可发现实际范围即为
        $$\arcsin\frac{2}{\sqrt5}-\pi\le\theta\le\pi-\arcsin\frac{2}{\sqrt5}$$

        \note 此范围的几何方法确定相比代数会简单很多。

        \item \textbf{化为定积分与计算}
        
        用下标$\theta$表示对$\theta$求导，直接计算有
        $$x_\theta=-\frac{\sqrt5}{2}\sin\theta,\quad y_\theta=\frac{\sqrt5}{2}\cos\theta$$
        代入$\dr x=x_\theta\dr\theta$、$\dr y=y_\theta\dr\theta$与$x$、$y$的参数方程可得原积分为
        $$\int_{\arcsin(2/\sqrt5)-\pi}^{\pi-\arcsin(2/\sqrt5)}\dr\theta=2\pi-2\arcsin\frac{2}{\sqrt5}$$

        \note 也可将$\arcsin\frac{2}{\sqrt5}$写成$\arctan2$以进一步化简。
    \end{itemize}
    也可以转化到\textbf{两条曲线}计算：
    \begin{itemize}
        \item \textbf{格林公式}
        
        在进行与之前相同的分析后，考虑$(0,-1)$到$(0,1)$的线段$L_1$与圆
        $$L_2\colon\bigg(x-\frac{1}{2}\bigg)^2+y^2=r^2$$
        我们取一个足够小的$r>0$使得$L_2$包含在$L$与$L_1$围成的区域中，设位于$L$与$L_1$拼接内侧、$L_2$外侧的区域为$D$。

        利用多连通区域的格林公式，我们可以得到(下方的符号可以直接通过几何关系确定，沿着$L$后走$L_1$的反向构成逆时针的闭曲线)
        $$\int_L(P\dr x+Q\dr y)-\int_{L_1}(P\dr x+Q\dr y)-\oint_{L_2^+}(P\dr x+Q\dr y)=\iint_D(Q_x-P_y)\dr x\dr y=0$$
        接下来只需计算$L_1$与$L_2$上$P\dr x+Q\dr y$的积分。

        \item \textbf{$L_1$积分计算}
        
        可直接写出线段的参数方程
        $$\begin{cases}x=0\\y=t\end{cases}(-1\le t\le1)$$
        用下标$t$表示对$t$求导，直接计算有
        $$x_t=0,\quad y_t=1$$
        代入$\dr x=x_t\dr t$、$\dr y=y_t\dr t$与$x$、$y$的参数方程可得原积分为
        $$\int_{-1}^1\frac{t-\frac{1}{2}}{\frac{1}{4}+t^2}\dr t$$
        利用对称性，将$t$换元为$-t$可发现
        $$\int_{-1}^1\frac{t-\frac{1}{2}}{\frac{1}{4}+t^2}\dr t=\int_{-1}^1\frac{-t-\frac{1}{2}}{\frac{1}{4}+t^2}\dr t$$
        从而$t$的部分积分为0，原积分可写为
        $$-\frac{1}{2}\int_{-1}^1\frac{1}{\frac{1}{4}+t^2}\dr t$$
        利用上册教材3.2节的积分公式，$c>0$时
        $$\int\frac{1}{c^2+z^2}\dr z=\frac{1}{c}\arctan\frac{z}{c}+C$$
        由此可直接得到积分结果
        $$-\frac{1}{2}(2\arctan 2-2\arctan(-2))=-2\arctan2$$

        \item \textbf{$L_2$积分计算}
        
        利用极坐标可以得到圆周的参数方程(注意这里$r$是常数，不再是极坐标的变量)
        $$\begin{cases}x=\frac{1}{2}+r\cos\theta\\y=r\sin\theta\end{cases}(0\le\theta\le2\pi)$$
        且利用极坐标几何意义，此参数方程已经是逆时针的定向，无需在结果中添加负号。

        用下标$\theta$表示对$\theta$求导，直接计算有
        $$x_\theta=-r\sin\theta,\quad y_\theta=r\cos\theta$$
        代入$\dr x=x_\theta\dr\theta$、$\dr y=y_\theta\dr\theta$与$x$、$y$的参数方程可得原积分为
        $$\int_0^{2\pi}1\dr\theta=2\pi$$

        \item \textbf{综合结果}

        综合以上讨论最终得到结果为
        $$\int_{L_1}(P\dr x+Q\dr y)+\int_{L_2}(P\dr x+Q\dr y)=2\pi-2\arctan2$$
    \end{itemize}

    \note 对本题来说，加两条曲线方法的\textbf{计算}会稍显复杂，但加一条曲线方法的\textbf{角度范围确定}又是十分易错的(很容易将范围误判成$-\arcsin\frac{2}{\sqrt5}$到$\arcsin\frac{2}{\sqrt5}$)，一般还是建议更稳定的两条曲线做法。
}
\end{framed}

\note 使用格林公式时务必注意围成的区域中是否有\textbf{无定义或不可导}的点，这时需要通过规避或挖洞等方法进行进一步的处理。

\note 与平面重积分类似，平面曲线积分也非常建议大家在有时间时\textbf{画出(带方向的)积分曲线}，避免范围选择等错误。

\begin{framed}
\begin{exercise}
    计算
    $$\oint_L\frac{-y}{4x^2+y^2}\dr x+\frac{x}{4x^2+y^2}\dr y+z\dr z$$
    其中$L$为曲面$4x^2+y^2=1$与平面$2x+y+z=1$的交线，且从$x$轴正向看去沿逆时针方向。
\end{exercise}
\sol{
    分为如下步骤：
    \begin{itemize}
        \item \textbf{积分化简}
        
        由于$L$上的点都满足$4x^2+y^2=1$，我们可以将$4x^2+y^2$改写为1以化简积分成
        $$\oint_L(-y)\dr x+x\dr y+z\dr z$$

        \item \textbf{参数化}
        
        $4x^2+y^2=1$是一个柱面，对$x$、$y$可直接参数化为
        $$\begin{cases}x=\frac{1}{2}\cos\theta\\y=\sin\theta\end{cases}(0\le\theta\le2\pi)$$
        代入平面方程就可以得到整体的参数方程
        $$\begin{cases}x=\frac{1}{2}\cos\theta\\y=\sin\theta\\z=1-\cos\theta-\sin\theta\end{cases}(0\le\theta\le2\pi)$$
        
        \item \textbf{定向}
        
        为了确定定向，用下标$\theta$代表对$\theta$求导，直接计算有
        $$x_\theta=-\frac{1}{2}\sin\theta,\quad y_\theta=\cos\theta,\quad z_\theta=\sin\theta-\cos\theta$$
        我们\textbf{任取}一个使得$x_\theta$、$y_\theta$、$z_\theta$\textbf{不全为0}的点，如取$\theta=\pi$，此处有$(x,y,z)=(-\frac{1}{2},0,2)$、$(x_\theta,y_\theta,z_\theta)=(0,-1,1)$，在曲线上的点$(x,y,z)$处画出向量$(x_\theta,y_\theta,z_\theta)$，可以发现它指向逆时针方向，与要求的定向相同，从而化为定积分无需加负号。

        \note 事实上还有更简单的无需画图的判断方法：根据极坐标参数化的性质，参数化后曲线在$Oxy$平面上的投影是逆时针方向的(注意这是指从上往下看)，而$x$更小时从参数方程可看出$z$更大，也即曲线在$x$增大时高度逐渐下降，$x$轴正向看去将是``从上往下''的，因此仍是逆时针。

        \item \textbf{化为定积分与计算}
        
        代入$\dr x=x_\theta\dr\theta$、$\dr y=y_\theta\dr\theta$、$\dr z=z_\theta\dr\theta$与$x$、$y$、$z$的参数方程可得原积分为
        $$\int_0^{2\pi}\bigg(\frac{1}{2}+\sin\theta-\cos\theta-\sin^2\theta+\cos^2\theta\bigg)\dr\theta$$
        将$\cos^2\theta-\sin^2\theta$写为$\cos(2\theta)$，利用积分范围0到$2\pi$中的对称性(或直接写出原函数计算)可以发现带三角函数的部分积分均为0，因此最终得到积分为
        $$\int_0^{2\pi}\frac{1}{2}\dr\theta=\pi$$
    \end{itemize}

    \note 本题在积分化简后的过程也可以通过\textbf{斯托克斯公式}解决，做法类似\exref{stolkes ellipse}。
}
\end{framed}

\note 应当随时注意被积函数是否可\textbf{由区域直接化简}，所有其他操作都建议在化简后进行(即使形式很类似可以求导化简)。

\note 必须掌握一种\textbf{曲线积分确定参数定向}的方式，本题的定向过程是个人推荐的方法：任取一个满足
$$(x_\theta,y_\theta,z_\theta)\ne\vec{0}$$
的$\theta$\ (否则\textbf{无法确定切向量})，在曲线上画出这点与对应的切向量$(x_\theta,y_\theta,z_\theta)$，观察指向顺时针还是逆时针。

\begin{framed}
\begin{exercise}\label{stolkes circle}
    计算
    $$\oint_Ly\dr x+z\dr y+x\dr z$$
    其中$L$为$x^2+y^2+z^2=a^2$与$x+y+z=0$交出的圆周，且从$x$轴正向看去沿逆时针方向，这里参数$a>0$。
\end{exercise}
\sol{
    分为如下步骤：
    \begin{itemize}
        \item \textbf{斯托克斯公式}
        
        将被积函数看作$P\dr x+Q\dr y+R\dr z$，由初等函数知$P$、$Q$、$R$在$\rb^3$上可导且导函数连续，直接计算得
        $$R_y-Q_z=-1,\quad P_z-R_x=-1,\quad Q_x-P_y=-1$$
        为了应用斯托克斯公式，考虑$x^2+y^2+z^2=a^2$在平面$x+y+z=0$上围成的部分$S$，则原积分可化为
        $$\iint_S-\dr y\dr z-\dr z\dr x-\dr x\dr y$$
        利用右手定则，$S$应选择指向$x$轴正向的法向作为正向。

        \item \textbf{化为第一型曲面积分}
        
        利用第一型曲面积分与第二型曲面积分的关系可知(这里$\vec{n}$代表第二型曲面积分对应定向的单位法向量)
        $$\vec{n}\dr S=(\dr y\dr z,\dr z\dr x,\dr x\dr y)$$
        利用平面的几何性质可知它的一个法向为$(1,1,1)$，指向$x$轴正向即代表$x$分量为正，从而对应分量的单位法向量为
        $$\frac{1}{\sqrt3}(1,1,1)$$
        由此即得原积分为
        $$\iint_S(-1,-1,-1)\cdot(\dr y\dr z,\dr z\dr x,\dr x\dr y)=\iint_S(-1,-1,-1)\cdot\frac{1}{\sqrt3}(1,1,1)\dr S=-\sqrt3\iint_S\dr S$$

        \item \textbf{积分计算}
        
        利用第一型曲面积分的几何意义，$\iint_S\dr S$即为曲面面积，而$S$是过原点的平面$x+y+z=0$被半径$a$的球$x^2+y^2+z^2=a^2$切割出的圆，半径为$a$，因此面积为$\pi a^2$，最终得到结果为
        $$-\sqrt3\pi a^2$$
    \end{itemize}

    \note 本题也可\textbf{参数化}后直接计算，参数化方法见本讲义21.1.1第4题。
}
\end{framed}

\note 斯托克斯公式可以将曲线积分转化成\textbf{任何一个}边界是此曲线的曲面的积分，最常用的即转化为平面。此外，曲线从哪个方向看去是\textbf{逆时针}的，利用斯托克斯后对应的定向就是那个方向的法向。

\note \textbf{平面}上对\textbf{常向量}$(a,b,c)$的第二型曲面积分$a\dr y\dr z+b\dr z\dr x+c\dr x\dr y$一定可以用相同方法化为面积的倍数。

\note 对于空间曲线，大家务必掌握一些常用的参数化方法，如\textbf{求解一个变量后化为二次曲线配方}(类似\exref{ball cont})。

\begin{framed}
\begin{exercise}\label{stolkes ellipse}
    计算
    $$\oint_L(y-z+\sin^2x)\dr x+(z-x+\sin^2y)\dr y+(x-y+\sin^2z)\dr z$$
    其中$L$为$x^2+y^2=1$与$x+z=1$交出的椭圆，且从$z$轴正向看去沿逆时针方向。
\end{exercise}
\sol{
    分为如下步骤：
    \begin{itemize}
        \item \textbf{斯托克斯公式}
        
        将被积函数看作$P\dr x+Q\dr y+R\dr z$，由初等函数知$P$、$Q$、$R$在$\rb^3$上可导且导函数连续，直接计算得
        $$R_y-Q_z=-2,\quad P_z-R_x=-2,\quad Q_x-P_y=-2$$
        为了应用斯托克斯公式，考虑$x^2+y^2=1$在平面$x+z=0$上围成的部分$S$，则原积分可化为
        $$\iint_S-2\dr y\dr z-2\dr z\dr x-2\dr x\dr y$$
        利用右手定则，$S$应选择指向$z$轴正向的法向作为正向。

        \item \textbf{化为第一型曲面积分}
        
        利用第一型曲面积分与第二型曲面积分的关系可知(这里$\vec{n}$代表第二型曲面积分对应定向的单位法向量)
        $$\vec{n}\dr S=(\dr y\dr z,\dr z\dr x,\dr x\dr y)$$
        利用平面的几何性质可知它的一个法向为$(1,0,1)$，指向$z$轴正向即代表$z$分量为正，从而对应分量的单位法向量为
        $$\frac{1}{\sqrt2}(1,0,1)$$
        由此即得原积分为
        $$\iint_S(-2,-2,-2)\cdot(\dr y\dr z,\dr z\dr x,\dr x\dr y)=\iint_S(-2,-2,-2)\cdot\frac{1}{\sqrt2}(1,0,1)\dr S=-2\sqrt2\iint_S\dr S$$

        \note 虽然这里只需计算面积就可以得到结果，但面积\textbf{并不容易确定}，因此还是应正常计算第一型曲面积分。

        \item \textbf{区域描述与参数化}
        
        由于$S$是$x^2+y^2=1$在平面$x+z=0$上围成的部分，这即是平面$x+z=0$上满足$x^2+y^2\le1$的点的集合。

        此平面上可以将$z$看作$x$的函数，从而有自然的参数化
        $$\vec{r}=(x,y,z)=(u,v,-u)$$
        代入约束可得到参数范围
        $$u^2+v^2\le1$$
        将上述条件确定的$(u,v)$区域记为$D$。

        \item \textbf{化为重积分}
        
        用下标$u$、$v$表示对$u$、$v$求偏导并保持另一个不变，直接计算有
        $$\rv_u=(x_u,y_u,z_u)=(1,0,-1)$$
        $$\rv_v=(x_v,y_v,z_v)=(0,1,0)$$
        代入公式可知
        $$\dr S=\sqrt{|\rv_u|^2|\rv_v|^2-(\rv_u\cdot\rv_v)^2}\dr u\dr v=\sqrt2\dr u\dr v$$
        由此原积分为
        $$\sqrt2\big(-2\sqrt2\big)\iint_D\dr u\dr v$$

        \item \textbf{积分计算}
        
        利用重积分的几何意义，$\iint_D\dr u\dr v$即为区域的面积，而此区域为单位圆，面积为1，由此最终得到积分结果为$-4\pi$。
    \end{itemize}
}
\end{framed}

\note 虽然本题进行了第二型曲线积分到第二型曲面积分到第一型曲面积分到重积分的转化，每一步转化本质上都是为了\textbf{化简问题}，\textbf{减少计算量}。事实上，完全可以在消除原本被积函数$\sin$相关项后直接计算第二型曲线积分，或是直接计算第二型曲面积分。

\begin{framed}
\begin{exercise}
    计算
    $$\oint_{L^+}\frac{\er^y}{x^2+y^2}\big((x\sin x+y\cos x)\dr x+(y\sin x-x\cos x)\dr y\big)$$
    其中$L$为椭圆周$\frac{x^2}{2}+y^2=1$。
\end{exercise}
\sol{
    分为如下步骤：
    \begin{itemize}
        \item \textbf{格林公式}
        
        将被积函数看作$P\dr x+Q\dr y$，由初等函数知$P$、$Q$在$\rb^2$去除原点外的区域内可导且导函数连续，直接计算得$Q_x$、$P_y$均为
        $$\frac{\er^y}{(x^2+y^2)^2}\big((x^2y+y^3+x^2-y^2)\cos x+(x^3+xy^2-2xy)\sin x\big)$$
        由此作差为0。不过，此区域并非单连通，无法得到积分为0。我们需要利用多连通区域的格林公式将积分曲线\textbf{转化为更好计算的曲线}。

        为了化简分母，考虑曲线
        $$L_r\colon x^2+y^2=r^2$$
        我们取一个足够小的$r>0$使得$L_0$包含在$L$中，设$L_0$与$L$之间的区域为$D$，利用\textbf{多连通区域的格林公式}有
        $$\oint_{L^+}(P\dr x+Q\dr y)-\oint_{L_r^+}(P\dr x+Q\dr y)=\iint_D(Q_x-P_y)\dr x\dr y=0$$
        由此积分转化为
        $$\oint_{L_r^+}\frac{\er^y}{x^2+y^2}\big((x\sin x+y\cos x)\dr x+(y\sin x-x\cos x)\dr y\big)$$

        \item \textbf{积分化简}
        
        由于$L_r$上满足$x^2+y^2=r^2$，我们可以将$x^2+y^2$改写为$r^2$以化简积分成
        $$\frac{1}{r^2}\oint_{L_r^+}\er^y\big((x\sin x+y\cos x)\dr x+(y\sin x-x\cos x)\dr y\big)$$

        \item \textbf{格林公式}
        
        由于此积分仍然难以直接看出计算方式，考虑重新应用格林公式。将新的被积函数看作$P_0\dr x+Q_0\dr y$，由初等函数知$P_0$、$Q_0$在$L_r$围成的单连通区域$D_r$内可导且导函数连续，进一步计算验证可知
        $$(Q_0)_x-(P_0)_y=\er^y(y\cos x-\cos x+x\sin x-(x\sin x+y\cos x+\cos x))=-2\er^y\cos x$$
        积分最终化为
        $$-\frac{2}{r^2}\iint_{D_r}\er^y\cos x\dr x\dr y$$

        \item \textbf{极限计算}
        
        直接换元计算会发现，上述的重积分虽然看起来简单，实际是非常难计算的。不过，一个重要的观察是，由于任取充分小的$r$都可以推出相同的式子，上式对任何$r$应当\textbf{相等}。也即记
        $$I(r)=-\frac{2}{r^2}\iint_{D_r}\er^y\cos x\dr x\dr y$$
        存在$\varepsilon>0$使得$0<r<\varepsilon$时$I(r)$为常数。由此，我们只需计算
        $$\lim_{r\to0^+}I(r)$$
        它必然与$I(r)$的值相等。

        \note 这样的\textbf{极限思路}事实上也是高等数学更一般的重要应用之一，也即通过极限将不精确的估算转化为精确的计算结果。

        由于$D_r$是满足$x^2+y^2\le r^2$的区域，在$r$接近0时应接近$(0,0)$，而被积函数在$(0,0)$处的值为1，且连续，因此$(0,0)$附近应接近1，这就得到积分结果约为$D_r$的面积$\pi r^2$，从而极限应当是
        $$-\frac{2}{r^2}\cdot\pi r^2=-2\pi$$
        这样就算出了形式上的结果。

        \item (\extra)\textbf{严谨性分析}
        
        为了证明上述过程的严谨性，我们需要使用重积分的\textbf{积分中值定理}(可见教材7.1节)，也即，由于$\er^y\cos x$是连续函数，存在$(x_0(r),y_0(r))\in D_r$使得
        $$\iint_{D_r}\er^y\cos x\dr x\dr y=\er^{y_0(r)}\cos(x_0(r))\pi r^2$$
        由此可知
        $$I(r)=-2\pi\er^{y_0(r)}\cos(x_0(r))$$
        利用$(x_0(r),y_0(r))\in D_r$可知两者平方和不超过$r^2$，从而绝对值不超过$r$，由\textbf{夹逼定理}得
        $$\lim_{r\to0^+}x_0(r)=\lim_{r\to0^+}y_0(r)=0$$
        由此可利用\textbf{复合函数极限}结论(可见上学期讲义14.2.2的``连续情况'')得到
        $$\lim_{r\to0^+}I(r)=\lim_{r\to0^+}-2\pi\er^{y_0(r)}\cos(x_0(r))=-2\pi$$
    \end{itemize}

    \note 本题也可以先用格林公式说明$(1+\sin x^2)\dr x+(1+\sin y^2)\dr y$在封闭曲线的上的积分为0，再作差后处理剩余部分。
}
\end{framed}

\note 只有\textbf{化为第一型曲线/曲面积分或重积分后}才可以进行极限的估算。此外，这类题若处理后发现无法直接算出积分，就可能需要通过极限考虑了。

\

\noindent\textbf{抽象函数}
\begin{framed}
\begin{exercise}\label{green normal}
    计算
    $$\oint_L\frac{\cos(\rv,\nv)}{|\rv|}\dr s$$
    其中$L$为平面上一条分段光滑的简单闭曲线，$\nv$为$L$对应位置的单位外法向量，$\rv=(x,y)$，$(\rv,\nv)$代表两向量夹角。
\end{exercise}
\sol{
    分为如下步骤：
    \begin{itemize}
        \item \textbf{化为第二型曲线积分}
        
        利用向量内积的性质可知
        $$\cos(\rv,\nv)=\frac{\rv\cdot\nv}{|\rv||\nv|}$$
        代入$\rv$的具体形式，并利用单位向量$|\nv|=1$可将积分写为
        $$\oint_L\frac{1}{x^2+y^2}(x,y)\cdot\nv\dr s$$
        由于并不知道曲线的具体形式，难以继续计算，由此想到需要\textbf{转化为第二型曲线积分}。

        设$\vec\tau=(\tau_1,\tau_2)$代表曲线逆时针方向的单位切向量，利用第一型曲线积分与第二型曲线积分的关系可知(这里曲线以逆时针为正定向)
        $$\vec\tau\dr s=(\dr x,\dr y)$$
        也即
        $$(\tau_1\dr s,\tau_2\dr s)=(\dr x,\dr y)$$
        不过，题中为\textbf{法向量}的相关式子，而非切向量。利用几何关系可知，$L$的逆时针方向单位切向量为单位外法向量\textbf{向逆时针方向旋转90度}的结果。设单位外法向量$\nv=(n_1,n_2)$，计算可发现
        $$n_1=\tau_2,\quad n_2=-\tau_1$$
        由此有
        $$n_1\dr s=\dr y,\quad n_2\dr s=-\dr x$$
        由此最终写出第二型曲线积分
        $$\oint_{L^+}\frac{x\dr y-y\dr x}{x^2+y^2}$$

        \item \textbf{格林公式}
        
        将被积函数看作$P\dr x+Q\dr y$，由初等函数知$P$、$Q$在$\rb^2$去除原点外的区域内可导且导函数连续，直接计算得$Q_x$、$P_y$均为
        $$\frac{y^2-x^2}{(x^2+y^2)^2}$$
        由此作差为0。

        接下来分为两种情况：
        \begin{enumerate}
            \item 若$L$不包含原点，直接利用格林公式可知积分结果为0。
            \item 若$L$包含原点，考虑曲线
            $$L_0\colon x^2+y^2=r^2$$
            我们取一个足够小的$r>0$使得$L_0$包含在$L$中，设$L_0$与$L$之间的区域为$D$，利用多连通区域的格林公式有
            $$\oint_{L^+}(P\dr x+Q\dr y)-\oint_{L_0^+}(P\dr x+Q\dr y)=\iint_D(Q_x-P_y)\dr x\dr y=0$$
            由此积分转化为
            $$\oint_{L_0^+}\frac{x\dr y-y\dr x}{x^2+y^2}$$
            由$L_0$上$x^2+y^2=r^2$可重新将其写为
            $$\frac{1}{r^2}\oint_{L_0^+}(x\dr y-y\dr x)$$
            验证条件并再次利用格林公式，设$L_0$围成的区域是$D_0$，即得到积分化为
            $$\frac{2}{r^2}\iint_{D_0}\dr x\dr y$$
            利用重积分的几何意义，重积分部分即为$D_0$的面积$\pi r^2$，从而结果为$2\pi$。
        \end{enumerate}
    \end{itemize}
}
\end{framed}

\note \textbf{三维版本}的第一型曲线积分与第二型曲线积分关系即为
$$\vec\tau\dr s=(\dr x,\dr y,\dr z)$$
这里$\vec\tau$代表第二型曲线积分定向下的单位切向量。

\note 务必熟悉\textbf{平面曲线的切向量法向量关系}，但注意空间曲线并没有这样的性质(因为曲线的法向并不唯一)。

\begin{framed}
\begin{exercise}\label{type 2 symm}
    证明
    $$\oint_{L^+}xf(y)\dr y-\frac{y}{f(x)}\dr x\ge2\pi$$
    其中$L$为圆周$(x-1)^2+(y-1)^2=1$，$f(x)$为$\rb$上恒正的连续函数。
\end{exercise}
\sol{
    记要估算的曲线积分为$I$。本题我们将介绍三种方法，分别是有严谨性问题的\textbf{格林公式}，需要非常谨慎使用的\textbf{第二型曲线积分换元}，与最标准的\textbf{定积分换元}做法：
    \begin{itemize}
        \item \textbf{格林公式}
        
        本题实际上\textbf{无法使用格林公式}，因为格林公式要求被积函数的各一阶偏导存在连续(即使是公式中未出现的$Q_y$与$P_x$)，但$f$只有连续的条件时$xf(y)$对$y$不可导。

        不过，若考虑$f$可导的特殊情况，可以使用格林公式。在$f$可导时，将被积函数看作$P\dr x+Q\dr y$，直接计算知
        $$Q_x-P_y=f(y)+\frac{1}{f(x)}$$
        在$\rb^2$上有定义，从而设$L$围成的区域为$D$，由格林公式有
        $$I=\iint_D\bigg(f(y)+\frac{1}{f(x)}\bigg)\dr x\dr y$$
        可以想到，若被积函数是$f(x)+\frac{1}{f(x)}$，就方便进行估算了。而这确实是可以进行的：由于$D$满足$(x,y)\in D$当且仅当$(y,x)\in D$，利用对称性交换$x$、$y$可发现(严谨证明与定理形式类似\exref{swap})
        $$\iint_Df(y)\dr x\dr y=\iint_Df(x)\dr x\dr y$$
        这就进一步得到了
        $$I=\iint_D\bigg(f(x)+\frac{1}{f(x)}\bigg)\dr x\dr y$$
        由于$f(x)>0$，利用基本不等式可知
        $$f(x)+\frac{1}{f(x)}\ge2$$
        由\textbf{重积分性质}有
        $$I\ge\iint_D2\dr x\dr y=2\iint_D\dr x\dr y$$
        根据重积分的几何意义，右侧的积分即为$D$的面积，而$D$是半径为1的圆，面积为$\pi$，这就得到了
        $$I\ge2\pi$$

        \note 实际考试中用这样的做法如何扣分取决于评卷标准。
        
        \item \textbf{第二型曲线积分换元}
        
        虽然格林公式方法无法直接使用，但它给出了一个非常重要的\textbf{思路}：配凑$f(x)+\frac{1}{f(x)}$。那么，直接对第二型曲线积分能否处理出类似的结果呢？我们接下来说明对本题中的曲线$L$与连续函数$\phi(x,y)$有
        $$\oint_{L^+}\phi(x,y)\dr y=-\oint_{L^+}\phi(y,x)\dr x$$
        在\exref{swap}中我们已经讨论过，交换$x$、$y$本质是一种反射对称。关于直线$y=x$反射对称后，原本的函数$\phi(x,y)$会变成$\phi(y,x)$，微元$\dr y$会变成$\dr x$，这就得到了被积函数的形式。另一方面，反射对称后，原本的逆时针会变成顺时针，由此想回到原本的方向需要添加负号。

        \note 严谨证明可以通过极坐标参数化后写出两边并进行定积分换元进行。

        利用此结论可直接得到
        $$\oint_{L^+}-\frac{y}{f(x)}\dr x=\oint_{L^+}\frac{x}{f(y)}\dr y$$
        由此
        $$I=\oint_{L^+}x\bigg(f(y)+\frac{1}{f(y)}\bigg)\dr y$$

        \note 注意即使曲线上$x\ge0$，此时也无法直接把$f(y)+\frac{1}{f(y)}$放缩为2，第二型曲线积分\textbf{并不具有被积函数越大积分越大的性质}。

        为了进行估算，我们仍需要写出定积分。利用非原点中心极坐标可以得到圆的参数方程
        $$\begin{cases}x=1+\cos\theta\\y=1+\sin\theta\end{cases}(0\le\theta\le2\pi)$$
        极坐标已经为逆时针定向，无需结果添加负号。直接计算有
        $$x_\theta=-\sin\theta,\quad y_\theta=\cos\theta$$    
        代入$\dr x=x_\theta\dr\theta$、$\dr y=y_\theta\dr\theta$与$x$、$y$的参数方程可得
        $$I=\int_0^{2\pi}(\cos\theta+\cos^2\theta)\bigg(f(1+\sin\theta)+\frac{1}{f(1+\sin\theta)}\bigg)\dr\theta$$
        注意到这步仍然无法直接放缩，因为$\cos\theta+\cos^2\theta$的符号无法保证。不过，\textbf{仔细观察}可以发现实际上有(由于$\sin0=\sin(2\pi)$，换元$t=\sin\theta$后是0到0积分，因此为0)
        $$\int_0^{2\pi}\cos\theta\bigg(f(1+\sin\theta)+\frac{1}{f(1+\sin\theta)}\bigg)\dr\theta=\int_0^{2\pi}\bigg(f(1+\sin\theta)+\frac{1}{f(1+\sin\theta)}\bigg)\dr\sin\theta=0$$

        \note 也可根据\textbf{对称性}分析，将$\theta$换元为$\pi-\theta$后考虑相互抵消的部分得到相同的结论。

        去除这一部分后，我们得到了
        $$I=\int_0^{2\pi}\cos^2\theta\bigg(f(1+\sin\theta)+\frac{1}{f(1+\sin\theta)}\bigg)\dr\theta$$
        此时被积函数非负，终于可以放缩出(计算与\exref{sin 2 and 4 to 2pi}完全类似)
        $$I\ge\int_0^{2\pi}2\cos^2\theta\dr\theta=2\pi$$

        \item \textbf{定积分换元}
        
        现在我们考虑最严谨的方法，直接写出与上一种方法相同的参数方程，并得到对应的定积分
        $$I=\int_0^{2\pi}\bigg((\cos\theta+\cos^2\theta)f(1+\sin\theta)+(\sin\theta+\sin^2\theta)\frac{1}{f(1+\cos\theta)}\bigg)\dr\theta$$
        为了配凑出$f(x)+\frac{1}{f(x)}$的形式，我们将$\theta$换元为$\frac{\pi}{2}-\theta$可以发现
        $$\int_0^{2\pi}(\sin\theta+\sin^2\theta)\frac{1}{f(1+\cos\theta)}\dr\theta=\int_{-3\pi/2}^{\pi/2}(\cos\theta+\cos^2\theta)\frac{1}{f(1+\sin\theta)}\dr\theta$$
        由于等号右侧的被积函数是以$2\pi$为周期的，利用上学期结论(可见上册教材3.4节)知它在任何一个周期积分相等，从而可以将积分移到$[0,2\pi]$，得到
        $$\int_0^{2\pi}(\sin\theta+\sin^2\theta)\frac{1}{f(1+\cos\theta)}\dr\theta=\int_0^{2\pi}(\cos\theta+\cos^2\theta)\frac{1}{f(1+\sin\theta)}\dr\theta$$
        将这部分与原本合并即得
        $$I=\int_0^{2\pi}\bigg((\cos\theta+\cos^2\theta)f(1+\sin\theta)+(\cos\theta+\cos^2\theta)\frac{1}{f(1+\sin\theta)}\bigg)\dr\theta$$
        此式就是上一种方法里参数化后得到的定积分，接下来的过程与上一种方法相同。
    \end{itemize}

    \note 考虑到其他方法的\textbf{复杂程度}，这类问题使用稍不严谨的格林公式方法(也即只要$Q_x$、$P_y$保证连续，就\textbf{形式上}使用格林公式)获得过程分或许是明智的选择。
}
\end{framed}

\note 注意定积分换元\textbf{无需单调}，而重积分\textbf{需要一一映射}。

\note 注意被积函数的大小关系推出积分的大小关系\textbf{不能对第二型曲线/曲面积分应用}，因为考虑方向时情况将变得非常复杂。

\subsubsection{曲面积分}
\noindent\textbf{第一型曲面积分}
\begin{framed}
\begin{exercise}
    计算
    $$\iint_Mx\dr S$$
    其中
    $$M=\{(x,y,z)\mid x^2+z^2=1,\ x^2+y^2\le1,\ x\ge0,\ y\ge0,\ z\ge0\}$$
\end{exercise}
\sol{
    分为如下步骤：
    \begin{itemize}
        \item \textbf{参数化}
        
        由于$x^2+z^2=1$可以自然用极坐标参数化$x$与$z$，而$y$可独立变化，从而能写出参数化
        $$\rv=(x,y,z)=(\cos u,v,\sin u)$$

        \note 这是\textbf{柱面}的通常参数化方式，写出关联的两个变量的参数化，第三个变量单独考虑。

        $x\ge0$、$z\ge0$对应极坐标中的第一象限，由此可写出区域的一个约束
        $$0\le u\le\frac{\pi}{2}$$
        剩下的约束代入即可得到
        $$v\ge0,\quad\cos^2u+v^2\le1$$
        第二式移项得到$v^2\le\sin^2u$，利用$\sin u$与$v$均正可以开平方根，最终将约束写为
        $$0\le u\le\frac{\pi}{2},\quad 0\le v\le\sin u$$
        将上述条件确定的$(u,v)$区域记为$D$。

        \item \textbf{化为重积分}
        
        用下标$u$、$v$表示对$u$、$v$求偏导并保持另一个不变，直接计算有
        $$\rv_u=(x_u,y_u,z_u)=(-\sin u,0,\cos u),\quad\rv_v=(x_v,y_v,z_v)=(0,1,0)$$
        代入公式可知
        $$\dr S=\sqrt{|\rv_u|^2|\rv_v|^2-(\rv_u\cdot\rv_v)^2}\dr u\dr v=\dr u\dr v$$
        被积函数即为$\cos u$，由此原积分为
        $$\iint_D\cos u\dr u\dr v$$

        \item \textbf{积分计算}
        
        根据区域描述可以直接写出累次积分
        $$\int_0^{\pi/2}\dr u\int_0^{\sin u}\cos u\dr v$$
        对$v$积分得结果为(换元$t=\sin u$)
        $$\int_0^{\pi/2}\sin u\cos u\dr u=\int_0^{\pi/2}\sin u\dr(\sin u)=\int_0^1t\dr t=\frac{1}{2}$$
    \end{itemize}
}
\end{framed}

\note 需要掌握一些空间曲面的\textbf{常用参数化方式}，最基本的方式是将一个变量看作另外两个的函数，不过善用球坐标、柱坐标可能会有更简单的形式。

\note 第一型曲面积分通常没有比直接参数化计算更好的方法。

\begin{framed}
\begin{exercise}
    计算
    $$\iint_M(x^2y^2+y^2z^2+z^2x^2)\dr S$$
    其中$M$是锥面$z=\sqrt{x^2+y^2}$被柱面$x^2+y^2=1$截下的部分。
\end{exercise}
\sol{
    \begin{itemize}
        \item \textbf{区域描述与参数化}
            
        根据``割下''的有界性要求，区域应为曲面$z=\sqrt{x^2+y^2}$中满足$x^2+y^2\le1$的部分(大于等于的部分无界)。

        虽然$z$是$x$、$y$的函数可以提供一个自然的参数化，由于曲面的旋转对称性，为了让计算更简洁可以考虑\textbf{极坐标}参数化
        $$\rv=(x,y,z)=(u\cos v,u\sin v,u)$$
        代入约束可得到参数范围
        $$u^2\le1$$
        将上述条件与极坐标自然范围确定的$(u,v)$区域记为$D$，可发现约束能写为
        $$0\le u\le1,\quad0\le v\le2\pi$$

        \item \textbf{化为重积分}
        
        用下标$u$、$v$表示对$u$、$v$求偏导并保持另一个不变，直接计算有
        $$\rv_u=(x_u,y_u,z_u)=(\cos v,\sin v,1),\quad\rv_v=(x_v,y_v,z_v)=(-u\sin v,u\cos v,0)$$
        代入公式可知
        $$\dr S=\sqrt{|\rv_u|^2|\rv_v|^2-(\rv_u\cdot\rv_v)^2}\dr u\dr v=\sqrt2u\dr u\dr v$$
        被积函数即为$u^4(1+\cos^2v\sin^2v)$，由此原积分为
        $$\iint_Du^5(1+\cos^2v\sin^2v)\sqrt2\dr u\dr v$$

        \item \textbf{积分计算}
        
        根据区域描述可以直接写出累次积分
        $$\sqrt2\int_0^1\dr u\int_0^{2\pi}u^5(1+\cos^2v\sin^2v)\dr v$$
        利用\exref{sin 2 and 4 to 2pi}对$I_0$的积分计算结果可知对$v$积分后为
        $$\frac{9\sqrt2\pi}{4}\int_0^1u^5\dr u=\frac{3\sqrt2\pi}{8}$$
    \end{itemize}
}
\end{framed}

\note 注意``截下''对应的区域确定方式与``围成''是完全类似的。

\note 若曲面由多个部分拼接而成，对应的曲面积分应\textbf{分别参数化并计算每个部分}，考试中这类题目往往大多数部分是容易计算的。

\

\noindent\textbf{第二型曲面积分}
\begin{framed}
\begin{exercise}
    计算
    $$\iint_Sx\dr y\dr z+y\dr z\dr x+z\dr x\dr y$$
    其中$S$为抛物面$z=x^2+y^2$在$z\le4$部分的下侧。
\end{exercise}
\sol{
    我们先介绍\textbf{直接计算}的方法：
    \begin{itemize}
        \item \textbf{参数化}
        
        $z$是$x$、$y$的函数可以提供一个自然的参数化
        $$\rv=(x,y,z)=(u,v,u^2+v^2)$$
        代入约束可得到参数范围
        $$u^2+v^2\le4$$
        将上述条件确定的$(u,v)$区域记为$D$。

        \item \textbf{定向}
        
        为了确定定向， 用下标$u$、$v$表示对$u$、$v$求偏导并保持另一个不变，直接计算有
        $$\rv_u=(x_u,y_u,z_u)=(1,0,2u),\quad\rv_v=(x_v,y_v,z_v)=(0,1,2v)$$
        我们\textbf{任取}一个使得$\rv_u$、$\rv_v$\textbf{不平行}(也可表述为线性无关/叉乘非零)的$D$中点，如取$(u,v)=(0,0)$，此处有$\rv=(0,0,0)$、$\rv_u=(1,0,0)$、$\rv_v=(0,1,0)$，在曲面上的点$\rv$处画出向量$\rv_u$与$\rv_v$，通过右手定则可发现$\rv_u\times\rv_v$指向上方，与要求的定向相反，因此\textbf{需要添加负号}。

        \note 事实上还有更简单的无需画图的判断方法：下侧代表取$z$分量小于0的法向量，直接计算可发现$\rv_u\times\rv_v$的$z$分量是1，大于0，因此与要求的定向相反。

        \item \textbf{化为重积分}
        
        代入
        $$\dr y\dr z=\frac{D(y,z)}{D(u,v)}\dr u\dr v,\quad\dr z\dr x=\frac{D(z,x)}{D(u,v)}\dr u\dr v,\quad\dr x\dr y=\frac{D(x,y)}{D(u,v)}\dr u\dr v$$
        与$x$、$y$、$z$的参数方程可得原积分为(最左侧的负号来自定向)
        $$-\iint_D\big(u(-2u)+v(-2v)+(u^2+v^2)\big)\dr u\dr v$$
        化简得
        $$\iint_D(u^2+v^2)\dr u\dr v$$

        \item \textbf{积分计算}
        
        考虑到区域与函数的旋转对称性，采用极坐标换元
        $$u=r\cos\theta,\quad v=r\sin\theta$$
        此时有
        $$\dr u\dr v=r\dr r\dr\theta$$
        约束化简为
        $$0\le r\le2$$
        被积函数为$r^2$，再结合极坐标自然范围可将积分化为
        $$\int_0^2\dr r\int_0^{2\pi}r^3\dr\theta$$
        对$\theta$积分得到结果为
        $$2\pi\int_0^2r^3\dr r=8\pi$$
    \end{itemize}
    当然，本题的形式也可以通过\textbf{高斯公式}计算：
    \begin{itemize}
        \item \textbf{高斯公式}
        
        将被积函数看作$P\dr y\dr z+Q\dr z\dr x+R\dr x\dr y$，由初等函数知$P$、$Q$、$R$在$\rb^3$上可导且导函数连续，直接计算得
        $$P_x+Q_y+R_z=3$$
        由于积分区域并非封闭曲面，我们补上顶部的
        $$S_0:\quad z=4,\quad x^2+y^2\le z$$
        并取上侧为正向。

        根据定义$S_0$在$S$的上方，因此$S_0$的上侧、$S$的下侧构成两者围成的封闭曲面的外侧。设围成的区域为$\Omega$，由高斯公式可知
        $$\iint_S(x\dr y\dr z+y\dr z\dr x+z\dr x\dr y)+\iint_{S_0}(x\dr y\dr z+y\dr z\dr x+z\dr x\dr y)=3\iiint_\Omega\dr V$$
        由此设
        $$I_1=\iint_{S_0}(x\dr y\dr z+y\dr z\dr x+z\dr x\dr y),\quad I_2=\iiint_\Omega\dr V$$
        结果即为$3I_2-I_1$。

        \item \textbf{$I_1$计算}
        
        由于$I_1$是平面区域上的第二类曲面积分，化为第一类可方便计算。利用第一型曲面积分与第二型曲面积分的关系可知(这里$\nv$代表第二型曲面积分对应定向的单位法向量)
        $$\nv\dr S=(\dr y\dr z,\dr z\dr x,\dr x\dr y)$$   
        根据几何可知$z=4$平面的一个上侧的单位法向量是$(0,0,1)$，由此即得原积分为(已代入$z=4$)
        $$\iint_{S_0}(x,y,z)\cdot(\dr y\dr z,\dr z\dr x,\dr x\dr y)=\iint_{S_0}(x,y,4)\cdot(0,0,1)\dr S=4\iint_{S_0}\dr S$$
        利用第一型曲面积分的几何意义，$\iint_{S_0}\dr S$即为曲面面积，而$S_0$的约束为$z=4$、$x^2+y^2\le4$，是一个半径为2的圆，从而面积为$4\pi$，最终得到
        $$I_1=16\pi$$
        
        \item \textbf{$I_2$计算}
        
        根据定义可发现$\Omega$为区域
        $$x^2+y^2\le z\le 4$$
        由于约束的旋转对称性，进行柱坐标换元
        $$x=r\cos\theta,\quad y=r\sin\theta,\quad z=z$$
        这样约束条件即化为了
        $$r^2\le z\le4$$
        直接计算Jacobi行列式可得
        而被积函数仍为1。由于$z$自然被$r$约束，且约束与被积函数关于$\theta$均为常数，考虑先$\theta$、再$z$、最后$r$的积分顺序。
        
        在$\theta$、$z$任取时，$r$的约束即$r^2\le4$，结合自然范围得到
        $$0\le r\le2$$
        再根据约束与自然范围有
        $$r^2\le z\le4,\quad0\le\theta\le2\pi$$
        这就可以写出累次积分
        $$\int_0^2\dr r\int_{r^2}^4\dr z\int_0^{2\pi}r\dr\theta$$
        依次积分得到结果为
        $$2\pi\int_0^2\dr r\int_{r^2}^4r\dr z=2\pi\int_0^2(4-r^2)r\dr r=8\pi$$

        \item \textbf{综合结果}

        综合以上最终可知积分为
        $$3I_2-I_1=8\pi$$
    \end{itemize}
}
\end{framed}

\note 必须掌握一种\textbf{曲面积分确定参数定向}的方式，本题的定向过程是个人推荐的方法：任取一组满足
$$\rv_u\times\rv_v\ne\vec{0}$$
的$u$与$v$\ (否则\textbf{无法确定法向量})，在曲线上画出这点与对应的法向量$\rv_u\times\rv_v$\ (也可以不计算，直接通过叉乘的右手定则确定方向)，观察指向哪一侧。

\note 当算出散度非零且曲面并非闭曲面时，高斯公式可能\textbf{并不能简化计算}。

\begin{framed}
\begin{exercise}
    计算
    $$\oiint_{S^-}xy\dr y\dr z+\big(y^2+\er^{xz^2}\big)\dr z\dr x+\sin(xy)\dr x\dr y$$
    其中$S$为抛物柱面$z=1-x^2$与平面$z=0$、$y=0$、$y+z=2$围成区域的表面，上标负号代表内侧。
\end{exercise}
\sol{
    分为如下步骤：
    \begin{itemize}
        \item \textbf{高斯公式}
        
        将被积函数看作$P\dr y\dr z+Q\dr z\dr x+R\dr x\dr y$，由初等函数知$P$、$Q$、$R$在$\rb^3$上可导且导函数连续，直接计算得
        $$P_x+Q_y+R_z=3y$$
        由于高斯公式要求外侧定向，使用内侧定向应相差负号，设$S$围成的区域为$\Omega$，利用高斯公式即可将积分化为
        $$-\iiint_\Omega3y\dr x\dr y\dr z$$
        利用``围成''的有界性可知$\Omega$为满足
        $$z\le1-x^2,\quad z\ge0,\quad y\ge0,\quad y+z\le2$$
        的$(x,y,z)$构成的集合。

        \note 更具体的代数分析：关于$x$的约束只有$z=1-x^2$，若取$z\ge1-x^2$则$x$可无穷增大，与有界矛盾，由此必须$z\le1-x^2$。若$y+z\ge2$，为使$y$有上界必须$y\le0$，但此时必然有$z\ge2$，与$z\le1-x^2$矛盾，由此必须$y+z\le2$。最终，为使$y$、$z$有下界，剩下两个约束必然为$y\ge0$、$z\ge0$。

        \note 实际上$y=0$、$z=0$、$y+z=2$构成一个三棱柱，为使区域有界几乎必然会取内部，也即$y\ge0$、$z\ge0$、$y+z\le2$，这样可以更容易得到区域。

        \item \textbf{积分计算}
        
        由于$z$自然被$x$约束，我们考虑先对$y$、再对$z$，最后对$x$的积分顺序。

        当$y$、$z$可以任取时，$x$的约束即$x^2+z\le1$，$z=0$时取到最值，范围是
        $$-1\le x\le1$$
        给定$x$、$y$任取时，$z$的范围即$z\ge0$、$z\le2-y$、$z\le1-x^2$。由于$y=0$可使第二个约束为$z\le2$，弱于$z\le1-x^2$，$z$的实际范围是
        $$0\le z\le1-x^2$$
        给定$x$、$z$后，$y$的范围即
        $$0\le y\le 2-z$$
        综合得到累次积分
        $$-3\int_{-1}^1\dr x\int_0^{1-x^2}\dr z\int_0^{2-z}y\dr y$$
        对$y$积分得到其为
        $$-\frac{3}{2}\int_{-1}^1\dr x\int_0^{1-x^2}(2-z)^2\dr z$$
        再对$z$积分即
        $$-\frac{1}{2}\int_{-1}^1(8-(1+x^2)^3)\dr x$$
        利用对称性可化为(也可直接计算)
        $$-\int_0^1(8-(1+x^2)^3)\dr x$$
        最终展开算出积分结果
        $$-\bigg(7-1-\frac{3}{5}-\frac{1}{7}\bigg)=-\frac{184}{35}$$

        \note 也可以选择$z$、$x$、$y$的顺序，得到的累次积分为
        $$-3\int_0^2\dr y\int_{-1}^1\dr x\int_0^{\min\{1-x^2,2-y\}}y\dr z$$
        积分过程中对$x$积分时需要讨论$1-x^2$与$2-y$何者更小。选择$z$、$y$、$x$的顺序做法也类似。
    \end{itemize}
}
\end{framed}

\note 高斯公式化为三重积分后，仍然可能需要\textbf{重积分计算技巧}才能得到正确的结果，但规避定向和曲面参数化的确可以简化计算。

\begin{framed}
\begin{exercise}
    计算
    $$\oiint_{S^+}\frac{x\dr y\dr z+y\dr z\dr x+z\dr x\dr y}{(x^2+4y^2+9z^2)^{3/2}}$$
    其中$S$为球面$x^2+y^2+z^2=1$。
\end{exercise}
\sol{
    分为如下步骤：
    \begin{itemize}
        \item \textbf{高斯公式}
        
        将被积函数看作$P\dr y\dr z+Q\dr z\dr x+R\dr x\dr y$，由初等函数知$P$、$Q$、$R$在$\rb^3$去除原点外的区域内可导且导函数连续，直接计算得$P_x+Q_y+R_z$为
        $$\frac{3}{(x^2+4y^2+9z^2)^{3/2}}-\frac{3x^2}{(x^2+4y^2+9z^2)^{5/2}}-\frac{12y^2}{(x^2+4y^2+9z^2)^{5/2}}-\frac{27z^2}{(x^2+4y^2+9z^2)^{5/2}}=0$$
        不过，此区域并非空间单连通，无法得到积分为0。我们需要利用空间多连通区域的高斯公式将积分曲面\textbf{转化为更好计算的曲面}。

        为了化简分母，考虑曲面
        $$S_0\colon x^2+4y^2+9z^2=r^2$$
        我们取一个足够小的$r>0$使得$S_0$包含在$S$中，设$S_0$与$S$之间的区域为$\Omega$，利用\textbf{空间多连通区域的高斯公式}有
        $$\begin{aligned}&\oiint_{S^+}(P\dr y\dr z+Q\dr z\dr x+R\dr x\dr y)-\oiint_{S_0^+}(P\dr y\dr z+Q\dr z\dr x+R\dr x\dr y)\\=&\iiint_\Omega(P_x+Q_y+R_z)\dr V=0\end{aligned}$$
        由此将积分转化为
        $$\oiint_{S_0^+}\frac{x\dr y\dr z+y\dr z\dr x+z\dr x\dr y}{(x^2+4y^2+9z^2)^{3/2}}$$

        \note 与\exref{green out}情况相同，$r$的存在性论证相对麻烦，可以无需严格论证，而保留$r$往往比写出具体的$r$更易于计算。

        \item \textbf{积分化简}
        
        由于$S_0$上满足$x^2+4y^2+9z^2=r^2$，我们可以将分母改写为$r^3$以化简积分成
        $$\frac{1}{r^3}\oiint_{S_0^+}x\dr y\dr z+y\dr z\dr x+z\dr x\dr y$$
        
        \item \textbf{高斯公式}
        
        虽然至此已经可以直接计算，我们还可以选择更简单的方式，也即\textbf{再用一次高斯公式}。将新的被积函数看作
        $$P_0\dr y\dr z+Q_0\dr z\dr x+R_0\dr x\dr y$$
        由初等函数知$P_0$、$Q_0$、$R_0$在$S_0$围成的区域$\Omega_0$内可导且导函数连续，进一步计算验证可知
        $$(P_0)_x+(Q_0)_y+(R_0)_z=3$$
        积分最终化为
        $$\frac{3}{r^3}\iiint_{\Omega_0}\dr x\dr y\dr z$$

        \item \textbf{积分计算}

        利用重积分的几何意义，$\iiint_{\Omega_0}\dr x\dr y\dr z$代表$\Omega_0$的体积，而此积分区域是椭球$x^2+4y^2+9z^2=r^2$，三个轴分别长$r$、$\frac{r}{2}$、$\frac{r}{3}$，从而体积为
        $$\frac{4}{3}\pi\cdot r\cdot\frac{r}{2}\cdot\frac{r}{3}$$
        综合以上最终得到积分结果
        $$\frac{2\pi}{3}$$
    \end{itemize}
}
\end{framed}

\note 由于对1进行第一型曲线积分、第一型曲面积分与重积分几何意义为曲线长度/曲面面积、区域面积或体积，建议\textbf{熟悉一些常用的周长/表面积/体积公式}。

\note 高斯公式可以与格林公式有完全类似的应用。

\begin{framed}
\begin{exercise}\label{column normal}
    计算
    $$I_1=\oiint_{S^+}(y-z)|x|\dr y\dr z+(z-x)|y|\dr z\dr x+(x-y)|z|\dr x\dr y$$
    $$I_2=\oiint_{S^+}(y-z)x^2\dr y\dr z+(z-x)y^2\dr z\dr x+(x-y)z^2\dr x\dr y$$
    $$I_3=\oiint_{S^+}(y-z)x^3\dr y\dr z+(z-x)y^3\dr z\dr x+(x-y)z^3\dr x\dr y$$
    其中$S$为柱体
    $$\Omega=\{(x,y,z)\mid x^2+y^2\le1,\ 0\le z\le1\}$$
    的表面。
\end{exercise}
\sol{
    由于含绝对值的项未必可导，$I_1$无法通过高斯公式直接计算。当然，我们仍然可以直接进行参数化后计算$I_1$、利用高斯公式计算$I_2$与$I_3$，不过，这里介绍一个\textbf{利用对称性直接得到结果}的方式，分为如下步骤：
    \begin{itemize}
        \item \textbf{区域描述}
        
        根据表面的要求，我们将三个约束中的每个分别取等，其他两个保留，就得到了$\Omega$的表面分为三个部分
        $$S_1:\quad z=0,\quad x^2+y^2\le1,\quad z\le1$$
        $$S_2:\quad z=1,\quad x^2+y^2\le1,\quad z\ge0$$
        $$S_3:\quad x^2+y^2=1,\quad 0\le z\le1$$
        可进一步化简为
        $$S_1:\quad z=0,\quad x^2+y^2\le1$$
        $$S_2:\quad z=1,\quad x^2+y^2\le1$$
        $$S_3:\quad x^2+y^2=1,\quad 0\le z\le1$$
        我们称这三个面为圆柱的底面、顶面、侧面。设积分$I_i$在$S_j$上部分的积分为$I_{ij}$，有(分析类似\exref{three column s})
        $$I_1=I_{11}+I_{12}+I_{13},\quad I_2=I_{21}+I_{22}+I_{23},\quad I_3=I_{31}+I_{32}+I_{33}$$
        下面按照$S_1$、$S_2$、$S_3$的顺序进行计算。
        
        \item \textbf{底面积分分析}
        
        为了能利用对称性，我们需要\textbf{先将第二型曲面积分转化为第一型曲面积分或重积分}。我们这里采用转化为第一型曲面的思路。底面$z=0$，对应圆柱的外侧应取向下的法向量，利用平面的几何性质可知其为$(0,0,-1)$，类似\exref{stolkes circle}的过程可化出
        $$I_{11}=-\iint_{S_1}(x-y)|z|\dr S$$
        $$I_{21}=-\iint_{S_1}(x-y)z^2\dr S$$
        $$I_{31}=-\iint_{S_1}(x-y)z^3\dr S$$
        由于$S_1$上$z$恒为0，这三个积分均为0。

        \item \textbf{顶面积分分析}
        
        仍然采用转化为第一型曲面的思路。底面$z=1$，对应圆柱的外侧应取向上的法向量，利用平面的几何性质可知其为$(0,0,1)$，类似\exref{stolkes circle}的过程可化出
        $$I_{12}=\iint_{S_2}(x-y)|z|\dr S$$
        $$I_{22}=\iint_{S_2}(x-y)z^2\dr S$$
        $$I_{32}=\iint_{S_2}(x-y)z^3\dr S$$
        由于$S_2$上$z$恒为1，这三个积分均为
        $$\iint_{S_2}(x-y)\dr x\dr y$$
        由于$S_2$为$z=1$、$x^2+y^2\le1$，满足$(x,y,z)\in S_2$当且仅当$(y,x,z)\in S_2$，与三重积分类似利用对称性(\exref{swap})可知
        $$\iint_{S_2}x\dr S=\iint_{S_2}y\dr S$$
        由此得到$x-y$在$S_2$上的积分为0。
        
        \item \textbf{侧面积分分析}
        
        仍然采用转化为第一型曲面的思路。类似\exref{stolkes circle}的过程，我们要找侧面
        $$x^2+y^2=1,\quad0\le z\le1$$
        在$(x,y,z)$处的外法向量。

        这里我们采用几何方式分析得到结果：由于$z$方向它是柱面，从几何来看法向量的$z$分量应为0，而平面上的圆$x^2+y^2=1$从几何可发现$(x,y)$就是外法向，从而可以得到柱面的外法向量
        $$(x,y,0)$$
        将其单位化得到
        $$\vec{n}=\frac{1}{\sqrt{x^2+y^2}}(x,y,0)$$

        利用它可以得到
        $$\begin{aligned}I_{13}&=\iint_{S_3}((y-z)|x|\dr y\dr z+(z-x)|y|\dr z\dr x+(x-y)|z|\dr x\dr y)\\ &=\iint_{S_3}((y-z)|x|,(z-x)|y|,(x-y)|z|)\cdot(\dr y\dr z,\dr z\dr x,\dr x\dr y)\\ &=\iint_{S_3}((y-z)|x|,(z-x)|y|,(x-y)|z|)\cdot\vec{n}\dr S\\ &=\iint_{S_3}\frac{(y-z)|x|x+(z-x)|y|y}{\sqrt{x^2+y^2}}\dr S\end{aligned}$$
        利用曲面定义可知$(x,y,z)\in S_3$当且仅当$(y,x,z)\in S_3$，从而利用对称性可知
        $$\iint_{S_3}\frac{(y-z)|x|x}{\sqrt{x^2+y^2}}\dr S=\iint_{S_3}\frac{(x-z)|y|y}{\sqrt{x^2+y^2}}\dr S$$
        由此将上式移项即可发现$I_{13}=0$，完全同理可得$I_{23}=I_{33}=0$。
        
        \item \textbf{综合结果}
        
        综合以上讨论，由于每个积分在表面每部分都是0，我们最终得到
        $$I_1=I_2=I_3=0$$
    \end{itemize}

    \note 对$I_2$使用高斯公式可得
    $$I_2=\iiint_\Omega2(xy-xz+yz-xy+zx-zy)\dr V=\iiint_\Omega0\dr V=0$$
    而对$I_3$使用高斯公式将化出
    $$I_3=\iiint_\Omega3(x^2y-zx^2+zy^2-xy^2+xz^2-yz^2)\dr V$$
    利用$\Omega$的对称性交换$x$、$y$即能得到其为0。因此，对于可以使用高斯公式的$I_2$、$I_3$，使用高斯公式后仍然可以分析对称性得到正确的结果。对于$I_1$，事实上按照$x$、$y$分为四个象限使用高斯公式(注意$z\ge0$，绝对值可去除)后也能得到正确的结果。

    \note 理论来说可以类似\exref{type 2 symm}的第二种方法，直接在第二型曲面积分时将$x$、$y$交换，但相关的分析较为复杂且不完全符合几何直观(形式上直接交换后变为选择内法向，再次添加负号可发现恰好为原积分的相反数，由此只能为0)，因此这里不介绍。
}
\end{framed}

\note 第一型曲线/曲面积分的\textbf{对称性分析化简}与重积分一样重要，可以降低后续的分析与计算难度。注意\textbf{不建议直接对第二型曲线或曲面积分进行对称性分析}，因为向量场对称的情况会非常复杂。

\note 第二型曲面积分大部分情况在能使用高斯公式时仍然\textbf{建议直接使用高斯公式}，因为高斯公式并不影响进行对称性分析。

\

\noindent\textbf{抽象函数}
\begin{framed}
\begin{exercise}
    设$f$是$\rb$上的连续函数，计算
    $$\iint_S\big(xf(xy)+2x−y\big)\dr y\dr z+\big(yf(xy)+2y+x\big)\dr z\dr x+\big(zf(xy)+z\big)\dr x\dr y$$
    其中$S$为锥面$z=\sqrt{x^2+y^2}$夹在平面$z=1$、$z=2$之间部分的下侧。
\end{exercise}
\sol{
    分为如下步骤：
    \begin{itemize}
        \item \textbf{参数化}
        
        \note 由于此区域并不封闭，被积函数且并不可导，无法使用高斯公式，只能采用基本的计算方式。

        虽然$z$是$x$、$y$的函数可以提供一个自然的参数化，由于曲面的旋转对称性，为了让计算更简洁可以考虑\textbf{极坐标}参数化
        $$\rv=(x,y,z)=(u\cos v,u\sin v,u)$$
        代入约束可得到参数范围
        $$1\le u\le 2$$
        将上述条件与极坐标自然范围确定的$(u,v)$区域记为$D$，可发现约束能写为
        $$1\le u\le2,\quad0\le v\le2\pi$$

        \item \textbf{定向}
        
        为了确定定向， 用下标$u$、$v$表示对$u$、$v$求偏导并保持另一个不变，直接计算有
        $$\rv_u=(x_u,y_u,z_u)=(\cos v,\sin v,1),\quad\rv_v=(x_v,y_v,z_v)=(-u\sin v,u\cos v,0)$$
        我们\textbf{任取}一个使得$\rv_u$、$\rv_v$\textbf{不平行}(也可表述为线性无关/叉乘非零)的$D$中点，如取$(u,v)=(\frac{3}{2},\pi)$，此处有$\rv=(-\frac{3}{2},0,\frac{3}{2})$、$\rv_u=(-1,0,1)$、$\rv_v=(0,-\frac{3}{2},0)$，在曲面上的点$\rv$处画出向量$\rv_u$与$\rv_v$，通过右手定则可发现$\rv_u\times\rv_v$指向上方，与要求的定向相反，因此\textbf{需要添加负号}。
        
        \note 事实上还有更简单的无需画图的判断方法：下侧代表取$z$分量小于0的法向量，直接计算可发现$\rv_u\times\rv_v$的$z$分量是$\frac{3}{2}$，大于0，因此与要求的定向相反。

        \item \textbf{化为重积分}
        
        代入
        $$\dr y\dr z=\frac{D(y,z)}{D(u,v)}\dr u\dr v,\quad\dr z\dr x=\frac{D(z,x)}{D(u,v)}\dr u\dr v,\quad\dr x\dr y=\frac{D(x,y)}{D(u,v)}\dr u\dr v$$
        与$x$、$y$、$z$的参数方程并化简可得原积分为(最左侧的负号来自定向)
        $$-\iint_D(-u^2)\dr u\dr v$$

        \item \textbf{积分计算}
        
        由区域描述可直接写出累次积分
        $$\int_1^2\dr u\int_0^{2\pi}u^2\dr v$$
        依次积分得到结果为
        $$2\pi\int_1^2u^2\dr u=\frac{14}{3}\pi$$
    \end{itemize}
}
\end{framed}

\note 对于无法应用高斯公式的情况，仍然需要掌握\textbf{直接参数化}与相应的\textbf{定向}技巧。

\begin{framed}
\begin{exercise}\label{ball cont}
    设$f(x)$是$\rb$上的连续函数，证明
    $$\iint_Sf(x+y+z)\dr S=2\pi\int_{-1}^1f(\sqrt{3}\xi)\dr\xi$$
    这里$S$为球面$x^2+y^2+z^2=1$。
\end{exercise}
\sol{
    本题最简单的做法是使用线性代数中的\textbf{正交变换}(也即进行适当的旋转)发现左侧积分等于
    $$\iint_Sf(\sqrt3z)\dr S$$
    后再化简。这里给出一个不要求线性代数基础但\textbf{更加复杂}的做法，分为如下步骤：
    \begin{itemize}
        \item \textbf{参数化}
        
        为了能够后续化简$f(x+y+z)$，我们想将球面用$u$、$v$表示，必须参数化出
        $$u=x+y+z$$
        并在第一型曲面积分化为重积分后先对$v$积分。

        为了得知此时$v$的可能形式，我们需要类似本讲义21.1.1第4题的思路进行参数化。

        我们需要用$u$、$v$描述$(x,y,z)$，且已知
        $$x+y+z=u,\quad x^2+y^2+z^2=1$$
        接下来只需将$u$看作常数，用$v$描述如上曲线即可。从第一式中解出$z=u-x-y$，代入第二式得到
        $$x^2+y^2+(x+y-u)^2=1$$
        将$u$看作常数时，这已经是一个二次曲线，对$x$配方得到
        $$2\bigg(x+\frac{y}{2}-\frac{u}{2}\bigg)^2+\frac{3}{2}y^2+\frac{1}{2}u^2-yu=1$$
        再对$y$配方得到
        $$2\bigg(x+\frac{y}{2}-\frac{u}{2}\bigg)^2+\frac{3}{2}\bigg(y-\frac{1}{3}u\bigg)^2=1-\frac{1}{3}u^2$$
        右侧需要大于等于才能存在$x$、$y$，这就得到了$u$的范围(利用柯西不等式可以直接得到相同的结果)
        $$-\sqrt3\le u\le\sqrt3$$
        而这已经与椭圆方程的形式类似，从而令
        $$\sqrt{2}\bigg(x+\frac{y}{2}-\frac{u}{2}\bigg)=\sqrt{1-\frac{1}{3}u^2}\cos v,\quad\sqrt{\frac{3}{2}}\bigg(y-\frac{1}{3}u\bigg)=\sqrt{1-\frac{1}{3}u^2}\sin v,\quad v\in[0,2\pi]$$
        反解出$x$、$y$就得到了
        $$x=\frac{u}{3}+\bigg(\frac{1}{\sqrt2}\cos v-\frac{1}{\sqrt6}\sin v\bigg)\sqrt{1-\frac{1}{3}u^2},\quad y=\frac{u}{3}+\sqrt{\frac{2}{3}}\sin v\sqrt{1-\frac{1}{3}u^2}$$
        最终代入$z$得到
        $$z=\frac{u}{3}-\bigg(\frac{1}{\sqrt2}\cos v+\frac{1}{\sqrt6}\sin v\bigg)\sqrt{1-\frac{1}{3}u^2}$$

        也即最终的参数化是
        $$\begin{cases}x=\frac{u}{3}+\big(\frac{1}{\sqrt2}\cos v-\frac{1}{\sqrt6}\sin v\big)\sqrt{1-\frac{1}{3}u^2}\\y=\frac{u}{3}+\sqrt{\frac{2}{3}}\sin v\sqrt{1-\frac{1}{3}u^2}\\z=\frac{u}{3}-\big(\frac{1}{\sqrt2}\cos v+\frac{1}{\sqrt6}\sin v\big)\sqrt{1-\frac{1}{3}u^2}\end{cases}(-\sqrt3\le u\le\sqrt3,\quad 0\le v\le2\pi)$$

        \item \textbf{化为重积分}
        
        我们设$\rv=(x,y,z)$，并设
        $$R(u)=\sqrt{1-\frac{1}{3}u^2}$$
        求导可发现
        $$R'(u)=-\frac{1}{3}\frac{u}{\sqrt{1-\frac{1}{3}u^2}}$$
        用下标$u$、$v$表示对$u$、$v$求偏导并保持另一个不变，直接计算有
        $$\rv_u=(x_u,y_u,z_u)=\bigg(\frac{1}{3},\frac{1}{3},\frac{1}{3}\bigg)+R'(u)\bigg(\frac{1}{\sqrt2}\cos v-\frac{1}{\sqrt6}\sin v,\sqrt{\frac{2}{3}}\sin v,-\frac{1}{\sqrt2}\cos v-\frac{1}{\sqrt6}\sin v\bigg)$$
        $$\rv_v=(x_v,y_v,z_v)=R(u)\bigg(-\frac{1}{\sqrt2}\sin v-\frac{1}{\sqrt6}\cos v,\sqrt{\frac{2}{3}}\cos v,\frac{1}{\sqrt2}\sin v-\frac{1}{\sqrt6}\cos v\bigg)$$
        代入公式可知
        $$\dr S=\sqrt{|\rv_u|^2|\rv_v|^2-(\rv_u\cdot\rv_v)^2}\dr u\dr v$$
        我们进行进一步的计算：设
        $$\vec{\alpha}=\bigg(\frac{1}{3},\frac{1}{3},\frac{1}{3}\bigg),\quad\vec{\beta}=\bigg(\frac{1}{\sqrt2}\cos v-\frac{1}{\sqrt6}\sin v,\sqrt{\frac{2}{3}}\sin v,-\frac{1}{\sqrt2}\cos v-\frac{1}{\sqrt6}\sin v\bigg)$$
        则$\rv_u=\vec{\alpha}+R'(u)\vec{\beta}$利用向量性质可知
        $$|\rv_u|^2=\rv_u\cdot\rv_u=|\vec{\alpha}|^2+(R'(u))^2|\vec{\beta}|^2+2R'(u)\vec{\alpha}\cdot\vec{\beta}$$
        直接计算可知$\vec{a}\cdot\vec{\beta}=0$、$|\vec{\alpha}|=\frac{\sqrt3}{3}$、$|\vec{\beta}|=1$，从而
        $$|\rv_u|=\sqrt{\frac{1}{3}+(R'(u))^2}=\frac{1}{\sqrt{3-u^2}}$$
        设
        $$\vec{\gamma}=\bigg(-\frac{1}{\sqrt2}\sin v-\frac{1}{\sqrt6}\cos v,\sqrt{\frac{2}{3}}\cos v,\frac{1}{\sqrt2}\sin v-\frac{1}{\sqrt6}\cos v\bigg)$$
        直接计算可知$|\vec{\gamma}|=1$，从而
        $$|\rv_v|=|R(u)||\vec{\gamma}|=\sqrt{1-\frac{1}{3}u^2}$$
        直接计算可知$\vec{\alpha}\cdot\vec{\gamma}=\vec{\beta}\cdot\vec{\gamma}=0$，从而
        $$\rv_u\cdot\rv_v=0$$
        综合以上可知
        $$\dr S=|\rv_u||\rv_v|\dr u\dr v=\frac{1}{\sqrt3}\dr u\dr v$$

        被积函数即为$f(u)$，由此原积分为
        $$\iint_D\frac{1}{\sqrt3}f(u)\dr u\dr v$$

        \item \textbf{积分计算}
        
        根据区域描述可以直接写出累次积分
        $$\int_{-\sqrt3}^{\sqrt3}\dr u\int_0^{2\pi}\frac{1}{\sqrt3}f(u)\dr v$$
        对$v$积分得其为
        $$2\pi\int_{-\sqrt3}^{\sqrt3}\frac{1}{\sqrt3}f(u)\dr u$$
        换元$\xi=\frac{1}{\sqrt3}u$最终得到结果为
        $$2\pi\int_{-1}^1f(\sqrt{3}\xi)\dr\xi$$
        这就是要证的结论。
    \end{itemize}
}
\end{framed}

\note 对于抽象函数问题的通用处理方法是将\textbf{抽象函数的参数看作整体换元}(包括``化为定积分''类型的题目)，不过实际处理可能相对复杂。

\subsubsection{更多证明}
本部分收录不完全为计算性的证明题。\exref{stolkes circle}中的第一/二型曲面积分关系、\exref{green normal}中的第一/二型曲线积分关系与平面曲线切向量法向量关系\textbf{常用来解决一些抽象函数相关的证明题}，\textbf{非常重要}，大家也可以在其他题中看到应用。一般有两种常用处理方法，以曲面积分的
$$\vec{n}\dr S=(\dr y\dr z,\dr z\dr x,\dr x\dr y)$$
为例介绍：
\begin{enumerate}
    \item 如果有一个第二型曲面积分
    $$\iint_S(P,Q,R)\cdot(\dr y\dr z,\dr z\dr x,\dr x\dr y)$$
    且$S$的单位法向量$\vec{n}$易于确定(如平面、柱面、球面，平面如\exref{stolkes ellipse}，柱面见\exref{column normal}，球面与圆类似，可见\exref{circle normal})，则可以将它改写为
    $$\iint_S(P,Q,R)\cdot\vec{n}\dr S$$
    \item 如果有一个第一型曲面积分
    $$\iint_SU\dr S$$
    需要化为三维区域上的积分，或可以相对自然地拆分出$U=(P,Q,R)\cdot\vec{n}$，则可以将它改写为(曲线积分的情况可见\exref{2d gauss})
    $$\iint_S(P,Q,R)\cdot(\dr y\dr z,\dr z\dr x,\dr x\dr y)$$
    若满足高斯公式条件还可以进一步改写为
    $$\iiint_\Omega(P_x+Q_y+R_z)\dr x\dr y\dr z$$
\end{enumerate}

\

\noindent\textbf{极限证明}

\begin{framed}
\begin{exercise}\label{circle normal}
    证明(这里对曲线$x^2+y^2=R^2$取逆时针定向)
    $$\lim_{R\to+\infty}\oint_{x^2+y^2=R^2}\frac{x\dr y-y\dr x}{(x^2+xy+y^2)^2}=0$$
\end{exercise}
\sol{
    记
    $$I(R)=\oint_{x^2+y^2=R^2}\frac{x\dr y-y\dr x}{(x^2+xy+y^2)^2}$$
    我们介绍两种方法：
    \begin{itemize}
        \item \textbf{化为定积分}
        
        最基本的方式即将其化为定积分。利用极坐标可得到曲线的参数化
        $$\begin{cases}x=R\cos\theta\\y=R\sin\theta\end{cases}(0\le\theta\le2\pi)$$
        且由极坐标几何意义可知定向为逆时针，无需添加符号。用下标$\theta$代表对$\theta$求导，直接计算有
        $$x_\theta=-R\sin\theta,\quad y_\theta=R\cos\theta$$
        代入$\dr x=x_\theta\dr\theta$、$\dr y=y_\theta\dr\theta$与$x$、$y$的参数方程可知
        $$I(R)=\frac{1}{R^2}\int_0^{2\pi}\frac{\dr\theta}{(1+\sin\theta\cos\theta)^2}$$
        由于
        $$1+\sin\theta\cos\theta=1+\frac{1}{2}\sin(2\theta)>0$$
        被积函数在$[0,2\pi]$连续，定积分存在，从而直接得到
        $$\lim_{R\to+\infty}I(R)=0$$

        \item \textbf{化为第一型曲线积分}

        另一种估算方法是化为第一型曲线积分。利用第一型曲线积分与第二型曲线积分的关系可知
        $$\vec\tau\dr s=(\dr x,\dr y)$$
        这里$\vec\tau=(\tau_1,\tau_2)$是曲线逆时针方向的单位切向量。对于圆$x^2+y^2=R^2$，利用几何知它在每点$(x,y)$处的一个外法向量是$(x,y)$，而逆时针方向切向量为其向左旋转$\frac{\pi}{2}$，由几何关系知为$(-y,x)$，从而单位切向量是
        $$\frac{1}{\sqrt{x^2+y^2}}(-y,x)$$
        代入可发现
        $$\begin{aligned}I(R)=&\oint_{x^2+y^2=R^2}\frac{1}{(x^2+xy+y^2)^2}(-y,x)\cdot(\dr x,\dr y)\\=&\oint_{x^2+y^2=R^2}\frac{1}{(x^2+xy+y^2)^2}(-y,x)\cdot\frac{1}{\sqrt{x^2+y^2}}(-y,x)\dr s\\=&\oint_{x^2+y^2=R^2}\frac{\sqrt{x^2+y^2}}{(x^2+xy+y^2)^2}\dr s\end{aligned}$$
        此时被积函数非负，从而可知$I(R)\ge0$。计算内积并利用曲线上$x^2+y^2=R^2$可化简得到$R>0$时积分为
        $$I(R)=R\oint_{x^2+y^2=R^2}\frac{1}{(x^2+xy+y^2)^2}\dr s$$
        接下来需要估算，利用基本不等式可知曲线上有
        $$x^2+xy+y^2\ge\frac{1}{2}(x^2+y^2)=\frac{R^2}{2}>0$$
        于是可将被积函数放大得到
        $$I(R)\le R\oint_{x^2+y^2=R^2}\bigg(\frac{R^2}{2}\bigg)^{-2}\dr s=\frac{4}{R^3}\oint_{x^2+y^2=R^2}\dr s$$
        根据第一型曲线积分的几何意义，右侧的积分即为圆$x^2+y^2=R^2$的周长$2\pi R$，最终得到
        $$I(R)\le\frac{8\pi}{R^2}$$
        之前已经估算了$I(R)\ge0$，利用\textbf{夹逼定理}即得$I(R)$在$R\to+\infty$时为0。
    \end{itemize}

    \note 本题用格林公式化为重积分无法得到结果，因为原点处无法定义，使用格林公式只能得到$I(R_1)-I(R_2)$的估算，无法说明其中一个在$R\to+\infty$时的极限。
}
\end{framed}
\note 与圆类似，球$x^2+y^2+z^2=R^2$在$(x,y,z)$处的一个外法向量即$(x,y,z)$。

\note 再次强调估算必须\textbf{化为第一型曲线/曲面积分或重积分后}再进行，第二型曲面积分无法直接放缩。

\begin{framed}
\begin{exercise}
    计算
    $$\lim_{n\to\infty}\int_{L_n}\er^{y^2-x^2}\cos(2xy)\dr x+\er^{y^2-x^2}\sin(2xy)\dr y$$
    其中$L_n$为(沿$t$增加方向)
    $$\begin{cases}x=t\\y=|\sin t|\end{cases}(0\le t\le n\pi)$$
\end{exercise}
\sol{
    分为如下步骤：
    \begin{itemize}
        \item \textbf{格林公式}
        
        将被积函数看作$P\dr x+Q\dr y$，由初等函数知$P$、$Q$在$\rb^2$上可导且导函数连续，直接计算得$Q_x$、$P_y$均为
        $$\er^{y^2-x^2}(2y\cos(2xy)-2x\sin(2xy))$$
        由此作差为0。由于$\rb^2$单连通，被积函数的积分与路径无关，$L_n$起点为$(0,0)$，终点为$(n\pi,0)$，我们可以将$L_n$上的积分转化为两点之间的线段$C_n$上的积分
        $$\int_{C_n}\er^{y^2-x^2}\cos(2xy)\dr x+\er^{y^2-x^2}\sin(2xy)\dr y$$
        
        \item \textbf{化为定积分}
        
        利用线段的参数方程可将$C_n$参数化为
        $$\begin{cases}x=t\\y=0\end{cases}(0\le t\le n\pi)$$
        用下标$t$代表对$t$求导，直接计算有
        $$x_t=1,\quad y_t=0$$
        代入$\dr x=x_t\dr t$、$\dr y=y_t\dr t$与$x$、$y$的参数方程可知积分化为
        $$\int_0^{n\pi}\er^{-t^2}\dr t$$

        \item \textbf{极限计算}
        
        由\exref{gauss}中已证
        $$\lim_{x\to+\infty}\int_0^x\er^{-t^2}\dr t=\frac{\sqrt\pi}{2}$$
        而$n\pi$在$n\to\infty$时将趋于$+\infty$，利用\textbf{归结原理}即得
        $$\lim_{n\to\infty}\int_0^{n\pi}\er^{-t^2}\dr t=\frac{\sqrt\pi}{2}$$
    \end{itemize}
}
\end{framed}

\note 本题中用到的$\er^{-x^2}$相关积分结论非常常用，\textbf{建议记忆}。

\begin{framed}
\begin{exercise}
    设$f(x,y)$为$\rb^2$上非负的连续函数，对$r>0$记
    $$I(r)=\iint_{D_r}f(x,y)\dr x\dr y,\quad J(r)=\iint_{S_r}f(x,y)\dr x\dr y$$
    其中
    $$D_r=\{(x,y)\mid x^2+y^2\le r^2\},\quad S_r=\{(x,y)\mid x\in[-r,r],\ y\in[-r,r]\}$$
    证明以下两极限
    $$\lim_{r\to+\infty}I(r),\quad\lim_{r\to+\infty}J(r)$$
    若其中一个存在，另一个必存在，且两者值相等。
\end{exercise}
\sol{
    证明分为两部分：
    \begin{itemize}
        \item $I(r)$\textbf{极限存在的情况}
        
        设$I(r)$极限为$A$，我们先证明对任何$r>0$有
        $$I(r)\le J(r)\le I(2r)$$

        \note 从之后的过程中可以看到，这是为了配凑出\textbf{夹逼定理}的形式进行估算。

        首先，无论是$I(r)$、$J(r)$还是$I(2r)$，被积函数都是$f(x,y)$，只是积分区域分别为
        $$D_r=\{(x,y)\in\rb^2\mid x^2+y^2\le r^2\}$$
        $$S_r=\{(x,y)\in\rb^2\mid-r\le x\le r,\ -r\le y\le r\}$$
        $$D_{2r}=\{(x,y)\in\rb^2\mid x^2+y^2\le 4r^2,\ x\ge0,\ y\ge0\}$$
        直接通过代数方式可验证$D_r\subset S_r\subset D_{2r}$。

        由于被积函数$f(x,y)$在$\rb^2$恒非负，它在$S_r$内$D_r$外的部分(记为$S_r\backslash D_r$)积分应当非负，在$D_{2r}$内$S_r$外的部分(记为$D_{2r}\backslash S_r$)积分也应当非负，由此利用重积分的性质即得
        $$\iint_{S_r}f(x,y)\dr x\dr y=\iint_{D_r}f(x,y)\dr x\dr y+\iint_{S_r\backslash D_r}f(x,y)\dr x\dr y\ge\iint_{D_r}f(x,y)\dr x\dr y$$
        $$\iint_{D_{2r}}f(x,y)\dr x\dr y=\iint_{S_r}f(x,y)\dr x\dr y+\iint_{D_{2r}\backslash S_r}f(x,y)\dr x\dr y\ge\iint_{S_r}f(x,y)\dr x\dr y$$
        这就得到了证明。

        接下来，利用\textbf{复合函数极限性质}可知(换元$t=2r$)
        $$\lim_{r\to+\infty}I(2r)=\lim_{t\to+\infty}I(t)=A$$
        由此结合条件利用\textbf{夹逼定理}即得
        $$\lim_{r\to+\infty}J(r)=A$$

        \item $J(r)$\textbf{极限存在的情况}
        
        设$J(r)$极限为$A$，由于已经证明了对任何$r>0$有
        $$I(r)\le J(r)\le I(2r)$$
        在右半边将$r$替换为$\frac{r}{2}$即得$r>0$时有
        $$J\bigg(\frac{r}{2}\bigg)\le I(r)\le J(r)$$
        同样利用复合函数极限性质可知(换元$t=\frac{r}{2}$)
        $$\lim_{r\to+\infty}J\bigg(\frac{r}{2}\bigg)=\lim_{t\to+\infty}J(t)=A$$
        由此结合条件利用\textbf{夹逼定理}即得
        $$\lim_{r\to+\infty}J(r)=A$$
    \end{itemize}

    \note 这里思路完全类似\exref{gauss}。
}
\end{framed}

\note 利用\textbf{重积分基本性质}的题目往往相对来说更难找到方向，但做起来更加简单，大家可以对重积分相关的定理有一些基本的概念。

\

\noindent\textbf{Green公式相关}
\begin{framed}
\begin{exercise}
    设函数$P(x,y)$、$Q(x,y)$在$\rb^2$有连续的偏导数，且对任何$(x_0,y_0)\in\rb^2$与$R>0$有
    $$\int_LP(x,y)\dr x+Q(x,y)\dr y=0$$
    这里$L$为函数曲线$y=y_0+\sqrt{R^2-(x-x_0)^2}$，$x$从$x_0-R$到$x_0+R$。证明对任何$(x,y)\in\rb^2$有
    $$P(x,y)=0,\quad\frac{\partial Q(x,y)}{\partial x}=0$$
\end{exercise}
\sol{
    我们将$Q$对第一个分量的偏导记作$Q_x$、$P$对第二个分量的偏导记作$P_y$。分为如下步骤：
    \begin{itemize}
        \item \textbf{分析}
        
        由于这种本题的条件形式有点类似保守场条件(沿任何闭曲线积分为0)，且从要证结论中确实可以看出$Q_x-P_y=0$，我们应当先尝试能否先证明上述条件。另一方面，证明此式以后就可以说明积分与路径无关，进一步化简条件。由此，我们可以将接下来的证明分为三步。
        
        \item \textbf{配凑格林公式}

        \note 思路：为了证明在任何$(x,y)$处$Q_x-P_y=0$，这里采用与教材8.3节类似的证明方式：若某个$(x_0,y_0)$处$Q_x-P_y\ne0$，我们考虑在它附近画一个小的封闭曲线，利用连续性可知在其中的积分就将非零。然而，条件为所有上半圆上的积分均为0，只要能用三个上半圆拼成一条封闭曲线，就可以得到矛盾。

        若$(x_0,y_0)$处$Q_x-P_y\ne0$，给定$r>0$，考虑如下三个半圆(定向均为$x$从小到大)：
        $$L_1:\quad y=y_0+\sqrt{r^2-(x-x_0)^2},\quad x\in[x_0-r,x_0+r]$$
        $$L_2:\quad y=y_0+\sqrt{\frac{r^2}{4}-\bigg(x-x_0-\frac{r}{2}\bigg)^2},\quad x\in[x_0,x_0+r]$$
        $$L_3:\quad y=y_0+\sqrt{\frac{r^2}{4}-\bigg(x-x_0+\frac{r}{2}\bigg)^2},\quad x\in[x_0-r,x_0]$$
        从几何上可以发现，$L_1$、$L_2$、$L_3$可以围成一个封闭区域$D$，且这个区域包含在$(x_0,y_0)$附近以$r$为半径的圆以内。

        以下省略函数的参数$(x,y)$。利用条件可知
        $$\int_{L_1}(P\dr x+Q\dr y)=\int_{L_2}(P\dr x+Q\dr y)=\int_{L_3}(P\dr x+Q\dr y)=0$$
        考虑$L_1$反向，$L_2$、$L_3$正向围成的封闭曲线$C$与它包含的区域$D$。利用$C$的定义可发现
        $$\oint_C(P\dr x+Q\dr y)=\int_{L_2}(P\dr x+Q\dr y)+\int_{L_3}(P\dr x+Q\dr y)-\int_{L_1}(P\dr x+Q\dr y)=0$$
        从几何可以看出$C$是逆时针方向，由此利用格林公式可知
        $$\oint_C(P\dr x+Q\dr y)=\iint_D(Q_x-P_y)\dr x\dr y$$
        也即
        $$\iint_D(Q_x-P_y)\dr x\dr y=0$$

        \item \textbf{连续性推矛盾}

        由于$Q_x(x_0,y_0)-P_y(x_0,y_0)\ne0$，不妨设其大于0，值为$c$。由条件$Q_x-P_y$连续，从而存在$r$使得
        $$\forall (x,y),\quad(x-x_0)^2+(y-y_0)^2<r^2\Longrightarrow|(Q_x(x,y)-P_x(x,y))-(Q_x(x_0,y_0)-P_x(x_0,y_0))|<\frac{c}{2}$$
        也就能得到
        $$\forall (x,y),\quad(x-x_0)^2+(y-y_0)^2<r^2\Longrightarrow Q_x(x,y)-P_x(x,y)\ge\frac{c}{2}$$
        由此可得
        $$\iint_D(Q_x-P_y)\dr x\dr y\ge\iint_D\frac{c}{2}\dr x\dr y$$
        利用重积分的几何意义，$\iint_D\dr x\dr y$即为区域面积，区域为一个大半圆去掉两个小半圆，因此面积是
        $$\frac{\pi r^2}{2}-\frac{\pi(r/2)^2}{2}-\frac{\pi(r/2)^2}{2}=\frac{1}{4}\pi r^2$$
        从而
        $$\iint_D(Q_x-P_y)\dr x\dr y\ge\frac{\pi cr^2}{8}>0$$
        与等于0矛盾。
        
        \item \textbf{进一步计算}
        
        综合以上，我们已经证明了$\rb^2$上任何点处$Q_x-P_y=0$，由此可知\textbf{积分与路径无关}。

        对任何两点$(x_0,y_0)$与$(x_1,y_0)$，其中$x_0<x_1$，可以构造一个从$(x_0,y_0)$走向$(x_1,y_0)$的上半圆，由此
        $$\int_{(x_0,y_0)}^{(x_1,y_0)}(P\dr x+Q\dr y)=0$$
        考虑$(x_0,y_0)$到$(x_1,y_0)$的线段，利用直线方程可参数化为
        $$\begin{cases}x=t\\y=y_0\end{cases}(x_0\le t\le x_1)$$
        用下标$t$表示对$t$求导，直接计算有$x_t=1$、$y_t=0$，代入$\dr x=x_t\dr t$、$\dr y=y_t\dr t$与$x$、$y$的参数方程可得
        $$\int_{x_0}^{x_1}P(t,y_0)\dr t=0$$
        若$(x_2,y_2)$处$P(x_2,y_2)\ne0$，取充分小的$\varepsilon>0$，并取$x_0=x_2-\varepsilon$、$x_1=x_2+\varepsilon$、$y_0=y_2$，与前一部分证明类似可以得到积分不为0，矛盾，由此最终得到
        $$\forall(x,y)\in\rb^2,\quad P(x,y)=0$$
        由此计算得到$P_y$恒为0，从而通过$Q_x-P_y$恒为0，于是
        $$\forall(x,y)\in\rb^2,\quad Q_x(x,y)=0$$
        这就证明了结论。
    \end{itemize}
}
\end{framed}

\note 这类条件相对奇怪的问题第一步是先\textbf{寻找熟悉的形式}，再设法向熟悉的形式配凑。对于平面相关的问题，\textbf{几何角度的分析}或许会很有帮助。

\begin{framed}
\begin{exercise}\label{2d gauss}
    设函数$P(x,y)$、$Q(x,y)$在有界闭区域$D$上有连续的偏导数，$D$有分段光滑边界$L$，证明(以下省略函数的参数$(x,y)$)：
    $$\oint_L(P,Q)\cdot\nv\dr s=\iint_D\bigg(\frac{\partial P}{\partial x}+\frac{\partial Q}{\partial y}\bigg)\dr x\dr y$$
    其中$\nv$为$L$对应位置的单位外法向量。
\end{exercise}
    \sol{
        分为如下步骤：
        \begin{itemize}
            \item \textbf{化为第二型曲线积分}
            
            由于需要让曲线积分与导数在区域的积分产生关联，需要应用格林公式，从而需要先将曲线积分化为第二型曲线积分。

            设$\vec\tau=(\tau_1,\tau_2)$代表曲线逆时针方向的单位切向量，利用第一型曲线积分与第二型曲线积分的关系可知(这里曲线以逆时针为正定向)
            $$\vec\tau\dr s=(\dr x,\dr y)$$
            不过，题中为\textbf{法向量}的相关式子，而非切向量。利用几何关系可知，$L$的逆时针方向单位切向量为单位外法向量向逆时针方向旋转90度的结果。设单位外法向量$\nv=(n_1,n_2)$，计算可发现
            $$n_1=\tau_2,\quad n_2=-\tau_1$$
            从而有
            $$n_1\dr s=\dr y,\quad n_2\dr s=-\dr x$$
            直接代入得到左侧的积分即
            $$\oint_{L^+}(P\dr y-Q\dr x)$$
            
            \item \textbf{格林公式}
            
            由条件$P$、$Q$在$D$上存在连续偏导数，且$D$以$L$为边界，直接利用格林公式可知
            $$\oint_{L^+}(P\dr y-Q\dr x)=\iint_D(P_x+Q_y)\dr x\dr y$$
            这就是所求的结果。
        \end{itemize}
    }
\end{framed}

\note 注意此形式事实上是\textbf{二维的高斯公式}，有时可以用来处理利用法向量而非法向量描述的题目。这类问题虽然是证明题的形式，但有着\textbf{鲜明的计算思路}，本质来说并不困难。

\begin{framed}
\begin{exercise}\label{unique}
    设开区域$U$为$\rb^2$去掉$x$轴正半轴(含原点)，$D$为$\rb^2$去掉原点，找到一个$T(x,y)$使其在$U$上满足
    $$\frac{\partial T(x,y)}{\partial x}=-\frac{y}{x^2+y^2},\quad\frac{\partial T(x,y)}{\partial y}=\frac{x}{x^2+y^2}$$
    证明不存在$T(x,y)$在$D$上满足上式。
\end{exercise}
\sol{
    分为如下步骤：
    \begin{itemize}
        \item \textbf{$U$上的$T(x,y)$计算}
        
        设
        $$P(x,y)=-\frac{y}{x^2+y^2},\quad Q(x,y)=\frac{x}{x^2+y^2}$$
        直接计算会发现
        $$\frac{\partial Q(x,y)}{\partial x}-\frac{\partial P(x,y)}{\partial y}=0$$
        由定义可发现$U$是单连通的，因此利用教材8.3节的定理即得存在符合要求的原函数$T$，下面进行计算。

        \note 这里$U$的单连通性质并不显然，但从几何角度可以发现成立。

        \note 若直接找$P(x,y)$对$x$的原函数可发现为$-\arctan\frac{x}{y}+C$，但其在$x$轴负半轴没有定义，需要由此时$x<0$用极限定义，说法相对复杂，我们这里采取\textbf{找路径积分}的方式。

        取区域$U$中一个点$(0,1)$，定义原函数
        $$T(p,q)=\int_{(0,1)}^{(p,q)}(P(x,y)\dr x+Q(x,y)\dr y)$$
        下面进行计算。

        由于$P$与$Q$存在分母$\frac{1}{x^2+y^2}$，而挖掉的$x$轴正半轴恰好为极轴，为了让积分更加简单，我们可以考虑\textbf{极坐标}下保持$\theta$不变$r$变化与保持$r$不变$\theta$变化的路径。

        具体来说，由于$(0,1)$的极坐标是$r=1$、$\theta=\frac{\pi}{2}$，设$(p,q)$的极坐标是$(r_0,\theta_0)$，也即
        $$p=r_0\cos\theta_0,\quad q=r_0\sin\theta_0$$
        且$\theta_0\in(0,2\pi)$，我们考虑由两条曲线拼成的积分路径，第一条$L_1$为
        $$\begin{cases}x=0\\y=t\end{cases}$$
        这里参数$t$从1走到$r_0$\ (无论谁大谁小)，第二条$L_2$为
        $$\begin{cases}x=r_0\cos t\\y=r_0\sin t\end{cases}$$
        这里参数$t$从$\frac{\pi}{2}$走到$\theta_0$\ (无论谁大谁小)。可以发现，由于$r$与$r_0$均为正，$\theta$与$\theta_0$都在$(0,2\pi)$中，积分路径一定恒在$U$里，不会经过挖掉的$\theta=0$。

        用下标$t$表示对$t$求导，直接计算有$L_1$上
        $$x_t=0,\quad y_t=1$$
        $L_2$上
        $$x_t=-r_0\sin t,\quad y_t=r_0\cos t$$
        代入$\dr x=x_t\dr t$、$\dr y=y_t\dr t$与$x$、$y$的参数方程可计算得
        $$T(p,q)=\int_1^{r_0}\frac{0}{t^2}\dr t+\int_{\pi/2}^\theta\frac{r_0^2}{r_0^2}\dr t=\theta-\frac{\pi}{2}$$
        由于在极坐标的上述定义下每个$(p,q)$对应唯一一个$\theta$，这已经是符合要求的$T$。

        \item \textbf{证明$D$上不存在}
        
        我们将上述$U$上符合要求的$T$记为$T_0$，假设$D$上存在$T$满足要求，我们下面来推出矛盾。

        根据定义可以发现$T$与$T_0$对$x$、$y$的偏导在$U$上均相同，从而
        $$\forall(x,y)\in U,\quad\frac{\partial T(x,y)-T_0(x,y)}{\partial x}=0,\quad \frac{\partial T(x,y)-T_0(x,y)}{\partial y}=0$$
        利用拉格朗日中值定理可证明$T-T_0$在$U$中为常数(详细证明可见上学期讲义15.1第10题，教材上也有这个结论)，从而可知存在$C_0\in\rb$使得
        $$\forall(x,y)\in U,\quad T(x,y)=\theta-\frac{\pi}{2}+C_0$$
        记$C_1=C_0-\frac{\pi}{2}$也即存在$C_1\in\rb$使得
        $$\forall(x,y)\in U,\quad T(x,y)=\theta+C_1$$
        这里$\theta$代表极坐标下$(x,y)$的极角，范围在$(0,2\pi)$。

        \note 从几何角度可以想象这样的$T(x,y)$不可能在$x$轴正半轴连续，因为从上方逼近是$C_1$，从下方逼近是$2\pi+C_1$，下面严谨证明。

        考虑$D$内$U$外的一个点$(1,0)$，并考虑过它的直线$x=1$，直接计算极角可以发现，直线上有
        $$T(1,y)=\begin{cases}\arctan y+C_1&y>0\\\arctan y+2\pi+C_1&y<0\end{cases}$$
        由此可发现
        $$\lim_{y\to0^+}T(1,y)=C_1,\quad\lim_{y\to0^-}T(1,y)=C_1+2\pi$$
        由于两者不等，$T$在$(1,0)$处\textbf{$y$方向不连续}，于是\textbf{对$y$偏导数不存在}，矛盾。
    \end{itemize}
}
\end{framed}

\note 这里用到了很特殊的性质，也即单连通区域上势函数相差常数的\textbf{唯一性}。

\begin{framed}
\begin{exercise}
    考虑在$\rb^2$上二阶连续可微的函数$f$，且满足
    $$\triangle f(x,y)=\er^{-(x^2+y^2)}$$
    计算
    $$\iint_D(x\partial_1f(x,y)+y\partial_2f(x,y))\dr x\dr y$$
    其中
    $$D=\{(x,y)\mid x^2+y^2\le1\}$$
\end{exercise}
\sol{
    设所求的积分为$I$，我们介绍两种做法，直接计算重积分与从二维高斯公式出发的做法：
    \begin{itemize}
        \item \textbf{直接计算-分部积分化简}
        
        为方便书写，以下有时会省略函数的参数$(x,y)$\ (或极坐标下写为$(r\cos\theta,r\sin\theta)$)根据区域的对称性，利用极坐标换元
        $$x=r\cos\theta,\quad y=r\sin\theta$$
        此时原本的被积函数为$r\cos\theta\partial_1f+r\sin\theta\partial_2f$，且计算Jacobi行列式得$\dr x\dr y=r\dr r\dr\theta$。此时区域要求仅$r^2\le t^2$，结合自然范围可知约束即化为
        $$0\le r\le1,\quad0\le\theta\le2\pi$$
        由此可以写出累次积分
        $$I=\int_0^1\dr r\int_0^{2\pi}r^2(\cos\theta\partial_1f(r\cos\theta,r\sin\theta)+\sin\theta\partial_2f(r\cos\theta,r\sin\theta))\dr\theta$$
        这里由于参数互相不影响，对$r$或$\theta$先积分均可。为了向已知的二阶导配凑，需要使用\textbf{分部积分}。尝试先对$\theta$积分并分部，计算有(两者的第三个等号都利用了多元复合函数求导)
        $$\begin{aligned}&\int_0^{2\pi}\cos\theta\partial_1f(r\cos\theta,r\sin\theta)\dr\theta\\=&\int_0^{2\pi}\partial_1f(r\cos\theta,r\sin\theta)\dr\sin\theta\\=&0-0-\int_0^{2\pi}\sin\theta\dr\partial_1f(r\cos\theta,r\sin\theta)\\=&-\int_0^{2\pi}\sin\theta(-r\sin\theta\partial_1^2f+r\cos\theta\partial_2\partial_1f)\dr\theta\end{aligned}$$
        $$\begin{aligned}&\int_0^{2\pi}\sin\theta\partial_2f(r\cos\theta,r\sin\theta)\dr\theta\\=&-\int_0^{2\pi}\partial_2f(r\cos\theta,r\sin\theta)\dr\cos\theta\\=&-\partial_2f(r,0)+\partial_2f(r,0)+\int_0^{2\pi}\cos\theta\dr\partial_2f(r\cos\theta,r\sin\theta)\\=&\int_0^{2\pi}\cos\theta(-r\sin\theta\partial_1\partial_2f+r\cos\theta\partial_2^2f)\dr\theta\end{aligned}$$
        由于$f$二阶连续可微，$\partial_1\partial_2f=\partial_2\partial_1f$，从而对$\theta$分部后积分化为了
        $$I=\int_0^1\dr r\int_0^{2\pi}r^3(\sin^2\theta\partial_1^2f-2\sin\theta\cos\theta\partial_1\partial_2f+\cos^2\theta\partial_2^2f)\dr\theta$$
        这依然无法解决问题，重新尝试先对$r$积分并分部，计算有(这里第三个等号仍然利用了多元复合函数求导，注意此时在对$r$求偏导)
        $$\begin{aligned}&\int_0^1r^2\partial_1f(r\cos\theta,r\sin\theta)\dr r=\frac{1}{3}\int_0^1\partial_1f(r\cos\theta,r\sin\theta)\dr r^3\\=&\frac{1}{3}\partial_1f(\cos\theta,\sin\theta)-\frac{1}{3}\int_0^1r^3\dr\partial_1f(r\cos\theta,r\sin\theta)\\=&\frac{1}{3}\partial_1f(\cos\theta,\sin\theta)-\frac{1}{3}\int_0^1r^3(\cos\theta\partial_1^2f+\sin\theta\partial_1\partial_2f)\dr r\end{aligned}$$
        $$\begin{aligned}&\int_0^1r^2\partial_2f(r\cos\theta,r\sin\theta)\dr r=\frac{1}{3}\int_0^1\partial_2f(r\cos\theta,r\sin\theta)\dr r^3\\=&\frac{1}{3}\partial_2f(\cos\theta,\sin\theta)-\frac{1}{3}\int_0^1r^3\dr\partial_2f(r\cos\theta,r\sin\theta)\\=&\frac{1}{3}\partial_2f(\cos\theta,\sin\theta)-\frac{1}{3}\int_0^1r^3(\cos\theta\partial_2\partial_1f+\sin\theta\partial_2^2f)\dr r\end{aligned}$$
        由于$f$二阶连续可微，$\partial_1\partial_2f=\partial_2\partial_1f$，从而对$r$分部后积分化为了
        $$\begin{aligned}I=&\frac{1}{3}\int_0^{2\pi}(\cos\theta\partial_1f(\cos\theta,\sin\theta)+\sin\theta\partial_2f(\cos\theta,\sin\theta))\\ &-\frac{1}{3}\int_0^{2\pi}\dr\theta\int_0^1r^3(\cos^2\theta\partial_1^2f+2\cos\theta\sin\theta\partial_1\partial_2f+\sin^2\theta\partial_2^2f)\dr r\end{aligned}$$

        \item \textbf{直接计算-对比化简}

        由于$r$与$\theta$互相不影响，可以将对$r$的分部积分结果改写为
        $$\begin{aligned}I=&\frac{1}{3}\int_0^{2\pi}(\cos\theta\partial_1f(\cos\theta,\sin\theta)+\sin\theta\partial_2f(\cos\theta,\sin\theta))\dr\theta\\ &-\frac{1}{3}\int_0^1\dr r\int_0^{2\pi}r^3(\cos^2\theta\partial_1^2f+2\cos\theta\sin\theta\partial_1\partial_2f+\sin^2\theta\partial_2^2f)\dr\theta\end{aligned}$$
        对比之前分部积分得到的
        $$I=\int_0^1\dr r\int_0^{2\pi}r^3(\sin^2\theta\partial_1^2f-2\sin\theta\cos\theta\partial_1\partial_2f+\cos^2\theta\partial_2^2f)\dr\theta$$
        第一式的三倍减去第二式即得
        $$2I=\int_0^{2\pi}(\cos\theta\partial_1f(\cos\theta,\sin\theta)+\sin\theta\partial_2f(\cos\theta,\sin\theta))\dr\theta-\int_0^1\dr r\int_0^{2\pi}r^3(\partial_1^2f+\partial_2^2f)\dr\theta$$
        这时终于可以应用条件得到
        $$2I=\int_0^{2\pi}(\cos\theta\partial_1f(\cos\theta,\sin\theta)+\sin\theta\partial_2f(\cos\theta,\sin\theta))\dr\theta-\int_0^1\dr r\int_0^{2\pi}r^3\er^{-r^2}\dr\theta$$
        我们拆分为两部分
        $$I_1=\int_0^{2\pi}(\cos\theta\partial_1f(\cos\theta,\sin\theta)+\sin\theta\partial_2f(\cos\theta,\sin\theta))\dr\theta$$
        $$I_2=\int_0^1\dr r\int_0^{2\pi}r^3\er^{-r^2}\dr\theta$$

        \item \textbf{直接计算-第一部分计算}
        
        虽然第一部分仍然难以直接计算，但此形式与几何意义(在$r$的边界，即单位圆上)让我们可以想到化为第二型曲线积分后重新利用格林公式。我们考虑单位圆$C$，用极坐标逆时针参数化为
        $$\begin{cases}x=\cos\theta\\y=\sin\theta\end{cases}(0\le\theta\le2\pi)$$
        则$\cos\theta$即为$y_\theta$、$\sin\theta$即为$-x_\theta$，由此可将其\textbf{改写回第二型曲线积分}
        $$\oint_{C^+}\partial_1f(x,y)\dr y-\partial_2f(x,y)\dr x$$
        这样就可以应用\textbf{格林公式}得到积分为
        $$\iint_D(\partial_1^2f(x,y)+\partial_2^2f(x,y))\dr x\dr y$$
        代入条件知为
        $$\iint_D\er^{-x^2-y^2}\dr x\dr y$$
        进行极坐标换元，完全类似\exref{gauss}中可计算得
        $$I_1=\pi(1-\er^{-1})$$

        \item \textbf{直接计算-第二部分计算与综合结果}
        
        直接对$\theta$积分有
        $$I_2=2\pi\int_0^1r^3\er^{-r^2}\dr r$$
        换元$s=r^2$得
        $$I_2=\pi\int_0^1s\er^{-s}\dr s$$
        分部积分有
        $$I_2=-\pi\int_0^1s\dr\er^{-s}=-\pi\er^{-1}+\pi\int_0^1\er^{-s}\dr s$$
        最终算出
        $$I_2=\pi-2\pi\er^{-1}$$
        综合以上有
        $$I=\frac{I_1-I_2}{2}=\frac{\pi}{2\er}$$

        \item \textbf{二维高斯公式}
        
        对于无法处理的区域积分，我们可以考虑使用\textbf{格林公式}将其重新化回曲线积分再进行处理。本题中由于$x$、$y$的对称性，利用\exref{2d gauss}的形式将更加有效，由此可以想到考虑函数$P(x,y)$与$Q(x,y)$使得$x\partial_1f(x,y)$接近$\partial_1P(x,y)$、$y\partial_2f(x,y)$接近$\partial_2Q(x,y)$。

        由于$\partial_1^2f+\partial_2^2f$是不难处理的，考虑构造
        $$P(x,y)=\frac{x^2+y^2}{2}\partial_1f(x,y),\quad Q(x,y)=\frac{x^2+y^2}{2}\partial_2f(x,y)$$
        根据单位圆的性质可知它$(x,y)$处的单位外法向量即为$(x,y)$，设$D$的边界为$C$，代入\exref{2d gauss}的结论计算可知
        $$\int_C(xP(x,y)+yQ(x,y))\dr s=\iint_D\bigg((x\partial_1f+y\partial_2f)+\frac{x^2+y^2}{2}(\partial_1^2f+\partial_2^2f)\bigg)\dr x\dr y$$
        由此代入$\partial_1^2f(x,y)+\partial_2^2f(x,y)$的条件并移项可发现
        $$I=\int_C\frac{x^2+y^2}{2}(x\partial_1f+y\partial_2f)\dr s-\iint_D\frac{x^2+y^2}{2}\er^{-x^2-y^2}\dr x\dr y$$
        由于$x^2+y^2$在$C$上恒为1，将其代换成1有
        $$I=\frac{1}{2}\int_C(x\partial_1f+y\partial_2f)\dr s-\frac{1}{2}\iint_D(x^2+y^2)\er^{-x^2-y^2}\dr x\dr y$$
        由于$(x,y)$是单位圆上的单位外法向量，对$C$上的积分再次使用\exref{2d gauss}的结论可知
        $$I=\frac{1}{2}\iint_D(\partial_1^2f+\partial_2^2f)\dr x\dr y-\frac{1}{2}\iint_D(x^2+y^2)\er^{-x^2-y^2}\dr x\dr y$$
        这里$I$的形式已经是$\frac{1}{2}I_1-\frac{1}{2}I_2$，极坐标换元计算可发现$I_1$、$I_2$与前一种做法中的完全相同，从而可算出相同的结果。
    \end{itemize}

    \note 本题有另一种相对简单(但较难想清楚)的做法：极坐标写出重积分后先对$\theta$再对$r$积分，将每个$r$处对$\theta$的积分类似\textbf{直接计算-第一部分计算}中方法改写为第二型曲线积分后利用格林公式计算。
}
\end{framed}

\note 很多题目即使只有基础的计算方法也能得到结果，只是会\textbf{相对复杂}。这里用分部积分将\textbf{函数与导数联系}是常用的，本质上在重积分情况与格林公式等价但更加自由。

\subsubsection{常微分方程}
常微分方程相关内容的基本要求是记忆如下的\textbf{所有12种}方程的\textbf{类型判断方式}与对应的解法：
\begin{itemize}
    \item 变量分离的方程、齐次方程、$\frac{\dr y}{\dr x}=f\big(\frac{a_1y+b_1x+c_1}{a_2y+b_2x+c_2}\big)$\ (最后这种情况不要硬记每一小类的解法，记忆\textbf{配凑思路})
    \item 一阶线性微分方程(最推荐记忆\textbf{积分因子解法})、伯努利方程
    \item 全微分方程、可观察积分因子的方程(这类方程的技巧相对复杂，可见\exref{find factor})、存在只和$x$或$y$有关的积分因子的方程
    \item 不含$y$的二阶微分方程、不含$x$的二阶微分方程
    \item 二阶常系数线性微分方程、欧拉方程
\end{itemize}
更多的证明与解法\textbf{几乎不作要求}，不过，很多时候更好的配凑方式可以大幅简化问题。

\

\noindent\textbf{一般问题}
\begin{framed}
\begin{exercise}
    求以下微分变量为$x$的微分方程通解：
    $$y'-y=xy^5$$
\end{exercise}
\sol{
    \note 此方程符合\textbf{伯努利方程}的形式，两边同除以$y^5$即可化为一阶线性微分方程。我们这里介绍更一般的出现$y$的高次幂时的分析方式。

    分为如下步骤：
    \begin{itemize}
        \item \textbf{分析}
        
        由于出现了$y$的高次方，我们希望能进行\textbf{降次}以方便操作。更具体来说，对实数$k$，两边同乘$y^k$得到
        $$y^ky'-y^{k+1}=xy^{k+5}$$
        在$k\ne-1$时可以换元得到(利用复合函数求导公式)
        $$\frac{1}{k+1}(y^{k+1})'-y^{k+1}=xy^{k+5}$$
        考虑$u=y^{k+1}$，为了其他部分能写为$u$的函数，只需$k+5$是$k+1$的\textbf{倍数}，而最简单的取法即令$k+5=0$，此时$k=-5$。

        \item \textbf{通解计算}
        
        在$y\ne0$时，两边同乘$\frac{1}{y^5}$并换元$u=\frac{1}{y^4}$，利用如上分析即得
        $$-\frac{1}{4}u'-u=x$$
        这已经是一阶线性微分方程，可以直接求解，我们下面进行适当的配凑以方便计算。由于其可以写为
        $$u'=-4u-4x$$
        构造$v=u+x$可发现问题简化为
        $$v'=u'+1=-4(u+x)+1=-4v+1$$
        此时再令$w=v-\frac{1}{4}$即得到
        $$w'=4w$$
        从而有($C$为任意常数)
        $$w=C\er^{-4x}$$
        回代得到
        $$\frac{1}{y^4}=C\er^{-4x}-x+\frac{1}{4}$$
        这就得到了通解的形式
        $$y^4\bigg(C\er^{-4x}-x+\frac{1}{4}\bigg)=1$$

        \item \textbf{特解获取}
        
        由于过程中除以了$y^5$，额外讨论$y=0$的情况可发现$y=0$为一个特解。

        \item (\extra)\textbf{严谨性分析}
        
        我们最后来论证上述通解与特解即为全部解。

        考虑某个解曲线$y=f(x)$。任给曲线上一点$(x_0,y_0)$，当$y_0\ne0$时根据通解部分的推导知这点附近一定存在$C_0$使得
        $$y^4\bigg(C_0\er^{-4x}-x+\frac{1}{4}\bigg)=1$$
        另一方面计算可发现不同的$C$对应的通解不交，且恒满足$y\ne0$，由此$y=f(x)$一定是上述$C_0$对应的曲线的一部分，被通解刻画。

        另一方面，上述推理已经证明了一点处$y_0\ne0$则$y$恒不为0，由此若一点处$y_0=0$，解曲线只能为$y=0$。

        \note 注意通解必须在$C\er^{-4x}-x+\frac{1}{4}>0$时才有定义，意味着给定的$x_0$处$C$并不能取全部实数。不过，对于初值问题，可以发现给定初值$(x_0,y_0)$后一定可以唯一求解出对应的$C_0$。

        \note 更严谨的论证需要用到\textbf{拓扑}，已经超出了大家需要了解的范围。
    \end{itemize}
}
\end{framed}

\note 注意对微分方程两边同乘$\frac{1}{f(x,y)}$时，本质上我们研究的是\textbf{初值$(x_0,y_0)$不在$f(x,y)=0$上的初值问题的局部解}。大部分情况下只需讨论$f(x,y)=0$是否为\textbf{特解}即可，无需更进一步的严谨性分析(如解可不可能在部分点满足$f(x,y)=0$)。

\begin{framed}
\begin{exercise}
    求以下微分变量为$x$的微分方程通解：
    $$xy'-y=(x+y)\ln\frac{x+y}{x}$$
\end{exercise}
\sol{
    \note 此方程符合\textbf{齐次方程}的形式，直接代换$y=ux$可化简得到$u'=(u+1)\ln(u+1)$，从而可直接用分离变量的方程结论积分得到结果(注意$\frac{u'}{u+1}=(\ln(u+1))'$)。我们这里介绍配凑的方法。
    
    我们介绍两种方法：
    \begin{itemize}
        \item \textbf{齐次方程}
        
        计算可验证此方程符合齐次方程的形式，从而令$y=ux$，代入可发现
        $$x^2u'=x(1+u)\ln(1+u)$$
        由于定义域已经保证了$\frac{x+y}{x}>0$，可知$x\ne0$且$1+u>0$，上述方程进一步化简为
        $$u'=(1+u)\ln(1+u)\cdot\frac{1}{x}$$
        当$u\ne0$时$\ln(1+u)\ne0$，它可以进一步写为
        $$\frac{u'}{(1+u)\ln(1+u)}=\frac{1}{x}$$
        这已经是分离变量方程的形式，直接积分算得
        $$\int\frac{\dr u}{(1+u)\ln(1+u)}=\int\frac{\dr\ln(1+u)}{\ln(1+u)}=\ln|\ln(1+u)|+C$$
        从而通解即为($C$为任意常数)
        $$\ln|\ln(1+u)|=\ln|x|+C$$
        两边同作$\er$的指数得
        $$|\ln(1+u)|=\er^C|x|$$
        考虑正负可去除两侧绝对值并将恒正的$\er^C$改写为新的任意常数$C$，且要求$C$\textbf{非零}。不过，进一步观察可发现$C=0$时恰好对应$\ln(1+u)=0$，可验证其为特解(此时$u$恒为0)，综合可将通解改写为
        $$\ln(1+u)=Cx$$
        也即
        $$\ln\frac{x+y}{x}=Cx$$
        可最终解出(这里$\er^C$不能改写为任意常数$C$，因为须保证为正)
        $$y=x(\er^{Cx}-1)$$

        \item \textbf{直接计算}
        
        由于此方程右侧$x+y$可作为整体，尝试换元$u=x+y$，代入$y=u-x$可发现左侧即
        $$xu'-x-u+x=xu'-u$$
        由此方程化为
        $$xu'-u=u\ln\frac{u}{x}$$
        当$u>0$时，根据条件也有$x>0$，此方程可化为
        $$x\frac{u'}{u}-1=\ln\frac{u}{x}=\ln u-\ln x$$
        此时$\frac{u'}{u}=(\ln u)'$，由此记$v=\ln u$有
        $$xv'-1=v-\ln x$$
        已经可以利用一阶线性微分方程的通解得到结果。不过，我们还有更好的化简方法。记$t=\ln x$，可发现
        $$\frac{\dr v}{\dr t}=\frac{\dr v}{\dr x}\frac{\dr x}{\dr t}=xv'$$
        从而有
        $$\frac{\dr v}{\dr t}-1=v-t$$
        记$w=v-t$即得
        $$\frac{\dr w}{\dr t}=w$$
        从而($C$为任意常数)
        $$w=C\er^t=Cx$$
        代回原式可发现
        $$\ln(x+y)=Cx+\ln x$$
        求解可得
        $$y=x(\er^{Cx}-1)$$
        当$x<0$时，记$v=\ln(-u)$可得到相同的结论。

        \item (\extra)\textbf{严谨性分析}
        
        本题的定义域保证了$x+y$与$x$同号且均非零。在此定义域内，上述通解对应的
        $$C=\frac{1}{x}\ln\frac{x+y}{x}$$
        可由初值\textbf{唯一确定}。由此可发现不同的$C$对应的曲线不交，从而根据过任一点$(x_0,y_0)$的解必须恒在$C_0=\frac{1}{x_0}\ln\frac{x_0+y_0}{x_0}$对应的曲线上，这就证明了通解是全部解。
    \end{itemize}
}
\end{framed}

\note 即使遗忘了某个类型的微分方程和解法，往往也可以通过\textbf{配凑微分}的形式一步步得到结果，且分析方法可能并不复杂。

\begin{framed}
\begin{exercise}
    求以下微分变量为$x$的微分方程满足$y(0)=1$的解：
    $$(xy-x^3y^3)\dr x+(1+x^2)\dr y=0$$
\end{exercise}
\sol{
    分为如下步骤：
    \begin{itemize}
        \item \textbf{初步化简}
        
        此方程符合\textbf{伯努利方程}的形式，当$y\ne0$时，两边同除以$y^3\dr x$得到
        $$\frac{x}{y^2}-x^3+(1+x^2)\frac{y'}{y^3}=0$$
        利用复合函数求导公式，设$u=\frac{1}{y^2}$即得到
        $$xu-x^3=\frac{1}{2}(1+x^2)u'$$
        由$1+x^2>0$可进一步写为
        $$u'-\frac{2x}{1+x^2}u=-\frac{2x^3}{1+x^2}$$

        \item \textbf{一阶线性微分方程}
        
        直接计算
        $$\int-\frac{2x}{1+x^2}\dr x=\int-\frac{\dr x^2}{1+x^2}=-\ln(1+x^2)+C$$
        由此可取一个积分因子
        $$\mu(x)=\er^{-\ln(1+x^2)}=\frac{1}{1+x^2}$$
        方程两边同乘$\mu(x)$即得
        $$\bigg(\frac{u}{1+x^2}\bigg)'=-\frac{2x^3}{(1+x^2)^2}$$
        对右侧不定积分有(换元$t=x^2$、$s=1+t$，$C$为任意常数)
        $$\int-\frac{2x^3}{(1+x^2)^2}\dr x=\int-\frac{t}{(1+t)^2}\dr t=\int\frac{1-s}{s^2}\dr s=-\frac{1}{s}-\ln|s|+C$$
        代入最终得到
        $$\frac{u}{1+x^2}=-\frac{1}{1+x^2}-\ln(1+x^2)+C$$

        \item \textbf{最终形式}
    
        在$x=0$、$y=1$时，$u$为1，由此可代入解得$C=2$，从而代入$u$与$C$有
        $$\frac{1}{y^2}=1+2x^2-(1+x^2)\ln(1+x^2)$$
        由于$y\ne0$，利用连续性可知其恒正，从而
        $$y=\frac{1}{\sqrt{1+2x^2-(1+x^2)\ln(1+x^2)}}$$

        \item (\extra)\textbf{严谨性分析}
        
        直接求导分析可发现存在$\alpha>0$使得$1+2x^2-(1+x^2)\ln(1+x^2)>0$当且仅当$x\in(-\alpha,\alpha)$，由此上述的解只能保证在区间$(-\alpha,\alpha)$确定。

        但是，由于$x\to-\alpha$与$x\to\alpha$时均有$y\to+\infty$，这个解的确保持了$y>0$，且进一步计算验证可发现它不会与其他的$C$对应的解曲线相交。这意味着，上述方程满足$y(0)=1$的解的\textbf{最大存在区间}就是$(-\alpha,\alpha)$，不可能延伸到此区间之外，且在区间内的解只能为上述形式。
    \end{itemize}
}
\end{framed}

\note 有的微分方程可能无法一眼看出类型，这就需要对\textbf{所有类型}有记忆并\textbf{逐个确认}。

\begin{framed}
\begin{exercise}\label{find factor}
    求以下微分变量为$x$的微分方程通解：
    $$(-y\sqrt{x^2+y^2+1}-x(x^2+y^2))\dr x+(x\sqrt{x^2+y^2+1}-y(x^2+y^2))\dr y=0$$
\end{exercise}
\sol{
    分为如下步骤：
    \begin{itemize}
        \item \textbf{初步化简}
        
        首先，由于题中$x^2+y^2$较清晰地构成一个整体，我们将方程整理为
        $$\sqrt{x^2+y^2+1}(x\dr y-y\dr x)-(x^2+y^2)(x\dr x+y\dr y)=0$$
        为了配凑出$x^2+y^2$相关，直接计算可发现
        $$x\dr x+y\dr y=\frac{1}{2}\dr(x^2+y^2)$$
        且$x\ne0$时利用\textbf{重要结论}
        $$\frac{x\dr y-y\dr x}{x^2+y^2}=\dr\arctan\frac{y}{x}$$
        可进一步化为
        $$\sqrt{x^2+y^2+1}(x^2+y^2)\dr\arctan\frac{y}{x}-\frac{1}{2}(x^2+y^2)\dr(x^2+y^2)=0$$
        设$u=x^2+y^2$即得
        $$\sqrt{u+1}u\dr\arctan\frac{y}{x}-\frac{1}{2}u\dr u=0$$
        
        \item \textbf{积分因子}
        
        \note 当$f$、$g$为连续函数时，$f(u)\dr u+g(v)\dr v$\textbf{一定为全微分}(考虑$f$的原函数$F$、$g$的原函数$G$，其原函数为$F(u)+G(v)$)，由此可以发现积分因子。

        考虑$\mu(u)=\frac{1}{\sqrt{u+1}u}$，之前的过程保证了$x\ne0$，此时$u$必然非零，且$u+1>0$，方程两侧乘$\mu(u)$得到
        $$\dr\arctan\frac{y}{x}-\frac{\dr u}{2\sqrt{u+1}}=0$$
        直接积分可发现
        $$\frac{\dr u}{2\sqrt{u+1}}=\dr\sqrt{u+1}$$
        由此这个方程已经成为全微分方程的形式，$\mu(u)$即为积分因子，解为($C$为任意常数)
        $$\arctan\frac{y}{x}-\sqrt{u+1}=C$$
        代入也即
        $$\arctan\frac{y}{x}-\sqrt{x^2+y^2+1}=C$$ 

        \item (\extra)\textbf{严谨性分析}
        
        关于此通解的严谨性分析相对复杂。可以想到，对于左半或右半平面的任何一点初值$(x_0,y_0)$\ (先考虑右半平面$x_0>0$)，都一定存在唯一一个$C_0$使得解在曲线$\arctan\frac{y_0}{x_0}-\sqrt{x_0^2+y_0^2+1}=C_0$上延伸(注意此曲线可能并非函数，因此可能只能延伸出一部分)。

        在这条曲线上令$x\to0$。设此时$y$极限为$y_1$，若$y_1>0$可以得到
        $$\frac{\pi}{2}-\sqrt{y_1^2+1}=C_0$$
        若$y_1<0$则可以得到
        $$-\frac{\pi}{2}-\sqrt{y_1^2+1}=C_0$$
        若能解出符合要求的$y_1$，曲线即可通过$y$轴上的对应点，进一步向左半平面延伸。由此我们可以得到全部不过原点的解(过原点的解讨论起来将更加复杂)。
    \end{itemize}
}
\end{framed}

\note \textbf{非常建议记忆}$\dr\theta=\frac{x\dr y-y\dr x}{x^2+y^2}$，这里$\theta$为极坐标下的极角，在$x\ne0$时也可以写为$\arctan\frac{y}{x}$。此结论在\exref{unique}也有应用。

\note 一般的积分因子寻找往往依赖一些特殊形式的\textbf{配凑}，最终往往需要不断\textbf{换元}并设法配出$f(u)\dr u+g(v)\dr v$。

\begin{framed}
\begin{exercise}
    求以下微分变量为$x$的微分方程通解：
    $$2xy^3+x^2y^2y'=y'$$
\end{exercise}
\sol{
    分为如下步骤：
    \begin{itemize}
        \item \textbf{积分因子}
        
        本题观察与尝试后会发现并不在任何熟悉的类型内，也没有易于分析的配凑方式($2xy^3+3x^2y^2y'$是$(x^2y^3)'$，但这里系数不对导致无法继续分析)，由此尝试构造\textbf{只与$x$或$y$有关的积分因子}。

        将方程写为
        $$2xy^3\dr x+(x^2y^2-1)\dr y=0$$
        并看作$M\dr x+N\dr y$，直接计算可发现$\frac{M_y-N_x}{N}$形式相对复杂，不过
        $$\frac{N_x-M_y}{M}=-\frac{2}{y}$$
        只与$y$有关，由此存在只与$y$有关的积分因子
        $$\mu(y)=\er^{\int_{y_0}^y-\frac{2}{y}\dr y}$$
        计算得可取$\mu(y)=\frac{1}{y^2}$。

        \note 若有很好的观察力，也可能直接注意到此积分因子。

        \item \textbf{通解计算}
        
        当$y\ne0$时，方程两边同乘$y^{-2}$即得
        $$2xy\dr x+x^2\dr y=\frac{\dr y}{y^2}$$
        此时左侧即为$\dr(x^2y)$，右侧即为$-\dr\frac{1}{y}$，这已经是一个全微分方程，从而可写出通解($C$为任意常数)
        $$x^2y+\frac{1}{y}=C$$
        也可移项后写出
        $$y(Cy-x^2y^2)=1$$

        \item \textbf{特解获取}
        
        由于过程中除以了$y^2$，额外讨论$y=0$的情况可发现$y=0$为一个特解。

        \item (\extra)\textbf{严谨性分析}
        
        考虑某个解曲线$y=f(x)$。任给曲线上一点$(x_0,y_0)$，当$y_0\ne0$时根据通解部分的推导知这点附近一定存在$C_0$使得
        $$y_0(C_0y_0-x_0^2y_0^2)=1$$
        另一方面计算可发现不同的$C$对应的通解不交，且恒满足$y\ne0$，由此$y=f(x)$一定是上述$C_0$对应的曲线的一部分，被通解刻画。

        另一方面，上述推理已经证明了一点处$y_0\ne0$则$y$恒不为0，由此若一点处$y_0=0$，解曲线只能为$y=0$。这就证明了通解与特解为全部解。
    \end{itemize}
}
\end{framed}

\note 若对于一般问题没有头绪，可以尝试写成微分的形式后直接计算\textbf{是否存在合适的积分因子}。

\begin{framed}
\begin{exercise}
    求以下微分变量为$x$的微分方程通解：
    $$(3x^2y+2xy+y^3)\dr x+(x^2+y^2)\dr y=0$$
\end{exercise}
\sol{
    分为如下步骤：
    \begin{itemize}
        \item \textbf{积分因子}
        
        本题观察与尝试后会发现并不在任何熟悉的类型内，也没有易于分析的配凑方式，由此尝试构造\textbf{只与$x$或$y$有关的积分因子}。

        将方程看作$M\dr x+N\dr y$，直接计算可发现
        $$\frac{M_y-N_x}{N}=3$$
        只与$x$有关，由此存在只与$y$有关的积分因子
        $$\mu(x)=\er^{\int_{x_0}^x3\dr x}$$
        计算得可取$\mu(x)=\er^{3x}$。

        \item \textbf{通解计算}
        
        方程两边同乘$\er^{3x}$即得
        $$\er^{3x}(3x^2y+2xy+y^3)\dr x+\er^{3x}(x^2+y^2)\dr y=0$$
        虽然已知这是一个全微分方程，并不容易直接看出对应的原函数。由此，利用它在$\rb^2$上有定义，可构造原函数
        $$u(p,q)=\int_{(0,0)}^{(p,q)}\er^{3x}(3x^2y+2xy+y^3)\dr x+\er^{3x}(x^2+y^2)\dr y$$
        取先从$(0,0)$到$(p,0)$、再从$(p,0)$到$(p,q)$的路径，分别参数化可发现
        $$u(p,q)=\int_0^p0\dr t+\int_0^q\er^{3p}(p^2+t^2)\dr t=\er^{3p}\bigg(p^2q+\frac{q^3}{3}\bigg)$$
        这已经得到了一个原函数，从而可写出通解($C$为任意常数)
        $$\er^{3x}\bigg(x^2y+\frac{y^3}{3}\bigg)=C$$

        \item (\extra)\textbf{严谨性分析}
        
        在$\rb^2$上，上述通解对应的$C$可由初值\textbf{唯一确定}。由此可发现不同的$C$对应的曲线不交，从而根据过任一点$(x_0,y_0)$的解必须恒在相应的$C_0$对应的曲线上，这就证明了通解是全部解。
    \end{itemize}    
}
\end{framed}

\note 必须掌握一种\textbf{从全微分还原出原函数}的方法，个人推荐第二型曲线积分，因为可以规避\textbf{定义域}的问题。

\begin{framed}
\begin{exercise}
    求所有$\rb$上可导的函数使得
    $$f'(x)=xf(x)+x\int_0^1tf(t)\dr t$$
\end{exercise}
\sol{
    分为如下步骤：
    \begin{itemize}
        \item \textbf{一阶线性微分方程}
        
        由于积分中为一个常数，我们设它为$a$，方程即
        $$f'(x)-xf(x)=ax$$
        这是一个一阶线性微分方程，直接计算可得$\int(-x)\dr x=-\frac{x^2}{2}+C$，由此可取一个积分因子
        $$\mu(x)=\er^{-x^2/2}$$
        方程两边同乘$\mu(x)$有
        $$\big(\er^{-x^2/2}f(x)\big)'=ax\er^{-x^2}$$
        对右侧不定积分有($C$为任意常数)
        $$\int ax\er^{-x^2/2}\dr x=\frac{a}{2}\int\er^{-x^2/2}\dr x^2=a\int\er^{-x^2/2}\dr\frac{x^2}{2}=-a\er^{-x^2/2}+C$$
        代入最终得到
        $$\er^{-x^2/2}f(x)=-a\er^{-x^2/2}+C$$
        也即
        $$f(x)=-a+C\er^{x^2/2}$$
        
        \item \textbf{代入计算}
        
        直接计算可知
        $$a=\int_0^1f(t)t\dr t=-\int_0^1at\dr t+C\int_0^1t\er^{t^2/2}\dr t$$
        右侧第一项即$-\frac{a}{2}$，第二项类似之前推导换元$s=\frac{t^2}{2}$即
        $$\int_0^{1/2}\er^s\dr s=\sqrt\er-1$$
        由此可得到
        $$\frac{3}{2}a=C(\sqrt\er-1)$$
        解出$a$以最终得到通解
        $$f(x)=C\bigg(\er^{x^2/2}-\frac{2}{3}(\sqrt\er-1)\bigg)$$
    \end{itemize}
}
\end{framed}

\note 包含更抽象部分的题目往往需要先关注本身的\textbf{方程类型}求解后再进一步研究。

\note 除了上面已经出现的问题类型外，注意熟悉教材9.2节的一些\textbf{基本换元方法}与两种特殊的二阶微分方程解法(尤其是不出现$x$时)。9.3节的\textbf{皮卡序列}也请记忆基本计算方式。

\

\noindent\textbf{线性问题}
\begin{framed}
\begin{exercise}
    求以下微分变量为$x$的微分方程通解：
    $$y'=xy+3x+2y+6$$
\end{exercise}
\sol{
    这是一个一阶线性微分方程，可以直接求解，不过我们可以进行一些配凑简化计算。可以发现右边为$(x+2)(y+3)$，由此换元$u=y+3$可知
    $$u'=(x+3)u$$
    这是一个齐次一阶线性微分方程，可以直接得到解为($C$为任意常数)
    $$u=C\er^{\int_0^x(t+3)\dr t}=C\er^{\frac{1}{2}x^2+2x}$$
    由此原方程解为
    $$y=C\er^{\frac{1}{2}x^2+2x}-3$$
}
\end{framed}

\note 注意对于$n$阶线性方程可以证明通解\textbf{一定为全部解}，无需进一步讨论。

\begin{framed}
\begin{exercise}
    求以下微分变量为$x$的微分方程通解：
    $$y''+3y'+2y=\frac{1}{2+\er^x}$$
\end{exercise}
\sol{
    我们介绍两种方法：
    \begin{itemize}
        \item \textbf{标准方法-齐次方程通解}
        
        考虑其特征方程$\lambda^2+3\lambda+2=0$，两根为$-1$、$-2$，由此对应的齐次方程通解为($C_1$、$C_2$为任意常数)
        $$C_1\er^{-x}+C_2\er^{-2x}$$

        \item \textbf{标准方法-常数变易法特解}
        
        由于右端项形式并不符合待定系数法，为计算特解，将$C_1$、$C_2$改写为$C_1(x)$、$C_2(x)$，可以使用常数变易法列出方程组
        $$\begin{cases}C_1'(x)\er^{-x}+C_2'(x)\er^{-2x}=0\\-C_1'(x)\er^{-x}-2C_2'(x)\er^{-2x}=\frac{1}{2+\er^x}\end{cases}$$
        也即得到
        $$C_1'(x)=\frac{\er^x}{2+\er^x},\quad C_2'(x)=-\frac{\er^{2x}}{2+\er^x}$$
        直接积分可得
        $$\int\frac{\er^x}{2+\er^x}\dr x=\int\frac{1}{2+\er^x}\dr\er^x=\ln(2+\er^x)+C$$
        $$\int-\frac{\er^{2x}}{2+\er^x}\dr x=\int-\frac{\er^x}{2+\er^x}\dr\er^x=\int\bigg(\frac{2}{2+\er^x}-1\bigg)\dr\er^x=2\ln(2+\er^x)-\er^x+C$$
        由此均取$C=0$可得到特解
        $$\ln(2+\er^x)\er^{-x}+(2\ln(2+\er^x)-\er^x)\er^{-2x}$$
        综合得通解为
        $$y=C_1\er^{-x}+C_2\er^{-2x}-\er^{-x}+\ln(2+\er^x)(\er^{-x}+2\er^{-2x})$$
        将$\er^{-x}$合并到$C_1$中，可得到更简洁的表达式
        $$y=C_1\er^{-x}+C_2\er^{-2x}+\ln(2+\er^x)(\er^{-x}+2\er^{-2x})$$

        \item \textbf{积分因子法-第一次积分}

        考虑其特征方程$\lambda^2+3\lambda+2=0$，两根为$-1$、$-2$，由此设中间函数
        $$u=y'+y$$
        可得左侧即为$u'+2u$。

        \note 对一般的两根$\lambda_1,\lambda_2\in\rb$，设$u=y'-\lambda_1y$可发现左侧为$u'-\lambda_2u$。

        现在方程化为
        $$u'+2u=\frac{1}{2+\er^x}$$
        左侧的一个积分因子是$\er^{2x}$，同乘$\er^{2x}$后方程变为
        $$(\er^{2x})'=\frac{\er^{2x}}{2+\er^x}$$
        与上一种方法相同可积分得到($C$为任意常数)
        $$\er^{2x}u=\er^x-2\ln(2+\er^x)+C$$
        也即
        $$u=\er^{-x}-2\ln(2+\er^x)\er^{-2x}+C\er^{-2x}$$

        \item \textbf{积分因子法-第二次积分}
        
        此时方程已经化为
        $$y'+y=\er^{-x}-2\ln(2+\er^x)\er^{-2x}+C\er^{-2x}$$
        左侧的一个积分因子是$\er^x$，同乘$\er^x$后方程变为
        $$(\er^xy)'=1-2\ln(2+\er^x)\er^{-x}+C\er^{-x}$$
        直接计算有(换元$t=\er^{-x}$，第二个等号为分部积分)
        $$\int\ln(2+\er^x)\er^{-x}\dr x=-\int\ln(2+t^{-1})\dr t=-t\ln(2+t^{-1})-\int\frac{\dr t}{2t+1}$$
        最终得到(利用$\ln(2\er^{-x}+1)=\ln(2+\er^x)-x$)
        $$\int\ln(2+\er^x)\er^{-x}\dr x=-\er^{-x}\ln(2+\er^x)-\frac{1}{2}\ln(2+\er^x)+\frac{x}{2}+C$$
        由此可算得($C_0$为任意常数)
        $$\er^xy=x+2\er^{-x}\ln(2+\er^x)+\ln(2+\er^x)-x-C\er^{-x}+C_0$$
        从而化简得到
        $$y=(2\er^{-2x}+\er^{-x})\ln(2+\er^x)-C\er^{-2x}+C_0\er^{-x}$$
        可以发现这个通解形式与上一种方法相同。
    \end{itemize}
}
\end{framed}

\note 用特征方程得到积分因子后计算有时比标准方法\textbf{更加简洁}，大家也可以熟悉。不过，一般情况下使用标准方法往往是更直接的。

\begin{framed}
\begin{exercise}
    求以下微分变量为$x$的微分方程通解：
    $$y''-2y'+y=\er^x$$
\end{exercise}
\sol{
    我们介绍两种方法：
    \begin{itemize}
        \item \textbf{标准方法-齐次方程通解}
        
        考虑其特征方程$\lambda^2-2\lambda+1=0$，两根为1、1，由此对应的齐次方程通解为($C_1$、$C_2$为任意常数)
        $$(C_1x+C_2)\er^x$$
        
        \item \textbf{标准方法-待定系数法特解}
        
        根据右端项的形式，查表得可以待定系数$\alpha$寻找特解
        $$y=\alpha x^2\er^x$$
        直接计算可知
        $$y''-2y'+y=2\alpha\er^x$$
        从而对比系数得到可取$\alpha=\frac{1}{2}$，综合得通解为
        $$y=\bigg(\frac{1}{2}x^2+C_1x+C_2\bigg)\er^x$$

        \item \textbf{积分因子法-第一次积分}

        考虑其特征方程$\lambda^2-2\lambda+1=0$，两根为1、1，由此设中间函数
        $$u=y'-y$$
        可得左侧即为$u'-u$。

        现在方程化为
        $$u'-u=\er^x$$
        左侧的一个积分因子是$\er^{-x}$，同乘$\er^{-x}$后方程变为
        $$(\er^{-x}u)'=1$$
        从而积分得到($C$为任意常数)
        $$u=\er^x(x+C)$$

        \item \textbf{积分因子法-第二次积分}
        
        此时方程已经化为
        $$y'+y=\er^(x+C)$$
        左侧的一个积分因子是$\er^{-x}$，同乘$\er^{-x}$后方程变为
        $$(\er^{-x}y)'=x+C$$
        直接计算有($C_0$为任意常数)
        $$y=\bigg(\frac{1}{2}x^2+Cx+C_0\bigg)\er^x$$
        可以发现这个通解形式与上一种方法相同。
    \end{itemize}
}
\end{framed}

\note 一般建议待定系数法记忆一些\textbf{较简单的情况}，其他情况使用常数变易法。

\begin{framed}
\begin{exercise}
    求以下微分变量为$x$的微分方程通解：
    $$y''+4y=\sin(3x)+\er^x$$
\end{exercise}
\sol{
    分为如下步骤：
    \begin{itemize}
        \item \textbf{齐次方程通解}
        
        考虑其特征方程$\lambda^2+4=0$，两根为$\pm2\ir$，由此对应的齐次方程通解为($C_1$、$C_2$为任意常数)
        $$C_1\sin(2x)+C_2\cos(2x)$$
        
        \item \textbf{待定系数法特解}
        
        由于右端分为$\sin(3x)$与$\er^x$两部分，我们考虑
        $$y_1''+4y_1=\sin(3x),\quad y_2''+4y_2=\er^x$$
        对$y_1$，查表得可以待定系数
        $$y_1=\alpha\sin(3x)+\beta\cos(3x)$$
        直接计算可发现
        $$-5\alpha\sin(3x)-5\beta\cos(3x)=\sin(3x)$$
        从而对比系数得到可取$\alpha=-\frac{1}{5}$、$\beta=0$，$y_1$的一个特解为
        $$y_1=-\frac{1}{5}\sin(3x)$$
        对$y_2$，查表得可以待定系数
        $$y_2=\gamma\er^x$$
        直接计算可发现
        $$5\gamma\er^x=\er^x$$
        从而对比系数得到可取$\gamma=\frac{1}{5}$，$y_2$的一个特解为
        $$y_2=\frac{1}{5}\er^x$$

        \item \textbf{综合结果}

        计算可发现$y_1+y_2$是原方程的一个特解，综合得通解为
        $$y=\frac{\er^x-\sin(3x)}{5}+C_1\sin(2x)+C_2\cos(2x)$$
    \end{itemize}
}
\end{framed}

\note 注意\textbf{线性方程}的特解可以\textbf{分部分构造}。

\note 当特征方程的根不为实根时，利用积分因子法构造通解将涉及\textbf{复函数}，不推荐使用。

\begin{framed}
\begin{exercise}
    求以下微分变量为$x$的微分方程通解：
    $$x^2y''+3xy'+2y=x+1$$
\end{exercise}
\sol{
    分为如下步骤：
    \begin{itemize}
        \item \textbf{换元方法}
        
        由于这是一个\textbf{欧拉方程}，在$x>0$时考虑换元$x=\er^t$，利用复合函数求导公式可发现
        $$\frac{\dr y}{\dr t}=y'\frac{\dr x}{\dr t}=xy'$$
        $$\frac{\dr^2y}{\dr t^2}=\frac{\dr xy'}{\dr t}=(xy')'\frac{\dr x}{\dr t}=x^2y''+xy'$$
        由此代入可知
        $$\frac{\dr^2y}{\dr t^2}+2\frac{\dr y}{\dr t}+2y=\er^t+1$$

        \item \textbf{齐次方程通解}
        
        将其看作微分变量为$t$的微分方程，考虑其特征方程$\lambda^2+2\lambda+2=0$，两根为$-1\pm\ir$，由此对应的齐次方程通解为($C_1$、$C_2$为任意常数)
        $$C_1\er^{-t}\cos t+C_2\er^{-t}\sin t$$

        \item \textbf{待定系数法特解}
        
        由于右端分为$\er^t$与1两部分，我们考虑
        $$\frac{\dr^2y_1}{\dr t^2}+2\frac{\dr y_1}{\dr t}+2y_1=\er^t,\quad \frac{\dr^2y_2}{\dr t^2}+2\frac{\dr y_2}{\dr t}+2y_2=1$$
        对$y_1$，查表得可以待定系数
        $$y_1=\alpha\er^t$$
        直接计算可发现
        $$5\alpha\er^t=\er^t$$
        从而对比系数得到可取$\alpha=\frac{1}{5}$，$y_1$的一个特解为
        $$y_1=\frac{1}{5}\er^t$$
        对$y_2$，查表得可以待定系数
        $$y_2=\beta$$
        直接计算可发现
        $$2\beta=1$$
        从而对比系数得到可取$\beta=\frac{1}{2}$，$y_2$的一个特解为
        $$y_2=\frac{1}{2}$$

        \item \textbf{综合结果}
        
        计算可发现$y_1+y_2$是原方程的一个特解，综合得通解为
        $$y=C_1\er^{-t}\cos t+C_2\er^{-t}\sin t+\frac{1}{5}\er^t+\frac{1}{2}$$
        代回$t=\ln x$可以得到$x>0$时
        $$y=C_1\frac{1}{x}\cos\ln x+C_2\frac{1}{x}\sin\ln x+\frac{1}{5}x+\frac{1}{2}$$
        同理，$x<0$时计算有
        $$y=-C_1\frac{1}{x}\cos\ln(-x)-C_2\frac{1}{x}\sin\ln(-x)+\frac{1}{5}x+\frac{1}{2}$$
        两种情况可以合并为通解
        $$y=C_1\frac{1}{|x|}\cos\ln|x|+C_2\frac{1}{|x|}\sin\ln|x|+\frac{1}{5}x+\frac{1}{2}$$

        \note 可以发现，欧拉方程的解往往$x>0$与$x<0$时存在一定\textbf{对称关系}，但具体情况仍然需要一些讨论。

        \item (\extra)\textbf{严谨性分析}
        
        二阶微分方程的严谨性分析将更加复杂，因为初值问题需要给定一点处的$y$与$y'$值，这里仅指出一点：由于$C_1$、$C_2$后的函数在$x\to0$时的极限均不存在，此微分方程过$y$轴的解仅有
        $$y=\frac{1}{5}x+\frac{1}{2}$$
        也就是说，当$x=0$时，只要给定的$y$不为$\frac{1}{2}$或$y'$不为$\frac{1}{5}$，初值问题都将\textbf{无解}。
    \end{itemize}
}
\end{framed}

\note 本题基本综合了这类问题的\textbf{常规处理方法}。

\begin{framed}
\begin{exercise}
    设$y=f(x)$，考虑方程
    $$y'+y=\phi(x)$$
    若$\phi(x)$连续且以$T$为周期($T>0$)，证明存在唯一以$T$为周期的$f$为方程的解。
\end{exercise}
\sol{
    分为如下步骤：
    \begin{itemize}
        \item \textbf{一阶线性微分方程}
        
        利用积分因子法可发现左侧的一个积分因子是$\er^x$，两边同乘$\er^x$得到
        $$(\er^xy)'=\phi(x)\er^x$$
        由此可知通解能写为
        $$f(x)=\er^{-x}\bigg(\int_0^x\phi(t)\er^t\dr t+C\bigg)$$
        要证明存在唯一的以$T$为周期的解，也即证明存在唯一的$C$。

        \item \textbf{进一步分析}
        
        对任何$x$，代入周期性条件$f(x+T)=f(x)$可发现
        $$\er^{-x-T}\bigg(\int_0^{x+T}\phi(t)\er^t\dr t+C\bigg)=\er^{-x}\bigg(\int_0^x\phi(t)\er^t\dr t+C\bigg)$$
        这是对$C$的一元一次方程，可解出
        $$C=\frac{1}{1-\er^{-T}}\bigg(\er^{-T}\int_0^{x+T}\phi(t)\er^t\dr t-\int_0^x\phi(t)\er^t\dr t\bigg)$$
        等式右边是一个关于$x$的函数，记为$g(x)$。

        可以发现，上述求解过程意味着，给定$x$后$f(x)=f(x+T)$当且仅当$C=g(x)$。若不同的$x$对应的$g(x)$不同，不可能存在$C$使得$C=g(x)$恒成立；若不同的$x$对应的$g(x)$相同，任取一个$x$得到的$C$就满足。由此，要证的结论等价于$g(x)$\textbf{是一个与$x$无关的常数}。

        \item \textbf{导数计算}
        
        利用上学期知识，一个函数在$\rb$上是常数当且仅当它的导数恒为0，由此可直接计算$g'(x)$。利用\textbf{变上限积分的导数}结论可知
        $$\frac{\dr}{\dr x}\int_0^{x+T}\phi(t)\er^t\dr t=\phi(x+T)\er^{x+T}$$
        $$\frac{\dr}{\dr x}\int_0^x\phi(t)\er^t\dr t=\phi(x)\er^x$$
        从而
        $$g'(x)=\frac{1}{1-\er^{-T}}\big(\er^{-T}\er^{x+T}\phi(x+T)-\er^x\phi(x)\big)$$
        利用$\phi$的周期性知其为0，原命题得证。

        计算$g(0)$可以发现，$f(x)$以$T$为周期当且仅当
        $$C=\frac{\er^{-T}}{1-\er^{-T}}\int_0^T\phi(t)\er^t\dr t$$
    \end{itemize}
}
\end{framed}

\note 微分方程相关的\textbf{计算性证明}往往是可以通过定义研究出结论的，不需要很复杂的分析。

\note 对于能解出通解的情况，相关的证明往往会需要\textbf{反解任意常数}。

\newpage
\section{下半简介}
\subsection{习题解答}
\note 教材第九章的解答只进行\textbf{通解与特解的获取}，不研究这是否构成全部解等严谨性问题，一些进一步的分析可见本讲义第二十三章或25.2.5。注意理论来说求出的通解可以\textbf{完全不进行化简与合并}，实际几乎只取决于改卷标准，解答中我们尽量化到相对简洁的形式。

\subsubsection{作业解答}
\begin{enumerate}
    \item (习题9.6.1)用常数变易法求以下微分变量为$x$的微分方程通解：
    $$y''+3y'+2y=\frac{1}{\er^x+1}$$

    \sol{
        分为如下步骤：
        \begin{itemize}
            \item \textbf{齐次方程通解}
            
            考虑其特征方程$\lambda^2+3\lambda+2=0$，两根为$-1$、$-2$，由此对应的齐次方程通解为($C_1$、$C_2$为任意常数)
            $$C_1\er^{-x}+C_2\er^{-2x}$$
            
            \item \textbf{常数变易法特解}
            
            为计算特解，将$C_1$、$C_2$改写为$C_1(x)$、$C_2(x)$，可以使用常数变易法列出方程组
            $$\begin{cases}C_1'(x)\er^{-x}+C_2'(x)\er^{-2x}=0\\-C_1'(x)\er^{-x}-2C_2'(x)\er^{-2x}=\frac{1}{1+\er^x}\end{cases}$$
            也即得到
            $$C_1'(x)=\frac{\er^x}{1+\er^x},\quad C_2'(x)=-\frac{\er^{2x}}{1+\er^x}$$
            直接积分可得
            $$\int\frac{\er^x}{1+\er^x}\dr x=\int\frac{1}{1+\er^x}\dr\er^x=\ln(1+\er^x)+C$$
            $$\int-\frac{\er^{2x}}{1+\er^x}\dr x=\int-\frac{\er^x}{1+\er^x}\dr\er^x=\int\bigg(\frac{1}{1+\er^x}-1\bigg)\dr\er^x=\ln(1+\er^x)-\er^x+C$$
            由此均取$C=0$可得到特解
            $$\ln(1+\er^x)\er^{-x}+(\ln(1+\er^x)-\er^x)\er^{-2x}$$
            综合得通解为
            $$y=C_1\er^{-x}+C_2\er^{-2x}-\er^{-x}+\ln(1+\er^x)(\er^{-x}+\er^{-2x})$$
            将$\er^{-x}$合并到$C_1$中，可得到更简洁的表达式
            $$y=C_1\er^{-x}+C_2\er^{-2x}+\ln(1+\er^x)(\er^{-x}+\er^{-2x})$$
        \end{itemize}

        \note 这类题目还可以采用\textbf{积分因子}的方法，见本讲义25.2.5。
    }

    \item (习题9.6.3)用常数变易法求以下微分变量为$x$的微分方程通解：
    $$y''+4y=2\tan x$$

    \sol{
        分为如下步骤：
        \begin{itemize}
            \item \textbf{齐次方程通解}
            
            考虑其特征方程$\lambda^2+4=0$，两根为$\pm2\ir$，由此对应的齐次方程通解为($C_1$、$C_2$为任意常数)
            $$C_1\cos(2x)+C_2\sin(2x)$$
            
            \item \textbf{常数变易法特解}
            
            为计算特解，将$C_1$、$C_2$改写为$C_1(x)$、$C_2(x)$，可以使用常数变易法列出方程组
            $$\begin{cases}C_1'(x)\cos(2x)+C_2'(x)\sin(2x)=0\\-2C_1'(x)\sin(2x)+2C_2'(x)\cos(2x)=2\tan x\end{cases}$$
            也即得到
            $$C_1'(x)=-\sin(2x)\tan x,\quad C_2'(x)=\cos(2x)\tan x$$
            直接积分可得(第二式计算中换元$t=\cos x$)
            $$\int-\sin(2x)\tan x\dr x=\int-2\sin^2x\dr x=\int(\cos(2x)-1)\dr x=\frac{1}{2}\sin(2x)-x+C$$
            $$\begin{aligned}\int\cos(2x)\tan x\dr x&=\int\frac{2\cos^2x-1}{\cos x}\frac{\sin x}{\cos x}\dr x=\int\frac{1-2\cos^2x}{\cos x}\dr\cos x\\ &=\int\bigg(\frac{1}{t}-2t\bigg)\dr t=\ln|t|-t^2+C=\ln|\cos x|-\cos^2x+C\end{aligned}$$
            由此均取$C=0$可得到特解
            $$\bigg(\frac{1}{2}\sin(2x)-x\bigg)\cos(2x)+(\ln|\cos x|-\cos^2x)\sin(2x)$$
            进一步计算可将其化简为
            $$-x\cos(2x)+\ln|\cos x|\sin(2x)-\sin(2x)$$
            由此综合得通解为
            $$y=C_1\cos(2x)+C_2\sin(2x)-x\cos(2x)+\ln|\cos x|\sin(2x)-\sin(2x)$$
            将$-\sin(2x)$合并到$C_2$中，可得到更简洁的表达式
            $$y=C_1\cos(2x)+C_2\sin(2x)-x\cos(2x)+\ln|\cos x|\sin(2x)$$
        \end{itemize}
    }

    \item (习题9.6.5)求以下微分变量为$x$的欧拉方程通解：
    $$x^2y''-4xy'+6y=0$$

    \sol{
        分为如下步骤：
        \begin{itemize}
            \item \textbf{换元方法}
            
            在$x>0$时考虑换元$x=\er^t$，利用复合函数求导公式可发现
            $$\frac{\dr y}{\dr t}=y'\frac{\dr x}{\dr t}=xy'$$
            $$\frac{\dr^2y}{\dr t^2}=\frac{\dr xy'}{\dr t}=(xy')'\frac{\dr x}{\dr t}=x^2y''+xy'$$
            由此代入可知
            $$\frac{\dr^2y}{\dr t^2}-5\frac{\dr y}{\dr t}+6y=0$$

            \item \textbf{通解计算}
            
            将其看作微分变量为$t$的微分方程，考虑其特征方程$\lambda^2-5\lambda+6=0$，两根为2、3，由此通解为($C_1$、$C_2$为任意常数)
            $$C_1\er^{2t}+C_2\er^{3t}$$

            \item \textbf{综合结果}
            
            代回$t=\ln x$可以得到$x>0$时
            $$y=C_1x^2+C_2x^3$$
            同理，$x<0$时计算有
            $$y=C_1x^2+C_2x^3$$
            进一步计算可验证上述表达式在$x=0$时也满足微分方程，从而是合理的通解形式。
        \end{itemize}
    }

    \item (习题9.6.8)求以下微分变量为$x$的欧拉方程通解：
    $$x^2y''+xy'+4y=10$$

    \sol{
        分为如下步骤：
        \begin{itemize}
            \item \textbf{换元方法}
            
            在$x>0$时考虑换元$x=\er^t$，利用复合函数求导公式可发现
            $$\frac{\dr y}{\dr t}=y'\frac{\dr x}{\dr t}=xy'$$
            $$\frac{\dr^2y}{\dr t^2}=\frac{\dr xy'}{\dr t}=(xy')'\frac{\dr x}{\dr t}=x^2y''+xy'$$
            由此代入可知
            $$\frac{\dr^2y}{\dr t^2}+4y=10$$

            \item \textbf{通解计算}
            
            \note 仍然可以通过先计算齐次方程通解再计算非齐次方程特解的方式得到上述方程通解，这里进行一些处理以\textbf{简化计算}。

            由于方程形式，换元$u=y-\frac{5}{2}$，则$u$的导函数与$y$相同，从而方程变为
            $$\frac{\dr^2u}{\dr t^2}+4u=0$$
            
            将其看作微分变量为$t$的微分方程，考虑其特征方程$\lambda^2+4=0$，两根为$\pm2\ir$，由此通解为($C_1$、$C_2$为任意常数)
            $$u=C_1\cos(2t)+C_2\sin(2t)$$
            于是
            $$y=C_1\cos(2t)+C_2\sin(2t)+\frac{5}{2}$$

            \item \textbf{综合结果}
            
            代回$t=\ln x$可以得到$x>0$时
            $$y=C_1\cos(2\ln x)+C_2\sin(2\ln x)+\frac{5}{2}$$
            同理，$x<0$时计算有
            $$y=C_1\cos(2\ln(-x))+C_2\sin(2\ln(-x))+\frac{5}{2}$$
            两种情况可以合并为通解
            $$y=C_1\cos(2\ln|x|)+C_2\sin(2\ln|x|)+\frac{5}{2}$$

            \note 由于此通解在$x=0$处无定义，无需额外验证。
        \end{itemize}
    }

    \item (习题9.7.1(4))用消元法求解
    $$\begin{cases}\frac{\dr x}{\dr t}=x\\\frac{\dr y}{\dr t}=-y+\sqrt2z\\\frac{\dr z}{\dr t}=\sqrt2y\end{cases}$$

    \sol{
        第一个方程与其他方程无关，可以直接解出($C_1$为任意常数)
        $$x=C_1\er^t$$
        为了消去第二个方程中的$z$，对第二个方程求导得到
        $$\frac{\dr^2y}{\dr t^2}=-\frac{\dr y}{\dr t}+\sqrt2\frac{\dr z}{\dr t}=-\frac{\dr y}{\dr t}+2y$$
        这是一个对$y$的常系数二阶线性齐次方程，考虑其特征方程$\lambda^2+\lambda-2=0$，两根为$-2$、1，由此通解为($C_2$、$C_3$为任意常数)
        $$y=C_2\er^{-2t}+C_3\er^t$$
        代入第二个方程可直接计算出
        $$z=\frac{1}{\sqrt2}\bigg(-C_2\er^{-2t}+2C_3\er^t\bigg)$$
    }

    \item (习题9.7.2(1))求以下微分变量为$t$的非齐次微分方程组的一个特解
    $$\begin{cases}x'=x+2y\\y'=x-5\sin t\end{cases}$$

    \sol{
        为了消去第一个方程中的$y$，对第一个方程求导得到
        $$x''=x'+2y'=x'+2(x-5\sin t)$$
        将其化为
        $$x''-x'-2x=-10\sin t$$
        计算可知$x''-x'-2x$的特征方程为$\lambda^2-\lambda-2=0$，两根为$-1$、2，查表得可以待定系数
        $$x=\alpha\cos t+\beta\sin t$$
        直接计算可发现
        $$(-3\alpha-\beta)\cos t+(\alpha-3\beta)\sin t=-10\sin t$$
        从而对比系数得到可取$\alpha=-1$、$\beta=3$，$x$的一个特解为
        $$x=-\cos t+3\sin t$$
        代入第一个方程可直接计算出对应的特解
        $$y=2\cos t-\sin t$$
    }
    
    \item (总练习题9.8)设$f(x)$在$[0,+\infty)$连续，$\varphi(x)$在$[0,+\infty)$可微且满足
    $$\forall x\ge0,\quad\varphi'(x)+f(x)\varphi(x)\le0$$
    证明
    $$\forall x\ge0,\quad\varphi(x)\le\varphi(0)\er^{-\int_0^xf(t)\dr t}$$

    \sol{
        我们先分析再证明：
        \begin{itemize}
            \item \textbf{分析形式}
            
            先观察条件本身的形式。条件中出现了一个特殊的函数$\er^{-\int_0^xf(t)\dr t}$，而这类型的表达式在\textbf{一阶常系数线性微分方程}中常见。利用相关的理论，可发现它和另一个条件$\varphi'(x)+f(x)\varphi(x)$的\textbf{积分因子}类似，具体来说，利用复合函数求导公式与变上限积分的导数结论，计算可以得到
            $$(\varphi'(x)+f(x)\varphi(x))\er^{\int_0^xf(t)\dr t}=\frac{\dr}{\dr x}\bigg(\varphi(x)\er^{\int_0^xf(t)\dr t}\bigg)$$
            这可以将条件转化为一个函数单调减，而这已经接近结论。

            \item \textbf{最终证明}
            
            利用一阶线性微分方程的积分因子解法，可以发现(将$\varphi$看作未知函数)
            $$(\varphi'(x)+f(x)\varphi(x))\er^{\int_0^xf(t)\dr t}=\frac{\dr}{\dr x}\bigg(\varphi(x)\er^{\int_0^xf(t)\dr t}\bigg)$$
            由于乘的积分因子$\er^{\int_0^xf(t)\dr t}$恒正，第一个不等式即可改为
            $$\forall x\ge0,\quad\frac{\dr}{\dr x}\bigg(\varphi(x)\er^{\int_0^xf(t)\dr t}\bigg)$$
            也即函数
            $$\Phi(x)=\varphi(x)\er^{\int_0^xf(t)\dr t}$$
            在$[0,+\infty)$单调递减。根据结论形式，我们考虑$\Phi(0)\le\Phi(x)$，将其展开得到
            $$\varphi(0)\le\varphi(x)\er^{\int_0^xf(t)\dr t}$$
            移项即得结论。
        \end{itemize}
    }
\end{enumerate}

\subsubsection{补充题解答}
\begin{enumerate}
    \item (习题9.4.2)设$\varphi_1(x)$、$\varphi_2(x)$是$(a,b)$上的微分方程
    $$y''(x)+q(x)y(x)=0$$
    的任意两个解(这里$q(x)$在$(a,b)$连续)，证明$\varphi_1(x)$与$\varphi_2(x)$的朗斯基行列式$W(x)$在$(a,b)$为常数。

    \sol{
        利用上学期知识，一个函数在$(a,b)$上是常数当且仅当它的导数恒为0，由此可直接计算$W'(x)$。直接利用乘积求导公式可知
        $$W'(x)=\varphi_1'(x)\varphi_2'(x)+\varphi_1(x)\varphi_2''(x)-\varphi_2'(x)\varphi_1'(x)-\varphi_2(x)\varphi_1''(x)=\varphi_1(x)\varphi_2''(x)-\varphi_2(x)\varphi_1''(x)$$
        由条件有
        $$\varphi_1''(x)=-q(x)\varphi_1(x),\quad\varphi_2''(x)=-q(x)\varphi_2(x)$$
        从而代入即得到$W'(x)=0$，得证。
    }

    \item (习题9.4.6)考虑$(a,b)$上的微分方程
    $$y''(x)+p(x)y'(x)+q(x)y(x)=f(x)$$
    这里$p(x)$、$q(x)$在$(a,b)$连续。若它对应的齐次方程的通解是
    $$C_1\varphi_1(x)+C_2\varphi_2(x),\quad C_1,C_2\in\rb$$
    且此方程的一个特解是$y^*(x)$，证明它的任何一个解$y(x)$都存在$C_1,C_2\in\rb$使得
    $$y(x)=C_1\varphi_1(x)+C_2\varphi_2(x)+y^*(x)$$

    \sol{
        设$u(x)=y(x)-y^*(x)$，直接计算可发现
        $$u''(x)+p(x)u'(x)+q(x)u(x)=f(x)-f(x)=0$$
        由此$u(x)$是原方程对应的齐次方程的一个解，由本节定理4，必然存在$C_1,C_2\in\rb$使得
        $$u(x)=C_1\varphi_1(x)+C_2\varphi_2(x)$$
        由此即有
        $$y(x)=C_1\varphi_1(x)+C_2\varphi_2(x)+y^*(x)$$
        得证。
    }
    
    \item (习题9.5.1(2))求以下微分变量为$x$的微分方程的通解：
    $$4y''+5y'+y=0$$

    \sol{
        考虑其特征方程$4\lambda^2+5\lambda+1=0$，两根为$-\frac{1}{4}$、$-1$，由此通解为($C_1$、$C_2$为任意常数)
        $$C_1\er^{-x/4}+C_2\er^{-x}$$
    }

    \item (习题9.5.2(1))求以下微分变量为$x$的初值问题的解：
    $$y''+2y'+4y=0,\quad y(0)=1,\quad y'(0)=-1$$

    \sol{
        考虑其特征方程$\lambda^2+2\lambda+4=0$，两根为$-1\pm\sqrt3\ir$，由此通解为($C_1$、$C_2$为任意常数)
        $$y(x)=\er^{-x}(C_1\cos(\sqrt3x)+C_2\sin(\sqrt3x))$$
        代入初值条件可知
        $$\begin{cases}C_1=1\\C_2-C_1=-1\end{cases}$$
        从而解得$C_1=1$、$C_2=0$，最终有
        $$y(x)=\er^{-x}\cos(\sqrt3x)$$
    }

    \item (习题9.5.3(3))用待定系数法求以下微分变量为$x$的微分方程的一个特解：
    $$y''-9y'+20y=x+1$$

    \sol{
        计算可知$y''-9y'+20y$的特征方程为$\lambda^2-9\lambda+20=0$，两根为4、5。查表得可以待定系数
        $$y=\alpha x+\beta$$
        直接计算可发现
        $$-9\alpha+20\alpha x+20\beta=x+1$$
        从而对比系数得到可取$\alpha=\frac{1}{20}$、$\beta=\frac{29}{400}$，一个特解为
        $$y=\frac{1}{20}x+\frac{29}{400}$$
    }

    \item (习题9.5.4(3))写出以下微分变量为$x$的微分方程待定系数的特解形式：
    $$y''+y=\er^{3x}(x-2)$$

    \sol{
        计算可知$y''+y$的特征方程为$\lambda^2+1=0$，两根为$\pm\ir$。查表得可以待定系数
        $$y=\er^{3x}(\alpha x+\beta)$$
    }

    \item (习题9.6.2)用常数变易法求以下微分变量为$x$的微分方程通解：
    $$y''+y=\frac{1}{\sin x}$$

    \sol{
        分为如下步骤：
        \begin{itemize}
            \item \textbf{齐次方程通解}
            
            考虑其特征方程$\lambda^2+1=0$，两根为$\pm\ir$，由此对应的齐次方程通解为($C_1$、$C_2$为任意常数)
            $$C_1\cos x+C_2\sin x$$
            
            \item \textbf{常数变易法特解}
            
            为计算特解，将$C_1$、$C_2$改写为$C_1(x)$、$C_2(x)$，可以使用常数变易法列出方程组
            $$\begin{cases}C_1'(x)\cos x+C_2'(x)\sin x=0\\-C_1'(x)\sin x+C_2'(x)\cos x=\frac{1}{\sin x}\end{cases}$$
            也即得到
            $$C_1'(x)=-1,\quad C_2'(x)=\frac{\cos x}{\sin x}$$
            直接积分可得
            $$\int-1\dr x=-x+C$$
            $$\int\frac{\cos x}{\sin x}\dr x=\int\frac{\dr\sin x}{\sin x}=\ln|\sin x|+C$$
            由此均取$C=0$可得到特解
            $$-x\cos x+\sin x\ln|\sin x|$$
            由此综合得通解为
            $$y=C_1\cos x+C_2\sin x-x\cos x+\sin x\ln|\sin x|$$
        \end{itemize}
    }

    \item (习题9.6.6)求以下微分变量为$x$的欧拉方程通解：
    $$x^2y''-xy'-3y=0$$

    \sol{
        分为如下步骤：
        \begin{itemize}
            \item \textbf{换元方法}
            
            在$x>0$时考虑换元$x=\er^t$，利用复合函数求导公式可发现
            $$\frac{\dr y}{\dr t}=y'\frac{\dr x}{\dr t}=xy'$$
            $$\frac{\dr^2y}{\dr t^2}=\frac{\dr xy'}{\dr t}=(xy')'\frac{\dr x}{\dr t}=x^2y''+xy'$$
            由此代入可知
            $$\frac{\dr^2y}{\dr t^2}-2\frac{\dr y}{\dr t}-3y=0$$

            \item \textbf{通解计算}
            
            将其看作微分变量为$t$的微分方程，考虑其特征方程$\lambda^2-2\lambda-3=0$，两根为$-1$、3，由此通解为($C_1$、$C_2$为任意常数)
            $$C_1\er^{-t}+C_2\er^{3t}$$

            \item \textbf{综合结果}
            
            代回$t=\ln x$可以得到$x>0$时
            $$y=\frac{C_1}{x}+C_2x^3$$
            同理，$x<0$时计算有
            $$y=\frac{C_1}{x}+C_2x^3$$
            这已经是合理的通解形式。
        \end{itemize}
    }

    \item (习题9.7.1(2))用消元法求解
    $$\begin{cases}\frac{\dr x}{\dr t}=x+y+2\er^t\\\frac{\dr y}{\dr t}=4x+y-\er^t\end{cases}$$

    \sol{
        从第一个方程可以得到
        $$y=\frac{\dr x}{\dr t}-x-2\er^t$$
        代入第二个方程计算出
        $$\frac{\dr^2x}{\dr t^2}-\frac{\dr x}{\dr t}-2\er^t=4x+\frac{\dr x}{\dr t}-x-2\er^t-\er^t$$
        整理得
        $$\frac{\dr^2x}{\dr t^2}-2\frac{\dr x}{\dr t}-3x=-\er^t$$
        此方程的求解分为两步：
        \begin{itemize}
            \item \textbf{齐次方程通解}
            
            考虑其特征方程$\lambda^2-2\lambda-3=0$，两根为$-1$、3，由此对应的齐次方程通解为($C_1$、$C_2$为任意常数)
            $$C_1\er^{-t}+C_2\er^{3t}$$

            \item \textbf{待定系数法特解}
            
            查表得可以待定系数
            $$x=\alpha\er^t$$
            直接计算可发现
            $$-4\alpha\er^t=-\er^t$$
            从而对比系数得到可取$\alpha=\frac{1}{4}$，一个特解为
            $$x=\frac{1}{4}\er^t$$
            综合得通解为
            $$x=C_1\er^{-t}+C_2\er^{3t}+\frac{1}{4}\er^t$$
        \end{itemize}
        代入第一个方程可以直接计算出
        $$y=-2C_1\er^{-t}+2C_2\er^{3t}-2\er^t$$
    }

    \item (习题9.7.2(4))求以下微分变量为$t$的非齐次微分方程组的一个特解：
    $$\begin{cases}x'=2x-y\\y'=x+2\er^t\end{cases}$$
    
    \sol{
        为了消去第一个方程中的$y$，对第一个方程求导得到
        $$x''=2x'-y'=2x'-(x+2\er^t)$$
        将其化为
        $$x''-2x'+x=-2\er^t$$
        计算可知$x''-2x'+x$的特征方程为$\lambda^2-2\lambda+1=0$，两根为1、1，查表得可以待定系数
        $$x=\alpha t^2\er^t$$
        直接计算可发现
        $$2\alpha\er^t=-2\er^t$$
        从而对比系数得到可取$\alpha=-1$，$x$的一个特解为
        $$x=-t^2\er^t$$
        代入第一个方程可直接计算出对应的特解
        $$y=(-t^2+2t)\er^t$$
    }

    \item (总练习题9.6)考虑$(a,b)$上的微分方程
    $$y''(x)+p(x)y'(x)+q(x)y(x)=0$$
    这里$p(x)$、$q(x)$在$(a,b)$\textbf{可微}。证明存在函数$u(x)$使得作变换$y(x)=u(x)z(x)$后方程化为了
    $$z''(x)+Q(x)z(x)=0$$
    的形式。

    \sol{
        直接代入$y(x)=u(x)z(x)$可计算得
        $$u''(x)z(x)+2u'(x)z'(x)+u(x)z''(x)+p(x)u(x)z'(x)+p(x)u'(x)z(x)+q(x)u(x)z(x)=0$$
        整理得
        $$u(x)z''(x)+(2u'(x)+p(x)u(x))z'(x)+(u''(x)+p(x)u'(x)+q(x)u(x))z(x)=0$$
        若要符合目标形式，$z'(x)$的系数必须为0，也即
        $$2u'(x)+p(x)u(x)=0$$
        这是一阶齐次线性微分方程，给定$(a,b)$中的一点$x_0$，可写出它的一个特解
        $$u(x)=\er^{-\frac{1}{2}\int_{x_0}^xp(t)\dr t}$$
        由于它非零，可以直接两侧同除以$u(x)$，由此代入计算可发现原方程化为了
        $$z''(x)+\bigg(-\frac{1}{2}p'(x)-\frac{1}{4}p^2(x)+q(x)\bigg)z(x)=0$$
        从而上述$u(x)$的确符合要求。
    }

    \item (总练习题9.7)给定地球半径$R$与重力加速度$g$，考虑垂直脱离地面的火箭，设其初速度为$v_0$，且只受地球引力，证明其高度$h$随时间的变化满足
    $$h''(t)=\frac{-g}{(1+\frac{h(t)}{R})^2}$$
    由此证明速度$v$随时间的变化满足
    $$\frac{1}{2}(v^2(t)-v_0^2)=\frac{Rg}{1+\frac{h(t)}{R}}-Rg$$
    若火箭发射后不再返回，也即
    $$\lim_{t\to+\infty}h(t)=+\infty$$
    求$v_0$的最小值。

    \sol{
        分为如下步骤：
        \begin{itemize}
            \item \textbf{方程构造}
            
            设向下的力为$F(t)$，根据牛顿第二定律可知
            $$mh''(t)=-F(t)$$
            设地球质量为$M$、引力常数为$G$，通过引力公式可知
            $$mg=\frac{GMm}{R^2},\quad F(t)=\frac{GMm}{(R+h(t))^2}$$
            两式相除得
            $$\frac{F(t)}{mg}=\frac{R^2}{h^2(t)}$$
            代入牛顿第二定律并化简即得方程。

            \item \textbf{性质证明}
            
            这里速度代表向上速度，由此有$v(t)=h'(t)$。由于原方程是一个不含未知数$t$的二阶微分方程，利用教材中的分析可知(这里右侧将$v$、$h$也同时看作$t$的表达式)
            $$h''(t)=v\frac{\dr v}{\dr h}$$
            由此有
            $$v\frac{\dr v}{\dr h}=\frac{-g}{(1+\frac{h}{R})^2}$$
            这已经是变量分离的方程，也即
            $$v\dr v=\frac{-g\dr h}{(1+\frac{h}{R})^2}$$
            同积分得到(右侧是$h$的幂函数复合一次函数)
            $$\frac{1}{2}v^2=\frac{Rg}{1+\frac{h}{R}}+C$$
            根据定义可知$v(0)=v_0$为初速度，$h(0)=0$为初始高度，考虑$t=0$可算出
            $$C=\frac{1}{2}v_0^2-Rg$$
            由此代入并移项最终得到要证明的
            $$\frac{1}{2}(v^2(t)-v_0^2)=\frac{Rg}{1+\frac{h(t)}{R}}-Rg$$

            \item \textbf{极限情况}
            
            由条件在上式两侧同取极限，右侧利用复合函数极限结论可知
            $$\lim_{t\to+\infty}\bigg(\frac{Rg}{1+\frac{h(t)}{R}}-Rg\bigg)=-Rg$$
            于是
            $$\lim_{t\to\infty}\frac{1}{2}(v^2(t)-v_0^2)=-Rg$$
            也即
            $$\lim_{t\to\infty}v^2(t)=v_0^2-2Rg$$
            利用极限保序性，此式可能成立至少需要$v_0^2\ge 2Rg$，即$v_0\ge\sqrt{2Rg}$。

            \item (\extra)\textbf{严谨性分析}
            
            上述过程只能说明至少需要$v_0\ge\sqrt{2Rg}$，我们下面证明$v_0\ge\sqrt{2Rg}$时的确能够保证$h(t)$趋于正无穷。

            首先，根据我们得到的$v(t)$与$h(t)$的函数关系，可以发现在$h(t)\ge0$的任何一点处都有
            $$v^2(t)=v_0^2+2Rg\bigg(\frac{1}{1+\frac{h(t)}{R}}-1\bigg)\ge2Rg-2Rg\bigg(\frac{h(t)}{R+h(t)}\bigg)=\frac{2R^2g}{R+h(t)}>0$$
            由于初速度$v_0>0$，利用极限保序性可知$t=0$向后的一小段时间$t\in[0,\delta]$内均有$h(t)>0$。利用\textbf{介值定理}可知$h(t)\ge0$时$v(t)$应当同号，从而$[0,\delta]$上$v(t)>0$恒成立，而这可以说明此段区间$h(t)$严格单调增，从而$h(\delta)>h(0)=0$。以此类推，即可证明$h(t)$在$[0,+\infty)$恒严格单调增。

            \note 本段的严谨说明需要利用拓扑，这里仅给出一个直观理解。

            与数列时完全类似可证明单调增且有界的函数必然在$+\infty$有极限。由此，若$h(t)$不趋于正无穷，必然存在$A>0$使得
            $$\lim_{t\to+\infty}h(t)=A$$
            代入计算可知
            $$\lim_{t\to\infty}v^2(t)=v_0^2+2\bigg(\frac{Rg}{1+\frac{A}{R}}-Rg\bigg)\ge\frac{2R^2g}{R+A}>0$$
            但利用极限保序性，这意味着当$t$充分大时恒有
            $$v(t)>\frac{1}{2}\sqrt{\frac{2R^2g}{R+A}}$$
            也即$h(t)$将大于等于斜率为$\frac{1}{2}\sqrt{\frac{2R^2g}{R+A}}$的一次函数，不可能存在有限极限，矛盾。
        \end{itemize}
    }
\end{enumerate}

\subsection{期中考试}
\subsubsection{试题}
\begin{enumerate}
    \item 求方程
    $$\frac{\dr y}{\dr x}=\frac{y}{x+y^2}$$
    的满足条件$y(1)=1$的解。

    \item 计算
    $$\iint_D(\sqrt2x+\sqrt3y)^2\dr x\dr y$$
    其中$D$是由曲线$(2x^2+3y^2)^2=2x^2-3y^2$围成的有界集合。

    \item 计算
    $$\iiint_\Omega\bigg(\frac{x}{1+x^2}+\frac{y}{1+y^2}\bigg)\dr x\dr y\dr z$$
    其中$\Omega$是由曲面$x=0$、$x=1$、$2(1+x^2)=y^2+\frac{z^2}{3}$围成的区域。

    \item 计算
    $$\iiint_\Omega(x+y+z)\cos((x+y+z)^2)\dr x\dr y\dr z$$
    其中
    $$\Omega=\big\{(x,y,z)\mid0\le x-y\le1,\ 0\le x-z\le1,\ 0\le x+y+z\le1\big\}$$

    \item 计算
    $$\oiint_\Sigma\Fv\cdot\nv\dr S$$
    其中$\Fv=(P,Q,R)$，且
    $$P=\frac{x}{(x^2+y^2+z^2)^{3/2}}+yz,\quad Q=\frac{y}{(x^2+y^2+z^2)^{3/2}}+xz,\quad R=\frac{z}{(x^2+y^2+z^2)^{3/2}}+xy$$
    $\Sigma$为椭球面$\frac{x^2}{4}+\frac{y^2}{9}+\frac{z^2}{25}=1$，$\nv$为$\Sigma$对应位置的单位外法向量。

    \item 计算
    $$\oint_\Gamma(y^2-z^2)\dr x+(z^2-x^2)\dr y+(x^2-y^2)\dr z$$
    其中$\Gamma$为平面$x+y+z=1$与球面$x^2+y^2+z^2=1$的交线，且从$z$轴正向看去沿逆时针方向。

    \item
    \begin{enumerate}[(1)]
        \item 设$\Omega\subset\rb^3$为区域，且$f(x,y,z)$在$\Omega$上连续。若在$\Omega$的任何有界子区域$\Omega_1$上均有
        $$\iiint_{\Omega_1}f(x,y,z)\dr x\dr y\dr z=0$$
        证明
        $$\forall(x,y,z)\in\Omega,\quad f(x,y,z)=0$$

        \item 设$f(x)$在$(0,+\infty)$有连续的一阶导数，且$\lim_{x\to0^+}f(x)$存在。若对$\Omega=\{(x,y,z)\mid x>0\}$中的任何一个光滑简单封闭曲面$S$均有
        $$\oiint_{S^+}xf(x)\dr y\dr z-xyf(x)\dr z\dr x-z\er^{2x}\dr x\dr y=0$$
        求$f(x)$的表达式。
    \end{enumerate}
\end{enumerate}

\subsubsection{解答}
\begin{enumerate}
    \item 介绍两种方法：
    \begin{itemize}
        \item \textbf{积分因子法}
        
        本题观察与尝试后会发现并不在任何熟悉的类型内，也没有易于分析的配凑方式，由此尝试构造\textbf{只与$x$或$y$有关的积分因子}。

        将方程写为
        $$y\dr x-(x+y^2)\dr y=0$$
        并看作$M\dr x+N\dr y$，直接计算可发现$\frac{M_y-N_x}{N}$形式相对复杂，不过
        $$\frac{N_x-M_y}{M}=-\frac{2}{y}$$
        只与$y$有关，由此存在只与$y$有关的积分因子
        $$\mu(y)=\er^{\int_{y_0}^y-\frac{2}{y}\dr y}$$
        计算得可取$\mu(y)=\frac{1}{y^2}$。当$y\ne0$时，方程两边同乘$\mu(y)$即得
        $$\frac{1}{y}\dr x-\frac{x}{y^2}\dr y-\dr y=0$$
        此时计算发现这即为
        $$\dr\frac{x}{y}-\dr y=\dr\bigg(\frac{x}{y}-y\bigg)=0$$
        这已经是一个全微分方程，从而可写出通解($C$为任意常数)
        $$\frac{x}{y}-y=C$$

        \note 也可采用其他\textbf{确定全微分原函数}的方法计算全微分方程的通解。

        代入初值条件$x=1$、$y=1$即得$C=0$，从而解最终为
        $$\frac{x}{y}-y=0$$
        将$y$重新写为$x$的函数即得结果($x$的范围是为了使其中$y\ne0$)
        $$y=\sqrt{x},\quad x\in(0,+\infty)$$
        
        \item \textbf{看作$x$的方程}
        
        \note 本讲义25.1.2第5题(2)给出了这样的处理思路。

        当$y\ne0$时(结合初值条件只需考虑$y>0$)，利用反函数求导公式可知
        $$\frac{\dr x}{\dr y}=\frac{x+y^2}{y}$$
        
        从而此方程可以看作关于$x$、微分变量为$y$的一阶线性微分方程
        $$\frac{\dr x}{\dr y}-\frac{x}{y}=y$$
        由$y>0$，直接计算可得$\int\frac{-\dr y}{y}=-\ln y+C$，由此可取一个积分因子
        $$\mu(y)=\er^{-\ln y}=\frac{1}{y}$$
        方程两边同乘$\mu(y)$有
        $$\frac{\dr}{\dr y}\frac{x}{y}=1$$
        对右侧不定积分有($C$为任意常数)
        $$\frac{x}{y}=y+C$$
        代入初值条件$x=1$、$y=1$即得$C=0$，从而解最终为
        $$\frac{x}{y}=y$$
        将$y$重新写为$x$的函数即得结果
        $$y=\sqrt{x},\quad x\in(0,+\infty)$$
    \end{itemize}

    \note 进一步严谨性分析(\extra)：由于在$x=0$处$y=\sqrt{x}$不可导，上述解的定义域$(0,+\infty)$已经是初值问题解的最大存在区间。

    \item 分为如下步骤：
    \begin{itemize}
        \item \textbf{区域描述}
            
        根据``围成''的有界性要求，有界集合应当对应
        $$(2x^2+3y^2)^2\le 2x^2-3y^2$$
        $D$即为满足上述不等式的$(x,y)$集合。

        \item \textbf{反射对称化简}
        
        将被积函数\textbf{展开}得到其为
        $$2x^2+3y^2+2\sqrt{6}xy$$
            
        由于区域$D$满足$(x,y)\in D$当且仅当$(-x,y)\in D$，且$xy=-(-x)y$，利用函数折半可发现
        $$\iint_Dxy\dr x\dr y=0$$
        由此可将原积分化为
        $$\iint_D(2x^2+3y^2)\dr x\dr y$$
            
        \item \textbf{旋转对称性换元}
            
        由于区域中$2x^2$与$3y^2$可看作整体，且区域描述、被积函数都与$2x^2+3y^2$有关，考虑\textbf{广义极坐标换元}
        $$x=\frac{1}{\sqrt2}r\cos\theta,\quad y=\frac{1}{\sqrt3}r\sin\theta$$
        此时，被积函数为$r^2$，且计算Jacobi行列式得
        $$\dr x\dr y=\frac{1}{\sqrt6}r\dr r\dr\theta$$

        区域的条件可写为
        $$r^2\le\cos^2\theta-\sin^2\theta=\cos(2\theta)$$
        将上述条件与极坐标自然范围确定的$(r,\theta)$区域记为$D_0$，原积分化为
        $$\frac{1}{\sqrt6}\iint_{D_0}r^3\dr r\dr\theta$$

        \item \textbf{累次积分化}
            
        由于$r$自然被$\theta$约束，选择先对$r$再对$\theta$积分。此时，我们需要将集合描述成
        $$\alpha\le\theta\le\beta,\quad\phi_1(\theta)\le r\le\phi_2(\theta)$$
        由于$r^2\le\cos(2\theta)$需要能确定合理的$r$，必然有
        $$\cos(2\theta)\ge0$$
        因此结合极坐标自然范围知$\theta$的范围对应
        $$0\le\theta\le\frac{\pi}{4},\quad\frac{3\pi}{4}\le\theta\le\frac{5\pi}{4},\quad\frac{7\pi}{4}\le\theta\le2\pi$$
        三部分，对每部分$r$的约束均为
        $$0\le r\le\sqrt{\cos(2\theta)}$$
        从而可最终写出积分
        $$\frac{1}{\sqrt6}\bigg(\int_0^{\pi/4}\dr\theta\int_0^{\sqrt{\cos(2\theta)}}r^3\dr r+\int_{3\pi/4}^{5\pi/4}\dr\theta\int_0^{\sqrt{\cos(2\theta)}}r^3\dr r+\int_{7\pi/4}^{2\pi}\dr\theta\int_0^{\sqrt{\cos(2\theta)}}r^3\dr r\bigg)$$

        \item \textbf{积分计算}
            
        由于$\cos(2\theta)=\cos(2(2\pi-\theta))$，在上述第三个积分中将$\theta$换元为$2\pi-\theta$可发现值与第一个积分相等。将第二个积分拆分为$3\pi/4$到$\pi$与$\pi$到$5\pi/4$两段，分别将$\theta$换元为$\pi-\theta$与$\theta-\pi$可发现值与第一个积分相等。综合以上，结果可化为
        $$\frac{4}{\sqrt6}\int_0^{\pi/4}\dr\theta\int_0^{\sqrt{\cos(2\theta)}}r^3\dr r$$
        直接对$r$积分得到其为
        $$\frac{1}{\sqrt6}\int_0^{\pi/4}\cos^2(2\theta)\dr\theta$$
        换元$r=2\theta$得到其为
        $$\frac{1}{2\sqrt6}\int_0^{\pi/2}\cos^2\theta\dr\theta$$
        利用对称性将$\theta$换元为$\frac{\pi}{2}-\theta$可得
        $$\int_0^{\pi/2}\cos^2\theta\dr\theta=\int_0^{\pi/2}\sin^2\theta\dr\theta$$
        从而
        $$\frac{1}{2\sqrt6}\int_0^{\pi/2}\cos^2\theta\dr\theta=\frac{1}{4\sqrt6}\int_0^{\pi/2}(\cos^2\theta+\sin^2\theta)\dr\theta=\frac{\sqrt6\pi}{48}$$
    \end{itemize}

    \item 分为如下步骤：
    \begin{itemize}
        \item \textbf{区域描述}
            
        根据``围成''需要以每个曲面为边界，$x=0$、$x=1$对应的约束必然为
        $$0\le x\le 1$$
        进一步根据有界性要求，第三个曲面对应的约束必然为(否则$y$、$z$将无界)
        $$y^2+\frac{z^2}{3}\le2(1+x^2)$$
        $\Omega$即为满足上述不等式的$(x,y,z)$集合。

        \item \textbf{反射对称化简}
            
        由于区域$\Omega$满足$(x,y,z)\in D$当且仅当$(x,-y,z)\in D$，且$\frac{y}{1+y^2}=-\frac{-y}{1+(-y)^2}$，利用函数折半可发现
        $$\iiint_\Omega\frac{y}{1+y^2}\dr x\dr y\dr z=0$$
        由此可将原积分化为
        $$\iiint_\Omega\frac{x}{1+x^2}\dr x\dr y\dr z$$

        \item \textbf{累次积分化}
            
        由于被积函数对$y$、$z$为常数，先对$y$、$z$积分，再对$x$积分，根据$x$的范围$[0,1]$将其写成
        $$\int_0^1\dr x\iint_{D_x}\frac{x}{1+x^2}\dr y\dr z$$
        $D_x$为给定$x$后$(y,z)$的区域$y^2+\frac{z^2}{3}\le2(1+x^2)$

        \item \textbf{积分计算}
            
        利用重积分的几何意义可知$\iint_{D_x}\dr y\dr z$代表$D_x$的面积，而$D_x$可以改写为
        $$\frac{y^2}{2(1+x^2)}+\frac{z^2}{6(1+x^2)}\le1$$
        为长轴长$\sqrt{6(1+x^2)}$、短轴长$\sqrt{2(1+x^2)}$的椭圆，面积即为
        $$2\sqrt3(1+x^2)\pi$$
        从而积分可化为
        $$\int_0^12\sqrt3\pi x\dr x=\sqrt3\pi$$
    \end{itemize}
    
    \note 本题反射对称化简后也可直接\textbf{广义柱坐标换元}再进行累次积分化。

    \item 分为如下步骤：
    \begin{itemize}
        \item \textbf{换元}
        
        由于被积函数与条件中$x+y+z$、$x-y$、$x-z$分别看作了整体，考虑换元
        $$u=x+y+z,\quad v=x-y,\quad w=x-z$$
        直接计算Jacobi行列式可知
        $$\frac{D(u,v,w)}{D(x,y,z)}=\begin{vmatrix}1&1&1\\1&-1&0\\1&0&-1\end{vmatrix}=3$$
        也即$\dr u\dr v\dr w=3\dr x\dr y\dr z$，从而
        $$\dr x\dr y\dr z=\frac{1}{3}\dr u\dr v\dr w$$
        且被积函数成为$u\cos(u^2)$，约束化为
        $$0\le u\le1,\quad0\le v\le1,\quad0\le w\le1$$
        将上方约束确定的$(u,v,w)$区域记为$\Omega_0$，积分即化为
        $$\frac{1}{3}\iiint_{\Omega_0}u\cos(u^2)\dr u\dr v\dr w$$

        \note 注意若计算的是$(u,v,w)$对$(x,y,z)$的Jacobi行列式，换元时需要\textbf{取倒数}。

        \item \textbf{累次积分化}
        
        根据区域描述可直接写出累次积分(由于被积函数只与$u$相关，最后对$u$积分)
        $$\frac{1}{3}\int_0^1\dr u\int_0^1\dr v\int_0^1u\cos(u^2)\dr w$$

        \item \textbf{积分计算}
        
        对$w$、$v$积分可得其为
        $$\frac{1}{3}\int_0^1\dr u\int_0^1u\cos(u^2)\dr v=\frac{1}{3}\int_0^1u\cos(u^2)\dr u$$
        换元$t=u^2$即得其为
        $$\frac{1}{6}\int_0^1\cos t\dr t=\frac{1}{6}\sin 1$$
    \end{itemize}

    \item 分为如下步骤：
    \begin{itemize}
        \item \textbf{化为第二型曲面积分}
        
        利用第一型曲面积分与第二型曲面积分的关系可知
        $$\nv\dr S=(\dr y\dr z,\dr z\dr x,\dr x\dr y)$$
        由此原积分即化为($S$以外侧为正定向)
        $$\iint_{\Sigma^+}(P\dr y\dr z+Q\dr z\dr x+R\dr x\dr y)$$

        \item \textbf{高斯公式}
        
        由初等函数知$P$、$Q$、$R$在$\rb^3$去除原点外的区域内可导且导函数连续，直接计算得$P_x+Q_y+R_z$为
        $$\frac{3}{(x^2+y^2+z^2)^{3/2}}-\frac{3x^2}{(x^2+y^2+z^2)^{5/2}}-\frac{3y^2}{(x^2+y^2+z^2)^{5/2}}-\frac{3z^2}{(x^2+y^2+z^2)^{5/2}}=0$$
        不过，此区域并非空间单连通，无法得到积分为0。我们需要利用空间多连通区域的高斯公式将积分曲面\textbf{转化为更好计算的曲面}。

        为了化简分母，考虑曲面
        $$S_0\colon x^2+y^2+z^2=1$$
        可以发现$S_0$包含在$S$中，设$S_0$与$S$之间的区域为$\Omega$，利用\textbf{空间多连通区域的高斯公式}有
        $$\begin{aligned}&\oiint_{S^+}(P\dr y\dr z+Q\dr z\dr x+R\dr x\dr y)-\oiint_{S_0^+}(P\dr y\dr z+Q\dr z\dr x+R\dr x\dr y)\\=&\iiint_\Omega(P_x+Q_y+R_z)\dr V=0\end{aligned}$$
        由此将积分转化为
        $$\oiint_{S_0^+}(P\dr y\dr z+Q\dr z\dr x+R\dr x\dr y)$$

        \item \textbf{积分化简}
        
        由于$S_0$上满足$x^2+y^2+z^2=1$，我们可以将分母改写为1以化简积分成
        $$\oiint_{S_0^+}\big((x+yz)\dr y\dr z+(y+xz)\dr z\dr x+(z+xy)\dr x\dr y\big)$$
        
        \item \textbf{高斯公式}
        
        虽然至此已经可以直接计算，我们还可以选择更简单的方式，也即\textbf{再用一次高斯公式}。将新的被积函数看作
        $$P_0\dr y\dr z+Q_0\dr z\dr x+R_0\dr x\dr y$$
        由初等函数知$P_0$、$Q_0$、$R_0$在$S_0$围成的区域$\Omega_0$内可导且导函数连续，进一步计算验证可知
        $$(P_0)_x+(Q_0)_y+(R_0)_z=3$$
        积分最终化为
        $$3\iiint_{\Omega_0}\dr x\dr y\dr z$$

        \item \textbf{积分计算}

        利用重积分的几何意义，$\iiint_{\Omega_0}\dr x\dr y\dr z$代表$\Omega_0$的体积，而此积分区域是半径为1的球，体积为$\frac{4}{3}\pi$，从而最终得到积分结果$4\pi$。
        

    \end{itemize}

    \item 分为如下步骤：
    \begin{itemize}
        \item \textbf{斯托克斯公式}
        
        将被积函数看作$P\dr x+Q\dr y+R\dr z$，由初等函数知$P$、$Q$、$R$在$\rb^3$上可导且导函数连续，直接计算得
        $$R_y-Q_z=-2y-2z,\quad P_z-R_x=-2z-2x,\quad Q_x-P_y=-2x-2y$$
        为了应用斯托克斯公式，考虑$\Gamma$在平面上围成的部分$S$，则原积分可化为
        $$-2\iint_S(y+z)\dr y\dr z+(z+x)\dr z\dr x+(x+y)\dr x\dr y$$
        利用右手定则，$S$应选择指向$z$轴正向的法向作为正向。

        \item \textbf{化为第一型曲面积分}
            
        利用第一型曲面积分与第二型曲面积分的关系可知(这里$\vec{n}$代表第二型曲面积分对应定向的单位法向量)
        $$\vec{n}\dr S=(\dr y\dr z,\dr z\dr x,\dr x\dr y)$$
        利用平面的几何性质可知它的一个法向为$(1,1,1)$，指向$z$轴正向即代表$z$分量为正，从而对应分量的单位法向量为
        $$\frac{1}{\sqrt3}(1,1,1)$$
        由此即得原积分为
        $$-2\iint_S(y+z,z+x,x+y)\cdot(\dr y\dr z,\dr z\dr x,\dr x\dr y)=-2\iint_S(y+z,z+x,x+y)\cdot\frac{1}{\sqrt3}(1,1,1)\dr S$$
        也即
        $$-\frac{4}{\sqrt3}\iint_S(x+y+z)\dr S$$

        \item \textbf{积分计算}
            
        由于平面上的点满足$x+y+z=1$，积分可以化为
        $$-\frac{4}{\sqrt3}\iint_S\dr S$$

        利用第一型曲面积分的几何意义，$\iint_S\dr S$即为曲面的面积。代入公式可知平面$x+y+z=1$与原点的距离为$\frac{1}{\sqrt3}$，从而根据几何关系得$S$是一个半径为
        $$\sqrt{1-\bigg(\frac{1}{\sqrt3}\bigg)^2}=\sqrt{\frac{2}{3}}$$
        的圆，面积为$\frac{2}{3}\pi$，最终得到结果为
        $$-\frac{8\sqrt3}{9}\pi$$
    \end{itemize}
    
    \item
    \begin{enumerate}[(1)]
        \item 若否，设
        $$f(x_0,y_0,z_0)\ne0,\quad(x_0,y_0,z_0)\in\Omega$$
        不妨设$f(x_0,y_0,z_0)=c>0$。

        利用$f$在$\Omega$上的\textbf{连续性}，存在$\delta>0$使得只要$|(x,y,z)-(x_0,y_0,z_0)|<\delta$且$(x,y,z)\in\Omega$，即有
        $$|f(x,y,z)-f(x_0,y_0,z_0)|<\frac{c}{2}$$
        另一方面，利用$\Omega$为\textbf{开集}，存在$\delta_0>0$使得只要$|(x,y,z)-(x_0,y_0,z_0)|<\delta_0$，即有$(x,y,z)\in\Omega$。

        由此，设$\gamma=\min\{\delta,\delta_0\}$，构造
        $$\Omega_1=\{(x,y,z)\in\rb^3\mid(x-x_0)^2+(y-y_0)^2+(z-z_0)^2<\gamma^2\}$$
        根据$\delta$与$\delta_0$的定义可知$\Omega_1\subset\Omega$，且其上$|f(x,y,z)-f(x_0,y_0,z_0)|<\frac{c}{2}$，从而
        $$f(x,y,z)>f(x_0,y_0,z_0)-\frac{c}{2}=\frac{c}{2}$$
        由此(等号是由于对常数重积分的几何意义为区域体积的倍数，而$\Omega_1$是半径为$\gamma$的球)
        $$\iiint_{\Omega_1}f(x,y,z)\dr x\dr y\dr z\ge\iiint_{\Omega_1}\frac{c}{2}\dr x\dr y\dr z=\frac{c}{2}\cdot\frac{4\pi}{3}\gamma^3>0$$
        矛盾。

        \item 分为如下步骤：
        \begin{itemize}
            \item \textbf{高斯公式}
            
            为了处理``对任何光滑曲面''的条件，需要应用高斯公式转化为区域。

            将被积函数看作$P\dr y\dr z+Q\dr z\dr x+R\dr x\dr y$，由条件可知$P$、$Q$、$R$在$\Omega$内可导且导函数连续，直接计算得
            $$P_x+Q_y+R_z=f(x)+xf'(x)-xf(x)-\er^{2x}$$
            由此设$S$围成的区域为$\Omega_1$，利用高斯公式可将积分化为
            $$\iiint_{\Omega_1}(f(x)+xf'(x)-xf(x)-\er^{2x})\dr V$$
            从而条件可变为
            $$\iiint_{\Omega_1}(f(x)+xf'(x)-xf(x)-\er^{2x})\dr V=0$$
            对任何有光滑简单边界的$\Omega$子区域$\Omega_1$成立。

            \item \textbf{结论应用}
            
            由于球是有光滑简单边界的区域，与第一问相同过程可证明必然有
            $$\forall x>0,\quad f(x)+xf'(x)-xf(x)-\er^{2x}=0$$

            \item \textbf{微分方程求解}
            
            由于$x>0$，我们可以将$f$满足的方程写为
            $$f'(x)+\frac{1-x}{x}f(x)=\frac{\er^{2x}}{x}$$
            这是一个一阶线性微分方程，直接计算可得(由于$x>0$可去除绝对值)
            $$\int\frac{1-x}{x}\dr x=\int\bigg(\frac{1}{x}-1\bigg)\dr x=\ln x-x+C$$
            由此可取一个积分因子
            $$\mu(x)=\er^{\ln x-x}=\frac{x}{\er^x}$$
            方程两边同乘$\mu(x)$有
            $$\frac{\dr}{\dr x}\bigg(\frac{x}{\er^x}f(x)\bigg)=\er^x$$
            对右侧不定积分有($C$为任意常数)
            $$\frac{x}{\er^x}f(x)=\er^x+C$$
            也即
            $$f(x)=\frac{\er^{2x}+C\er^x}{x}$$

            \item \textbf{常数确定}

            由条件可知下方极限存在：
            $$\lim_{x\to0^+}\frac{\er^{2x}+C\er^x}{x}$$
            若分子在0处的极限非零，则利用上学期知识可知上述极限为无穷，不存在，从而为使其存在必然
            $$\lim_{x\to0^+}(\er^{2x}+C\er^x)=0$$
            由初等函数连续性即得这是$1+C=0$，于是$C=-1$。进一步利用泰勒展开(也可使用洛必达法则)计算可发现
            $$\lim_{x\to0^+}\frac{\er^{2x}-\er^x}{x}=\lim_{x\to0^+}\frac{1+2x+o(x)-1-x-o(x)}{x}=1$$
            于是此时极限的确存在，最终得到
            $$f(x)=\frac{\er^{2x}-\er^x}{x}$$

            \note 最后的验证过程是\textbf{必须}的，因为我们仅从极限存在推出了$C=1$，但推出过程并非等价。
        \end{itemize}
    \end{enumerate}
\end{enumerate}

\subsection{下半学期建议}
\subsubsection{知识与处理技术}

\subsubsection{级数知识基础}

\subsubsection{积分进阶知识基础}

\newpage
\section{数项级数}
\subsection{习题解答}
\subsubsection{作业解答}
\begin{enumerate}
    \item (习题10.1.1(3))利用柯西收敛原理证明，设两个级数$\suminf a_n$、$\suminf b_n$均收敛，且存在正整数$N$使得
    $$\forall n>N,\quad a_n\le u_n\le b_n$$
    则级数$\suminf u_n$收敛。

    \sol{
        根据柯西收敛原理，对任何$\varepsilon>0$，存在$N_1$、$N_2$使得
        $$\forall m>n>N_1,\quad\bigg|\sum_{k=n+1}^ma_k\bigg|<\varepsilon$$
        $$\forall m>n>N_2,\quad\bigg|\sum_{k=n+1}^mb_k\bigg|<\varepsilon$$
        由此，取$N_0=\max\{N,N_1,N_2\}$，则$m>n>N_0$时有
        $$\sum_{k=n+1}^ma_k\le\sum_{k=n+1}^mu_k\le\sum_{k=n+1}^mb_k$$
        由条件左侧右侧的绝对值都小于$\varepsilon$，从而
        $$-\varepsilon<\sum_{k=n+1}^ma_k\le\sum_{k=n+1}^mu_k\le\sum_{k=n+1}^mb_k<\varepsilon$$
        于是
        $$\bigg|\sum_{k=n+1}^mu_k\bigg|<\varepsilon$$
        由于对任何$\varepsilon>0$都能找到对应的$N_0$，这已经符合柯西收敛原理要求。

        \note 注意柯西收敛原理是\textbf{充要}的，可以两方向使用。
    }


    \item (习题10.1.3)判断下列级数是否收敛：
    \begin{enumerate}
        \item[(3)] $$\suminf\frac{1}{(5n-4)(5n+1)}$$
        
        \sol{
            利用分子的形式可直接\textbf{裂项}计算得到
            $$\sum_{n=1}^m\frac{1}{(5n-4)(5n+1)}=\frac{1}{5}\sum_{n=1}^m\frac{5}{(5n-4)(5n+1)}=\frac{1}{5}\sum_{n=1}^m\bigg(\frac{1}{5n-4}-\frac{1}{5n+1}\bigg)$$
            这也就是
            $$\frac{1}{5}\bigg(1-\frac{1}{6}+\frac{1}{6}-\frac{1}{11}+\dots-\frac{1}{5m+1}\bigg)=\frac{1}{5}\bigg(1-\frac{1}{5m+1}\bigg)=\frac{m}{5m+1}$$
            利用级数定义与多项式相关极限结论(可见上学期讲义2.4.1)可知
            $$\suminf\frac{1}{(5n-4)(5n+1)}=\lim_{m\to\infty}\frac{m}{5m+1}=\frac{1}{5}$$
            从而收敛。
        }

        \item[(4)] $$\suminf\cos^2\frac{\pi}{n}$$
        
        \sol{
            由于$\frac{\pi}{n}$在$n\to\infty$时极限为0，利用\textbf{归结原理}与初等函数连续性可知
            $$\lim_{n\to\infty}\cos^2\frac{\pi}{n}=\lim_{x\to0}\cos^2x=\cos^20=1$$
            因此原级数通项极限不为0，发散。

            \note 注意换元等操作需要先用归结原理\textbf{转化为函数极限}后才能进行，可见上学期讲义3.2.2。
        }

        \item[(6)] $$\suminf\frac{\ln^n5}{3^n}$$
        
        \sol{
            设$q=\frac{\ln 5}{3}$，估算可知$0<q<1$，通过等比数列求和公式可计算出
            $$\sum_{n=1}^mq^m=\frac{q-q^{m+1}}{1-q}$$
            利用级数定义与指数函数极限性质可知原级数为
            $$\suminf q^n=\lim_{m\to\infty}\frac{q-q^{m+1}}{1-q}=\frac{q}{1-q}$$
            从而收敛。

            \note 准确来说，从$m\to\infty$时$q^m\to0$计算出上述结果还需要利用\textbf{四则运算}的极限结论。

            \note 也可直接利用教材中等比级数的敛散性结论。
        }
    \end{enumerate}

    \item (习题10.1.5)设级数$\suminf u_n$收敛，且对任何正整数$n$有$u_n\ge u_{n+1}\ge0$，证明
    $$\lim_{n\to\infty}nu_n=0$$

    \sol{
        介绍两种方法：
        \begin{itemize}
            \item \textbf{子列处理}
            
            \note 这是教材中提示的方法，想法本质上来自教材本节中对调和级数发散的\textbf{证明过程}。

            由条件利用柯西收敛原理可知(由$u_n$非负可去除绝对值)
            $$\forall\varepsilon>0,\quad \exists N,\quad\forall m>n>N,\quad\sum_{k=n+1}^mu_k<\varepsilon$$
            取$m=2n$，则此时利用$u_n$单调性有
            $$nu_{2n}\le\sum_{k=n+1}^{2n}u_k<\varepsilon$$
            又由$nu_{2n}$恒大于等于0，其绝对值必然小于$\varepsilon$，这已经符合
            $$\lim_{n\to\infty}nu_{2n}=0$$
            的\textbf{定义}。利用极限四则运算可以进一步得到
            $$\lim_{n\to\infty}2nu_{2n}=0$$
            直接根据$u_n$的单调性放缩得
            $$0\le(2n+1)u_{2n+1}\le(2n+1)u_{2n}$$
            而再次利用极限的四则运算与多项式相关极限结论可知
            $$\lim_{n\to\infty}(2n+1)u_{2n}=\lim_{n\to\infty}\frac{2n+1}{2n}\cdot2nu_{2n}=0$$
            于是利用\textbf{夹逼定理}最终得到
            $$\lim_{n\to\infty}(2n+1)u_{2n+1}=0$$
            由于$2nu_{2n}$与$(2n+1)u_{2n+1}$是$nu_n$的奇数项与偶数项构成的子列，且它们收敛于同一极限0，利用上学期结论(或见本讲义27.1.2第3题)可知$nu_n$也收敛于0，得证。

            \note 级数相关的问题时常会用到\textbf{子列极限}的相关结论，详见本讲义27.2.1。
            
            \item \textbf{直接放缩}
            
            由于此形式与调和级数接近，我们尝试直接将其向调和级数放缩。

            采用\textbf{反证法}。若要证的极限不存在或存在且非零，利用极限定义，存在$\varepsilon>0$使得
            $$\forall N,\quad\exists n>N,\quad |nu_n|\ge\varepsilon$$
            结合$u_n$非负也即
            $$u_n\ge\frac{\varepsilon}{n}$$
            直接估算是相对困难的，为了说明矛盾，我们仍然需要考虑通过\textbf{柯西收敛原理}证明发散。柯西收敛原理的否定即
            $$\exists\varepsilon_0>0,\quad\forall N_0,\quad\exists m_0>n_0>N_0,\quad\bigg|\sum_{k=n_0+1}^{m_0}u_k\bigg|\ge\varepsilon_0$$
            由于已知$u_k$单调减且大于0，只需证明
            $$\exists\varepsilon_0>0,\quad\forall N_0,\quad\exists m_0>n_0>N_0,\quad(m_0-n_0)u_{m_0}>\varepsilon_0$$
            即有$|\sum_{k=n_0+1}^{m_0}u_k|\ge(m_0-n_0)u_{m_0}\ge\varepsilon_0$成立。

            为了确定这里的$\varepsilon_0$，并在$N_0$给定后确定$n_0$与$m_0$，我们先进行分析：首先，取定$n_0=N_0+1$，这是由于$u_{m_0}>0$，$n_0$越小可以使得$(m_0-n_0)u_{m_0}$越大；其次，为了能使$u_{m_0}$能够被控制，我们必须取$m_0$为条件中的$n$，这样才能得到估算
            $$(m_0-n_0)u_{m_0}>\frac{m_0-N_0-1}{m_0}\varepsilon=\bigg(1-\frac{N_0+1}{m_0}\bigg)\varepsilon$$
            可以看出，只要$m_0$充分大就可以使右侧永远大于$\frac{\varepsilon}{2}$，这样就得到了取法。

            取定$\varepsilon_0=\frac{\varepsilon}{2}$，对任何$N_0$，取$n_0=N_0+1$，在条件中令$N=2N_0+2$，则存在$n>N$使得
            $$u_n>\frac{\varepsilon}{n}$$
            将此$n$记为$m_0$，则最终由$m_0>2N_0+2$有
            $$\bigg|\sum_{k=n_0+1}^{m_0}u_k\bigg|\ge(m_0-n_0)u_{m_0}>\bigg(1-\frac{N_0+1}{m_0}\bigg)\varepsilon>\frac{\varepsilon}{2}=\varepsilon_0$$
            由于对任何$N_0$都存在上述的$m_0$、$n_0$，利用柯西收敛原理可知原级数发散，矛盾。

            \note 可以看出，本题的放缩做法需要\textbf{相对复杂的估算技巧}，但直接通过柯西收敛原理估算的思想是更加简单的。
        \end{itemize}
    }

    \item (总练习题10.1)对数列$a_n$、$b_n$，若存在常数$C>0$使得$n$充分大时有$|a_n|\le C|b_n|$，则记为
    $$a_n=O(|b_n|)\quad(n\to\infty)$$
    这是\textbf{阶估算}的一个常用记号。证明如下结论：
    \begin{enumerate}
        \item[(3)] 对任意实数$k$有
        $$n^{1+\frac{k}{\ln n}}=O(n)\quad(n\to\infty)$$

        \sol{
            利用对数的换底公式$\frac{\log_ab}{\log_ac}=\log_cb$可计算得
            $$n^{1+\frac{k}{\ln n}}=nn^{\frac{k}{\ln n}}=nn^{\log_n(\er^k)}=\er^kn$$
            从而取$C=\er^k>0$即对任何正整数$n$有
            $$n^{1+\frac{k}{\ln n}}=\er^kn\le Cn$$
            符合定义要求。
        }
        
        \item[(4)] $$\frac{n^3(\sqrt2+(-1)^n)^n}{3^n}=O\bigg(\bigg(\frac{\sqrt2+1}{3}\bigg)^nn^3\bigg)\quad(n\to\infty)$$
        
        \sol{
            由$0<\sqrt2-1<\sqrt2+1$可知
            $$(\sqrt2-1)^n<(\sqrt2+1)^n$$
            而$(\sqrt2+(-1)^n)^n$为$(\sqrt2-1)^n$或$(\sqrt2+1)^n$，从而它不超过$(\sqrt2+1)^n$，由此即得对任何正整数$n$有
            $$\frac{n^3(\sqrt2+(-1)^n)^n}{3^n}\le\bigg(\frac{\sqrt2+1}{3}\bigg)^nn^3$$
            取定义中的常数$C=1$可发现上式已经符合定义。
        }
    \end{enumerate}

    \item (习题10.2.1)判断下列级数是否收敛：
    \begin{enumerate}
        \item[(1)] $$\suminf2^n\sin\frac{\pi}{4^n}$$
        
        \sol{
            \note 利用等价无穷小结论可以感受到$\sin\frac{\pi}{4^n}$与$\frac{\pi}{4^n}$接近，由此可想到构造。

            考虑级数$\suminf\frac{1}{2^n}$，利用\textbf{等比级数}敛散性结论可知其收敛，且由于$\lim_{n\to\infty}\frac{1}{4^n}=0$，利用\textbf{推广的归结原理}与等价无穷小替换可得
            $$\lim_{n\to\infty}\frac{2^n\sin\frac{\pi}{4^n}}{\frac{1}{2^n}}=\lim_{n\to\infty}4^n\sin\frac{\pi}{4^n}=\lim_{x\to0}\frac{\sin(\pi x)}{x}=\pi$$
            从而根据\textbf{比较判别法}得原级数也收敛。
        }

        \item[(3)] $$\suminf\frac{1}{\sqrt[n]{n}}$$
        
        \sol{
            利用\textbf{函数极限阶估算}结论(可见上学期讲义3.1第11题)有
            $$\lim_{x\to+\infty}\frac{\ln x}{x}=0$$
            从而通过推广的\textbf{归结原理}与\textbf{复合函数极限定理}可算出
            $$\lim_{n\to\infty}\frac{1}{\sqrt[n]{n}}=\lim_{x\to+\infty}\frac{1}{\sqrt[x]{x}}=\lim_{x\to+\infty}\frac{1}{\er^{\frac{\ln x}{x}}}=\frac{1}{\er^0}=1$$
            因此原级数通项极限不为0，发散。
        }

        \item[(4)] $$\suminf\frac{4n}{n^2+4n-3}$$
        
        \sol{
            \note 由于分母为二阶、分子为一阶，此级数通项应当为一阶无穷小，从而与$\frac{1}{n}$比较。

            考虑级数$\suminf\frac{1}{n}$，利用\textbf{$p$级数}敛散性结论可知其发散，通过多项式相关极限结论(可见上学期讲义2.4.1)可证明
            $$\lim_{n\to\infty}\frac{\frac{4n}{n^2+4n-3}}{\frac{1}{n}}=\lim_{n\to\infty}\frac{4n^2}{n^2+4n-3}=4$$
            从而根据\textbf{比较判别法}得原级数也发散。
        }
    
        \item[(7)] $$\suminf n\tan\frac{1}{3^n}$$
        
        \sol{
            与第一问类似，利用\textbf{推广的归结原理}与等价无穷小替换可得
            $$\lim_{n\to\infty}\frac{n\tan\frac{1}{3^n}}{\frac{n}{3^n}}=\lim_{n\to\infty}3^n\tan\frac{1}{3^n}=\lim_{x\to0}\frac{\tan x}{x}=1$$
            由此可以将原级数的敛散性问题转化为研究如下级数是否收敛：
            $$\suminf\frac{n}{3^n}$$
            
            \note 由于此级数比$\suminf\frac{1}{3^n}$更大，可以采用\textbf{放缩后比较}的策略将分母适当缩小。
            
            考虑级数$\suminf\frac{1}{2^n}$，利用\textbf{等比级数}敛散性结论可知其收敛，利用\textbf{数列极限阶估算}结论(可见上学期讲义2.4.1)得
            $$\lim_{n\to\infty}\frac{\frac{n}{3^n}}{\frac{1}{2^n}}=\lim_{n\to\infty}\frac{n}{(3/2)^n}=0$$
            从而根据\textbf{比较判别法}得原级数也收敛。

            \note 转化问题后也可使用\textbf{比值判别法}从$\lim_{n\to\infty}\frac{u_{n+1}}{u_n}=\frac{1}{3}$得到收敛。

        }
    \end{enumerate}

    \item (习题10.2.2)判断下列级数是否收敛：
    \begin{enumerate}
        \item[(2)] $$\suminf\frac{n!}{3n^2}$$
        
        \sol{
            利用\textbf{数列极限阶估算}结论(可见上学期讲义2.4.1)得
            $$\lim_{n\to\infty}\frac{n^2}{n!}=0$$
            从而利用无穷极限性质可知
            $$\lim_{n\to\infty}\frac{n!}{3n^2}=\infty$$
            因此原级数通项极限不为0，发散。
        }

        \item[(4)] $$\suminf\frac{1}{n^{1+1/n}}$$
        
        \sol{
            考虑级数$\suminf\frac{1}{n}$，利用\textbf{$p$级数}敛散性结论可知其发散，由本部分第5题(3)的过程可计算出
            $$\lim_{n\to\infty}\frac{\frac{1}{n^{1+1/n}}}{\frac{1}{n}}=\lim_{n\to\infty}\frac{1}{\sqrt[n]{n}}=1$$
            从而根据\textbf{比较判别法}得原级数也发散。
        }

        \item[(6)] $$\suminf[n=2]\frac{n}{\ln^nn}$$
        
        \sol{
            \note 若分母是公比大于1的等比数列，与本部分第5题(7)相同可证明收敛，而此分母实际上更大，由此想到放缩后比较。

            当$n>100$时直接估算可知$\ln n>3$，从而此时有
            $$\frac{n}{\ln^nn}<\frac{n}{3^n}$$
            由本部分第5题(7)的过程可知$\suminf\frac{n}{3^n}$收敛，从而根据\textbf{比较判别法}得原级数也收敛。
            
            \note 也可使用\textbf{柯西判别法}从$\lim_{n\to\infty}\sqrt[n]{u_n}=0$得到收敛。
        }

        \item[(8)] $$\suminf\frac{(n!)^2}{(2n)!}$$
        
        \sol{
            为了进行判断，我们需要先估算原式。直接消去可知
            $$\frac{(n!)^2}{(2n)!}=\frac{1\cdot2\cdots n}{(n+1)\cdot(n+2)\cdots(2n)}$$
            由于上方下方的差恒为$n$，利用$a>b>0$、$c>0$时$\frac{a}{b}<\frac{a+c}{b+c}$\ (也称为\textbf{糖水不等式})将每个乘积放大到$\frac{n}{2n}$可知
            $$\frac{(n!)^2}{(2n)!}=\frac{1}{n+1}\cdot\frac{2}{n+2}\cdots\frac{n}{2n}\le\frac{n}{2n}\cdot\frac{n}{2n}\cdots\frac{n}{2n}=\frac{1}{2^n}$$
            利用\textbf{等比级数}敛散性结论可知放大后的新级数收敛，从而根据\textbf{比较判别法}得原级数也收敛。

            \note 也可使用\textbf{比值判别法}从$\lim_{n\to\infty}\frac{u_{n+1}}{u_n}=\lim_{n\to\infty}\frac{(n+1)^2}{(2n+1)(2n+2)}=\frac{1}{4}$得到收敛。
        }

        \item[(11)] $$\suminf[n=3]\frac{1}{n(\ln\ln n)^q}$$
        这里参数$q>0$。

        \sol{
            当$n\ge3$时$\ln\ln n>0$，从而此级数为正项级数。分为如下步骤：
            \begin{itemize}
                \item \textbf{判断$\suminf[n=2]\frac{1}{n\ln^pn}$的敛散性}
                
                由定义可以发现$\suminf[n=2]\frac{1}{n\ln^pn}$是通项单调下降且恒正，从而设$f(x)=\frac{1}{x\ln^px}$，利用\textbf{积分判别法}得原级数敛散性等价于极限
                $$\lim_{x\to+\infty}\int_2^xf(t)\dr t$$
                的存在性。

                直接换元$s=\ln t$计算可知
                $$\int_2^x\frac{1}{t\ln^pt}\dr t=\int_2^x\frac{\dr\ln t}{\ln^pt}=\int_{\ln 2}^{\ln x}\frac{1}{s^p}\dr s=\begin{cases}\frac{1}{1-p}(\ln^{1-p}2-\ln^{1-p}x)&p\ne1\\\ln\ln x-\ln\ln 2&p=1\end{cases}$$
                由于$x\to+\infty$时$\ln x\to+\infty$，换元可知$\ln\ln x\to+\infty$，从而利用推广的复合函数极限结论可计算得
                $$\lim_{x\to+\infty}\int_2^x\frac{1}{t\ln^pt}\dr t=\begin{cases}+\infty&0<p\le1\\\frac{1}{1-p}\ln^{1-p}2&p>1\end{cases}$$

                也即级数$\suminf[n=2]\frac{1}{n\ln^pn}$在$p>1$时收敛，在$p\in(0,1]$时发散。

                \note 利用数列极限的阶估算结论可知$n\ln^qn$是比$n$高阶的无穷大，但比$c>1$时的$n^c$低阶，因此此级数是对$p$级数\textbf{边界情况的更精细讨论}。难以用$p$级数比较时可考虑使用此级数。
                
                \item \textbf{原题证明}
                
                利用阶估算可以感受到对任何$c>0$，$\frac{1}{n(\ln\ln n)^q}$都比$\frac{1}{n\ln^cn}$更大，而后者对充分小的$c$发散，由此前者发散，下面对此想法进行严谨化。

                由于$n\to\infty$时$\ln n\to+\infty$，通过推广的\textbf{归结原理}与\textbf{复合函数极限定理}可算出
                $$\lim_{n\to\infty}\frac{\frac{1}{n\ln n}}{\frac{1}{n(\ln\ln n)^q}}=\lim_{n\to\infty}\frac{(\ln\ln n)^q}{\ln n}=\lim_{x\to+\infty}\frac{\ln^qx}{x}=0$$
                这里最后一步使用了\textbf{函数极限阶估算}结论(可见上学期讲义3.1第11题)。

                前一部分已经证明$\suminf[n=2]\frac{1}{n\ln n}$发散，从而根据\textbf{比较判别法}得原级数也发散。
            \end{itemize}
        }

        \item[(12)] $$\suminf[n=3]\frac{1}{n\ln^pn(\ln\ln n)^q}$$
        这里参数$p>0$、$q>0$。

        \sol{
            当$n\ge3$时$\ln\ln n>0$，从而此级数为正项级数。当$p\ne1$时，可以用与(11)相同的方法得到结论，否则需要更精细的讨论，由此我们分为三种情况研究：
            \begin{itemize}
                \item $p>1$\textbf{时}
                
                $n>100$时直接估算可知$\ln\ln n>1$，从而此时有
                $$\frac{1}{n\ln^pn(\ln\ln n)^q}<\frac{1}{n\ln^pn}$$
                我们在(11)中已经证明$\suminf[n=2]\frac{1}{n\ln^pn}$收敛，从而根据\textbf{比较判别法}得原级数也收敛。

                \item $p<1$\textbf{时}
                
                与(11)完全相同可计算得到
                $$\lim_{n\to\infty}\frac{\frac{1}{n\ln n}}{\frac{1}{n\ln^pn(\ln\ln n)^q}}=\lim_{n\to\infty}\frac{(\ln\ln n)^q}{\ln^{1-p}n}=\lim_{x\to+\infty}\frac{\ln^qx}{x^{1-p}}=0$$

                我们在(11)中已经证明$\suminf[n=2]\frac{1}{n\ln^pn}$发散，从而根据\textbf{比较判别法}得原级数也发散。

                \item $p=1$\textbf{时}
                
                由定义可以发现$\suminf[n=3]\frac{1}{n\ln n(\ln\ln n)^q}$通项单调下降且恒正，从而设$f(x)=\frac{1}{x\ln x(\ln\ln x)^q}$，利用\textbf{积分判别法}得原级数的敛散性等价于极限
                $$\lim_{x\to+\infty}\int_3^xf(t)\dr t$$
                的存在性。

                直接换元$s=\ln t$计算可知
                $$\int_3^x\frac{1}{t\ln t(\ln\ln t)^q}\dr t=\int_3^x\frac{\dr\ln t}{\ln t(\ln\ln t)^q}=\int_{\ln 3}^{\ln x}\frac{1}{s\ln^qs}\dr s$$
                
                \note 注意此形式与用积分判别法判断$\suminf[n=2]\frac{1}{n\ln^pn}$敛散性得到的积分形式\textbf{完全相同}，只是积分范围从2到$x$变为了$\ln 3$到$\ln x$。
                
                我们在(11)中已经证明右侧积分在$\ln x\to+\infty$时极限存在当且仅当$q>1$，由于$x\to+\infty$时$\ln x\to+\infty$，利用\textbf{复合函数极限}结论可知此积分在$x\to+\infty$时极限存在也当且仅当$q>1$。由此，我们最终得到$q>1$时原级数收敛，$q\in(0,1]$时原级数发散。
            \end{itemize}
            
            \note 与(11)类似，此级数是对$\suminf[n=2]\frac{1}{n\ln^pn}$边界情况的更精细讨论，使用相同的方法可以讨论更精细的级数$\suminf[n=2]\frac{1}{n\ln n\ln\ln n(\ln\ln\ln n)^p}$等。
        }
    \end{enumerate}

    \item (习题10.2.4)证明若级数$\suminf a_n^2$、$\suminf b_n^2$收敛，则以下三个级数都收敛：
    $$\suminf|a_nb_n|,\quad\suminf(a_n+b_n)^2,\quad\suminf\frac{|a_n|}{n}$$

    \sol{
        根据定义可发现三个级数均为正项级数，我们分别证明三个级数的敛散性：
        \begin{itemize}
            \item 利用收敛级数的性质可知
            $$\suminf(a_n^2+b_n^2)$$
            收敛，而对任何正整数$n$，利用基本不等式有
            $$|a_nb_n|\le\frac{1}{2}(a_n^2+b_n^2)$$
            从而根据\textbf{比较判别法}得原级数也收敛。

            \item 对任何正整数$n$，利用基本不等式有
            $$(a_n^2+b_n^2)\le2(a_n^2+b_n^2)$$
            从而与第一种情况相同根据\textbf{比较判别法}得原级数也收敛。

            \item 取定$b_n=\frac{1}{n}$，利用$p$级数敛散性结论可知
            $$\suminf b_n^2=\suminf\frac{1}{n^2}$$
            收敛，此时$|a_nb_n|=\frac{|a_n|}{n}$，由第一种情况得原级数收敛。
        \end{itemize}
    }

    \item (习题10.2.6)设存在实数$l>0$使得
    $$\lim_{n\to\infty}nu_n=l$$
    证明级数$\suminf u_n^2$收敛，$\suminf u_n$发散。
    
    \sol{
        利用极限的\textbf{保序性}，由$nu_n$极限大于0可知对充分大的$n$有$nu_n>0$，从而对充分大的$n$有$u_n>0$，忽略有限项后可以看作正项级数。

        \note 注意使用比较判别法等判别法时时需要先\textbf{确认其为正项级数}。

        利用四则运算的极限结论可知
        $$\lim_{n\to\infty}\frac{u_n}{\frac{1}{n}}=\lim_{n\to\infty}nu_n=l>0$$
        $$\lim_{n\to\infty}\frac{u_n^2}{\frac{1}{n^2}}=\lim_{n\to\infty}(nu_n)\cdot(nu_n)=l^2>0$$
        利用\textbf{$p$级数}敛散性结论可知$\suminf\frac{1}{n}$发散、$\suminf\frac{1}{n^2}$收敛，从而根据\textbf{比较判别法}得$\suminf u_n$发散、$\suminf u_n^2$收敛。
    }

    \item (总练习题10.2)将比较判别法写为阶估算的形式(记号$O$含义见本部分第4题)：若
    $$a_n=O(|b_n|)\quad(n\to\infty)$$
    且$\suminf|b_n|$收敛，则$\suminf|a_n|$收敛。以此判断下列级数是否收敛：
    \begin{enumerate}
        \item[(1)] $$\suminf[n=2](\sqrt{n+1}-\sqrt{n})^p\ln\frac{n-1}{n+1}$$
        这里参数$p>0$。

        \sol{
            我们利用上学期学习的技巧分别处理乘积的两部分：对第一部分进行\textbf{分子有理化}得到
            $$\sqrt{n+1}-\sqrt{n}=\frac{1}{\sqrt{n+1}+\sqrt{n}}<\frac{1}{2\sqrt n}$$
            对第二部分，利用\textbf{归结原理}与\textbf{等价无穷小替换可知}
            $$\lim_{n\to\infty}n\ln\frac{n-1}{n+1}=\lim_{n\to\infty}n\ln\bigg(1-\frac{2}{n+1}\bigg)=\lim_{x\to+\infty}x\ln\bigg(1-\frac{2}{x+1}\bigg)=\lim_{x\to+\infty}-\frac{2x}{x+1}=-2$$
            从而利用极限\textbf{保序性}可知$n$充分大时
            $$n\ln\frac{n-1}{n+1}\in(-1,-3)$$
            取绝对值可发现此时
            $$\bigg|\ln\frac{n-1}{n+1}\bigg|<\frac{3}{n}$$
            综合以上两部分，我们最终得到了$n$充分大时
            $$\bigg|(\sqrt{n+1}-\sqrt{n})^p\ln\frac{n-1}{n+1}\bigg|\le\frac{1}{(2\sqrt{n})^p}\cdot\frac{3}{n}=\frac{3}{2^p}\frac{1}{n^{1+p/2}}$$
            于是
            $$\bigg|(\sqrt{n+1}-\sqrt{n})^p\ln\frac{n-1}{n+1}\bigg|=O\bigg(\frac{1}{n^{1+p/2}}\bigg)\quad(n\to\infty)$$
            利用\textbf{$p$级数}敛散性结论与\textbf{比较判别法}得$\suminf|(\sqrt{n+1}-\sqrt{n})^p\ln\frac{n-1}{n+1}|$收敛，而直接计算可发现
            $$\suminf\bigg|(\sqrt{n+1}-\sqrt{n})^p\ln\frac{n-1}{n+1}\bigg|=\suminf\bigg(-(\sqrt{n+1}-\sqrt{n})^p\ln\frac{n-1}{n+1}\bigg)$$
            从而根据收敛级数的性质知原级数收敛。
        }

        \item[(4)] $$\suminf\frac{n^3(\sqrt2+(-1)^n)^n}{3^n}$$
        
        \sol{
            在本部分第4题中已经证明
            $$\frac{n^3(\sqrt2+(-1)^n)^n}{3^n}=O\bigg(\bigg(\frac{\sqrt2+1}{3}\bigg)^nn^3\bigg)\quad(n\to\infty)$$
            且左侧为正项级数，从而为证明左侧级数收敛只需证明右侧级数收敛，下设$c=\frac{\sqrt2+1}{3}<1$。

            由于比等比级数多出了$n^3$，需要进行\textbf{放缩估算}。任取$s\in(c,1)$，考虑级数$\suminf s^n$，利用\textbf{等比级数}敛散性结论可知其收敛，利用\textbf{数列极限阶估算}结论(可见上学期讲义2.4.1)得
            $$\lim_{n\to\infty}\frac{n^3c^n}{s^n}=\lim_{n\to\infty}\frac{n^3}{(s/c)^n}=0$$
            从而根据\textbf{比较判别法}得原级数也收敛。

            \note 转化问题后也可使用\textbf{比值判别法}从$\lim_{n\to\infty}\frac{u_{n+1}}{u_n}=c$得到收敛。
        }
    \end{enumerate}

    \item (总练习题10.3)对正数列$a_n$，设存在$\alpha>0$使得
    $$\exists N,\quad\forall n\ge N,\quad\ln\frac{1}{a_n}\ge(1+\alpha)\ln n$$
    证明$\suminf a_n$收敛。

    \sol{
        利用对数性质有$(1+\alpha)\ln n=\ln n^{1+\alpha}$。由指数函数单调性，两边同作$\er$的指数可将不等式转化为
        $$\frac{1}{a_n}\ge n^{1+\alpha},\quad a_n\le\frac{1}{n^{1+\alpha}}$$
        利用\textbf{$p$级数}敛散性结论可知$\suminf\frac{1}{n^{1+\alpha}}$收敛，从而根据\textbf{比较判别法}得原级数也收敛。

    }
\end{enumerate}

\subsubsection{补充题解答}
\begin{enumerate}
    \item (习题10.1.1)利用柯西收敛原理判断下列级数是否收敛：
    \begin{enumerate}[(1)]
        \item $$\suminf\frac{\cos n}{n(n+1)}$$
        
        \sol{
            由于需要应用柯西收敛原理，我们先对任何$m>n$估算部分和
            $$\bigg|\sum_{k=n+1}^m\frac{\cos k}{k(k+1)}\bigg|$$
            直接放缩裂项有
            $$\bigg|\sum_{k=n+1}^m\frac{\cos k}{k(k+1)}\bigg|\le\sum_{k=n+1}^m\frac{|\cos k|}{k(k+1)}\le\sum_{k=n+1}^m\frac{1}{k(k+1)}=\sum_{k=n+1}^m\bigg(\frac{1}{k}-\frac{1}{k+1}\bigg)=\frac{1}{n+1}-\frac{1}{m+1}$$
            由此，对任何$\varepsilon>0$，取正整数$N>\frac{1}{\varepsilon}$，即有$m>n>N$时
            $$\bigg|\sum_{k=n+1}^m\frac{\cos k}{k(k+1)}\bigg|\le\frac{1}{n+1}-\frac{1}{m+1}<\frac{1}{n+1}<\varepsilon$$
            利用柯西收敛原理可知其收敛。
        }
    
        \item $$\suminf\frac{1}{\sqrt n}$$
        
        \sol{
            由于需要应用柯西收敛原理，我们先对任何$m>n$估算部分和(由于$\sqrt n>0$恒成立可去除绝对值)
            $$\sum_{k=n+1}^m\frac{1}{\sqrt k}$$
            将其全部缩小至最后一项可发现
            $$\sum_{k=n+1}^m\frac{1}{\sqrt k}>\sum_{k=n+1}^m\frac{1}{\sqrt m}=\frac{m-n}{\sqrt m}$$
            从而$m=2n$时它将超过$\frac{1}{\sqrt2}\sqrt n$。

            由此，取$\varepsilon=1$，对任何正整数$N$，考虑$n=N+1$、$m=2N+2$，则$m>n>N$，且
            $$\sum_{k=n+1}^m\frac{1}{\sqrt k}>\frac{m-n}{\sqrt m}=\frac{1}{\sqrt2}\sqrt{N+1}>\varepsilon$$
            这与柯西收敛原理矛盾，从而其发散。
        }
    \end{enumerate}

    \item (习题10.1.3)判断下列级数是否收敛：
    \begin{enumerate}
        \item[(2)] $$\suminf\frac{1}{(2n-1)(2n+1)}$$
        
        \sol{
            利用分子的形式可直接\textbf{裂项}计算得到
            $$\sum_{n=1}^m\frac{1}{(2n-1)(2n+1)}=\frac{1}{2}\sum_{n=1}^m\frac{2}{(2n-1)(2n+1)}=\frac{1}{2}\sum_{n=1}^m\bigg(\frac{1}{2n-1}-\frac{1}{2n+1}\bigg)$$
            这也就是
            $$\frac{1}{2}\bigg(1-\frac{1}{3}+\frac{1}{3}-\frac{1}{5}+\dots-\frac{1}{2m+1}\bigg)=\frac{1}{2}\bigg(1-\frac{1}{2m+1}\bigg)=\frac{m}{2m+1}$$
            利用级数定义与多项式相关极限结论(可见上学期讲义2.4.1)可知
            $$\suminf\frac{1}{(2n-1)(2n+1)}=\lim_{m\to\infty}\frac{m}{2m+1}=\frac{1}{2}$$
            从而收敛。
        }

        \item[(5)] $$\suminf\frac{n}{2n-1}$$
        
        \sol{
            利用多项式相关极限结论(可见上学期讲义2.4.1)可知
            $$\lim_{n\to\infty}\frac{n}{2n-1}=\frac{1}{2}$$
            因此原级数通项极限不为0，发散。
        }

        \item[(7)] $$\suminf\sqrt[n]{0.0001}$$
        
        \sol{
            利用\textbf{归结原理}与指数函数\textbf{连续性}可知
            $$\lim_{n\to\infty}\sqrt[n]{0.0001}=\lim_{n\to\infty}0.0001^{\frac{1}{n}}=\lim_{x\to0}0.0001^x=0.0001^0=1$$
            因此原级数通项极限不为0，发散。
        }
    \end{enumerate}

    \item (习题10.1.4)设级数$\suminf u_n$的部分和数列为$S_n$。若存在实数$A$使得
    $$\lim_{n\to\infty}S_{2n}=\lim_{n\to\infty}S_{2n+1}=A$$
    证明$\suminf u_n$收敛。

    \sol{
        我们下面证明
        $$\lim_{n\to\infty}S_n=A$$
        这就已经符合级数收敛定义。

        利用极限定义可知，对任何$\varepsilon>0$，存在正整数$N_1$、$N_2$使得
        $$\forall n>N_1,\quad|S_{2n}-A|<\varepsilon$$
        $$\forall n>N_2,\quad|S_{2n+1}-A|<\varepsilon$$
        由此，取$N=\max\{2N_1,2N_2+1\}$，对任何$n>N$，它或是一个大于$2N_1$的偶数，或是一个大于$2N_2+1$的奇数，因此有
        $$|S_n-A|<\varepsilon$$
        由于对任何$\varepsilon>0$都能取出符合要求的$N$，这已经符合极限定义。
    }

    \item (总练习题10.1)证明如下结论(记号$O$含义见本讲义27.1.1第4题)：
    \begin{enumerate}[(1)]
        \item $$(\sqrt{n+1}-\sqrt n)\ln\frac{n-1}{n+1}=O(n^{-3/2})\quad(n\to\infty)$$
        
        \sol{
            对第一部分进行\textbf{分子有理化}得到
            $$\sqrt{n+1}-\sqrt{n}=\frac{1}{\sqrt{n+1}+\sqrt{n}}<\frac{1}{2\sqrt n}$$
            对第二部分，利用\textbf{归结原理}与\textbf{等价无穷小替换可知}
            $$\lim_{n\to\infty}n\ln\frac{n-1}{n+1}=\lim_{n\to\infty}n\ln\bigg(1-\frac{2}{n+1}\bigg)=\lim_{x\to+\infty}x\ln\bigg(1-\frac{2}{x+1}\bigg)=\lim_{x\to+\infty}-\frac{2x}{x+1}=-2$$
            从而利用极限\textbf{保序性}可知$n$充分大时
            $$n\ln\frac{n-1}{n+1}\in(-1,-3)$$
            取绝对值可发现此时
            $$\bigg|\ln\frac{n-1}{n+1}\bigg|<\frac{3}{n}$$
            综合以上两部分，我们最终得到了$n$充分大时
            $$\bigg|(\sqrt{n+1}-\sqrt{n})\ln\frac{n-1}{n+1}\bigg|\le\frac{1}{2\sqrt{n}}\cdot\frac{3}{n}=\frac{3}{2}\frac{1}{n^{3/2}}$$
            取定义中的常数$C=\frac{3}{2}$可发现上式已经符合定义。
        }
        
        \item $$\ln\sec\frac{\pi}{n}=O\bigg(\frac{1}{n^2}\bigg)\quad(n\to\infty)$$
        
        \sol{
            利用\textbf{归结原理}可直接计算极限
            $$\lim_{n\to\infty}n^2\ln\sec\frac{\pi}{n}=\lim_{x\to+\infty}x^2\ln\frac{1}{\cos\frac{\pi}{x}}=\lim_{x\to+\infty}x^2\ln\bigg(1+\frac{1}{\cos\frac{\pi}{x}}-1\bigg)$$
            利用推广的复合函数极限结论可知$x\to+\infty$时$\cos\frac{\pi}{x}\to1$，从而$\frac{1}{\cos\frac{\pi}{x}}-1\to0$，进行\textbf{等价无穷小替换}得到极限为
            $$\lim_{x\to+\infty}x^2\bigg(\frac{1}{\cos\frac{\pi}{x}}-1\big)=\lim_{x\to\infty}\frac{x^2\big(1-\cos\frac{\pi}{x}\big)}{\cos\frac{\pi}{x}}$$
            利用乘除极限可知将分母替换为1不影响极限情况，再次利用等价无穷小替换可得其为
            $$\lim_{x\to+\infty}x^2\bigg(1-\cos\frac{\pi}{x}\bigg)=\lim_{x\to+\infty}x^2\cdot\frac{1}{2}\bigg(\frac{\pi}{x}\bigg)^2=\frac{\pi^2}{2}$$
            利用极限\textbf{保序性}可知$n$充分大时
            $$n^2\ln\sec\frac{\pi}{n}\in(0,100)$$
            也即此时
            $$\bigg|\ln\sec\frac{\pi}{n}\bigg|<\frac{100}{n^2}$$
            取定义中的常数$C=100$可发现上式已经符合定义。
        }
    \end{enumerate}
\end{enumerate}

\subsection{数列与级数}
\subsubsection{基本性质}
有限重排、丢弃有限项

级数收敛则数列极限为0

加括号对应子列

两种加括号等价的情况(正项、化交错)

\subsubsection{单调有界}

正项级数收敛判别

积分判别-连续化-归结

比较判别-放缩-夹逼

柯西强于比值

\subsubsection{阶估算 V}
多项式的情况

利用阶估算确定敛散性

\subsection{一般级数}
\subsubsection{条件收敛}
用正项控制

拆分单调

条件收敛的含义

黎曼重排

\subsubsection{柯西收敛}
数列的柯西收敛

级数版本

\subsubsection{级数乘积}
莱布尼茨判别法

阿贝尔、迪利克雷-分部-复合、分部

\newpage
\section{一致收敛}
\subsection{习题解答}
\subsubsection{作业解答}
\begin{enumerate}
    \item (习题10.3.1)判断下列级数是条件收敛、绝对收敛还是发散：
    \begin{enumerate}
        \item[(2)] $$\suminf\frac{(-1)^{n+1}}{(2n-1)^p}$$
        
        \sol{
            分为如下步骤：
            \begin{itemize}
                \item \textbf{绝对收敛判定}
                
                其通项取绝对值得到的级数为$\suminf\frac{1}{(2n-1)^p}$，考虑级数$\suminf\frac{1}{n^p}$，利用\textbf{$p$级数}敛散性结论可知其收敛当且仅当$p>1$，通过推广的\textbf{归结原理}与\textbf{复合函数极限}结论可证明
                $$\lim_{n\to\infty}\frac{\frac{1}{(2n-1)^p}}{\frac{1}{n^p}}=\lim_{n\to\infty}\frac{n^p}{(2n-1)^p}=\lim_{x\to+\infty}\bigg(\frac{x}{2x-1}\bigg)^p=\frac{1}{2^p}$$
                从而根据\textbf{比较判别法}得原级数绝对收敛当且仅当$p>1$。
                
                \item \textbf{条件收敛判定}
                
                由于$p>1$时已经绝对收敛，只考虑$p\le1$的情况。利用前一部分类似的方法可计算出
                $$\lim_{n\to\infty}\frac{1}{(2n-1)^p}=\begin{cases}+\infty&p<0\\1&p=0\\0&p>0\end{cases}$$
                在$p\le0$时，通项绝对值的极限不为0，于是通项极限不为0，发散。

                在$p>0$时，通项绝对值的极限为0，进一步由原级数为交错级数与$\frac{1}{(2n-1)^p}$对$n$单调递减(可以直接从指数函数单调性看出)，利用\textbf{莱布尼茨判别法}得此时必然收敛，从而原级数条件收敛当且仅当$p\in(0,1]$。
            \end{itemize}
        }

        \item[(4)] $$\suminf(-1)^n\frac{\sqrt{n}-1}{n}$$
        
        \sol{
            分为如下步骤：
            \begin{itemize}
                \item \textbf{绝对收敛判定}
                
                其通项取绝对值得到的级数为$\suminf\frac{\sqrt{n}-1}{n}$，通过阶估算可发现它的量级应为$n^{-1/2}$，因此考虑级数$\suminf\frac{1}{\sqrt{n}}$，利用\textbf{$p$级数}敛散性结论可知其发散，直接计算得到
                $$\lim_{n\to\infty}\frac{\frac{\sqrt n-1}{n}}{\frac{1}{\sqrt n}}=\lim_{n\to\infty}\frac{\sqrt{n}-1}{\sqrt{n}}=\lim_{n\to\infty}\bigg(1-\frac{1}{\sqrt n}\bigg)=1$$
                从而根据\textbf{比较判别法}得原级数不绝对收敛。
                
                \item \textbf{条件收敛判定}
                
                由于此为交错级数，可直接尝试采用\textbf{莱布尼茨判别法}。

                先验证通项绝对值趋于0：通过前一部分的计算结果，利用乘除极限结论即得
                $$\lim_{n\to\infty}\frac{\sqrt{n}-1}{n}=\lim_{n\to\infty}\frac{\frac{\sqrt n-1}{n}}{\frac{1}{\sqrt n}}\cdot\frac{1}{\sqrt{n}}=1\cdot0=0$$
                再验证通项绝对值单调递减：直接计算有
                $$\frac{\sqrt{n}-1}{n}-\frac{\sqrt{n+1}-1}{n+1}=\frac{(n+1)\sqrt{n}-n\sqrt{n+1}-1}{n(n+1)}$$
                对此分子进行有理化得到其为
                $$\sqrt{n(n+1)}(\sqrt{n+1}-\sqrt{n})-1=\frac{\sqrt{n(n+1)}}{\sqrt{n}+\sqrt{n+1}}-1>\frac{\sqrt{n(n+1)}}{2\sqrt{n+1}}-1=\frac{1}{2}\sqrt{n}-1$$
                由此$n>4$时分子大于0，也即此时通项绝对值单调递减。

                综合以上，利用莱布尼茨判别法可知原级数条件收敛。

                \note 由于级数改变有限项不影响敛散性，所有判别法只需在$n$\textbf{充分大}时验证成立即可。
            \end{itemize}
        }

        \item[(6)]  $$\suminf(-1)^{n+1}\frac{n!}{3^{n^2}}$$
        
        \sol{
            分为如下步骤：
            \begin{itemize}
                \item \textbf{绝对收敛判定}
                
                其通项取绝对值得到的级数为$\suminf\frac{n!}{3^{n^2}}$，尝试计算或感受可发现它的下降速度超过了等比数列，由此考虑进行\textbf{放缩}。

                通过归纳可发现$3^n>2n$对正整数$n$恒成立，由此有
                $$\frac{n!}{3^{n^2}}=\frac{1\cdot2\cdots n}{3^n\cdot 3^n\cdots 3^n}\le\frac{1\cdot2\cdots n}{3^1\cdot 3^2\cdots 3^n}<\frac{1\cdot2\cdots n}{2\cdot 4\cdots(2n)}=\frac{1}{2^n}$$

                利用\textbf{等比级数}敛散性结论可知放大后的新级数收敛，从而根据\textbf{比较判别法}得原级数绝对收敛。

                \item \textbf{条件收敛判定}
                
                由于此级数已经绝对收敛，无需进一步判定条件收敛。
            \end{itemize}
        }

        \item[(8)] $$\suminf(-1)^{n+1}\tan\frac{\varphi}{n}$$
        这里参数$\varphi\in(-\frac{\pi}{2},\frac{\pi}{2})$。
        
        \sol{
            分为如下步骤：
            \begin{itemize}
                \item \textbf{绝对收敛判定}
                
                由$\varphi$的范围，其通项取绝对值得到的级数为$\suminf\tan\frac{|\varphi|}{n}$，由等价无穷小可以想到与$\frac{1}{n}$比较。

                当$\varphi=0$时，其通项恒为0，绝对收敛，否则考虑级数$\suminf\frac{1}{n}$，利用\textbf{$p$级数}敛散性结论可知其发散，利用\textbf{归结原理}与基本等价无穷小得到
                $$\lim_{n\to\infty}\frac{\tan\frac{|\varphi|}{n}}{\frac{1}{n}}=\lim_{x\to0}\frac{\tan(|\varphi|x)}{x}=|\varphi|>0$$
                从而根据\textbf{比较判别法}得原级数也发散。
                
                综合以上，原级数绝对收敛当且仅当$\varphi=0$。
                
                \item \textbf{条件收敛判定}
                
                由于$\varphi=0$时已绝对收敛，只考虑$\varphi\ne0$的情况。

                若$\varphi>0$，原级数为交错级数，尝试使用\textbf{莱布尼茨判别法}。
                
                先验证通项绝对值趋于0：通过前一部分的计算结果，利用乘除极限结论即得
                $$\lim_{n\to\infty}\tan\frac{|\varphi|}{n}=\lim_{n\to\infty}\frac{\tan\frac{|\varphi|}{n}}{\frac{1}{n}}\cdot\frac{1}{n}=|\varphi|\cdot0=0$$
                此外，由$\varphi\in(0,\frac{\pi}{2})$，根据$\tan$的单调性可知通项绝对值单调递减，从而利用莱布尼茨判别法可知原级数收敛。

                若$\varphi<0$，可发现原级数能够改写为$\suminf(-1)^{n+2}\tan\frac{|\varphi|}{n}$，这仍然是一个交错级数，于是由$|\varphi|\in(0,\frac{\pi}{2})$仍然可以用莱布尼茨判别法得到原级数收敛。

                综合以上，原级数条件收敛当且仅当$\varphi\in(-\frac{\pi}{2},0)\cup(0,\frac{\pi}{2})$。
            \end{itemize}
        }

        \item[(10)] $$\suminf\sin(\sqrt{n^2+1}\pi)$$
        
        \sol{
            分为如下步骤：
            \begin{itemize}
                \item \textbf{分析与变形}
                
                可以发现，我们通常对数列的阶估算是利用\textbf{归结原理}将数列转化为了函数进行处理，但本题中，由于$\sin(\sqrt{x^2+1}\pi)$会不断在$-1$到1之间振荡，直接转化为函数是没有意义的。

                为了能够进行有意义的估算，我们需要观察数列本身的特殊性。可以发现$\sqrt{n^2+1}\pi$事实上与$n\pi$接近，而$\sin(n\pi)=0$，由此可将通项转化为
                $$\sin((\sqrt{n^2+1}-n)\pi+n\pi)=(-1)^n\sin((\sqrt{n^2+1}-n)\pi)$$
                进一步进行分子有理化，此即为
                $$(-1)^n\sin\frac{\pi}{\sqrt{n^2+1}+n}$$
                这样已经可以看出分母趋于0，转化为函数后仍然有意义，从而可以进行进一步估算。经过以上分析，我们最终将级数转化为了
                $$\suminf(-1)^n\sin\frac{\pi}{\sqrt{n^2+1}+n}$$
                由于$n\ge1$时$\sqrt{n^2+1}+n>2$，此形式的确是一个\textbf{交错级数}。

                \note 上学期讲义9.2.1介绍了类似的操作。

                \item \textbf{绝对收敛判定}
                
                由$n\ge1$时$\sqrt{n^2+1}+n>2$，其通项取绝对值得到的级数为
                $$\suminf\sin\frac{\pi}{\sqrt{n^2+1}+n}$$
                由于此形式与$\sin\frac{\pi}{2n}$接近，可以想到与$\frac{1}{n}$比较。考虑级数$\suminf\frac{1}{n}$，利用\textbf{$p$级数}敛散性结论可知其发散，利用\textbf{归结原理}与\textbf{等价无穷小替换}得到
                $$\lim_{n\to\infty}\frac{\sin\frac{\pi}{\sqrt{n^2+1}+n}}{\frac{1}{n}}=\lim_{x\to+\infty}x\sin\frac{\pi}{\sqrt{x^2+1}+x}=\lim_{x\to+\infty}\frac{\pi x}{\sqrt{x^2+1}+x}$$
                分子分母同除以$x$后由初等函数\textbf{连续性}与复合函数极限结论即可得到其为
                $$\lim_{x\to+\infty}\frac{\pi}{\sqrt{1+\frac{1}{x^2}}+1}=\frac{\pi}{2}$$
                从而根据\textbf{比较判别法}得原级数不绝对收敛。

                \item \textbf{条件收敛判定}

                由于此为交错级数，可直接尝试采用\textbf{莱布尼茨判别法}。

                先验证通项绝对值趋于0：通过前一部分的计算结果，利用乘除极限结论即得
                $$\lim_{n\to\infty}\sin\frac{\pi}{\sqrt{n^2+1}+n}=\lim_{n\to\infty}\frac{\sin\frac{\pi}{\sqrt{n^2+1}+n}}{\frac{1}{n}}\cdot\frac{1}{n}=\frac{\pi}{2}\cdot0=0$$
                再验证通项绝对值单调递减：由于$n$增加时$\sqrt{n^2+1}+n$增加，且恒大于2，根据三角函数单调性可知$\sin\frac{\pi}{\sqrt{n^2+1}+n}$单调递减。

                综合以上，利用莱布尼茨判别法可知原级数条件收敛。

                \note 由于级数改变有限项不影响敛散性，所有判别法只需在$n$\textbf{充分大}时验证成立即可。
            \end{itemize}
        }
    \end{enumerate}

    \item (习题10.3.2)已知级数$\suminf u_n$收敛，证明级数
    $$\suminf\frac{u_n}{n^p},\quad\suminf\frac{n}{n+1}u_n$$
    均收敛，这里参数$p>0$。

    \sol{
        将两级数改写为
        $$\suminf\frac{1}{n^p}\cdot u_n,\quad\suminf\frac{n}{n+1}\cdot u_n$$
        由于$\frac{1}{n^p}$随$n$增加单调减且在$(0,1)$中，$\frac{n}{n+1}=1-\frac{1}{n}$随$n$增加单调增且在$(0,1)$中，而$\suminf u_n$收敛，利用\textbf{阿贝尔判别法}得两者均收敛。
    }

    \item (习题10.3.3)证明级数
    $$\suminf\frac{\cos(n\varphi)}{n^p}$$
    在$\varphi\in(0,2\pi)$、$p>1$时绝对收敛，在$\varphi\in(0,2\pi)$、$p\in(0,1]$时条件收敛。

    \sol{
        我们将介绍两种方法说明$p\in(0,1]$时不绝对收敛。分为如下步骤：
        \begin{itemize}
            \item \textbf{$p>1$绝对收敛判定}
            
            其通项取绝对值得到的级数为$\suminf\frac{|\cos(n\varphi)|}{n^p}$，由于
            $$\frac{|\cos(n\varphi)|}{n^p}\le\frac{1}{n^p}$$
            在$p>1$时，利用\textbf{$p$级数}敛散性结论可知$\suminf\frac{1}{n^p}$收敛，从而根据\textbf{比较判别法}得原级数绝对收敛。
            
            \item \textbf{$p\in(0,1]$绝对收敛判定-拆分级数}
            
            当$\varphi=\pi$时，原级数的通项取绝对值得到的级数即为$\suminf\frac{1}{n^p}$，在$p\in(0,1]$时利用\textbf{$p$级数}敛散性结论可知发散。下面只考虑$\varphi\in(0,\pi)\cup(\pi,2\pi)$的情况。
            
            我们进行如下的放缩：
            $$\frac{|\cos(n\varphi)|}{n^p}\ge\frac{\cos^2(n\varphi)}{n^p}=\frac{1}{2n^p}+\frac{\cos(2n\varphi)}{2n^p}$$
            在$p\in(0,1]$时利用\textbf{$p$级数}敛散性结论可知$\suminf\frac{1}{n^p}$发散，从而$\suminf\frac{1}{2n^p}$发散。另一方面，由$\varphi\in(0,\pi)\cup(\pi,2\pi)$可知$\sin\varphi\ne0$，从而利用\textbf{等差数列的正弦求和公式}(可见上学期讲义7.2.1)得
            $$\sum_{n=1}^m\cos(2n\varphi)=\sum_{n=1}^m\sin\bigg(2n\varphi+\frac{\pi}{2}\bigg)=\frac{\sin((n+1)\varphi)\cos(n\varphi)}{\sin\varphi}$$
            由此可进一步得到
            $$\bigg|\sum_{n=1}^m\cos(2n\varphi)\bigg|\le\frac{1}{|\sin\varphi|}$$
            将级数$\suminf\frac{\cos(2n\varphi)}{2n^p}$写为
            $$\suminf\frac{1}{2n^p}\cdot\cos(2n\varphi)$$
            之前的推导已证明$\cos(2n\varphi)$的部分和有界，而$\frac{1}{2n^p}$随$n$增加单调递减趋于0，利用\textbf{迪利克雷判别法}得其收敛。

            综合以上，利用\textbf{级数的性质}可知发散级数与收敛级数之和发散，因此$\suminf\big(\frac{1}{2n^p}+\frac{\cos(2n\varphi)}{2n^p}\big)$发散，再由正项级数的\textbf{比较判别法}最终得到$\suminf\frac{|\cos(n\varphi)|}{n^p}$发散，原级数不绝对收敛。

            \note 建议记忆三角函数等差数列的正弦求和公式，它常用于一些三角函数相关的级数估算。

            \note 此方法相对简洁，但\textbf{难以适用一般情况}，如将分子的$\cos(n\varphi)$改为$\cos^9(n\varphi)$后即难以处理。下一种方法虽然书写相对复杂，但应用范围实际上更广。
            
            \item \textbf{$p\in(0,1]$绝对收敛判定-直接估算}
                
            与上一种情况同理，我们只考虑$\varphi\in(0,\pi)\cup(\pi,2\pi)$的情况。进一步地，由于$|\cos(n\varphi)|=|\cos(n(\varphi-\pi))|$，当$\varphi\in(\pi,2\pi)$时可将其减去$\pi$，因此事实上只需考虑$\varphi\in(0,\pi)$的情况。

            我们采用一个可以称为\textbf{猎人兔子原理}的思路进行估算(可见上学期讲义2.3.1)：给定$\delta\in(0,\frac{\pi}{2})$，我们将所有如下区间
            $$\bigg(k\pi+\frac{\pi}{2}-\delta,k\pi+\frac{\pi}{2}+\delta\bigg),\quad k\in\mathbb{Z}$$
            的并集记为集合$A$，利用三角函数的性质可得$x\in A$当且仅当$|\cos x|<\sin\delta$。

            可以发现，集合$A$中的两个点若属于同一个区间，则距离一定小于$2\delta$，而若属于不同区间，距离一定大于$\pi-2\delta$。由此，我们取$\delta$使得$2\delta<\varphi<\pi-2\delta$，也即
            $$0<\delta<\min\bigg\{\frac{\pi-\varphi}{2},\frac{\varphi}{2}\bigg\}$$
            由$\varphi$的范围，这样的$\delta$一定存在。此时，由于$n\varphi$与$(n+1)\varphi$的差恰好为$\varphi$，根据上方的讨论，它们不可能同时落在集合$A$中，于是必有一个的余弦绝对值大于等于$\sin\delta$。综合以上的讨论，我们找到了一个$\delta\in(0,\frac{\pi}{2})$使得对任何正整数$n$都有
            $$\max\{|\cos(n\varphi)|,|\cos((n+1)\varphi)|\}\ge\sin\delta$$

            接下来的过程就相对直接了：直接将$\frac{|\cos(n\varphi)|}{n^p}$两两分组，得到(无论$|\cos((2n-1)\varphi)|\ge\sin\delta$还是$|\cos(2n\varphi)|\ge\sin\delta$都能如此放缩)
            $$\sum_{n=1}^{2M}\frac{|\cos(n\varphi)|}{n^p}=\sum_{n=1}^M\bigg(\frac{|\cos((2n-1)\varphi)|}{(2n-1)^p}+\frac{|\cos(2n\varphi)|}{(2n)^p}\bigg)\ge\sum_{n=1}^M\frac{\sin\delta}{(2n)^p}=\frac{\sin\delta}{2^p}\sum_{n=1}^M\frac{1}{n^p}$$
            利用\textbf{$p$级数}敛散性结论可知
            $$\lim_{M\to\infty}\frac{\sin\delta}{2^p}\sum_{n=1}^M\frac{1}{n^p}=+\infty$$
            从而根据\textbf{推广的夹逼定理}
            $$\lim_{M\to\infty}\sum_{n=1}^{2M}\frac{|\cos(n\varphi)|}{n^p}=+\infty$$
            而这是$\suminf\frac{|\cos(n\varphi)|}{n^p}$的一个\textbf{子列极限}，由子列极限不存在即可知原数列极限不存在，最终得到原级数不绝对收敛。

            \item \textbf{条件收敛判定}    
            
            综合之前的讨论，我们已经得到原级数绝对收敛当且仅当$p>1$，下面研究$p\in(0,1]$时的敛散性。

            由$\varphi\in(0,2\pi)$，利用\textbf{等差数列的正弦求和公式}(可见上学期讲义7.2.1)得
            $$\sum_{n=1}^m\cos(n\varphi)=\sum_{n=1}^m\sin\bigg(n\varphi+\frac{\pi}{2}\bigg)=\frac{\sin(\frac{n+1}{2}\varphi)\cos\frac{n\varphi}{2}}{\sin\frac{\varphi}{2}}$$
            由此可进一步得到
            $$\bigg|\sum_{n=1}^m\cos(n\varphi)\bigg|\le\frac{1}{|\sin\frac{\varphi}{2}|}$$
            将级数$\suminf\frac{\cos(n\varphi)}{n^p}$写为
            $$\suminf\frac{1}{n^p}\cdot\cos(n\varphi)$$
            之前的推导已证明$\cos(n\varphi)$的部分和有界，而$\frac{1}{n^p}$随$n$增加单调递减趋于0，利用\textbf{迪利克雷判别法}得其收敛。

            由此，原级数条件收敛当且仅当$p\in(0,1]$。
        \end{itemize}
    }

    \item (习题10.3.7)证明收敛的级数
    $$\suminf\frac{(-1)^{n-1}}{\sqrt{n}}$$
    重排后的级数
    $$1+\frac{1}{\sqrt3}-\frac{1}{\sqrt2}+\cdots+\frac{1}{\sqrt{4k-3}}+\frac{1}{\sqrt{4k-1}}-\frac{1}{\sqrt{2k}}+\cdot$$
    发散。

    \sol{
        原级数直接利用莱布尼茨判别法可验证收敛，我们现在证明重排后的级数发散。

        将重排后的第$k$项记为$v_k$，并记
        $$u_k=v_{3k-2}+v_{3k-1}+v_{3k}=\frac{1}{\sqrt{4k-3}}+\frac{1}{\sqrt{4k-1}}-\frac{1}{\sqrt{2k}}$$
        首先，直接放缩可发现
        $$u_k>\frac{1}{\sqrt{4k}}+\frac{1}{\sqrt{4k}}-\frac{1}{\sqrt{2k}}=\bigg(1-\frac{1}{\sqrt2}\bigg)\frac{1}{\sqrt{k}}>0$$
        从而这是一个\textbf{正项级数}。利用\textbf{$p$级数}敛散性结论可知$\suminf[k=0]\frac{1}{\sqrt{k}}$发散，从而$\suminf[k=0](1-\frac{1}{\sqrt2})\frac{1}{\sqrt{k}}$发散，由\textbf{比较判别法}得$\suminf[k=0]u_k$也发散。

        由定义可知$\suminf[k=0]u_k$是$\suminf[k=0]v_k$的一个\textbf{子列极限}，由子列极限不存在即可知原数列极限不存在，这就证明了结论。

        \note 注意不要忘记最后一步。$\suminf[k=0]u_k$相当于$\suminf[k=0]v_k$每三项加括号，由此是一个子列极限。
    }
\end{enumerate}
\subsubsection{补充题解答}
\begin{enumerate}
    \item (习题10.2.1)判断下列级数是否收敛：
    \begin{enumerate}
        \item[(2)] $$\suminf\frac{1}{\sqrt{2n^3+1}}$$
        
        \sol{
            通过阶估算可发现它的量级应为$n^{-3/2}$，因此考虑级数$\suminf\frac{1}{n^{3/2}}$，利用\textbf{$p$级数}敛散性结论可知其收敛，利用\textbf{归结原理}与初等函数\textbf{连续性}可知
            $$\lim_{n\to\infty}\frac{\frac{1}{\sqrt{2n^3+1}}}{\frac{1}{n^{3/2}}}=\lim_{n\to\infty}\sqrt{\frac{n^3}{2n^3+1}}=\lim_{n\to\infty}\sqrt{\frac{1}{2+\frac{1}{n^3}}}=\lim_{x\to0}\sqrt{\frac{1}{2+x}}=\frac{1}{\sqrt2}$$
            从而根据\textbf{比较判别法}得原级数也收敛。
        }

        \item[(6)] $$\suminf[n=2]\frac{n}{(\ln n)^{\ln n}}$$
        
        \sol{
            当$n\ge2$时$\ln n>0$，从而此级数为正项级数。先尝试化简分母，取对数得到
            $$(\ln n)^{\ln n}=\er^{\ln n\cdot\ln(\ln n)}=(\er^{\ln n})^{\ln\ln n}=n^{\ln\ln n}$$
            估算可发现$n^{\ln\ln n}$比$n$的多项式阶数更高，从而可以考虑采用\textbf{放缩}的思路：利用\textbf{归结原理}与推广的\textbf{复合函数极限定理}可知(换元$t=\ln x$)
            $$\lim_{n\to\infty}\ln\ln n=\lim_{x\to+\infty}\ln\ln x=\lim_{t\to+\infty}\ln t=+\infty$$
            由此利用极限\textbf{保序性}，存在正整数$N$使得$n>N$时$\ln\ln n>3$，此时即有
            $$\frac{n}{(\ln n)^{\ln n}}=\frac{n}{n^{\ln\ln n}}<\frac{1}{n^2}$$
            利用\textbf{$p$级数}敛散性结论可知$\suminf\frac{1}{n^2}$收敛，从而根据\textbf{比较判别法}得原级数也收敛。
        }
    \end{enumerate}

    \item (习题10.2.2)判断下列级数是否收敛：
    \begin{enumerate}
        \item[(1)] $$\suminf\frac{n^5}{n!}$$
        
        \sol{
            由于$n!$比$n$的多项式高阶，考虑级数$\suminf\frac{1}{n^2}$，利用\textbf{$p$级数}敛散性结论可知其收敛，利用\textbf{数列极限阶估算}结论(可见上学期讲义2.4.1)得
            $$\lim_{n\to\infty}\frac{\frac{n^5}{n!}}{\frac{1}{n^2}}=\lim_{n\to\infty}\frac{n^7}{n!}=0$$
            从而根据\textbf{比较判别法}得原级数也收敛。
        }

        \item[(5)] $$\suminf\frac{n^2}{\big(3-\frac{1}{n}\big)^n}$$
        
        \sol{
            由于分母恒大于$2^n$，其又比$n$的多项式高阶，考虑级数$\suminf\frac{1}{n^2}$，利用\textbf{$p$级数}敛散性结论可知其收敛，而由$n$为正整数时$3-\frac{1}{n}>2$可知
            $$0\le\frac{n^4}{\big(3-\frac{1}{n}\big)^n}\le\frac{n^4}{2^n}$$
            右侧利用\textbf{数列极限阶估算}结论(可见上学期讲义2.4.1)可知极限为0，根据\textbf{夹逼定理}知中间极限为0，而这就代表
            $$\lim_{n\to\infty}\frac{\frac{n^2}{\big(3-\frac{1}{n}\big)^n}}{\frac{1}{n^2}}=\lim_{n\to\infty}\frac{n^4}{\big(3-\frac{1}{n}\big)^n}=0$$
            从而根据\textbf{比较判别法}得原级数也收敛。
        }
        
        \item[(10)] $$\suminf[n=2]\frac{1}{n\ln^pn}$$
        这里参数$p>0$。
        
        \sol{
            见本讲义27.1.1第6题(11)第一部分证明。
        }
    \end{enumerate}

    \item (习题10.2.3)证明若正项级数$\suminf u_n$收敛，则$\suminf u_n^2$一定收敛，反之不然。
    
    \sol{
        分两部分说明：
        \begin{itemize}
            \item \textbf{原级数收敛则平方收敛}
            
            由$\suminf u_n$收敛可知
            $$\lim_{n\to\infty}u_n=0$$
            从而利用极限\textbf{保序性}，存在正整数$N$使得$n>N$时
            $$u_n<1$$
            又由$u_n$为正可知此时
            $$u_n^2<u_n$$
            从而根据\textbf{比较判别法}得$\suminf u_n^2$也收敛。
            
            \item \textbf{平方收敛原级数未必收敛}
            
            考虑$u_n=\frac{1}{n}$，利用\textbf{$p$级数}敛散性结论可知$\suminf u_n$收敛、$\suminf u_n^2$发散，从而得证。
        \end{itemize}
    }

    \item (习题10.2.5)若正项级数$\suminf u_n$、$\suminf v_n$均发散，判断下列级数是必然发散、必然收敛还是可能发散可能收敛：
    \begin{enumerate}[(1)]
        \item $$\suminf(u_n+v_n)$$
        
        \sol{
            由正项级数定义可知$u_n+v_n\ge u_n$，从而根据\textbf{比较判别法}得到其必然发散。
        }

        \item $$\suminf(u_n-v_n)$$
        
        \sol{
            由级数定义可直接算出若级数通项为常数，当且仅当通项\textbf{恒为0}时级数收敛。

            当$u_n=v_n=1$时$u_n-v_n=0$，从而$\suminf(u_n-v_n)$收敛；当$u_n=2$、$v_n=1$时$u_n-v_n=1$，从而$\suminf(u_n-v_n)$发散。

            综合以上，原级数可能发散可能收敛。
        }

        \item $$\suminf u_nv_n$$
        
        \sol{
            利用\textbf{$p$级数}敛散性结论，取$u_n=v_n=\frac{1}{\sqrt{n}}$符合条件，此时$u_nv_n=\frac{1}{n}$，再次利用$p$级数敛散性结论可知$\suminf u_nv_n$发散；利用\textbf{$p$级数}敛散性结论，取$u_n=v_n=\frac{1}{n}$符合条件，此时$u_nv_n=\frac{1}{n^2}$，再次利用$p$级数敛散性结论可知$\suminf u_nv_n$收敛。

            综合以上，原级数可能发散可能收敛。
        }
    \end{enumerate}
\end{enumerate}

\subsection{和函数}
\subsubsection{函数项级数}
函数逼近

级数逼近优势

性质刻画

一般函数项级数

\subsubsection{和函数的性质}
连续性

可导性

可积性

\subsubsection{一致收敛}
从连续性到一致

一致收敛下的可导与可积

\subsection{收敛判定}
\subsubsection{基本判别}
比较与数项级数

条件与绝对

阿贝尔、迪利克雷

柯西准则

\subsubsection{一致收敛的否定}
找点列

利用一致收敛性质

\subsection{补充习题}
\subsubsection{数项级数}
\begin{framed}
\begin{exercise}
    设$a_n$为正数列，其部分和数列为$S_n$，且$S_n$极限为$A$\ (可能为有限数或正无穷)，对$\alpha>0$，判断
    $$\suminf\frac{a_n}{S_n^\alpha}$$
    的敛散性。
\end{exercise}
\sol{

}
\end{framed}

\begin{framed}
\begin{exercise}
    设正项级数$\suminf a_n$发散，证明存在发散的正向级数$\suminf b_n$使得
    $$\lim_{n\to\infty}\frac{b_n}{a_n}=0$$

    \note 这说明没有\textbf{发散最慢}的级数。
\end{exercise}
\sol{

}
\end{framed}

\begin{framed}
\begin{exercise}
    设$a_n$为单调递减的正数列，证明
    $$\suminf a_n,\quad\suminf[n=0]2^na_{2^n}$$
    同敛散。
\end{exercise}
\sol{

}
\end{framed}

\begin{framed}
\begin{exercise}
    设$a_n=\frac{1}{n}$，$b_n$为$a_n$中去除所有分母含有数字``9''的项后构成的数列，证明$\suminf b_n$收敛。
\end{exercise}
\sol{

}
\end{framed}

\begin{framed}
\begin{exercise}
    判断
    $$\suminf\frac{(-1)^{[\sqrt{n}]}}{n}$$
    的敛散性。
\end{exercise}
\sol{

}
\end{framed}

\begin{framed}
\begin{exercise}
    证明存在数列$a_n$使得$\suminf a_n$收敛但$\suminf a_n^3$发散。
\end{exercise}
\sol{

}
\end{framed}

\begin{framed}
\begin{exercise}
    设级数$\suminf a_n$收敛，证明
    $$\lim_{n\to\infty}\frac{a_1+2a_2+\dots+na_n}{n}=0$$
\end{exercise}
\sol{

}
\end{framed}

\begin{framed}
\begin{exercise}
    若级数$\suminf a_n$绝对收敛，且
    $$\lim_{n\to\infty}b_n=0$$
    证明
    $$\lim_{n\to\infty}(a_1b_n+a_2b_{n-1}+\dots+a_nb_1)=0$$
\end{exercise}
\sol{

}
\end{framed}

\subsubsection{函数项级数}
\begin{framed}
\begin{exercise}
    对函数列
    $$f_n(x)=xn^{-x}\ln^\alpha n$$
    证明它在$[0,+\infty)$一致收敛当且仅当参数$\alpha<1$。
\end{exercise}
\sol{

}
\end{framed}

\begin{framed}
\begin{exercise}
    设$f_1(x)$为$[a,b]$上的连续函数，证明函数列
    $$f_n(x)=\int_a^xf_{n-1}(t)\dr t$$
    在$[a,b]$一致收敛于0。
\end{exercise}
\sol{

}
\end{framed}

\begin{framed}
\begin{exercise}
    证明关于$x$的函数项级数
    $$\suminf[n=2]\frac{\cos(nx)}{n\ln x}$$
    在$(0,\pi]$中不一致收敛。
\end{exercise}
\sol{

}
\end{framed}

\begin{framed}
\begin{exercise}
    证明函数项级数
    $$\suminf u_n(x),\quad u_n(x)=\begin{cases}\frac{1}{n}&x=\frac{1}{n}\\0&x\ne\frac{1}{n}\end{cases}$$
    在$[0,1]$上一致收敛，但无法使用强级数判别法判别。
\end{exercise}
\sol{

}
\end{framed}

\begin{framed}
\begin{exercise}
    设
    $$f(x)=\suminf\frac{x^n}{(1+2x)^n}\cos\frac{n\pi}{x}$$
    计算$f(x)$在1处与无穷处的极限。
\end{exercise}
\sol{

}
\end{framed}

\begin{framed}
\begin{exercise}[\extra]
    在$(0,1)$中任取一个两两不同的数列$a_n$，设
    $$f(x)=\suminf\frac{|x-a_n|}{2^n}$$
    证明它在$[0,1]$连续，并判断$f(x)$在哪些点可导，计算导函数。
\end{exercise}
\sol{

}
\end{framed}

\begin{framed}
\begin{exercise}[\extra]
    微分方程解存在唯一性    
\end{exercise}
\end{framed}

\newpage
\section{特殊函数项级数}
\subsection{习题解答}
\subsubsection{作业解答}
\begin{enumerate}
    \item (习题10.4.1)计算下列关于$x$的函数项级数的收敛域：
    \begin{enumerate}
        \item[(1)] $$\suminf(\ln x)^n$$
        
        \sol{

        }

        \item[(3)] $$\suminf\frac{1}{x^n}\sin\frac{\pi}{3^n}$$
    \end{enumerate}

    \item (习题10.4.2)判断下列关于$x$的函数列在给定区间是否一致收敛；
    \begin{enumerate}
        \item[(2)] $$f_n(x)=\sqrt{x^4+\er^{-n}},\quad x\in\rb$$
        
        \sol{

        }

        \item[(4)] $$f_n(x)=\frac{n^2x}{1+n^2x},\quad x\in(0,1)$$
        
        \sol{

        }

        \item[(6)] $$f_n(x)=\frac{\sin\frac{x}{n}}{\frac{x}{n}},\quad x\in(0,1)$$
        
        \sol{

        }
    \end{enumerate}

    \item (习题10.4.3)判断下列关于$x$的函数项级数在给定区间是否一致收敛；
    \begin{enumerate}
        \item[(2)] $$\suminf\bigg(\frac{x^n}{n}-\frac{x^{n+1}}{n+1}\bigg),\quad x\in[-1,1]$$
        
        \sol{

        }

        \item[(4)] $$\suminf\frac{x}{1+4n^4x^2},\quad x\in\rb$$
        
        \sol{

        }

        \item[(7)] $$\suminf(-1)^{n-1}x^2\er^{-nx^2},\quad x\in\rb$$
        
        \sol{
            直接计算可知
            $$$$
        }
    \end{enumerate}

    \item (习题10.4.5)证明函数项级数
    $$g(x)=\suminf 2^n\sin\frac{x}{3^n}$$
    对任何正数$M$在$[-M,M]$上一致收敛，在$\rb$上不一致收敛，且有连续的导函数。

    \sol{

    }

    \item (习题10.4.7)证明函数项级数
    $$f(x)=\suminf\frac{\sin(nx)}{n^{2+p}}$$
    在$\rb$上有连续的导函数，这里参数$p>0$。

    \sol{

    }

    \item (习题10.4.8)设$\suminf a_n$收敛，证明函数项级数
    $$f(x)=\suminf\frac{a_n}{n^x}$$
    在$[0,+\infty)$一致收敛，且
    $$\lim_{x\to0^+}f(x)=\suminf a_n$$

    \sol{

    }

    \item (总练习题10.7)对复数列$a_n$，考虑它的部分和序列
    $$S_n=\sum_{k=1}^na_k$$
    若复数$\alpha$满足$S_n$的实部收敛于$\alpha$的实部，$S_n$的复部收敛于$\alpha$的复部，则称级数$\suminf a_n$收敛于$\alpha$。证明复数项级数的柯西收敛原理(\extra)，并以此证明对任意复数$z$，复数项级数
    $$\suminf\frac{z^n}{n!}$$
    收敛。

    \note 教材原题并未要求证明柯西收敛原理，这里将证明部分补充完整。

    \sol{

    }

    \item (习题10.5.1)求下列关于$x$的幂级数的收敛半径：
    \begin{enumerate}
        \item[(2)] $$\suminf\frac{n^k}{n!}x^n$$
        这里参数$k$为正整数。

        \sol{

        }
        
        \item[(4)] $$\suminf\frac{(n!)^2}{(2n)!}x^n$$
        
        \sol{

        }
    \end{enumerate}

    \item (习题10.5.2)求下列关于$x$的幂级数的收敛区间与收敛域：
    \begin{enumerate}
        \item[(2)] $$\suminf\frac{x^n}{na^n}$$
        这里参数$a>0$。

        \sol{

        }

        \item[(4)] $$\suminf\frac{n}{n+1}\bigg(\frac{x}{2}\bigg)^n$$
        
        \sol{

        }

        \item[(8)] $$\suminf(3^n+5^n)x^n$$
        
        \sol{

        }

        \item[(11)] $$\suminf\frac{n!}{n^n}(x-2)^n$$
        
        \sol{

        }
    \end{enumerate}

    \item (习题10.5.3)求下列关于$x$的幂级数的和函数：
    \begin{enumerate}
        \item[(1)] $$\suminf[n=0](n+1)x^n$$
        
        \sol{

        }

        \item[(3)] $$\suminf(-1)^{n-1}\frac{x^{n+1}}{n(n+1)}$$
        
        \sol{

        }

        \item[(5)] $$\suminf[n=0]\frac{2n+1}{n!}x^{2n}$$
        
        \sol{

        }

        \item[(7)] $$\suminf\frac{x^n}{n(n+1)}$$
        
        \sol{

        }
    \end{enumerate}
\end{enumerate}

\subsubsection{补充题解答}
\begin{enumerate}
    \item (习题10.3.1)判断下列级数是条件收敛、绝对收敛还是发散：
    \begin{enumerate}
        \item[(9)] $$\suminf\frac{(-1)^{n+1}}{n^t\ln^sn}$$
        这里参数$t>0$、$s>0$。

        \sol{

        }

        \item[(11)] $$\suminf(-1)^n\frac{1\cdot3\cdot5\cdots(2n-1)}{2\cdot4\cdot6\cdots(2n)}$$
        
        \sol{

        }
    \end{enumerate}

    \item (习题10.3.4)给定参数$\varphi\in(0,2\pi)$、$p>0$，判断以下列级数是条件收敛、绝对收敛还是发散：
    $$\suminf\frac{\cos(n\varphi)}{n^p}\bigg(1+\frac{1}{n}\bigg)^n$$

    \sol{

    }

    \item (习题10.3.5)给定数列$a_n$，形如
    $$\suminf\frac{a_n}{n^x}$$
    的关于$x$的函数项级数称为Dirichlet级数。

    证明若Dirichlet级数在$x_0$收敛，则对任何$x>x_0$也收敛；若在$x_0$发散，则对任何$x<x_0$也发散。

    \sol{

    }

    \item (习题10.3.6)设级数$\suminf u_n$绝对收敛，证明级数
    $$\suminf\frac{2n-1}{n}u_n$$
    绝对收敛。

    \sol{

    }

    \item (总练习题10.2(2))用阶估算形式的比较判别法(定义见本讲义27.1.1第9题)判断以下级数是否收敛：
    $$\suminf\ln^p\sec\frac{\pi}{n}$$
    这里参数$p>0$。

    \sol{

    }


    \item (习题10.4.1(2))计算以下关于$x$的函数项级数的收敛域：
    $$\suminf\frac{1}{2n-1}\bigg(\frac{1-x}{1+x}\bigg)^n$$

    \sol{

    }

    \item (习题10.4.2)判断下列关于$x$的函数列在给定区间是否一致收敛；
    \begin{enumerate}
        \item[(1)] $$f_n(x)=\frac{1}{2^n+x^2},\quad x\in\rb$$
        
        \sol{

        }

        \item[(3)] $$f_n(x)=\ln\bigg(1+\frac{x^2}{n^2}\bigg)$$
        区间为给定参数$l>0$的$(-l,l)$与$\rb$。

        \sol{

        }

        \item[(5)] $$f_n(x)=\frac{\sin x^n}{3+x^n}$$
        区间为给定参数$\delta\in(0,1)$的$[0,1-\delta]$与$(0,1)$。
        
        \sol{

        }
    \end{enumerate}
    \item (习题10.4.3)判断下列关于$x$的函数项级数在给定区间是否一致收敛；
    \begin{enumerate}
        \item[(1)] $$\suminf(-1)^n\frac{\sqrt n}{x^2+n^2},\quad x\in\rb$$
        
        \sol{

        }

        \item[(3)] $$\suminf\frac{\sin(nx)}{\sqrt{1+(x^2+n^2)^3}},\quad x\in\rb$$
        
        \sol{

        }

        \item[(5)] $$\suminf\frac{x^2}{(1+x)^n},\quad x\in(0,+\infty)$$
        
        \sol{

        }
    \end{enumerate}

    \item (习题10.4.6)给定参数$\delta>0$，证明函数项级数
    $$\zeta(x)=\suminf\frac{1}{n^x},\quad\suminf\frac{\ln n}{n^x}$$
    都在$[1+\delta,+\infty)$一致收敛，由此证明$\zeta(x)$在$(1,+\infty)$有连续的导函数。

    \note 这里的$\zeta(x)$称为\textbf{黎曼Zeta函数}。

    \sol{

    }

    \item (总练习题10.4)证明Dirichlet级数(定义见本部分第3题)的收敛域必然为$(\alpha,+\infty)$或$[\alpha,+\infty)$的形式，且它在收敛域中的任何闭区间上一致收敛。

    
    \sol{

    }

    \item (总练习题10.5)证明Dirichlet级数(定义见本部分第3题)在收敛域中有连续的导函数。
    
    \sol{

    }
    
\end{enumerate}

\subsection{幂级数}
\subsubsection{局部逼近}
刻画局部

求导确定系数

\subsubsection{敛散性}
弱化到内闭一致

逐项求导与逐项积分

\subsubsection{初等函数}
证明初等函数泰勒级数存在处收敛

向初等函数泰勒级数配凑的计算技巧

\subsection{傅里叶级数简介}
\subsubsection{整体逼近}
刻画整体频率

积分确定系数

\subsubsection{性质简介}
弱化到$L^2$

\newpage
\section{积分敛散性}
\subsection{习题解答}
\subsubsection{作业解答}
\begin{enumerate}
    \item (习题10.6.1)利用初等函数的展开式写出下列函数在$x=0$处的幂级数展开式，并求收敛域：
    \begin{enumerate}
        \item[(5)] $$(1+x)\er^{-x}$$
        
        \sol{

        }

        \item[(7)] $$\sin^3x$$
        
        \sol{

        }

        \item[(10)] $$\frac{5x-2}{x^2+5x-6}$$
        
        \sol{

        }
    \end{enumerate}

    \item (习题10.6.2(2))利用逐项微分/积分求如下函数在$x=0$处的幂级数展开式：
    $$\ln(x+\sqrt{1+x^2})$$

    \sol{

    }
    
    \item (习题10.6.3)证明
    $$\forall x\ne0,\quad\suminf[n=0]\frac{x^n}{(n+1)!}=\frac{1}{x}(\er^x-1)$$
    并证明
    $$\suminf\frac{n}{(n+1)!}=1$$

    \sol{

    }
    
    \item (总练习题10.9)给定$x_0\in\rb$、$a>0$，若$f(x)$在$(x_0-a,x_0+a)$有定义，有任意阶导函数，且存在$M>0$使得
    $$\forall n\in\mathbb{N},\quad |f^{(n)}(x)|\le M$$
    证明
    $$\forall x\in(x_0-a,x_0+a),\quad f(x)=\suminf[n=0]\frac{1}{n!}f^{(n)}(x_0)(x-x_0)^n$$

    \sol{

    }
    
    \item (习题11.1.1)判断下列广义积分敛散性，若收敛，求出其值：
    \begin{enumerate}
        \item[(2)] $$\int_0^{+\infty}\frac{\dr x}{(x+1)(x+2)}$$
        
        \sol{

        }

        \item[(6)] $$\int_4^{+\infty}\frac{\dr x}{x\sqrt{x+1}}$$
        
        \sol{

        }

        \item[(10)] $$\int_{-1}^1\frac{\dr x}{\sqrt{1-x^2}}$$
        
        \sol{

        }
    \end{enumerate}

    \item (习题11.1.2(2))证明$\sigma>0$时
    $$\frac{1}{\sigma\sqrt{2\pi}}\int_{-\infty}^{+\infty}x\er^{-\frac{(x-a)^2}{2\sigma^2}}\dr x=a$$
    
    \item (习题11.1.3)判断下列广义积分敛散性：
    \begin{enumerate}
        \item[(3)] $$\int_0^2\frac{\dr x}{\sqrt[3]{x}+3\sqrt[4]{x}+x^3}$$
        
        \sol{

        }

        \item[(8)] $$\int_0^1\frac{\ln x}{1-x}\dr x$$
        
        \sol{

        }

        \item[(9)] $$\int_0^{\pi/2}\frac{\dr x}{\sin^2x\cos^2x}$$
        
        \sol{

        }
    \end{enumerate}

    \item (习题11.1.4(1))判断如下积分是条件收敛、绝对收敛还是发散：
    $$\int_0^{+\infty}\frac{\sqrt{x}\cos x}{x+3}\dr x$$

    \sol{

    }

    \item (习题11.1.6)设$f(x)$定义在$(a,b]$上并以$a$为瑕点。通过换元将瑕积分
    $$\int_a^bf(x)\dr x$$
    转换为同敛散性的无穷积分，并证明敛散性相同。

    \sol{

    }
\end{enumerate}

\subsubsection{补充题解答}
\begin{enumerate}
    \item (习题10.5.1)求下列关于$x$的幂级数的收敛半径：
    \begin{enumerate}
        \item[(1)] $$\suminf\frac{x^n}{n2^n}$$
        
        \sol{

        }

        \item[(3)] $$\suminf\frac{n!}{n^n}x^n$$
        
        \sol{

        }
    \end{enumerate}

    \item (习题10.5.2)求下列关于$x$的幂级数的收敛区间与收敛域：
    \begin{enumerate}
        \item[(3)] $$\suminf[n=0](-1)^n\frac{x^{2n+1}}{(2n+1)(2n+1)!}$$
        
        \sol{

        }

        \item[(5)] $$\suminf[n=0]x^{2n+1}$$
        
        \sol{

        }

        \item[(9)] $$\suminf\bigg(1+\frac{1}{2}+\dots+\frac{1}{n}\bigg)x^n$$
        
        \sol{

        }
    \end{enumerate}

    \item (习题10.5.3)求下列关于$x$的幂级数的和函数：
    \begin{enumerate}
        \item[(2)] $$\suminf(-1)^n(2n-1)x^{2n-2}$$
        
        \sol{

        }

        \item[(6)] $$\suminf\frac{2n-1}{2^n}x^{2n-2}$$
        
        \sol{

        }
    \end{enumerate}

    \item (习题10.6.1)利用初等函数的展开式写出下列函数在$x=0$处的幂级数展开式，并求收敛域：
    \begin{enumerate}
        \item[(8)] $$\sin\bigg(\frac{\pi}{4}+x\bigg)$$
        
        \sol{

        }

        \item[(9)] $$\ln(1+x-2x^2)$$
        
        \sol{

        }
    \end{enumerate}

    \item (习题10.6.2(3))利用逐项微分/积分求如下函数在$x=0$处的幂级数展开式：
    $$\int_0^x\frac{\dr t}{\sqrt{1-t^4}}$$

    \sol{

    }
    
    \item (习题10.6.4)证明
    $$\suminf\frac{n}{(n+2)!}=3-\er$$

    \sol{
        介绍两种方法：
        \begin{itemize}
            \item \textbf{直接变形}
            
            由于已知$\suminf[n=0]\frac{1}{n!}=\er$，我们试图变形出结果。拆分可发现
            $$\suminf\frac{n}{(n+2)!}=\suminf\frac{n+2}{(n+2)!}-\suminf\frac{2}{(n+2)!}$$
            也即原式为
            $$\suminf\frac{1}{(n+1)!}-2\suminf\frac{1}{(n+2)!}$$
            而将$n$替换为$n-1$或$n-2$即可进一步写为
            $$\suminf[n=2]\frac{1}{n!}-2\suminf[n=3]\frac{1}{n!}=(\er-1-1)-2\bigg(\er-1-1-\frac{1}{2}\bigg)=3-\er$$
            
            \item \textbf{构造级数}
        \end{itemize}
    }
    
    \item (习题10.6.5)求下列定积分的近似值，精确到小数点后第三位：
    \begin{enumerate}[(1)]
        \item $$\int_0^2\frac{\sin x}{x}\dr x$$
        
        \sol{

        }
        
        \item $$\int_0^1\cos(x^2)\dr x$$
        
        \sol{

        }
    \end{enumerate}
    
    \item (总练习题10.8)定义复函数$\er^z$为(敛散性见本讲义29.1.1第7题)
    $$\er^z=\suminf[n=0]\frac{1}{n!}z^n$$
    证明欧拉公式
    $$\forall x\in\rb,\quad \er^{\ir x}=\cos x+\ir\sin x$$

    \sol{

    }
\end{enumerate}

\subsection{广义积分}
\subsubsection{无穷积分与瑕积分}
定义

直接计算方法(微积分基本定理)

\subsubsection{正函数的广义积分}
单调有界函数必存极限

比较-阶估算

各种估算方法

\subsubsection{一般广义积分}
绝对收敛推收敛

Dirichlet、Abel

柯西收敛原理

\subsection{含参积分}
\subsubsection{含参的常义积分}
类比级数有限项，讨论有限区间常义积分

有限区间的连续、可微、可积

\subsubsection{一致收敛}
一致收敛将有限放大到无限

思考：瑕点不公共可否谈论？

连续、可微、可积

\subsubsection{收敛判别}
直接计算

比较(与普通广义积分)

绝对一致收敛

Dirichlet、Abel

柯西收敛原理

\newpage
\section{积分技巧补充}
\subsection{习题解答}
\subsubsection{作业解答}
\begin{enumerate}
    \item (习题11.2.1(4))
    \item (习题11.2.2)
    \begin{enumerate}
        \item[(1)] 
        \item[(3)] 
    \end{enumerate}
    \item (习题11.2.3(3))
    \item (习题11.3.1)
    \begin{enumerate}
        \item[(2)] 
        \item[(4)] 
        \item[(6)] 
    \end{enumerate}
    \item (习题11.3.2)
    \begin{enumerate}
        \item[(1)] 
        \item[(3)] 
        \item[(5)] 
    \end{enumerate}
    \item (习题11.3.3)
    \item (习题11.3.4)
    \begin{enumerate}
        \item[(2)] 
        \item[(6)] 
        \item[(7)] 
        \item[(9)] 
    \end{enumerate}
\end{enumerate}
\subsubsection{补充题解答}
\begin{enumerate}
    \item (习题11.1.1)
    \begin{enumerate}
        \item[(3)] 
        \item[(4)] 
        \item[(13)] 
    \end{enumerate}
    \item (习题11.1.2(1))
    \item (习题11.1.3(10))
    \item (习题11.1.4(2))
    \item (习题11.1.5)
    \item (习题11.2.1)
    \begin{enumerate}
        \item[(3)] 
        \item[(5)] 
    \end{enumerate}
    \item (习题11.2.2)
    \begin{enumerate}
        \item[(2)] 
        \item[(4)] 
    \end{enumerate}
    \item (习题11.2.3(2))
    \item (习题11.2.4)
    \item (习题11.3.1)
    \begin{enumerate}
        \item[(1)] 
        \item[(3)] 
        \item[(5)] 
    \end{enumerate}
\end{enumerate}

\subsection{敛散性}
\subsubsection{广义积分}
\begin{framed}
\begin{exercise}
    构造在$[0,+\infty)$上连续的函数$f$满足积分
    $$\int_0^{+\infty}f(x)\dr x$$
    收敛，但
    $$\lim_{x\to\infty}f(x)$$
    不存在。
\end{exercise}
\sol{

}
\end{framed}

\begin{framed}
\begin{exercise}
    设在$[0,+\infty)$上连续的函数$f$满足积分
    $$\int_0^{+\infty}f(x)\dr x$$
    收敛，证明
    $$\lim_{x\to\infty}f(x)$$
    若存在则为0。
\end{exercise}
\sol{

}
\end{framed}

\begin{framed}
\begin{exercise}
    证明
    $$\int_0^{+\infty}\frac{\sin x\dr x}{\sqrt{x}+\sin x}$$
    发散。
\end{exercise}
\sol{

}
\end{framed}

\begin{framed}
\begin{exercise}
    计算
    $$\lim_{n\to\infty}\frac{1^{\alpha-1}+2^{\alpha-1}+\cdots+n^{\alpha-1}}{n^\alpha}$$
    这里参数$\alpha>0$。
\end{exercise}
\sol{

}
\end{framed}

\subsubsection{含参积分}

\begin{framed}
\begin{exercise}
    证明含有参数$u$的广义积分
    $$\int_0^{+\infty}\frac{\sin(3x)}{x+u}\er^{-ux}\dr x$$
    在$[0,+\infty)$上一致收敛。
\end{exercise}
\sol{

}
\end{framed}

\begin{framed}
\begin{exercise}
    证明
    $$\int_0^{+\infty}\frac{x\dr x}{1+x^\alpha\cos^2x}$$
    收敛，这里参数$\alpha>4$。判断它在$(4,+\infty]$是否一致收敛。
\end{exercise}
\sol{
    我们介绍两种方法：
    \begin{itemize}
        \item \textbf{分段放缩}
        
        \item \textbf{整体估算}
        
        
    \end{itemize}
}
\end{framed}

\begin{framed}
\begin{exercise}
    判断
    $$\int_0^{+\infty}\bigg(\bigg(1-\frac{\sin x}{x}\bigg)^{-\alpha}-1\bigg)\dr x$$
    的敛散性，这里参数$\alpha>0$。进一步研究它对参数$\alpha$在何区间一致收敛。
\end{exercise}
\sol{

}
\end{framed}

\subsection{计算}
\subsubsection{直接计算}
\begin{framed}
\begin{exercise}
    计算
    $$\int_0^{\pi/2}\ln\sin x\dr x$$
\end{exercise}
\sol{

}
\end{framed}

\begin{framed}
\begin{exercise}
    计算
    $$\int_0^{+\infty}\frac{\sin x^2}{x}\dr x$$
\end{exercise}
\sol{

}
\end{framed}

\begin{framed}
\begin{exercise}
    证明
    $$\int_0^{+\infty}f\bigg(\bigg(ax-\frac{b}{x}\bigg)^2\bigg)\dr x=\frac{1}{a}\int_0^{+\infty}f(x^2)\dr x$$
    这里$f$在$[0,+\infty)$连续且右侧积分存在，参数$a>0$、$b>0$。
\end{exercise}
\sol{

}
\end{framed}

\begin{framed}
\begin{exercise}
    设在$[0,+\infty)$上连续的函数$f$满足积分
    $$\int_1^{+\infty}\frac{f(x)}{x}\dr x$$
    收敛，证明
    $$\int_0^{+\infty}\frac{f(ax)-f(bx)}{x}\dr x=f(0)\ln\frac{b}{a}$$
    这里参数$a>0$、$b>0$。
\end{exercise}
\sol{

}
\end{framed}

\subsubsection{一致性技巧}
\begin{framed}
\begin{exercise}
    计算
    $$\int_0^{+\infty}\sin(x^2)\dr x$$
\end{exercise}
\sol{

}
\end{framed}

\begin{framed}
\begin{exercise}
    计算
    $$\int_0^{+\infty}\frac{\er^{-ax}-\er^{-bx}}{x}\sin x\dr x$$
    这里参数$a>0$、$b>0$。
\end{exercise}
\sol{

}
\end{framed}

\begin{framed}
\begin{exercise}
    计算
    $$\int_0^{+\infty}\frac{x^{p-1}}{1+x}\dr x$$
    这里参数$p\in(0,1)$。
\end{exercise}
\sol{

}
\end{framed}

\subsubsection{特殊函数}
\begin{framed}
\begin{exercise}
    证明余元公式
    $$\forall p\in(0,1),\quad\Gamma(p)\Gamma(1-p)=\frac{\pi}{\sin(p\pi)}$$
\end{exercise}
\end{framed}

\begin{framed}
\begin{exercise}
    计算
    $$\suminf\frac{1}{nC_{2n}^n}$$
\end{exercise}
\end{framed}

\begin{framed}
\begin{exercise}
    证明Stirling公式
    $$n!\sim\sqrt{2\pi n}\bigg(\frac{n}{\er}\bigg)^n$$
\end{exercise}
\end{framed}

\newpage
\section{傅里叶级数}
\subsection{习题解答}
\subsubsection{作业解答}
\begin{enumerate}
    \item (习题12.1.2)
    \item (习题12.2.1)
    \begin{enumerate}
        \item[(1)] 
        \item[(3)] 
        \item[(5)] 
    \end{enumerate}
    \item (习题12.2.3)
    \item (习题12.2.5)
    \item (习题12.2.7)
    \item (习题12.2.9)
    \item (习题12.2.11)
    \item (总练习题12.1)
    \item (总练习题12.3)
    \item (总练习题12.5)
    \item (总练习题12.6)
\end{enumerate}

\subsubsection{补充题解答}
\begin{enumerate}
    \item (习题11.3.2)
    \begin{enumerate}
        \item[(2)] 
        \item[(4)] 
    \end{enumerate}
    \item (习题11.3.4)
    \begin{enumerate}
        \item[(1)]
        \item[(3)]  
        \item[(4)] 
        \item[(5)] 
        \item[(8)] 
    \end{enumerate}
    \item (习题12.1.1)
    \item (习题12.1.3)
    \item (习题12.2.1)
    \begin{enumerate}
        \item[(2)]
        \item[(4)]
        \item[(6)]   
    \end{enumerate}
    \item (总练习题12.2)
    \item (总练习题12.4)
    \item (习题12.2.2)
    \item (习题12.2.4)
    \item (习题12.2.8)
    \item (习题12.3.1)
    \item (习题12.3.2)
\end{enumerate}

\subsection{敛散性}
\subsubsection{基本性质}
三角函数系正交性

系数趋于0

\subsubsection{平方积分敛散性}

\subsubsection{逐点敛散性}

\subsection{计算}
\subsubsection{一般区间的展开}
周期延拓

奇/偶延拓与对应级数

\subsubsection{级数计算}
计算数项级数

\subsubsection{傅里叶变换}

\newpage
\section{期末复习}
本章为针对教材第十章到第十二章的期末考试复习题，请大家\textbf{务必关注每题解答后的注释}，强调的重点往往在注释中。

\subsection{级数复习题}
\subsubsection{数项级数}
\begin{framed}
\begin{exercise}
    若数列$a_n$分别满足以下三个条件：
    $$\lim_{n\to\infty}na_n=0,\quad\lim_{n\to\infty}na_n=+\infty,\quad\lim_{n\to\infty}|na_n|=+\infty$$
    能否判断$\suminf a_n$的敛散性？对每个条件，若能则证明，若不能则举出例子。
\end{exercise}
\sol{
    我们对三个条件分别说明：
    \begin{itemize}
        \item \textbf{极限为0}
        
        无法判断敛散性。对$n>1$考虑
        $$a_n=\frac{1}{n\ln n},\quad b_n=\frac{1}{n\ln^2n}$$
        由定义可发现(利用$n\to\infty$时$\ln n\to+\infty$与乘除极限性质)
        $$\lim_{n\to\infty}na_n=\lim_{n\to\infty}\frac{1}{\ln n}=0,\quad\lim_{n\to\infty}nb_n=\lim_{n\to\infty}\frac{1}{\ln^2n}=0$$
        利用$\suminf[n=2]\frac{1}{n\ln^rn}$当且仅当$r>1$时收敛(见本讲义27.1.1第6题(11))可知$\suminf[n=2]a_n$发散，$\suminf[n=2]b_n$收敛。

        \item \textbf{极限为$+\infty$}
        
        必定发散。利用极限的\textbf{保序性}，存在$N$使得$n>N$时$na_n>1$，也即
        $$\exists N,\quad\forall n>N,\quad a_n>\frac{1}{n}$$
        由\textbf{$p$级数}敛散性$\suminf\frac{1}{n}$发散，从而根据\textbf{比较判别法}得$a_n$也发散。

        \item \textbf{绝对值极限为$+\infty$}
        
        无法判断敛散性。考虑
        $$a_n=\frac{1}{\sqrt{n}},\quad b_n=\frac{(-1)^n}{\sqrt{n}}$$
        由定义可发现
        $$\lim_{n\to\infty}|na_n|=\lim_{n\to\infty}|nb_n|=\sqrt{n}=+\infty$$
        利用\textbf{$p$级数}敛散性结论可知$\suminf a_n$发散；由于$b_n$为交错级数，且通项绝对值$\frac{1}{\sqrt{n}}$单调递减趋于0，利用\textbf{莱布尼茨判别法}可得$\suminf b_n$收敛。
    \end{itemize}
}
\end{framed}

\note \textbf{比较判别法}常用的级数为\textbf{等比级数}$q^n$、\textbf{$p$级数}$\frac{1}{n^p}$与$p$级数的\textbf{精细化}$\frac{1}{n\ln^rn}$\ (证明可见本讲义27.1.1第6题(11)，此结论在复杂的题目中一般可\textbf{直接使用}，但仍然建议学习\textbf{证明过程}作为积分判别法的应用实例)，收敛条件分别为$q\in(-1,1)$、$p>1$、$r>1$。

\note 由于丢弃有限项不影响级数敛散性，比较判别法可\textbf{丢弃有限项}后使用。这结合极限的\textbf{保序性}可以得到更强的结论，我们将在之后看到。

\begin{framed}
\begin{exercise}
    若正数列$a_n$\ (下标$n\ge2$)满足
    $$\lim_{n\to\infty}\log_n(a_n)=b$$
    给出使得$\suminf[n=2]a_n$必然收敛、必然发散、可能收敛可能发散的实数$b$取值范围。
\end{exercise}
\sol{
    由于已知$a_n$是正数列，我们可以采用\textbf{比较判别法}。而为了确定$b$的取值范围，我们需要先考虑最基本的$\log_n(a_n)=b$的情况。可以发现，这当且仅当$a_n=n^b$，由于\textbf{$p$级数}$\suminf\frac{1}{n^p}$当且仅当$p>1$时收敛，$\suminf n^b$收敛当且仅当$b<-1$。与比值判别法/柯西判别法类似，我们可以猜测$b<-1$时必然收敛，$b>-1$时必然发散，边界情况$b=-1$可能收敛可能发散。下面给出证明：
    \begin{itemize}
        \item $b<-1$\textbf{时收敛}
        
        设$c=\frac{b-1}{2}$，则$b<c<-1$，利用极限的\textbf{保序性}可知存在$N$使得$n>N$时
        $$\log_n(a_n)<c$$
        利用对数函数单调性即可以转化为
        $$a_n<n^c$$
        由\textbf{$p$级数}敛散性结论，由$c<-1$可知$\suminf[n=2]n^c$收敛，从而根据\textbf{比较判别法}可得$\suminf[n=2]a_n$也收敛。


        \item $b=-1$\textbf{时可能收敛可能发散}
        
        考虑
        $$a_n=\frac{1}{n},\quad b_n=\frac{1}{n\ln^2n}$$
        利用\textbf{$p$级数}敛散性结论可知$\suminf[n=2]a_n$发散，利用$\suminf[n=2]\frac{1}{n\ln^rn}$敛散性结论(见本讲义27.1.1第6题(11))可知$\suminf[n=2]b_n$收敛，且
        $$\lim_{n\to\infty}\log_n(a_n)=\lim_{n\to\infty}(-1)=-1$$
        下面证明
        $$\lim_{n\to\infty}\log_n(b_n)=-1$$
        这就说明了极限为$-1$时原级数可能收敛可能发散。

        直接利用对数函数性质计算可发现(第二个等号运用了换底公式)
        $$\log_n(b_n)=-1-2\log_n(\ln n)=-1-2\frac{\ln\ln n}{\ln n}$$
        而通过\textbf{归结原理}与\textbf{函数极限阶估算}结论(可见上学期讲义3.1第11题)可知(由$\ln n$在$n\to\infty$时极限为$+\infty$)
        $$\lim_{n\to\infty}\frac{\ln\ln n}{\ln n}=\lim_{x\to+\infty}\frac{\ln x}{x}=0$$
        于是根据四则运算的极限结论即得到
        $$\lim_{n\to\infty}\log_n(b_n)=-1$$
        从而证明了结论。
        
        \item $b>-1$\textbf{时发散}
        
        设$c=\frac{b-1}{2}$，则$b>c>-1$，利用极限的\textbf{保序性}可知存在$N$使得$n>N$时
        $$\log_n(a_n)>c$$
        利用对数函数单调性即可以转化为
        $$a_n>n^c$$
        由\textbf{$p$级数}敛散性结论，由$c>-1$可知$\suminf[n=2]n^c$发散，从而根据\textbf{比较判别法}可得$\suminf[n=2]a_n$也发散。
    \end{itemize}
}
\end{framed}

\note 这类问题本质上是考察对教材中基于比较判别法而产生的比值判别法、柯西判别法\textbf{证明过程}的理解，注意其中通过\textbf{保序性}放宽常数与\textbf{丢弃有限项}后再比较的思路。

\begin{framed}
\begin{exercise}
    判断如下级数是否收敛：
    $$\suminf[n=3]\frac{1}{(\ln\ln n)^{\ln\ln n}}$$
\end{exercise}
\sol{
    当$n\ge 3$时$\ln\ln n>0$，从而此级数为正项级数。先尝试化简分母，取对数得到
    $$(\ln\ln n)^{\ln\ln n}=\er^{(\ln\ln n)\cdot(\ln\ln\ln n)}$$
    为了能够与\textbf{$p$级数}进行比较，我们需要将它化为$n$的次方，由此进一步化为
    $$\er^{\ln n\cdot\frac{(\ln\ln n)\cdot(\ln\ln\ln n)}{\ln n}}=(\er^{\ln n})^{\frac{(\ln\ln n)\cdot(\ln\ln\ln n)}{\ln n}}=n^{\frac{(\ln\ln n)\cdot(\ln\ln\ln n)}{\ln n}}$$
    接下来的过程分为两步，先估算次方的大小，再进行比较：
    \begin{itemize}
        \item \textbf{处理}$\frac{(\ln\ln n)\cdot(\ln\ln\ln n)}{\ln n}$
        
        利用\textbf{函数极限的阶估算}结论(可见上学期讲义3.1第11题)可知
        $$\lim_{x\to\infty}\frac{\ln^2x}{x}=0$$
        而$n\ge3$时有$\ln\ln n>0$，于是$\ln\ln\ln n<\ln\ln n$，从而
        $$\frac{(\ln\ln n)\cdot(\ln\ln\ln n)}{\ln n}<\frac{\ln^2(\ln n)}{\ln n}$$
        利用\textbf{归结原理}即得到(由$\ln n$在$n\to\infty$时极限为$+\infty$)
        $$\lim_{n\to\infty}\frac{\ln^2(\ln n)}{\ln n}=\lim_{x\to+\infty}\frac{\ln^2x}{x}=0$$
        另一方面，由$\ln n$在$n\to\infty$时极限为$+\infty$与归结原理也可得到
        $$\lim_{n\to\infty}\ln\ln n=\lim_{x\to+\infty}\ln x=+\infty$$
        由此存在$N$使得$n>N$时$\ln\ln n>1$，于是$\ln\ln\ln n>0$，从而此时有
        $$\frac{(\ln\ln n)\cdot(\ln\ln\ln n)}{\ln n}>0$$
        利用\textbf{夹逼定理}可知
        $$\lim_{n\to\infty}\frac{(\ln\ln n)\cdot(\ln\ln\ln n)}{\ln n}=0$$
        
        \item \textbf{原级数判断}
        
        利用极限的保序性，存在$N$使得$n>N$时$\frac{(\ln\ln n)\cdot(\ln\ln\ln n)}{\ln n}<1$，此时有
        $$\frac{1}{(\ln\ln n)^{\ln\ln n}}=n^{\frac{(\ln\ln n)\cdot(\ln\ln\ln n)}{\ln n}}>\frac{1}{n}$$
        由\textbf{$p$级数}敛散性可知$\suminf[n=3]\frac{1}{n}$发散，从而根据\textbf{比较判别法}可得原级数也发散。
    \end{itemize}
}
\end{framed}

\note 注意\textbf{比较判别法只能对正项级数使用}，由此使用前需要先判断级数是否为正项(或去除有限项后为正项)。比较判别法的三个常用标准$\suminf q^n$、$\suminf\frac{1}{n^p}$、$\suminf[n=2]\frac{1}{n\ln^rn}$是逐渐\textbf{精细化}的，也即如果能与排在前的比较得出结果，一定能与排在后的比较得出结果。由此，建议优先尝试在中间的$\suminf\frac{1}{n^p}$。

\note 本学期仍然可能用到上学期的利用\textbf{对数性质}化简指数的技巧，更简单的应用可见本讲义28.1.2第1题(6)。

\begin{framed}
\begin{exercise}
    判断如下级数是否收敛：
    $$\suminf\bigg(\ln\frac{1}{n}-\ln\sin\frac{1}{n}\bigg)$$
\end{exercise}
\sol{
    直接计算可知其为
    $$\ln\frac{1}{n\sin\frac{1}{n}}$$
    由$n>0$可知$\sin\frac{1}{n}<\frac{1}{n}$，从而此级数为正项级数。

    为了对其进行估算，将它写为
    $$\ln\bigg(1+\frac{1-n\sin\frac{1}{n}}{n\sin\frac{1}{n}}\bigg)$$
    将$\frac{1}{n}$替换为趋于0的$x$，上式即为
    $$\ln\bigg(1+\frac{x-\sin x}{\sin x}\bigg)$$
    根据\textbf{基本等价无穷小}结论，$x\to0$时$x\sim\sin x$，于是$\frac{x-\sin x}{\sin x}$在$x\to0$时极限为0，从而根据\textbf{等价无穷小替换}(先整体看作$\ln(1+t)$，再替换分母)可知
    $$\ln\bigg(1+\frac{x-\sin x}{\sin x}\bigg)\sim\frac{x-\sin x}{\sin x}\sim\frac{x-\sin x}{x}\quad(x\to0)$$
    利用三角函数的\textbf{泰勒展开}结论可知
    $$\sin x=x-\frac{x^3}{6}+o(x^3)$$
    从而有(涉及的$o$运算性质可见上学期讲义11.3.2)
    $$\frac{x-\sin x}{x}=\frac{\frac{x^3}{6}+o(x^3)}{x}=\frac{x^2}{6}+o(x^2)\quad(x\to0)$$
    由此利用$o$的定义可知
    $$\lim_{x\to0}\frac{\frac{x-\sin x}{x}}{x^2}=\lim_{x\to0}\frac{\frac{x^2}{6}+o(x^2)}{x^2}=\frac{1}{6}$$
    也即
    $$\frac{x-\sin x}{x}\sim\frac{x^2}{6}\quad(x\to0)$$
    综合以上的讨论，我们最终得到了
    $$\ln\bigg(1+\frac{x-\sin x}{\sin x}\bigg)\sim\frac{x^2}{6}\quad(x\to0)$$
    由此，利用\textbf{归结原理}与\textbf{等价无穷小替换}可得到
    $$\lim_{n\to\infty}\frac{\ln\frac{1}{n\sin\frac{1}{n}}}{\frac{1}{n^2}}=\lim_{x\to0}\frac{1}{x^2}\cdot\ln\bigg(1+\frac{x-\sin x}{\sin x}\bigg)=\lim_{x\to0}\frac{1}{x^2}\cdot\frac{x^2}{6}=\frac{1}{6}$$
    利用\textbf{$p$级数}敛散性结论可知$\suminf\frac{1}{n^2}$收敛，从而根据\textbf{比较判别法}得原级数也收敛。
}
\end{framed}

\note 利用\textbf{归结原理}转化为函数极限后由泰勒展开、等价无穷小替换等方法(注意这些方法只能对函数极限使用)进行\textbf{阶估算}常与\textbf{比较判别法}一起出现。

\begin{framed}
\begin{exercise}
    判断如下级数是否收敛：
    $$\suminf\bigg(\cos\frac{1}{\sqrt{n}}\bigg)^{n^2}$$
\end{exercise}
\sol{
    对正整数$n$，$0<\frac{1}{\sqrt{n}}<1$，从而此级数为正项级数。仍然采用先\textbf{估算}后\textbf{证明}的思路，并提供两种证明思路：
    \begin{itemize}
        \item \textbf{估算}
        
        为了对其进行估算，将\textbf{指数转化为对数}写为
        $$\er^{n^2\ln\cos\frac{1}{\sqrt{n}}}$$
        将$\frac{1}{n}$替换为趋于$0^+$的$x$\ (注意此处必须考虑\textbf{右极限}，也即保证$x$恒为正以在根号定义域内)，上式即为
        $$\er^{\frac{\ln\cos\sqrt{x}}{x^2}}$$
        我们先估算指数部分，由于$x\to0$时$\cos\sqrt{x}\to1$，利用$\ln$在\textbf{1处的泰勒展开}可知
        $$\ln\cos\sqrt{x}=\ln(1+(\cos\sqrt{x}-1))=\cos\sqrt{x}-1+o(\cos\sqrt{x}-1)\quad(x\to0^+)$$
        再对$\cos t$在0处进行泰勒展开，并代入$t=\sqrt{x}$可知(由于$x>0$，$(\sqrt{x})^2$成立)
        $$\cos\sqrt{x}-1=-\frac{x}{2}+o(x)\quad(x\to0^+)$$
        从而有$o(\cos\sqrt{x}-1)=o(x)$，进一步综合得到(涉及的$o$运算性质可见上学期讲义11.3.2)
        $$\ln\cos\sqrt{x}=-\frac{x}{2}+o(x)\quad(x\to0^+)$$
        这也就说明
        $$\lim_{x\to0^+}\frac{\ln\cos\sqrt{x}}{x}=\lim_{x\to0^+}\bigg(-\frac{1}{2}+\frac{o(x)}{x}\bigg)=-\frac{1}{2}$$

        \item \textbf{原级数判断-直接估算}
        
        由于原级数为
        $$\suminf\big(\er^{n\ln\cos\frac{1}{\sqrt{n}}}\big)^n$$
        可以考虑适当底数放大。具体来说，利用\textbf{归结原理}与\textbf{复合函数极限}结论可以得到
        $$\lim_{n\to\infty}\er^{n\ln\cos\frac{1}{\sqrt{n}}}=\lim_{x\to0^+}\er^{\frac{\ln\cos\sqrt{x}}{x}}=\er^{-1/2}$$
        由此，取$c\in(\er^{-1/2},1)$根据极限\textbf{保序性}可知存在$N$使得$n>N$时
        $$\er^{n\ln\cos\frac{1}{\sqrt{n}}}<c$$
        从而此时
        $$\er^{n^2\ln\cos\frac{1}{\sqrt{n}}}<c^n$$
        利用\textbf{等比级数}敛散性结论可知$\suminf c^n$收敛，从而根据\textbf{比较判别法}得原级数也收敛。

        \item \textbf{原级数判断-柯西判别}
        
        与上一种方法相同可得
        $$\lim_{n\to\infty}\er^{n\ln\cos\frac{1}{\sqrt{n}}}=\er^{-1/2}$$
        而此即为
        $$\bigg(\cos\frac{1}{\sqrt{n}}\bigg)^n=\sqrt[n]{\bigg(\cos\frac{1}{\sqrt{n}}\bigg)^{n^2}}$$
        由于$\er^{-1/2}<1$，利用\textbf{柯西判别法}得原级数收敛。
    \end{itemize}
}
\end{framed}

\note 一般利用泰勒展开估算阶时只需要估算出\textbf{非零的最低次项}，其他部分可放缩解决。

\note 当级数的通项包含$n$次方时\textbf{未必可以直接通过柯西判别法判别}，由此仍然需要\textbf{先估算}，根据估算结果进一步处理。

\begin{framed}
\begin{exercise}
    判断以下三个级数是否收敛：
    $$\suminf[n=2]\frac{1}{n\sqrt{\ln n}},\quad\suminf[n=2]\frac{1}{n\ln\sqrt{n}},\quad\suminf\frac{1}{n\ln(1+n+n^2)}$$
\end{exercise}
\sol{
    根据定义可发现三个级数均为正项级数。我们分别判断：
    \begin{itemize}
        \item \textbf{第一个级数}
        
        利用$\suminf[n=2]\frac{1}{n\ln^rn}$的敛散性结论(见本讲义27.1.1第6题(11))可知原级数发散。

        \item \textbf{第二个级数}
        
        由于
        $$\suminf[n=2]\frac{1}{n\ln\sqrt{n}}=\suminf[n=2]\frac{2}{n\ln n}$$
        根据数项级数的基本性质，它的敛散性与$\suminf[n=2]\frac{1}{n\ln n}$相同，利用$\suminf[n=2]\frac{1}{n\ln^rn}$的敛散性结论(见本讲义27.1.1第6题(11))可知原级数发散。

        \item \textbf{第三个级数}
        
        由于$\ln(1+n+n^2)$与$\ln(n^2)$接近，而后者为$2\ln n$，我们尝试用此级数进行比较。

        考虑级数$\suminf[n=2]\frac{1}{n\ln n}$，利用$\suminf[n=2]\frac{1}{n\ln^rn}$的敛散性结论(见本讲义27.1.1第6题(11))可知其发散，而利用对数函数的性质可得
        $$\lim_{n\to\infty}\frac{\frac{1}{n\ln(1+n+n^2)}}{\frac{1}{n\ln n}}=\lim_{n\to\infty}\frac{\ln n}{\ln(1+n+n^2)}=\lim_{n\to\infty}\frac{\ln n}{2\ln n+\ln\big(\frac{1}{n^2}+\frac{1}{n}+1\big)}$$
        由此通过\textbf{归结原理}与\textbf{四则运算的极限}结论进一步计算知上式为
        $$\lim_{x\to+\infty}\frac{\ln x}{2\ln x+\ln\big(\frac{1}{x^2}+\frac{1}{x}+1\big)}=\lim_{x\to+\infty}\frac{1}{2+\ln(\frac{1}{x^2}+\frac{1}{x}+1)\cdot\frac{1}{\ln x}}=\frac{1}{2+0\cdot0}=\frac{1}{2}$$
        这里$x\to+\infty$时$\ln(\frac{1}{x^2}+\frac{1}{x}+1)\to0$是由于\textbf{复合函数极限}结论(将$x$换元为$\frac{1}{t}$)与初等函数\textbf{连续性}。上述计算最终得到了
        $$\lim_{n\to\infty}\frac{\frac{1}{n\ln(1+n+n^2)}}{\frac{1}{n\ln n}}=\frac{1}{2}$$
        从而根据\textbf{比较判别法}得原级数也发散。
        
        \note 也可在$n$充分大时直接将通项放大到$\frac{1}{n\ln(n^3)}$后比较。
    \end{itemize}
}
\end{framed}
\note 注意对数中的次方可以视为\textbf{常系数}，不影响它的\textbf{阶}。

\begin{framed}
\begin{exercise}\label{mono plus unconv}
    判断如下级数是否收敛：
    $$\suminf\frac{n!\cdot2^n}{n^n},\quad\suminf\frac{n!\cdot\er^n}{n^n},\quad\suminf\frac{n^n}{(n^2+3n+1)^{\frac{n+2}{2}}}$$
\end{exercise}
\sol{
    根据定义可发现三个级数均为正项级数。我们分别判断：
    \begin{itemize}
        \item \textbf{第一个级数}
        
        设其第$n$项为$u_n$，直接计算有
        $$\frac{u_{n+1}}{u_n}=\frac{(n+1)!2^{n+1}}{(n+1)^{n+1}}\cdot\frac{n^n}{n!2^n}=\frac{2(n+1)}{(n+1)^{n+1}}\cdot n^n=2\bigg(\frac{n}{n+1}\bigg)^n=2\frac{1}{\big(1+\frac{1}{n}\big)^n}$$
        利用$\er$的定义与四则运算的极限结论可知
        $$\lim_{n\to\infty}2\frac{1}{\big(1+\frac{1}{n}\big)^n}=\frac{2}{\er}<1$$
        从而根据\textbf{比值判别法}可得原级数收敛。

        \note 也可直接将其放大为公比更大的等比数列证明收敛。
        
        \item \textbf{第二个级数}
        
        \note 尝试使用比值判别法会发现$\lim_{n\to\infty}\frac{u_{n+1}}{u_n}=1$，因此无法通过此方法判断。
        
        设其第$n$项为$u_n$，与上一个级数类似计算得
        $$\frac{u_{n+1}}{u_n}=\frac{\er}{\big(1+\frac{1}{n}\big)^n}$$
        利用数列极限结论可知$\big(1+\frac{1}{n}\big)^n$单调递增趋于$\er$，从而其恒小于等于$\er$，也即得到
        $$\forall n\in\mathbb{N}^*,\quad\frac{u_{n+1}}{u_n}>1$$
        这说明$u_n$随着$n$增加而\textbf{单调递增}，且由定义$u_1=\er$。

        若此级数收敛，应当有
        $$\lim_{n\to\infty}u_n=0$$
        从而利用极限\textbf{保序性}，存在$N$使得$n>N$时$u_n<\er=u_1$，与其单调递增矛盾。由此可知原级数通项极限不为0，发散。
        
        \item \textbf{第三个级数}
        
        \note 尝试使用柯西判别法会发现$\lim_{n\to\infty}\sqrt[n]{u_n}=1$，因此无法通过此方式判断。

        由于分母的底数接近$n^2$，可考虑先将分母的底数缩小至$n^2$观察，此时可以得到
        $$\frac{n^n}{(n^2+3n+1)^{\frac{n+2}{2}}}<\frac{n^n}{(n^2)^\frac{n+2}{2}}=\frac{1}{n^2}$$
        利用\textbf{$p$级数}敛散性结论可知$\suminf\frac{1}{n^2}$收敛，从而根据\textbf{比较判别法}已经可以得到原级数也收敛。
    \end{itemize}
}
\end{framed}

\note 注意比值判别法与柯西判别法一般必须\textbf{极限小于或大于1}才能得到结论，每项小于1、每项大于1并没有意义。不过，$\frac{u_{n+1}}{u_n}$恒大于1是\textbf{足以说明发散}的，类似上方第二个级数。

\note 总的来说，对于无法直接算出部分和的正项级数，或直接进行\textbf{积分判别}，或在阶估算后进行\textbf{比较判别}(柯西判别与比值判别均为其衍生，使用范围相对狭窄)。如果二者都行不通，可以尝试使用\textbf{柯西收敛原理}进行直接的估算，但这种情况往往过于复杂。

\begin{framed}
\begin{exercise}
    判断如下级数是否收敛：
    $$\suminf[n=2]\frac{(-1)^n}{\sqrt{n}+(-1)^n}$$
\end{exercise}
\sol{
    \note 虽然此为交错级数，但计算可发现它的通项\textbf{并不单调递减}，由此无法使用莱布尼茨判别法。 

    由于$n\to\infty$时$\sqrt{n}$的大小将远大于$(-1)^n$，此级数的通项接近$\frac{(-1)^n}{\sqrt{n}}$，由此将过程分为两步：
    \begin{itemize}
        \item $\suminf[n=2]\frac{(-1)^n}{\sqrt{n}}$\textbf{敛散性分析}
        
        由于$\frac{(-1)^n}{\sqrt{n}}$正负交替，此级数为交错级数，且通项绝对值$\frac{1}{\sqrt{n}}$单调递减趋于0，利用\textbf{莱布尼茨判别法}可得$\suminf[n=2]\frac{(-1)^n}{\sqrt{n}}$收敛。
        
        \item \textbf{估算误差}
        
        利用数项级数的基本性质(收敛与收敛之和收敛、收敛与发散之和发散)，\textbf{原级数通项减去收敛级数的通项并不影响敛散性}，由此原级数敛散性与
        $$\suminf[n=2]\bigg(\frac{(-1)^n}{\sqrt{n}+(-1)^n}-\frac{(-1)^n}{\sqrt{n}}\bigg)$$
        相同，直接计算可知其为
        $$\suminf[n=2]\frac{(-1)^n\cdot(-1)^{n+1}}{\sqrt{n}(\sqrt{n}+(-1)^n)}=\suminf[n=2]\frac{-1}{\sqrt{n}(\sqrt{n}+(-1)^n)}$$
        利用数项级数基本性质可知分子的$-1$改为1不影响敛散性，此时这已经是一个\textbf{正项级数}，由此可以使用阶估算的方法。由于分母的阶与$n$相同，可直接进行比较。

        考虑级数$\suminf[n=2]\frac{1}{n}$，利用\textbf{$p$级数}的敛散性结论可知其发散，而利用\textbf{有界量乘无穷小是无穷小}(从而$\frac{(-1)^n}{\sqrt{n}}$在$n\to\infty$时极限为0)与\textbf{四则运算的极限}结论直接计算可知
        $$\lim_{n\to\infty}\frac{\frac{1}{\sqrt{n}(\sqrt{n}+(-1)^n)}}{\frac{1}{n}}=\lim_{n\to\infty}\frac{1}{1+\frac{(-1)^n}{\sqrt{n}}}=\frac{1}{1+0}=1$$
        从而根据\textbf{比较判别法}得原级数也发散。
    \end{itemize}
}
\end{framed}

\note 通项减去收敛级数的通项不影响敛散性保证了\textbf{先判断主体再估算误差}的可行性。这在高等数学中是非常重要的技巧，上学期讲义第九章中给出了很多的应用。

\note 注意\textbf{比较判别法只能对正项级数使用}。计算可知本题中的级数通项与$\frac{(-1)^n}{\sqrt{n}}$等价，但两者敛散性并不相同。

\begin{framed}
\begin{exercise}
    判断如下级数是条件收敛、绝对收敛还是发散：
    $$\suminf[n=2]\frac{\sin(\sqrt{n^2+a^2}\pi)}{\ln^pn}$$
    这里参数$a\ne0$、$p\ge0$。
\end{exercise}
\sol{
    分为如下步骤：
    \begin{itemize}
        \item \textbf{分析与变形}
        
        \note 本段分析与本讲义28.1.1第1题(10)完全类似，只是需要一些额外的分析。
                
        可以发现，我们通常对数列的阶估算是利用\textbf{归结原理}将数列转化为了函数进行处理，但本题中，由于$\sin(\sqrt{x^2+a^2}\pi)$会不断在$-1$到1之间振荡，直接转化为函数是没有意义的。

        为了能够进行有意义的估算，我们需要观察数列本身的特殊性。可以发现$\sqrt{n^2+a^2}\pi$事实上与$n\pi$接近，而$\sin(n\pi)=0$，由此可将其转化为
        $$\sin((\sqrt{n^2+a^2}-n)\pi+n\pi)=(-1)^n\sin((\sqrt{n^2+a^2}-n)\pi)$$
        进一步进行分子有理化，此即为
        $$(-1)^n\sin\frac{a^2\pi}{\sqrt{n^2+a^2}+n}$$
        这样已经可以看出分母趋于0，转化为函数后仍然有意义，从而可以进行进一步估算。经过以上分析，我们最终将级数转化为了
        $$\suminf(-1)^n\frac{1}{\ln^pn}\sin\frac{a^2\pi}{\sqrt{n^2+a^2}+n}$$
        
        由于$n$为正整数时
        $$0<\frac{a^2\pi}{\sqrt{n^2+a^2}+n}<\frac{a^2\pi}{2n}$$
        且$\lim_{n\to\infty}\frac{a^2\pi}{2n}=0$，利用\textbf{夹逼定理}可知
        $$\lim_{n\to\infty}\frac{a^2\pi}{\sqrt{n^2+a^2}+n}=0$$
        由此，根据极限\textbf{保序性}，存在$N$使得$n>N$时
        $$\frac{a^2\pi}{\sqrt{n^2+a^2}+n}<\frac{\pi}{2}$$
        进一步利用幂函数性质可发现$\frac{a^2\pi}{\sqrt{n^2+a^2}+n}$单调减且恒大于0，由此利用三角函数的单调性结论，在$n>N$时
        $$\sin\frac{a^2\pi}{\sqrt{n^2+a^2}+n}$$
        单调减且恒大于0。

        \item \textbf{绝对收敛判定}
        
        利用\textbf{去除有限项不影响敛散性}，可只考虑$n>N$的项。根据上述讨论，当$n>N$时通项的绝对值即为
        $$\frac{1}{\ln^pn}\sin\frac{a^2\pi}{\sqrt{n^2+a^2}+n}$$
        由于此形式与$\frac{1}{\ln^pn}\sin\frac{a^2\pi}{2n}$接近，可以想到与$\frac{1}{n\ln^pn}$比较。
        
        考虑级数$\suminf[n=2]\frac{1}{n\ln^pn}$，利用\textbf{$p$级数}敛散性结论可知其收敛当且仅当$p>1$，利用\textbf{归结原理}与\textbf{等价无穷小替换}得到
        $$\lim_{n\to\infty}\frac{\frac{1}{\ln^pn}\sin\frac{a^2\pi}{\sqrt{n^2+a}+n}}{\frac{1}{n\ln^pn}}=\lim_{x\to+\infty}x\sin\frac{a^2\pi}{\sqrt{x^2+a^2}+x}=\lim_{x\to+\infty}\frac{a^2\pi x}{\sqrt{x^2+a^2}+x}$$
        分子分母同除以$x$后由初等函数\textbf{连续性}与\textbf{复合函数极限}结论即可得到其为
        $$\lim_{x\to+\infty}\frac{a^2\pi}{\sqrt{1+\frac{a^2}{x^2}}+1}=\frac{a^2\pi}{2}$$
        从而根据\textbf{比较判别法}得原级数绝对收敛当且仅当$p>1$。

        \item \textbf{条件收敛判定}
        
        由于$p>1$时已经绝对收敛，我们只考虑$p\in[0,1]$的情况。由于当$n>N$时这是一个交错级数，尝试去除前$N$项后使用\textbf{莱布尼茨判别法}。

        先验证通项绝对值趋于0：通过前一部分的计算结果，利用乘除极限结论即得
        $$\lim_{n\to\infty}\frac{1}{\ln^pn}\sin\frac{a^2\pi}{\sqrt{n^2+a^2}+n}=\lim_{n\to\infty}\frac{\frac{1}{\ln^pn}\sin\frac{a^2\pi}{\sqrt{n^2+a}+n}}{\frac{1}{n\ln^pn}}\cdot\frac{1}{n\ln^pn}=\frac{a^2\pi}{2}\cdot0=0$$
        此外，第一部分的讨论已经说明$n>N$时通项绝对值单调递减，从而利用莱布尼茨判别法可知原级数条件收敛当且仅当$p\in[0,1]$。
    \end{itemize}
}
\end{framed}

\note 利用数项级数去除有限项不影响敛散性，所有判别法都可以在\textbf{去除有限项后使用}。

\note 注意对数列直接的阶估算都是建立在\textbf{函数极限与归结原理}上的，而化为函数极限无法直接操作时就必须对原数列进行化简。

\begin{framed}
\begin{exercise}\label{ln alter}
    判断如下级数是条件收敛、绝对收敛还是发散：
    $$\suminf[n=2]\ln\bigg(1+\frac{(-1)^n}{n^\alpha}\bigg)$$
    这里参数$\alpha>0$。
\end{exercise}
\sol{
    由于$x\to0$时$\ln(1+x)\sim x$，此级数的通项应当与$\frac{(-1)^n}{n^\alpha}$接近，由此我们讨论绝对收敛性与条件收敛性时均试图向其靠拢：
    \begin{itemize}
        \item \textbf{绝对收敛判定}
        
        利用有界量乘无穷小是无穷小可知$\lim_{n\to\infty}\frac{(-1)^n}{n^\alpha}=0$，由此利用\textbf{归结原理}即得到
        $$\lim_{n\to\infty}\frac{\ln\big(1+\frac{(-1)^n}{n^\alpha}\big)}{\frac{(-1)^n}{n^\alpha}}=\lim_{x\to0}\frac{\ln(1+x)}{x}=1$$
        利用数列极限性质可知
        $$\lim_{n\to\infty}\bigg|\frac{\ln\big(1+\frac{(-1)^n}{n^\alpha}\big)}{\frac{(-1)^n}{n^\alpha}}\bigg|=1$$
        而这就代表
        $$\lim_{n\to\infty}\frac{\big|\ln\big(1+\frac{(-1)^n}{n^\alpha}\big)\big|}{\frac{1}{n^\alpha}}=1$$
        由此，利用\textbf{比较判别法}得原级数的通项取绝对值后得到的级数(这已经保证了是\textbf{正项级数})与$\suminf[n=2]\frac{1}{n^\alpha}$敛散性相同。由\textbf{$p$级数}敛散性结论最终得到原级数绝对收敛当且仅当$\alpha>1$。

        \item \textbf{条件收敛判定-引理}
        
        为了进行后续的判定，我们将先证明如下的引理：对级数$\suminf u_n$与\textbf{正项级数}$\suminf v_n$，若
        $$\lim_{n\to\infty}\frac{u_n}{v_n}=c\ne0$$
        则$\suminf u_n$与$\suminf v_n$敛散性相同。

        当$c>0$时，利用极限\textbf{保序性}可知存在$N$使得$n>N$时$\frac{u_n}{v_n}>0$，又由$v_n>0$可知$n>N$时$u_n>N$，从而$u_n$去除前$N$项(这不影响敛散性)后即成为正项级数，利用正项级数比较判别法可知与$v_n$敛散性相同。

        当$c<0$时，设$w_n=-u_n$，则$\lim_{n\to\infty}\frac{w_n}{v_n}=-c>0$，利用前一种情况可知$w_n$与$v_n$敛散性相同，再由数项级数\textbf{基本性质}可知$w_n$与$u_n$敛散性相同，从而$u_n$与$v_n$敛散性相同。

        
        \item \textbf{条件收敛判定-结论}
        
        由于$\alpha>1$时已经绝对收敛，我们只考虑$\alpha\in(0,1]$的情况。

        \note 利用对数函数的性质可知此为交错级数，但计算可发现它的通项\textbf{并不单调递减}，由此无法使用莱布尼茨判别法。
        
        由于此级数通项与$\frac{(-1)^n}{n^\alpha}$接近，我们先考虑$\suminf[n=2]\frac{(-1)^n}{n^\alpha}$的敛散性：这也是一个交错级数，且通项绝对值$\frac{1}{n^\alpha}$单调递减趋于0，利用\textbf{莱布尼茨判别法}可得其收敛。

        利用数项级数的基本性质(收敛与收敛之和收敛、收敛与发散之和发散)，\textbf{原级数通项减去收敛级数的通项并不影响敛散性}，由此原级数敛散性与
        $$\suminf\bigg(\ln\bigg(1+\frac{(-1)^n}{n^\alpha}\bigg)-\frac{(-1)^n}{n^\alpha}\bigg)$$
        相同。
    
        设$f(n)=\frac{(-1)^n}{n^\alpha}$，前一部分已经证明$n\to\infty$时$f(n)\to0$，而上方通项为$\ln(1+f(n))-f(n)$，对其估算必须考虑$\ln$在1处\textbf{展开到二阶}。写出二阶泰勒展开
        $$\ln(1+x)=x-\frac{x^2}{2}+o(x^2)\quad(x\to0)$$
        可以利用$o(x^2)$的定义将其写为函数极限形式
        $$\lim_{x\to0}\frac{\ln(1+x)-x}{x^2}=\lim_{x\to0}\frac{-\frac{x^2}{2}+o(x^2)}{x^2}=-\frac{1}{2}$$
        由此利用\textbf{归结原理}有(注意$1=((-1)^n)^2$)
        $$\lim_{n\to\infty}\frac{\ln\big(1+\frac{(-1)^n}{n^\alpha}\big)-\frac{(-1)^n}{n^\alpha}}{\frac{1}{n^{2\alpha}}}=\lim_{x\to0}\frac{\ln(1+x)-x}{x^2}=-\frac{1}{2}$$
        由引理，原级数与$\suminf[n=2]\frac{1}{n^{2\alpha}}$敛散性相同，从而根据\textbf{$p$级数}敛散性结论可知原级数收敛当且仅当$\alpha>\frac{1}{2}$。

        综合以上，原级数条件收敛当且仅当$\alpha\in(\frac{1}{2},1]$。
    \end{itemize}
}
\end{framed}

\note 仍然注意比较判别法只能对正项级数使用，对于非正项级数，即使通项为等价无穷小也未必敛散性相同。

\note 本题中的引理是\textbf{加强的极限形式比较判别法}，非常建议掌握，考试需要用到时可以写成``利用极限保序性与比较判别法得''。

\begin{framed}
\begin{exercise}
    判断如下级数是条件收敛、绝对收敛还是发散：
    $$\suminf\bigg(\sqrt[n]{\alpha}-\sqrt{1+\frac{1}{n}}\bigg)$$
    这里参数$\alpha>0$。
\end{exercise}
\sol{
    由于此处通项形式较为复杂，我们先进行\textbf{估算}再证明：
    \begin{itemize}
        \item \textbf{估算}
        
        将$\frac{1}{n}$替换为趋于0的$x$，通项即为
        $$\alpha^x-\sqrt{1+x}$$
        由于$x\to0$时$\alpha^x-\sqrt{1+x}=0$，利用\textbf{归结原理}可以得到原级数通项趋于0。为了进一步估算大小，将其在0处展开到一阶，由于
        $$\alpha^x=1+(\ln\alpha)x+o(x),\quad\sqrt{1+x}=1+\frac{x}{2}+o(x)\quad(x\to0)$$
        作差即得
        $$\alpha^x-\sqrt{1+x}=\bigg(\ln\alpha-\frac{1}{2}\bigg)x+o(x)\quad(x\to0)$$
        由此利用$o(x)$定义有
        $$\lim_{x\to0}\frac{\alpha^x-\sqrt{1+x}}{x}=\frac{\big(\ln\alpha-\frac{1}{2}\big)x+o(x)}{x}=\ln\alpha-\frac{1}{2}$$
        上述极限非零当且仅当$\alpha\ne\sqrt\er$，由此进行分类讨论。

        \item \textbf{$\alpha\ne\sqrt\er$情况判定}
        
        此时利用\textbf{归结原理}可知
        $$\lim_{n\to\infty}\frac{\sqrt[n]{\alpha}-\sqrt{1+\frac{1}{n}}}{\frac{1}{n}}=\lim_{x\to0}\frac{\alpha^x-\sqrt{1+x}}{x}=\ln\alpha-\frac{1}{2}\ne0$$
        通过\exref{ln alter}中的引理即得到原级数与$\suminf\frac{1}{n}$敛散性相同，从而根据\textbf{$p$级数}敛散性结论可知原级数发散。

        \item \textbf{$\alpha=\sqrt\er$情况判定}
        
        此时，由于$\ln\alpha-\frac{1}{2}=0$，我们需要进行更进一步的估算。将$\alpha^x=\er^{\frac{x}{2}}$与$\sqrt{1+x}$在0处展开到二阶，得到
        $$\er^{\frac{x}{2}}=1+\frac{x}{2}+\frac{x^2}{8}+o(x^2),\quad\sqrt{1+x}=1+\frac{x}{2}-\frac{x^2}{8}+o(x^2)\quad(x\to0)$$
        作差即得
        $$\er^{\frac{x}{2}}-\sqrt{1+x}=\frac{x^2}{4}+o(x^2)\quad(x\to0)$$
        由此利用$o(x^2)$定义有
        $$\lim_{x\to0}\frac{\er^{\frac{x}{2}}-\sqrt{1+x}}{x^2}=\frac{\frac{x^2}{4}+o(x^2)}{x^2}=\frac{1}{4}$$
        利用\textbf{归结原理}可知此时
        $$\lim_{n\to\infty}\frac{\sqrt[n]{\alpha}-\sqrt{1+\frac{1}{n}}}{\frac{1}{n^2}}=\lim_{x\to0}\frac{\er^{\frac{x}{2}}-\sqrt{1+x}}{x^2}=\frac{1}{4}$$
        通过\exref{ln alter}中的引理即得到原级数与$\suminf\frac{1}{n^2}$敛散性相同，从而根据\textbf{$p$级数}敛散性结论可知原级数收敛。进一步地，根据引理的证明过程，$n$充分大时原级数为正项级数，因此从收敛可以推出绝对收敛。
    \end{itemize}
}
\end{framed}

\note 仍然注意\textbf{等价无穷小只能在乘除中替换}，无论加减还是指数都无法直接替换，一般情况需要泰勒展开求解。若需要用到泰勒展开而没有记住，可以直接利用定义\textbf{求导计算系数}。

\begin{framed}
\begin{exercise}
    判断如下级数是条件收敛、绝对收敛还是发散：
    $$\suminf\frac{\sin(n\varphi)}{(n+\frac{1}{n})^p}\bigg(1+\frac{1}{n}\bigg)^n$$
    这里参数$\varphi\in(0,\pi)$、$p>0$。
\end{exercise}
\sol{
    本题中的级数通项与$\er\frac{\sin(n\varphi)}{n^p}$接近，由此我们讨论绝对收敛性与条件收敛性时均试图向其靠拢：
    \begin{itemize}
        \item \textbf{绝对收敛判定-化简}
        
        其通项取绝对值得到的级数为
        $$\suminf\frac{|\sin(n\varphi)|}{(n+\frac{1}{n})^p}\bigg(1+\frac{1}{n}\bigg)^n$$
        我们先证明它的收敛性与
        $$\suminf\frac{|\sin(n\varphi)|}{n^p}$$
        相同。将两个级数分别记作$\suminf u_n$、$\suminf v_n$。

        直接计算可知
        $$\lim_{n\to\infty}\frac{\frac{1}{(n+\frac{1}{n})^p}(1+\frac{1}{n})^n}{\frac{1}{n^p}}=\lim_{n\to\infty}\frac{n^p}{(n+\frac{1}{n})^p}\cdot\bigg(1+\frac{1}{n}\bigg)^n=\lim_{n\to\infty}\frac{1}{(1+\frac{1}{n^2})^p}\cdot\bigg(1+\frac{1}{n}\bigg)^n$$
        利用$\er$的定义可知第二部分极限为$\er$，利用\textbf{归结原理}与初等函数\textbf{连续性}可知第一部分极限为
        $$\lim_{x\to0}\frac{1}{(1+x)^p}=1$$
        从而由乘积极限结论得
        $$\lim_{n\to\infty}\frac{\frac{1}{(n+\frac{1}{n})^p}(1+\frac{1}{n})^n}{\frac{1}{n^p}}=\er$$
        利用极限的\textbf{保序性}，存在$N$使得$n>N$时
        $$1<\frac{\frac{1}{(n+\frac{1}{n})^p}(1+\frac{1}{n})^n}{\frac{1}{n^p}}<3$$
        也即
        $$\frac{1}{n^p}<\frac{1}{(n+\frac{1}{n})^p}\bigg(1+\frac{1}{n}\bigg)^n<3\cdot\frac{1}{n^p}$$
        两边同乘非负的$|\sin(n\varphi)|$即得到$n>N$时
        $$v_n\le u_n\le 3v_n$$
        由于$\suminf u_n$、$\suminf v_n$均为正项级数，利用\textbf{比较判别法}，从$\suminf v_n$发散可得$\suminf u_n$发散；另一方面，从$\suminf v_n$收敛利用数项级数\textbf{基本性质}可得$\suminf(3v_n)$收敛，于是再次利用\textbf{比较判别法}得$\suminf u_n$收敛。这就证明了二者敛散性相同。

        \note 注意由于$\sin(n\varphi)$可能为0，$\frac{u_n}{v_n}$的极限是\textbf{无法讨论}的(除非进行等价无穷小的\textbf{推广}，见上学期讲义15.2.2，但结论无法直接使用)，因此必须采用这样的方式。

        \item \textbf{绝对收敛判定-结论}
        
        通过上一部分的讨论，我们最终证明了原级数是否绝对收敛等价于
        $$\suminf\frac{|\sin(n\varphi)|}{n^p}$$
        是否收敛。当$p>1$时，由
        $$\frac{|\sin(n\varphi)|}{n^p}\le\frac{1}{n^p}$$
        通过\textbf{$p$级数}的敛散性结论与比较判别法可知$\suminf\frac{|\sin(n\varphi)|}{n^p}$收敛，从而原级数收敛。

        当$p\in(0,1]$时，类似本讲义28.1.1习题3可证明$\suminf\frac{|\sin(n\varphi)|}{n^p}$发散，从而原级数不绝对收敛。我们这里采用\textbf{直接估算}方法，可以自行仿照本讲义28.1.1习题3书写\textbf{拆分级数}方法。

        我们采用一个可以称为\textbf{猎人兔子原理}的思路进行估算(可见上学期讲义2.3.1)：给定$\delta\in(0,\frac{\pi}{2})$，我们将所有如下区间
        $$(k\pi-\delta,k\pi+\delta),\quad k\in\mathbb{Z}$$
        的并集记为集合$A$，利用三角函数的性质可得$x\in A$当且仅当$|\sin x|<\sin\delta$。

        可以发现，集合$A$中的两个点若属于同一个区间，则距离一定小于$2\delta$，而若属于不同区间，距离一定大于$\pi-2\delta$。由此，我们取$\delta$使得$2\delta<\varphi<\pi-2\delta$，也即
        $$0<\delta<\min\bigg\{\frac{\pi-\varphi}{2},\frac{\varphi}{2}\bigg\}$$
        由$\varphi$的范围，这样的$\delta$一定存在。此时，由于$n\varphi$与$(n+1)\varphi$的差恰好为$\varphi$，根据上方的讨论，它们不可能同时落在集合$A$中，于是必有一个的正弦绝对值大于等于$\sin\delta$。综合以上的讨论，我们找到了一个$\delta\in(0,\frac{\pi}{2})$使得对任何正整数$n$都有
        $$\max\{|\sin(n\varphi)|,|\sin((n+1)\varphi)|\}\ge\sin\delta$$

        接下来，将$\frac{|\sin(n\varphi)|}{n^p}$两两分组，得到(无论$|\sin((2n-1)\varphi)|\ge\sin\delta$还是$|\sin(2n\varphi)|\ge\sin\delta$都能如此放缩)
        $$\sum_{n=1}^{2M}\frac{|\sin(n\varphi)|}{n^p}=\sum_{n=1}^M\bigg(\frac{|\sin((2n-1)\varphi)|}{(2n-1)^p}+\frac{|\sin(2n\varphi)|}{(2n)^p}\bigg)\ge\sum_{n=1}^M\frac{\sin\delta}{(2n)^p}=\frac{\sin\delta}{2^p}\sum_{n=1}^M\frac{1}{n^p}$$
        利用\textbf{$p$级数}敛散性结论可知
        $$\lim_{M\to\infty}\frac{\sin\delta}{2^p}\sum_{n=1}^M\frac{1}{n^p}=+\infty$$
        从而根据\textbf{推广的夹逼定理}
        $$\lim_{M\to\infty}\sum_{n=1}^{2M}\frac{|\sin(n\varphi)|}{n^p}=+\infty$$
        而这是$\suminf\frac{|\sin(n\varphi)|}{n^p}$的一个\textbf{子列极限}，由子列极限不存在即可知原数列极限不存在，最终得到发散。

        综合以上讨论，原级数绝对收敛当且仅当$p>1$。

        \item \textbf{条件收敛判定-化简}
        
        由于$p>1$时已经绝对收敛，我们只考虑$p\in(0,1]$的情况。
        
        由于原级数的通项与
        $$\suminf\frac{\sin(n\varphi)}{(n+\frac{1}{n})^p}\er$$
        接近，我们先证明此级数收敛。

        利用等差数列的正弦求和公式(可见上学期讲义7.2.1)得
        $$\sum_{n=1}^m\sin(n\varphi)=\frac{\sin(\frac{n+1}{2}\varphi)\sin\frac{n\varphi}{2}}{\sin\frac{\varphi}{2}}$$
        由此可进一步得到(由条件$\sin\frac{\varphi}{2}>0$)
        $$\bigg|\sum_{n=1}^m\sin(n\varphi)\bigg|\le\frac{1}{\sin\frac{\varphi}{2}}$$
        将级数$\suminf\frac{\sin(n\varphi)}{(n+\frac{1}{n})^p}\er$写为
        $$\suminf\frac{\er}{(n+\frac{1}{n})^p}\cdot\sin(n\varphi)$$
        之前的推导已证明$\cos(n\varphi)$的部分和有界，下面证明$\frac{\er}{(n+\frac{1}{n})^p}$随$n$增加单调递减趋于0。
        
        单调递减：对正整数$n$，作差可知
        $$n+1+\frac{1}{n+1}-n-\frac{1}{n}=1-\frac{1}{n(n+1)}>0$$
        由此$n+\frac{1}{n}$随$n$增加单调递增，利用幂函数单调性即得$\frac{\er}{(n+\frac{1}{n})^p}$单调递减。

        极限为0：直接放缩有
        $$0<\frac{\er}{(n+\frac{1}{n})^p}<\frac{\er}{n^p}$$
        由于左右极限均为0，利用\textbf{夹逼定理}可知中间极限也为0。

        综合以上，利用\textbf{迪利克雷判别法}得$\suminf\frac{\sin(n\varphi)}{(n+\frac{1}{n})^p}\er$收敛。

        \note 证明这部分收敛后也可以利用\textbf{阿贝尔判别法}由$(1+\frac{1}{n})^n$单调有界直接得到结论。这里介绍更为通用的\textbf{作差估算}方法。

        \item \textbf{条件收敛判定-估算}
        
        利用数项级数的基本性质(收敛与收敛之和收敛、收敛与发散之和发散)，\textbf{原级数通项减去收敛级数的通项并不影响敛散性}，由此根据上一部分讨论，原级数敛散性与
        $$\suminf\frac{\sin(n\varphi)}{(n+\frac{1}{n})^p}\bigg(\bigg(1+\frac{1}{n}\bigg)^n-\er\bigg)$$
        相同。

        为了进行估算，我们先考虑$\big(1+\frac{1}{n}\big)^n-\er$。将$\frac{1}{n}$替换为趋于0的$x$，此式即为
        $$(1+x)^{\frac{1}{x}}-\er=\er^{\frac{\ln(1+x)}{x}}-\er=\er\big(\er^{\frac{\ln(1+x)-x}{x}}-1\big)$$
        利用基本等价无穷小结论可知$x\to0$时$\ln(1+x)\sim x$，从而此时$\frac{\ln(1+x)-x}{x}\to0$，则利用$t\to0$时$\er^t-1\sim t$可进行\textbf{等价无穷小替换}得到
        $$(1+x)^{\frac{1}{x}}-\er\sim\er\frac{\ln(1+x)-x}{x}\quad(x\to0)$$
        为了进一步计算，将$\ln(1+x)$在0处展开到二阶可得
        $$\ln(1+x)=x-\frac{x^2}{2}+o(x^2)\quad(x\to0)$$
        从而(涉及的$o$运算性质可见上学期讲义11.3.2)
        $$\er\frac{\ln(1+x)-x}{x}=\er\frac{-\frac{x^2}{2}+o(x^2)}{x}=-\frac{\er}{2}x+o(x)\quad(x\to0)$$
        由此利用$o(x)$的定义可得到
        $$\er\frac{\ln(1+x)-x}{x}\sim-\frac{\er}{2}x\quad(x\to0)$$
        也即最终得到
        $$(1+x)^{\frac{1}{x}}-\er\sim-\frac{\er}{2}x\quad(x\to0)$$
        将其写为
        $$\lim_{x\to0}\frac{(1+x)^{\frac{1}{x}}-\er}{x}=-\frac{\er}{2}$$
        利用\textbf{归结原理}可以得到
        $$\lim_{n\to\infty}\frac{(1+\frac{1}{n})^n-\er}{\frac{1}{n}}=\lim_{x\to0}\frac{(1+x)^{\frac{1}{x}}-\er}{x}=-\frac{\er}{2}$$
        而第一部分证明中已经证明了
        $$\lim_{n\to\infty}\frac{\frac{1}{(n+\frac{1}{n})^p}}{\frac{1}{n^p}}=1$$
        综合得到
        $$\lim_{n\to\infty}\frac{\frac{(1+\frac{1}{n})^n-\er}{(n+\frac{1}{n})^p}}{\frac{1}{n^{1+p}}}=-\frac{\er}{2}\cdot1=-\frac{\er}{2}$$

        \item \textbf{条件收敛判定-结论}
        
        利用\exref{ln alter}中的引理，$\suminf\frac{(1+\frac{1}{n})^n-\er}{(n+\frac{1}{n})^p}$与$\suminf\frac{1}{n^{1+p}}$敛散性相同，由条件$p>0$，从而根据\textbf{$p$级数}敛散性结论可知$\suminf\frac{(1+\frac{1}{n})^n-\er}{(n+\frac{1}{n})^p}$收敛。

        利用数列极限结论可知$(1+\frac{1}{n})^n$单调递增趋于$\er$，从而其恒小于等于$\er$，也即
        $$\bigg|\suminf\frac{\sin(n\varphi)}{(n+\frac{1}{n})^p}\bigg(\bigg(1+\frac{1}{n}\bigg)^n-\er\bigg)\bigg|\le\bigg|\frac{(1+\frac{1}{n})^n-\er}{(n+\frac{1}{n})^p}\bigg|=-\frac{(1+\frac{1}{n})^n-\er}{(n+\frac{1}{n})^p}$$
        利用数项级数的基本性质，右侧级数收敛，因此左侧级数\textbf{绝对收敛}，于是它收敛，这就得到了原级数在$p>0$时恒收敛。

        综合以上讨论，原级数条件收敛当且仅当$p\in(0,1]$。
    \end{itemize}
}
\end{framed}

\note 建议熟悉\textbf{等差数列的正弦求和公式}，它常用于说明一些正弦/余弦函数相关的级数求和的有界性。

\note 本题虽然过程细节相对复杂，但思路仍然是之前题目思路的\textbf{直接组合}，建议仔细阅读过程。对于一般项级数，判别的顺序为：先判断是否\textbf{绝对收敛}，绝对收敛已经可以推出收敛；再尝试利用\textbf{迪利克雷判别法}或\textbf{阿贝尔判别法}进行拆分；若仍然无法做出，则需要考虑更复杂的技巧，例如\textbf{分部分估算}；可以将\textbf{柯西收敛原理}作为最终手段。

\subsubsection{函数项级数}
\begin{framed}
\begin{exercise}\label{func uni}
    判断函数列
    $$f_n(x)=\bigg(1-\frac{1}{\sqrt{n}}\bigg)^{x^2}$$
    在$(0,10]$、$(0,+\infty)$是否一致收敛到某个$f$。
\end{exercise}
\sol{
    利用\textbf{归结原理}与\textbf{初等函数连续性}直接计算可知$x>0$时
    $$\lim_{n\to\infty}\bigg(1-\frac{1}{\sqrt{n}}\bigg)^{x^2}=\lim_{t\to0}(1-t)^{x^2}=1$$
    由此利用\textbf{一致收敛结果必为逐点收敛结果}可知若一致收敛，只能收敛到
    $$f(x)=1$$

    利用定义，一致收敛成立当且仅当
    $$\forall\varepsilon>0,\quad\exists N\in\mathbb{N},\quad\forall x\in(0,+\infty),\quad\forall n>N,\quad\bigg|\bigg(1-\frac{1}{\sqrt{n}}\bigg)^{x^2}-1\bigg|<\varepsilon$$

    为了能够估算，我们固定$n$，考虑$x$变化时$\big|\big(1-\frac{1}{\sqrt{n}}\big)^{x^2}-1\big|$的尽可能大取值。下面分别讨论两种情况：
    \begin{itemize}
        \item \textbf{有界区间情况}
        
        由于$1-\frac{1}{\sqrt{n}}<1$，利用指数函数的\textbf{单调性}可发现$x\in(0,10]$时
        $$\bigg|\bigg(1-\frac{1}{\sqrt{n}}\bigg)^{x^2}-1\bigg|=1-\bigg(1-\frac{1}{\sqrt{n}}\bigg)^{x^2}\le1-\bigg(1-\frac{1}{\sqrt{n}}\bigg)^{100}$$
        进一步利用\textbf{归结原理}与\textbf{初等函数连续性}可发现
        $$\lim_{n\to\infty}1-\bigg(1-\frac{1}{\sqrt{n}}\bigg)^{100}=\lim_{t\to0}1-(1-t)^{100}=0$$
        下面证明这已经足以说明一致收敛。

        对任何$\varepsilon>0$，由已算出的极限可取$N$使得$n>N$时
        $$\bigg|1-\bigg(1-\frac{1}{\sqrt{n}}\bigg)^{100}\bigg|<\varepsilon$$
        根据上方讨论可知此时对任何$x\in(0,10]$有
        $$\bigg|\bigg(1-\frac{1}{\sqrt{n}}\bigg)^{x^2}-1\bigg|\le1-\bigg(1-\frac{1}{\sqrt{n}}\bigg)^{100}<\varepsilon$$
        从而已经符合定义。

        \item \textbf{无界区间情况}
        
        直接计算可发现由于$1-\frac{1}{\sqrt{n}}<1$，对正整数$n$有
        $$\lim_{x\to+\infty}\bigg|\bigg(1-\frac{1}{\sqrt{n}}\bigg)^{x^2}-1\bigg|=\lim_{x\to+\infty}1-\bigg(1-\frac{1}{\sqrt{n}}\bigg)^{x^2}=1$$
        从而利用极限的\textbf{保序性}，对任何正整数$n$都存在$x>0$使得$|f_n(x)-f(x)|>\frac{1}{2}$。
        
        这已经符合一致收敛的\textbf{否定}：
        $$\exists\varepsilon>0,\quad\forall N\in\mathbb{N},\quad\exists x\in(0,+\infty),\quad\exists n>N,\quad\bigg|\bigg(1-\frac{1}{\sqrt{n}}\bigg)^{x^2}-1\bigg|\ge\varepsilon$$
        具体来说，取$\varepsilon=\frac{1}{2}$，对任何$N\in\mathbb{N}$，取$n=N+1$，并取出对应的使得$|f_n(x)-f(x)|>\frac{1}{2}$的$x$，对这个点即有
        $$|f_n(x)-f(x)|>\frac{1}{2}=\varepsilon$$
        由此函数列在$(0,+\infty)$不一致收敛。
    \end{itemize}
}
\end{framed}

\note 函数列$f_n(x)$在集合$D$上一致收敛于$f(x)$当且仅当$|f_n(x)-f(x)|$在$D$上的\textbf{上确界}(含义为最小上界，可以粗略理解为最大值)随$n\to\infty$而趋于0。由此，对于能够计算出$f_n$与$f$的问题，可以直接\textbf{估算上确界}以判断是否一致收敛。更具体来说，若$f_n(x)-f(x)$的某个\textbf{上界}(注意\textbf{最大值}也是上界)在$n\to\infty$时趋于0，即可推出$f_n(x)$一致收敛于$f(x)$，此结论可直接使用。

\begin{framed}
\begin{exercise}
    求如下关于$x$的函数项级数的收敛域，并判断每点是条件收敛还是绝对收敛：
    $$\suminf(-1)^n\frac{1}{n^x+n}$$
\end{exercise}
\sol{
    由于对正整数$n$与实数$x$有$n^x>0$，此函数项级数在任何点处均为\textbf{交错级数}。分为如下步骤：
    \begin{itemize}
        \item \textbf{绝对收敛判定}
        
        其通项取绝对值得到的级数为
        $$\suminf\frac{1}{n^x+n}$$
        由于$x$已经给定，我们可以利用类似\textbf{多项式阶估算}(可见上学期讲义2.4.1)的方式处理。

        具体来说，分为三种情况讨论：
        
        若$x<1$，直接计算有(倒数第二步用了四则运算的极限结论)
        $$\lim_{n\to\infty}\frac{\frac{1}{n^x+n}}{\frac{1}{n}}=\lim_{n\to\infty}\frac{1}{n^{x-1}+1}=\frac{1}{0+1}=1$$
        于是此级数与$\suminf\frac{1}{n}$敛散性相同，从而根据\textbf{$p$级数}敛散性结论可知其发散。

        若$x=1$，此时此级数为$\suminf\frac{1}{2n}$，利用\textbf{$p$级数}敛散性结论与数项级数基本性质可知其发散。

        若$x>1$，直接计算有(倒数第二步用了四则运算的极限结论)
        $$\lim_{n\to\infty}\frac{\frac{1}{n^x+n}}{\frac{1}{n^x}}=\lim_{n\to\infty}\frac{1}{1+n^{1-x}}=\frac{1}{1+0}=1$$
        于是此级数与$\suminf\frac{1}{n^x}$敛散性相同，从而根据\textbf{$p$级数}敛散性结论可知其收敛。

        综合以上，原级数绝对收敛当且仅当$x>1$。

        \item \textbf{条件收敛判定}
        
        由于$x>1$时已经绝对收敛，我们只考虑$x\in(-\infty,1]$的情况，由于此为交错级数尝试使用\textbf{莱布尼茨判别法}。

        先验证通项绝对值趋于0：由于
        $$0<\frac{1}{n^x+n}<\frac{1}{n}$$
        且右侧在$n\to\infty$时极限为0，利用\textbf{夹逼定理}可得中间极限也为0。

        再验证通项绝对值单调递减：相邻两项分母作差得
        $$(n+1)^x+(n+1)-(n^x+n)=1+(n+1)^x-n^x$$
        当$x\in[0,1]$时，$(n+1)^x-n^x>0$，从而分母单调增，因此通项绝对值单调减；当$x<0$时，$(n+1)^x-n^x>-n^x\ge-1$，仍有分母单调增、通项绝对值单调减。

        综合以上，利用莱布尼茨判别法可知原级数在$x\in(-\infty,1]$均收敛。
        
        \item \textbf{综合结果}
        
        根据以上讨论，原级数收敛域为$\rb$，其中在$(1,+\infty)$绝对收敛，$(-\infty,1]$条件收敛。
    \end{itemize}
}
\end{framed}

\note 由于只关心函数项级数的\textbf{逐点}性质，这本质是一个\textbf{数项级数问题}，由此采用通常的先判断绝对收敛、再判断条件收敛的顺序(即使这与问题中先问收敛域的顺序不一致)。

\begin{framed}
\begin{exercise}
    求如下关于$x$的函数项级数的收敛域，并判断每点是条件收敛还是绝对收敛：
    $$\suminf\frac{2^n\sin^nx}{1+4^n\sin^{2n}x}$$
\end{exercise}
\sol{
    由于我们只关心此级数的逐点性质，可以设$t=2\sin x$，则原级数化为新的级数
    $$\suminf\frac{t^n}{1+t^{2n}},\quad t\in[-2,2]$$
    只要知道新的级数对怎样的$t$收敛，即可得到原级数对怎样的$x$收敛。下面均考虑新的级数。分为如下步骤：
    \begin{itemize}
        \item \textbf{绝对收敛判定}
        
        由于分母恒为正，其通项取绝对值得到的级数为
        $$\suminf\frac{|t|^n}{1+t^{2n}}$$
        当$|t|=1$时，此为$\suminf\frac{1}{2}$发散，其他情况可以想到与适当的\textbf{等比级数}比较：

        若$|t|<1$，有
        $$\frac{|t|^n}{1+t^{2n}}<|t|^n$$
        根据等比级数敛散性结论可知右侧部分和收敛，从而利用\textbf{比较判别法}得此时新级数绝对收敛。

        若$|t|>1$，有
        $$\frac{|t|^n}{1+t^{2n}}<\frac{|t|^n}{t^{2n}}=\frac{1}{|t|^n}=\bigg(\frac{1}{|t|}\bigg)^n$$
        根据等比级数敛散性结论可知右侧部分和收敛，从而利用\textbf{比较判别法}得此时新级数绝对收敛。

        综合以上，新级数绝对收敛当且仅当$|t|\ne1$。

        \item \textbf{条件收敛判定}
        
        由于$|t|\ne1$时已经绝对收敛，我们只考虑$t=\pm1$的情况。$t=1$时新级数为$\suminf\frac{1}{2}$，$t=-1$时新级数为$\suminf\frac{(-1)^n}{2}$，两种情况均可由\textbf{通项极限不为0}得到发散。

        \item \textbf{综合结果}
        
        利用三角函数性质可知$|t|=1$当且仅当
        $$x=k\pi\pm\frac{\pi}{6},\quad k\in\mathbb{Z}$$
        由此，原级数的收敛域为
        $$\bigg\{x\in\rb\ \bigg|\ x\ne k\pi\pm\frac{\pi}{6},\quad k\in\mathbb{Z}\bigg\}$$
        在收敛域中处处绝对收敛。
    \end{itemize}
}
\end{framed}

\note 注意对于\textbf{一致收敛问题不能换元后判断}，但逐点收敛问题并没有影响。

\begin{framed}
\begin{exercise}\label{not uni}
    判断如下关于$x$的函数项级数在$[1,2]$、$[0,1]$是否一致收敛：
    $$\suminf\frac{\sin(nx)}{n^p}$$
    这里参数$p>0$。
\end{exercise}
\sol{
    分两个区间判断：
    \begin{itemize}
        \item \textbf{不含0的区间}
        
        由于本身形式可以拆分为
        $$\suminf\frac{1}{n^p}\cdot\sin(nx)$$
        我们尝试证明第一部分通项一致收敛于0，第二部分部分和一致有界。

        第一部分：由于其不含$x$，给定$n$后，$x\in[1,2]$时其恒等于$\frac{1}{n^p}$，最大值为$\frac{1}{n^p}$。由$p>0$可知
        $$\lim_{n\to\infty}\bigg|\frac{1}{n^p}-0\bigg|=\lim_{n\to\infty}\frac{1}{n^p}=0$$
        根据\exref{func uni}后注释即得第一部分通项一致收敛于0。

        第二部分：利用\textbf{等差数列的正弦求和公式}(可见上学期讲义7.2.1)得
        $$\sum_{n=1}^m\sin(nx)=\frac{\sin(\frac{m+1}{2}x)\sin\frac{mx}{2}}{\sin\frac{x}{2}}$$
        利用三角函数性质可发现$x\in[1,2]$时分母绝对值最小为$\sin\frac{1}{2}$，由此对任意正整数$m$有
        $$\forall x\in[1,2],\quad\bigg|\sum_{n=1}^m\sin(nx)\bigg|=\bigg|\frac{\sin(\frac{m+1}{2}x)\sin\frac{mx}{2}}{\sin\frac{x}{2}}\bigg|\le\frac{1}{\big|\sin\frac{x}{2}\big|}\le\frac{1}{\sin\frac{1}{2}}$$
        这已经符合一致有界的定义。

        综合以上，利用\textbf{迪利克雷判别法}得原级数在$[1,2]$一致收敛。
        
        \item \textbf{含0的区间-初步分析}

        当$p>1$时，由于
        $$\forall x\in[0,1],\quad\bigg|\frac{\sin(nx)}{n^p}\bigg|<\frac{1}{n^p}$$
        根据\textbf{$p$级数}敛散性结论可知右侧部分和收敛，从而利用\textbf{强级数判别法}得左侧部分和一致收敛。

        当$p\le1$时，由于$[1,2]$一致收敛性证明中第二部分的一致有界性无法再成立，需要尝试直接通过\textbf{柯西收敛原理}判断是否一致收敛。

        根据柯西收敛原理，此函数项级数在$[0,1]$一致收敛当且仅当
        $$\forall\varepsilon>0,\quad\exists N,\quad\forall m>n>N,\quad\forall x\in[0,1],\quad\bigg|\sum_{k=n+1}^m\frac{\sin(kx)}{k^p}\bigg|<\varepsilon$$

        \item \textbf{含0的区间-估算}
        
        \note 本部分\textbf{仅为分析过程}，实际考试并不需要写出，但希望大家能掌握分析的思路。

        由于$x\to0$时$\sin x\sim x$，当$kx$不太大时(如$0<kx<1$)，应当存在$c>0$使得
        $$cx\le\sin(kx)\le x$$
        由此，我们可以在$x\le\frac{1}{m}$时将要估算的式子看作(注意$k$为正整数、$1-p\ge0$)
        $$\sum_{k=n+1}^m\frac{kx}{k^p}=\sum_{k=n+1}^mk^{1-p}x\ge(m-n)x$$
        可以发现，无论$n$多大，取$m=2n$、$x=\frac{1}{m}$，均可以得
        $$\sum_{k=n+1}^m\frac{\sin(kx)}{k^p}\ge c\sum_{k=n+1}^m\frac{kx}{k^p}\ge(m-n)x=\frac{c}{2}$$
        而这已经可以证明柯西收敛原理的否定。

        \item \textbf{含0的区间-引理}
        
        我们先严谨证明上一部分中提到的引理：存在$c>0$使得
        $$\forall\varphi\in[0,1],\quad\sin\varphi\ge c\varphi$$

        设$f(x)=\frac{\sin x}{x}$，利用基本等价无穷小，补充定义$f(0)=1$后$f(x)$即成为了$[0,1]$上的连续函数，从而利用闭区间连续函数的性质，它在$[0,1]$上存在\textbf{最小值}。

        直接计算可知$(0,1)$上有(利用$(0,1)$上$x$、$\sin x$、$\cos x$均正与三角函数的性质)
        $$f'(x)=\frac{x\cos x-\sin x}{x^2}=\frac{x-\tan x}{x^2\cos x}<0$$
        由此$f(x)$在$[0,1]$单调减，从而其最小值在$x=1$取到，由此即证明了取$c=\sin1$已经符合要求。

        \item \textbf{含0的区间-结论}
        
        我们现在来利用柯西收敛原理证明当$p\in(0,1]$时原级数在$[0,1]$不一致收敛。这也即要证明
        $$\exists\varepsilon>0,\quad\forall N,\quad\exists m>n>N,\quad\exists x\in[0,1],\quad\bigg|\sum_{k=n+1}^m\frac{\sin(kx)}{k^p}\bigg|\ge\varepsilon$$
        设$c$的定义与前一部分证明相同。取$\varepsilon=\frac{c}{2}$，对任何$N$，考虑
        $$n=N+1,\quad m=2N+2,\quad x=\frac{1}{2N+2}$$
        此时，对任何$n+1\le k\le m$均有$kx\in[0,1]$，因此(第二个不等号是由于$k^{1-p}\ge1$)
        $$\bigg|\sum_{k=n+1}^m\frac{\sin(kx)}{k^p}\bigg|=\sum_{k=n+1}^m\frac{\sin(kx)}{k^p}\ge\sum_{k=n+1}^m\frac{ckx}{k^p}=\sum_{k=n+1}^mck^{1-p}x\ge(m-n)cx=\frac{c}{2}=\varepsilon$$
        这已经符合一致收敛的否定。

        综合以上，原级数在$[0,1]$一致收敛当且仅当$p>1$。

        \note 教材10.4节例9后也给出了$p=1$时原级数在$[0,1]$上不一致收敛的\textbf{估算证明思路}，但由于缺乏构造过程，教材中看似更简洁的做法是更难\textbf{想到}的。
    \end{itemize}
}
\end{framed}

\note 函数项级数一致收敛的\textbf{基本判别流程}与数项级数是一致的：先看是否能\textbf{直接计算部分和}后用定义判断，然后考虑比较判别法与绝对一致收敛性判定合并而成的\textbf{强级数判别法}，若无法使用则尝试\textbf{迪利克雷判别法}或\textbf{阿贝尔判别法}，还是无法进行则考虑使用\textbf{柯西收敛原理}直接估算。证明不一致收敛的\textbf{最基本方法}即是通过柯西收敛原理，因为它是\textbf{充要}的。

\begin{framed}
\begin{exercise}
    证明如下关于$x$的函数项级数
    $$\sum_{n=1}^\infty(-1)^nn^2\er^{-nx}$$
    对任何正数$r$在$[r,+\infty)$一致收敛，在$(0,+\infty)$不一致收敛。
\end{exercise}
\sol{
    分两种情况判断：
    \begin{itemize}
        \item \textbf{一致收敛证明}
        
        由于$r>0$，利用指数函数的单调性，对$x\in[r,+\infty)$有
        $$|(-1)^nn^2\er^{-nx}|\le n^2\er^{-nr}$$
        由于右侧为指数下降，阶估算可知其求和应当收敛，从而尝试运用\textbf{强级数判别法}。

        考虑$\suminf\frac{1}{n^2}$，利用\textbf{$p$级数}敛散性结论可知其收敛，利用\textbf{数列极限阶估算}结论(可见上学期讲义2.4.1)得
        $$\lim_{n\to\infty}\frac{n^2\er^{-rn}}{\frac{1}{n^2}}=\lim_{n\to\infty}\frac{n^4}{(\er^r)^n}=0$$
        从而根据\textbf{比较判别法}得$\suminf n^2\er^{-nr}$也收敛，最终通过强级数判别法得到原级数在$[r,+\infty)$一致收敛。

        \item \textbf{不一致收敛证明}
        
        由于$x\to0$时$n^2\er^{-nx}\to n^2$，在$(0,+\infty)$上$|(-1)^nn^2\er^{-nx}|$可任意接近$n^2$，由此可以直接估算：

        对任意$N$，取$n=N+1$、$m=n+1$\ (相当于$n+1$到$m$求和只保留一项)，则
        $$\bigg|\sum_{k=n+1}^m(-1)^kk^2\er^{-kx}\bigg|=m^2\er^{-mx}$$
        取$x=\frac{1}{m}$即有
        $$\bigg|\sum_{k=n+1}^m(-1)^kk^2\er^{-kx}\bigg|=\frac{m^2}{\er}\ge\frac{1}{\er}$$
        因此取$\varepsilon=\frac{1}{\er}$已经符合一致收敛的\textbf{柯西收敛原理}的否定(可见\exref{not uni})，这就证明了不一致收敛。
    \end{itemize}
}
\end{framed}

\note 注意在任意闭区间一致收敛\textbf{无法推进到开区间的边界}，由此内闭一致收敛才有讨论的意义。

\note 虽然用柯西收敛原理证明不一致收敛的说明相对复杂，但这是\textbf{必须熟悉}的基本方式。总的来说，可以先对自然数$m$、$n$估算求和$|\sum_{k=n+1}^mf_k(x)|$，再考虑如何取$x$使得$m$、$n$充分大时能使其大于等于某非零常数，最后再将常数记为$\varepsilon$。

\begin{framed}
\begin{exercise}
    判断以下三个关于$x$的函数项级数在$[0,1]$是否一致收敛：
    $$\suminf n(1-x)x^n,\quad\suminf n(1-x)x^{n^2},\quad\suminf n(1-x)x^{n^3}$$
\end{exercise}
\sol{
    我们按照从简单到难的顺序进行判断：
    \begin{itemize}
        \item \textbf{第一个级数}
        
        直接裂项相消计算可知
        $$\sum_{n=1}^mn(1-x)x^n=(x-x^2)+(2x^2-2x^3)+\dots+(mx^m-mx^{m+1})=x+x^2+\dots+x^m-mx^{m+1}$$
        由此，利用等比数列求和公式进一步计算得到
        $$\sum_{n=1}^mn(1-x)x^n=\begin{cases}\frac{x-x^{m+1}}{1-x}-mx^{m+1}&x\in[0,1)\\0&x=1\end{cases}$$
        令$m$趋于无穷，利用\textbf{数列极限阶估算}结论(可见上学期讲义2.4.1)与\textbf{四则运算的极限}结论，$x\in[0,1)$时
        $$\lim_{m\to\infty}\frac{x-x^{m+1}}{1-x}-mx^{m+1}=\frac{x-0}{1-x}-0=\frac{x}{1-x}$$
        由此得到
        $$\suminf n(1-x)x^n=\begin{cases}\frac{x}{1-x}&x\in[0,1)\\0&x=1\end{cases}$$
        直接计算可发现
        $$\lim_{x\to1^-}\frac{x}{1-x}=+\infty\ne0$$
        从而$\suminf n(1-x)x^n$在$[0,1]$不连续。
        
        若原级数在$[0,1]$一致收敛，由于通项均为连续函数，应当有\textbf{和函数连续}，矛盾。这就证明了原级数在$[0,1]$不一致收敛。

        \item \textbf{第三个级数}
        
        由于难以直接看出是否一致收敛，我们尝试先估算\textbf{最大值}以判断能否使用强级数判别法。设
        $$f_n(x)=n(1-x)x^{n^3}$$
        直接计算可知
        $$f_n(0)=f_n(1)=0$$
        $$\forall x\in(0,1),\quad f_n'(x)=n(n^3x^{n^3-1}-(n^3+1)x^{n^3})$$
        由此比较0、1与导函数为0的点可发现$f_n(x)$的最大值在$x=\frac{n^3}{n^3+1}$取到，为
        $$\frac{n}{n^3+1}\bigg(\frac{n^3}{n^3+1}\bigg)^{n^3}$$
        利用数列极限知识可知右侧部分极限为$\frac{1}{\er}$，以此可以估算，不过我们完全可以进行更直接的\textbf{放缩}得到
        $$\frac{n}{n^3+1}\bigg(\frac{n^3}{n^3+1}\bigg)^{n^3}<\frac{n}{n^3+1}<\frac{1}{n^2}$$
        另一方面，计算得$f_n(x)$在$[0,1]$恒非负，这就得到了
        $$\forall x\in[0,1],\quad |f_n(x)|=f_n(x)\le\frac{1}{n^2}$$
        根据\textbf{$p$级数}敛散性结论可知右侧部分和收敛，从而利用\textbf{强级数判别法}得左侧部分和一致收敛。
        
        \item \textbf{第二个级数}
        
        与第三个级数类似思路可发现第二个级数第$n$项在$[0,1]$的最大值为
        $$\frac{n}{n^2+1}\bigg(\frac{n^2}{n^2+1}\bigg)^{n^2}$$
        估算可发现它与$\frac{1}{n}$同阶，由此\textbf{无法使用强级数判别法}。

        由于最大值点$x_n=\frac{n^2}{n^2+1}$在$n\to\infty$时极限为1，我们尝试在1附近\textbf{估算证明不一致收敛}。

        仍然尝试估算
        $$\bigg|\sum_{k=n+1}^mk(1-x)x^{k^2}\bigg|=(1-x)\sum_{k=n+1}^mkx^{k^2}$$
        为了能够控制，我们尝试取$x$为$f_m(x)$处的最大值点$\frac{m^2}{m^2+1}$，此时上式化为
        $$\frac{1}{m^2+1}\sum_{k=n+1}^mk\bigg(\frac{m^2}{m^2+1}\bigg)^{k^2}$$
        
        利用\textbf{阶估算}可以想到，当$m$与$n$\textbf{同量级}时，所有$(\frac{m^2}{m^2+1})^{k^2}$可以被常数控制，而所有$k$求和则拥有$m^2$阶，可以与分母抵消，这就足以说明它大于等于某常数。由此，我们可以取定$m=2n$。

        \note 这步需要\textbf{较好的阶估算能力}，但即使不会阶估算，\textbf{尝试}取$m=2n$也是相对常见的。

        取$m=2n$后，我们可以进一步估算上式得到其为(不等号是将前面的$k$与后面的$k^2$次方分别缩小至最小)
        $$\frac{1}{4n^1+1}\sum_{k=n+1}^{2n}k\bigg(\frac{4n^2}{4n^2+1}\bigg)^{k^2}\ge\frac{1}{4n^2+1}\sum_{k=n+1}^{2n}n\bigg(\frac{4n^2}{4n^2+1}\bigg)^{(2n)^2}=\frac{n^2}{4n^2+1}\bigg(\frac{4n^2}{4n^2+1}\bigg)^{4n^2}$$
        分别放缩有(第二部分放缩利用了数列极限结论$(1+\frac{1}{k})^k$恒小于$\er$)
        $$\frac{n^2}{4n^2+1}\ge\frac{n^2}{5n^2}=\frac{1}{5}$$
        $$\bigg(\frac{4n^2}{4n^2+1}\bigg)^{4n^2}=\frac{1}{\big(1+\frac{1}{4n^2}\big)^{4n^2}}>\frac{1}{\er}$$
        由此最终得到
        $$\frac{1}{4n^1+1}\sum_{k=n+1}^{2n}k\bigg(\frac{4n^2}{4n^2+1}\bigg)^{k^2}\ge\frac{1}{5\er}$$
        由此，我们取定$\varepsilon=\frac{1}{5\er}$，对任何正整数$N$，取$n=N+1$、$m=2n$、$x=\frac{m^2}{m^2+1}$，即有
        $$\bigg|\sum_{k=n+1}^mk(1-x)x^{k^2}\bigg|=\frac{1}{4n^2+1}\sum_{k=n+1}^{2n}k\bigg(\frac{4n^2}{4n^2+1}\bigg)^{k^2}\ge\frac{1}{5\er}$$
        已经符合一致收敛的\textbf{柯西收敛原理}的否定(可见\exref{not uni})，从而证明了原级数不一致收敛。
    \end{itemize}
}
\end{framed}

\note 注意\textbf{直接算出部分和}有时是有效的手段，此外，证明不一致收敛的另一个可行的方式是证明\textbf{和函数的性质不满足}，如通项连续但和函数不连续。

\begin{framed}
\begin{exercise}
    证明
    $$f(x)=\suminf[n=2]\bigg(\frac{x\sin x}{\ln n}\bigg)^n$$
    在$\rb$上连续。
\end{exercise}
\sol{
    设$f_n(x)=\big(\frac{x\sin x}{\ln n}\big)^n$，利用\textbf{初等函数连续性}可发现固定正整数$n\ge2$时$f_n(x)$是$\rb$上的连续函数。

    可以证明$\suminf[n=2]f_n(x)$在$\rb$上不一致收敛于$f(x)$，但我们仍然能够说明连续性，分为如下两步：
    \begin{itemize}
        \item \textbf{内闭一致收敛性}
        
        我们先证明，对任何$M>0$，$\suminf[n=2]f_n(x)$在$[-M,M]$上一致收敛。

        利用幂函数单调性结论可发现$x\in[-M,M]$时
        $$|f_n(x)|=\frac{|x\sin x|^n}{\ln^nn}\le\frac{|x|^n}{\ln^nn}\le\frac{M^n}{\ln^nn}$$
        由于右侧$\ln n$将任意大，阶估算可知其应比指数下降更快，求和收敛收敛，从而尝试运用\textbf{强级数判别法}。

        由于$n>\er^{2M}$时$\ln n>2M$，$n>\er^{2M}$时有
        $$\frac{M^n}{\ln^nn}<\frac{1}{2^n}$$
        根据\textbf{等比级数}敛散性结论可知$\suminf[n=2]\frac{1}{2^n}$收敛，从而根据\textbf{比较判别法}得$\suminf[n=2]\frac{M^n}{\ln^nn}$也收敛，最终通过强级数判别法得到原级数在$[-M,M]$一致收敛。

        \item \textbf{连续性证明}
        
        根据题中定义的形式，$f(x)$代表原级数在$x$处的\textbf{逐点收敛}结果。由定义可发现\textbf{一致收敛必然逐点收敛}，由此既然原级数在$[-M,M]$上一致收敛到某个函数，它在$[-M,M]$上必然一致收敛到$f(x)$。

        对任何$x_0\in\rb$，取$M=|x_0|+1$，由于原级数在$[-M,M]$一致收敛于$f(x)$，且通项连续，$f$必然在$[-M,M]$连续，从而在$x_0$连续。根据连续性的定义，在$\rb$上连续等价于在$\rb$上\textbf{任何一点都连续}。由于上述过程可对任何$x_0\in\rb$进行，我们就证明了$f$在$\rb$上的连续性。

        \note 这部分证明在考试时可\textbf{省略大部分}。
    \end{itemize}
}
\end{framed}

\note 只需要\textbf{内闭一致收敛}即可说明连续性、导数等\textbf{局部}性质，因为这可以通过在\textbf{一点附近}取区间证明，本题即为一个例子。

\begin{framed}
\begin{exercise}
    计算
    $$\lim_{x\to\infty}\suminf\frac{x}{n^3+x^2}$$
\end{exercise}
\sol{
    由于本题的形式较为复杂，我们尝试在分析后进行计算：
    \begin{itemize}
        \item \textbf{分析与变形}
        
        首先，本题为一个函数项级数的和函数在$x\to+\infty$处的极限问题。但是，我们学过的关于和函数连续性等定理只足以确定和函数在\textbf{一点附近的极限}，而不足以处理\textbf{无穷处的极限}。

        为了将问题转化为一点附近的情况，我们设和函数为$f(x)$，换元$t=\frac{1}{x}$得到
        $$\lim_{x\to\infty}f(x)=\lim_{t\to0}f\bigg(\frac{1}{t}\bigg)$$
        
        \note 这里的等号是指，若右侧极限存在，则左侧极限存在且与右侧相等，本质上利用了\textbf{推广的复合函数极限}结论(可见上学期讲义4.3.1，注意这是教材并未介绍的第二种情况)。

        也即我们将原极限转化为了
        $$\lim_{t\to0}\suminf\frac{\frac{1}{t}}{n^3+\frac{1}{t^2}}=\lim_{t\to0}\suminf\frac{t}{n^3t^2+1}$$
        记
        $$g(t)=\suminf g_n(t),\quad g_n(t)=\frac{t}{n^3t^2+1}$$
        所求极限即为
        $$\lim_{t\to0}g(t)$$

        \item \textbf{结论证明}
        
        可以发现，换元后的和函数$g(t)$在0处有定义：由于任何$g_n(0)$均为0，$g(t)$在0处为0。由于所有$g_n(t)$都是$\rb$上的\textbf{连续函数}，我们只要证明了$g(t)$在0处的连续性，即直接证明了极限为0。

        利用\textbf{基本不等式}得对任何$t\ne0$有
        $$|g_n(t)|=\frac{|t|}{n^3t^2+1}\le\frac{|t|}{2\sqrt{n^3t^2}}=\frac{1}{2n^{3/2}}$$
        由于$t=0$时$g_n(t)=0$，上式对$t=0$也成立。根据\textbf{$p$级数}敛散性结论可知右侧部分和收敛，从而利用\textbf{强级数判别法}得左侧部分和在$\rb$上一致收敛。

        根据题中定义的形式，$g(t)$代表原级数在$t$处的\textbf{逐点收敛}结果。由定义可发现\textbf{一致收敛必然逐点收敛}，由此既然原级数在$\rb$上一致收敛到某个函数，它在$\rb$上必然一致收敛到$g(t)$。

        由于原级数通项连续，$g(t)$在$\rb$上连续，从而所求极限为
        $$\lim_{t\to0}g(t)=g(0)=0$$
    \end{itemize}
}
\end{framed}

\note \textbf{和函数的性质}本身可能结合\textbf{换元技巧}使用，但注意必须\textbf{换元后再判断一致收敛性}。之后我们将默认\textbf{一致收敛结果必然为逐点收敛结果}，不再强调。

\subsubsection{幂级数}
\begin{framed}
\begin{exercise}
    计算幂级数
    $$\suminf\bigg(\sum_{k=1}^n\frac{1}{k}\bigg)x^n,\quad\suminf\frac{(-1)^n+(-2)^n+3^n}{n(n+2)(n+4)}x^n$$
    的收敛域。
\end{exercise}
\sol{
    分为如下步骤：
    \begin{itemize}
        \item \textbf{第一个级数-收敛半径确定}
        
        直接将求和中每项都缩小到$\frac{1}{n}$或放大到1可知
        $$1\le\sqrt[n]{\sum_{k=1}^n\frac{1}{k}}\le\sqrt[n]{n}$$
        利用数列极限知识(可见本讲义27.1.1第5题(3))得
        $$\lim_{n\to\infty}\sqrt[n]{n}=1$$
        而左侧恒为1，从而根据\textbf{夹逼定理}可知中间在$n\to\infty$时的极限也为1。由\textbf{柯西判别法}，原幂级数的收敛半径为1。
        
        \item \textbf{第一个级数-收敛域确定}
        
        利用收敛半径的性质，只需再讨论$\pm1$处是否收敛。利用\textbf{$p$级数}的收敛性结论可知
        $$\lim_{n\to\infty}\sum_{k=1}^n\frac{1}{k}=+\infty$$
        无论1处还是$-1$处，通项的极限均为无穷，不为0，从而发散。

        通过上述讨论，我们最终得到原幂级数的收敛域为$(-1,1)$。
        
        \item \textbf{第二个级数-收敛半径确定}
        
        由于分子、分母形式都较为复杂，我们先进行适当的放缩。由于$(-1)^n$、$(-2)^n$均不会超过$3^n$，我们可以放大得到
        $$\frac{(-1)^n+(-2)^n+3^n}{n(n+2)(n+4)}\le\frac{3^{n+1}}{n(n+2)(n+4)}$$
        同理，由于$(-2)^n\ge-2^n\ge-2\cdot 3^{n-1}$，且$n>1$时$(-1)^n\ge-1\ge -2\cdot 3^{n-2}$，我们可以将分子缩小得到
        $$\frac{(-1)^n+(-2)^n+3^n}{n(n+2)(n+4)}\ge\frac{3^n-2\cdot 3^{n-1}-2\cdot 3^{n-2}}{n(n+2)(n+4)}=\frac{3^{n-2}}{(n+4)^3}=\frac{3^{n-2}}{n(n+2)(n+4)}$$
        综合以上放缩有
        $$\frac{3}{\sqrt[n]{9n(n+2)(n+4)}}\le\sqrt[n]{\frac{(-1)^n+(-2)^n+3^n}{n(n+2)(n+4)}}\le\frac{3\sqrt[n]{3}}{\sqrt[n]{n(n+2)(n+4)}}$$
        利用数列极限知识(后三部分类似本讲义27.1.1第5题(3))得
        $$\lim_{n\to\infty}\sqrt[n]{3}=\lim_{n\to\infty}\sqrt[n]{9}=\lim_{n\to\infty}\sqrt[n]{n}=\lim_{n\to\infty}\sqrt[n]{n+2}=\lim_{n\to\infty}\sqrt[n]{n+4}=1$$
        由此通过\textbf{乘除极限}结论可知
        $$\lim_{n\to\infty}\frac{3}{\sqrt[n]{9n(n+2)(n+4)}}=\lim_{n\to\infty}\frac{3\sqrt[n]{3}}{\sqrt[n]{n(n+2)(n+4)}}=3$$
        从而根据\textbf{夹逼定理}可知中间部分在$n\to\infty$时的极限也为3。由\textbf{柯西判别法}，原幂级数的收敛半径为$\frac{1}{3}$。
        
        \item \textbf{第二个级数-收敛域确定}
        
        利用收敛半径的性质，只需再讨论$\pm\frac{1}{3}$处是否收敛。无论$\frac{1}{3}$处还是$-\frac{1}{3}$处，计算可发现
        $$\bigg|\frac{(-1)^n+(-2)^n+3^n}{n(n+2)(n+4)}x^n\bigg|=\bigg|\frac{(-\frac{1}{3})^n+(-\frac{2}{3})^n+1}{n(n+2)(n+4)}\bigg|\le\frac{3}{n^3}$$
        由\textbf{$p$级数}敛散性结论与数项级数基本性质可知$\suminf\frac{3}{n^3}$收敛，于是根据\textbf{比较判别法}得原幂级数绝对收敛，从而收敛。

        通过上述讨论，我们最终得到原幂级数的收敛域为$[-\frac{1}{3},\frac{1}{3}]$。
    \end{itemize}
}
\end{framed}

\note 幂级数的收敛域本质是数项级数问题，但由于幂级数的特殊性质，一般采用\textbf{先确定收敛半径再讨论端点}的方法。

\note 利用柯西判别法，幂级数的收敛半径只与系数的\textbf{指数部分}有关，因此放缩只需保证指数部分的\textbf{阶}不变即可。

\begin{framed}
\begin{exercise}
    计算幂级数
    $$\suminf\frac{(n-1)!\cdot n!}{(2n-1)!}x^n$$
    的收敛域。
\end{exercise}
\sol{
    分为如下步骤：
    \begin{itemize}
        \item \textbf{收敛半径确定}
        
        记$u_n=\frac{(n-1)!\cdot n!}{(2n-1)!}$，直接计算可知(最后一步利用了\textbf{多项式阶估算}(可见上学期讲义2.4.1)结论)
        $$\lim_{n\to\infty}\frac{u_{n+1}}{u_n}=\lim_{n\to\infty}\frac{(2n-1)!}{(n-1)!\cdot n!}\cdot\frac{n!\cdot(n+1)!}{(2n+1)!}=\lim_{n\to\infty}\frac{n(n+1)}{2n(2n+1)}=\frac{1}{4}$$
        由\textbf{比值判别法}，原幂级数的收敛半径为4。
        
        \item \textbf{收敛域确定}
        
        利用收敛半径的性质，只需再讨论$\pm4$处是否收敛。

        当$x=\pm4$时，直接计算得
        $$\bigg|\frac{u_{n+1}x^{n+1}}{u_nx^n}\bigg|=4\cdot\frac{n(n+1)}{2n(2n+1)}=\frac{2n+2}{2n+1}>1$$
        由此，4处还是$-4$处，通项的绝对值都单调递增，且$|u_1x|=4$，类似\exref{mono plus unconv}即可说明原幂级数通项极限不为0，发散。

        通过上述讨论，我们最终得到原幂级数的收敛域为$(-4,4)$。
    \end{itemize}
}
\end{framed}

\note 对于阶乘常用\textbf{比值判别法}处理，但仍然需要估算的手段辅助。可以考虑记忆\textbf{斯特林公式}$n!\sim\frac{n^n}{\er^n}\sqrt{2\pi n}$以猜测结果，但此公式不在我们的学习范围内，一般\textbf{不能直接使用}。

\begin{framed}
\begin{exercise}
    计算幂级数
    $$\suminf\bigg(\frac{n-1}{n}\bigg)^{n^2}x^n$$
    的收敛域。
\end{exercise}
\sol{
    分为如下步骤：
    \begin{itemize}
        \item \textbf{收敛半径确定}
        
        由于这具有$n$次方的形式，可尝试使用\textbf{柯西判别法}。利用数列极限知识可知(可见上册教材1.3节例17)
        $$\lim_{n\to\infty}\sqrt[n]{\bigg(\frac{n-1}{n}\bigg)^{n^2}}=\lim_{n\to\infty}\bigg(1-\frac{1}{n}\bigg)^n=\frac{1}{\er}$$
        由柯西判别法，原幂级数的收敛半径为$\er$。
        
        \item \textbf{收敛域确定}
        
        利用收敛半径的性质，只需再讨论$\pm\er$处是否收敛。

        先考虑$\er$处，此时对应的数项级数为
        $$\suminf\bigg(\frac{n-1}{n}\bigg)^{n^2}\er^n$$
        这是一个\textbf{正项级数}，由此可以尝试进行阶的估算。

        将$\frac{1}{n}$替换为趋于0的$x$，通项可以写为
        $$(1-x)^{1/x^2}\er^{1/x}$$
        为了能够处理指数，我们将其取\textbf{对数}写为
        $$\er^{\frac{1}{x^2}\ln(1-x)+\frac{1}{x}}=\er^{\frac{\ln(1-x)+x}{x^2}}$$
        至此，我们已经可以考虑$\ln(1-x)$在0处的\textbf{泰勒展开}(可从$\ln(1+x)$的泰勒展开换元得到)
        $$\ln(1-x)=-x-\frac{x^2}{2}-\frac{x^3}{3}+o(x^3)\quad(x\to0)$$
        由此这可以写为(注意$o(1)$即代表$x\to0$时这项趋于0)
        $$\er^{\frac{-\frac{x^2}{2}+o(x^2)}{x^2}}=\er^{-\frac{1}{2}+o(1)}\quad(x\to0)$$
        综合以上分析并利用\textbf{复合函数极限}结论，我们可以得到
        $$\lim_{x\to0}(1-x)^{1/x^2}\er^{1/x}=\er^{-\frac{1}{2}}$$
        通过\textbf{归结原理}有
        $$\lim_{n\to\infty}\bigg(\frac{n-1}{n}\bigg)^{n^2}\er^n=\er^{-\frac{1}{2}}$$
        因此原幂级数在$\er$处通项极限不为0，发散。

        在$-\er$处，原幂级数通项的\textbf{绝对值}不变，由此在$\er$处通项绝对值极限不为0，通项极限也不为0，发散。

        通过上述讨论，我们最终得到原幂级数的收敛域为$(-\er,\er)$。
    \end{itemize}
}
\end{framed}

\note 幂级数在端点处的情况需要作为\textbf{数项级数}进行分析，而分析过程仍然可能相对\textbf{复杂}。不过，总体来说还是遵循之前介绍的数项级数\textbf{基本判别流程}。

\begin{framed}
\begin{exercise}
    计算幂级数
    $$\suminf\frac{x^{n^3}}{10^n}$$
    的收敛域。
\end{exercise}
\sol{
    本题由于只有$n$为立方数的系数非零，无法通过比值或柯西判别法判断，需要直接分析。

    直观来看，$x^{n^3}$会比$x^n$放大下降或增长，因此我们分为两种情况考虑：
    \begin{itemize}
        \item \textbf{下降情况}
        
        当$|x|\le1$时，有
        $$\bigg|\frac{x^{n^3}}{10^n}\bigg|\le\frac{1}{10^n}$$
        根据\textbf{等比级数}敛散性结论可知$\suminf\frac{1}{10^n}$收敛，于是根据\textbf{比较判别法}得原幂级数绝对收敛，从而收敛。

        \item \textbf{增长情况}
        
        当$|x|>1$时，由于$x^{n^3}=(x^{n^2})^n$，而$x^{n^2}$为$x^n$的子列，利用指数函数性质即得
        $$\lim_{n\to\infty}|x|^{n^2}=+\infty$$
        由此，利用极限\textbf{保序性}，存在$N$使得$n>N$时$|x^{n^2}|>20$，于是此时
        $$\bigg|\frac{x^{n^3}}{10^n}\bigg|>\frac{(20)^n}{10^n}=2^n$$
        利用\textbf{夹逼定理}即得
        $$\lim_{n\to\infty}\bigg|\frac{x^{n^3}}{10^n}\bigg|=+\infty$$
        从而原幂级数通项极限不为0，发散。
    \end{itemize}

    通过上述讨论，我们最终得到原幂级数的收敛域为$[-1,1]$。
}
\end{framed}

\note 幂级数的柯西判别法和比值判别法未必有效，此时基本的做法仍然是当作\textbf{数项级数}分析收敛域。幂级数的\textbf{收敛半径}可能使得分析更加简单，不过也未必一定需要先确定收敛半径。

\begin{framed}
\begin{exercise}
    设
    $$f(x)=\ln(1+x+x^2+x^3+x^4)+\arctan x$$
    计算$f$在0处的幂级数展开，并求其收敛域。
\end{exercise}
\sol{
    分为如下步骤：
    \begin{itemize}
        \item \textbf{展开计算}
        
        由于$\ln$中的形式相对复杂，我们利用等比数列求和公式将其化为
        $$\ln\frac{1-x^5}{1-x}=\ln(1-x^5)-\ln(1-x)$$

        \note 这样的转化会导致\textbf{定义域变化}，但由于我们只关心$f$在\textbf{$0$附近}的情况，0附近上述等式是严格成立的。

        利用初等函数在0附近的泰勒级数结论
        $$\ln(1+t)=\suminf(-1)^{n-1}\frac{t^n}{n}$$
        $$\arctan x=\suminf(-1)^n\frac{x^{2n+1}}{2n+1}$$
        由于$x\to0$时$(-x)\to0$、$(-x^5)\to0$，我们可以直接用$-x$与$-x^5$代换$\ln(1+t)$泰勒级数中的$t$得到
        $$\ln(1-x)=\suminf-\frac{x^n}{n},\quad\ln(1-x^5)=\suminf-\frac{x^{5n}}{n}$$

        \note 与复合函数的泰勒展开需要考虑$o$的运算不同，由于幂级数展开两侧是相等的，可以直接进行\textbf{换元}。

        为了得到最终的幂级数展开，我们需要整理三个部分的系数：由于存在$-\ln(1-x)$，对任何正整数$n$都有系数$\frac{1}{n}$；当$n$为奇数时，会额外添加$\arctan x$中的系数$(-1)^{\frac{n-1}{2}}\frac{1}{n}$\ (注意要将展开式中的$2n+1$还原为$n$)；当$n$为5的倍数时，会额外添加$\ln(1-x^5)$中的系数$-\frac{5}{n}$。由此，$f(x)$的幂级数展开最终为(这里$a\mid n$表示$n$为$a$的倍数)
        $$f(x)=\suminf a_nx^n,\quad a_n=\begin{cases}\frac{1}{n}(1+(-1)^{\frac{n-1}{2}})&2\nmid n,\ 5\nmid n\\\frac{1}{n}&2\mid n,\ 5\nmid n\\\frac{1}{n}(-4+(-1)^{\frac{n-1}{2}})&2\nmid n,\ 5\mid n\\-\frac{4}{n}&2\mid n,\ 5\mid n\end{cases}$$

        \item \textbf{收敛半径确定}
        
        \note 由于系数$a_n$均与$\frac{1}{n}$相关，而乘除多项式很难影响等比级数的收敛性，由此可以\textbf{猜测}收敛半径与$\suminf x^n$相同，为1。
        
        由定义可发现
        $$|a_n|\le\frac{5}{n}$$
        由此，当$|x|<1$时有
        $$|a_nx^n|\le\frac{5|x^n|}{n}\le 5|x|^n$$
        根据\textbf{等比级数}敛散性结论可知$\suminf5|x|^n$收敛，于是根据\textbf{比较判别法}得原幂级数绝对收敛，从而收敛。

        另一方面，当$|x|>1$时若原幂级数收敛，通项极限应为0，从而通项中所有是10的倍数的$n$构成的\textbf{子列}极限应为0，但利用\textbf{数列极限阶估算}结论(可见上学期讲义2.4.1)有
        $$\lim_{m\to\infty}-\frac{4}{(10m)}x^{10m}=\infty$$
        从而此子列的极限为无穷，矛盾。

        综合以上，原幂级数的收敛半径为1。

        \item \textbf{收敛域确定}
        
        利用收敛半径的性质，只需再讨论$\pm1$处是否收敛。由于原函数实际上分为$\ln(1+x+x^2+x^3+x^4)$与$\arctan x$两部分，我们先研究简单的$\arctan x$部分，这部分对应的级数为
        $$\suminf(-1)^n\frac{x^{2n+1}}{2n+1}$$
        无论$x$取1还是$-1$，此级数均为\textbf{交错级数}(因为$x$取$-1$相当于每项乘$-1$)，且通项绝对值$\frac{1}{2n+1}$单调递减趋于0。由此，利用\textbf{莱布尼茨判别法}得这部分在$x=\pm1$均收敛。

        \note 注意若这部分级数在$-1$或1发散则无法进行下面的操作，因为\textbf{减去发散级数}难以确定性质。

        利用数项级数的基本性质(收敛与收敛之和收敛、收敛与发散之和发散)，\textbf{原级数通项减去收敛级数的通项并不影响敛散性}，由此原幂级数在$x=\pm1$的敛散性与
        $$\suminf b_nx^n,\quad b_n=\begin{cases}\frac{1}{n}&5\nmid n\\-\frac{4}{n}&5\mid n\end{cases}$$
        相同。

        为了最终判断敛散性，我们将上述幂级数写为
        $$\suminf c_nd_n,\quad c_n=\frac{1}{n},\quad d_n=\begin{cases}x^n&5\nmid n\\(-4)x^n&5\mid n\end{cases}$$
        $\frac{1}{n}$随$n$增加单调递减趋于0，而直接计算可发现无论是$x=1$还是$x=-1$，$d_n$都以10为周期循环，且每个周期求和为0，因此可以估算出对任何正整数$m$有
        $$\bigg|\sum_{n=1}^md_n\bigg|\le\sum_{n=1}^{10}|d_n|=16$$
        从而其部分和有界，利用\textbf{迪利克雷判别法}得原幂级数在$-1$与1均收敛。

        通过上述讨论，我们最终得到原幂级数的收敛域为$[-1,1]$。
    \end{itemize}
}
\end{framed}

\note 注意幂级数展开式中一定\textbf{只能有一个求和符号}，多个函数相加的幂级数展开需要\textbf{整理系数}。

\note 幂级数逐项求导/求原函数后的收敛半径不变可以直观理解为乘除多项式几乎不改变等比级数的收敛性，但注意\textbf{临界情况}(端点处)可能受到影响。

\begin{framed}
\begin{exercise}
    设
    $$f(x)=\frac{\sqrt{|x|}}{2}\ln\frac{1+\sqrt{|x|}}{1-\sqrt{|x|}}$$
    计算$f$在0处关于$|x|$的幂级数展开，并求其收敛域。
\end{exercise}
\sol{
    分为如下步骤：
    \begin{itemize}
        \item \textbf{展开计算}
        
        我们先将原级数转化为
        $$\frac{\sqrt{|x|}}{2}\ln(1+\sqrt{|x|})-\frac{\sqrt{|x|}}{2}\ln(1-\sqrt{|x|})$$

        利用初等函数在0附近的泰勒级数结论
        $$\ln(1+t)=\suminf(-1)^{n-1}\frac{t^n}{n}$$
        由于$x\to0$时$\sqrt{|x|}\to0$、$-\sqrt{|x|}\to0$，我们可以直接用它们代换$\ln(1+t)$泰勒级数中的$t$得到
        $$\ln(1-\sqrt{|x|})=\suminf-\frac{\sqrt{|x|}^n}{n},\quad\ln(1+\sqrt{|x|})=\suminf(-1)^{n-1}\frac{\sqrt{|x|}^n}{n}$$
        可以发现，当$n$为偶数时，上述两个级数在$\sqrt{|x|}^n$前的系数恰好抵消，而$n$为奇数时则成为相等，因此可写出展开式
        $$\ln(1+\sqrt{|x|})-\ln(1-\sqrt{|x|})=\suminf2\cdot\frac{\sqrt{|x|}^{2n-1}}{2n-1}=\suminf\frac{2}{2n-1}\sqrt{|x|}^{2n-1}$$
        乘上系数以后，我们最终得到$f(x)$关于$|x|$的幂级数展开式
        $$f(x)=\suminf\frac{1}{2n-1}\sqrt{|x|}^{2n}=\suminf\frac{1}{2n-1}|x|^n$$

        \note 注意这\textbf{并不是关于$x$的幂级数}。

        \item \textbf{收敛半径确定}
        
        设$t=|x|$，则我们即要判断
        $$\suminf\frac{1}{2n-1}t^n$$
        对$t$的收敛域，且只需考虑$t\ge0$的情况。
        
        我们仍然先确定此幂级数的收敛半径，直接计算可知
        $$\lim_{n\to\infty}\frac{\frac{1}{2n+1}}{\frac{1}{2n-1}}=\lim_{n\to\infty}\frac{2n-1}{2n+1}=1$$
        由\textbf{比值判别法}，此幂级数对$t$的收敛半径为1。

        \item \textbf{收敛域确定}
        
        利用收敛半径的性质，由于$t\ge0$，只需再讨论$t=1$处是否收敛。此时幂级数的通项为
        $$\frac{1}{2n-1}\ge\frac{1}{2n}$$
        由\textbf{$p$级数}敛散性结论与数项级数基本性质可知$\suminf\frac{1}{2n}$发散，于是根据\textbf{比较判别法}得左侧也发散。

        通过上述讨论，我们最终得到原幂级数对$|x|$的收敛域为$[0,1)$，而$|x|\in[0,1)$当且仅当$x\in(-1,1)$，从而它对$x$的收敛域为$(-1,1)$。
    \end{itemize}
}
\end{framed}

\note 只要保证0附近展开换元后仍然是0附近展开，我们就可以在展开式中\textbf{任意换元}，即使换元完不再是幂级数也能保证相等。

\note 仍然注意收敛域对应\textbf{逐点收敛性}，换元并不影响结论。

\begin{framed}
\begin{exercise}
    设
    $$f(x)=\frac{1}{x^2-2x-3}$$
    计算$f$在0处的幂级数展开，并求其收敛域。计算$f$在1处的各阶导数。
\end{exercise}
\sol{
    分为如下步骤：
    \begin{itemize}
        \item \textbf{展开计算}
        
        由于$f(x)$分母相对复杂，我们先进行\textbf{有理函数的化简}：将$x^2-2x-3$因式分解为$(x+1)(x-3)$，并直接裂项或待定系数得到
        $$f(x)=\frac{1}{4}\bigg(\frac{1}{x-3}-\frac{1}{x+1}\bigg)$$
        为了将它化为更易于展开的形式，我们将其写成
        $$f(x)=-\frac{1}{12}\cdot\frac{1}{1-\frac{x}{3}}-\frac{1}{4}\cdot\frac{1}{1+x}$$
        利用初等函数在0附近的泰勒级数结论
        $$\frac{1}{1-t}=\suminf[n=0]t^n$$
        由于$x\to0$时$\frac{x}{3}\to0$、$-x\to0$，我们可以直接用它们代换$\frac{1}{1-t}$泰勒级数中的$t$得到
        $$\frac{1}{1-\frac{x}{3}}=\suminf[n=0]\frac{x^n}{3^n}$$
        $$\frac{1}{1+x}=\suminf[n=0](-1)^nx^n$$
        合并并\textbf{整理系数}最终得到
        $$f(x)=\suminf[n=0]\bigg(-\frac{1}{12\cdot3^n}-\frac{(-1)^n}{4}\bigg)x^n$$
        
        \item \textbf{收敛半径确定}
        
        由于系数为两个指数函数求和，其应与底数较大的相近。由此我们直接进行放缩(注意这里$n$为自然数)
        $$\bigg|-\frac{1}{12\cdot 3^n}-\frac{(-1)^n}{4}\bigg|\le\frac{1}{4}+\frac{1}{12\cdot3^n}\le1$$
        $$\bigg|-\frac{1}{12\cdot 3^n}-\frac{(-1)^n}{4}\bigg|\ge\frac{1}{4}-\frac{1}{12\cdot3^n}\ge\frac{1}{6}$$
        由此，利用$\lim_{n\to\infty}\sqrt[n]{\frac{1}{6}}$与\textbf{夹逼定理}即得
        $$\lim_{n\to\infty}\sqrt[n]{\bigg|-\frac{1}{12\cdot 3^n}-\frac{(-1)^n}{4}\bigg|}=1$$
        由\textbf{柯西判别法}，原幂级数的收敛半径为1。
        
        \item \textbf{收敛域确定}
        
        无论在1处还是$-1$处，均有
        $$\frac{1}{4}-\frac{1}{12\cdot3^n}\le\bigg|\bigg(-\frac{1}{12\cdot3^n}-\frac{(-1)^n}{4}\bigg)x^n\bigg|\le\frac{1}{4}+\frac{1}{12\cdot3^n}$$
        由于左侧右侧在$n\to\infty$时的极限均为$\frac{1}{4}$，利用\textbf{夹逼定理}即得
        $$\lim_{n\to\infty}\bigg|\bigg(-\frac{1}{12\cdot3^n}-\frac{(-1)^n}{4}\bigg)x^n\bigg|=\frac{1}{4}$$
        从而原幂级数通项极限不为0，发散。

        通过上述讨论，我们最终得到原幂级数的收敛域为$(-1,1)$。

        \item \textbf{导数计算}
        
        注意由于幂级数收敛域不包含1，\textbf{无法通过幂级数计算导数}。我们回到展开式
        $$f(x)=\frac{1}{4}\bigg(\frac{1}{x-3}-\frac{1}{x+1}\bigg)$$
        直接归纳或逐次计算可以证明对任何$c\in\rb$有
        $$\frac{\dr^n}{\dr x^n}\frac{1}{x+c}=\frac{(-1)^nn!}{(x+c)^{n+1}}$$
        由此即得
        $$f^{(n)}(x)=\frac{(-1)^nn!}{4}\bigg(\frac{1}{(x-3)^{n+1}}-\frac{1}{(x+1)^{n+1}}\bigg)$$
        代入$x=1$最终得到
        $$f^{(n)}(1)=\frac{(-1)^nn!}{4}\cdot\frac{(-1)^{n+1}-1}{2^{n+1}}=-\frac{n!}{2^{n+3}}(1+(-1)^n)$$
    \end{itemize}
}
\end{framed}

\note 幂级数本质上可以给出\textbf{0处的各阶导数结果}，但其他点处若不能保证收敛性，需要考虑其他方法。

\begin{framed}
\begin{exercise}
    计算幂级数
    $$\suminf n^2x^n$$
    的收敛域与和函数。
\end{exercise}
\sol{
    分为如下步骤：
    \begin{itemize}
        \item \textbf{收敛半径确定}
        
        利用数列极限知识(可见本讲义27.1.1第5题(3))与\textbf{乘除极限}结论得
        $$\lim_{n\to\infty}\sqrt[n]{n^2}=\lim_{n\to\infty}\sqrt[n]{n}\cdot\sqrt[n]{n}=1\cdot1=1$$
        由\textbf{柯西判别法}，原幂级数的收敛半径为1。

        \item \textbf{收敛域确定}
        
        无论在1处还是$-1$处，通项绝对值的极限均为
        $$\lim_{n\to\infty}n^2=+\infty$$
        从而原幂级数通项极限不为0，发散。

        通过上述讨论，我们最终得到原幂级数的收敛域为$(-1,1)$。

        \item \textbf{和函数计算}
        
        设原幂级数在$(-1,1)$上的和函数为$S(x)$。为了能够进行\textbf{逐项积分}，我们需要先提出一个$x$得到
        $$S(x)=x\suminf n^2x^{n-1}$$
        这步\textbf{不影响收敛域}：$\suminf n^2x^n$与$\suminf n^2x^{n-1}$在$x=0$时均收敛，而$x\ne0$时利用数项级数基本性质可知相差非零常数倍的两个级数同敛散。

        由此，$\suminf n^2x^{n-1}$的收敛域也为$(-1,1)$，利用幂级数的\textbf{内闭一致收敛性}可知收敛域内
        $$S(x)=x\suminf n^2x^{n-1}=x\suminf n\frac{\dr x^n}{\dr x}=x\frac{\dr}{\dr x}\suminf nx^n$$
        与上方完全类似，提出$x$后计算可知$(-1,1)$中有
        $$\suminf nx^n=x\suminf nx^{n-1}=x\suminf\frac{\dr x^n}{\dr x}=x\frac{\dr}{\dr x}\suminf x^n$$
        利用\textbf{基本初等函数的泰勒级数}可知
        $$\suminf x^n=\suminf[n=0]x^n-1=\frac{1}{1-x}-1=\frac{x}{1-x}$$
        $$\suminf nx^n=x\frac{\dr}{\dr x}\frac{x}{1-x}=\frac{x}{(1-x)^2}$$
        $$S(x)=x\frac{\dr}{\dr x}\frac{x}{(1-x)^2}=\frac{x(1+x)}{(1-x)^3}$$
        上述结果对$x\in(-1,1)$成立。
    \end{itemize}
}
\end{framed}

\note 计算和函数时务必验证每一步在\textbf{收敛半径内}进行。对于开区间，这件事相对简单，对于闭区间则需要特别讨论\textbf{边界}。

\begin{framed}
\begin{exercise}\label{poly close}
    计算幂级数
    $$\suminf\frac{(-1)^n}{(2n-1)(2n+1)}x^{2n+1}$$
    的收敛域与和函数。
\end{exercise}
\sol{
    分为如下步骤：
    \begin{itemize}
        \item \textbf{换元处理}
        
        由于此幂级数偶数次项的系数均为0，我们考虑处\textbf{换元}使其能通过判别法处理。将它写为
        $$x\suminf\frac{(-1)^n}{(2n-1)(2n+1)}x^{2n}$$
        $x=0$处原级数收敛，而除了$x=0$外的点由数项级数基本性质可知乘常数$x$不影响收敛性，从而原幂级数收敛域与幂级数$\suminf\frac{(-1)^n}{(2n-1)(2n+1)}x^{2n}$相同。

        进一步地，换元$t=x^2$，得到新的关于$t$的幂级数
        $$\suminf\frac{(-1)^n}{(2n-1)(2n+1)}t^n$$
        这样，原问题转化为新幂级数关于$t$的收敛域，且只需考虑$t\ge0$的情况，

        \item \textbf{收敛半径确定}
        
        利用数列极限知识(类似本讲义27.1.1第5题(3))得
        $$\lim_{n\to\infty}\sqrt[n]{2n-1}=\lim_{n\to\infty}\sqrt[n]{2n+1}=1$$
        由此通过\textbf{乘除极限}结论可知
        $$\lim_{n\to\infty}\sqrt[n]{\frac{1}{(2n-1)(2n+1)}}=\lim_{n\to\infty}\frac{1}{\sqrt[n]{2n-1}\cdot\sqrt[n]{2n+1}}=\frac{1}{1\cdot1}=1$$
        由\textbf{柯西判别法}，新幂级数的收敛半径为1。

        \item \textbf{收敛域确定}
        
        利用收敛半径的性质，由于$t\ge0$，只需再讨论$t=1$处是否收敛。此时有
        $$\bigg|\frac{(-1)^n}{(2n-1)(2n+1)}\bigg|=\frac{1}{4n^2-1}<\frac{1}{3n^2}$$
        根据\textbf{$p$级数}敛散性结论与数项级数基本性质可知$\suminf\frac{1}{4n^2}$收敛，于是根据\textbf{比较判别法}得$t=1$时新幂级数绝对收敛，从而收敛。

        通过上述讨论，我们最终得到新幂级数对$x^2$的收敛域为$[0,1]$，而$x^2\in[0,1]$当且仅当$x\in[-1,1]$，从而原幂级数对$x$的收敛域为$[-1,1]$。

        \item \textbf{和函数计算-初步化简}
        
        设原幂级数在$[-1,1]$上的和函数为$S(x)$。我们先考虑\textbf{开区间}$(-1,1)$中的情况。将原级数裂项为
        $$\suminf\bigg((-1)^n\frac{1}{2(2n-1)}-(-1)^n\frac{1}{2(2n+1)}\bigg)x^{2n+1}$$
        与本题第一部分完全类似可证明
        $$\suminf\frac{(-1)^nx^{2n+1}}{2n-1},\quad\suminf\frac{(-1)^nx^{2n+1}}{2n+1}$$
        收敛半径均为1，由此在$(-1,1)$中均收敛，从而有
        $$\forall x\in(-1,1),\quad S(x)=\frac{1}{2}\bigg(\suminf\frac{(-1)^nx^{2n+1}}{2n-1}-\suminf\frac{(-1)^nx^{2n+1}}{2n+1}\bigg)$$

        \item \textbf{和函数计算-第二部分}
        
        我们先计算
        $$\suminf\frac{(-1)^nx^{2n+1}}{2n+1}$$
        利用幂级数的\textbf{内闭一致收敛性}逐项求导可知$(-1,1)$中
        $$\frac{\dr}{\dr x}\suminf\frac{(-1)^nx^{2n+1}}{2n+1}=\suminf(-1)^n\frac{\dr}{\dr x}\frac{x^{2n+1}}{2n+1}=\suminf(-1)^nx^{2n}$$
        利用\textbf{基本初等函数的泰勒级数}可知
        $$\forall t\in(-1,1),\quad\suminf t^n=\suminf[n=0]t^n-1=\frac{1}{1-t}-1=\frac{t}{1-t}$$
        代入$t=-x^2$\ ($x\in(-1,1)$时也有$t\in(-1,1)$)即知
        $$\suminf(-1)^nx^{2n}=-\frac{x^2}{1+x^2}$$
        由于要算的幂级数在0处值为0，利用\textbf{微积分基本定理}可知
        $$\forall x\in(-1,1),\quad\suminf\frac{(-1)^nx^{2n+1}}{2n+1}=\int_0^x\frac{-t^2\dr t}{1+t^2}=\int_0^x\bigg(\frac{1}{1+t^2}-1\bigg)\dr t=\arctan x-x$$

        \item \textbf{和函数计算-第一部分}
        
        接下来计算
        $$\suminf\frac{(-1)^nx^{2n+1}}{2n-1}$$
        我们先将$n$更换为$n+1$\ (这样求和范围将变为0到无穷)得到
        $$\suminf[n=0]\frac{(-1)^{n+1}x^{2n+3}}{2n+1}$$
        再从中提出$-x^2$可以得到其为(与证明第一部分相同分析可知不改变收敛域，从而在$(-1,1)$仍收敛)
        $$(-x^2)\suminf[n=0]\frac{(-1)^nx^{2n+1}}{2n+1}$$
        由于这里的幂级数已经是和函数第二部分加上第一项，其对任何$x\in(-1,1)$的结果为
        $$(-x^2)(\arctan x-x+x)=-x^2\arctan x$$

        \item \textbf{综合结果}

        综合以上，我们最终得到
        $$\forall x\in(-1,1),\quad S(x)=\frac{1}{2}(-x^2\arctan x-\arctan x+x)$$
        利用幂级数的\textbf{内闭一致收敛性}，幂级数的和函数在收敛域必然\textbf{连续}，再由\textbf{初等函数连续性}结论即得
        $$\forall x\in[-1,1],\quad S(x)=\frac{1}{2}(-x^2\arctan x-\arctan x+x)$$
        这已经得到了整个收敛域的和函数表达式。

        \note 准确来说，由连续可知$S(1)$与$S(-1)$为对应点处的左右极限，再用初等函数连续性结论即知这恰好是$\frac{1}{2}(-x^2\arctan x-\arctan x+x)$在两点处的值。
    \end{itemize}
}
\end{framed}

\note 对于一些难以直接计算的情况也可以考虑\textbf{先换元}再计算。本题展示了先计算内部的和函数，再由\textbf{连续性}推到边界的方法。

\note 计算和函数过程务必注意\textbf{确认每一步的收敛性}再操作。

\begin{framed}
\begin{exercise}
    计算
    $$\suminf\frac{(-1)^n}{3n-1}$$
\end{exercise}
\sol{
    分为如下步骤，我们提供两种方法进行和函数计算：
    \begin{itemize}
        \item \textbf{收敛性与构造}
        
        由于这是一个交错级数，且通项绝对值$\frac{1}{3n-1}$单调递减趋于0，利用\textbf{莱布尼茨判别法}得其收敛。我们考虑幂级数
        $$S(x)=\suminf\frac{x^n}{3n-1}$$
        则所求的结果即为$S(-1)$。

        \note 也可构造幂级数$\suminf\frac{x^{3n-1}}{3n-1}$，过程没有本质区别。

        由于我们已经证明了$S(-1)$收敛，利用幂级数的性质可知$S(x)$必然在$[-1,1)$收敛。下面计算$S(x)$在$(-1,1)$的和函数。
        
        \item \textbf{和函数计算-方程方法}
        
        由于分母为一次，虽然无法直接消去，我们利用幂级数的\textbf{内闭一致收敛性}逐项求导仍然可以得到
        $$\forall x\in(-1,1),\quad S'(x)=\suminf\frac{nx^{n-1}}{3n-1}$$
        可以发现，此时在$(-1,1)$有(右侧级数同乘$x$的严谨性类似\exref{poly close}第一部分分析)
        $$xS'(x)=\suminf\frac{nx^n}{3n-1}$$
        由此利用数项级数基本性质可知
        $$3xS'(x)-S(x)=\suminf\bigg(\frac{3nx^n}{3n-1}-\frac{x^n}{3n-1}\bigg)=\suminf x^n$$
        利用\textbf{基本初等函数的泰勒级数}可知
        $$\suminf x^n=\suminf[n=0]x^n-1=\frac{1}{1-x}-1=\frac{x}{1-x}$$
        从而我们得到了$(-1,1)$中
        $$3xS'(x)-S(x)=\frac{x}{1-x}$$
        在$x\ne0$时(由于我们最终要算$S(-1)$，可只考虑$x\in(-1,0)$)，它可以化为
        $$S'(x)-\frac{1}{3x}S(x)=\frac{1}{3(1-x)}$$

        这是一个\textbf{一阶线性常微分方程}，从而可以考虑利用\textbf{积分因子法}求解。直接计算可得(由$x<0$)
        $$\int-\frac{1}{3x}\dr x=-\frac{1}{3}\ln(-x)+C$$
        由此可取一个积分因子
        $$\mu(x)=\er^{-\frac{1}{3}\ln(-x)}=(-x)^{-1/3}=-x^{-1/3}$$
        方程两边同乘$\mu(x)$并消去负号有
        $$\big(x^{-1/3}S(x)\big)'=\frac{x^{-1/3}}{3(1-x)}$$
        由于$S(x)$的首项为$x$，有
        $$\lim_{x\to0}x^{-1/3}S(x)=0$$

        在0到$x$积分可得(这里涉及\textbf{广义积分}，但本身收敛性可以保证)
        $$x^{-1/3}S(x)=\int_0^x\frac{t^{-1/3}}{3(1-t)}\dr t$$
        换元$t=s^3$可知
        $$\int_0^x\frac{t^{-1/3}}{3(1-t)}\dr t=\int_0^{x^{1/3}}\frac{s}{1-s^3}\dr s$$
        于是两侧同乘$x^{1/3}$得到(0处由$S(x)=0$可验证相等)
        $$S(x)=x^{1/3}\int_0^{x^{1/3}}\frac{s}{1-s^3}\dr s$$

        由幂级数的和函数\textbf{连续性}与\textbf{变上限积分连续性}结论代入$x=-1$最终得到
        $$S(-1)=\int_{-1}^0\frac{s}{1-s^3}\dr s$$

        \item \textbf{和函数计算-直接方法}
        
        由于分母为一次，我们希望$S(x)$可以\textbf{求导消去分母}。由此，我们尝试从$S(x)$中提取出$x^{1/3}$\ (严谨性类似\exref{poly close}第一部分分析)得到
        $$S(x)=x^{1/3}\suminf\frac{x^{n-1/3}}{3n-1}$$
        不过，这里右侧并不是一个幂级数，没有直接的方法保证逐项求导的可行性。

        利用与幂级数相似的方法，事实上可以证明$\suminf\frac{x^{n-1/3}}{3n-1}$在$(-1,1)$\textbf{内闭一致收敛}。逐项求导形式上得到
        $$\frac{\dr}{\dr x}\suminf\frac{x^{n-1/3}}{3n-1}=\suminf\frac{\dr}{\dr x}\frac{x^{n-1/3}}{3n-1}=\suminf\frac{x^{n-4/3}}{3}$$
        但是，$x^{-1/3}$在0处并不存在，由此将$S(x)$的第一项拿出，转化为
        $$S(x)=\frac{1}{2}x+x^{1/3}\suminf[n=2]\frac{x^{n-1/3}}{3n-1}$$
        此时可以保证$\suminf[n=2]\frac{x^{n-3/4}}{3}$在$(-1,1)$内闭一致收敛，由此我们最终得到了逐项求导结果
        $$\frac{\dr}{\dr x}\suminf[n=2]\frac{x^{n-1/3}}{3n-1}=\suminf[n=2]\frac{x^{n-4/3}}{3}=\suminf\frac{x^{n-1/3}}{3}$$
        利用\textbf{基本初等函数的泰勒级数}可知
        $$\suminf x^n=\suminf[n=0]x^n-1=\frac{1}{1-x}-1=\frac{x}{1-x}$$
        由此我们最终可以算出
        $$\suminf\frac{x^{n-1/3}}{3}=\frac{x^{2/3}}{3(1-x)}$$
        由于$\suminf[n=2]\frac{x^{n-1/3}}{3n-1}$在0处的值为0，利用\textbf{微积分基本定理}可知
        $$\suminf[n=2]\frac{x^{n-1/3}}{3n-1}=\int_0^x\frac{t^{2/3}}{3(1-t)}\dr t$$
        这样我们就最终写出了结果
        $$S(x)=\frac{1}{2}x+x^{1/3}\int_0^x\frac{t^{2/3}}{3(1-t)}\dr t$$

        换元$t=s^3$可知
        $$\int_0^x\frac{t^{2/3}}{3(1-t)}\dr t=\int_0^{x^{1/3}}\frac{s^4}{1-s^3}\dr s$$

        由此
        $$S(x)=\frac{1}{2}x+x^{1/3}\int_0^{x^{1/3}}\frac{s^4}{1-s^3}\dr s$$

        由幂级数的和函数\textbf{连续性}与\textbf{变上限积分连续性}结论代入$x=-1$最终得到
        $$S(-1)=-\frac{1}{2}+\int_{-1}^0\frac{s^4}{1-s^3}\dr s$$

        \note 将$\frac{s^4}{1-s^3}$写为$\frac{s^4-s+s}{1-s^3}=-s+\frac{s}{1-s^3}$，即可得到与第一种方法完全相同的形式。

        \item \textbf{积分计算}
        
        综合以上讨论，我们最终将结果化为了
        $$\int_{-1}^0\frac{s}{1-s^3}\dr s$$
        换元$r=-s$得到其为
        $$-\int_0^1\frac{r}{1+r^3}\dr r$$
        下面计算此积分。

        可以采用有理函数的\textbf{标准做法}计算，这里我们介绍更简单的配凑做法(可参考上学期讲义7.2.3)。由于(第三个积分通过$\frac{1}{a^2+r^2}$的积分公式得到)
        $$I_1=\int_0^1\frac{3r^2}{1+r^3}\dr r=\int_0^1\frac{\dr r^3}{1+r^3}=\ln(1+r^3)\big|^1_0=\ln2$$
        $$I_2=\int_0^1\frac{1-r+r^2}{1+r^3}\dr r=\int_0^1\frac{1}{1+r}\dr r=\ln(1+r)\big|^1_0=\ln2$$
        $$I_3=\int_0^1\frac{1+r}{1+r^3}\dr r=\int_0^1\frac{1}{1-r+r^2}\dr r=\int_0^1\frac{1}{(r-\frac{1}{2})^2+\frac{3}{4}}=\frac{2}{\sqrt3}\arctan\frac{2r-1}{\sqrt3}\bigg|^1_0=\frac{2\pi}{3\sqrt3}$$
        而待定系数即可发现
        $$r=\frac{1}{2}((1+r)-(1-r+r^2)+r^2)$$
        于是所求积分为
        $$-\bigg(\frac{1}{2}I_3-\frac{1}{2}I_2+\frac{1}{6}I_1\bigg)=\frac{1}{3}\ln2-\frac{\pi}{3\sqrt3}$$
    \end{itemize}
}
\end{framed}

\note 对于数项级数求和，采用\textbf{构造幂级数}技巧可以将它转化为\textbf{积分问题}，但转化后的问题也未必简单，甚至可能不存在初等原函数。

\note 对于这类过程很长的问题，可以\textbf{省略复杂的严谨性说明}，但需要\textbf{保证步骤的正确性}。

\begin{framed}
\begin{exercise}
    考虑函数项级数
    $$S_1(x)=\suminf[n=0]\bigg(\frac{x^{4n+1}}{4n+1}+\frac{x^{4n+3}}{4n+3}-\frac{x^{4n+4}}{2n+2}\bigg),\quad S_2(x)=\suminf[n=0]\bigg(\frac{x^{4n+1}}{4n+1}+\frac{x^{4n+3}}{4n+3}-\frac{x^{2n+2}}{2n+2}\bigg)$$
    计算它们的收敛域，并计算
    $$\suminf[n=0]\bigg(\frac{1}{4n+1}+\frac{1}{4n+3}-\frac{1}{2n+2}\bigg)$$
    进一步计算它们的和函数。
\end{exercise}
\sol{
    我们需要先分析再进行计算：
    \begin{itemize}
        \item \textbf{分析与幂级数构造}
        
        首先，从$S_1(x)$的形式中可以提取出幂级数
        $$S(x)=\suminf a_nx^n,\quad a_n=\begin{cases}\frac{1}{4k+1}&n=4k+1\\0&n=4k+2\\\frac{1}{4k+3}&n=4k+3\\-\frac{1}{2k+2}&n=4k+4\end{cases}$$
        可以发现，$S_1(x)$即为将$S(x)$的每四项加上括号，由此$S_1(x)$的部分和是$S(x)$的部分和的\textbf{子列}。

        与之相对，$S_2(x)$对应的$x$次数并不是单调增加，因此\textbf{无法看作幂级数的子列}。

        不过，两者在形式上都与幂级数接近，由此可以用类似的方法考虑收敛域。对幂级数$S(x)$来说，它几乎是$\suminf x^n$除以多项式，因此可以估算得到收敛半径仍为1。由此，考虑$S_1(x)$与$S_2(x)$也分为$(-1,1)$、$(-\infty,-1)\cup(1,+\infty)$与端点处。

        \item \textbf{$(-1,1)$收敛情况}
        
        此时对$S_1(x)$直接放缩可知
        $$\bigg|\frac{x^{4n+1}}{4n+1}+\frac{x^{4n+3}}{4n+3}-\frac{x^{4n+4}}{2n+2}\bigg|\le\frac{|x|^{4n+1}}{4n+1}+\frac{|x|^{4n+3}}{4n+3}+\frac{|x|^{4n+4}}{2n+2}\le\frac{3\cdot|x|^{4n+1}}{2n+2}\le\frac{3}{2}(x^4)^n$$
        根据\textbf{等比级数}敛散性结论可知$\suminf\frac{3}{2}(x^4)^n$收敛，于是根据\textbf{比较判别法}得$S_1(x)$绝对收敛，从而收敛。

        对$S_2(x)$直接放缩可知
        $$\bigg|\frac{x^{4n+1}}{4n+1}+\frac{x^{4n+3}}{4n+3}-\frac{x^{2n+2}}{2n+2}\bigg|\le\frac{|x|^{4n+1}}{4n+1}+\frac{|x|^{4n+3}}{4n+3}+\frac{|x|^{2n+2}}{2n+2}\le\frac{3\cdot|x|^{2n+2}}{2n+2}\le\frac{3}{2}(x^2)^n$$
        根据\textbf{等比级数}敛散性结论可知$\suminf\frac{3}{2}(x^2)^n$收敛，于是根据\textbf{比较判别法}得$S_2(x)$绝对收敛，从而收敛。
        
        \item \textbf{$(-\infty,-1)\cup(1,+\infty)$收敛情况}
        
        此时，我们希望能够说明\textbf{通项极限非零}。

        对$S_1(x)$，提取出$x^{4n}$并通分整理得到通项为
        $$x^{4n}\frac{8(x+x^3-2x^4)n^2+2(7x+5x^3-8x^4)n+(6x+2x^3-3x^4)}{(4n+1)(4n+3)(2n+2)}$$
        可以估算得到在$|x|>1$时$x+x^3-2x^4\ne0$，由此利用\textbf{多项式阶估算}(可见上学期讲义2.4.1)可知
        $$\lim_{n\to\infty}|8(x+x^3-2x^4)n^2+2(7x+5x^3-8x^4)n+(6x+2x^3-3x^4)|=+\infty$$
        另一方面，利用数列极限的\textbf{等价替换}与\textbf{数列极限阶估算}结论(可见上学期讲义2.4.1)有
        $$\lim_{n\to\infty}\bigg|\frac{x^{4n}}{(4n+1)(4n+3)(2n+2)}\bigg|=\lim_{n\to\infty}\bigg|\frac{x^{4n}}{32n^3}\bigg|=+\infty$$
        再利用\textbf{无穷极限性质}，最终得到$S_1(x)$通项绝对值的极限为$+\infty$，因此通项极限不为0，发散。

        对$S_2(x)$，估算更加容易：由于$x^{4n+1}$与$x^{4n+3}$同号，必然有
        $$\bigg|\frac{x^{4n+1}}{4n+1}+\frac{x^{4n+3}}{4n+3}-\frac{x^{2n+2}}{2n+2}\bigg|\ge\frac{|x|^{4n+1}}{4n+1}+\frac{|x|^{4n+3}}{4n+3}-\frac{|x|^{2n+2}}{2n+2}$$
        由于
        $$\lim_{n\to\infty}x^{2n}=+\infty$$
        存在$N$使得$n>N$时$x^{2n}>2$，此时
        $$\frac{|x|^{4n+3}}{4n+3}-\frac{|x|^{2n+2}}{2n+2}\ge x^{2n+2}\bigg(\frac{2|x|}{4n+3}-\frac{1}{2n+2}\bigg)>0$$
        于是
        $$\bigg|\frac{x^{4n+1}}{4n+1}+\frac{x^{4n+3}}{4n+3}-\frac{x^{2n+2}}{2n+2}\bigg|\ge\frac{|x|^{4n+1}}{4n+1}$$
        利用\textbf{数列极限阶估算}结论(可见上学期讲义2.4.1)有
        $$\lim_{n\to\infty}\frac{|x|^{4n+1}}{4n+1}=+\infty$$
        利用\textbf{夹逼定理}得到$S_2(x)$通项绝对值的极限为$+\infty$，因此通项极限不为0，发散。

        \note 虽然这里的说明细节相对复杂，但总体放缩是很\textbf{松}的，估算并不困难。

        \item \textbf{最终收敛域确定}
        
        最后，我们还需要确定端点$-1$、1处的收敛情况。可以发现，这两点处$S_1(x)$、$S_2(x)$均相同，具体来说
        $$S_1(-1)=S_2(-1)=\suminf[n=0]\bigg(-\frac{1}{4n+1}-\frac{1}{4n+3}-\frac{1}{2n+2}\bigg)$$
        $$S_1(1)=S_2(1)=\suminf[n=0]\bigg(\frac{1}{4n+1}+\frac{1}{4n+3}-\frac{1}{2n+2}\bigg)$$
        对第一个数项级数，利用\textbf{数项级数基本性质}可知它的敛散性与
        $$\suminf[n=0]\bigg(\frac{1}{4n+1}+\frac{1}{4n+3}+\frac{1}{2n+2}\bigg)$$
        相同，这是一个\textbf{正项级数}，且有$n\ge1$时
        $$\frac{1}{4n+1}+\frac{1}{4n+3}+\frac{1}{2n+2}\ge\frac{1}{2n+2}\ge\frac{1}{4n}$$
        由\textbf{$p$级数}敛散性结论与数项级数基本性质可知$\suminf\frac{1}{4n}$发散，于是根据\textbf{比较判别法}得左侧也发散，这就说明了$S_1(-1)$、$S_2(-1)$发散。

        对第二个数项级数，通分得到其为
        $$\suminf[n=0]\frac{8n+5}{(4n+1)(4n+3)(2n+2)}$$
        这也是一个\textbf{正项级数}，且有$n\ge1$时
        $$\frac{8n+5}{(4n+1)(4n+3)(2n+2)}\le\frac{13n}{(4n)^3}\le\frac{13}{64}\cdot\frac{1}{n^2}$$
        由\textbf{$p$级数}敛散性结论与数项级数基本性质可知$\suminf\frac{13}{64n^2}$收敛，于是根据\textbf{比较判别法}得左侧也收敛，这就说明了$S_1(1)$、$S_2(1)$收敛。
        
        综合以上讨论，$S_1(x)$与$S_2(x)$的收敛域均为$(-1,1]$。

        \item \textbf{幂级数收敛域确定}
        
        由于$S_1(x)$的部分和是$S(x)$的部分和的\textbf{子列}，在$S_1(x)$发散的部分$S(x)$必然发散，由此$S(x)$的收敛域在$(-1,1]$中。

        下面证明$S(x)$在1处收敛，结合\textbf{幂级数性质}即知它在$(-1,1]$必然收敛，于是收敛域即为$(-1,1]$。

        我们构造如下两个数列，$b_n$、$c_n$满足
        $$b_n=\frac{1}{n},\quad c_n=\begin{cases}1&n=4k+1\\0&n=4k+2\\1&n=4k+3\\-2&n=4k+4\end{cases}$$
        可以发现此时有
        $$S(1)=\suminf b_nc_n$$
        $\frac{1}{n}$随$n$增加单调递减趋于0，而$c_n$以4为周期循环，且每个周期求和为0，因此可以估算出对任何正整数$m$有
        $$\bigg|\sum_{n=1}^mc_n\bigg|\le\sum_{n=1}^{4}|c_n|=4$$
        从而其部分和有界，利用\textbf{迪利克雷判别法}得原幂级数1处收敛。

        通过上述讨论，我们最终得到原幂级数的收敛域为$(-1,1]$。

        \item \textbf{$S_1(x)$与幂级数和函数计算}
        
        由于$S_1(x)$的部分和是$S(x)$的部分和的子列，且两者收敛域相同，我们实质上证明了
        $$\forall x\in(-1,1],\quad S_1(x)=S(x)$$
        我们先在$(-1,1)$上计算这个函数。在$(-1,1)$上，与证明第二部分类似可估算出幂级数
        $$\suminf[n=0]\frac{x^{4n+1}}{4n+1},\quad\suminf[n=0]\frac{x^{4n+3}}{4n+3},\quad\suminf[n=0]\frac{x^{4n+4}}{2n+2}$$
        均收敛，由此有
        $$S(x)=\suminf[n=0]\frac{x^{4n+1}}{4n+1}+\suminf[n=0]\frac{x^{4n+3}}{4n+3}-\suminf[n=0]\frac{x^{4n+4}}{2n+2}$$
        利用幂级数的\textbf{内闭一致收敛性}，在$(-1,1)$可\textbf{逐项求导}，由此
        $$\begin{aligned}S'(x)&=\frac{\dr}{\dr x}\suminf[n=0]\frac{x^{4n+1}}{4n+1}+\frac{\dr}{\dr x}\suminf[n=0]\frac{x^{4n+3}}{4n+3}-\frac{\dr}{\dr x}\suminf[n=0]\frac{x^{4n+4}}{2n+2}\\ &=\suminf[n=0]\frac{\dr}{\dr x}\frac{x^{4n+1}}{4n+1}+\suminf[n=0]\frac{\dr}{\dr x}\frac{x^{4n+3}}{4n+3}-\suminf[n=0]\frac{\dr}{\dr x}\frac{x^{4n+4}}{2n+2}\\ &=\suminf[n=0]x^{4n}+\suminf[n=0]x^{4n+2}-\suminf[n=0]2x^{4n+3}\end{aligned}$$

        利用\textbf{基本初等函数的泰勒级数}可知
        $$\forall t\in(-1,1),\quad\suminf[t=0]t^n=\frac{1}{1-t}$$
        代入$t=x^4$\ ($x\in(-1,1)$时也有$t\in(-1,1)$)即知
        $$\suminf[n=0]x^{4n}=\frac{1}{1-x^4}$$
        再乘1、$x^2$与$2x^3$\ (严谨性类似\exref{poly close}第一部分分析)即得到$(-1,1)$中
        $$S'(x)=\frac{1+x^2-2x^3}{1-x^4}$$
        可将其因式分解后进一步化简为
        $$S'(x)=\frac{2x^2+x+1}{(x+1)(x^2+1)}=\frac{1}{x+1}+\frac{x}{x^2+1}$$

        由定义可发现$S(0)=0$，积分即得到$(-1,1)$中有(第二部分换元$s=t^2$)
        $$\int_0^x\frac{1}{t+1}\dr t=\ln(t+1)$$
        $$\int_0^x\frac{t\dr t}{t^2+1}=\int_0^{x^2}\frac{\dr s}{2(s+1)}=\frac{1}{2}\ln(x^2+1)$$
        综合以上即得
        $$S(x)=\ln(x+1)+\frac{1}{2}\ln(x^2+1)$$
        利用幂级数的\textbf{内闭一致收敛性}，幂级数的和函数在收敛域必然\textbf{连续}，再由\textbf{初等函数连续性}结论即得
        $$\forall x\in(-1,1],\quad S_1(x)=S(x)=\ln(x+1)+\frac{1}{2}\ln(x^2+1)$$

        \note 注意不引入幂级数就\textbf{无法说明1处连续性}，因此我们才需要额外引入$S(x)$。

        题中要算的数项级数即为$S(x)=S_1(x)$在1处的值，由此可知
        $$\suminf[n=0]\bigg(\frac{1}{4n+1}+\frac{1}{4n+3}-\frac{1}{2n+2}\bigg)=\frac{3}{2}\ln2$$
        
        \item \textbf{$S_2(x)$计算}
        
        最后，我们来计算$S_2(x)$。在$(-1,1)$上，与证明第二部分类似可估算出幂级数
        $$\suminf[n=0]\frac{x^{4n+1}}{4n+1},\quad\suminf[n=0]\frac{x^{4n+3}}{4n+3},\quad\suminf[n=0]\frac{x^{2n+2}}{2n+2}$$
        $$\suminf[n=0]\frac{x^{4n+1}}{4n+1},\quad\suminf[n=0]\frac{x^{4n+3}}{4n+3},\quad\suminf[n=0]\frac{x^{4n+4}}{2n+2}$$
        均收敛，由此有
        $$S_2(x)=\suminf[n=0]\frac{x^{4n+1}}{4n+1}+\suminf[n=0]\frac{x^{4n+3}}{4n+3}-\suminf[n=0]\frac{x^{2n+2}}{2n+2}$$
        利用幂级数的\textbf{内闭一致收敛性}，在$(-1,1)$可\textbf{逐项求导}，由此
        $$\begin{aligned}S_2'(x)&=\frac{\dr}{\dr x}\suminf[n=0]\frac{x^{4n+1}}{4n+1}+\frac{\dr}{\dr x}\suminf[n=0]\frac{x^{4n+3}}{4n+3}-\frac{\dr}{\dr x}\suminf[n=0]\frac{x^{2n+2}}{2n+2}\\ &=\suminf[n=0]\frac{\dr}{\dr x}\frac{x^{4n+1}}{4n+1}+\suminf[n=0]\frac{\dr}{\dr x}\frac{x^{4n+3}}{4n+3}-\suminf[n=0]\frac{\dr}{\dr x}\frac{x^{2n+2}}{2n+2}\\ &=\suminf[n=0]x^{4n}+\suminf[n=0]x^{4n+2}-\suminf[n=0]x^{2n+1}\end{aligned}$$

        利用\textbf{基本初等函数的泰勒级数}可知
        $$\forall t\in(-1,1),\quad\suminf[t=0]t^n=\frac{1}{1-t}$$
        代入$t=x^4$与$t=x^2$\ ($x\in(-1,1)$时也有$t\in(-1,1)$)即知
        $$\suminf[n=0]x^{4n}=\frac{1}{1-x^4},\quad\suminf[n=0]x^{2n}=\frac{1}{1-x^2}$$
        再乘$x$的适当次方\ (严谨性类似\exref{poly close}第一部分分析)即得到$(-1,1)$中
        $$S_2'(x)=\frac{1+x^2}{1-x^4}-\frac{x}{1-x^2}=\frac{1}{1+x}$$
        由定义可发现$S_2(0)=0$，积分得到$(-1,1)$中有
        $$S_2(x)=\int_0^x\frac{1}{t+1}\dr t=\ln(1+x)$$
        但由于它不为幂级数，\textbf{无法保证在1处连续}，1处的值需要由$S_1(x)=S_2(x)$计算，最终得到
        $$S_2(x)=\begin{cases}\ln(1+x)&x\in(-1,1)\\\frac{3}{2}\ln2&x=1\end{cases}$$
    \end{itemize}
}
\end{framed}

\note 注意\textbf{换序}对收敛性的影响：对数项级数条件收敛而言，它可能\textbf{改变收敛性与收敛结果}，对函数项级数而言则可能影响收敛的一致性，从而影响和函数的形式。

\subsubsection{傅里叶级数}
\begin{framed}
\begin{exercise}
    求如下函数在$(-\pi,\pi]$上的傅里叶级数：
    $$f(x)=\sin^2x,\quad x\in[-\pi,\pi]$$
\end{exercise}
\sol{
    本题可以用常规方法计算，我们这里介绍更简便的方法：由于
    $$\forall x\in(-\pi,\pi],\quad\sin^2x=\frac{1-\cos(2x)}{2}=\frac{1}{2}-\frac{1}{2}\cos(2x)$$
    而右侧已经为傅里叶级数的形式，这必然为$\sin^2x$的傅里叶级数(此结论称为傅里叶级数的\textbf{唯一性})。

    这里给出严格的证明：利用三角函数系的\textbf{正交性}可知当$n$不为0或2时
    $$a_n=\frac{1}{\pi}\bigg(\frac{1}{2}\int_{-\pi}^\pi\cos(nx)\dr x-\frac{1}{2}\int_{-\pi}^{\pi}\cos(nx)\cos(2x)\dr x\bigg)=0$$
    而对任何正整数$n$有
    $$b_n=\frac{1}{\pi}\bigg(\frac{1}{2}\int_{-\pi}^\pi\sin(nx)\dr x-\frac{1}{2}\int_{-\pi}^{\pi}\sin(nx)\cos(2x)\dr x\bigg)=0$$
    由此傅里叶系数中只有$a_0$、$a_2$非零，进一步利用正交性可知
    $$a_0=\frac{1}{\pi}\bigg(\frac{1}{2}\int_{-\pi}^\pi\dr x-\frac{1}{2}\int_{-\pi}^{\pi}\cos(2x)\dr x\bigg)=\frac{1}{2\pi}\int_{-\pi}^\pi1\dr x=1$$
    $$a_2=\frac{1}{\pi}\bigg(\frac{1}{2}\int_{-\pi}^\pi\cos(2x)\dr x-\frac{1}{2}\int_{-\pi}^{\pi}\cos(2x)\cos(2x)\dr x\bigg)=-\frac{1}{2\pi}\int_{-\pi}^\pi\cos^2(2x)\dr x=-\frac{1}{2}$$
    这就得到了结论。
}
\end{framed}

\note 傅里叶级数的\textbf{唯一性}与\textbf{收敛性}结论有时可以简化计算，考试时仍然推荐\textbf{使用常规方法}，但可以通过唯一性验证正确性。

\begin{framed}
\begin{exercise}\label{fourier n4}
    设以$2\pi$为周期的函数$f(x)$满足$f(x)=x^2$在$(-\pi,\pi]$成立，将它在$(-\pi,\pi]$展开为傅里叶级数，计算其和函数，并计算
    $$\suminf\frac{1}{n^2},\quad\suminf\frac{1}{(2n-1)^2},\quad\suminf\frac{1}{n^4}$$
\end{exercise}
\sol{
    分为如下步骤：
    \begin{itemize}
        \item \textbf{系数计算}
        
        设
        $$x^2\sim\frac{a_0}{2}+\suminf(a_n\cos(nx)+b_n\sin(nx))$$
        由于$x^2$在$(-\pi,\pi)$为\textbf{偶函数}，由傅里叶系数性质可知所有$b_n=0$，另一方面，由傅里叶级数定义可知
        $$a_0=\frac{1}{\pi}\int_{-\pi}^\pi x^2\dr x=\frac{2}{3}\pi^2$$
        利用\textbf{分部积分}计算可知$n>0$时
        $$a_n=\frac{1}{\pi}\int_{-\pi}^\pi x^2\cos(nx)\dr x=\frac{1}{n\pi}\int_{-\pi}^\pi x^2\dr\sin(nx)=\frac{1}{n\pi}\bigg(x^2\sin(nx)\big|^\pi_{-\pi}-\int_{-\pi}^\pi2x\sin(nx)\dr x\bigg)$$
        进一步化简得到
        $$a_n=-\frac{2}{n\pi}\int_{-\pi}^\pi x\sin(nx)\dr x=\frac{2}{n^2\pi}\int_{-\pi}^\pi x\dr\cos(nx)=\frac{2}{n^2\pi}\bigg(x\cos(nx)\big|^\pi_{-\pi}-\int_{-\pi}^\pi\cos(nx)\dr x\bigg)$$
        由于$n>0$，第二部分积分为0，由此最终得到系数
        $$a_n=\frac{4(-1)^n}{n^2}$$
        也即$[-\pi,\pi]$中
        $$x^2\sim\frac{1}{3}\pi^2+\suminf\frac{4(-1)^n}{n^2}\cos(nx)$$

        \item \textbf{和函数计算}
        
        利用周期性可以发现$f(-\pi)=\pi^2=(-\pi)^2$，由此$f(x)=x^2$在$[-\pi,\pi]$恒成立，它在$\rb$上连续。又由于$f(x)$在$[-\pi,\pi]$分段单调，利用傅里叶级数\textbf{逐点收敛结论}即得
        $$\forall x\in\rb,\quad\frac{1}{3}\pi^2+\suminf\frac{4(-1)^n}{n^2}\cos(nx)=f(x)$$
        
        \item \textbf{级数计算}
        
        在和函数表达式中代入$x=\pi$即得
        $$\frac{1}{3}\pi^2+\suminf\frac{4(-1)^n\cdot(-1)^n}{n^2}=\pi^2$$
        利用数项级数基本性质化简即得
        $$\suminf\frac{1}{n^2}=\frac{\pi^2}{6}$$
        利用数项级数基本性质可知
        $$\suminf\frac{1}{4n^2}=\frac{\pi^2}{24}$$
        而根据定义可发现
        $$\sum_{n=1}^m\frac{1}{(2n-1)^2}=\sum_{n=1}^{2m}\frac{1}{n^2}-\sum_{n=1}^m\frac{1}{(2n)^2}=\sum_{n=1}^{2m}\frac{1}{n^2}-\sum_{n=1}^m\frac{1}{4n^2}$$
        在上式中令$m\to\infty$，则第一个级数趋于$\frac{\pi^2}{6}$\ (它的部分和是$\suminf\frac{1}{n^2}$部分和的\textbf{子列})，第二个级数趋于$\frac{\pi^2}{24}$，由此最终得到
        $$\suminf\frac{1}{(2n-1)^2}=\frac{\pi^2}{6}-\frac{\pi^2}{24}=\frac{\pi^2}{8}$$
        为了从系数为二阶的级数中算出$\frac{1}{n^4}$求和，需要应用\textbf{帕塞瓦尔等式}。由$f(x)$在$[-\pi,\pi]$连续可知
        $$\frac{1}{\pi}\int_{-\pi}^\pi x^4\dr x=\frac{1}{2}\bigg(\frac{2}{3}\pi^2\bigg)^2+\suminf\frac{16}{n^4}$$
        计算积分并利用数项级数基本性质化简可知
        $$\suminf\frac{1}{n^4}=\frac{1}{16}\bigg(\frac{2}{5}\pi^4-\frac{2}{9}\pi^4\bigg)=\frac{\pi^4}{90}$$
        综合以上，我们已经算出了三个级数的结果。
    \end{itemize}
}
\end{framed}

\note 计算傅里叶系数常利用\textbf{分部积分}的方式，注意过程中边界项的计算。非连续情况计算和函数需要考虑\textbf{左右极限平均值}，若给出了$f(x)$在$(-\pi,\pi)$的表达式$g(x)$，且$g(x)$在$[-\pi,\pi]$连续、分段单调，和函数应为(这里只给出一个周期的值，再利用以$2\pi$为周期可知任何点处值)
$$S(x)=\begin{cases}g(x)&x\in(-\pi,\pi)\\\frac{g(-\pi)+g(\pi)}{2}&x=\pm\pi\end{cases}$$

\note 注意记忆在\textbf{一般区间}$[-l,l]$计算傅里叶级数的表达式，一般区间的\textbf{收敛性结论不变}。

\note 利用傅里叶级数计算数项级数有\textbf{直接代入}与\textbf{帕塞瓦尔等式}两种方法，一般通过$n$的次数即可确定。

\begin{framed}
\begin{exercise}\label{fourier n6}
    设以$2\pi$为周期的奇函数$f(x)$满足
    $$f(x)=x(\pi-x),\quad x\in[0,\pi]$$
    求$f$在$(-\pi,\pi]$上的傅里叶级数并计算
    $$\suminf\frac{1}{n^6}$$
\end{exercise}
\sol{
    分为如下步骤：
    \begin{itemize}
        \item \textbf{系数计算}
        
        设
        $$f(x)\sim\frac{a_0}{2}+\suminf(a_n\cos(nx)+b_n\sin(nx))$$
        由于$f(x)$为\textbf{奇函数}，由傅里叶系数性质可知所有$a_n=0$，且
        $$b_n=\frac{1}{\pi}\int_{-\pi}^\pi\sin(nx)f(x)\dr x=\frac{2}{\pi}\int_0^\pi\sin(nx)f(x)\dr x$$
        代入$f(x)$可得
        $$b_n=\frac{2}{\pi}\int_0^\pi x(\pi-x)\sin(nx)\dr x=-\frac{2}{n\pi}\int_0^\pi x(\pi-x)\dr\cos(nx)$$
        进行分部积分，由于$x(\pi-x)$在0与$\pi$均为0，边界项为0，从而
        $$b_n=\frac{2}{n\pi}\int_0^\pi(\pi-2x)\cos(nx)\dr x=\frac{2}{n^2\pi}\int_0^\pi(\pi-2x)\dr\sin(nx)$$
        再次分部积分，由于$\sin(nx)$在0处与$\pi$处均为0即得
        $$b_n=-\frac{2}{n^2\pi}\int_0^\pi\sin(nx)\dr(\pi-2x)=\frac{4}{n^2\pi}\int_0^\pi\sin(nx)\dr x=-\frac{4}{n^3\pi}\cos(nx)\big|^\pi_0=\frac{4(1-(-1)^n)}{n^3\pi}$$
        由于系数在$n$为偶数时为0，$n$为奇数时为2，我们最终可以写出傅里叶级数
        $$f(x)\sim\suminf\frac{8}{(2n-1)^3\pi}\sin((2n-1)x)$$
        
        \item \textbf{级数计算}
        
        由于$x(\pi-x)$在0处与$\pi$处均为0，可发现$f(x)$在$\rb$上连续，由此通过\textbf{帕塞瓦尔等式}得到
        $$\frac{1}{\pi}\int_{-\pi}^\pi f^2(x)\dr x=\suminf\frac{64}{(2n-1)^6\pi^2}$$
        利用$f(x)$为奇函数可得左侧积分为(这里直接展开为多项式计算即可)
        $$\frac{2}{\pi}\int_0^\pi x^2(\pi-x)^2\dr x=\frac{\pi^4}{15}$$
        由此利用数项级数基本性质化简可知
        $$\suminf\frac{1}{(2n-1)^6}=\frac{\pi^6}{960}$$
        为了计算最终结果，设(利用\textbf{$p$级数}敛散性结论可知其收敛)
        $$S=\suminf\frac{1}{n^6}$$
        我们先根据定义写出
        $$\sum_{n=1}^m\frac{1}{(2n-1)^6}=\sum_{n=1}^{2m}\frac{1}{n^6}-\sum_{n=1}^m\frac{1}{(2n)^6}=\sum_{n=1}^{2m}\frac{1}{n^6}-\sum_{n=1}^m\frac{1}{64n^6}$$
        在上式中令$m\to\infty$，则第一个级数趋于$S$\ (它的部分和是$\suminf\frac{1}{n^6}$部分和的\textbf{子列})，第二个级数趋于$\frac{S}{64}$，由此最终得到
        $$\frac{\pi^6}{960}=S-\frac{1}{64}S$$
        求解得
        $$S=\frac{\pi^6}{945}$$
    \end{itemize}
}
\end{framed}

\note 注意只给定$f(x)$在$[0,\pi]$中的值时，\textbf{奇偶性}会带来完全不同的性质与计算结果。

\note 注意当两级数均收敛时对$\suminf f(n)$的计算与对$\suminf f(2n-1)$的计算可\textbf{相互转化}。

\begin{framed}
\begin{exercise}
    构造合适的函数$f(x)$，通过$f$在$(-\pi,\pi]$上的傅里叶级数与帕塞瓦尔等式计算
    $$\suminf\frac{1}{n^8}$$
\end{exercise}
\sol{
    分为如下步骤：
    \begin{itemize}
        \item \textbf{基础计算}
        
        由于傅里叶级数本身的计算复杂性，我们往往希望利用\textbf{多项式}进行配凑。由此，我们先进行分部积分得到一些结果，对\textbf{正整数}$n$有(计算过程类似\ref{fourier n4}与\exref{fourier n6})：
        $$\int_0^\pi x\cos(nx)\dr x=\frac{(-1)^n-1}{n^2}$$
        $$\int_0^\pi x\sin(nx)\dr x=\frac{(-1)^{n-1}\pi}{n}$$
        $$\int_0^\pi x^2\cos(nx)\dr x=\frac{2(-1)^n\pi}{n^2}$$
        $$\int_0^\pi x^2\sin(nx)\dr x=\frac{\pi^2n^2(-1)^{n-1}+2((-1)^n-1)}{n^3}$$
        $$\int_0^\pi x^3\cos(nx)\dr x=\frac{3\pi^2n^2(-1)^n+6(1-(-1)^n)}{n^4}$$
        $$\int_0^\pi x^3\sin(nx)\dr x=\frac{\pi^3n^2(-1)^{n-1}+6(-1)^n\pi}{n^3}$$

        \item \textbf{函数构造}
        
        由于我们希望能计算$\frac{1}{n^8}$求和，系数应当尽量接近$\frac{1}{n^4}$。通过以上的计算结果可以发现，$x^3\cos(nx)$是可能符合要求的。由于
        $$\int_0^\pi x^3\cos(nx)\dr x=\frac{3\pi^2(-1)^n}{n^2}+\frac{6(1-(-1)^n)}{n^4}$$
        可利用$x^2$消去第一项得到
        $$\int_0^\pi(2x^3-3\pi x^2)\cos(nx)=\frac{12(1-(-1)^n)}{n^4}$$
        由此，我们令$f(x)$是以$2\pi$为周期的\textbf{偶函数}，且
        $$f(x)=2x^3-3\pi x^2,\quad x\in[0,\pi)$$
        由于其为偶函数，其傅里叶级数为
        $$f(x)\sim\frac{a_0}{2}+\suminf a_n\cos(nx)$$
        且
        $$a_0=\frac{2}{\pi}\int_0^\pi(2x^3-3\pi x^2)=-\pi^3$$
        $$a_n=\frac{2}{\pi}\int_0^\pi(2x^3-3\pi x^2)\cos(nx)=\frac{24(1-(-1)^n)}{\pi n^4}$$
        由于$1-(-1)^n$在$n$为偶数时为0，$n$为奇数时为2，我们最终可以写出傅里叶级数
        $$f(x)\sim-\frac{\pi^3}{2}+\suminf\frac{48}{\pi(2n-1)^4}\cos((2n-1)\pi)$$
        
        \item \textbf{最终结果}
        
        由定义$f(x)$在$[0,\pi)$可积，进一步由偶函数可知$f(x)$在$[-\pi,\pi]$上可积，由此通过\textbf{帕塞瓦尔等式}得到
        $$\frac{1}{\pi}\int_{-\pi}^\pi f^2(x)\dr x=\frac{\pi^6}{2}+\suminf\frac{2304}{(2n-1)^8\pi^2}$$
        利用$f(x)$为偶函数可得左侧积分为(这里直接展开为多项式计算即可)
        $$\frac{2}{\pi}\int_0^\pi(2x^3-3\pi x^2)^2\dr x=\frac{26\pi^6}{35}$$
        由此利用数项级数基本性质化简可知
        $$\suminf\frac{1}{(2n-1)^8}=\frac{\pi^8}{2304}\bigg(\frac{26}{35}-\frac{1}{2}\bigg)=\frac{17}{161280}\pi^8$$
        为了计算最终结果，设(利用\textbf{$p$级数}敛散性结论可知其收敛)
        $$S=\suminf\frac{1}{n^8}$$
        我们先根据定义写出
        $$\sum_{n=1}^m\frac{1}{(2n-1)^8}=\sum_{n=1}^{2m}\frac{1}{n^6}-\sum_{n=1}^m\frac{1}{(2n)^8}=\sum_{n=1}^{2m}\frac{1}{n^6}-\sum_{n=1}^m\frac{1}{256n^8}$$
        在上式中令$m\to\infty$，则第一个级数趋于$S$\ (它的部分和是$\suminf\frac{1}{n^6}$部分和的\textbf{子列})，第二个级数趋于$\frac{S}{256}$，由此最终得到
        $$\frac{17}{161280}\pi^8=S-\frac{1}{256}S$$
        求解得
        $$S=\frac{\pi^8}{9450}$$
    \end{itemize}
}
\end{framed}

\note 上述方法事实上可以计算出所有$k$为正整数的$\suminf\frac{1}{n^{2k}}$。

\note 部分数项级数的问题也存在从\textbf{傅里叶级数}出发的构造计算方法，但由于过程往往\textbf{非常复杂}，仍然建议\textbf{优先考虑幂级数}。

\begin{framed}
\begin{exercise}
    设以$2\pi$为周期的连续函数$f(x)$各傅里叶系数为$a_n,b_n$，计算
    $$F(x)=\frac{1}{\pi}\int_{-\pi}^\pi f(t-x)f(2x-t)\dr t$$
    的傅里叶级数。
\end{exercise}
\sol{
    分为如下步骤：
    \begin{itemize}
        \item \textbf{分析与变形}
        
        设
        $$F(x)\sim\frac{A_0}{2}+\suminf(A_n\cos(nx)+B_n\sin(nx))$$
        由于$F(x)$的傅里叶级数需要计算
        $$\int_{-\pi}^\pi F(x)\sin(nx)\dr x=\frac{1}{\pi}\int_{-\pi}^\pi\bigg(\int_{-\pi}^\pi f(t-x)f(2x-t)\dr t\bigg)\sin(nx)\dr x$$
        等式子，即使能够利用连续性交换积分次序，$f(t-x)f(2x-t)\sin(nx)$这样的式子仍然无法积分。由此，我们需要先行化简。

        设$F(x)$中的被积函数$g(t)=f(t-x)f(2x-t)$，由于$f$以$2\pi$为周期，可以发现$g(t)$也是以$2\pi$为周期的：
        $$g(t+2\pi)=f(t-x+2\pi)f(2x-t-2\pi)=f(t-x)f(2x-t)$$
        由此，$F(x)$为$g(t)$在一个周期的积分除以$\pi$，利用周期函数的性质，\textbf{周期函数在一个周期的积分值与起点无关}，因此我们可以对$t$作\textbf{平移}。具体来说，我们换元$s=t-x$以消除一个$x$，得到(这里第三个等号即利用了周期函数的性质)
        $$F(x)=\frac{1}{\pi}\int_{-\pi}^\pi g(t)\dr t=\frac{1}{\pi}\int_{-\pi-x}^{\pi-x}g(s+x)\dr s=\frac{1}{\pi}\int_{-\pi}^\pi g(s+x)\dr s$$
        于是有
        $$F(x)=\frac{1}{\pi}\int_{-\pi}^\pi f(s)f(x-s)\dr s$$

        \item \textbf{系数计算}
        
        利用复合函数与初等函数的\textbf{连续性}结论，$f(s)f(x-s)\sin(nx)$与$f(s)f(x-s)\cos(nx)$都是$\rb^2$上的\textbf{连续函数}，由此可\textbf{交换积分次序}得到(这里将$f(s)$提到左侧是利用它在对$x$积分时为常数)
        $$A_n=\frac{1}{\pi^2}\int_{-\pi}^\pi\cos(nx)\dr x\int_{-\pi}^\pi f(s)f(x-s)\dr s=\frac{1}{\pi^2}\int_{-\pi}^\pi f(s)\dr s\int_{-\pi}^\pi f(x-s)\cos(nx)\dr x$$
        为了计算积分，我们需要再次进行\textbf{平移}：与之前可类似验证$f(x-s)\cos(nx)$对$x$是以$2\pi$为周期的函数，由此可换元$y=x-s$并将积分范围重新写为$[-\pi,\pi]$得到
        $$\int_{-\pi}^\pi f(x-s)\cos(nx)\dr x=\int_{-\pi}^\pi f(y)\cos(ny+ns)\dr y$$
        由于$\cos(ny+ns)=\cos(ny)\cos(ns)-\sin(ny)\sin(ns)$，利用$f$的傅里叶系数\textbf{定义}最终得到(这里设$b_0=0$，为$f(x)\sin(0x)$的积分结果)
        $$\begin{aligned}\int_{-\pi}^\pi f(x-s)\cos(nx)\dr x&=\int_{-\pi}^\pi(\cos(ns)\cdot f(y)\cos(ny)-\sin(ns)\cdot f(y)\sin(ny))\dr y\\ &=\pi(\cos(ns)a_n-\sin(ns)b_n)\end{aligned}$$
        再次积分即得
        $$A_n=\frac{1}{\pi}\int_{-\pi}^\pi(a_n\cdot f(s)\cos(ns)-b_n\cdot f(s)\sin(ns))=a_n^2-b_n^2$$
        具体来说，由于规定了$b_0=0$，有
        $$A_0=a_0^2,\quad\forall n\in\mathbb{N}^*,\quad A_n=a_n^2-b_n^2$$
        完全类似计算可知
        $$B_n=\frac{1}{\pi^2}\int_{-\pi}^\pi\sin(nx)\dr x\int_{-\pi}^\pi f(s)f(x-s)\dr s=\frac{1}{\pi^2}\int_{-\pi}^\pi f(s)\dr s\int_{-\pi}^\pi f(x-s)\sin(nx)\dr x$$
        $$\begin{aligned}\int_{-\pi}^\pi f(x-s)\sin(nx)\dr x&=\int_{-\pi}^\pi f(y)\sin(ny+ns)\dr y\\ &=\int_{-\pi}^\pi(\cos(ns)\cdot f(y)\sin(ny)+\sin(ns)\cdot f(y)\cos(ny))\dr y\\ &=\pi(b_n\cos(ns)+a_n\sin(ns))\end{aligned}$$
        于是
        $$B_n=\frac{1}{\pi}\int_{-\pi}^\pi(b_n\cdot f(s)\cos(ns)+a_n\cdot f(s)\sin(ns))=2a_nb_n$$
        由此我们最终可以写出傅里叶级数
        $$F(x)\sim\frac{a_0^2}{2}+\suminf((a_n^2-b_n^2)\cos(nx)+2a_nb_n\sin(nx))$$
    \end{itemize}
}
\end{framed}

\note 计算傅里叶系数的过程中也常用到\textbf{周期函数积分}相关的积分技巧。

\begin{framed}
\begin{exercise}
    设$p\in\rb$且$p$不是整数，以$2\pi$为周期的函数$f(x)$满足
    $$f(x)=\cos(px),\quad x\in(-\pi,\pi]$$
    求$f$在$(-\pi,\pi]$上的傅里叶级数与和函数，以此证明$\sin t\ne0$时
    $$\frac{1}{\sin t}=\frac{1}{t}+\suminf(-1)^n\bigg(\frac{1}{t+n\pi}+\frac{1}{t-n\pi}\bigg)$$
    并计算
    $$\int_0^{+\infty}\frac{\sin x}{x}\dr x$$
\end{exercise}
\sol{
    分为如下步骤：
    \begin{itemize}
        \item \textbf{系数计算}
        
        设
        $$f(x)\sim\frac{a_0}{2}+\suminf(a_n\cos(nx)+b_n\sin(nx))$$
        由于$f(x)$为\textbf{偶函数}，由傅里叶系数性质可知所有$b_n=0$，且
        $$a_n=\frac{1}{\pi}\int_{-\pi}^\pi\cos(nx)f(x)\dr x=\frac{2}{\pi}\int_0^\pi\cos(nx)f(x)\dr x$$
        代入$f(x)$可得
        $$a_n=\frac{2}{\pi}\int_0^\pi\cos(nx)\cos(px)\dr x$$
        利用\textbf{积化和差}公式可得
        $$a_n=\frac{1}{\pi}\int_0^\pi(\cos((n-p)x)+\cos((n+p)x))\dr x=\frac{1}{\pi}\bigg(\frac{\sin((n-p)\pi)}{n-p}+\frac{\sin((n+p)\pi)}{n+p}\bigg)$$
        利用三角函数性质可将结果化简为
        $$a_n=(-1)^n\frac{\sin(p\pi)}{\pi}\bigg(\frac{1}{n+p}-\frac{1}{n-p}\bigg)$$

        代入$n=0$，由正弦函数为奇函数可知
        $$a_0=\frac{2}{\pi}\cdot\frac{\sin(p\pi)}{p}$$
        于是有傅里叶级数(这里提出公共的系数)
        $$\cos(px)\sim\frac{\sin(p\pi)}{\pi}\bigg(\frac{1}{p}+\suminf(-1)^n\bigg(\frac{1}{n+p}-\frac{1}{n-p}\bigg)\cos(nx)\bigg)$$

        \item \textbf{和函数计算}
        
        利用周期性可以发现$f(-\pi)=\cos(p\pi)=\cos(p(-\pi))$，由此$f(x)=\cos(p\pi)$在$[-\pi,\pi]$恒成立，它在$\rb$上连续。又由于$f(x)$在$[-\pi,\pi]$分段单调，利用傅里叶级数\textbf{逐点收敛结论}即得
        $$\forall x\in\rb,\quad\frac{\sin(p\pi)}{\pi}\bigg(\frac{1}{p}+\suminf(-1)^n\bigg(\frac{1}{n+p}-\frac{1}{n-p}\bigg)\cos(nx)\bigg)=f(x)$$

        \item \textbf{级数计算}
        
        为了配凑出题目要求的形式，代入$x=0$可知
        $$\frac{\sin(p\pi)}{\pi}\bigg(\frac{1}{p}+\suminf(-1)^n\bigg(\frac{1}{n+p}-\frac{1}{n-p}\bigg)\bigg)=f(0)=1$$
        令$t=p\pi$，化简即得到
        $$\sin t\bigg(\frac{1}{t}+\suminf(-1)^n\bigg(\frac{1}{n\pi+t}-\frac{1}{n\pi-t}\bigg)\bigg)=\sin t\bigg(\frac{1}{t}+\suminf(-1)^n\bigg(\frac{1}{t+n\pi}+\frac{1}{t-n\pi}\bigg)\bigg)=1$$
        由$p$不为整数可知$\sin t\ne0$，两边同除以$\frac{1}{\sin t}$即得结论。

        \item \textbf{积分计算-变形与分析}
        
        记$g(x)=\frac{\sin x}{x}$。为了能够计算积分，我们在已证明的级数两边同乘$\sin t$得到
        $$\frac{\sin t}{t}+\suminf(-1)^n\bigg(\frac{\sin t}{t+n\pi}+\frac{\sin t}{t-n\pi}\bigg)=1$$
        可以发现，由于$(-1)^n\sin t=\sin(t-n\pi)=\sin(t+n\pi)$，我们可以将上式凑成与结果更近的形式
        $$\frac{\sin t}{t}+\suminf\bigg(\frac{\sin(t+n\pi)}{t+n\pi}+\frac{\sin(t-n\pi)}{t-n\pi}\bigg)=1$$
        也即
        $$g(t)+\suminf(g(t+n\pi)+g(t-n\pi))=1$$
        由于上述求和中出现了所有$t+k\pi$，$k$为任意整数，我们可以\textbf{形式上}写为(注意此式实际上是\textbf{没有意义}的，只能作为分析，因为无穷求和需要指定顺序，不可能从$-\infty$开始)
        $$\sum_{k=-\infty}^{+\infty}g(t+k\pi)=1$$
        并进行如下的分析：

        由于上式对$t\in(0,\pi)$均成立，对它在0到$\pi$积分得到
        $$\int_0^\pi\bigg(\sum_{k=-\infty}^{+\infty}g(t+k\pi)\bigg)\dr t=\pi$$
        如果能够\textbf{交换积分与求和的次序}，我们可以得到
        $$\sum_{k=-\infty}^{+\infty}\int_0^\pi g(t+k\pi)\dr t=\pi$$
        将$t+k\pi$换元为$t$进一步得到
        $$\sum_{k=-\infty}^{+\infty}\int_{k\pi}^{(k+1)\pi}g(t)\dr t=\pi$$
        由于所有这些区间恰好覆盖整个$\rb$，此求和可以改写为
        $$\int_{-\infty}^{+\infty}g(t)\dr t=\pi$$
        最终，由$g$为偶函数即得
        $$\int_0^{+\infty}g(t)\dr t=\frac{1}{2}\int_{-\infty}^{+\infty}g(t)\dr t=\frac{\pi}{2}$$
        这就是所求的结果。
    
        接下来，我们需要\textbf{严谨证明}上述过程的合理性。

        \item \textbf{积分计算-严谨性分析}
            
        首先，我们希望将$g(t+n\pi)+g(t-n\pi)$\textbf{拆分}开以方便后续说明，也即要证
        $$\suminf(g(t+n\pi)+g(t-n\pi))=\suminf g(t+n\pi)+\suminf g(t-n\pi)$$
        由于$t\ne k\pi$时
        $$\suminf g(t+n\pi)=\suminf\frac{(-1)^n\sin t}{t+n\pi},\quad\suminf g(t-n\pi)=\suminf\frac{(-1)^n\sin t}{t-n\pi}$$
        在$n\pi>|t|$时分母不变号，两级数均为\textbf{交错级数}，且此后通项绝对值$\frac{1}{n\pi\pm t}$均单调递减趋于0，于是由\textbf{莱布尼茨判别法}均收敛，这就得到了上式成立。

        补充定义$g(0)=1$，$g$即成为了$\rb$上的连续函数，由此可得到对任何$t\in\rb$有($t=k\pi$时只有一项非零，可直接算出结果)
        $$g(t)+\suminf g(t+n\pi)+\suminf g(t-n\pi)=1$$
        接下来，为了能够进行逐项积分，我们需要先验证收敛的一致性，准确来说，$\suminf g(t+n\pi)$与$\suminf g(t-n\pi)$在$t\in[0,\pi]$均\textbf{一致收敛}。

        直接估算可知对任何$n\ge2$\ (这是为了避免$n=1$、$t=\pi$时$g(t-n\pi)$的特殊情况)有
        $$\forall t\in[0,\pi],\quad|g(t+n\pi)|\le\frac{1}{t+n\pi}\le\frac{1}{n\pi}$$
        $$\forall t\in[0,\pi],\quad|g(t-n\pi)|\le\frac{1}{n\pi-t}\le\frac{1}{(n-1)\pi}$$
        且计算得此时有
        $$g(t+n\pi)=|g(t+n\pi)|\cdot(-1)^n,\quad g(t-n\pi)=|g(t-n\pi)|\cdot(-1)^{n-1}$$
        $(-1)^n$或$(-1)^{n-1}$的部分和一致有界，利用估算结果可知$|g(t\pm n\pi)|=\frac{\sin t}{n\pi\pm t}$均一致收敛于0，且$n\ge2$时随$n$增加单调递减，从而由\textbf{迪利克雷判别法}可得一致收敛。

        由于所有$g(t\pm n\pi)$均为\textbf{连续函数}，它们在$[0,\pi]$可积。在等式两边对0到$\pi$积分，并利用一致收敛性\textbf{交换积分与求和}的次序，得到
        $$\int_0^\pi g(t)\dr t+\suminf\int_0^\pi g(t+n\pi)\dr t+\suminf\int_0^\pi g(t-n\pi)\dr t=\pi$$
        将积分中的$t\pm n\pi$重新换元为$t$，可进一步得到
        $$\int_0^\pi g(t)\dr t+\suminf\int_{n\pi}^{(n+1)\pi}g(t)\dr t+\suminf\int_{-n\pi}^{-(n-1)\pi}g(t)\dr t=\pi$$
        由于$g(t)$为\textbf{偶函数}，将$t$换元为$-t$可将求和中的第三项变形
        $$\begin{aligned}\suminf\int_{-n\pi}^{-(n-1)\pi}g(t)\dr t&=\suminf\int_{(n-1)\pi}^{n\pi}g(-t)\dr t=\suminf\int_{(n-1)\pi}^{n\pi}g(t)\dr t\\ &=\int_0^\pi g(t)\dr t+\suminf[n=2]\int_{(n-1)\pi}^{n\pi}g(t)\dr t\\ &=\int_0^\pi g(t)\dr t+\suminf\int_{n\pi}^{(n+1)\pi}g(t)\dr t\end{aligned}$$
        由此，上方求和中第三项实际上等于前两项之和，且它们都等于$\suminf\int_{(n-1)\pi}^{n\pi}g(t)\dr t$，于是我们可以化简得到
        $$2\suminf\int_{(n-1)\pi}^{n\pi}g(t)\dr t=\pi$$
        也即
        $$\suminf\int_{(n-1)\pi}^{n\pi}g(t)\dr t=\frac{\pi}{2}$$
        最后，记
        $$G(x)=\int_0^xg(t)\dr t$$
        则直接计算部分和可知上方级数求和意味着
        $$\lim_{n\to\infty}G(n\pi)=\frac{\pi}{2}$$
        另一方面，利用广义积分的\textbf{定义}与\textbf{迪利克雷判别法}得
        $$\lim_{x\to+\infty}G(x)=\int_0^{+\infty}g(t)\dr t$$
        存在(可见教材11.1节例8)，根据\textbf{归结原理}，必然有
        $$\lim_{x\to+\infty}G(x)=\frac{\pi}{2}$$
        代入$G$与$g$的表达式，我们最终证明了
        $$\int_0^{+\infty}\frac{\sin t}{t}\dr t=\frac{\pi}{2}$$
    \end{itemize}
}
\end{framed}

\note 本题的严谨性分析可以看作\textbf{一致收敛和函数性质}的一个很好的应用。通过一致收敛，我们保证了\textbf{积分与求和可以交换}，从而可以计算出结果。

\subsection{积分复习题}
\subsubsection{广义积分}
\begin{framed}
\begin{exercise}\label{xlnx}
    判断如下积分是否收敛：
    $$\int_0^{1/2}\frac{\ln x}{\sqrt{x}(1-x)^2}\dr x$$
\end{exercise}
\sol{
    可以发现被积函数在$[0,\frac{1}{2}]$上只有0处无定义，其他点均有定义且连续，从而讨论0处即可。
    
    由于
    $$\lim_{x\to0^+}\ln x=-\infty,\quad\lim_{x\to0^+}\frac{1}{\sqrt{x}}=+\infty,\quad\lim_{x\to0}(1-x)^2=1$$
    利用\textbf{无穷极限性质}可知
    $$\lim_{x\to0^+}\frac{\ln x}{\sqrt{x}(1-x)^2}=-\infty$$
    从而其为瑕点。

    由于被积函数在$(0,\frac{1}{2})$恒小于0，我们利用广义积分的基本性质将原本的收敛性问题转化为\textbf{非负函数}瑕积分
    $$\int_0^{1/2}\frac{-\ln x}{\sqrt{x}(1-x)^2}\dr x$$
    的收敛性问题。

    为了估算这点的情况，我们需要利用\textbf{推广的复合函数极限}结论与\textbf{函数极限阶估算}结论(可见上学期讲义3.1第11题)得到如下的重要极限，对任何$a>0$有(换元$t=\frac{1}{x}$)
    $$\lim_{x\to0^+}x^a\ln x=\lim_{t\to+\infty}\frac{-\ln t}{t^a}=0$$
    由此，取定指数$\frac{1}{2}<c<1$，直接计算可发现(由于$(1-x)^2$在0处极限为1，利用\textbf{乘除极限}结论可替换为1)
    $$\lim_{x\to0^+}\frac{\frac{-\ln x}{\sqrt{x}(1-x)^2}}{\frac{1}{x^c}}=\lim_{x\to0^+}-x^{c-\frac{1}{2}}\ln x=0$$
    由\textbf{幂函数瑕积分}敛散性结论可知$\int_0^{1/2}\frac{1}{x^c}\dr x$收敛，于是根据\textbf{比较判别法}得原广义积分在0处也收敛，从而收敛。
}
\end{framed}

\note 给定$a>0$，$x^a\ln x$在$0^+$处极限为0是$\ln x$常用的\textbf{放缩}方式。

\note 广义积分与数项级数有着完全类似的基本判别流程：先观察能否\textbf{直接算出}，无法算出时非负情况考虑\textbf{比较判别法}(无穷积分常用\textbf{幂函数}$\frac{1}{x^p}$与\textbf{指数函数}$q^x$，瑕积分常用\textbf{幂函数}$\frac{1}{(x-a)^p}$、$\frac{1}{(b-x)^p}$)，一般情况先判断是否\textbf{绝对收敛}，不绝对收敛则尝试\textbf{迪利克雷判别法}或\textbf{阿贝尔判别法}，还是无法进行则考虑使用\textbf{柯西收敛原理}直接估算。

\begin{framed}
\begin{exercise}
    判断以下两个积分是否收敛：
    $$\int_0^1\frac{\arctan x\cdot\ln x}{x^2}\dr x,\quad\int_1^{+\infty}\frac{\arctan x\cdot\ln x}{x^2}\dr x$$
\end{exercise}
\sol{
    我们分别讨论两个积分：
    \begin{itemize}
        \item \textbf{有界区间}
        
        可以发现被积函数在$[0,1]$上只有0处无定义，其他点均有定义且连续，从而讨论0处即可。
    
        利用等价无穷小替换可知
        $$\lim_{x\to0^+}\ln x=-\infty,\quad\lim_{x\to0^+}\frac{\arctan x}{x^2}=\lim_{x\to0^+}\frac{x}{x^2}=+\infty$$
        利用\textbf{无穷极限性质}可知
        $$\lim_{x\to0^+}\frac{\arctan x\cdot\ln x}{x^2}=-\infty$$
        从而其为瑕点。

        由于被积函数在$(0,1)$恒小于0，我们利用广义积分的基本性质将原本的收敛性问题转化为\textbf{非负函数}瑕积分
        $$\int_0^1\frac{-\arctan x\cdot\ln x}{x^2}\dr x$$
        的收敛性问题。

        下面对被积函数在0处的阶进行估算。直接利用\textbf{等价无穷小替换}可知
        $$\lim_{x\to0^+}\frac{\frac{-\arctan x\cdot\ln x}{x^2}}{\frac{-\ln x}{x}}=\lim_{x\to0^+}\frac{\arctan x}{x}=1$$
        且$\frac{-\ln x}{x}$在$(0,1)$恒正，由此利用\textbf{比较判别法}得原广义积分敛散性与
        $$\int_0^1\frac{-\ln x}{x}\dr x$$
        相同。

        直接计算可知(由$x\to0^+$时$\ln x\to-\infty$，$\ln^2x\to+\infty$)
        $$\int_0^1-\frac{\ln x}{x}\dr x=\int_0^1-(\ln x)\dr(\ln x)=\bigg(-\frac{1}{2}\ln^2x\bigg)\bigg|^1_0=+\infty$$
        由此原广义积分发散。

        \item \textbf{无穷区间}
        
        可以发现被积函数在$[1,+\infty)$连续，从而讨论$+\infty$处即可。由于被积函数在$(1,+\infty)$恒大于0，这是一个\textbf{非负函数}无穷积分。

        估算得$\frac{\arctan x}{x^2}$在$+\infty$处与$\frac{1}{x^2}$同阶，超过了临界情况$\frac{1}{x}$，因此可以考虑适当\textbf{放大}$\ln x$到某个幂函数。取定指数$1<c<2$，由\textbf{函数极限阶估算}结论(可见上学期讲义3.1第11题)可知(由于$\arctan x$在$+\infty$处极限为$\frac{\pi}{2}$，利用\textbf{乘除极限}结论可替换为$\frac{\pi}{2}$)
        $$\lim_{x\to+\infty}\frac{\frac{\arctan x\cdot\ln x}{x^2}}{\frac{1}{x^c}}=\lim_{x\to+\infty}\frac{\pi}{2}\cdot\frac{\ln x}{x^{2-c}}=0$$
        由\textbf{幂函数瑕积分}敛散性结论可知$\int_1^{+\infty}\frac{1}{x^c}\dr x$收敛，于是根据\textbf{比较判别法}得原广义积分在$+\infty$处也收敛，从而收敛。
    \end{itemize}
}
\end{framed}

\note 由于瑕积分在一点处往往趋于正无穷或负无穷，利用极限\textbf{保序性}可知局部恒正或恒负，可以当作非负函数处理，从而\textbf{几乎只需考虑比较判别法}。

\note 若一点处并不趋于正无穷或负无穷，一般建议将其换元为\textbf{无穷积分}再处理，如分析$\int_0^1x^p\sin\frac{1}{x}\dr x$的收敛性$(p\le0)$可以换元$t=\frac{1}{x}$将它转化为$\int_1^{+\infty}-t^{-p-2}\sin t\dr t$。

\begin{framed}
\begin{exercise}
    判断如下积分是否收敛：
    $$\int_0^{+\infty}\frac{\sin x}{\sqrt{x}}\arctan x\dr x,\quad\int_0^{+\infty}\frac{\sin x}{\sqrt{x}+\sin x}\arctan x\dr x$$
\end{exercise}
\sol{
    我们分别讨论两个积分：
    \begin{itemize}
        \item \textbf{分母为幂函数}
        
        可以发现被积函数$[0,+\infty)$上只有0处无定义，其他点均有定义且连续，从而需要讨论0处与$+\infty$处。

        直接利用\textbf{等价无穷小替换}计算可知0处有
        $$\lim_{x\to0^+}\frac{\sin x}{\sqrt{x}}\arctan x=\lim_{x\to0^+}\frac{x^2}{\sqrt{x}}=0$$
        由此补充定义0处为0即连续，这点\textbf{不为瑕点}。

        在$+\infty$处，由于被积函数正负交替，尝试使用\textbf{迪利克雷判别法}，将它拆分为
        $$\int_0^{+\infty}\frac{\arctan x}{\sqrt{x}}\cdot\sin x$$

        第一部分：直接求导可知$x>0$时
        $$\frac{\dr}{\dr x}\frac{\arctan x}{\sqrt{x}}=\frac{\frac{\sqrt{x}}{x^2+1}-\frac{\arctan x}{2\sqrt{x}}}{x}=\frac{\frac{x}{x^2+1}-\frac{\arctan x}{2}}{x\sqrt{x}}$$
        由于分子满足
        $$\lim_{x\to+\infty}\bigg(\frac{x}{x^2+1}-\frac{\arctan x}{2}\bigg)=0-\frac{\pi}{4}=-\frac{\pi}{4}<0$$
        利用极限\textbf{保序性}可知存在$M>0$使得$x>M$时其小于0，而$x\sqrt{x}>0$恒成立，从而$x>M$时$\frac{\arctan x}{\sqrt{x}}$单调递减，且
        $$\lim_{x\to+\infty}\frac{\arctan x}{\sqrt{x}}=\frac{\pi}{2}\cdot0=0$$
        于是$x>M$时其\textbf{单调递减趋于0}。

        第二部分：直接计算可得$x>M$时
        $$\bigg|\int_M^x\sin t\dr t\bigg|=|\cos M-\cos x|\le2$$
        于是$x>M$时其\textbf{积分有界}。

        利用广义积分的基本性质，\textbf{去除正常积分不影响广义积分敛散性}，因此
        $$\int_0^{+\infty}\frac{\arctan x}{\sqrt{x}}\sin x,\quad\int_M^{+\infty}\frac{\arctan x}{\sqrt{x}}\sin x$$
        敛散性相同。由迪利克雷判别法已经得到后者收敛，因此前者也收敛。

        \item \textbf{分母不为幂函数} 
        
        可以发现被积函数$[0,+\infty)$上只有0处无定义，其他点均有定义且连续，从而需要讨论0处与$+\infty$处。

        直接利用\textbf{等价无穷小替换}计算可知0处有(由于$x\to0$时$\sin x\sim x$，它在0处为$o(\sqrt{x})$，因此$\sqrt{x}$与$\sqrt{x}+\sin x$等价)
        $$\lim_{x\to0^+}\frac{\sin x}{\sqrt{x}+\sin x}\arctan x=\lim_{x\to0^+}\frac{x^2}{\sqrt{x}+o(\sqrt{x})}=\lim_{x\to0^+}\frac{x^2}{\sqrt{x}}=0$$
        由此补充定义0处为0即连续，这点\textbf{不为瑕点}，只需考虑$+\infty$处。

        求导可以发现$\frac{\arctan x}{\sqrt{x}+\sin x}$并不单调递减，由此\textbf{难以直接估算}。不过，利用广义积分的基本性质(收敛与收敛之和收敛、收敛与发散之和发散)，\textbf{被积函数减去收敛积分的被积函数并不影响敛散性}，由此我们可以利用第一部分证明将问题转化为研究
        $$\int_0^{+\infty}\bigg(\frac{\sin x}{\sqrt{x}+\sin x}\arctan x-\frac{\sin x}{\sqrt{x}}\arctan x\bigg)\dr x$$
        的敛散性。

        计算差并化简可得其为(等号运用了广义积分的基本性质)
        $$\int_0^{+\infty}-\frac{\sin^2x\cdot\arctan x}{\sqrt{x}(\sqrt{x}+\sin x)}=-\int_0^{+\infty}\frac{\sin^2x\cdot\arctan x}{\sqrt{x}(\sqrt{x}+\sin x)}$$
        估算可以发现右侧的被积函数在$(0,+\infty)$恒正，这是一个\textbf{非负函数}无穷积分。

        设
        $$g(x)=\frac{\sin^2x\cdot\arctan x}{\sqrt{x}(\sqrt{x}+\sin x)},\quad G(y)=\int_0^yg(x)\dr x$$
        由于$\sin x$以$2\pi$为周期，我们尝试将积分以$2\pi$分组，写出
        $$G(2\pi n)=\sum_{k=0}^{n-1}\int_{2k\pi}^{2(k+1)\pi}g(x)\dr x$$
        由于$g(x)$非负，它在0到$2\pi$的积分大于等于0，而当$k\ge1$时将$\sin^2x$以外的项\textbf{放缩至端点}有(将分子缩小至最小，分母的每项放大至最大)
        $$\forall x\in[2k\pi,2(k+1)\pi],\quad\frac{\sin^2x\cdot\arctan x}{\sqrt{x}(\sqrt{x}+\sin x)}\ge\frac{\sin^2x\cdot\arctan(2k\pi)}{\sqrt{2(k+1)\pi}(\sqrt{2(k+1)\pi}+1)}$$
        两侧同积分即得对正整数$k$有($\sin^2x$在$[2k\pi,2(k+1)\pi]$的积分为$\pi$，计算可见\exref{sin 2 and 4 to 2pi})
        $$\int_{2k\pi}^{2(k+1)\pi}g(x)\dr x\ge\frac{\arctan(2k\pi)}{\sqrt{2(k+1)\pi}(\sqrt{2(k+1)\pi}+1)}\pi$$
        因此
        $$G(2\pi n)\ge\sum_{k=1}^{n-1}\frac{\arctan(2k\pi)}{\sqrt{2(k+1)\pi}(\sqrt{2(k+1)\pi}+1)}\pi$$
        此时右侧已经是一个\textbf{正项级数}，我们可以直接进行\textbf{阶估算}：观察可发现它的分母为1次，分子为0次，由此可以与$\frac{1}{k}$比较，利用\textbf{归结原理}与\textbf{初等函数连续性}、$\arctan x$极限结论计算可知
        $$\begin{aligned}\lim_{k\to\infty}\frac{\frac{\arctan(2k\pi)}{\sqrt{2(k+1)\pi}(\sqrt{2(k+1)\pi}+1)}\pi}{\frac{1}{k}}&=\lim_{k\to\infty}\arctan(2k\pi)\cdot\frac{\pi}{\sqrt{2(\frac{1}{k}+1)\pi}(\sqrt{2(\frac{1}{k}+1)\pi}+\frac{1}{\sqrt{k}})}\\ &=\lim_{x\to0^+}\arctan\frac{2\pi}{x}\cdot\frac{\pi}{\sqrt{2(x+1)\pi}(\sqrt{2(x+1)\pi}+\sqrt{x})}\\ &=\frac{\pi}{2}\cdot\frac{\pi}{\sqrt{2\pi}\cdot\sqrt{2\pi}}\\ &=\frac{\pi}{4}\end{aligned}$$
        由\textbf{$p$级数}敛散性结论可知$\suminf\frac{1}{n}$发散，于是根据\textbf{比较判别法}得(正项级数发散即为求和正无穷)
        $$\suminf[k=1]\frac{\arctan(2k\pi)}{\sqrt{2(k+1)\pi}(\sqrt{2(k+1)\pi}+1)}\pi=+\infty$$
        由此即得
        $$\lim_{n\to\infty}G(2\pi n)=+\infty$$
        若$G(x)$在$x\to\infty$时极限存在，利用\textbf{归结原理}可知$G(2\pi n)$应趋于同一极限，矛盾，由此最终得到原广义积分发散。
    \end{itemize}
}
\end{framed}

\note 类似数项级数去除有限项不影响敛散性，广义积分\textbf{去除正常积分不影响敛散性}。

\note 广义积分的很多技巧都与数项级数\textbf{完全类似}，例如阶估算、与收敛积分作差等。对于含周期函数的广义积分，可以考虑\textbf{分段估算}方法得到\textbf{数项级数}后判断，更多严谨性讨论可见\exref{func to series}。

\begin{framed}
\begin{exercise}
    判断如下积分是否收敛：
    $$\int_1^{+\infty}\frac{\ln x\cdot\ln(x+1)}{\sqrt{x^p-1}}\dr x$$
    这里参数$p>0$。
\end{exercise}
\sol{
    可以发现被积函数在$[1,+\infty)$上只有1处无定义，其他点均有定义且连续，从而需要讨论1处与$+\infty$处。分别讨论后综合结果：
    \begin{itemize}
        \item \textbf{1处收敛性}
        
        利用基本等价无穷小$t\to0$时$\ln(1+t)\sim t$、$(t+1)^p-1\sim pt$，我们可以直接计算得到(换元$t=x-1$以\textbf{移至0处}，利用\textbf{乘除极限}结论将$\ln(2+t)$替换为$\ln2$)
        $$\lim_{x\to1^+}\frac{\ln x\ln(1+x)}{\sqrt{(x+1)^p-1}}=\lim_{t\to0^+}\frac{\ln(1+t)\ln(2+t)}{\sqrt{(1+t)^p-1}}=\lim_{t\to0^+}\frac{t\ln 2}{\sqrt{pt}}=0$$
        由此补充定义1处为0即连续，这点\textbf{不为瑕点}。
        
        \item \textbf{无穷处收敛性-转化}

        由于被积函数在$(1,+\infty)$恒大于0，这是一个\textbf{非负函数}无穷积分，从而可以直接进行阶估算。当$x\to+\infty$时，它接近$\frac{\ln^2x}{x^{p/2}}$，因此我们先进行转化：直接计算得
        $$\lim_{x\to+\infty}\frac{\frac{\ln x\ln(1+x)}{\sqrt{(x+1)^p-1}}}{\frac{\ln^2x}{x^{p/2}}}=\lim_{x\to\infty}\frac{\ln(1+x)}{\ln x}\cdot\frac{1}{\sqrt{(1+x^{-1})^p-x^{-p}}}$$
        第二部分利用\textbf{初等函数连续性}与\textbf{推广的复合函数极限}结论即得极限为$\frac{1}{\sqrt{(1+0)^p-0^p}}=1$。对第一部分，可先计算
        $$\lim_{x\to+\infty}\frac{\ln(1+x)-\ln x}{\ln x}=\lim_{x\to+\infty}\frac{\ln(1+\frac{1}{x})}{\ln x}$$
        由于分子极限为0、分母极限为$+\infty$，利用无穷极限性质可知上述极限为0，由此第一部分极限为1。

        综合以上讨论，我们最终得到了
        $$\lim_{x\to+\infty}\frac{\frac{\ln x\ln(1+x)}{\sqrt{(x+1)^p-1}}}{\frac{\ln^2x}{x^{p/2}}}=1$$
        由此根据\textbf{比较判别法}得原广义积分敛散性与
        $$\int_1^{+\infty}\frac{\ln^2x}{x^{p/2}}\dr x$$
        相同。

        \item \textbf{无穷处收敛性-分类}
        
        我们按照幂函数的临界情况分类，也即$p>2$与$p\le2$分别讨论(严格来说$p=2$的情况应单独讨论，但本题可以合并)。

        当$p\le2$时，由于
        $$\lim_{x\to+\infty}\frac{\frac{\ln^2x}{x^{p/2}}}{\frac{1}{x^{p/2}}}=\lim_{x\to+\infty}\ln^2x=+\infty$$
        由\textbf{幂函数无穷积分}敛散性结论可知$\int_1^{+\infty}\frac{1}{x^{p/2}}\dr x$发散，于是根据\textbf{比较判别法}得原广义积分在$+\infty$处也发散。

        当$p>2$时，取定指数$1<c<\frac{p}{2}$，由\textbf{函数极限阶估算}结论(可见上学期讲义3.1第11题)可知
        $$\lim_{x\to+\infty}\frac{\frac{\ln^2x}{x^{p/2}}}{\frac{1}{x^c}}=\lim_{x\to+\infty}\frac{\ln^2x}{x^{p/2-c}}=0$$
        由\textbf{幂函数无穷积分}敛散性结论可知$\int_1^{+\infty}\frac{1}{x^c}\dr x$收敛，于是根据\textbf{比较判别法}得原广义积分在$+\infty$处也收敛。

        \item \textbf{综合结果}

        综合以上讨论，原广义积分收敛当且仅当$p>2$。
    \end{itemize}
}
\end{framed}

\note 注意$\ln(1+x)$在$x\to+\infty$时与$\ln x$而非$x$等价，\textbf{等价无穷小结论对无穷大不适用}。

\note 同时含有$\ln x$与$x$的广义积分的收敛性讨论往往与$\suminf[n=2]\frac{1}{n^p\ln^qn}$的收敛性讨论\textbf{类似}。

\begin{framed}
\begin{exercise}
    判断如下积分是否收敛：
    $$\int_0^{+\infty}\frac{x^n(1-\er^{-x})}{(1+x)^m}\dr x$$
    这里参数$m$、$n$为任意实数。
\end{exercise}
\sol{
    可以发现被积函数$[0,+\infty)$上只有0处可能无定义，其他点均有定义且连续，从而需要讨论0处与$+\infty$处。分别讨论：
    \begin{itemize}
        \item \textbf{0处收敛性}
        
        若$n\ge0$，被积函数在0处连续，0附近积分收敛。
        
        若$-1\le n<0$，利用\textbf{等价无穷小替换}($x\to0$时有$1-\er^{-x}\sim-(-x)=x$)与乘除极限结论有
        $$\lim_{x\to0^+}\frac{x^n(1-\er^{-x})}{(1+x)^m}=\lim_{x\to0^+}x^{n+1}\cdot\frac{1}{(1+x)^m}=\begin{cases}0&n<-1\\1&n=-1\end{cases}$$
        由此补充定义0处为0或1即连续，这点\textbf{不为瑕点}。

        若$n<-1$，原广义积分在0附近的收敛性可以等价于
        $$\int_0^1\frac{x^n(1-\er^{-x})}{(1+x)^m}\dr x$$
        的收敛性。由于被积函数在$(0,1)$恒大于0，这是一个\textbf{非负函数}瑕积分。直接进行阶估算：利用\textbf{等价无穷小替换}与乘除极限结论可知
        $$\lim_{x\to0^+}\frac{\frac{x^n(1-\er^{-x})}{(1+x)^m}}{x^{n+1}}=\lim_{x\to0^+}\frac{1-\er^{-x}}{x}\cdot\frac{1}{(1+x)^m}=1$$
        由\textbf{幂函数瑕积分}敛散性结论可知$\int_0^1x^{n+1}\dr x$收敛当且仅当$-1<n+1<0$，于是根据\textbf{比较判别法}得原广义积分为瑕积分时在0处收敛当且仅当$-2<n<-1$。
        
        综合以上的各种情况，原广义积分在0处收敛当且仅当$n>-2$。
        
        \item \textbf{无穷处收敛性}
        
        为了避免0处的影响，我们将原广义积分在无穷处的收敛性等价于
        $$\int_1^{+\infty}\frac{x^n(1-\er^{-x})}{(1+x)^m}\dr x$$
        的收敛性。由于被积函数在$(1,+\infty)$恒大于0，这是一个\textbf{非负函数}无穷积分。
        
        直接进行阶估算：利用\textbf{等价无穷小替换}、初等函数连续性与乘除极限结论可知
        $$\lim_{x\to+\infty}\frac{\frac{x^n(1-\er^{-x})}{(1+x)^m}}{x^{n-m}}=\lim_{x\to+\infty}(1-\er^{-x})\cdot\frac{1}{\big(\frac{1}{x}+1\big)^m}=1\cdot1=1$$
        由\textbf{幂函数广义积分}敛散性结论可知$\int_0^1x^{n-m}\dr x$收敛当且仅当$n-m<-1$，于是根据\textbf{比较判别法}得原广义积分在无穷处收敛当且仅当$n-m<-1$。

        \item \textbf{综合结果}

        综合以上讨论，原广义积分收敛当且仅当$n>-2$且$n-m<-1$。
    \end{itemize}
}
\end{framed}

\note 当出现多个瑕点或无穷点时，基本的做法是\textbf{分别讨论}，收敛当且仅当其中\textbf{每一部分均收敛}。注意幂级数瑕积分与无穷积分收敛的\textbf{不同指数范围}。

\begin{framed}
\begin{exercise}
    设
    $$\theta(x)=\int_0^x\sqrt{(t+1)(t+2)(t+3)}\dr t$$
    证明
    $$\int_0^{+\infty}\cos(\theta(x))\dr x$$
    收敛。
\end{exercise}
\sol{
    利用\textbf{变上限积分}性质可知$\theta(x)$在$[0,+\infty)$连续，从而讨论$+\infty$处即可。不过，此形式难以看出收敛性，我们先进行分析再证明：
    \begin{itemize}
        \item \textbf{分析与变形}
        
        利用变上限积分的性质，在$(0,+\infty)$有
        $$\theta'(x)=\sqrt{(x+1)(x+2)(x+3)}>0$$
        且
        $$\theta(x)\ge\int_0^xt^{3/2}\dr t=\frac{2}{5}x^{5/2}$$
        由此可发现$\theta(x)$单调递增趋于无穷。

        可以联想到，证明$\int_0^{+\infty}\sin x^2\dr x$收敛性时，我们考虑换元$t=x^2$，这样积分化为
        $$\int_0^{+\infty}\frac{\sin t}{2\sqrt{t}}\dr t$$
        就可以使用\textbf{迪利克雷判别法}判别了。

        本题我们希望能进行相同的操作，因此换元$t=\theta(x)$。由于$\theta(x)$严格单调增且连续，它存在\textbf{严格单调增且连续的反函数}$x=\psi(t)$。利用反函数求导公式可知
        $$\psi'(t)=\frac{\dr x}{\dr t}=\bigg(\frac{\dr t}{\dr x}\bigg)^{-1}=\frac{1}{\sqrt{(x+1)(x+2)(x+3)}}=\frac{1}{\sqrt{(\psi(t)+1)(\psi(t)+2)(\psi(t)+3)}}$$
        
        由于$\theta(0)=0$，且
        $$\lim_{x\to+\infty}\theta(x)=+\infty$$
        换元后$t$的范围也为$[0,+\infty)$，从而换元后得到新的积分
        $$\int_0^{+\infty}(\cos t)\cdot\psi'(t)\dr t=\int_0^{+\infty}\frac{\cos t}{\sqrt{(\psi(t)+1)(\psi(t)+2)(\psi(t)+3)}}\dr t$$

        \item \textbf{收敛性证明}

        换元后，由于$\psi(t)$在$[0,+\infty)$严格单调增且连续，$\psi(0)=0$，仍然只需要讨论$+\infty$处。由于被积函数正负交替，尝试将它拆分为
        $$\int_0^{+\infty}\frac{1}{\sqrt{(\psi(t)+1)(\psi(t)+2)(\psi(t)+3)}}\cdot\cos t\dr t$$
        并使用\textbf{迪利克雷判别法}。

        第一部分：之前的讨论中已经证明了$\psi(t)$严格单调增，从而其单调递减。为了证明其趋于0，我们还需要证明
        $$\lim_{t\to+\infty}\psi(t)=+\infty$$
        若否，由于其单调增，不趋于无穷可以说明其\textbf{有界}。设它小于等于$M$，则利用$\theta(x)$\textbf{严格单调增}与反函数的\textbf{定义}有
        $$t=\theta(\psi(t))\le\theta(M)$$
        但上式不可能对任何$t\in[0,+\infty)$成立，矛盾。由此，再通过\textbf{推广的复合函数极限}结论即可说明第一部分单调递减趋于0。

        第二部分：直接计算可得
        $$\bigg|\int_0^x\cos t\dr t\bigg|=|\sin x|\le1$$
        于是其\textbf{积分有界}。

        由此，利用迪利克雷判别法最终得到换元后的积分收敛，因此原积分收敛。
    \end{itemize}

    \note 本题也可以考虑将积分看作
    $$\int_0^{+\infty}\frac{1}{\theta'(x)}\cdot\cos(\theta(x))\theta'(x)\dr x$$
    第一部分单调递减趋于0，第二部分直接计算可知在$[0,m]$的积分有界，从而通过\textbf{迪利克雷判别法}可得结果。
}
\end{framed}

\note 对广义积分，必须保证换元后的变量对$x$是\textbf{严格单调}的才能进行换元。

\begin{framed}
\begin{exercise}
    判断如下积分是否收敛：
    $$\int_0^{+\infty}\frac{1}{(1+x^2)(\sin^2x)^\alpha}\dr x$$
    这里参数$\alpha>0$。
\end{exercise}
\sol{
    可以发现被积函数$[0,+\infty)$上在所有$k\pi$\ ($k$为自然数)处无定义，其他点均有定义且连续，从而需要讨论$k\pi$处与$+\infty$处。分别讨论：
    \begin{itemize}
        \item \textbf{$k\pi$处收敛性-确定范围}
        
        由于被积函数在有定义的点恒正，在$k\pi$处为\textbf{非负函数}的瑕积分。

        先考虑0处，为了避免其他瑕点的影响，将原广义积分在0处的收敛性等价于
        $$\int_0^{\pi/2}\frac{1}{(1+x^2)(\sin^2x)^\alpha}\dr x$$
        的收敛性。

        利用\textbf{等价无穷小替换}直接计算可知
        $$\lim_{x\to0^+}\frac{\frac{1}{(1+x^2)(\sin^2x)^\alpha}}{\frac{1}{x^{2\alpha}}}=\lim_{x\to0^+}\bigg(\frac{x}{\sin x}\bigg)^{2\alpha}\cdot\frac{1}{1+x^2}=1$$
        由此根据\textbf{比较判别法}得原广义积分在0处的敛散性与
        $$\int_0^{\pi/2}\frac{1}{x^{2\alpha}}\dr x$$
        相同，从而由\textbf{幂函数瑕积分}敛散性结论可知0处收敛当且仅当$\alpha<\frac{1}{2}$。

        若$\alpha\ge\frac{1}{2}$，原广义积分在0处发散，因此已经发散，我们下面证明$\alpha<\frac{1}{2}$时它在任何$k\pi$处均收敛。

        \item \textbf{$k\pi$处收敛性-收敛证明}
        
        在任何$k\pi$\ ($k$为正整数)处，我们需要分别考虑左侧与右侧。对左侧，为了避免其他瑕点的影响，原广义积分的收敛性等价于
        $$\int_{k\pi-\pi/2}^{k\pi}\frac{1}{(1+x^2)(\sin^2x)^\alpha}\dr x$$
        的收敛性。

        换元$t=k\pi-x$，利用\textbf{等价无穷小替换}直接计算可知
        $$\lim_{x\to(k\pi)^-}\frac{\frac{1}{(1+x^2)(\sin^2x)^\alpha}}{\frac{1}{(k\pi-x)^{2\alpha}}}=\lim_{t\to0^+}\bigg(\frac{t}{\sin t}\bigg)^{2\alpha}\cdot\frac{1}{1+(k\pi-t)^2}=\frac{1}{1+k^2\pi^2}$$
        由\textbf{幂函数瑕积分}敛散性结论可知$\int_{k\pi-\pi/2}^{k\pi}\frac{1}{(k\pi-x)^{2\alpha}}\dr x$收敛，于是根据\textbf{比较判别法}得原广义积分在$k\pi$左侧也收敛。

        对右侧，为了避免其他瑕点的影响，原广义积分的收敛性等价于
        $$\int_{k\pi}^{k\pi+\pi/2}\frac{1}{(1+x^2)(\sin^2x)^\alpha}\dr x$$
        的收敛性。

        换元$t=x-k\pi$，利用\textbf{等价无穷小替换}直接计算可知
        $$\lim_{x\to(k\pi)^+}\frac{\frac{1}{(1+x^2)(\sin^2x)^\alpha}}{\frac{1}{(k\pi+x)^{2\alpha}}}=\lim_{t\to0^+}\bigg(\frac{t}{\sin t}\bigg)^{2\alpha}\cdot\frac{1}{1+(k\pi+t)^2}=\frac{1}{1+k^2\pi^2}$$
        由\textbf{幂函数瑕积分}敛散性结论可知$\int_{k\pi}^{k\pi+\pi/2}\frac{1}{(x-k\pi)^{2\alpha}}\dr x$收敛，于是根据\textbf{比较判别法}得原广义积分在$k\pi$右侧也收敛。
        
        \item \textbf{无穷处收敛性}
        
        最后，我们将证明$\alpha<\frac{1}{2}$原积分在无穷处也收敛。记
        $$f(x)=\frac{1}{(1+x^2)(\sin^2x)^\alpha},\quad F(y)=\int_0^yf(x)\dr x$$
        由于$\sin^2x$以$\pi$为周期，我们尝试将积分以$\pi$分组，写出
        $$F(n\pi)=\sum_{k=0}^{n-1}\int_{k\pi}^{(k+1)\pi}f(x)\dr x$$
        由于$f(x)$非负，它在0到$2\pi$的积分大于等于0，将$\sin^2x$以外的项\textbf{放缩至端点}有
        $$\forall x\in[k\pi,(k+1)\pi],\quad\frac{1}{(1+x^2)(\sin^2x)^\alpha}\le\frac{1}{1+k^2}\cdot\frac{1}{(\sin^2x)^\alpha}$$
        两侧同积分即得对自然数$k$有
        $$\int_{k\pi}^{(k+1)\pi}g(x)\dr x\le\frac{1}{k^2+1}\int_{k\pi}^{(k+1)\pi}\frac{1}{(\sin^2x)^\alpha}=\frac{1}{k^2+1}\int_0^\pi\frac{1}{(\sin^2x)^\alpha}$$
        这里瑕积分的收敛性证明与上一部分完全类似。由此，我们可以得到估算
        $$F(n\pi)\le\int_0^\pi\frac{1}{(\sin^2x)^\alpha}\sum_{k=0}^{n-1}\frac{1}{k^2+1}$$
        由$k$为正整数时$\frac{1}{k^2+1}\le\frac{1}{k^2}$，利用\textbf{$p$级数}敛散性结论与\textbf{比较判别法}即可得到右侧求和收敛，从而$F(n\pi)$有界。由$f(x)$非负可知$F(n\pi)$对$n$单调增，因此存在$A$使得
        $$\lim_{n\to\infty}F(n\pi)=A$$
        利用\exref{func to series}的结论(逆命题在$f(x)$非负时成立)，这已经足以推出
        $$\lim_{x\to+\infty}F(x)=A$$
        从而原广义积分收敛。

        \note 此结论在较复杂的题目中往往可以\textbf{直接使用}，它等价于非负函数广义积分\textbf{有界则收敛}。

        \note 事实上，这里的无穷处收敛性已经代表了某种\textbf{整体收敛性}，因为其中包含了无穷多个瑕点。

        \item \textbf{综合结果}
        
        综合以上讨论，原广义积分收敛当且仅当$\alpha<\frac{1}{2}$。
    \end{itemize}
}
\end{framed}

\note 非负函数广义积分有界则收敛基本可以看作\textbf{函数极限的单调有界定理}。

\note 出现多个瑕点与无穷点时，需要结合每个位置的实际情况构造\textbf{合适的放缩}以证明结论。

\begin{framed}
\begin{exercise}
    计算
    $$\int_0^{+\infty}\sqrt{x^3}\er^{-x}\dr x,\quad\int_0^1\sqrt{\frac{x^3}{1-x}}\dr x$$
\end{exercise}
\sol{
    我们分别计算：
    \begin{itemize}
        \item \textbf{第一个积分}
        
        根据$\Gamma$函数定义可知此积分即为$\Gamma(\frac{5}{2})$，利用$\Gamma$函数性质$\Gamma(\alpha+1)=\alpha\Gamma(\alpha)$与$\Gamma(\frac{1}{2})=\sqrt\pi$有
        $$\Gamma\bigg(\frac{5}{2}\bigg)=\frac{3}{2}\Gamma\bigg(\frac{3}{2}\bigg)=\frac{3}{2}\cdot\frac{1}{2}\Gamma\bigg(\frac{1}{2}\bigg)=\frac{3}{4}\sqrt\pi$$

        \item \textbf{第二个积分}

        根据$B$函数定义可知此积分为
        $$\int_0^1x^{3/2}(1-x)^{-1/2}\dr x=B\bigg(\frac{5}{2},\frac{1}{2}\bigg)$$
        由$B$函数的性质可知其为($\Gamma(n+1)=n!$，而分子的两项已在上一部分算出)
        $$\frac{\Gamma(\frac{5}{2})\Gamma(\frac{1}{2})}{\Gamma(3)}=\frac{\frac{3}{4}\sqrt\pi\cdot\sqrt\pi}{2}=\frac{3}{8}\pi$$
    \end{itemize}
}
\end{framed}

\note 对于$\Gamma$函数与$B$函数，需要记忆基本公式：
$$\Gamma(n+1)=n!,\quad\Gamma\bigg(\frac{1}{2}\bigg)=\sqrt{\pi},\quad\Gamma(\alpha+1)=\alpha\Gamma(\alpha),\quad B(p,q)=\frac{\Gamma(p)\Gamma(q)}{\Gamma(p+q)}$$
其他相关问题需要应用\textbf{积分变换}转化为符合要求的形式，如最常见的将\textbf{三角函数}换元转化为$B$函数得到的
$$B(p,q)=2\int_0^{\pi/2}\cos^{2p-1}\theta\sin^{2q-1}\theta\dr\theta$$

\note 注意熟悉$\Gamma$函数与$B$函数的\textbf{定义域}，分别为$(0,+\infty)$与$(0,+\infty)\times(0,+\infty)$。

\note 另一个常用但教材并未介绍的公式是\textbf{余元公式}
$$\forall p\in(0,1),\quad\Gamma(p)\Gamma(1-p)=\frac{\pi}{\sin(p\pi)}$$
它可以\textbf{辅助得到结果}，但不建议直接使用。

\subsubsection{含参积分}
\begin{framed}
\begin{exercise}
    证明含有参数$t$的无穷积分
    $$\int_1^{+\infty}t\er^{-tx}\frac{\cos x}{x}\dr x$$
    在$[0,+\infty)$一致收敛。
\end{exercise}
\sol{
    可以发现被积函数$[1,+\infty)$上有定义且连续，从而只需考虑$+\infty$处的一致收敛性。我们在给出一个估算后介绍两种方法：
    \begin{itemize}
        \item \textbf{函数估算}
        
        由于这里包含$t$的因子为$t\er^{-tx}$，在$t\in[0,+\infty)$、$x\in[1,+\infty)$时它非负，我们尝试计算给定$x$、$t$可以任取时它的\textbf{最大值}以方便估算。

        求导可知
        $$\frac{\partial}{\partial t}(t\er^{-tx})=(1-xt)\er^{-tx}$$
        由于$t\er^{-tx}$在$t=0$与$t\to+\infty$时均为0，分析可发现它在导数为0的点$t=\frac{1}{x}$时取到最大值
        $$\frac{1}{\er x}$$

        \item \textbf{迪利克雷判别法思路}
        
        将积分写为
        $$\int_1^{+\infty}\frac{t\er^{-tx}}{x}\cdot\cos x\dr x$$

        第一部分：利用形式可发现给定$t$后它对$x$单调递减，且由估算有
        $$\forall t\in[0,+\infty),\quad\bigg|\frac{t\er^{-tx}}{x}\bigg|\le\frac{1}{\er x^2}$$
        由于
        $$\lim_{x\to+\infty}\frac{1}{\er x^2}=0$$
        对任何$\varepsilon>0$，存在$M$使得$x>M$时
        $$\forall t\in[0,+\infty),\quad\bigg|\frac{t\er^{-tx}}{x}\bigg|\le\frac{1}{\er x^2}<\varepsilon$$
        这已经符合$\frac{t\er^{-tx}}{x}$在$x\to+\infty$时对$t\in[1,+\infty)$\textbf{一致收敛于0}的定义。

        第二部分：由于它与$t$无关，直接计算可知对任何$m>1$有
        $$\forall t\in[0,+\infty),\quad\bigg|\int_1^m\cos x\dr x\bigg|=|\sin m-\sin1|<2$$

        由\textbf{迪利克雷判别法}可知原含参积分一致收敛。       

        \item \textbf{阿贝尔判别法思路}
        
        将积分写为
        $$\int_1^{+\infty}t\er^{-tx}\cdot\frac{\cos x}{x}\dr x$$

        第一部分：利用形式可发现给定$t$后它对$x$单调递减，且由估算有
        $$\forall t\in[0,+\infty),\quad\forall x\in[1,+\infty),\quad\bigg|\frac{t\er^{-tx}}{x}\bigg|\le\frac{1}{\er x}\le\frac{1}{\er}$$
        从而$t\er^{-tx}$也\textbf{一致有界}。

        第二部分：由于它与$t$无关，利用广义积分的\textbf{迪利克雷判别法}($\frac{1}{x}$单调减趋于0，$\cos x$的积分与前一种方法相同得有界)可知其收敛，从而由定义它对$t\in[0,+\infty)$一致收敛。

        由\textbf{阿贝尔判别法}可知原含参积分一致收敛。 
    \end{itemize}
}
\end{framed}

\note 由于\textbf{一致收敛必然逐点收敛}，使用判别法验证区间上一致收敛后，\textbf{无需再验证固定参数时广义积分的收敛性}。

\note \textbf{直接算出}、\textbf{$M$判别法}、\textbf{迪利克雷判别法}或\textbf{阿贝尔判别法}、\textbf{柯西收敛原理}的\textbf{基本判别流程}仍然是一致的。

\note 对一般的题目，可以考虑被积函数进行\textbf{不同拆分}尝试迪利克雷判别法和阿贝尔判别法，一般来说，\textbf{迪利克雷判别法}的应用范围更广。

\begin{framed}
\begin{exercise}
    证明含有参数$x$的无穷积分
    $$\int_0^{+\infty}\er^{-tx^2}\cos t\dr t$$
    对任何正数$r$在$[r,+\infty)$一致收敛，在$(0,+\infty)$不一致收敛。
\end{exercise}
\sol{
    可以发现被积函数在$[0,+\infty)$上有定义且连续，从而只需考虑$+\infty$处的一致收敛性。分两种情况证明：
    \begin{itemize}
        \item \textbf{一致收敛证明}
        
        由于被积函数自然分为
        $$\int_0^{+\infty}\er^{-tx^2}\cdot\cos t\dr t$$
        可以考虑使用\textbf{迪利克雷判别法}。

        第一部分：利用形式可发现给定$t$后它对$x$单调递减，且
        $$\forall x\in[r,+\infty),\quad|\er^{-tx^2}|\le\er^{-r^2t}$$
        由于
        $$\lim_{t\to+\infty}\er^{-r^2t}=0$$
        对任何$\varepsilon>0$，存在$M$使得$t>M$时
        $$\forall x\in[r,+\infty),\quad|\er^{-tx^2}|\le\er^{-r^2t}<\varepsilon$$
        这已经符合$\er^{-tx^2}$在$t\to+\infty$时对$x\in[r,+\infty)$\textbf{一致收敛于0}的定义。

        第二部分：由于它与$x$无关，直接计算可知对任何$m>1$有
        $$\forall x\in[r,+\infty),\quad\bigg|\int_0^m\cos t\dr x\bigg|=|1-\sin m|<2$$

        综合上述讨论，迪利克雷判别法可知原含参积分在$[r,+\infty)$一致收敛。  
        
        \item \textbf{不一致收敛证明}
        
        为了证明不一致收敛，我们需要写出\textbf{柯西收敛原理}的否定：
        $$\exists\varepsilon>0,\quad\forall M>0,\quad\exists b>a>M,\quad x\in(0,+\infty),\quad\bigg|\int_a^b\er^{-tx^2}\cos t\dr t\bigg|\ge\varepsilon$$
        由于我们可以让$\er^{-tx^2}$任意接近1，而$\cos t$在一个周期中非负部分的积分为2，可以取任何$\varepsilon<2$，如取$\varepsilon=1$。

        此时，对任何$M>0$，取正整数$k$使得$a=2k\pi-\frac{\pi}{2}>M$、$b=2k\pi+\frac{\pi}{2}$，此时区间中被积函数为正，绝对值为
        $$\bigg|\int_{2k\pi-\pi/2}^{2k\pi+\pi/2}\er^{-x^2t}\cos t\dr t\bigg|=\int_{2k\pi-\pi/2}^{2k\pi+\pi/2}\er^{-x^2t}\cos t\dr t$$
        取$x=\sqrt{\frac{\ln2}{2k\pi+\pi/2}}$，可发现它为正数，且$\er^{-x^2(2k\pi+\pi/2)}=\frac{1}{2}$，此时有
        $$\int_{2k\pi-\pi/2}^{2k\pi+\pi/2}\er^{-x^2t}\cos t\dr t\ge\int_{2k\pi-\pi/2}^{2k\pi+\pi/2}\er^{-x^2(2k\pi+\pi/2)}\cos t\dr t=\frac{1}{2}\int_{2k\pi-\pi/2}^{2k\pi+\pi/2}\cos t\dr t=1$$
        这已经符合柯西收敛原理的否定。
    \end{itemize}
}
\end{framed}

\note 务必分清一致收敛证明中的\textbf{被积变量}与\textbf{参数}。

\note 不一致收敛的证明仍然往往依赖\textbf{柯西收敛原理}直接构造，思路与函数项级数时类似，可以参考对应的证明。

\begin{framed}
\begin{exercise}
    计算
    $$I(t)=\int_0^{+\infty}\er^{-x^2}\cos(tx)\dr x$$
    可直接使用$I(0)=\frac{\sqrt\pi}{2}$。
\end{exercise}
\sol{
    分为如下步骤：
    \begin{itemize}
        \item \textbf{求导积分交换性}
        
        被积函数对$t$求偏导可得$-x\er^{-x^2}\sin(tx)$，下面证明
        $$\forall t\in\rb,\quad I'(t)=\int_0^{+\infty}(-x\er^{-x^2}\sin(tx))\dr x$$

        由于
        $$|-x\er^{-x^2}\sin(tx)|\le x\er^{-x^2}$$
        而
        $$\int_0^{+\infty}x\er^{-x^2}\dr x=\frac{1}{2}\int_0^{+\infty}\er^{-x^2}\dr x^2=\frac{1}{2}(-\er^{-x^2})\big|^{+\infty}_0=\frac{1}{2}$$
        收敛，利用\textbf{$M$判别法}可知$\int_0^{+\infty}(-x\er^{-x^2}\sin(tx))\dr x$在$t\in\rb$\textbf{一致收敛}。

        由于
        $$|\er^{-x^2}\cos(tx)|\le\er^{-x^2}$$
        而
        $$\int_0^{+\infty}\er^{-x^2}\dr x=\frac{\sqrt\pi}{2}$$
        收敛，利用\textbf{$M$判别法}可知$\int_0^{+\infty}\er^{-x^2}\cos(tx)\dr x$在$t\in\rb$\textbf{一致收敛}，从而\textbf{逐点收敛}。

        利用含参积分一致收敛的性质即得$I'(t)$的表达式成立。

        \item \textbf{结果计算}

        为了能够计算$I'(t)$，利用\textbf{分部积分}可得(直接计算知$\dr\er^{-x^2}=-2x\er^{-x^2}$)
        $$I'(t)=\frac{1}{2}\int_0^{+\infty}\sin(tx)\dr\er^{-x^2}=\frac{1}{2}\bigg(\sin(tx)\er^{-x^2}\big|^{+\infty}_0-\int_0^{+\infty}t\cos(xt)\er^{-x^2}\dr x\bigg)$$
        由于$|\sin(tx)\er^{-x^2}|\le\er^{-x^2}$，利用\textbf{夹逼定理}可知它在无穷处的极限为0，由此
        $$I'(t)=-\frac{1}{2}\int_0^{+\infty}t\cos(xt)\er^{-x^2}\dr x$$
        对比形式即可发现
        $$I'(t)=-\frac{t}{2}I(t)$$
        由于$I(0)=\frac{\sqrt\pi}{2}>0$，在$I(t)>0$时将它写为
        $$\frac{\dr I'(t)}{I(t)}=-\frac{t}{2}\dr t$$
        两侧同积分得到($C$为任意常数)
        $$\ln I(t)=-\frac{1}{4}t^2+C$$
        也即
        $$I(t)=\er^C\er^{-\frac{1}{4}t^2}$$
        再利用$I(0)$的值最终得到
        $$I(t)=\frac{\sqrt\pi}{2}\er^{-\frac{1}{4}t^2}$$

        \note 注意解出的$I(t)$保证了在$\rb$上恒大于0，从而符合化简中用到的条件。
    \end{itemize}
}
\end{framed}

\begin{framed}
\begin{exercise}\label{inner deri}
    计算
    $$I(a)=\int_0^{\pi/2}\ln(a^2\sin^2x+\cos^2x)\dr x$$
\end{exercise}
\sol{
    分为如下步骤：
    \begin{itemize}
        \item \textbf{$I(0)$收敛性}
        
        可以发现，$a>0$时在$[0,\frac{\pi}{2}]$有
        $$a^2\sin^2x+\cos^2x\ge\min\{a^2,1\}>0$$
        原积分为\textbf{含参正常积分}。而$a=0$时，被积函数$\ln\cos^2x$以$\frac{\pi}{2}$为\textbf{瑕点}。

        我们先证明$a=0$时积分
        $$\int_0^{\pi/2}\ln\cos^2x\dr x$$
        收敛，以方便之后的推导。

        由于被积函数在$(0,\frac{\pi}{2})$恒小于0，我们利用广义积分的基本性质将原本的收敛性问题转化为\textbf{非负函数}瑕积分
        $$\int_0^{\pi/2}-\ln\cos^2x\dr x=\int_0^{\pi/2}-2\ln\cos x\dr x$$
        的收敛性问题。

        下面对被积函数在$\frac{\pi}{2}$处进行估算。换元$t=\frac{\pi}{2}-x$可发现被积函数为
        $$-2\ln\sin t$$
        而$x\to(\frac{\pi}{2})^-$可转化为$t\to0^+$。由于$t\to0^+$时$\sin t\sim t$，利用极限\textbf{保序性}可发现存在$\delta>0$使得$t\in(0,\delta)$时有$\frac{\sin t}{t}>\frac{1}{2}$，也即$$\sin t>\frac{t}{2}$$
        由此在$(0,\delta)$可放缩
        $$0<-2\ln\sin t<-2\ln\frac{t}{2}=-2\ln t+2\ln2$$
        
        利用$\ln t$的阶估算性质(可见\exref{xlnx})，任取$0<c<1$可计算得
        $$\lim_{t\to0^+}\frac{-2\ln t+2\ln 2}{\frac{1}{t^c}}=\lim_{t\to0^+}(-2t^c\ln t+2t^c\ln 2)=0$$
        从而在不等式每个部分同除以$\frac{1}{t^c}$，由\textbf{夹逼定理}有
        $$\lim_{t\to0^+}\frac{-2\ln\sin t}{\frac{1}{t^c}}=0$$
        代换回$x$得到
        $$\lim_{x\to(\frac{\pi}{2})^-}\frac{-\ln\cos^2x}{\frac{1}{(\frac{\pi}{2}-x)^c}}=0$$
        由\textbf{幂函数瑕积分}敛散性结论可知$\int_0^{\pi/2}\frac{1}{(\frac{\pi}{2}-x)^c}\dr x$收敛，于是根据\textbf{比较判别法}得原瑕积分在$\frac{\pi}{2}$处也收敛，从而收敛。

        \item \textbf{$I(a)$连续性}
        
        由定义$I(a)$是一个\textbf{偶函数}，由此只需证明$I(a)$在$[0,+\infty)$的连续性。
        
        当$a>0$时，利用\textbf{初等函数连续性}可知$\ln(a^2\sin^2x+\cos^2x)$是$a\in(0,+\infty)$、$x\in[0,\frac{\pi}{2}]$上的连续函数，从而由\textbf{含参正常积分}结论可知$I(a)$在$a\in(0,+\infty)$连续。

        \note 教材11.2节的定理只给出了$a$、$x$都在闭区间中的情况，对任何$a_0\in(0,+\infty)$，考虑闭区间$[a_0-\varepsilon,a_0+\varepsilon]\subset(0,+\infty)$并应用教材定理即可得到在$a_0$点的连续性。

        为了证明$a=0$处的连续性，考虑闭区间$[-1,1]$。当$a\in[-1,1]$、$x\in[0,\frac{\pi}{2}]$时恒有$a^2\sin^2x+\cos^2x\le1$，从而有
        $$|\ln(a^2\sin^2x+\cos^2x)|\le|\ln\cos^2x|=-\ln\cos^2x$$
        由于第一部分已经证明了
        $$\int_0^{\pi/2}-\ln\cos^2x\dr x$$
        收敛，从而通过\textbf{$M$判别法}可知$I(a)$在$[-1,1]$上\textbf{一致收敛}(这里的一致收敛性为将$\frac{\pi}{2}$视为瑕点的一致收敛，虽然$a\ne0$时其不为瑕积分，但可完全类似定义并具有类似性质)，这就用含参积分一致收敛的性质证明了$I(a)$在$[-1,1]$连续，将它与之前的结果结合可得$I(a)$在$\rb$上连续。

        \item \textbf{求导积分交换性}
        
        被积函数对$a$偏导可得$\frac{2a\sin^2x}{a^2\sin^2x+\cos^2x}$，我们下面证明，对任何$a\in(0,+\infty)$有
        $$I'(a)=\int_0^{\pi/2}\frac{2a\sin^2x}{a^2\sin^2x+\cos^2x}\dr x$$
        对任何$a_0\in(0,+\infty)$，取闭区间$[a_0-\varepsilon,a_0+\varepsilon]\subset(0,+\infty)$，与上一部分完全相同可知闭区间中
        $$\int_0^{\pi/2}\frac{2a\sin^2x}{a^2\sin^2x+\cos^2x}\dr x$$
        是含参正常积分，且利用\textbf{初等函数连续性}结论可知被积函数是$a\in[a_0-\varepsilon,a_0+\varepsilon]$、$x\in[0,\frac{\pi}{2}]$上的连续函数。
        
        由此，根据\textbf{含参正常积分}结论可知$a\in[a_0-\varepsilon,a_0+\varepsilon]$时
        $$I'(a)=\int_0^{\pi/2}\frac{2a\sin^2x}{a^2\sin^2x+\cos^2x}\dr x$$
        由于$a_0$可取任何正数，我们证明了求导与积分可交换。
            
        \item \textbf{结果计算}
            
        我们先计算$I'(a)$。将被积函数改写为
        $$I'(a)=\int_0^{\pi/2}\frac{2a\tan^2x}{a^2\tan^2x+1}\dr x$$
        换元$t=\tan x$得到(这里过程中会出现\textbf{广义积分}，但由于换元过程保证了单调性，不会出现收敛性问题)
        $$I'(a)=\int_0^{+\infty}\frac{2at^2}{a^2t^2+1}\frac{\dr t}{t^2+1}=\int_0^{+\infty}\frac{2at^2\dr t}{(a^2t^2+1)(t^2+1)}$$
        在$a>0$的基础上，我们进一步假设$a\ne1$，将被积函数\textbf{裂项}得到
        $$\frac{2at^2}{(a^2t^2+1)(t^2+1)}=\frac{2a}{a^2-1}\bigg(\frac{1}{t^2+1}-\frac{1}{a^2t^2+1}\bigg)$$
        直接计算可知$a>0$时
        $$\int_0^{+\infty}\frac{\dr t}{t^2+1}=\arctan t\big|^{+\infty}_0=\frac{\pi}{2}$$
        $$\int_0^{+\infty}\frac{\dr t}{a^2t^2+1}=\frac{1}{a}\int_0^{+\infty}\frac{\dr(at)}{a^2t^2+1}=\frac{1}{a}\arctan(at)\big|^{+\infty}_0=\frac{\pi}{2a}$$

        从而$a>0$且$a\ne1$时
        $$I'(a)=\frac{2a}{a^2-1}\bigg(\frac{\pi}{2}-\frac{\pi}{2a}\bigg)=\frac{\pi}{a+1}$$
        另一方面，直接代入可知$I(1)=0$，利用$I$在$\rb$上的\textbf{连续性}，根据微积分基本定理可知
        $$\forall a\in[0,+\infty),\quad I(a)=\int_1^a\frac{\pi}{b+1}\dr b=\pi\ln(a+1)-\pi\ln2=\pi\ln\frac{a+1}{2}$$
        由于$I(a)$为偶函数，最终得到
        $$\forall a\in\rb,\quad I(a)=\pi\frac{|a|+1}{2}$$
    \end{itemize}
}
\end{framed}

\note 对一般点处的瑕积分收敛性判断往往会先\textbf{移至0处}以方便应用等价无穷小等结论。

\note 与函数项级数类似，如果无法说明整体的一致收敛性，可以通过\textbf{内闭一致收敛}(内部闭区间上的一致收敛)证明导函数的性质。

\note 本题展示了\textbf{含参正常积分}与\textbf{含参广义积分}的性质结合以确定$I(a)$在$\rb$上的性质。

\begin{framed}
\begin{exercise}
    计算
    $$\int_0^{+\infty}\frac{\sin(ax)-\sin(bx)}{x}\er^{-x}\dr x$$
    这里参数$a>0$、$b>0$。
\end{exercise}
\sol{
    分为如下步骤：
    \begin{itemize}
        \item \textbf{构造含参积分}
        
        我们考虑
        $$I(c)=\int_0^{+\infty}\frac{\sin(ax)-\sin(bx)}{x}\er^{-cx}\dr x$$
        利用\textbf{乘除极限}结论与\textbf{洛必达法则}可算出
        $$\lim_{x\to0}\frac{\sin(ax)-\sin(bx)}{x}\er^{-cx}=\lim_{x\to0}\frac{a\cos(ax)-b\cos(bx)}{1}\cdot1=a-b$$
        由此，补充定义0处为$a-b$即连续，这点\textbf{不为瑕点}。

        我们先证明$I(c)$在$(0,+\infty)$的存在性。由三角函数的性质可知$x\in(0,+\infty)$时
        $$\bigg|\frac{\sin(ax)-\sin(bx)}{x}\bigg|\le\frac{|\sin(ax)|}{x}+\frac{|\sin(bx)|}{x}\le a+b$$
        于是有
        $$\bigg|\frac{\sin(ax)-\sin(bx)}{x}\er^{-cx}\bigg|\le(a+b)\er^{-cx}$$
        而
        $$\int_0^{+\infty}(a+b)\er^{-cx}\dr x=\bigg(-\frac{a+b}{c}\er^{-cx}\bigg)\bigg|^{+\infty}_0=\frac{a+b}{c}$$
        收敛，利用\textbf{比较判别法}可知$I(c)$对应的广义积分\textbf{绝对收敛}，从而收敛。

        上述过程也证明了
        $$|I(c)|\le\frac{a+b}{c}$$
        于是由\textbf{夹逼定理}可知
        $$\lim_{c\to+\infty}I(c)=0$$
        为了计算题中要求的$I(1)$，我们接下来考虑区间$[1,+\infty)$。

        \note 不考虑$[0,1]$是由于$[0,1]$上的\textbf{一致收敛性}难以证明。

        \item \textbf{求导积分交换性}
        
        被积函数对$c$求偏导可得$(\sin(bx)-\sin(ax))\er^{-cx}$，下面证明
        $$\forall c\in[1,+\infty),\quad I'(c)=\int_0^{+\infty}(\sin(bx)-\sin(ax))\er^{-cx}\dr x$$

        由于$c\in[1,+\infty)$时
        $$|(\sin(bx)-\sin(ax))\er^{-cx}|\le2\er^{-x}$$
        利用\textbf{指数函数无穷积分}收敛性结论可知右侧收敛，从而利用\textbf{$M$判别法}得$\int_0^{+\infty}(\sin(bx)-\sin(ax))\er^{-cx}\dr x$在$c\in[1,+\infty)$\textbf{一致收敛}。

        上一部分已经证明$I(c)$在$c\in(0,+\infty)$\textbf{逐点收敛}，利用含参积分一致收敛的性质即得$I'(t)$的表达式成立。
            
        \item \textbf{结果计算}
        
        我们先计算$I'(c)$。直接\textbf{分部积分}可得对$c\in[1,+\infty)$时
        $$\begin{aligned}\int_0^{+\infty}\sin(bx)\er^{-cx}\dr x&=-\frac{1}{c}\int_0^{+\infty}\sin(bx)\dr\er^{-cx}\\ &=-\frac{1}{c}\bigg(\sin(bx)\er^{-cx}\big|^{+\infty}_0-\int_0^{+\infty}b\cos(bx)\er^{-cx}\dr x\bigg)\end{aligned}$$
        由于$|\sin(bx)\er^{-cx}|\le\er^{-cx}$，利用\textbf{夹逼定理}可知它在无穷处的极限为0，由此
        $$\int_0^{+\infty}\sin(bx)\er^{-cx}\dr x=\frac{b}{c}\int_0^{+\infty}\cos(bx)\er^{-cx}\dr x$$
        再次\textbf{分部积分}可得
        $$\int_0^{+\infty}\cos(bx)\er^{-cx}\dr x=-\frac{1}{c}\bigg(\cos(bx)\er^{-cx}\big|^{+\infty}_0+\int_0^{+\infty}b\sin(bx)\er^{-cx}\dr x\bigg)$$
        化简也即
        $$\int_0^{+\infty}\cos(bx)\er^{-cx}\dr x=\frac{1}{c}-\frac{b}{c}\int_0^{+\infty}\sin(bx)\er^{-cx}\dr x$$
        将此式代入上一个分部积分得结果可以发现
        $$\int_0^{+\infty}\sin(bx)\er^{-cx}\dr x=\frac{b}{c^2}-\frac{b^2}{c^2}\int_0^{+\infty}\sin(bx)\er^{-cx}\dr x$$
        由此可得
        $$\int_0^{+\infty}\sin(bx)\er^{-cx}\dr x=\frac{b}{b^2+c^2}$$
        同理可得
        $$\int_0^{+\infty}\sin(ax)\er^{-cx}\dr x=\frac{a}{a^2+c^2}$$
        因此
        $$\forall c\in[1,+\infty),\quad I'(c)=\frac{b}{b^2+c^2}-\frac{a}{a^2+c^2}$$
        利用\textbf{广义积分的微积分基本定理}可知
        $$\lim_{c\to+\infty}I(c)-I(1)=\int_1^{+\infty}\bigg(\frac{b}{b^2+c^2}-\frac{a}{a^2+c^2}\bigg)\dr c$$
        利用$\frac{1}{a^2+x^2}$对$x$的不定积分公式(可见上册教材3.2节)可直接得到右侧被积函数的一个原函数
        $$\arctan\frac{c}{b}-\arctan\frac{c}{a}$$
        再由第一部分的讨论$I(c)$在$+\infty$处的极限为0，综合得到
        $$0-I(1)=\bigg(\frac{\pi}{2}-\frac{\pi}{2}\bigg)-(\arctan\frac{1}{b}-\arctan\frac{1}{a})$$
        也即
        $$I(1)=\arctan\frac{1}{b}-\arctan\frac{1}{a}$$
        这就是所求的结果。

        \note 结果事实上可以利用\textbf{反三角函数性质}化简为$\arctan a-\arctan b$。
    \end{itemize}
}
\end{framed}

\note 构造含参积分并对参数求导往往需要以\textbf{求导之后易于计算}且\textbf{保证收敛性}为标准，具体技巧相对灵活。

\note 注意学完广义积分后，微积分基本定理可以对\textbf{无穷处}使用，因此计算特殊点的值时可以考虑将导函数积分到无穷处。

\begin{framed}
\begin{exercise}
    计算
    $$\int_0^{+\infty}\frac{1}{(t+x^2)^n}\dr x$$
    其中参数$n$为正整数、$t$为正数。
\end{exercise}
\sol{
    分为如下步骤：
    \begin{itemize}
        \item \textbf{收敛性分析}
        
        设本题要计算的含参广义积分为$I(n,t)$，我们先证明它的收敛性。可以发现，被积函数在$[0,+\infty)$连续，从而只需要讨论$+\infty$处。
        
        利用广义积分的基本性质，\textbf{去除正常积分不影响广义积分敛散性}，因此$I(n,t)$敛散性与
        $$\int_1^{+\infty}\frac{1}{(t+x^2)^n}\dr x$$
        相同。由于被积函数在$[1,+\infty)$恒正，这是一个\textbf{非负函数}无穷积分，直接估算可知对正整数$n$有
        $$\forall x\in[1,+\infty),\quad\frac{1}{(t+x^2)^n}\le\frac{1}{t+x^2}<\frac{1}{x^2}$$
        由\textbf{幂函数无穷积分}敛散性结论可知$\int_1^{+\infty}\frac{1}{x^2}\dr x$收敛，于是根据\textbf{比较判别法}得原广义积分在$+\infty$处也收敛，从而收敛。

        \item \textbf{递推计算}
        
        为了构建$n$之间的递推，我们尝试直接进行\textbf{分部积分}，此时可得
        $$\int_0^{+\infty}\frac{1}{(t+x^2)^n}\dr x=\frac{x}{(t+x^2)^n}\bigg|^{+\infty}_0-\int_0^{+\infty}x\dr\frac{1}{(t+x^2)^n}$$
        由分母的次数高于分子可知$\frac{x}{(t+x^2)^n}$在$+\infty$处的极限为0，因此计算导数可知
        $$\int_0^{+\infty}\frac{1}{(t+x^2)^n}\dr x=n\int_0^{+\infty}\frac{2x^2}{(t+x^2)^{n+1}}\dr x$$
        由于已经证明了$I(n,t)$均收敛，我们可以拆分被积函数得到
        $$\int_0^{+\infty}\frac{2x^2}{(t+x^2)^{n+1}}\dr x=\int_0^{+\infty}\frac{2(t+x^2)-2t}{(t+x^2)^{n+1}}\dr x=2\int_0^{+\infty}\frac{\dr x}{(t+x^2)^n}-2t\int_0^{+\infty}\frac{\dr x}{(t+x^2)^{n+1}}$$
        也即我们证明了对正整数$n$与$t>0$有
        $$I(n,t)=n(2I(n,t)-2tI(n+1,t))$$
        化简得到
        $$I(n+1,t)=\frac{2n-1}{2nt}I(n,t)$$
        直接计算可知
        $$I(1,t)=\int_0^{+\infty}\frac{1}{t+x^2}\dr x=\bigg(\frac{1}{\sqrt{t}}\arctan\frac{x}{\sqrt{t}}\bigg)\bigg|^{+\infty}_0=\frac{\pi}{2\sqrt{t}}$$
        由此最终得到
        $$I(n,t)=\frac{\pi}{2t^{n-1/2}}\cdot\frac{1\cdot3\cdots(2n-3)}{2\cdot4\cdots(2n-2)}$$
        也可写成
        $$I(n,t)=\frac{\pi}{2t^{n-1/2}}\cdot\frac{(2n-3)!!}{(2n-2)!!}$$
    \end{itemize}

    \note 本题对$t$求导也可以从$I(n,t)$计算出$I(n+1,t)$，但需要\textbf{更复杂的一致收敛性分析}，并没有上述方法简便。

    \note 本题也可以换元$x=\sqrt{t}\tan p$\ (换元成对$p$的积分)后利用\textbf{$B$函数性质}得到结论。
}
\end{framed}

\note 注意含参的积分未必一定要对参数求导或积分，也可能\textbf{直接递推}得到结果。只有直接计算难以进行时再尝试引入参数。

\begin{framed}
\begin{exercise}
    计算
    $$\lim_{n\to\infty}\int_0^{+\infty}\er^{-x^n}\dr x$$
\end{exercise}
\sol{
    换元$t=x^n$可知
    $$\int_0^{+\infty}\er^{-x^n}\dr x=\int_0^{+\infty}\er^{-t}\dr\sqrt[n]{t}=\frac{1}{n}\int_0^{+\infty}t^{1/n-1}\er^{-t}\dr t$$
    利用$\Gamma$函数的定义，这即为
    $$\frac{1}{n}\Gamma\bigg(\frac{1}{n}\bigg)$$

    \note 注意换元后计算出显式表达式已经证明了原积分的\textbf{收敛性}。

    由$\Gamma$函数的性质，它可以进一步化为$\Gamma(1+\frac{1}{n})$，从而利用$\Gamma$函数的\textbf{连续性}与\textbf{归结原理}即得所求的极限为
    $$\lim_{n\to\infty}\Gamma\bigg(1+\frac{1}{n}\bigg)=\lim_{x\to0}\Gamma(1+x)=\Gamma(1)=1$$
}
\end{framed}

\note 对于$\er^{-x^n}$相关的问题，往往可以换元$t=x^n$化为$\Gamma$函数或类似形式处理。本题若是想不到这样的算法，\textbf{难以直接估算出结果}，因为结果难以猜到。

\begin{framed}
\begin{exercise}
    计算
    $$\int_a^b\frac{x}{t^2+x^2}\dr t$$
    这里参数$b>a\ge0$、$x>0$。以此计算
    $$\lim_{x\to+\infty}\suminf\frac{x}{n^2+x^2}$$
\end{exercise}
\sol{
    分为如下步骤：
    \begin{itemize}
        \item \textbf{积分计算}
        
        利用$\frac{1}{a^2+t^2}$对$t$的不定积分公式(可见上册教材3.2节)可知被积函数的一个原函数为
        $$\arctan\frac{t}{x}$$
        由此积分结果为
        $$\arctan\frac{b}{x}-\arctan\frac{a}{x}$$

        \item \textbf{收敛性证明}
        
        为了计算极限，我们先证明对任何$x>0$都有
        $$\suminf\frac{x}{n^2+x^2}$$
        收敛。由于这是一个正项级数，直接放缩可知
        $$\suminf\frac{x}{n^2+x^2}<\suminf\frac{x}{n^2}$$
        由\textbf{$p$级数}敛散性结论与数项级数基本性质可知$\suminf\frac{x}{n^2}$收敛，于是根据\textbf{比较判别法}得左侧也收敛。

        \item \textbf{极限计算}
        
        为了利用之前计算的\textbf{积分}结果，可以将$\frac{x}{n^2+x^2}$在\textbf{两个方向}放缩为积分，具体来说，对正整数$n$与$x>0$，由于$\frac{x}{t^2+x^2}$随$t$增加\textbf{单调递减}，必然有
        $$\int_n^{n+1}\frac{x}{t^2+x^2}\dr t\le\frac{x}{n^2+x^2}\le\int_{n-1}^n\frac{x}{t^2+x^2}\dr t$$
        对$n$从1到$M$求和有
        $$\int_1^{M+1}\frac{x}{t^2+x^2}\dr t\le\sum_{n=1}^M\frac{x}{n^2+x^2}\le\int_0^M\frac{x}{t^2+x^2}\dr t$$
        利用前一部分计算结果也即
        $$\arctan\frac{M+1}{x}-\arctan\frac{1}{x}\le\sum_{n=1}^M\frac{x}{n^2+x^2}\le\arctan\frac{M}{x}$$
        令$M\to+\infty$，由于已知中间部分极限存在，由极限\textbf{保序性}与$\arctan$的极限性质可知
        $$\frac{\pi}{2}-\arctan\frac{1}{x}\le\suminf\frac{x}{n^2+x^2}\le\frac{\pi}{2}$$
        换元$t=\frac{1}{x}$，由初等函数\textbf{连续性}直接计算可知
        $$\lim_{x\to+\infty}\frac{\pi}{2}-\arctan\frac{1}{x}=\lim_{t\to0}\frac{\pi}{2}-\arctan t=\frac{\pi}{2}$$
        由此利用\textbf{夹逼定理}最终得到
        $$\lim_{x\to+\infty}\suminf\frac{x}{n^2+x^2}=\frac{\pi}{2}$$
    \end{itemize}
}
\end{framed}

\note 注意本题的求和与极限\textbf{不可交换}。学完下半学期后，大家必须熟悉求和、极限、积分、导数\textbf{何时可以交换}相关的定理。

\note 学完含参积分后，我们可能面对\textbf{各种情况的积分与极限计算题}，这时需要\textbf{综合之前的学习的所有技巧}。本题的思路实际上来自\textbf{积分判别法的证明}，参数\textbf{仅在放缩中使用}。

\subsubsection{更多积分问题}

\begin{framed}
\begin{exercise}\label{func to series}
    给定满足
    $$a_0=0,\quad\lim_{n\to\infty}a_n=+\infty$$
    且严格单调增的数列$a_n$与$[0,+\infty)$上的连续函数$f(x)$，证明
    $$\int_0^{+\infty}f(x)\dr x$$
    收敛可以推出
    $$\suminf\int_{a_{n-1}}^{a_n}f(x)\dr x$$
    收敛到相同结果。

    证明其逆命题在$f(x)$非负时成立，否则可能不成立。
\end{exercise}
\sol{
    我们对$t\ge0$记
    $$F(t)=\int_0^tf(x)\dr x$$
    由$f(x)$在$[0,+\infty)$连续，利用\textbf{变上限积分}的性质可知$F(x)$也在$[0,+\infty)$连续。分为三个部分证明：
    \begin{itemize}
        \item \textbf{广义积分推级数}
        
        若广义积分收敛，设其结果为$A$，利用定义可知
        $$\lim_{t\to+\infty}F(t)=A$$
        另一方面，由定积分性质可得
        $$\sum_{n=1}^M\int_{a_{n-1}}^{a_n}f(x)\dr x=\int_{a_0}^{a_M}$$
        由$a_0=0$即得其为$F(a_M)$，由此，级数的极限可以写为
        $$\lim_{M\to\infty}F(a_M)$$
        由于$n\to\infty$时$a_n\to+\infty$，利用\textbf{推广的归结原理}可知
        $$\lim_{M\to\infty}F(a_M)=\lim_{t\to+\infty}F(t)=A$$
        这就证明了结论。
        
        \item \textbf{级数无法推出广义积分}
        
        考虑$a_n=2n\pi$、$f(x)=\sin x$，此时直接计算可知
        $$F(t)=\int_0^t\sin x\dr x=1-\cos t$$
        而
        $$F(a_M)=\int_0^{2M\pi}\sin x\dr x=0$$
        由此可知
        $$\lim_{M\to\infty}F(a_M)=0$$
        但
        $$\lim_{t\to+\infty}F(t)$$
        不存在，也即级数收敛但广义积分不收敛。
        
        \item \textbf{非负情况级数推广义积分}
        
        若$f(x)$非负，利用定积分性质可发现$F(x)$\textbf{单调递增}。我们设级数的极限
        $$\lim_{M\to\infty}F(a_M)=A$$
        利用定义，对任何$\varepsilon>0$，存在正整数$N$使得
        $$\forall n>N,\quad|F(a_n)-A|<\varepsilon$$
        由此，取$t=a_{N+1}$。对任何$x>t$，由于$n\to\infty$时$a_n\to+\infty$，必然存在$N_0$使得$x<a_{N_0}$，由此可知
        $$a_{N+1}<x<a_{N_0}$$
        再由$F(x)$的单调性可知
        $$F(a_{N+1})\le F(x)\le F(a_{N_0})$$
        由$a_n$单调递增，从$a_{N+1}<a_{N_0}$可看出$N_0>N+1>N$，由此必然有
        $$|F(a_{N+1})-A|<\varepsilon,\quad|F(a_{N_0})-A|<\varepsilon$$
        于是
        $$F(x)-A\le F(a_{N_0})-A<\varepsilon,\quad A-F(x)\ge A-F(a_{N+1})>-\varepsilon$$
        这就得到了
        $$|F(x)-A|<\varepsilon$$
        由于对任何$\varepsilon>0$都存在$t$使得$x>t$时上式成立，这已经符合极限
        $$\lim_{x\to+\infty}F(x)=A$$
        的定义。

        由此，我们证明了此时若级数的极限存在，广义积分也收敛到相同的结果。
        
        \note 结合第一部分结论可知$f(x)$非负时两者\textbf{等价}，收敛到相同结果或同时发散。
    \end{itemize}
}
\end{framed}

\note 本题是一个关键的\textbf{引理}，实际做题时$a_n$常取为$nT$，$T$为$f(x)$中涉及的周期函数的某个\textbf{周期}。

\note 本题的证明核心是利用单调性结论与\textbf{极限定义}估算。下半学期的各种收敛性问题中，由定义与极限语言直接操作仍然是一个基本的做题方式。

\begin{framed}
\begin{exercise}
    判断
    $$\int_0^{+\infty}\bigg(\suminf[n=0]\frac{(Ex)^n}{n!}\er^{-x}\bigg)\dr x,\quad\suminf[n=0]\int_0^{+\infty}\frac{(Ex)^n}{n!}\er^{-x}\dr x$$
    是否收敛，这里参数$E\in\rb$。
\end{exercise}
\sol{
    我们分别判断两式：
    \begin{itemize}
        \item \textbf{第一式}
        
        利用\textbf{基本初等函数的泰勒级数}结论可知
        $$\forall t\in\rb,\quad\er^t=\suminf[n=0]\frac{t^n}{n!}$$
        由此，代入$t=Ex$可算出
        $$\suminf[n=0]\frac{(Ex)^n}{n!}\er^{-x}=\er^{Ex}\cdot\er^{-x}=\er^{(E-1)x}$$
        由此，原积分也即
        $$\int_0^{+\infty}\er^{(E-1)x}\dr x$$
        利用\textbf{指数函数的无穷积分}敛散性结论，此积分收敛当且仅当$E-1<0$，也即$E<1$。

        \item \textbf{第二式}
        
        利用$\Gamma$函数定义可知
        $$\int_0^{+\infty}\frac{(Ex)^n}{n!}\er^{-x}\dr x=\frac{E^n}{n!}\int_0^{+\infty}x^n\er^{-x}\dr x=\frac{E^n}{n!}\Gamma(n+1)$$
        由正整数$\Gamma$的函数性质可知
        $$\frac{E^n}{n!}\Gamma(n+1)=\frac{E^n}{n!}\cdot n!=E^n$$
        由此，原级数也即
        $$\suminf[n=0]E^n$$
        利用\textbf{等比级数}敛散性结论，此级数收敛当且仅当$E\in(-1,1)$。
    \end{itemize}
}
\end{framed}

\note 注意所有的收敛判定框架中\textbf{直接计算}都是优先级最高的，需要\textbf{尽量计算化简}后再判断收敛性。

\note 进一步计算可发现本题交换积分、求和顺序后在收敛域交集的\textbf{收敛结果相同}，但\textbf{收敛域不同}。

\begin{framed}
\begin{exercise}
    设$f(x)$在$[0,+\infty)$连续，且$x\to+\infty$时的极限存在，记为$f(+\infty)$。设$0<a<b$，计算
    $$\int_0^{+\infty}\frac{f(ax)-f(bx)}{x}\dr x$$
\end{exercise}
\sol{
    分为如下步骤：
    \begin{itemize}
        \item \textbf{正常积分计算}
        
        \note 由于未知$\frac{f(ax)}{x}$、$\frac{f(bx)}{x}$的积分收敛，我们\textbf{无法将积分拆分为两部分}。

        由于此广义积分存在可能的瑕点0与无穷，而这两处的情况均难以确定，我们先计算它在\textbf{不含0的有限区间}的积分结果(设$\alpha<\beta$)
        $$F(\alpha,\beta)=\int_\alpha^\beta\frac{f(ax)-f(bx)}{x}\dr x$$
        此时，由于$\frac{f(ax)}{x}$、$\frac{f(bx)}{x}$均连续，它们均\textbf{可积}，积分可以写为
        $$\int_\alpha^\beta\frac{f(ax)}{x}\dr x-\int_\alpha^\beta\frac{f(bx)}{x}\dr x$$
        第一部分换元$t=ax$、第二部分换元$t=bx$，由$a$、$b$均为正可知
        $$F(\alpha,\beta)=\int_{a\alpha}^{a\beta}\frac{f(x)}{x}\dr x-\int_{b\alpha}^{b\beta}\frac{f(x)}{x}\dr x$$
        利用定积分性质，我们可按照$\alpha$、$\beta$\textbf{重新整理}为
        $$F(\alpha,\beta)=\int_{a\alpha}^{b\alpha}\frac{f(x)}{x}\dr x-\int_{a\beta}^{b\beta}\frac{f(x)}{x}\dr x$$

        \note 可记$G(t)=\int_1^t\frac{f(x)}{x}\dr x$，则上下的表达式分别为$(G(a\beta)-G(a\alpha))-(G(b\beta)-G(b\alpha))$、$(G(b\alpha)-G(a\alpha))-(G(b\beta)-G(a\beta))$，由此可看出相等。

        \item \textbf{极限计算}
        
        为了计算最终的广义积分，我们需要在$F(\alpha,\beta)$中令$\alpha\to0^+$、$\beta\to+\infty$。

        先考虑$\alpha$相关的部分
        $$\int_{a\alpha}^{b\alpha}\frac{f(x)}{x}\dr x$$
        由于未知$f(x)$的更多性质，我们可以考虑利用\textbf{积分中值定理}(由于$f(x)$在区间\textbf{连续}、$\frac{1}{x}$在区间\textbf{不变号})取出$\xi(\alpha)\in[a\alpha,b\alpha]$使得
        $$\int_{a\alpha}^{b\alpha}\frac{f(x)}{x}\dr x=f(\xi(\alpha))\int_{a\alpha}^{b\alpha}\frac{1}{x}\dr x=f(\xi(\alpha))\ln\frac{b}{a}$$
        利用\textbf{夹逼定理}有$\alpha\to0^+$时$\xi(\alpha)\to0^+$，再通过\textbf{复合函数极限定理}(由条件$f$在0处\textbf{右连续})即得
        $$\lim_{\alpha\to0^+}\int_{a\alpha}^{b\alpha}\frac{f(x)}{x}\dr x=\lim_{\alpha\to0^+}f(\xi(\alpha))\ln\frac{b}{a}=\lim_{\xi\to0^+}f(\xi)\ln\frac{b}{a}=f(0)\ln\frac{b}{a}$$

        \note 这里复合函数极限定理本质上需要上学期讲义4.1.2中推广到极限数的版本，这样才能保证换元后$\xi\to0^+$而不是$\xi\to0$。

        同理，对$\beta$相关的部分，利用\textbf{积分中值定理}可取出$\xi(\beta)\in[a\beta,b\beta]$使得
        $$\int_{a\beta}^{b\beta}\frac{f(x)}{x}\dr x=f(\xi(\beta))\int_{a\beta}^{b\beta}\frac{1}{x}\dr x=f(\xi(\beta))\ln\frac{b}{a}$$

        因此类似通过\textbf{夹逼定理}与\textbf{推广的复合函数极限定理}得到
        $$\lim_{\beta\to+\infty}\int_{a\beta}^{b\beta}\frac{f(x)}{x}\dr x=\lim_{\beta\to+\infty}f(\xi(\beta))\ln\frac{b}{a}=\lim_{\xi\to+\infty}f(\xi)\ln\frac{b}{a}=f(+\infty)\ln\frac{b}{a}$$

        综合以上两部分，我们要求的极限最终为
        $$f(0)\ln\frac{b}{a}-f(+\infty)\ln\frac{b}{a}=(f(0)-f(+\infty))\ln\frac{b}{a}$$
    \end{itemize}
}
\end{framed}

\note \textbf{积分中值定理}是比放缩到边界\textbf{更一般的积分估算}方法。之前利用放缩到边界进行的估算事实上也都可以通过中值定理得到，反过来，本题由于未知$f$的单调性，只能使用积分中值定理。

\begin{framed}
\begin{exercise}
    设
    $$f(t)=\bigg(\int_0^t\er^{-x^2}\dr x\bigg)^2,\quad g(t)=\int_0^1\frac{\er^{-t^2(1+x^2)}}{1+x^2}\dr x$$
    证明$f(t)+g(t)$为常数，并计算
    $$\int_0^{+\infty}\er^{-x^2}\dr x$$
\end{exercise}
\sol{
    证明一个量为常数往往需要通过\textbf{求导}计算，由此分为如下步骤：
    \begin{itemize}
        \item \textbf{计算$f'(t)$}
        
        利用\textbf{复合函数求导公式}与\textbf{变上限积分的导数}性质可知
        $$f'(t)=2\int_0^t\er^{-x^2}\dr x\cdot\frac{\dr}{\dr t}\int_0^t\er^{-x^2}\dr x=2\er^{-t^2}\int_0^t\er^{-x^2}\dr x$$
        
        \item \textbf{计算$g'(t)$}

        可以发现$g(t)$为\textbf{含参正常积分}，对参数求导得到的积分为
        $$\int_0^1(-2t)\er^{-t^2(1+x^2)}\dr x$$
        由于被积函数在$x\in[0,1]$、$t\in\rb$连续，与\exref{inner deri}完全类似，研究$t$处导数时考虑闭区间$[t-1,t+1]$可得
        $$\forall t\in\rb,\quad g'(t)=\int_0^1(-2t)\er^{-t^2(1+x^2)}\dr x$$

        \item \textbf{常数证明与计算}
        
        利用上学期\textbf{微分中值定理}结论，为了证明$f(t)+g(t)$为常数，我们只需要证明$f'(t)+g'(t)=0$，也即$g'(t)=-f'(t)$

        在$g'(t)$的表达式中换元$y=tx$，可发现
        $$g'(t)=-2\int_0^1\er^{-t^2-(tx)^2}\dr(tx)=-2\int_0^t\er^{-t^2-y^2}\dr y=-2\er^{-t^2}\int_0^t\er^{-y^2}\dr y$$
        再将定积分积分变量从$y$\textbf{替换}为$x$即得到$g'(t)=-f'(t)$，从而得证其为常数。

        进一步地，取$t=0$可发现
        $$f(0)+g(0)=\int_0^1\frac{1}{1+x^2}\dr x=\arctan x\big|^1_0=\frac{\pi}{4}$$
        由此我们可以得到
        $$\forall t\in\rb,\quad f(t)+g(t)=\frac{\pi}{4}$$
        
        \item \textbf{极限计算}

        为了计算无穷积分，我们需要令$t\to+\infty$，直接估算有
        $$\forall x\in[0,1],\quad\er^{-t^2(1+x^2)}\le\er^{-t^2}$$
        又由$g(t)$被积函数非负可知
        $$0\le g(t)\le\int_0^1\frac{\er^{-t^2}}{1+x^2}\dr x=\frac{\pi}{4}\er^{-t^2}$$
        利用复合函数极限结论换元$s=t^2$可知
        $$\lim_{t\to+\infty}\er^{-t^2}=\lim_{s\to+\infty}\er^{-s}=0$$
        由\textbf{夹逼原理}即得
        $$\lim_{t\to+\infty}g(t)=0$$
        而
        $$\lim_{t\to+\infty}(f(t)+g(t))=\lim_{t\to+\infty}\frac{\pi}{4}=\frac{\pi}{4}$$
        因此有
        $$\lim_{t\to+\infty}f(t)=\frac{\pi}{4}$$

        \item \textbf{无穷积分计算}
        
        直接估算可知题中所求的无穷积分\textbf{收敛}(可见教材11.1节例5)，设
        $$\int_0^{+\infty}\er^{-x^2}\dr x=\lim_{t\to+\infty}\int_0^t\er^{-x^2}\dr x=A$$
        利用乘除极限结论可知
        $$\lim_{t\to+\infty}f(t)=\lim_{t\to+\infty}\int_0^t\er^{-x^2}\dr x\cdot\int_0^t\er^{-x^2}\dr x=A^2$$
        由此可知$A^2=\frac{\pi}{4}$。由于$A$是非负函数的无穷积分，它必然\textbf{非负}，因此最终得到
        $$A=\frac{\sqrt\pi}{2}$$
        这就是要计算的无穷积分结果。

        \note 注意必须先证明$\int_0^{+\infty}p(x)\dr x$的\textbf{收敛性}，才能从
        $$\lim_{t\to\infty}\bigg(\int_0^tp(x)\dr x\bigg)^2$$
        的极限中得到无穷积分的结果。
    \end{itemize}
}
\end{framed}

\note 结合含参正常积分、含参广义积分与变上限积分，我们可以计算一般的\textbf{积分的导数}相关问题。

\begin{framed}
\begin{exercise}
    证明
    $$\int_0^1\frac{1}{x^x}\dr x=\suminf\frac{1}{n^n}$$
\end{exercise}
\sol{
    分为如下步骤：
    \begin{itemize}
        \item \textbf{初步分析}
        
        本题左右的形式看上去缺乏联系，因此我们只能尝试先把左侧的被积函数化为\textbf{指数函数}
        $$\int_0^1\er^{-x\ln x}\dr x$$
        虽然被积函数在0处无定义，但利用\exref{xlnx}中的重要估算与\textbf{复合函数极限}结论(换元$t=x\ln x$)可知
        $$\lim_{x\to0^+}\er^{-x\ln x}=\lim_{t\to0}\er^{-t}=1$$
        由此补充定义0处为1即连续，此积分实际上是正常积分。

        为了凑出一个级数，利用\textbf{基本初等函数的泰勒级数}可知
        $$\forall t\in\rb,\quad\er^t=\suminf[n=0]\frac{t^n}{n!}$$
        代入$t=-x\ln x$，可以得到
        $$\forall x\in(0,1],\quad\er^{-x\ln x}=\suminf[n=0]\frac{(-x\ln x)^n}{n!}$$
        于是左侧积分为
        $$\int_0^1\suminf[n=0]\frac{(-x\ln x)^n}{n!}\dr x$$

        \item \textbf{积分计算}
        
        为了确认这样的展开是否有意义，我们先尝试计算(与前一部分讨论完全类似可得$I_n$是正常积分)
        $$I_n=\int_0^1\frac{(-x\ln x)^n}{n!}\dr x=\frac{(-1)^n}{n!}\int_0^1(x\ln x)^n\dr x$$
        直接计算可知$I_0=1$，当$n$为正整数时\textbf{分部积分}可得
        $$I_n=\frac{(-1)^n}{(n+1)!}\int_0^1\ln^nx\dr x^{n+1}=\frac{(-1)^n}{(n+1)!}\bigg((x^{n+1}\ln^nx)\big|^1_0-\int_0^1x^{n+1}\dr\ln^nx\bigg)$$
        利用\exref{xlnx}中的重要估算可知$x^{n+1}\ln^nx$在0处的极限为0，而由$n$为正整数可知$\ln^nx$在1处为0，由此第一项为0，对第二项计算导数可发现
        $$I_n=\frac{(-1)^{n-1}}{(n-1)!}\cdot\frac{1}{n+1}\int_0^1x^n\ln^{n-1}x\dr x$$
        当$n=1$时这已经足以算出$I_1=\frac{1}{4}$，否则再次分部积分可得
        $$I_n=\frac{(-1)^{n-2}}{(n-2)!}\cdot\frac{1}{(n+1)^2}\int_0^1x^n\ln^{n-2}x\dr x$$
        由此再重复$n-2$次，直到$\ln$的次数为0，最终得到
        $$I_n=\frac{1}{(n+1)^n}\int_0^1x^n\dr x=\frac{1}{(n+1)^{n+1}}$$

        \note 更严谨的写法是设$I_{m,n}=\int_0^1x^m\ln^nx\dr x$并\textbf{归纳}计算。

        此表达式对任何$n\in\mathbb{N}$成立，由此，我们可以将右侧级数写为
        $$\suminf[n=0]I_n=\suminf[n=0]\int_0^1\frac{(-x\ln x)^n}{n!}\dr x$$

        \item \textbf{交换次序}

        至此，我们可以发现要证明的等式左右两侧实际上是交换了积分与求和的次序。为了证明次序可交换，我们需要证明级数的\textbf{一致收敛}性。

        定义(这里$n$为正整数)
        $$f(x)=\begin{cases}1&x=0\\\er^{-x\ln x}&x\in(0,1]\end{cases},\quad f_0(x)=1,\quad f_n(x)=\begin{cases}0&x=0\\\frac{(-x\ln x)^n}{n!}&x\in(0,1]\end{cases}$$
        之前的讨论已经证明了(0处直接计算即可)
        $$\forall x\in[0,1],\quad f(x)=\suminf[n=0]f_n(x)$$
        且每个$f_n(x)$都是$[0,1]$上的连续函数。此时要证的式子化为
        $$\int_0^1f(x)\dr x=\suminf[n=0]\int_0^1f_n(x)\dr x$$
        利用一致收敛的性质，只需再证明$\suminf[n=0]f_n(x)$在$[0,1]$一致收敛于$f(x)$。

        直接求导可发现
        $$\frac{\dr}{\dr x}x\ln x=\ln x+1$$
        又由它在0处极限为0、1处为0，可发现它在$[0,1]$上\textbf{绝对值最大}在导数为0点$x=\er^{-1}$取到，因此
        $$\forall x\in(0,1],\quad |x\ln x|\le|\er^{-1}\ln\er^{-1}|=\frac{1}{\er}$$
        由此可发现($n$为正整数时利用上述估算，$n=0$时直接计算验证)
        $$|f_n(x)|\le\frac{1}{n!\er^n}\le\frac{1}{\er^n}$$
        根据\textbf{等比级数}敛散性结论可知右侧部分和收敛，从而利用\textbf{强级数判别法}得左侧部分和在$[0,1]$上一致收敛，得证。
    \end{itemize}
}
\end{framed}

\note 本题展示了\textbf{通过展开为级数计算积分}的技巧，此技巧对于幂级数也有一定的应用。

\begin{framed}
\begin{exercise}
    设对正整数$n$有
    $$f_1(x)=\int_0^x\sqrt{1+t^4}\dr x,\quad f_{n+1}(x)=\int_0^xf_n(x)\dr x$$
    证明
    $$\suminf f_n(x)$$
    在$[0,2]$上一致收敛，判断它在$[0,+\infty)$是否一致收敛。
\end{exercise}
\sol{
    由于$f_1(x)$为非负函数积分，它在$[0,+\infty)$非负，同理递推可得到所有$f_n(x)$在$[0,+\infty)$非负。我们分别判断两种情况：
    \begin{itemize}
        \item \textbf{有界区间一致收敛性}
        
        由于我们无法算出$f_n(x)$的显式表达式，尝试直接进行\textbf{放大}。

        在有界区间$[0,2]$，取$c=\sqrt{17}$，则恒有$\sqrt{1+t^2}\le c$，于是对$x\ge0$有
        $$f_1(x)\le\int_0^xc\dr t=cx$$
        $$f_2(x)\le\int_0^x(ct)\dr t=\frac{1}{2}cx^2$$
        以此归纳可得对任何正整数$n$有
        $$f_n(x)\le c\frac{x^n}{n!}$$
        由此再结合非负性可知
        $$\forall x\in[0,2],\quad|f_n(x)|\le c\frac{2^n}{n!}$$
        利用$\er^t$的\textbf{泰勒级数}结论代入$t=2$可知
        $$\suminf[n=0]\frac{2^n}{n!}=\er^2$$
        由此通过数项级数基本性质可知右侧部分和收敛，从而利用\textbf{强级数判别法}得左侧部分和在$[0,2]$上一致收敛，得证。

        \item \textbf{无界区间情况}
        
        对于无界区间，我们尝试进行\textbf{反向}的估算。由于$\sqrt{1+t^4}\ge1$，于是对$x\ge0$有
        $$f_1(x)\ge\int_0^x1\dr t=x$$
        $$f_2(x)\ge\int_0^xt\dr t=\frac{1}{2}x^2$$
        以此归纳可得对任何正整数$n$有
        $$f_n(x)\ge\frac{x^n}{n!}$$
        从而
        $$\lim_{x\to+\infty}f_n(x)=+\infty$$
        我们下面证明满足上述条件的$f_n(x)$在$(0,+\infty)$\textbf{不可能一致收敛}。

        考虑\textbf{柯西收敛原理的否定}(见\exref{not uni})，取$\varepsilon=1$，对任何正整数$N$取$n=N+1$、$m=n+1$，此时对应的求和
        $$\bigg|\sum_{k=n+1}^mf_k(x)\bigg|=|f_m(x)|$$
        利用极限的\textbf{保序性}，给定$m$后存在$x_0$使得$f_m(x_0)\ge1$，此时即有
        $$\bigg|\sum_{k=n+1}^mf_k(x_0)\bigg|=|f_m(x_0)|=f_m(x_0)\ge\varepsilon$$
        这已经符合柯西收敛原理的否定。
    \end{itemize}
}
\end{framed}

\note 无法得到显式表达式时的\textbf{估算}往往依赖更多\textbf{对阶的直觉}。

\begin{framed}
\begin{exercise}
    设正项级数
    $$\suminf a_n$$
    收敛，对于任意$x>0$，定义$L(x)$为数列$a_n$中大于$x$的项的个数，证明$L(x)$对任何$x>0$均可定义，并证明广义积分
    $$\int_0^{+\infty}L(x)\dr x=\suminf a_n$$
\end{exercise}
\sol{
    分为如下步骤：
    \begin{itemize}
        \item \textbf{$L(x)$存在性}
        
        我们先证明$L(x)$对任何$x>0$有意义。

        由于正项级数$\suminf a_n$收敛，利用数项级数基本性质有
        $$\lim_{n\to\infty}a_n=0$$
        由此利用极限定义可知存在$N$使得
        $$\forall n>N,\quad|a_n|<x$$
        由此大于$x$的项只能在前$N$项中，$L(x)$即为$a_1$到$a_N$中大于$x$的数的个数，必然存在。
        
        \item \textbf{$a_n$最大值存在性}
        
        我们现在证明，若$a_n$均为正且
        $$\lim_{n\to\infty}a_n=0$$
        则它必然存在\textbf{最大值}，也即存在正整数$m$使得
        $$\forall n\in\mathbb{N}^*,\quad a_m\ge a_n$$
        且取到最大值的元素个数有限。

        考虑\textbf{反证法}。若最大值不存在，对每个$a_s$，一定存在$t>s$使得$a_t>a_s$\ (否则$a_s$比之后的项都大，$a_1$到$a_s$中的最大值即为数列的最大值)，从$a_1$出发，可以如此构造一个\textbf{严格单调递增}的子列，又由于$a_1>0$，它不可能趋于0，与收敛数列\textbf{任何子列收敛到相同极限}矛盾。

        同样，若取到最大值的元素个数无限，存在一个恒为最大值的子列，其极限也为最大值，不为0\ (由于$a_1>0$，最大值不可能为0)，矛盾。

        \item \textbf{$L(x)$计算}
        
        利用上述性质，我们可以先找出数列$a_n$的最大值，设为$w_1$，并假设有$k_1$项取到了最大值。根据定义有
        $$\forall x\ge w_1,\quad L(x)=0$$
        接下来，我们把这$k_1$项从数列$a_n$中\textbf{删除}，由于删除有限项不影响极限，剩下的项构成的数列$b_n$仍然恒正且极限为0。由此，我们可以找出数列$b_n$的最大值，设为$w_2$，并假设有$k_2$项取到了最大值。

        由于$b_n$中已经删除了所有等于$w_1$的项，所有$b_n$中的项必然小于$w_1$，从而有$w_2<w_1$。

        在$x\in[w_2,w_1)$时，根据定义可发现只有取到$w_1$的项比$x$大，因此有
        $$\forall x\in[w_2,w_1),\quad L(x)=k_1$$
        不断重复此操作，我们可以得到一个\textbf{严格单调递减}的$w_n$与对应的$k_n$，满足$w_n$是数列$a_n$中删除等于$w_1,\dots,w_{n-1}$的项后剩下的数列中的\textbf{最大值}，且有$k_n$项取到。

        与之前类似讨论，可以发现对正整数$n$有
        $$\forall x\in[w_{n+1},w_n),\quad L(x)=k_1+\dots+k_n$$

        \item \textbf{$w_n$极限计算}

        由于$w_n$严格单调下降，且恒大于0，利用\textbf{单调有界定理}可知它必然存在极限。我们现在证明
        $$\lim_{n\to\infty}w_n=0$$
        若否，利用\textbf{极限保序性}可知必然存在$c>0$使得所有$w_n>c$\ (存在$N$使得$n>N$时$w_n>c$，再由其单调递减可知对所有项均成立)。但是，对每个$w_i$，根据定义可发现$a_n$中至少有一项取到了$w_i$，因此$a_n$中存在无穷多项大于$c$，与\textbf{极限定义}矛盾。

        \item \textbf{无穷处收敛性}
        
        由于$L(x)$在$x\ge w_1$时恒为0，利用定义可发现
        $$\int_{w_1}^{+\infty}L(x)\dr x=0$$
        由此$L(x)$在$+\infty$处积分收敛，我们可以直接将积分
        $$\int_0^{+\infty}L(x)\dr x$$
        改写为
        $$\int_0^{w_1}L(x)\dr x$$
        
        \item \textbf{瑕积分引理}
        
        在\exref{func to series}中，我们证明了对于非负函数，若将无穷积分\textbf{拆分为无穷级数}后收敛，则原无穷积分必然收敛到相同结果，反之亦然。

        对于瑕积分，由于$w_n$严格单调下降趋于0，我们可以通过\textbf{完全相同}的方法证明如下的级数
        $$\suminf\int_{w_{n+1}}^{w_n}L(x)\dr x$$
        与积分
        $$\int_0^{w_1}L(x)\dr x$$
        等价，两者或\textbf{同时不存在}，或同时存在且收敛到\textbf{相同结果}。由此，问题可以转化为研究
        $$\suminf\int_{w_{n+1}}^{w_n}L(x)\dr x$$
        利用之前的计算结果可知其为
        $$\suminf(w_n-w_{n+1})(k_1+\dots+k_n)$$
        也即我们要证明的式子最终化为
        $$\suminf(w_n-w_{n+1})(k_1+\dots+k_n)=\suminf a_n$$

        \item \textbf{最终计算-拆分}
        
        直接计算可知
        $$\begin{aligned}&\sum_{n=1}^m(w_n-w_{n+1})(k_1+\dots+k_n)\\=&k_1(w_1-w_2+\dots+w_m-w_{m+1})+k_2(w_2-w_3+\dots+w_m-w_{m+1})+\dots+k_m(w_m-w_{m+1})\\=&k_1(w_1-w_{m+1})+\dots+k_m(w_m-w_{m+1})\\=&\sum_{n=1}^mk_n(w_n-w_{m+1})\\=&\sum_{n=1}^mk_nw_n-\sum_{n=1}^mk_nw_{m+1}\end{aligned}$$
        我们接下来证明
        $$\lim_{m\to\infty}\sum_{n=1}^mk_nw_n=\suminf k_nw_n=\suminf a_n$$
        $$\lim_{m\to\infty}\sum_{n=1}^mk_nw_{m+1}=\lim_{m\to\infty}w_{m+1}\sum_{n=1}^mk_n=0$$
        综合两部分即可得到最终结论。

        \item \textbf{最终计算-第一部分}
        
        由于$\sum_{n=1}^mk_nw_n$中至少包含$m$个数列$a_n$中的项，且它是\textbf{从大向小}取出，必然有
        $$\sum_{n=1}^mk_nw_n\ge\sum_{n=1}^ma_n$$
        另一方面，由于$\sum_{n=1}^mk_nw_n$中取出的所有项都在数列$a_n$中，且每项最多被计入了一次，必然有(利用$a_n$均为正)
        $$\sum_{n=1}^mk_nw_n\le\suminf a_n$$
        由此，我们得到了估算
        $$\sum_{n=1}^ma_n\le\sum_{n=1}^mk_nw_n\le\suminf a_n$$
        由于左侧在$m\to\infty$时等于右侧，利用\textbf{夹逼定理}可知
        $$\lim_{m\to\infty}\sum_{n=1}^mk_nw_n=\suminf a_n$$

        \note 这其实本质上是将$a_n$的所有项进行了\textbf{重新排列}。

        \item \textbf{最终计算-第二部分}
        
        利用定义可发现$\sum_{n=1}^mk_n$是所有大于$w_{m+1}$的项的个数。若此极限非零，由定义可知存在$\varepsilon>0$与无穷多个$m$满足(由其为正可去除绝对值)
        $$w_{m+1}\sum_{n=1}^mk_n\ge\varepsilon$$
        将这无穷多个$m$排成子列$m_1,m_2,\dots$，并记$r_i=w_{m_i}$，利用数列极限的基本性质，$r_i$也单调递减趋于0。

        由此，对任何正整数$i$，数列$a_n$中大于$r_i$项的至少有$\frac{\varepsilon}{r_i}$个。利用此结论可估算求和：考虑所有大于$r_1$的项可知
        $$\suminf a_n\ge r_1\cdot\frac{\varepsilon}{r_1}=\varepsilon$$
        进一步考虑所有大于$r_2$的项，其中至少有$\frac{\varepsilon}{r_1}$项大于$r_1$，可知
        $$\suminf a_n\ge r_1\cdot\frac{\varepsilon}{r_1}+r_2\bigg(\frac{\varepsilon}{r_2}-\frac{\varepsilon}{r_1}\bigg)=(r_1-r_2)\frac{\varepsilon}{r_1}+\varepsilon$$
        以此类推，我们可以得到对正整数$m$，求和
        $$\varepsilon+\sum_{n=1}^m(r_n-r_{n+1})\frac{\varepsilon}{r_n}$$
        有上界。

        然而，由于$r_i$单调递减趋于0，可以用\textbf{积分}估算
        $$(r_n-r_{n+1})\frac{\varepsilon}{r_n}\ge\int_{r_{n+1}}^{r_n}\frac{\varepsilon}{x}\dr x$$
        由此可知
        $$\varepsilon+\sum_{n=1}^m(r_n-r_{n+1})\frac{\varepsilon}{r_n}\ge\varepsilon+\int_{r_{m+1}}^{r_1}\frac{\varepsilon}{x}\dr x$$
        在$m\to\infty$时趋于正无穷，不存在上界，矛盾。
    \end{itemize}

    \note 本题的另一个方法是构造函数列
    $$f_n(x)=\begin{cases}a_n-x&x\in[0,a_n]\\0&x>a_n\end{cases}$$
    则$f_n(x)$在$[0,+\infty)$连续，且由\textbf{强级数判别法}可知$\suminf f_n(x)$一致收敛。进一步计算可验证
    $$\suminf f_n(x)=\int_x^{+\infty}L(t)\dr t$$
    再令$x\to0^+$即可。此方法相对巧妙，但\textbf{很难直接想到构造}(本质上是定义$f_n(x)$为某种特征函数的积分)。

}
\end{framed}

\note \textbf{构造级数解决问题}可能可以避免一些复杂的估算。

\end{document}